\documentclass[a4paper,12pt]{article}

\usepackage[top=2.0cm,bottom=2.0cm,left=2.5cm,right=2.0cm,headheight=14.5pt]{geometry}
\usepackage{amsmath,amssymb,amsthm}
\usepackage{bm}
\usepackage{xspace}
\usepackage{booktabs}
\usepackage{longtable}
\usepackage[modulo,pagewise]{lineno}
\modulolinenumbers[5]
\usepackage{tabularx}
\usepackage{enumitem}
\usepackage{setspace}
\usepackage{fancyhdr}
% hyperref loaded last (LaTeX convention) so that footnote destinations
% from lineno's per-page insertion interleave cleanly with hyperref's
% Hfootnote.N anchors. Resolves "Hfootnote.N has been referenced but
% does not exist" pdfTeX warnings.
\usepackage{hyperref}

\newtheorem{definition}{Definition}
\newtheorem{remark}{Remark}
\frenchspacing
\emergencystretch=1.5em

\newcolumntype{Y}{>{\raggedright\arraybackslash}X}
\newcommand{\E}{\mathbb{E}}
\newcommand{\KL}{\mathrm{KL}}
\newcommand{\R}{\mathbb{R}}
\newcommand{\TE}{\mathrm{TE}}
\newcommand{\CLE}{\mathrm{CLE}}
% Canonical terminology macros (enforce consistency)
\newcommand{\cb}{convolution bundle\xspace}
\newcommand{\cbs}{convolution bundles\xspace}  % plural form
\newcommand{\cba}{convolution-bundle\xspace}  % adjectival form
\newcommand{\Cb}{Convolution bundle\xspace}   % sentence-initial
\newcommand{\Cba}{Convolution-bundle\xspace}  % sentence-initial adjectival
\newcommand{\yd}{yoke depth\xspace}
\newcommand{\yda}{yoke-depth\xspace}          % adjectival form
\newcommand{\Yd}{Yoke depth\xspace}           % sentence-initial
% Regime name macros (enforce consistency)
\newcommand{\TB}{Truth Beam\xspace}
\newcommand{\LI}{Limager\xspace}
\newcommand{\RT}{Reality Transform\xspace}
\newcommand{\RTa}{Reality-Transform\xspace}      % adjectival
\newcommand{\RK}{Reality Kernel\xspace}
\newcommand{\RKs}{Reality Kernels\xspace}

% IPOI formatting: 1.5 line spacing, page numbers top centre
\onehalfspacing
\pagestyle{fancy}
\fancyhf{}
\fancyhead[C]{\thepage}
\renewcommand{\headrulewidth}{0pt}

\title{Physical Markov-Channel Apparatus for Verification,\\
Perception, and Controllable Rendering}
\author{Cathal Ryan Hynes\\P.I.G.M.I.E. Ltd.}
\date{May 2026}

\begin{document}
\linenumbers
\thispagestyle{empty}
\nolinenumbers

% ======================================================================
% DEDICATION (separable — does not affect technical disclosure)
% ======================================================================

\vspace*{\fill}

\begin{center}
{\Large\bfseries Ring of Light: Metric Might}
\end{center}

\bigskip

For my parents:

\medskip

\textbf{Mary Alacoque Ryan}---the fast loop: architect of my earliest
reactor environments, designer of the toughest probes and loss
functions, eternal Eve and devoted Demeter whose corrections arrived
instantly and whose signal never dropped.

\medskip

\textbf{Martin Hynes}---the slow loop: provider of new modalities and
training data that persisted quietly beneath the fast dynamics,
enriching the \cb with signals the fast loop could not
generate on its own.

\medskip

\textbf{Everyone who left a trace in the bundle}---the Iris channel.

\medskip

\begin{center}
$\alpha \to 1$.
\end{center}

\vspace*{\fill}

\begin{center}
\small\emph{This page may be detached without affecting the technical
disclosure, claims, or filing status.}
\end{center}

\clearpage
\linenumbers
\thispagestyle{fancy}

% ======================================================================
% DESCRIPTION
% ======================================================================

\begin{center}
{\Large\bfseries Physical Markov-Channel Apparatus for Verification,
Perception, and Controllable Rendering\footnote{AI language models
assisted in drafting and formatting this disclosure; the named human
inventor reviewed, selected, and is responsible for the inventive
content.}}
\end{center}

\bigskip

% ======================================================================
\section{Technical Field}
% ======================================================================

The present disclosure relates to sensing, projection, optical
computing, and verification systems. More particularly, it relates to
projector--detector assemblies that implement a parameterised physical
channel, represented as a Markov kernel, and that can be operated under
multiple objective families including verification, perception, and
controllable rendering, with optional bidirectional dynamical coupling
between the device and its environment.

% ======================================================================
\section{Background Art}
% ======================================================================

The following discussion provides context for the present disclosure
and is not intended as an admission that any referenced work
constitutes citable prior art or common general knowledge against the
claims to be filed.

Projector--camera and structured-illumination systems have been
proposed for
metrology, projection mapping, active sensing, and computational
imaging. Such systems are typically configured for a single
objective such as reconstruction quality, visual effects, or
measurement throughput.

Separately, digital generative models enable synthesis of plausible
images and videos that can be used for benign creation or for spoofing.
As a result, there is a growing need for physically grounded mechanisms
that produce auditable records of what illumination was applied and what
responses were measured, that permit objective evaluation of whether a
recording is consistent with a physical device and scene, and that
reuse a common hardware front-end across verification, sensing, and
stylised rendering tasks.

\paragraph{Digital provenance and content attestation.}
The Coalition for Content Provenance and Authenticity (C2PA) defines
a signed-manifest structure for tamper-evident media provenance,
binding cryptographic signatures to content metadata (C2PA Technical
Specification, Version~2.1, 2024, \texttt{c2pa.org}).  Related device
attestation approaches use trusted computing components to certify
camera or sensor state; commercial implementations include chip-level
content credentials at capture (e.g., Qualcomm/Truepic
\emph{Snapdragon}-based signing) and camera-integrated signing (e.g.,
Leica M11-P with embedded C2PA support).  Complementary provenance
frameworks include Project Origin and the Starling Framework, which
focus on chain-of-custody and distributed verification respectively.  These systems provide digital commitment
(metadata integrity via key management), but the trust root is a
cryptographic key, not a physical interaction trace: provenance can
be stripped, re-signed, or fabricated if keys are compromised, and
the signed manifest does not bind to the empirically hard-to-reproduce physics of the
capture event itself.  Existing digital provenance systems identified in the prior art survey do not produce
committed records whose authenticity rests on the empirical
difficulty of emulating a physical light--matter interaction.

\paragraph{Optical PUFs and physical authentication.}
Physically unclonable functions (PUFs) exploit manufacturing
variability to produce device-specific fingerprints.  Silicon PUFs
based on delay-circuit variability were introduced by Gassend et al.\
(``Silicon physical random functions,'' \emph{ACM~CCS},
pp.~148--160, 2002), establishing the challenge--response pair
framework.  Pappu et al.\
demonstrated that laser illumination of a disordered scattering token
produces speckle patterns suitable for one-way authentication
(``Physical One-Way Functions,'' \emph{Science}, 2002).  Suh and
Devadas established the canonical two-phase enrolment/verification
protocol for PUF-based device authentication (``Physical unclonable
functions for device authentication and secret key generation,''
\emph{DAC}, pp.~9--14, 2007).  Subsequent
work extended optical PUFs to nonlinear electro-optic media (Hui et
al., ``Non-linear optical scattering PUF,''
\emph{Optics Express}~31(24), 40646--40657, 2023), to quantum-secure single-photon
authentication (Goorden et al., ``Quantum-secure authentication of a
physical unclonable key,'' \emph{Optica}~1, 421--424, 2014),
and to integrated SLM--PUF--camera
apparatus (Davis, Letz, Mosk, and Pinkse, EP3252740B1, Schott AG and
Universiteit Twente).
Published patent literature also describes combined sensor-and-display
devices
used to interrogate a PUF, co-locating emission and sensing in a
single assembly for challenge--response readout.
However, many reported optical PUF systems operate open-loop: a
challenge is presented,
the medium responds, and the response is captured.  To the best of the
inventors' knowledge, no optical PUF
system closes a feedback loop in which the medium's evolving state
participates in a continuous control protocol, and none produces a
time-ordered, tamper-evident record of the full emission--observation
sequence.  Separately, Naughton, Hennelly, and Dowling adapted
cryptographic modes of operation (including cipher-feedback and
output-feedback modes) to optical encryption systems using
double-random-phase encoding (``Introducing secure modes of operation
for optical encryption,'' \emph{J.\ Opt.\ Soc.\ Am.~A}~25(10),
2608--2617, 2008), demonstrating chained optical transformations
for security purposes; however, those systems use fixed optical
elements rather than closed-loop feedback with evolving reactor state.
Robust key extraction from noisy PUF measurements was addressed by
\v{S}kori\'{c}, Tuyls, and Ophey (``Robust key extraction from
physical uncloneable functions,'' \emph{ACNS}, LNCS~3531,
pp.~407--422, 2005), establishing helper-data schemes that tolerate
measurement noise while preserving entropy---a concern shared by any
system that derives verification signals from physical observables.
Buchanan et al.\ demonstrated that laser speckle patterns from
everyday paper and plastic surfaces can serve as intrinsic
fingerprints (``Forgery: `fingerprinting' documents and packaging,''
\emph{Nature}~436, 475, 2005), extending microstructure-unclonable authentication
beyond bespoke tokens to ambient materials.
Mesaritakis et al.\ explored photonic PUFs based on multimode
waveguide interference (``Physical unclonable function based on a
multi-mode optical waveguide,'' \emph{Sci.\ Rep.}~8, 9653, 2018).  More recently, photonic neuromorphic PUFs
have been proposed that combine optical neural-network architectures
with PUF characteristics (Dermanis, Mesaritakis, et al., ``Photonic
physical unclonable function based on integrated neuromorphic
devices,'' \emph{J.\ Lightwave Technol.}~40(22), 7333--7341, 2022).
In the electronic domain, Wu et al.\ demonstrated a feedback-loop
arbiter PUF (FLAM-PUF) that uses output feedback to amplify response
entropy (``FLAM-PUF: a response--feedback-based lightweight
anti-machine-learning-attack PUF,''
\emph{IEEE Trans.\ Comput.-Aided Des.\ Integr.\ Circuits Syst.}~41(11),
4433--4444, 2022; DOI 10.1109/TCAD.2022.3197696).
However, these feedback and neuromorphic PUF systems operate in the
electronic or guided-wave optical domain, not in free-space
scene-coupled configurations, and none of the cited systems produces committed time-ordered
observation records.

\paragraph{Projector--camera feedback and dynamical systems.}
Optical feedback loops and multi-pass arrangements can
generate nonlinear input--output transformations from individually
linear or weakly nonlinear components.  Crutchfield analysed
camera--monitor feedback as a dynamical system producing emergent
spatiotemporal patterns (``Space-time dynamics in video feedback,''
\emph{Physica~D}, 1984).  Amano analysed projector--camera feedback
via reaction-diffusion equations (Amano, ``Analyzing the behavior of
projector-camera systems based on reaction-diffusion equations,''
\emph{ICAT-EGVE}, pp.~11--18, 2023).  Closed-cavity
optical accelerators circulate light between mirrors or modulators for
matrix-style inference.  However, these systems operate on internally
generated patterns without coupling to an external physical scene,
without multiple selectable operating regimes, and without the
bidirectional dynamical coupling described herein.

\paragraph{Time-varying metamaterials and temporal diffraction.}
Recent work on time-varying metamaterials has demonstrated that
thin-film epsilon-near-zero (ENZ) materials---such as indium tin
oxide---can switch optical state on timescales approaching an optical
cycle, producing temporal diffraction in the frequency domain (Tirole
et al., ``Double-slit time diffraction at optical frequencies,''
\emph{Nature Physics}~19, 999--1002, 2023).  Time-varying media can
also enable nonreciprocal responses under appropriate spatiotemporal
modulation conditions (Galiffi et al., ``Photonics of time-varying
media,'' \emph{Adv.\ Photonics}~4, 014002, 2022; Engheta,
``Metamaterials with high degrees of freedom: space, time, and more,''
\emph{Nanophotonics}~10, 639--642, 2021).  These
results demonstrate the feasibility of ultrafast temporal modulation
in materials already present in display hardware, but do not describe
closed-loop feedback systems, trainable physical transforms,
multi-regime operation, or committed evidence records.
Photonic neural networks perform matrix--vector multiplication using
interferometric meshes, with in-situ training via adjoint light
propagation (Hughes et al., ``Training of photonic neural networks
through in situ backpropagation and gradient measurement,''
\emph{Optica}~5, 864--871, 2018; Pai et al., ``Experimentally
realized in situ backpropagation for deep learning in photonic neural
networks,'' \emph{Science}~380, 398--404, 2023).  Spall, Guo, and
Lvovsky demonstrated optical backpropagation through both linear and
nonlinear layers (``Training neural networks with end-to-end optical
backpropagation,'' \emph{Advanced Photonics}~7, 016004, 2025).
Pai et al.\ demonstrated a digitally verifiable photonic hash
primitive for blockchain verification on programmable photonic
interferometer meshes
(LightHash, ``Experimental evaluation of digitally-verifiable photonic
computing for blockchain and cryptocurrency,'' \emph{Optica}~10,
552--560, 2023),
achieving computation and verification on shared hardware.  However,
LightHash uses linear systems for both computation and verification;
no cited system pairs linear compute media with nonlinear attestation
media, and no cited photonic computing system produces committed
\cba records of its computational trajectory.
Separately, proposals for remote attestation of optical systems
using optical identification and blockchain-style immutability have
been described, but these bind identity to a fixed optical
fingerprint rather than to the full temporal trajectory of a
computation through a physical medium.

\paragraph{Camera-sensor fingerprinting and dual-use hardware.}
Photo-response non-uniformity (PRNU) fingerprinting has shown
that the same camera data used for perception simultaneously contains
device-identity information (Luk\'{a}\v{s}, Fridrich, and Goljan,
``Digital camera identification from sensor pattern noise,''
\emph{IEEE TIFS}~1(2), 205--214, 2006).  The Ed-PUF integrates microstructure-based authentication
into dynamic-vision-sensor pixel circuits (Zheng et al.,
``Ed-PUF: Event-driven physical unclonable function for camera
authentication in reactive monitoring system,''
\emph{IEEE~TIFS}~15, 2824--2839, 2020), achieving dual-use sensing
and verification on shared hardware.  These systems share the
observation that a single data stream can serve multiple objectives,
but neither cited system provides
controllable rendering, reactor-mediated
dynamics, or a unified regime framework.

\paragraph{PUF security, attacker models, and hardness.}
Machine-learning attacks on PUFs demonstrate that challenge--response
data can train models that predict device responses (R\"{u}hrmair et
al., ``Modeling attacks on physical unclonable functions,''
\emph{ACM~CCS}, pp.~237--249, 2010).  R\"{u}hrmair and van~Dijk
proposed formal attacker models conditioning security evaluations on
attacker capabilities (``PUFs in security protocols: Attack models and
security evaluations,'' \emph{IEEE~S\&P}, pp.~286--300, 2013).
Ganji, Tajik, and Seifert formalised PUF learnability via PAC learning
(``PAC learning of arbiter PUFs,''
\emph{J.~Cryptographic Engineering}~6(3),
249--258, 2016).  These works implicitly treat security as
dependent on attacker advancement, but none formalises hardness as
an explicit two-dimensional index managed as a race condition with
periodic re-evaluation.


\paragraph{Self-configuring photonic systems and adaptive reservoirs.}
Self-configuring universal linear optical components use progressive
local feedback loops (detectors and controls at each stage) to
configure and stabilise programmable photonic circuits against drift
and cross-talk (Miller, ``Self-configuring universal linear optical
component,'' \emph{Photon.\ Res.}~1(1), 1--15, 2013; subsequent
demonstrations of management automation for programmable PICs).
Reservoir computing and in-materio computing study nonlinear
dynamical substrates whose internal dynamics can be tuned via
plasticity rules or homeostatic mechanisms to improve task
performance.  These systems anticipate ``inward sensing for
calibration'' and ``internal parameter adaptation,'' but they treat
internal trajectories as transient control signals, not as committed,
auditable evidence artifacts integrated into a trust or attestation
framework.  Conditional computation and mixture-of-experts
architectures have been proposed for photonic and optical systems,
including routed optical reservoirs and switchable scattering
elements, but existing proposals do not commit the routing path,
expert identity, or execution transcript into a provenance-bound
evidence record.

\paragraph{Stigmergy and indirect multi-agent communication.}
Stigmergy---indirect coordination via persistent environmental
modifications---has been studied in swarm robotics, including
engineered pheromone-release systems and formalised
environment-mediated memory.  These systems demonstrate that physical
traces can serve as a shared communication substrate among agents.
However, stigmergic traces are typically low-dimensional coordination
variables (scalar intensities, binary markers), not high-dimensional
encodings of internal dynamical state designed to be decoded through a
specific perception channel and to influence the receiver's internal
attractor landscape.

\paragraph{Runtime assurance and evaluation partitioning.}
The Simplex architecture and related runtime assurance (RTA)
frameworks separate an untrusted advanced controller from a trusted
safety controller, with a monitor switching control authority when
safety properties may be violated.  This architectural partition
between performance optimisation and safety evaluation is a
conceptual precursor to the meter partition described herein.
However, Simplex-style systems enforce binary switching (safe
controller vs advanced controller) rather than a continuous partition
between optimisation-visible and held-out evaluation metrics, and
they do not address self-modifying physical systems whose evaluated
properties (empirical hardness, attractor topology) can change under
the agent's own control.

\paragraph{Calibration, model validation, and distributed quality improvement.}
Bayesian calibration frameworks include an explicit model discrepancy
term to account for systematic differences between computer models and
physical observations, treating discrepancy as a nuisance to be
estimated and marginalised (Kennedy and O'Hagan, ``Bayesian calibration
of computer models,'' \emph{J.\ R.\ Stat.\ Soc.\ Ser.~B}~63(3),
425--464, 2001).
Adversarial benchmarking platforms such as Dynabench reward
participants for discovering inputs that cause trained models to fail,
but operate entirely in the digital domain on software models.
Metrological proficiency testing under ISO~17043 distributes
circulated artefacts among laboratories to assess measurement
competence, but uses identical test items rather than distinct
unclonable devices.  Recent work has demonstrated that physically
unclonable optical scattering devices can produce measurement data
that reveals systematic model errors when compared across
independently manufactured instances.  However, none of these cited systems uses
physical reproduction on distinct unclonable devices as a verification
mechanism for model improvement, and none treats model error as a
multi-role resource serving simultaneous model improvement, device
identification, and security functions.

\paragraph{Gap addressed by the present disclosure.}
None of the technologies identified in the prior art survey provides a single physical operator
that supports all of verification, perception, controllable rendering,
and dynamic scene coupling under a unified mathematical framework, with
logged protocols and auditable evidence.  Digital provenance systems
(C2PA and related standards) commit metadata via cryptographic keys
but do not bind authenticity to the physics of the capture event.
Self-configuring photonic systems calibrate internal parameters but do
not produce committed inner-state records for calibration or audit.
Conditional computation and mixture-of-experts architectures route
inputs through physical substrates but do not commit the routing path,
expert identity, or execution transcript into a provenance-bound
evidence record.
Runtime assurance architectures
partition safety from performance but do not address continuous
evaluation metrics held out from physically evolving systems.
Nor does any existing system
unify a physical reactor medium (providing empirical hardness and
device-specific signatures) with committed evidence records, multiple
operating regimes selectable by objective function, inward-facing
self-calibration modes, routed physical ensemble computation with
committed execution paths, and fleet-level model
improvement through physically verified discrepancy discovery.

% ======================================================================

\section{Claim Tier Structure}
\label{sec:claim-tier-structure}

The claims in this disclosure are organised into three tiers for which
support and proof requirements are non-transitive across tiers. A
limitation, disproof, or narrowing of claims in one tier does not
propagate to the other tiers, and proof of a limitation in one tier does
not, without separately recited and separately supported limitations,
establish a limitation in another tier. Facts disclosed in a single
embodiment may be relevant to more than one tier, but the
written-description support and empirical proof required for a
tier-specific limitation must be satisfied within that tier.

\paragraph{Tier A --- Attestation and provenance claims.}
These claims concern commitment infrastructure, protocol digests,
reactor-microstructure-unclonability binding, hash-chain linkage, fleet corroboration,
proof-of-discrepancy, and semantic \TB verification. Proof of a
Tier A limitation does not, without separately recited and separately
supported limitations, establish a Tier B or Tier C limitation.

\paragraph{Tier B --- Performance claims for specific computational primitives.}
These claims concern the speed, energy efficiency, and native
parallelism of the physical substrate for matrix-vector multiplication,
two-dimensional correlation, Fourier and fractional-Fourier transforms,
random feature projection, energy-based sampling, in-sensor
preprocessing, and active perception under meter-bounded protocols.
Proof of a Tier B limitation does not, without separately recited and
separately supported limitations, establish a Tier A or Tier C limitation.

\paragraph{Tier C --- Capacity and training-scale claims.}
These claims concern trainable parameter counts, task-useful effective
rank after readout training, and gradient-extraction throughput.
Capability assertions in Tier C are tied to measured task-useful rank
per benchmark and per embodiment, and are explicitly bounded by the
capacity accounting rule of Section~\ref{sec:capacity-accounting-rule}.

No claim, recital, or drafting language in this disclosure shall be
construed as allowing a conclusion in one tier to be supported by
evidence or assertions proper to another tier. In particular,
attestation support (Tier A) does not support capacity claims (Tier C),
and primitive-performance support (Tier B) does not support open-ended
semantic scaling claims (Tier C).

\section{Summary of the Invention}
% ======================================================================

\subsection{The physical picture}

A \RK module~100 (FIG.~1) is, at its simplest, a device that directs
structured emission (typically light, but more generally any
controllable physical carrier) toward a physical channel and records
what comes back. In representative embodiments it may include a reactor emitter~20, a reactor detector~40, and a reactor~30---a
physical medium, optionally nonlinear, with memory and
feedback---which may be integrated in a single assembly with a shared controller~10 or
distributed across modular components.
The controller~10 drives a scan protocol: it decides where to point, how
bright to shine, what wavelength and polarisation to use, how long to
integrate. The scene~50---whatever is out there---reflects, scatters,
absorbs, fluoresces, or otherwise transforms the emitted light,
and is interrogated by a scene-facing emitter, shown
schematically in FIG.~1 as a scene emitter~52, and observed by a
scene-facing detector, shown schematically in FIG.~1 as a scene
detector~51. The reactor~30 inside module~100 is interrogated by
the reactor emitter~20 and observed by the reactor detector~40,
and the scene-side and reactor-side observations are coupled
within a single closed cycle as detailed in
Section~\ref{sec:architectural-alternatives}. The pair
$(u(t), \mathbf{y}_t)$---what was emitted and what was observed
across the cycle---is one sample of the \cb.
A run produces a time-ordered sequence $C_{0:T}$ of such samples, forming
a \cb~60: a
physical record of a structured conversation between the device and
its environment.

In some embodiments, what distinguishes the assembly from a projector--camera
pair is the inclusion of a reactor. The
reactor is a physical medium---a scattering plate, a nonlinear
crystal, a phosphor screen, a fibre-delay loop, a CRT, a biological
tissue, or any combination---that sits in the optical path and
contributes its own dynamics. A reactor may contribute dynamics
through memory (phosphor persistence, charge traps, thermal states),
through nonlinearity (saturation, bistability, chaotic mixing),
and through manufacturing idiosyncrasies
(grain structure, defects, wear patterns) that make each reactor
physically unique---singly or in combination.
Empirical hardness arises primarily from nonlinearity and
microstructure; a linear or identity reactor remains a legitimate
operating configuration for applications where hardness is not the
objective (see the reactor-complexity discussion in
Section~\ref{sec:regimes}).
The reactor transforms the light in ways that
depend on its current state, its history, and its configuration
parameters $\theta$---and in ways that may be impractical to reproduce
faithfully by a digital simulation without reproducing the full physical
complexity.

The result is a device whose output at any moment depends on (a)~what
the controller chose to emit, (b)~what the scene did to the light,
(c)~what the reactor did to the light, and (d)~the reactor's
accumulated history. This four-way dependence is the source of
everything the device can do: verification (because the reactor's
response is empirically hard to clone under declared attacker
families---see Section~\ref{sec:security-theory}), sensing (because the scene's
contribution is information-rich), and rendering (because the
controller can steer the output toward a target).

\subsection{Architectural alternatives --- parallel and cascade embodiments}
\label{sec:architectural-alternatives}

The physical picture of the preceding subsection deliberately does not
commit to a particular topology for how the emitter, reactor, detector,
and scene are interconnected. Two distinct architectural embodiments
are disclosed, both within the scope of this specification, and both
non-limiting:

\paragraph{Embodiment (a): parallel-subsystem architecture with
A(t)-only coupling.}
In some embodiments, the apparatus is implemented as two physically
separate subsystems --- a \emph{scene subsystem} comprising a
scene-facing emitter, an outward optical path to the scene, and a
scene detector that observes the scene's response; and a
\emph{reactor subsystem} comprising a reactor-facing emitter (which
may be a separate light source or a derived signal source), the
reactor element, and a reactor detector that observes the reactor's
response. In this embodiment, the two subsystems are independent
optical loops: the scene-emitter does not illuminate the reactor, and
the reactor-emitter does not illuminate the scene. The single
coupling between the two subsystems is a closed-loop feedback path
from the reactor detector through a low-latency conditioning chain
back to the scene-emitter modulation, herein denoted $A(t)$. The
$A(t)$ feedback path closes the cycle: in each sample period, the
controller's scan protocol synchronises both the next scene
interrogation and the next reactor interrogation across the two
subsystems, and the reactor's response is fed back via $A(t)$ to
modulate the next scene-emitter excitation. The scene-detector
output in embodiment (a) is recorded into the \cb
but is not delivered into the reactor side of the cycle by any
forward-coupling path; rather, the temporal alignment of scene
and reactor interrogations is established by the shared scan
protocol acting as a synchronisation mechanism, not as an
architectural coupling. Embodiment (a) is described in further
detail at the 1D anchor of Section~\ref{sec:anchor-1d} and the 2D
anchor of Section~\ref{sec:anchor-2d}. The parallel two-subsystem
embodiment (a) is obtained from the layouts of FIGs.~4A and~6 by
omitting the forward low-latency coupling arrow shown in those
figures and relying on the $A(t)$ feedback path as the sole
architectural coupling between the scene and reactor sides; all
other components of FIGs.~4A and~6 are common to both embodiments.

\paragraph{Embodiment (b): cascade architecture with forward
low-latency coupling.}
In some embodiments, the apparatus is implemented as a single closed
cycle in which the scene observation is used as the direct, live
input to the reactor stage. Specifically, the cycle proceeds:
scene-emitter $\to$ scene $\to$ scene-detector $\to$
\emph{forward low-latency coupling} $\to$ reactor-emitter $\to$
reactor $\to$ reactor-detector $\to$ $A(t)$ feedback $\to$
scene-emitter (next cycle). The forward low-latency coupling is the
load-bearing element distinguishing embodiment (b) from embodiment
(a): it is a direct, bounded-latency signal path from the scene
detector to the reactor emitter, with end-to-end latency declared
and committed to the protocol digest, and with no per-period wait,
no frame buffer, and no semantic scene-decoding, frame-level
processing, or runtime-adaptive processing stage on that forward
signal path. In declared mixed-signal implementations,
digitisation, the fixed deterministic transfer function, and any
re-synthesis are part of the forward low-latency coupling and are
included in the committed latency budget. The forward low-latency
coupling may be implemented in any
of the following non-limiting ways:
\begin{itemize}[nosep]
  \item \emph{Optical implementation, gain-pumped media.} The
    scene-detector output drives a gain-pumped active optical
    medium (for example a semiconductor optical amplifier, a doped
    fibre amplifier, or a pumped solid-state gain medium) whose
    output illuminates or otherwise drives the reactor emitter.
    The gain-pumping permits the forward coupling to operate at
    optical signal levels without intermediate electrical
    re-representation.
  \item \emph{Optical implementation, direct relay.} The
    scene-detector output drives a directly-modulated optical
    source whose emission is relayed onto the reactor emitter
    surface or input aperture. Suitable for embodiments where
    optical-to-optical relay is preferred over gain-pumped
    amplification.
  \item \emph{Electrical implementation.} The scene-detector
    photocurrent or transimpedance-amplifier output drives the
    reactor emitter's modulation input directly, optionally
    through a passive impedance-matching or a bandwidth-shaping
    network whose transfer characteristic is declared and
    committed to the protocol digest. Suitable for low-cost
    embodiments where optical relay is not justified.
  \item \emph{Mixed-signal implementation.} The scene-detector
    output is digitised at the sample period of the cycle and the
    digital sample drives the reactor emitter through a fixed,
    deterministic, declared transfer function whose tap weights
    and timing parameters are committed to the protocol digest.
    The digitisation depth, sample rate, and total digital-domain
    latency are declared as part of the latency budget.
\end{itemize}
The reactor in embodiment (b) is typically the primary trainable
element of the cycle, with reactor parameters (microstructure,
biasing, internal state) updated according to the training protocols
of Section~\ref{sec:meters-training} and the anchor-specific
training pathways of Section~\ref{sec:anchor-training} (1D anchor)
and Section~\ref{sec:anchor-2d-trainable} (2D anchor) during
enrolment phases and held fixed during inference phases. The
forward low-latency coupling and the $A(t)$ feedback path are
themselves typically not the primary trainable elements; rather,
their transfer characteristics are declared and committed at
deployment time, and the reactor carries the learnable degrees of
freedom of the cycle.

\paragraph{Single closed cycle in both embodiments.}
Both embodiments (a) and (b) implement a single closed cycle in
which scene observations are causally bound to reactor responses
within a declared per-cycle latency budget. The two embodiments
differ in \emph{how} the cycle is closed: in (a), closure is via
the $A(t)$ feedback path alone, with the controller's scan protocol
serving as the synchronisation mechanism between the otherwise
independent scene and reactor subsystems; in (b), closure is via
both the forward low-latency coupling and the $A(t)$ feedback
path, with the forward coupling carrying the scene observation
directly into the reactor stage in real time. In both embodiments,
the \cb of Definition~\ref{def:convolution-bundle}
records the cycle as a multi-tap trace covering scene-side and
reactor-side observations and, where physically tappable, the
intermediate states of the forward coupling and the $A(t)$
feedback path.

\paragraph{Common controller, common scan protocol.}
In both embodiments, a common controller drives all active
elements of the cycle (scene-emitter, scene-detector, reactor-emitter,
reactor-detector, and any tunable elements of the forward coupling
or $A(t)$ feedback path). Where the controller distributes a
shared scan program to multiple servos --- as in the 2D-anchor
embodiment of Section~\ref{sec:anchor-2d} --- the shared scan
program is a synchronisation mechanism, and is structurally
distinct from the architectural coupling paths (forward low-latency
coupling in (b); $A(t)$ feedback in both): the shared scan
program does not, on its own, constitute a forward coupling
within the meaning of embodiment (b).

\paragraph{Implementation modalities of the A(t) feedback path
(non-limiting).}
The $A(t)$ feedback path may be implemented in any of the
following non-limiting ways, paralleling the implementation
variants disclosed above for the forward low-latency coupling:
\begin{itemize}[nosep]
  \item \emph{Electrical implementation.} The reactor-detector
    photocurrent or transimpedance-amplifier output is
    conditioned by a deterministic low-latency analogue chain
    --- including the frozen committed reactor-detector gain
    $G_{\mathrm{det}}$ --- whose output drives the scene-emitter
    modulation input directly. Suitable for bench-top and
    early-stage embodiments where electrical conditioning meets
    the per-cycle latency budget.
  \item \emph{Mixed-signal implementation.} The reactor-detector
    output is digitised at the sample period of the cycle, the
    digital sample is processed by a fixed deterministic transfer
    function (for example the look-up-table or sampled
    $G_{\mathrm{det}}$ implementations of
    Section~\ref{sec:anchor-2d-coupling}), and the result is
    re-synthesised as an analogue scene-emitter modulation drive.
    Digitisation depth, sample rate, and total digital-domain
    latency are declared as part of the latency budget.
  \item \emph{Optical implementation, gain-pumped media.} The
    reactor-detector output drives a gain-pumped active optical
    medium whose output illuminates or otherwise drives the
    scene-emitter input. Suitable for embodiments where optical
    conditioning is preferred over intermediate electrical
    re-representation.
  \item \emph{Optical implementation, direct relay.} The
    reactor-detector output drives a directly-modulated optical
    source whose emission is relayed onto the scene-emitter input
    aperture. Suitable for embodiments where optical-to-optical
    relay is preferred over gain-pumped amplification.
  \item \emph{Hybrid optical-electrical implementation.} The
    $A(t)$ feedback path comprises an electrical conditioning
    stage and an optical conditioning stage in series, in either
    order, with the boundary between stages declared and committed
    to the protocol digest.
\end{itemize}
The bench-top embodiments described elsewhere in this disclosure,
including the 1D anchor of Section~\ref{sec:anchor-1d} and the 2D
anchor of Section~\ref{sec:anchor-2d}, and as illustrated in
FIGs.~4A and~6, are non-limiting electrical instances of the
$A(t)$ feedback path. Hybrid optical-electrical and fully-optical
$A(t)$ feedback embodiments are within the scope of this
disclosure and represent contemplated implementations for
embodiments where end-to-end optical operation is achievable
within the declared latency budget. The forward low-latency
coupling and the $A(t)$ feedback path are independent in their
implementation modalities: a given apparatus may, for example,
combine an optical forward coupling with an electrical $A(t)$,
or an electrical forward coupling with an optical $A(t)$, as
declared and committed to the protocol digest.

\paragraph{Tiny neural-network element on the A(t) path
(non-limiting hybrid embodiment).}
In some hybrid embodiments of either (a) or (b), the $A(t)$
feedback path includes a small trainable neural-network element
with a memory window of one cycle period (or a small integer
multiple thereof), arranged to apply a local gain or a local
non-linear shaping to the conditioning chain output before the
signal is delivered to the scene-emitter modulation input. The
neural-network element is typically constrained to a small
parameter count and a fixed inference-time budget, so that its
contribution to the per-cycle latency is bounded within the
declared latency budget of Section~\ref{sec:wm-latency-budget}.
The neural-network parameters are declared and committed to the
protocol digest and are not modified during inference phases.
The neural-network element is non-limiting and may be omitted
in embodiments where deterministic linear or piecewise-linear
conditioning suffices.

\paragraph{Mapping to figures.}
FIG.~1 is a high-level block diagram of the closed-cycle signal
relationships. The $A(t)$ feedback path from the reactor
detector~40 to the scene emitter~52 is depicted as a solid arrow
and is present in both embodiments (a) and (b). The forward
low-latency coupling from the scene detector~51 to the reactor
emitter~20 is depicted as a dashed arrow and is present only in
the cascade embodiment (b); it is omitted in the parallel-subsystem
embodiment (a). FIGs.~4, 4A, and~6 each show the cascade
embodiment (b) as the primary illustrated topology, with the
forward low-latency coupling drawn as a heavy-stroke arrow from
the scene detector to the reactor emitter, labelled
\emph{LOW-LATENCY COUPLING}. The parallel embodiment (a) of these
same anchor configurations is obtained by removing the forward
low-latency coupling arrow and relying on the $A(t)$ feedback
path as the sole architectural coupling between the scene and
reactor sides; both embodiments share all other components of
FIGs.~4, 4A, and~6.


\subsection{The mathematical abstraction}

The physical picture can be formalised by treating the entire assembly as
a parameterised Markov kernel:
\[
  \mathsf{P}_\theta : \mathcal{S} \times \mathcal{U}^{T+1}
  \;\to\; \mathcal{P}(\mathcal{C}),
\]
mapping a scene $S \in \mathcal{S}$ and a control protocol
$U_{0:T} = \{u(t)\}_{t=0}^T$ to a probability law over \cba
trajectories $C_{0:T} \in \mathcal{C}$. The configuration parameter
$\theta = (\theta_{\mathrm{hw}}, \theta_{\mathrm{sw}})$ includes hardware parameters $\theta_{\mathrm{hw}}$ (lens positions, gain settings,
reactor state) and optional learned parameters $\theta_{\mathrm{sw}}$. This abstraction
unifies diverse physical instantiations---optical, acoustic, RF,
biological---under a common formalism and supports modular reactor
assemblies.

\subsection{Three operating regimes}

\paragraph{Point of novelty (plain-language summary for prosecution purposes, non-limiting).}
In some embodiments, the disclosure provides a \textbf{physical
Markov channel apparatus} that produces a \textbf{canonical,
tamper-evident \cb} --- a time-ordered record pairing
each emitted probe with its measured physical response --- under a
\textbf{committed protocol digest and meter envelope}, such that the
same physical hardware and output record format supports three
selectable operating objectives (physical verification, active
sensing, and constrained rendering) and an optional closed-loop
dynamical coupling modifier (Yoked operation), all governed by a
common fleet-level model improvement and discrepancy verification
protocol (Proof of Discrepancy).  In some embodiments, this
combination is distinct from
static optical PUF verifiers (which produce snapshot challenge--response
pairs under no committed protocol), structured-light perception
systems (which do not produce authentication-grade committed
records), projection mapping systems (which do not enforce
meter-bounded admissibility or produce tamper-evident logs), and
optical computing accelerators (which do not treat the
emission--observation transcript as a first-class committed
security artefact).

\subsection{Key aspects}

In one aspect, the principal deployment of the apparatus and methods of
this disclosure is a plurality of \RK modules operating as a
\emph{witness mesh}: multiple modules, each comprising a fast closed
loop in causal feedback with a portion of physical reality observable to
that module, mutually coupled through continuous, low-latency analogue
inter-module signal paths, such that the joint \cba record across the
mesh is empirically hard to forge under bounded adversaries. The
single-module aspects recited below are non-limiting building blocks of
the witness-mesh architecture and may also be used as standalone
instantiations. The witness-mesh deployment is described in detail in
Section~\ref{sec:witness-mesh}.

In one aspect, a system comprises a \RK module including at
least one emitter, at least one scanner or pattern-address mechanism, at
least one scene-facing optical path, at least one detector, and a
controller that executes a scan protocol $U_{0:T}$. The module produces
a \cba sequence $C_{0:T}$ in which each sample includes
(a)~a scan coordinate or pattern index and emission parameters, and
(b)~one or more measured responses.

In another aspect, the \RK is represented as a parameterised
Markov kernel $\mathsf{P}_\theta(\cdot\mid S,U_{0:T})$ from
scenes and control protocols to \cba trajectories, where
$\theta = (\theta_{\mathrm{hw}}, \theta_{\mathrm{sw}})$ includes hardware configuration parameters $\theta_{\mathrm{hw}}$ and optional
learned parameters $\theta_{\mathrm{sw}}$.

In another aspect, the same physical module supports multiple operating
regimes (FIG.~2) differentiated by objective families applied to the \cb,
including:
\begin{enumerate}
  \item \textbf{\TB}~110: verification and authenticity assessment
    using empirical hardness of forging \cba traces under
    bounded adversaries;
  \item \textbf{\LI}~120: active sensing and perception using task
    losses and/or information-theoretic objectives;
  \item \textbf{\RT}~130: controllable rendering and style
    control using style or perceptual losses while maintaining
    constraint to the physical channel.
\end{enumerate}


Empirical hardness reported in this disclosure is a measured, time-indexed,
attacker-indexed quantity, conditioned on a declared attacker family,
compute budget, and meter envelope recorded in the protocol digest. The
disclosed framework includes re-evaluation, reactor-refresh, and
query-throttle mechanisms that maintain useful hardness over time;
hardness values are not asserted as permanent or unconditional.

In another aspect, each of the above regimes can be operated in a Yoked
mode~200 (FIG.~3) in which the feedback loop is configured so that device and scene
are driven toward a coupled dynamical state, targeted by
balanced transfer entropy, conditional Lyapunov exponent approaching
zero from below (negative values indicating stable synchronisation,
per the Pecora--Carroll convention), and a joint attractor whose
structure is empirically assessed via meters rather than asserted
\emph{a priori}, so that the \cb encodes dynamical
properties of the scene that are not accessible in static operation.

\paragraph{Modular swept-kernel assemblies (non-limiting).}
In some embodiments, a \RK is assembled from modular
components---including, without limitation, lasers, galvanometers,
spatial light modulators, CRTs, phosphor screens, diffusers, nonlinear
media, and detector arrays---arranged in a variety of topologies rather
than a fixed optical bench. In these embodiments, the parameterised
swept kernel $\mathsf{P}_\theta(\cdot \mid S, U_{0:T})$, which governs $C_{0:T} \sim \mathsf{P}_\theta(\cdot \mid S, U_{0:T})$, arises from
the aggregate behaviour of the assembled elements under a shared or
coordinated scan protocol. Protocol digests and meter envelopes adapt
to the current configuration, and the \cb remains the canonical
output regardless of assembly topology.

\begin{remark}[Cross-channel structure as a unifying primitive
(non-limiting)]
In some embodiments, all three regimes exploit a shared primitive:
statistical dependencies between channels with different invariance
properties. The regimes differ in which dependencies they optimise and
how they use them---\LI maximises reconstruction or task
performance from joint observations, \RT separates
style-bearing from content-bearing channels, and \TB treats the
full joint record as an empirically hard-to-forge signature under bounded
adversaries---but the cross-channel dependence structure is shared.
Yoked operation adds temporal cross-channel coupling: not only do
different spatial or spectral channels exhibit dependencies, but the
time-lagged causal dependencies between device output and scene
response (and vice versa) become first-class observables.
\end{remark}

\begin{table}[t]
  \centering
  \begin{tabularx}{\textwidth}{p{0.13\textwidth} p{0.10\textwidth}
    p{0.17\textwidth} p{0.22\textwidth} Y}
    \toprule
    Regime & Essence & Optimised quantity & Static mode &
      Yoked mode \\
    \midrule
    \TB
      & Illumination
      & $I(\theta)$
      & Fingerprint verification
      & Continuous tamper detection; coupling-based auth \\[4pt]
    \LI
      & Perception
      & $I(S; Y)$
      & Reflectance, geometry, spectra
      & Resonances, damping, transfer functions \\[4pt]
    \RT
      & Creation
      & $d(Y, Y^*)$
      & Projection mapping; style transfer
      & Entrainment-driven display; responsive environments \\
    \bottomrule
  \end{tabularx}
  \caption{Three regimes on a single \RK, each operating in
    static or Yoked mode. The underlying physical operator is shared;
    objectives, training, and coupling depth differ.}
  
\end{table}

\begin{table}[ht]
  \centering
  \begin{tabularx}{\textwidth}{p{0.15\textwidth} p{0.14\textwidth}
    p{0.30\textwidth} Y}
    \toprule
    Regime & Question & Static design dial & Yoked design dial \\
    \midrule
    \TB
      & \emph{Who are you?}
      & Increase microstructure sensitivity; widen
        verifier--adversary gap
      & Exploit coupled dynamics as security primitive;
        scene as continuous unpredictable challenge \\[0.3em]
    \LI
      & \emph{What is out there?}
      & Increase invertibility and SNR; maximise $I(S;Y)$
      & Sweep coupling parameters to map dynamical
        response; construct Arnold tongue profiles \\[0.3em]
    \RT
      & \emph{Make it look like this.}
      & Balance ambiguity vs control; respect meter
        envelopes
      & Steer joint attractor toward target behaviour;
        specify coupling profile as rendering target \\
    \bottomrule
  \end{tabularx}
  \caption{Static and Yoked design dials for each regime.}
  
\end{table}

\begin{remark}[Selective dissipation across regimes (non-limiting)]
In some embodiments, both \LI and \TB configurations employ
selective dissipation; they differ in selection criteria rather than in
the presence or absence of information loss. \LI dissipation
profiles preserve task-relevant mutual information while suppressing
nuisance variables, noise, and irrelevant degrees of freedom. \TB dissipation profiles preserve correlation structure tied to reactor
microstructure while destroying information an adversary would need to
invert or forge the channel under bounded resources and disclosure
policies. \RT profiles balance controllability with
ambiguity by selecting which aspects of the \cba record are
treated as controllable style degrees of freedom versus constrained
content degrees of freedom. In Yoked operation, dissipation profiles
additionally preserve temporal coupling structure---the time-lagged
causal dependencies between device and scene that encode dynamical
properties.
\end{remark}

\paragraph{Continuous regime mixing (non-limiting).}
In some embodiments, the operating regime is set by continuously varying
weights rather than selecting a discrete mode. A composite objective
combines \TB, \LI, and \RT terms with weights
$\lambda_{\mathrm{TB}}$, $\lambda_{\mathrm{L}}$, and
$\lambda_{\mathrm{RT}}$ on the two-simplex $\Delta_2$, and a \yd
$\alpha \in [0,1]$ controls the strength of bidirectional dynamical
coupling. The full operating point
$(\lambda_{\mathrm{TB}}, \lambda_{\mathrm{L}},
\lambda_{\mathrm{RT}}, \alpha)$ specifies a point in a triangular
prism: the TB--L--RT triangle at each value of $\alpha$, extruded along
the \yda axis. In preferred embodiments, the weights, \yd,
and any transition schedule are recorded in the protocol digest so that
meters and audits condition on the declared regime mix and coupling mode.

In another aspect, the system includes one or more reactor subsystems
and/or memory media that introduce history dependence and increase
diversity of response statistics. In some embodiments, the system
further includes multi-channel sensing with distinct invariance profiles,
and evaluation is performed on cross-channel dependence structure.

In another aspect, networked deployments exchange commitments or
summaries of \cba traces to support transitive or multi-party
verification, including proof-of-projection patterns (empirically
calibrated attestation protocols; see terminological note in
Section~\ref{sec:networks}).

In another aspect, a fleet of devices jointly improves a shared
predictive model through a proof-of-discrepancy protocol
(Section~\ref{sec:proof-of-discrepancy}) in which participants submit
reproducible physical methods identifying where the predictive model
diverges from measured device responses, and independent devices verify
claimed discrepancies by physical re-execution.  The predictive model
accepts a challenge configuration and a device configuration
representation inferred from calibration measurements, and the
discrepancy between prediction and measurement is decomposed into
shared model error (driving model improvement), device-specific
captured variation (enabling fleet management and task allocation), and
an irreducible residual whose physical basis supports clone resistance under the declared attacker model.

\subsection{Notation preview}
Throughout this specification, $\mathsf{P}_\theta$ denotes the
parameterised stochastic kernel (\RK), $S$ denotes the
scene tuple, $U_{0:T}$ denotes the control protocol over the
observation interval, $C_{0:T}$ denotes the resulting \cb trajectory,
$\theta = (\theta_{\mathrm{hw}}, \theta_{\mathrm{sw}})$ denotes
hardware and optional software configuration parameters, and
$\alpha \in [0,1]$ denotes the \yd (coupling strength between device
and scene).  The symbols $\TE$ and $\CLE$ denote transfer entropy and
conditional Lyapunov exponent, respectively.  Formal definitions of
all terms and a complete symbol table are provided in Section~\ref{sec:definitions} (Definitions
and Glossary).

% ======================================================================
\section{Brief Description of the Drawings}
% ======================================================================

The accompanying drawings are schematic and non-limiting.

\textbf{FIG.~1} is a schematic block diagram of a \RK
module~100 showing the controller~10, reactor emitter~20, reactor~30
(physical medium), and reactor detector~40 disposed within
module~100, an external scene~50 comprising at least one scene
detector~51 and at least one scene emitter~52, and the resulting
\cb~60 $C_{0:T}$ including protocol digest and meter outputs.
Reactor emitter~20, reactor~30, and reactor detector~40 are shown
in horizontal arrangement within module~100, with reactor~30
disposed in the device-internal signal path between reactor
emitter~20 and reactor detector~40. Two cross-couplings between
the device and the external scene are shown: an $A(t)$ feedback
path from the reactor detector~40 to the scene emitter~52,
depicted as a solid arrow and present in both the
parallel-subsystem embodiment~(a) and the cascade embodiment~(b)
of Section~\ref{sec:architectural-alternatives}; and an optional
forward low-latency coupling from the scene detector~51 to the
reactor emitter~20, depicted as a dashed arrow and present only
in the cascade embodiment~(b). In the parallel-subsystem
embodiment~(a), the scene detector~51 is recorded and
time-aligned into the \cb~60 but does not drive
the reactor emitter~20; the dashed forward arrow is accordingly
omitted. The \cb~60 records, in some embodiments,
multiple taps along the cycle (scene-side excitation and
observation; reactor-side excitation and observation; and, where
physically tappable, intermediate states of the forward
low-latency coupling and the $A(t)$ feedback path), and is not
limited to endpoint emission--observation pairs.

\textbf{FIG.~2} is a schematic diagram showing how a single shared
physical operator~100 supports three operating regimes (\TB~110, \LI~120, \RT~130) selected by objective
weights, with the operating point parameterised by
$(\lambda_{\mathrm{TB}}, \lambda_{\mathrm{L}},
\lambda_{\mathrm{RT}}, \alpha)$.

\textbf{FIG.~3} is a schematic diagram illustrating the contrast
between static operation ($\alpha = 0$, separable device and scene
contributions) and the maximum-commanded-coupling side of the diagram
($\alpha \to 1$, coupled dynamical state with joint attractor structure),
with Yoked operation~200 declared only when the applicable CLE/TE
co-condition and confidence conditions recorded in the protocol
digest are satisfied; the \yd $\alpha$ parameterises coupling depth
rather than constituting the declaration by itself.

\textbf{FIG.~4} is a schematic diagram of the CRT/phosphor
analogue-memory extension embodiment
(Section~\ref{sec:crt-phosphor-extension}), comprising a fast
optical loop (laser source~21, AOD/EOM deflector~22, relay
optics~23, reactor medium~30, photodetector~41) and a slow
analogue-memory loop (CRT reactor~31 with electron gun, yoke,
and phosphor screen~32, observed by a rolling-shutter
camera~42), under a common controller~10. The cascade
embodiment~(b) of
Section~\ref{sec:architectural-alternatives} is illustrated in
FIG.~4 by a forward
low-latency coupling drawn as a heavy-stroke arrow from
photodetector~41 to the slow-loop reactor stage; this forward
coupling carries the scene-side observation directly into the
analogue-memory stage and distinguishes the cascade embodiment
from the parallel-subsystem embodiment~(a) of
Section~\ref{sec:architectural-alternatives},
which is obtained by removing the forward coupling arrow. This
extension adds spatially indexed analogue memory atop the 1D
anchor of FIG.~4A, the 2D generalisation of FIG.~6, or another
declared anchor-compatible reactor-to-scene coupling topology.

\textbf{FIG.~4A} is a schematic diagram of the bench-top 1D
\RK of Section~\ref{sec:anchor-1d}, comprising a
scene subsystem with a scene laser, a mirror on a single-axis
micromotor servo, and a scene detector, and a sealed reactor
subsystem comprising an aluminium cup or can with a mixed
fluorescent and scattering bed, for example a bed including
fluorescent glass beads or fragments such as uranium-glass
elements where permitted by applicable safety and regulatory
requirements, a UV laser, a tunable emitter-side focus lens
on a micromotor servo, a steering mirror on a separate
micromotor servo, a tunable detector-side focus lens on a
separate micromotor servo, and a reactor detector behind a
green-pass filter, all driven by a common controller. The
reactor detector is mounted near the UV laser, approximately
co-aligned with it, or at an angle to the specular axis, as
declared; precise alignment is not required because the
wavelength-shifted fluorescence response is robust to
alignment imprecision. The cascade embodiment~(b) of
Section~\ref{sec:architectural-alternatives} is
illustrated in FIG.~4A by a forward low-latency coupling drawn
as a heavy-stroke arrow from the scene detector to the UV
laser of the reactor subsystem; this forward coupling carries
the scene-side observation directly into the reactor stage and
distinguishes the cascade embodiment from the parallel-subsystem
embodiment~(a) of Section~\ref{sec:architectural-alternatives},
which is obtained by removing the forward
coupling arrow. Reactor detector output is delivered, in both
embodiments, through a declared low-latency conditioning chain
including frozen committed reactor detector gain
$G_{\mathrm{det}}$ to the scene laser amplitude under the
per-sample feedback schedule.

\textbf{FIG.~5} is a schematic flow diagram of the
proof-of-discrepancy protocol showing fleet device
representations~300 training a predictive model~310, a discovery
step~320 comparing observed and predicted responses, method
submission~330, and physical re-execution~340 on a distinct device
with model update.

\textbf{FIG.~6} is a schematic diagram of the 2D
generalisation of the anchor embodiment, comprising a scene
subsystem~210 (eye-safe emitter~24, galvanometer pair~25,
scan lens, scene detector~43) observing scene~50 within
module~100, and a sealed reactor subsystem~220 (reactor
source, reactor galvanometer pair~26, tunable emitter-side
focus lens, scattering element, nonlinear medium~30, tunable
detector-side focus lens, detector-side aperture, reactor
detector~44), driven from a common controller~10 by a shared
scan program, with reactor detector output delivered through a
declared low-latency conditioning chain including frozen
committed reactor detector gain $G_{\mathrm{det}}$ to the
scene emitter amplitude as $A(t)$. The cascade embodiment~(b) of
Section~\ref{sec:architectural-alternatives} is
illustrated in FIG.~6 by a forward low-latency coupling drawn
as a heavy-stroke arrow from scene detector~43 to the reactor
source of subsystem~220; this forward coupling carries the
scene-side observation directly into the reactor stage and
distinguishes the cascade embodiment from the parallel-subsystem
embodiment~(a) of Section~\ref{sec:architectural-alternatives},
which is obtained by removing the forward
coupling arrow. The shared scan program of FIG.~6 is a
synchronisation mechanism distributed by controller~10 to
galvanometer pairs~25 and~26 and is structurally distinct from
the architectural coupling paths (forward low-latency coupling
in embodiment~(b); $A(t)$ feedback in both embodiments). The
illustrated ordering of the detector-side focus lens and
aperture is non-limiting; alternative declared detector-path
orders are described in Section~\ref{sec:anchor-2d-reactor}.

% ======================================================================
\section{Definitions and Glossary}
\label{sec:definitions}
% ======================================================================

This section defines terms used throughout. These definitions are
intended to clarify, not limit, the scope of the claims. Unless
otherwise stated, references to optical or projector--detector
embodiments are illustrative and non-limiting.

As used herein, the terms ``apparatus,'' ``device,'' ``module,''
``assembly,'' and ``system'' refer interchangeably to physical
embodiments of the disclosed subject matter, including any
combination of optical, electronic, mechanical, and computational
components arranged to perform the recited functions. Nothing in
the choice of any such term is intended to imply a structural,
functional, or scope distinction or to limit the breadth of the
claims, except where the surrounding context expressly so indicates
(for example, where a ``module'' is identified by a specific
reference numeral in conjunction with a specific embodiment).
Likewise, references to a particular component (for example, an
``emitter,'' a ``detector,'' a ``reactor,'' or a ``controller'') do
not require that component to be physically distinct from other
components; in some embodiments two or more named components are
realised within a single physical element, and in some embodiments
a single named component is distributed across multiple physical
elements.

\subsection{Core terms}

\begin{definition}[\RK]
A \emph{\RK} refers to a parameterised Markov kernel
\[
  \mathsf{P}_\theta : \mathcal{S} \times \mathcal{U}^{T+1}
  \to \mathcal{P}(\mathcal{C}),
\]
mapping a scene $S \in \mathcal{S}$ and a control protocol $U_{0:T}$
to a probability law over \cba trajectories
$C_{0:T} \in \mathcal{C}$, where $\theta$ denotes configuration
parameters of a physical module and optional learned parameters.
When the distinction matters, $\theta = (\theta_{\mathrm{hw}},
\theta_{\mathrm{sw}})$, where $\theta_{\mathrm{hw}}$ denotes
physical hardware parameters (lens positions, gain settings, reactor
state, mechanical configuration) and $\theta_{\mathrm{sw}}$ denotes
optional learned or digitally set parameters (neural-network weights,
look-up tables, calibration offsets).  In embodiments without learned
parameters, $\theta_{\mathrm{sw}}$ is empty and
$\theta = \theta_{\mathrm{hw}}$.
\end{definition}

\begin{remark}[Physical Authentication Principle]
\label{rem:physical-auth-principle}
The \RK is a physical authentication architecture.  Its
primary authentication act is a declared physical channel event:
known input light enters a reactor-operated system, observed output
light exits, and the measured relation between input, output,
operating conditions, circuit configuration, and coupled scene
constitutes the primary authentication event.  Cryptographic
identifiers, signed transcripts, data-integrity proofs, attestation
evidence, and portable credentials exist only to bind, preserve, and
transport records of that physical event for parties who were not
present.  \textbf{Physical channel first; cryptographic carrier
second.}

Equivalently: in the \RK, cryptography authenticates
records of physical events; it does not replace the physical event
as the primary object.  The default authentication mechanism
throughout the system is light in, light out.  The reactor's
characteristic response to known inputs under declared operating
conditions is what authentication means.  DID-style identifiers,
verifiable credential proofs, and remote-attestation evidence
structures are secondary enabling infrastructure that make the
primary physical act legible to parties not present at the
original event.
\end{remark}

\begin{remark}[Markov kernel terminology]
As used herein, ``Markov kernel'' refers to a regular conditional
probability (stochastic kernel) from the input space
$\mathcal{S} \times \mathcal{U}^{T+1}$ into the measurable space of
\cba trajectories, and does not imply memorylessness of
the raw observations.  Where the reactor exhibits persistence (for
example phosphor afterglow or thermal relaxation), the extended state
$X_t$ is intended to absorb the relevant history so that the kernel
satisfies an \emph{approximate} Markov property over the extended state
space.

In some embodiments, this approximation is formalised via a declared
\emph{finite effective memory horizon}: the assumption that all
embodiments satisfy a mixing condition of the form
$\nu_{\mathrm{mix}}(u; L) \leq \varepsilon_{\mathrm{mem}}$, where
$\nu_{\mathrm{mix}}(u; L)$ is the conditional mutual information
between the bundle at time $t$ and the trajectory before time $t - L$,
given the state at time $t - L$ and the protocol $u$; and $L$ and
$\varepsilon_{\mathrm{mem}}$ are declared embodiment-specific
parameters recorded in the protocol digest.  The memory-depth bound
$\tau_{\mathrm{mem}}$ (for example the lag at which the empirical
autocorrelation of residuals falls below a declared threshold, such as
99\,\% of total autocorrelation mass) is the practical estimator for
$L$; the Markov approximation is understood to hold up to
$\tau_{\mathrm{mem}}$, and a dedicated meter flags bundles in which
residual autocorrelation at lag $\tau_{\mathrm{mem}}$ exceeds the
declared threshold.

\paragraph{Channel-with-memory interpretation (non-limiting).}
In some embodiments, the finite-effective-memory-horizon analysis is
interpreted through the framework of channels with memory and
information-spectrum coding theory (Gallager 1968; Verd\'u--Han 1994),
with the declared mixing condition
$\nu_{\mathrm{mix}}(u;L) \leq \varepsilon_{\mathrm{mem}}$ acting as the
operative ergodicity or mixing assumption for the declared protocol
class. Capacity, mutual-information, and error-exponent quantities are
therefore evaluated over the committed bundle horizon, using
block-structured estimators and confidence bounds that account for the
declared memory depth $\tau_{\mathrm{mem}}$, with the estimator family,
block length, and confidence level recorded in the protocol digest.

No claim of exact finite-dimensional Markovianity is made.  Reactor
media with power-law relaxation (such as some photorefractive crystals
and biological tissue with multi-scale thermal memory) may have no
finite-dimensional Markovian embedding; for such media, the operative
claim is that $\nu_{\mathrm{mix}}(u; L) \leq \varepsilon_{\mathrm{mem}}$
for a declared $L$ and $\varepsilon_{\mathrm{mem}}$, not that the
embedding is exact.  Slow reactor drifts (phosphor burn-in, thermal
history, charge-trap accumulation) that exceed $\tau_{\mathrm{mem}}$
are explicitly acknowledged as potential sources of non-Markovian
behaviour, and their effect on verification and sensing claims is
bounded by the meter envelope.
\end{remark}

\begin{definition}[\RK module]
A \emph{\RK module} is a physical apparatus that realises a
\RK, including at least one emitter, at least one detector,
one or more optical or physical paths coupling emission to measurement,
a controller that executes a control protocol, and optionally
additional components. The term ``\RK'' refers to the induced stochastic mapping, while ``\RK module'' refers to hardware.
\end{definition}

\begin{definition}[Scene tuple]
A \emph{scene} is represented as
\[
  S = \bigl(S^{\mathrm{scene}},\; S^{\mathrm{reactor}}_\theta\bigr),
\]
where $S^{\mathrm{scene}}$ denotes a physical scene being interrogated,
including, without limitation, (i)~an external environment directly
coupled to the module, (ii)~an indirect or proxy scene surface such as
an internal screen or other surface onto which an external environment
is imaged (for example camera-obscura or relay-imaging coupling) and/or
onto which recorded or synthetic imagery is projected or displayed,
(iii)~a continuous abstract or computational process whose state
trajectory constitutes the scene dynamics (as in PoliePuter mode;
physical reality mediates the coupling in all cases), or
(iv)~a null or termination state
$S^{\mathrm{scene}} = S_{\mathrm{term}}$ in which the scene path is
shuttered, beam-dumped, or otherwise bypassed (the degenerate
open-loop sub-case of PoliePuter mode).
$S^{\mathrm{reactor}}_\theta$ denotes an engineered internal
subsystem, when present, that contributes dynamics under configuration
$\theta$.
\end{definition}

\begin{definition}[Control protocol and per-step control]
A \emph{control protocol} is a sequence $U_{0:T}=\{u(t)\}_{t=0}^T$. A
per-step control may be written as
\[
  u(t) = \bigl(\mathbf{s}(t),\; \mathbf{e}(t),\; g_{\mathrm{det}}(t)\bigr),
\]
where $\mathbf{s}(t)$ denotes a scan coordinate or address, $\mathbf{e}(t)$
denotes an emission envelope or pattern parameter (optionally including
wavelength band, polarisation, phase, or other emission settings), and
$g_{\mathrm{det}}(t)$ denotes detector configuration such as exposure,
gain, or gating.
\end{definition}

\begin{definition}[Scan law]
A \emph{scan law} is a mapping $\gamma(t;\eta)$ from time (or discrete
index) to scan coordinates, with parameters $\eta$ (for example
frequencies, amplitudes, phases, envelopes, or basis coefficients). The
scan law may be fixed, selected from a library, or adapted online. Scan
laws include, without limitation, raster, Lissajous, spiral, stepped,
or arbitrary waveforms.
\end{definition}

\begin{definition}[\Cb]
\label{def:convolution-bundle}
For a run under $(S,U_{0:T})$, the \emph{\cb} is the
time-indexed sequence
\[
  C_{0:T} = (c_0,\ldots,c_T),
\]
where each sample $c_t$ is an emission--observation pair
\[
  c_t = \bigl(u(t),\; \mathbf{y}_t\bigr),
\]
and $\mathbf{y}_t$ is an observation vector comprising one or more
measured responses from one or more detector ports (optical, acoustic,
RF, auxiliary sensors, and optional quantum-sensitive observables).
In some embodiments, the observation vector $\mathbf{y}_t$ is a
\emph{multi-tap} record covering not only endpoint emission and
detection but also intermediate states of the cycle that are
physically tappable, including, without limitation, scene-side
emitter command, scene-side detector response, the output of any
forward low-latency coupling between the scene side and reactor
side (per
Section~\ref{sec:architectural-alternatives}), reactor-side emitter
command, reactor-side detector response, and intermediate states
of any conditioning chain on the $A(t)$ feedback path. In further
embodiments where the reactor element exposes intermediate physical
layers --- for example a multi-pass reactor stack, a layered
non-linear medium, or an intermediate phosphor state in a
CRT/phosphor reactor --- the observation vector additionally
includes one or more taps from those intermediate layers, where
physically tappable. The \cb is therefore not
limited to endpoint emission--observation pairs; it is, in some
embodiments, a multi-tap trace of the cycle whose tap set is
declared and committed to the protocol digest.

The name reflects the data structure's characteristic property: at each
time step, what the device emitted and what the device observed are
bundled inseparably---contributions from the device and the scene are
typically difficult to disentangle because the observation $\mathbf{y}_t$
is the physical convolution of the emission $u(t)$ with the scene's
response. The \cb is not a recording of the scene, nor
a recording of the device; it is the record of their joint
interaction through the physical channel.
\end{definition}

\begin{definition}[Reactor latent]
A \emph{reactor latent} $Z_\theta(t)$ is a modelling state adapted to
the filtration generated by $\{c_\tau\}_{\tau \le t}$ and may be
updated by an operator
\[
  Z_\theta(t) = \mathcal{R}_\theta\bigl(Z_\theta(t-1),\;c_t\bigr).
\]
The latent is a modelling or control convenience. A hardware
implementation need not expose $Z_\theta(t)$ as a separate physical
port.
\end{definition}

\begin{definition}[Reactor signature]
A \emph{reactor signature} is a summary of a run, for example a scalar
\[
  r_T = f(C_{0:T}),
\]
where $f$ is a fixed functional applied to all or part of a
\cba trajectory. In some verification embodiments, a
verifier depends on $r_T$ rather than on full $C_{0:T}$.
\end{definition}

\begin{definition}[Media state]
A \emph{media state} $M_t$ is a representation of internal memory media
or persistent state affected by operation (for example persistence,
charge trapping, phase-change, photorefractive response, or other
history-dependent media). The media state may be modelled as part of the
extended state and contributes to the induced channel
$\mathsf{P}_\theta(\cdot\mid S,U_{0:T})$.
\end{definition}

\begin{definition}[Extended state]
An \emph{extended state} may be written as
\[
  X_t = \bigl(S^{\mathrm{scene}}_t,\;
  S^{\mathrm{reactor}}_t,\; M_t,\; q_t\bigr),
\]
where $q_t$ denotes fast internal controller or reactor state, and where
additional slow degrees of freedom may be incorporated into $M_t$ or
into the scene tuple.
\end{definition}

\begin{definition}[Invariance profile]
An \emph{invariance profile} is a pattern of sensitivities and
insensitivities of an observation channel to aspects of a scene,
influenced by configuration knobs including, without limitation,
integration time, spectral band, coherence, numerical aperture,
polarisation, and scan schedule.
\end{definition}

\begin{definition}[Cross-channel statistics]
When multiple channels observe the same scene under distinct invariance
profiles, \emph{cross-channel statistics} include the joint distribution
and dependence structure of channel outputs, including pairwise mutual
informations, conditional distributions, and higher-order dependencies.
\end{definition}

\begin{definition}[\Yd]
The \emph{\yd} $\alpha \in [0,1]$ parameterises the strength of
bidirectional dynamical coupling between the \RK module and
its scene. For example, at $\alpha = 0$, the device operates in static
mode (scene treated as having fixed properties); at $\alpha = 1$, the
device commands maximum coupling depth (the resulting joint dynamics
are empirically assessed for coupling quality via meters; whether the
run satisfies Yoked operation is determined by the CLE/TE co-condition
recorded in the protocol digest, not by $\alpha$ alone). The
\yd is recorded in the protocol digest.
\end{definition}

\begin{definition}[Protocol digest]
A \emph{protocol digest} is a compact record of the control
configuration actually used during a run or window, including, without
limitation, the scan law parameters, emission schedule, detector
settings, gain and alignment values, \yd, regime weights, and
any conditioning signals (for example prompt digests or model
identifiers), and any derived-meter estimation parameters (for example
transfer entropy estimator family, window length, embedding dimension,
and lag selection). In preferred embodiments, the protocol digest is
committed cryptographically (for example by inclusion in the hash-chain
state $\chi_t$ or by binding into a Merkle commitment) so that the
protocol used can be verified after the fact and any discrepancy between
declared and actual protocol is detectable.
\end{definition}

\begin{definition}[Meter vector and meter envelope]
\label{def:meter-envelope}
A \emph{meter vector} $\mathbf{m}(t)$ is a tuple of scalar quantities
computed from the \cb and protocol state at step or window~$t$,
measuring properties such as alignment, gain consistency, signal level,
verisimilitude, empirical hardness, calibration drift, and timing
integrity. A \emph{meter envelope} is a set of admissibility bounds on
the meter vector (for example per-component intervals or a convex
acceptance region) that define the \emph{admissibility regime}---the
region of meter space within which the system is considered to be
operating correctly. (The admissibility regime is distinct from the
\emph{objective regime}---\TB, \LI, or \RT---which specifies the objective family applied to the
\cb.  A single objective regime may have multiple
admissibility regimes, and vice versa.)  A run or window
is considered within its meter envelope when
$\mathbf{m}(t) \in \mathcal{M}_{\mathrm{accept}}$ for all steps in
the window. Meter envelopes are declared in protocol digests and bound
into commitments.

In some embodiments, the meter vector serves as a
\emph{declared-fidelity summary} of the full \cb: it approximates the
informational content of the bundle about specified quantities (regime
classification, hardness index, Yoked declaration) to within a declared
fidelity $\eta_{\mathrm{meter}}$, empirically verified on the anchor
embodiment.  This is not formal statistical sufficiency
($P(\theta \mid \mathbf{m}) = P(\theta \mid C_{0:T})$), which is a
stronger condition that is not asserted; the meter is a
declared-accuracy compression of the bundle, and any claims stated in
terms of the meter vector carry the corresponding fidelity qualification.
\end{definition}

\begin{definition}[Disclosure atom and evidence window]
A \emph{disclosure atom} is the smallest unit of the \cb and
associated metadata that can be independently committed and selectively
opened for audit. In some embodiments, an atom corresponds to a
contiguous time window; in others, to a structured slice (for example a
channel subset or spatial region). An \emph{evidence window} is one or
more disclosure atoms that are grouped for a specific verification,
audit, or cross-attestation purpose. Per-atom digests are aggregated
into Merkle trees so that a verifier can request openings of selected
atoms without requiring disclosure of the full run.
\end{definition}

\begin{definition}[\TB (verification regime)]
\emph{\TB} is the operating regime in which the \RK
is configured to maximise the empirical difficulty of forging or
simulating the \cb, so that the record serves as physically
grounded evidence of a genuine device--scene interaction.  The
optimised quantity is, in some embodiments, Fisher information
$I(\theta)$ about the device's configuration parameters, making the
device's physical identity detectable in the bundle.  For
multi-parameter $\theta$, the optimised quantity may be a scalar
functional of the Fisher information matrix $\mathcal{F}(\theta)$,
such as its trace, determinant, or minimum eigenvalue, or a
directional Fisher information along a task-relevant direction.
Detailed treatment is in Section~\ref{sec:truth-beam}.
\end{definition}

\begin{definition}[\LI (perception and sensing regime)]
\emph{\LI} is the operating regime in which the \RK is
configured to maximise information gain or task-relevant
reconstruction quality from active probing of a scene.  The device
steers its emission protocol to extract scene information efficiently,
subject to meter-envelope constraints.  Detailed treatment is in
Section~\ref{sec:limager}.
\end{definition}

\begin{definition}[\RT (controllable rendering regime)]
\emph{\RT} is the operating regime in which the \RK steers the visual or physical appearance of the scene while
remaining constrained by the physical channel, so that the resulting
\cb records both the intended effect and the scene's actual
response.  Detailed treatment is in
Section~\ref{sec:reality-transform}.
\end{definition}

\begin{definition}[Yoked operation (dynamic coupling modifier)]
\emph{Yoked operation} is a modifier applicable to any of the three
base regimes (\TB, \LI, \RT) in which the
feedback coupling between the device and the scene is tuned so that
both are driven toward a coupled dynamical state.  The \emph{\yd}
$\alpha \in [0,1]$ parameterises coupling strength: $\alpha = 0$ is
static (open-loop), $\alpha = 1$ is the maximum commanded coupling
depth (with the CLE/TE co-condition determining whether a particular
run is classified as Yoked operation).  Yoked
operation is not a fourth regime but a continuous modifier that
enriches any base regime with dynamic coupling structure.  Detailed
treatment is in Section~\ref{sec:yoked}.
\end{definition}

\subsection{Fleet model and proof-of-discrepancy terms}

\begin{definition}[Device configuration representation]
A \emph{device configuration representation} $\mathbf{d}_i$ is a
learned vector or structured object encoding device-specific properties
(manufacturing variability, material composition, calibration state) of
device~$i$, inferred from challenge--response data.  A predictive model
$\mathcal{Q}(X, \mathbf{d}_i)$ accepts the device configuration
representation alongside a challenge configuration $X$ and produces a
predicted response.  See Section~\ref{sec:device-config-rep}.
\end{definition}

\begin{definition}[Proof-of-discrepancy]
A \emph{proof-of-discrepancy} is a reproducible physical method---a
challenge configuration and associated operating conditions---together
with evidence that the predictive model's prediction diverges from
measured physical response, confirmed by independent physical
re-execution on distinct devices.  The protocol by which such proofs are
submitted, verified, and incorporated into the predictive model is
described in Section~\ref{sec:pod-protocol}.
\end{definition}

\begin{definition}[Practical learnability-exclusion condition]
A device satisfies a \emph{practical learnability-exclusion condition}
when no known feasible learning procedure achieves prediction accuracy
above the verification threshold given bounded observations from the
device.  This condition provides the physical basis for the irreducible
residual and the security margin against model-based cloning.  See
Section~\ref{sec:learnability-exclusion}.
\end{definition}

\begin{definition}[PoliePuter mode]
\label{def:poliputer}
\emph{PoliePuter mode} refers to operation of a \RK module
in which the scene partner is a continuous abstract or computational
process.  In the coupled case, the reactor's response at each step
is returned to the computational process as an input before its next
state update, closing a bidirectional loop; the \cba then constitutes
a physically attested record of the joint trajectory.  The degree of
coupling is empirically characterised by the \cba under the declared
Yoked conditions; no coupling class is pre-declared.  In the
open-loop (degenerate) case, the scene path is shuttered or the
return channel is absent ($S^{\mathrm{scene}} = S_{\mathrm{term}}$),
and the reactor operates as a compute substrate without active
feedback.  In all cases, the reactor's nonlinear dynamics, media
states, and feedback loops are exploited for computation, generation,
or optimisation tasks using the \cba logging and meter
infrastructure of the \RK.
\end{definition}

\begin{definition}[PolieBot]
\label{def:poliebot}
A \emph{PolieBot} is a non-limiting example of a mobile or deployable
agent platform that integrates one or more \RK modules with
locomotion, communication, and decision-making subsystems. PolieBots
interact with scenes and with other agents through their \RK
channels and are subject to the verification, evidence-commitment, and
meter-bounded protocols described herein.  Optional multi-agent
coordination extensions are described in the related Filing 2 application.
See Section~\ref{sec:agents}.
\end{definition}

\begin{definition}[Run seed and opening seed]
\label{def:runseed}
A \emph{run seed} $s_{\mathrm{run}}$ is a value derived from the
protocol digest and a pre-committed nonce before a run begins; it
determines the emission schedule, scan law, and any protocol
randomisation for that run, and is logged in the protocol digest.
An \emph{opening seed} $s_{\mathrm{open}}$ is a value supplied by a
verifier or independent beacon \emph{after} the run commitment has
been published; it determines which atoms or windows are to be
selectively disclosed.  In embodiments using the two-seed commitment
structure (Section~\ref{par:two-seed}), the run seed and opening seed
are generated independently, so that the device is not expected to
predict which atoms will be opened at the time the run is executed
under the declared randomness and timing assumptions.
\end{definition}

\subsection{Reactors and ponderable cells}

In this specification, a \emph{reactor} is a subsystem whose dynamics
are dominated by electromagnetic field interactions and that
participates in a \RK. A reactor receives electromagnetic
control inputs, interacts with a scene, and produces electromagnetic
outputs for one or more sensors. Non-limiting examples include optical
elements (diffusers, nonlinear films, scattering plates, spatial light
modulators), CRT phosphor screens, RF and microwave resonators and
waveguides, magnetic coil assemblies driven by swept fields, and
free-space or fibre-coupled optical cavities.

A \emph{Ponderable Cell} is any subsystem whose dominant internal
degrees of freedom are governed by massive-carrier physics---mechanical
displacement, acoustic vibration, fluid flow, thermal diffusion,
chemical concentration, biological metabolism, or hybrid
combinations---rather than by electromagnetic field dynamics. The term
\emph{ponderable} follows classical usage (ponderable matter =~matter
possessing weight and inertia). Ponderable Cells may couple to one or
more reactors via transducers (for example piezoelectric, magnetostrictive,
electro-thermal, or opto-mechanical interfaces) and thereby modulate
behaviour of the associated channel.

The distinction is functional, not ontological: a given physical
component may contain both electromagnetic-dominant and
massive-carrier-dominant degrees of freedom. For example, a
ferrite-loaded coil assembly is classified as a reactor because the
dynamics exploited for \RK operation are electromagnetic
(permeability, hysteresis, field coupling), even though the core
material has rest mass. Both reactors and Ponderable Cells are treated
as parts of the same parameterised channel
$\mathsf{P}_\theta(\cdot\mid S,U_{0:T})$.
The distinction groups electromagnetic-field-dominant embodiments
separately from massive-carrier-dominant embodiments for clarity.

\subsection{Markov-model interpretations (non-limiting)}

In some embodiments, the \RK viewpoint admits several
standard Markov-model lenses, depending on what is treated as state,
what is treated as observation, and whether control is present. In a
non-limiting interpretation, the same underlying physical loop can be
viewed as (i)~a sampler whose long-run statistics approximate a target
distribution (MCMC-style), (ii)~a latent-state time series in which
internal physical state is hidden and the \cb is an emission
(HMM/state-space style), (iii)~a metastable switching process in which
dwell-time and transition statistics are signatures (biophysics-style
channel switching), or (iv)~a controlled process in which emissions are
actions and meter-bounded evidence quality shapes objectives
(MDP/POMDP-style).

\subsection{Adversary models and empirical hardness (non-limiting)}

In some verification embodiments, ``hardness'' is evaluated empirically
relative to specified adversary families and budgets.  Non-limiting adversary
families are parameterised by three dimensions: (a)~\emph{query access
type}: offline CRP collection, online adaptive querying, or physical
access to the device; (b)~\emph{information access}: black-box
(input-output only), grey-box (with knowledge of device architecture
and material composition), or white-box (with full model $\mathcal{Q}$
and representation $\mathbf{d}_i$ access); and (c)~\emph{resource
budget}: computational scale $k$, physical scale $m$, query count $Q$,
and time horizon.  Representative non-limiting instances include:
(i)~a purely digital simulator trained on logged traces (black-box,
offline, bounded compute); (ii)~an adversary with bounded physical
co-illumination capability (physical access, grey-box);
(iii)~an adversary who later obtains control of a previously trusted
device (white-box, physical access); and (iv)~an adversary with
knowledge of device architecture and material composition who builds a
physics-informed surrogate model (grey-box, offline, potentially
unbounded compute).  For each declared embodiment, the relevant
adversary instance and its parameter bounds are recorded in the
protocol digest.  The learnability-exclusion condition
(Section~\ref{sec:learnability-exclusion}) is evaluated separately
against each declared adversary instance, including instance~(iv),
by verifying that prediction accuracy plateaus on held-out challenges
even when the adversary's model is seeded with declared material
composition and geometry parameters.

\subsection{Substrate properties supporting cross-layer coupling and
  governance evidence}

\begin{definition}[Optical anchor]
\label{def:optical-anchor}
An \emph{optical anchor} refers to a cross-layer continuous coupling
mechanism by which two or more coupled \RK layers are
bound to one another in apparatus-level real time, such that an
attacker compromising one layer cannot substitute its contribution
to a recursive cascade without simultaneously satisfying the anchor's
declared latency, modulation, and challenge-response constraints
across all coupled layers. The optical anchor comprises in some
embodiments a declared optical band, a declared modulation or
framing scheme carrying challenge-response content, declared
emitter and detector elements, a declared timing source, declared
allowed path lengths, declared cross-channel isolation, declared
failure modes, and an anomaly-detection response ladder, each
committed to the protocol digest as anchor-layer parameters. The
optical anchor is a security property complementary to per-layer
hardness: per-layer hardness bounds local recoverability of state
from a single layer, while the optical anchor binds layers
continuously to one another. The full apparatus-level disclosure
appears in Section~\ref{sec:optical-anchor-mechanism}.
\end{definition}

\begin{definition}[Optical long-term memory]
\label{def:optical-long-term-memory}
An \emph{optical long-term memory} refers to a storage subsystem in
which committed evidence records are written into a physical optical
substrate configured for durable, individually-retrievable storage,
where the substrate has a write-once-read-many (WORM) physical
property such that, once a record is written, the underlying
physical state of the substrate at the written location does not
support in-place modification or overwriting through the apparatus's
normal write pathway. The WORM property is the apparatus-level
substrate property on which the integrity of durable governance
evidence relies, and is committed to the protocol digest as part
of the storage-layer parameter family. The optical long-term memory
complements channelised optical memory of the kind disclosed in
Section~\ref{sec:crt-phosphor} (which provides controller-managed
short- and medium-term storage) by providing a substrate appropriate
for the durable retention of selected committed records, including
in some embodiments hard-floor violation records, rescue-derived
ghost records, and attestation tokens. The full apparatus-level
disclosure appears in
Section~\ref{sec:optical-long-term-memory}.
\end{definition}

\begin{definition}[Governance-partition isolation]
\label{def:governance-partition-isolation}
\emph{Governance-partition isolation} refers to a substrate-level
property of a \RK apparatus by which the substrate that
produces evaluation meters and other governance-side signals
operates in a partition of the apparatus that is physically and
logically isolated from the agent's behaviour-generation pathways,
such that the agent cannot read, write, or directly influence the
partition's operation through its normal operating pathways.
Governance-partition isolation is a structural property of the
apparatus rather than a runtime monitoring discipline: isolation
forecloses agent influence on the governance partition as a
property of the apparatus's construction, rather than relying on
detection of attempted influence after the fact. The full
apparatus-level disclosure appears at
Section~\ref{sec:governance-partition-isolation}.
\end{definition}

\subsection{Symbols}

\begin{longtable}{@{}p{0.25\textwidth}p{0.70\textwidth}@{}}
    \toprule
    Symbol & Meaning \\
    \midrule
    \endfirsthead
    \toprule
    Symbol & Meaning \\
    \midrule
    \endhead
    \midrule
    \multicolumn{2}{r}{\emph{continued on next page}} \\
    \endfoot
    \bottomrule
    \endlastfoot
    $S$ & Scene tuple
      $S=(S^{\mathrm{scene}},S^{\mathrm{reactor}}_\theta)$ \\
    $U_{0:T}$ & Control protocol $\{u(t)\}_{t=0}^T$ \\
    $u(t)$ & Per-step control, including scan coordinate
      $\mathbf{s}(t)$, emission parameter $\mathbf{e}(t)$, and detector
      setting $g_{\mathrm{det}}(t)$ \\
    $\gamma(t;\eta)$ & Scan law parameterised by $\eta$ \\
    $C_{0:T}$ & \Cb (full run) \\
    $c_t$ & Per-step emission--observation pair
      $(u(t),\mathbf{y}_t)$ \\
    $\mathbf{y}_t$ & Observation vector \\
    $Z_\theta(t)$ & Reactor latent \\
    $r_T$ & Reactor signature summarising a run \\
    $M_t$ & Media or memory state \\
    $q_t$ & Fast internal controller or reactor state \\
    $X_t$ & Extended state \\
    $\mathbf{m}(t)$ & Meter vector \\
    $\alpha$ & \Yd $\in [0,1]$ \\
    $\mathbf{e}(t)$ & Emission envelope or pattern parameter \\
    $\chi_t$ & One-way chain state (cryptographic) \\
    $\mathsf{P}_\theta(\cdot\mid S,U_{0:T})$ & Conditional law
      of the \cb \\
    $\theta = (\theta_{\mathrm{hw}}, \theta_{\mathrm{sw}})$ &
      Configuration parameters: hardware ($\theta_{\mathrm{hw}}$) and
      optional learned ($\theta_{\mathrm{sw}}$) \\
    $\TE_{S \to R}$, $\TE_{R \to S}$ & Transfer entropy from scene
      to reactor and vice versa \\
    $\CLE$ & Conditional Lyapunov exponent \\
    $\sigma_{ij}(t)$ & Trust score from agent $i$ toward $j$
      (see Section on agent layer; distinguished from verifier
      timing fidelity $\sigma^{\mathrm{ver}}_{\tau}$ and from
      conventional statistics-notation $\sigma$ for standard
      deviation) \\
    $\nu_{\mathrm{idemp}},\nu_{\mathrm{contr}},\ldots$ &
      Algebraic deviation metrics (RK algebra) \\
    $K_2 \triangleright K_1$ & Containment (nesting) composition
      of RK modules \\
    $\Pi^{(i)}$ & Protocol digest of agent or instance $i$
      (per-agent digest) \\
    $\Pi_{\mathrm{dig}}$ & Per-run protocol digest: the governing
      protocol-template digest of the current run, committed before
      the capture window opens; binds the run to the declared
      configuration, scan law, exposure schedule, verifier cadence,
      meter envelope, and any selective-opening policy. Distinct
      from $\Pi^{(i)}$, which is a per-agent digest \\
    $K_{\mathrm{spot}}(\cdot,\cdot)$ & Spot point-spread function
      (CRT embodiment only) \\
    $\mathcal{I}_{t_0:t_1}$ & Inner-trace summary (inner
      \cba tuple) \\
    $\lambda_{\mathrm{TE}}$ & Transfer-entropy penalty weight
      (in combined regime/Yoked objectives); units are protocol-digest
      declared to render all objective terms commensurate \\\
    $\lambda_{\CLE}$ & Conditional Lyapunov exponent penalty weight
      (Yoked objective); units are protocol-digest declared to render
      all objective terms commensurate \\
    $\rho_{\mathrm{AR}}$ & Autoregressive coefficient
      (drift/persistence modelling) \\
    $\mu_{\mathrm{mix}}$ & Trust weight mixing rate
      (fleet pairing dynamics) \\
    $r_{\mathrm{rep}}$ & Resource replenishment rate
      (fleet game dynamics) \\
    \midrule
    \multicolumn{2}{l}{\emph{Fleet, PoD, and substrate-theory symbols
      (non-limiting)}} \\
    \midrule
    $\mathcal{Q}(X, \mathbf{d}_i)$ & Predictive model mapping
      challenge $X$ and device configuration representation
      $\mathbf{d}_i$ to predicted response \\
    $\mathbf{d}_i$ & Device configuration representation for
      device $i$ \\
    $\mathbf{d}_0$ & Population-level device configuration
      representation \\
    $e_i(X)$, $\bar{e}_i(X)$ & Per-device residual and
      noise-reduced residual \\
    $\delta_{\mathrm{sys}}(X)$ & Shared model error (fleet-wide
      discrepancy) \\
    $\delta_{\mathrm{dev}}(X,i)$ & Device-specific captured
      variation (device signature) \\
    $\delta_{\mathrm{irr}}(X,i)$ & Irreducible residual (security
      margin) \\
    $\Delta_2$ & Two-simplex of regime weights
      $(\lambda_{\mathrm{TB}},\lambda_{\mathrm{L}},\lambda_{\mathrm{RT}})$ \\
    \multicolumn{2}{p{0.95\textwidth}}{\emph{Convention:
      $\lambda_{\mathrm{subscript}}$ denotes a non-negative scalar
      weight coefficient throughout; distinguished from optical
      wavelength $\lambda$ (unsubscripted, standard optics) by
      subscript presence.}} \\
    $I(\theta)$ & Scalar Fisher information about $\theta$
      (one-dimensional or marginal) \\
    $\mathcal{F}(\theta)$ & Fisher information matrix (used in
      natural gradient methods (ascent or descent depending on
      objective sign); pseudoinverse where singular) \\
    $\mathcal{F}_Q(\theta)$ & Quantum Fisher information
      (Bures metric; see quantum coupling section) \\
    $I(X;Y)$, $I(X;Y \mid Z)$ & Mutual information (conditional
      mutual information) between random variables; distinguished from
      scalar Fisher information $I(\theta)$ by argument structure \\
    $n$ & Non-limiting count of reliably extractable bits per
      response (security parameter context) \\
    $n_{\mathrm{sec}}$ & Effective security parameter
      (scales with transcript information content; see
      cryptographic-primitive substrate) \\
    $k^*_{\mathrm{dig}}(\theta;\varepsilon)$ & Digital hardness index:
      minimal model capacity to drive balanced classification error
      below $\varepsilon$ (see Security Theory) \\
    $m^*_{\mathrm{ana}}(\theta;\varepsilon, r)$ & Analogue hardness index:
      minimal physical emulator scale to drive distinguishability
      advantage below $\varepsilon$ against a verifier/distinguisher
      class of scale $r$ (see Security Theory) \\
    $\mathrm{err}^*_k(\theta)$ & Best achievable balanced
      classification error at model scale $k$ \\
    $\mathrm{Adv}^{\mathrm{ana}}_{m, r}(\theta)$ & Analogue
      distinguishability advantage at emulator scale $m$ against a
      verifier/distinguisher class of scale $r$ \\
    $\varepsilon$ & Hardness threshold (declared; conditions all
      hardness index statements) \\
    $\phi \in \Phi$ & Style parameter: formal input to the \RT objective specifying the target behaviour class
      (see Section~\ref{sec:style-parameter}) \\
    $\varphi(t)$ & Scan-law angular coordinate in helical/cylindrical
      scan embodiments (distinct from RT style parameter $\phi$) \\
    $\Gamma_{\mathrm{addr}}(t;\eta)$ & CRT address-space map from
      timebase and scan parameters $\eta$ to
      internal address coordinate $(\rho,\varphi)$ \\
    $\Sigma_{\mathrm{sum}}$ & Summary map compressing a run record
      $(C_{0:T},\Pi,\mathbf{m})$ to a compact summary vector $\mathbf{s}$ \\
    $\operatorname{agg}$ & Declared and pre-committed aggregation
      functional (e.g.\ median, trimmed mean); used in
      $\delta_{\mathrm{sys}}$ and $T(X^*)$; recorded in protocol digest \\
    $\mathcal{V}$ & Set of participating verifiers in a PoD
      confirmation round \\
    $\tau_{\mathrm{device}}$ & Per-device exceedance threshold for
      optional quorum rule in PoD fleet confirmation \\
    $\CLE_{\mathrm{target}}$ & Target conditional Lyapunov exponent
      for Yoked operation (declared; negative, close to zero) \\
    $\varepsilon_{\mathrm{CLE}}$ & CLE convergence window half-width
      (declared protocol parameter for Yoked convergence criterion) \\
    $\tau_{\TE}$ & TE balance threshold (declared; Yoked convergence
      co-condition) \\
    $\tau_{\mathrm{mag}}$ & Minimum TE magnitude threshold (declared;
      Yoked convergence co-condition; ensures balance criterion cannot
      be trivially satisfied when both TE values are near zero) \\
    $N$ & Number of consecutive convergence windows required to declare
      Yoked operation (declared protocol parameter) \\
    $f_{\mathrm{mod}}$ & Modulation envelope frequency in lock-in /
      sweeping-reactor embodiments (the Arnold tongue axis frequency;
      distinguished from optical carrier frequency, AOD/EOM RF
      frequency, and scan repetition rate). In sweeping-reactor
      embodiments, $f_{\mathrm{mod}}$ denotes the modulation frequency
      of the swept physical parameter as the same envelope quantity
      applied to the reactor parameter sweep \\
    $\Delta f$ & Frequency detuning $f_{\mathrm{mod}} -
      f_{\mathrm{scene}}$ (Arnold tongue horizontal axis) \\
    $\phi_{\mathrm{ref}}(t)$ & Scene-derived reference phase signal
      (lock-in reference for scene-phased full-parameter sweeping;
      see Section~\ref{sec:full-theta-sweep}) \\
    $\delta\theta(t)$ & Parameter modulation perturbation vector
      (scene-phased sweep; small relative to $\theta_0$) \\
    $\delta\theta_k$, $\psi_k$ & Declared amplitude and phase offset
      of the $k$-th parameter component in scene-phased modulation \\
    $S_{\mathrm{IP}}$ & In-phase (real) demodulated lock-in output,
      proportional to $\nabla_\theta T|_{\theta_0}\cdot\delta\theta$ \\
    $J_{\mathrm{train}}(\theta,\pi)$ & Training-phase proxy objective
      for Yoked TE-gradient phase (component of $J_{\mathrm{Yoked}}$) \\
    $T(X^*)$ & Fleet-level PoD confirmation statistic: robust aggregate
      of per-verifier residuals at submitted challenge $X^*$ \\
    $\tau_{\mathrm{confirm}}$ & Pre-registered confirmation threshold
      for the PoD fleet statistic; committed in protocol digest \\
    $\mathcal{K}$ & Index set for composite agent evaluation metrics
      (distinct from meter envelope $\mathcal{M}_{\mathrm{accept}}$) \\
    $\mathsf{tc}$ & Tool-call specification (canonical 5-tuple
      \texttt{(name, args-schema, provider-policy,
      canonicalisation-rule, derivation-rule)} committed to the
      protocol digest before a capture window opens; canonical
      serialization is the UTF-8 byte sequence of the named-field
      JSON encoding with sorted keys; selective opening of
      $\mathsf{tc}$ discloses any subset of its fields under the
      declared opening policy. Distinct from trust score
      $\sigma_{ij}(t)$ and from any conventional statistics-notation
      $\sigma$) \\
\end{longtable}

\paragraph{Declared optical-radiation safety classification
(non-limiting).}
In some embodiments that use lasers or high-intensity optical emitters,
the apparatus includes power monitors, duty-cycle clamps, scan-failure
detection, beam-stop or enclosure controls, service-mode controls, and
hardware or firmware interlocks configured to maintain accessible
emission within a declared optical-radiation safety classification
under the declared operating conditions. In some embodiments, the
declared classification is IEC 60825-1 Class 1, Class 1M, or another
applicable class; in some embodiments, ANSI Z136.1 safe-use controls
are declared in addition to, or in place of, an IEC product-
classification statement. The protocol digest records the applicable
standard and edition, hazard class, wavelength range,
accessible-emission-limit or maximum-permissible-exposure assumptions,
beam divergence, aperture assumptions, exposure time, interlock state,
service-mode exceptions, and any downgrade or shutdown event. The
protocol record is an engineering and audit record; formal regulatory
certification, if any, is performed under the applicable regulatory
process.

% ======================================================================
\section{Detailed Description}
% ======================================================================

Unless otherwise stated or required by context, embodiments described
in this Detailed Description are non-limiting and illustrative.
Alternative configurations, parameter ranges, materials, and
architectures consistent with the definitions herein are within the
scope of the disclosure.

Where this Detailed Description references companion-filing content,
those references are optional and non-essential: no embodiment in
this Detailed Description requires the companion filing for enablement, and
all references to the companion filing may be disregarded without affecting
the enablement of any embodiment described outside a companion filing.


% ----------------------------------------------------------------------
\subsection{Networked Closed-Loop Reactor Witness Mesh}
\label{sec:witness-mesh}

\subsubsection{Overview and primary deployment}
\label{sec:wm-overview}

The principal deployment of the apparatus and methods of this
disclosure is a plurality of \RK modules operating as a
\emph{witness mesh}: multiple modules, each comprising a fast closed
loop in causal feedback with a portion of physical reality observable
to that module, mutually coupled through continuous, low-latency
analogue inter-module signal paths. Single-module embodiments
disclosed elsewhere in this specification --- including the 1D anchor
of Section~\ref{sec:anchor-1d}, the 2D anchor of Section~\ref{sec:anchor-2d}, and the various
reactor variants of Section~\ref{sec:layered-stacks} and subsequent reactor sections ---
are non-limiting building blocks of the witness-mesh architecture and
may also be used as standalone instantiations.

The witness-mesh architecture is the primary intended deployment
because the security and provenance claims of greatest practical value
arise from the joint physical reproduction problem across multiple
modules bound to their respective observations of physical reality,
rather than from any single module's per-device hardness considered
in isolation.

\paragraph{Architectural keystone --- minimum-latency mutual
connection.}
The architectural keystone of the witness mesh is the
\emph{minimum-latency mutual connection} of the constituent modules'
fast loops, bounded by the declared latency budget of
Section~\ref{sec:wm-latency-budget}: each module's fast-loop output
is delivered to its peer modules' coupling-network inputs over a
continuous-analogue, frame-free, buffer-free signal path, with
end-to-end inter-module latency bounded by the propagation delay of
the carrier and the hardware latency of the analogue front ends, and
with no contribution from packet framing, error-correction stages,
queues, buffers, operating-system-mediated transport, or any other
digital re-representation stage. The latency budget so defined is a
load-bearing component of the security primitive of
Section~\ref{sec:wm-security-primitive}: any adversarial intervention
in the inter-module signal path must complete within the same
latency budget, and the absence of intermediate digital
representations removes framing-level insertion targets and raises
the cost of analogue-domain insertion under the declared latency
budget and verifier timing fidelity. The latency budget operates in
conjunction with the witnessing requirement of
Section~\ref{sec:wm-witnessing} and the joint-trace verifier of
Section~\ref{sec:wm-verifier-binding}; the joint security claim
derives from all three components, not from the latency budget alone.

\paragraph{Construction of ``minimum-latency mutual connection''
(non-limiting).}
For purposes of this section, a \emph{minimum-latency mutual
connection} between two modules $M_i$ and $M_j$ means a live
inter-module signal path whose allowed end-to-end latency is the
declared propagation delay of the identified signal path plus
characterised analogue-front-end latency and declared measurement
tolerance, and that does not store a complete verifier-openable
sample window or protocol frame in the live path. The term is a
structural element relative to the declared latency budget of
Section~\ref{sec:wm-latency-budget}; it does not require zero
physical delay or an absolute minimum over all possible routes,
and is not a term of degree relative to the speed of light or any
other physical constant. The declared latency budget, the declared
verifier timing fidelity $\sigma^{\mathrm{ver}}_{\tau}$, and the
declared measurement tolerance are committed to the protocol
digest at each deployment version per
Section~\ref{sec:wm-lifetime}.

\paragraph{Construction of ``verifier-openable sample window''
(non-limiting).}
For purposes of this section, a \emph{verifier-openable sample
window} means the smallest committed atom or atom sequence that
the verifier can select for opening under the two-seed selective-
opening protocol of Section~\ref{par:two-seed}. Storage of a complete such
atom or atom sequence in the live inter-module path constitutes
buffering for purposes of claims reciting the no-buffering
limitation of this section. The atom granularity is committed to
the protocol digest as part of the verifier specification and is,
in some preferred embodiments, at the per-sample level of the
fast-loop sample rate $\{f_i\}$; in other embodiments, the atom
is a contiguous sequence of samples whose duration is committed
to the protocol digest and whose verifier-side timing tolerance
is calibrated against $\sigma^{\mathrm{ver}}_{\tau}$.
Sub-atom-duration analogue storage --- including without
limitation continuous-time analogue delay lines, sample-and-hold
elements within the declared analogue-front-end latency, and
event-driven analogue memory whose retention time is shorter than
the atom duration --- does not, on its own, constitute buffering
for purposes of the no-buffering limitation, provided the
storage is declared and committed to the protocol digest as part
of the inter-module path characterisation.

\subsubsection{Architectural primitive}
\label{sec:wm-architectural-primitive}

A witness mesh comprises a plurality $\{M_i\}_{i\in I}$ of \RK modules, indexed by a finite set $I$ with $|I| \ge 2$, each
module $M_i$ comprising:
\begin{enumerate}[nosep]
  \item an emitter $E_i$ configured to drive a portion of physical
    reality $\Sigma_i$ observable to module $M_i$;
  \item a reactor element comprising a physical medium with a transfer
    function dependent on the medium's microstructure and
    instantaneous state;
  \item a detector $D_i$ configured to record the response of
    $\Sigma_i$ and the reactor to the emission of $E_i$ on a
    continuous or near-continuous timebase, where
    \emph{near-continuous} excludes ordinary frame-buffered
    capture in which a complete verifier-openable sample window or
    protocol frame is accumulated before inter-module coupling, and
    includes sampled analogue or sampled mixed-signal
    implementations whose sample period and conversion latency are
    declared and committed to the protocol digest as part of the
    latency budget of Section~\ref{sec:wm-latency-budget};
  \item a fast closed loop in causal feedback with $\Sigma_i$, in
    which the detector output, after a declared deterministic live
    conditioning chain, drives the emitter modulation at the same
    timebase, with no per-period wait, no frame buffer, and no
    intermediate digital decoding stage in the live coupling path;
    and
  \item one or more inter-module coupling ports configured to
    broadcast a derived signal of the module's fast-loop state to
    one or more peer modules and to receive corresponding signals
    from peer modules, over the inter-module signal path of
    Section~\ref{sec:wm-inter-module-coupling}.
\end{enumerate}

The fast loop of each module $M_i$ operates at a declared sample or
modulation rate $f_i$ and is coupled to its physical reality
$\Sigma_i$ at a declared coupling latency $\tau_i^{\mathrm{loop}}$,
bounded by the propagation delay of light or other carrier through
the optical or signal path of the module and the hardware latency of
the live conditioning chain. The fast loop contains no
runtime-adaptive element and no runtime-learned element: its
behaviour is fully determined by the trained committed parameters of
the deployment version and by the instantaneous physical state of
$\Sigma_i$ and the reactor.

\paragraph{Inter-module coupling --- continuous-analogue signal path
(load-bearing).}
\label{sec:wm-inter-module-coupling}
The inter-module coupling is mediated by a \emph{configured
deterministic inter-module transfer function}, which is a frozen,
deterministic transfer function from peer fast-loop signals to
per-module modulation input, optionally trained, and which is
committed to the protocol digest $\Pi_{\mathrm{dig}}$ as one of
the trainable parameter families of the deployment version. (This
function is referred to elsewhere in this section as the
``trained inter-module coupling network'', the ``inter-module
coupling network'', or simply the ``coupling network''; the terms
are interchangeable in this disclosure and apparatus claims may
recite either form. The ``configured deterministic'' formulation
is preferred for apparatus claims that should not depend on a
training-process limitation.)

\paragraph{Inter-module coupling as versioned parameter family
(non-limiting).}
For purposes of apparatus claims and infringement analysis, the
configured deterministic inter-module transfer function is, in
some preferred embodiments, an element of a \emph{versioned
parameter family}: a deployment version comprises (a) a frozen
parameter family identifier, (b) a frozen modality identifier
selecting one or more of the modulation modalities of
Section~\ref{sec:wm-coupling-modality}, (c) frozen structural
parameters identifying the coupling graph $\mathcal{G}$ and per-
edge propagation/hardware latency declarations, and (d) frozen
trained or hand-designed parameter values realising the
deterministic transfer function from peer fast-loop signals to
per-module modulation input. The frozen parameter family
identifier and the frozen parameter values are committed to the
protocol digest $\Pi_{\mathrm{dig}}$ at each deployment version.
The element is claimable as the versioned parameter family
implementing the configured deterministic transfer function,
without dependence on the training process, on external
declared quantities such as the empirical hardness estimate, or
on whether a corresponding deployment is operational at the
time of infringement analysis. The
inter-module signal path between modules and through the coupling
network is \emph{continuous-analogue} and \emph{frame-free}: the path
contains no analogue-to-digital conversion, no packet framing, no
forward-error-correction stage, no buffer or queue, no digital
protocol decoding, and no operating-system-mediated transport,
between a peer module's fast-loop output and the receiving module's
modulation input mixer. Each module's instantaneous emission is
therefore a direct analogue function of (a) its own
immediately-prior detector output, (b) its peers' immediately-prior
fast-loop outputs as received over the declared inter-module signal
paths, and (c) the frozen committed coupling network, up to a
bounded and declared inter-module latency $\tau_{ij}$.

The continuous-analogue, frame-free property of the inter-module
coupling path is load-bearing for the security primitive of
Section~\ref{sec:wm-security-primitive} as a component of the
combined architectural primitive; the reasons are disclosed in
Section~\ref{sec:wm-latency-budget}.

\paragraph{Coupling network modulation modality (non-limiting).}
\label{sec:wm-coupling-modality}
The trained inter-module coupling network may modulate the receiving
module's input via any one or more of the following modalities, each
of which is a non-limiting embodiment of the inter-module coupling
of this disclosure:

\begin{itemize}[nosep]
  \item \emph{Amplitude modulation}, in which the coupling network
    output drives the amplitude of the receiving module's input
    mixer. Amplitude modulation is a preferred canonical embodiment,
    consistent with the live amplitude drive $A(t)$ of the 2D anchor
    of Section~\ref{sec:anchor-2d} and the per-sample reactor-detector gain
    $G_{\mathrm{det}}$ of the 1D anchor of Section~\ref{sec:anchor-1d} transposed to
    the inter-module case.
  \item \emph{Phase modulation}, in which the coupling network output
    modulates the phase of an optical, RF, or other carrier
    presented at the receiving module's input mixer.
  \item \emph{Frequency modulation}, in which the coupling network
    output modulates the carrier frequency.
  \item \emph{Polarization modulation}, in which the coupling network
    output modulates the polarization state of the carrier.
  \item \emph{Wavelength modulation}, in which the coupling network
    output selects or modulates the carrier wavelength, including
    embodiments using wavelength-division multiplexing of multiple
    peer signals onto a shared inter-module carrier.
  \item \emph{Spatial-mode modulation}, in which the coupling network
    output modulates the spatial mode of the carrier, including
    embodiments using mode-division multiplexing.
  \item \emph{Mixed modality coupling}, combining two or more of the
    above modalities on the same or distinct carriers, with each
    modality contributing distinct physical and timing
    characteristics committed to the protocol digest.
\end{itemize}

The selected modulation modality or modalities, together with the
coupling network's structural parameters, are committed to
$\Pi_{\mathrm{dig}}$ at each declared deployment version. The
witness-mesh primitive is not limited to amplitude modulation; the
choice of modality is an embodiment-level decision driven by the
available physical layer, the carrier modality of the inter-module
signal path, and the verifier's measurement modality.

\paragraph{Representative enabled species per modality
(non-limiting).}
For written-description support of the modulation-modality genus,
representative non-limiting enabled species include:
\begin{itemize}[nosep]
  \item \emph{Amplitude-modulation species:} a transimpedance-
    amplified photodetector at the broadcasting module produces a
    voltage proportional to fast-loop intensity; this voltage
    drives, through the configured deterministic inter-module
    transfer function, the analogue input to a laser-diode bias
    or electro-optic intensity modulator at the receiving module;
    the receiving module's input mixer integrates the resulting
    optical intensity over its declared input bandwidth.
  \item \emph{Phase-modulation species:} the broadcasting module's
    fast-loop output drives an electro-optic phase modulator
    impressing phase modulation on a coherent carrier
    (free-space or fibre-optic) presented at the receiving
    module's coupling-network input mixer; the receiving module
    employs phase-resolving detection (heterodyne, homodyne, or
    interferometric) calibrated to the declared phase reference.
  \item \emph{Frequency-modulation species:} the broadcasting
    module's fast-loop output drives a voltage-controlled
    oscillator or current-driven distributed-feedback laser
    producing frequency modulation of the inter-module carrier;
    the receiving module employs frequency discrimination through
    a declared discriminator filter or wavelength-resolving
    detection.
  \item \emph{Polarization-modulation species:} the broadcasting
    module's fast-loop output drives a polarization controller
    (electro-optic, liquid-crystal, or rotating wave plate)
    impressing polarization modulation on the inter-module
    carrier; the receiving module employs polarization-resolving
    detection through a polarizing beamsplitter and balanced
    photodetectors.
  \item \emph{Wavelength-modulation species:} the broadcasting
    module's fast-loop output controls a tunable laser or
    wavelength-selective filter, producing wavelength modulation
    of the inter-module carrier; the receiving module employs
    wavelength-resolving detection through a wavelength-division-
    multiplexing demultiplexer or grating-coupled detector array.
  \item \emph{Spatial-mode-modulation species:} the broadcasting
    module's fast-loop output drives a spatial light modulator
    or mode-converter array producing spatial-mode modulation of
    the inter-module carrier; the receiving module employs
    mode-resolving detection through a mode-division-
    multiplexing demultiplexer or spatially-resolved photodetector
    array.
\end{itemize}

Each species is a non-limiting enabled embodiment of the
corresponding modulation modality. Combinations and variants
within each species are within the scope of the witness-mesh
primitive provided that the resulting inter-module signal path
satisfies the continuous-analogue, frame-free property of this
section and the latency budget of
Section~\ref{sec:wm-latency-budget}.

\paragraph{Coupling network functional form (non-limiting).}
Independent of the selected modulation modality, the trained
inter-module coupling network's mathematical transfer function may
be implemented as one or more of:
\begin{itemize}[nosep]
  \item a frozen scalar damping coefficient per peer pair;
  \item a frozen vector of per-peer weights applied to each module's
    input mixer;
  \item a frozen matrix or low-rank matrix mapping peer-broadcast
    vectors to each module's modulation input;
  \item a frozen analogue filter with declared transfer function;
  \item a frozen small neural network performing nonlinear weighted
    combination of peer signals, implemented in analogue hardware
    where consistent with the continuous-analogue property of the
    inter-module signal path; or
  \item any other declared deterministic transfer function from peer
    fast-loop outputs to per-module modulation input that is
    implementable within the selected analogue or mixed-signal
    coupling modality of Section~\ref{sec:wm-coupling-modality},
    with declared input/output dimensionality, declared
    bandwidth, declared end-to-end latency within the latency
    budget of Section~\ref{sec:wm-latency-budget}, declared
    parameter values, and declared residual tolerances, all
    committed to the protocol digest, and without storage of a
    complete verifier-openable sample window in the live path.
\end{itemize}
The structure, dimensionality, and parameter values of the coupling
network are committed to $\Pi_{\mathrm{dig}}$ and are public under
the declared threat model. The commitment provides version binding
and anti-malleability of the deployed coupling under the declared
digest scheme; the witness-mesh security primitive does not assume
secrecy of the coupling network's parameters or structure.

\paragraph{Coupling graph.}
The set of inter-module coupling links forms a directed graph
$\mathcal{G} = (I, \mathcal{E}, \tau, w)$, where
$\mathcal{E} \subseteq I \times I$ is the set of directed coupling
edges, $\tau : \mathcal{E} \to \mathbb{R}_{>0}$ assigns to each edge
$(i,j)$ a declared inter-module latency $\tau_{ij}$ corresponding to
the propagation and hardware-latency budget of the analogue path
from $M_i$ to $M_j$, and $w$ parameterises the contribution of edge
$(i,j)$ to the coupling-network input at $M_j$. The graph
$\mathcal{G}$, including its topology and per-edge declared
latencies, is committed to $\Pi_{\mathrm{dig}}$.

\paragraph{Slow-loop and analogue-memory extensions (non-limiting).}
In some embodiments, each module additionally comprises a slow
analogue-memory loop as disclosed in Section~\ref{sec:crt-phosphor} and Section~\ref{sec:dual-loop-appliance};
the slow loop is an optional extension and is not the operative
mechanism of the witness-mesh security primitive. The operative
mechanism is the network of fast loops, each independently closed
with reality, mutually constrained via the continuous-analogue
inter-module coupling network.

\paragraph{Heterogeneous-module embodiments (non-limiting).}
In some embodiments, the witness mesh comprises modules of
heterogeneous construction --- for example, one module implementing
the 1D fluorescent-bed anchor of Section~\ref{sec:anchor-1d}, another the 2D
galvo-based continuous-coupled anchor of Section~\ref{sec:anchor-2d}, a third the
SLM--scattering-medium--camera reactor of Section~\ref{sec:slm-scattering-camera}, and so on
--- coupled through a common inter-module coupling network.
Heterogeneity of module construction contributes to the joint
hardness index by requiring an adversary to model multiple distinct
physical reactor classes simultaneously under the declared meter
envelope.

\subsubsection{Latency budget and the no-buffering requirement}
\label{sec:wm-latency-budget}

The witness-mesh security primitive depends on a declared latency
budget for inter-module coupling that is bounded by the optical-delay
logical clocks of Section~\ref{sec:distributed-time}. The budget is composed of:

\begin{enumerate}[nosep]
  \item the propagation delay
    $\tau^{\mathrm{prop}}_{ij} = d_{ij}/v_{ij}$ of the carrier from
    $M_i$ to $M_j$ along the declared signal path of length $d_{ij}$
    at carrier velocity $v_{ij}$ (for free-space optical, $v_{ij}=c$;
    for fibre, $v_{ij} = c/n$ at the operating wavelength; for RF or
    other carriers, the appropriate carrier velocity is declared);
  \item the hardware latency $\tau^{\mathrm{hw}}_{ij}$ of the
    transmit and receive analogue front ends, including transducers,
    conditioning amplifiers, filters, and the analogue
    coupling-network input stage;
  \item zero contribution from packet framing,
    error-correction stages, or buffer queues, by construction in
    the preferred embodiment.
\end{enumerate}

The total declared inter-module latency
$\tau_{ij} = \tau^{\mathrm{prop}}_{ij} + \tau^{\mathrm{hw}}_{ij}$ is
committed to $\Pi_{\mathrm{dig}}$, and verifier-side checks of
optical-delay logical-clock consistency (Section~\ref{sec:distributed-time}) compare
$\tau_{ij}$ against recorded inter-module signal arrival times under
the declared meter envelope.

\paragraph{Verifier timing fidelity (declared parameter).}
The verifier-side timing fidelity
$\sigma^{\mathrm{ver}}_{\tau}$ at which inter-module signal arrival
times are checked against the declared $\tau_{ij}$ is itself a
declared parameter of the deployment, committed to
$\Pi_{\mathrm{dig}}$ at each declared deployment version. The
witness-mesh security claim of
Section~\ref{sec:wm-security-primitive} is indexed to
$\sigma^{\mathrm{ver}}_{\tau}$: the latency-budget contribution to
the joint hardness scales with the ratio of the adversary's minimum
achievable intervention latency to $\sigma^{\mathrm{ver}}_{\tau}$,
and deployments with coarser $\sigma^{\mathrm{ver}}_{\tau}$ have
correspondingly weaker latency-budget contribution under the declared
adversary class. In some embodiments, $\sigma^{\mathrm{ver}}_{\tau}$
is selected based on (i) the carrier-medium propagation tolerance,
(ii) the available timestamping infrastructure including without
limitation White-Rabbit-class, PTP-HA-class, or higher-fidelity
time-to-digital converter or streak-camera timing, and (iii) the
declared adversary class's minimum achievable intervention latency
for the inter-module signal modality, with margin selected to
support the per-version hardness commitment of
Section~\ref{sec:wm-lifetime}. The relationship between
$\sigma^{\mathrm{ver}}_{\tau}$ and the empirical hardness index is
established at each declared re-evaluation cycle.

\paragraph{Why frame-free continuous-analogue coupling matters
(non-limiting).}
The continuous-analogue, frame-free property serves four distinct
functions in the security primitive.

\textit{First}, it removes framing-level intervention targets. A
framed digital protocol stack exposes a sequence of structured
intermediate representations --- packets, frames, decoded messages
--- each of which constitutes a discrete target for substitution,
replay, or synthesis by an adversary with access to the inter-module
link. A continuous-analogue path with no framing has no such
discrete framing-level targets; the adversary's substitution, if
attempted, must operate in the analogue waveform domain, the
hardness of which is addressed in
Section~\ref{sec:wm-security-primitive}.

\textit{Second}, it raises the cost of waveform-faithful intervention
under the declared latency budget. Any adversarial relay or
substitution of an inter-module signal must itself complete within
the declared $\tau_{ij}$, and any intervention path imposes an
irreducible latency floor consisting of the adversary's own
propagation, capture, processing, and re-transmission stages.
Insertion of a digital decoding, processing, and re-encoding stage
between modules imposes a latency floor of typically tens of
nanoseconds to tens of microseconds depending on hardware, which
is detectable as a violation of the optical-delay logical clock
under verifier-side timing checks only when the added latency,
jitter, and path asymmetry exceed the declared verifier timing
fidelity $\sigma^{\mathrm{ver}}_{\tau}$ and the declared
path-calibration uncertainty. Lower-latency analogue or
mixed-signal intervention paths --- including without limitation
passive optical relay, electro-optic modulation, FPGA-assisted
narrow-feature relay, or analogue waveform substitution at a relay
point --- impose lower latency floors and are correspondingly harder
to detect via the latency budget alone; the witness-mesh primitive
therefore relies on the latency budget in conjunction with the
witnessing requirement of Section~\ref{sec:wm-witnessing} and the
joint-trace verifier of Section~\ref{sec:wm-verifier-binding}, as
disclosed in Section~\ref{sec:wm-security-primitive}, rather than
on the latency budget alone.

\textit{Third}, it preserves the closed-with-reality property
across the network. In the preferred embodiment, each module's fast
loop is fully analogue from detector to reactor to emitter, and the
inter-module coupling extends this fully-analogue signal path
across the network without an intermediate digital re-representation.
The joint dynamics of the coupled mesh therefore evolve as a
continuous physical system, and the verifier's joint-trace check
(Section~\ref{sec:wm-verifier-binding}) evaluates a continuous joint
trajectory rather than a sequence of discrete frame-level decisions.

\textit{Fourth}, it eliminates buffering-induced ambiguity in
causal attribution. A buffered or framed inter-module path admits
queue-occupancy variations, retransmission, and protocol-stack
scheduling jitter, all of which loosen the binding between a sender
module's instantaneous fast-loop state and a receiver module's
instantaneous coupling input. The continuous-analogue path binds
the two with a single declared $\tau_{ij}$ and a verifier-side
fidelity $\sigma^{\mathrm{ver}}_{\tau}$, making the optical-delay
logical clocks of Section~\ref{sec:distributed-time} a tight rather than loose
constraint.

\paragraph{Synchronisation of fast-loop sample rates (non-limiting).}
In some embodiments, the fast-loop sample or modulation rates
$\{f_i\}$ of the modules are synchronised to a common physical
reference --- including without limitation a shared optical clock,
an injection-locked frequency reference, or a phase-locked loop
driven by a distributed timing signal --- so that peer signal
arrivals at each module's coupling-network input are aligned to the
receiving module's fast-loop sampling phase up to a declared phase
tolerance. The shared reference, the synchronisation mechanism, and
the phase tolerance are committed to $\Pi_{\mathrm{dig}}$. In some
embodiments, fast-loop rates are heterogeneous across modules with
declared rate ratios and inter-rate phase relationships committed
to the protocol digest.

\paragraph{Declared exceptions and degraded modes (non-limiting).}
In some embodiments, a portion of the inter-module coupling path is
implemented digitally for engineering reasons. In such embodiments,
the digital-stage latency is itself a declared, characterised, and
committed quantity included in $\tau^{\mathrm{hw}}_{ij}$, and the
verifier-side timing check is calibrated accordingly. Such
embodiments are degraded modes of the primary deployment in the
sense that the hardness contribution from the no-buffering property
is reduced; the security claim of
Section~\ref{sec:wm-security-primitive} is correspondingly indexed
to the declared latency budget, the analogue fraction of the
inter-module path, and the verifier timing fidelity
$\sigma^{\mathrm{ver}}_{\tau}$. The preferred embodiment is the
fully-analogue, frame-free inter-module coupling.

\paragraph{Continuous-analogue carve-out for claims (non-limiting).}
Claims reciting an uninterrupted continuous-analogue inter-module
coupling path apply only to embodiments omitting any digital stage
in the live path, including without limitation the framing,
forward-error-correction, buffer/queue, digital-decoding, and
operating-system-mediated transport stages disclosed as excluded
in Section~\ref{sec:wm-inter-module-coupling}. Embodiments
containing a declared digital stage in the live path remain within
broader architectural mesh claims of the witness-mesh primitive
that do not require uninterrupted continuous-analogue coupling,
but do not satisfy claims that recite the uninterrupted
continuous-analogue limitation.

\paragraph{Latency-budget contribution against analogue and mixed-
signal relay attacks (non-limiting).}
The latency contribution of the continuous-analogue, frame-free
property does not by itself exclude passive optical relay,
electro-optic relay, FPGA-assisted analogue or narrow-feature
relay, mixed-signal relay, or analogue waveform substitution
attacks at a relay point on the inter-module path. Such attacks
are evaluated under the committed verifier timing fidelity
$\sigma^{\mathrm{ver}}_{\tau}$, the declared physical path model
(including signal-path length $d_{ij}$, carrier velocity
$v_{ij}$, and analogue-front-end transfer characteristics), and
the verifier coverage tuple of
Section~\ref{sec:wm-verifier-coverage}. The witness-mesh
primitive's empirical hardness against analogue and mixed-signal
relay attacks is a declared, time-indexed, attacker-indexed
quantity per Section~\ref{sec:wm-lifetime}, established at each
declared re-evaluation cycle against the strongest available
analogue-domain attack class, and is committed to the protocol
digest as part of the per-version hardness commitment.

\subsubsection{Security primitive}
\label{sec:wm-security-primitive}

The witness-mesh security primitive is the joint consistency of the
network's \cb: each module's contribution to the joint record
$\{C^{(i)}_{0:T}\}_{i\in I}$ must simultaneously be consistent with:
\begin{enumerate}[label=(\alph*),nosep]
  \item that module's actual observation of its portion of physical
    reality $\Sigma_i$ during the evidence window $[0,T]$;
  \item that module's reception of every peer module's fast-loop
    broadcast at the declared inter-module latencies $\{\tau_{ij}\}$
    over the continuous-analogue inter-module signal paths of
    Section~\ref{sec:wm-latency-budget}, evaluated to verifier-side
    fidelity $\sigma^{\mathrm{ver}}_{\tau}$; and
  \item the committed trained inter-module coupling network
    mediating those receptions, as committed to
    $\Pi_{\mathrm{dig}}$.
\end{enumerate}

To produce a forged joint record that the verifier of
Section~\ref{sec:wm-verifier-binding} accepts, an adversary must
reproduce all three constraints simultaneously across all
networked modules, under the verifier's declared coverage parameters
of Section~\ref{sec:wm-verifier-binding}.

\paragraph{Adversary's joint reproduction problem.}
Specifically, the adversary must:

(i) Reproduce each module's fast-loop response to the actual
physical reality observable from that module's vantage point at the
relevant times during $[0,T]$, which requires either physical
possession of the module and physical access to its vantage point
during the evidence window, or reconstruction of the physical
reality at that vantage from external sources of equivalent
fidelity under the verifier's declared meter envelope.

(ii) Reproduce each module's reception of every peer module's
fast-loop broadcast at the propagation-delay-bounded times disclosed
by the optical-delay logical clocks of Section~\ref{sec:distributed-time} and the
declared latency budget of Section~\ref{sec:wm-latency-budget},
within $\sigma^{\mathrm{ver}}_{\tau}$, which requires knowing what
each peer module broadcast --- and therefore requires reproducing
the peer modules' responses under constraints (i) and (iii) applied
transitively to those peers.

(iii) Reproduce the network's coupled dynamics under the specific
committed inter-module coupling network, which is frozen and
committed to the protocol digest. Although the coupling network's
parameters are public under the declared threat model, reproduction
requires evaluation of the coupling network on the actual peer
fast-loop signals each module would receive in genuine operation.
The adversary cannot freely choose alternative peer signals to ease
reproduction, because the peer signals are themselves constrained
by (i) and (ii) applied to the peer modules.

\paragraph{Mutually-recursive structure of the forgery problem.}
The constraints (i)--(iii) are mutually recursive across the
network: forging module $M_A$'s record requires the records of all
peers connected to $M_A$ over the coupling graph $\mathcal{G}$,
each of which requires the records of their peers under the same
constraints, transitively across the connected component of
$\mathcal{G}$ containing $M_A$. The adversary's problem is therefore
the joint reproduction of the entire connected component's coupled
behaviour under the committed coupling, given the actual physical
reality each constituent module observed during the evidence window
and given the declared latency budget on inter-module coupling.

This mutually-recursive structure does not constitute a formal
proof of joint reproduction hardness in the sense of cryptographic
soundness; consistent with Section~\ref{par:proof-of-projection-empirical}, the witness-mesh primitive is an empirical hardness construction, indexed to a
declared adversary class and a declared meter envelope. The
structural contribution of this section is to articulate why the
joint reproduction problem is \emph{capable of being} empirically
harder, and in preferred deployments is measured as harder, than
any per-module reproduction problem under the declared adversary
model and verifier coverage: the joint problem is not separable
into per-module subproblems under the continuous-analogue
inter-module coupling, the declared latency budget evaluated at
$\sigma^{\mathrm{ver}}_{\tau}$, and the verifier-side joint-trace
check at the declared coverage parameters of
Section~\ref{sec:wm-verifier-binding}. The realised hardness of
any specific deployment is established empirically at the
declared re-evaluation cycle and is committed per-version under
Section~\ref{sec:wm-lifetime}.

\paragraph{Latency-budget alone is insufficient (non-limiting).}
The latency budget of Section~\ref{sec:wm-latency-budget} is a
necessary but not sufficient component of the witness-mesh security
primitive. An adversary with access to the inter-module signal path
may, in principle, intervene in the analogue domain within the
declared latency budget, by characterising the inter-module
waveform from observation and injecting a learned analogue
perturbation through an electro-optic, acousto-optic, or equivalent
modulator at a relay point on the inter-module path, or at the
receiving module's coupling-network input. Such analogue-domain
intervention is not excluded by the continuous-analogue, frame-free
property alone. The witness-mesh primitive does not rely on the
latency budget alone to exclude adversarial intervention. Rather,
the primitive derives its empirical hardness from the
\emph{conjunction} of:
\begin{enumerate}[label=(\alph*),nosep]
  \item the latency budget on inter-module coupling, evaluated at
    verifier-side fidelity $\sigma^{\mathrm{ver}}_{\tau}$
    (Section~\ref{sec:wm-latency-budget});
  \item the witnessing requirement of
    Section~\ref{sec:wm-witnessing}, which constrains each module's
    fast-loop response to that module's observation of physical
    reality during the evidence window; and
  \item the joint-trace verifier of
    Section~\ref{sec:wm-verifier-binding}, which checks per-edge
    directed causal consistency at sample-level granularity under
    the declared coverage parameters and the declared meter
    envelope.
\end{enumerate}
An analogue-domain adversarial perturbation that satisfies the
latency budget must additionally produce, at each module, a
fast-loop response consistent with that module's observation of its
physical reality and consistent with the committed coupling network
operating on the genuine peer signals; the empirical hardness of
the joint reproduction problem of this section is the joint
hardness of constraints (a)--(c) under the declared adversary
class. The continuous-analogue, frame-free property contributes to
(a) by removing framing-level intervention targets and by raising
the cost of waveform-faithful intervention under the declared
latency budget; the witnessing requirement (b) and joint-trace
verifier (c) contribute the remaining structural barriers to
forgery.

\subsubsection{Witnessing requirement}
\label{sec:wm-witnessing}

The witness-mesh security primitive incorporates a \emph{witnessing
requirement}: each module's fast-loop response is bound to its
observation of physical reality $\Sigma_i$ during the evidence
window, so that the joint record is rejected by the verifier of
Section~\ref{sec:wm-verifier-binding} to the extent that opened
verifier meters depend on residual scene variables not
reconstructable by an adversary, under the adversary's declared
surveillance budget and the verifier's declared meter envelope.
The witnessing contribution to joint hardness is, accordingly, a
deployment-indexed quantity bounded by per-vantage residual
unobserved entropy as set out in the following paragraphs, and
not an absolute property of the architecture.

\paragraph{Witnessing contribution and per-vantage residual entropy.}
The witnessing contribution to the joint hardness of the security
primitive of Section~\ref{sec:wm-security-primitive} is bounded
above by the per-vantage residual unobserved entropy of $\Sigma_i$
under the adversary's surveillance budget within the declared
threat model. In some embodiments, $\Sigma_i$ comprises a sealed,
occluded, hidden, or otherwise non-externally-observable region of
physical space, in which case the per-vantage residual unobserved
entropy is large and the witnessing contribution is correspondingly
large; in some embodiments, $\Sigma_i$ comprises a publicly
observable scene under known geometry, in which case the
per-vantage residual unobserved entropy may be small under a
sufficiently capable adversary surveillance budget, and the
witnessing contribution is correspondingly reduced. The
witness-mesh primitive's overall hardness profile is therefore
deployment-specific and is established at each declared
re-evaluation cycle and committed to the protocol digest as part
of the per-version hardness commitment of
Section~\ref{sec:wm-lifetime}.

\paragraph{Vantage-point independence and amplification.}
Vantage-point independence between modules can amplify the
witnessing contribution by adding constraints that, under the
declared adversary surveillance budget and the verifier's
declared meter envelope, are not necessarily reconstructable
from a single vantage. Vantage-point independence
does not by itself constitute the witnessing contribution: even a
witness mesh in which all modules observe overlapping or coincident
regions of $\Sigma$ retains a per-module witnessing contribution to
the extent that each module's individual physical-state evolution
during the evidence window is not externally reconstructable to
the verifier's declared meter envelope. Conversely, a witness mesh
deployed entirely in publicly observable scenes may derive most of
its joint hardness from vantage-point independence, the latency
budget, and the joint coupled dynamics, with reduced per-module
witnessing contribution. Both regimes are within the scope of the
witness-mesh primitive; the relative contributions of witnessing,
vantage-point independence, latency budget, and coupled dynamics
to the joint hardness are deployment-specific and are committed
per-version under Section~\ref{sec:wm-lifetime}.

\paragraph{Per-module reactor microstructure (non-limiting).}
In some preferred embodiments, each module's reactor incorporates
microstructure-derived per-device transfer-function variability of
reactor-microstructure-unclonable as disclosed in Section~\ref{sec:layered-stacks}, contributing an additional
per-module hardness component independent of and complementary to
the witnessing contribution above. Reactor-microstructure hardness
and witnessing-derived hardness are independent contributions to
the empirical hardness profile, with the relative weights
deployment-specific and committed per-version under
Section~\ref{sec:wm-lifetime}. An adversary with clones of every
module but without scene access at the verifier's declared meter
envelope is rejected to the extent that opened meters depend on
residual scene variables not reconstructable within the verifier's
amplitude and phase residual tolerances; an adversary with full
scene access but without per-module reactor characterisation is
rejected to the extent that opened meters depend on reactor
transfer features not reproducible from the adversary's reactor
model under the verifier's tolerances; an adversary lacking both
is rejected to the extent that opened meters depend on either
class of variables. The relative weights are deployment-specific,
empirically established, and committed per-version.

\paragraph{Effect of dynamic scene reconstruction on
vantage-point independence (non-limiting).}
Modern dynamic scene-reconstruction methodologies, including
without limitation neural radiance fields, four-dimensional
Gaussian splatting, sparse-view generative completion, and
multi-modal scene priors, can reduce the effective vantage-point
independence between modules in publicly observable scenes by
inferring or hallucinating unseen views from sparse external
observations. In deployments where the witness mesh observes
publicly observable scenes under known geometry, the per-vantage
residual unobserved entropy is correspondingly reduced, and the
witnessing contribution to joint hardness is correspondingly
reduced. This effect is the structural reason that the
witness-mesh primitive's overall hardness profile is
deployment-specific: deployments with sealed, occluded, or
otherwise non-externally-reconstructable per-vantage observations
retain a large witnessing contribution, while deployments in
publicly reconstructable scenes derive most of their joint
hardness from other components of the security primitive
(latency budget, joint coupled dynamics, verifier coverage,
per-module reactor microstructure where applicable).

\paragraph{Residual-entropy evaluation procedure (non-limiting).}
The per-vantage residual unobserved entropy of $\Sigma_i$ under
the declared adversary surveillance budget is, in some preferred
embodiments, estimated by the following procedure for each
module $M_i$ at each declared re-evaluation cycle:
\begin{enumerate}[nosep]
  \item enumerate the adversary observation channels available
    under the declared adversary class, including without
    limitation external surveillance modalities, ambient
    sensors, scene priors, and any leaked legitimate joint-record
    windows;
  \item define the meter variables of $\Sigma_i$ that are
    checked by the verifier through per-module reality binding
    of Section~\ref{sec:wm-verifier-binding}, including
    representative non-limiting meter variables such as
    optical amplitude residuals at the per-module detector,
    phase residuals at the per-module detector where
    applicable, local illumination measurements, scene-pose
    variables, temporal derivatives of the per-module
    detector output, scene occlusion state, and where
    applicable acoustic, RF, thermal, or other side-meter
    variables associated with the module's vantage of
    $\Sigma_i$;
  \item fit the best declared reconstruction model under the
    adversary's observation budget that predicts the meter
    variables from the adversary's available observations,
    using the strongest adversary-model class committed to the
    protocol digest;
  \item measure the conditional prediction error of the fitted
    model against the verifier's amplitude and phase residual
    tolerances $\epsilon_{\mathrm{amp}}$ and
    $\epsilon_{\mathrm{phase}}$ on a held-out evaluation
    corpus;
  \item commit the resulting residual-entropy estimate, with a
    declared confidence interval and the corpus identifier, as
    part of the per-version hardness commitment of
    Section~\ref{sec:wm-lifetime}.
\end{enumerate}
The procedure provides an operational method for a person of
ordinary skill to evaluate the witnessing contribution to the
joint hardness in a specific deployment. Variants of the
procedure --- for example, using multiple reconstruction-model
classes, employing adaptive adversary models, or stratifying
across time windows --- are non-limiting and may be selected
based on the declared adversary class and meter envelope.

\subsubsection{Verifier binding}
\label{sec:wm-verifier-binding}

The witness-mesh security primitive presupposes a verifier that
checks the joint \cb record at sample-level granularity under the
two-seed selective-opening protocol of Section~\ref{par:two-seed}. The
verifier operates under selective opening rather than real-time
streaming: each module's full sample-level \cb is committed under
a Merkle root $\mathrm{com}^{(i)}_{\mathrm{full}}$ at the close of
the evidence window, and the verifier subsequently requests
selective opening of randomised atoms or windows under the
two-seed protocol; the verifier does not require continuous receipt
or processing of the full joint sample stream during the evidence
window. The sample-level granularity refers to the resolution at
which selective openings, when requested, are evaluated --- not to
a real-time streaming requirement.

The verifier performs:

\begin{enumerate}[nosep]
  \item per-edge directed causal consistency: for each coupling
    edge $(i,j) \in \mathcal{E}$ and each opened atom or window,
    verification that the receiving module $M_j$'s coupling input
    at sample time $t$ is consistent with the broadcasting module
    $M_i$'s fast-loop output at sample time $t - \tau_{ij}$ under
    the declared latency budget and verifier timing fidelity
    $\sigma^{\mathrm{ver}}_{\tau}$ of
    Section~\ref{sec:wm-latency-budget};
  \item joint-trace consistency: for each opened window,
    verification that the joint sample sequence
    $\{(u^{(i)}(t), \mathbf{y}^{(i)}_t)\}_{i\in I, t\in[0,T]}$ is
    consistent with the committed coupling network operating on
    the recorded peer signals under the declared meter envelope;
  \item per-module reality binding: verification that each module's
    response $\mathbf{y}^{(i)}_t$ is consistent with that module's
    independently-recorded observation of its physical reality
    $\Sigma_i$ at time $t$, via the meter envelope and declared
    verisimilitude meters of the security-theory section
    (Section~\ref{sec:security-theory}); and
  \item full-record commitment: verification of the Merkle
    commitment $\mathrm{com}^{(i)}_{\mathrm{full}}$ of each
    module's full sample-level \cb under the protocol of
    Section~\ref{sec:two-seed-protocol}, with selective opening of randomised atoms
    under the two-seed protocol.
\end{enumerate}

\paragraph{Verifier coverage parameters (declared parameters).}
\label{sec:wm-verifier-coverage}
The verifier's coverage of the joint record is characterised by a
tuple of declared parameters committed to $\Pi_{\mathrm{dig}}$ at
each deployment version:

\begin{itemize}[nosep]
  \item $\rho_{\mathrm{open}} \in (0,1]$, the selective-opening
    probability or fraction of committed atoms or windows opened
    by the verifier under the two-seed protocol per evidence
    window;
  \item $\kappa_{\mathrm{edge}} \in (0,1]$, the edge-check coverage
    or fraction of coupling edges $(i,j) \in \mathcal{E}$ checked
    for directed causal consistency per opened window, where the
    coverage is stratified over \emph{directed} edges of
    $\mathcal{E}$ rather than aggregated over undirected pairs or
    aggregate edge count, so that for each directed coupling
    $i \to j$ the verifier's per-edge causal-consistency check is
    a distinct coverage event;
  \item $\delta_{\mathrm{dec}} \in [1, \infty)$, the decimation
    factor relating the verifier's checked sample rate to each
    module's fast-loop sample rate $f_i$, with
    $\delta_{\mathrm{dec}} = 1$ corresponding to checking every
    sample;
  \item $\epsilon_{\mathrm{amp}}$, the verifier's amplitude
    residual tolerance for joint-trace consistency under the
    declared meter envelope;
  \item $\epsilon_{\mathrm{phase}}$, the verifier's phase residual
    tolerance for joint-trace consistency under the declared meter
    envelope, where applicable to the inter-module coupling
    modality.
\end{itemize}

The witness-mesh security claim of
Section~\ref{sec:wm-security-primitive} is indexed to the verifier
coverage tuple
$(\rho_{\mathrm{open}}, \kappa_{\mathrm{edge}},
\delta_{\mathrm{dec}}, \epsilon_{\mathrm{amp}},
\epsilon_{\mathrm{phase}})$ together with
$\sigma^{\mathrm{ver}}_{\tau}$ from
Section~\ref{sec:wm-latency-budget}: the joint hardness profile is
committed per-version under
Section~\ref{sec:wm-lifetime} as a function of these declared
parameters under the declared adversary class. Deployments with
reduced verifier coverage --- for example, lower
$\rho_{\mathrm{open}}$, lower $\kappa_{\mathrm{edge}}$, larger
$\delta_{\mathrm{dec}}$, or looser $\epsilon_{\mathrm{amp}}$ or
$\epsilon_{\mathrm{phase}}$ --- have correspondingly reduced
joint hardness, and the relationship between the coverage tuple
and the empirical hardness index is established at each declared
re-evaluation cycle.

\paragraph{Detection-probability relation (non-limiting,
first-order).}
For purposes of relating the verifier coverage tuple to the
detection of a forged joint record, the verifier's detection probability against a forgery in which a fraction $q_{\mathrm{bad}}$ of \emph{committed} directed edge-atoms would exhibit a residual exceeding the declared tolerances
$\epsilon_{\mathrm{amp}}$ and $\epsilon_{\mathrm{phase}}$ admits
the non-limiting first-order analytical bound
\[
P_{\mathrm{detect}} \ge 1 - \bigl(1 - q_{\mathrm{bad}}\,
\rho_{\mathrm{open}}\,\kappa_{\mathrm{edge}}\bigr)^{m},
\]
where $m$ is the number of \emph{committed} directed edge-atoms in the evidence window over which independent opening and edge-check trials are applied, $\rho_{\mathrm{open}}$ is the per-atom opening probability, and $\kappa_{\mathrm{edge}}$ is the per-opened-window edge-check probability. The bound
assumes independence of opened atoms and uniformity of
$q_{\mathrm{bad}}$ across the opened atoms; both assumptions are
imperfect in practice (stratification creates dependencies, and
$q_{\mathrm{bad}}$ varies across atoms) and the bound is
accordingly a first-order analytical anchor rather than an
exact relation. Refinements --- including stratified-sampling
corrections, edge-specific $q_{\mathrm{bad}}^{(i,j)}$
weightings, and decimation-aliased detection terms --- are
within the scope of this disclosure as non-limiting
embodiments and are committed, where applicable, to the
protocol digest as part of the per-version hardness commitment
of Section~\ref{sec:wm-lifetime}.

\paragraph{Opening-seed mechanics and stratification (non-limiting).}
In some preferred embodiments, the verifier's coverage parameters
are realised under the following non-limiting mechanics:
\begin{enumerate}[nosep]
  \item \emph{Opening seed.} Opened atoms or windows are selected
    after full-record commitment using a verifier-generated or
    jointly-generated unpredictable opening seed
    $s_{\mathrm{open}}$ derived from a declared unpredictability
    source independent of the run seed of the two-seed protocol of
    Section~\ref{par:two-seed}. The opening seed is not disclosed to the
    deployment prior to commitment.
  \item \emph{Stratification.} Selection of opened atoms is
    stratified over (a) time within the evidence window, (b)
    modules in $I$, and (c) coupling edges in $\mathcal{E}$, so
    that no module, edge, or temporal sub-window is
    systematically excluded or systematically over-represented
    across selective openings. The stratification scheme is
    committed to $\Pi_{\mathrm{dig}}$.
  \item \emph{Decimation aliasing constraint.} The decimation
    factor $\delta_{\mathrm{dec}}$ is selected so that aliasing
    against the declared fast-loop sample rates $\{f_i\}$ and any
    declared modulation frequencies in the inter-module coupling
    is excluded or bounded by a declared aliasing margin
    committed to $\Pi_{\mathrm{dig}}$.
  \item \emph{Tolerance calibration.} The amplitude and phase
    residual tolerances $\epsilon_{\mathrm{amp}}$ and
    $\epsilon_{\mathrm{phase}}$ are calibrated against the
    declared meter envelope and the empirical noise floor of the
    measurement chain, with the calibration procedure and
    measurement chain committed to $\Pi_{\mathrm{dig}}$.
  \item \emph{Failure accounting.} The commitment protocol does
    not permit selective discard of windows after capture; all
    captured evidence-window records are committed under
    $\mathrm{com}^{(i)}_{\mathrm{full}}$, and any declared
    exclusions (for example for hardware-fault windows under a
    declared health-monitor protocol) are themselves committed.
\end{enumerate}

\paragraph{Adversarial-discriminator-ensemble check (optional,
non-limiting).}
In some embodiments, the verifier additionally performs an
adversarial-discriminator check, in which one or more learned
discriminator models trained against declared benchmark forgery
model families available at the evaluation date are evaluated
against the joint record; the discriminator ensemble's output is
incorporated into the meter envelope as a declared evaluation
meter under the meter-partition discipline of the security-theory
section. The discriminator
ensemble is not a standalone acceptance criterion of the
witness-mesh primitive, and any false-acceptance rate is
version-specific and committed only for the declared model
families, training corpus, verifier tuple, and adversary
budget. The discriminator ensemble has no structural security
guarantee outside its committed receiver-operating-characteristic
evaluation, and is not a substitute for the per-edge directed
causal consistency, joint-trace consistency, per-module reality
binding, or full-record commitment checks of the verifier
specification of this section. In some embodiments, the
discriminator ensemble's training data, held-out adversary
families, false-acceptance threshold, and confidence interval
are committed to $\Pi_{\mathrm{dig}}$ as part of the per-version
hardness commitment of Section~\ref{sec:wm-lifetime}.

\paragraph{Reduced-verifier embodiments and non-conforming
verifiers (non-limiting).}
Reduction of the verifier to per-module summary statistics,
pairwise correlations, low-dimensional aggregates, coarse
geometric consistency checks, non-stratified opening, biased
opening-seed generation, decimation that aliases against the
declared fast-loop rates, or tolerances loose enough to admit
generative-surrogate forgeries yields a weaker primitive whose
security claim is correspondingly reduced and is not the
security claim of this section. A verifier that does not check
the per-edge directed causal consistency, joint-trace
consistency, per-module reality binding, and full-record
commitment of this section under the declared coverage tuple is
\emph{non-conforming} with respect to claims of this disclosure
that recite the joint-trace verifier with declared coverage
tuple. Such a non-conforming verifier may satisfy broader
architectural claims of the witness-mesh primitive that do not
recite the joint-trace verifier element, but does not satisfy
claims that recite the joint-trace verifier limitation. The
full-joint-record sample-level verifier under selective opening
with the declared coverage tuple is the preferred embodiment for
the witness-mesh security claim. In some embodiments,
intermediate-fidelity verifiers are employed in resource-
constrained deployments, with the meter envelope and declared
adversary class adjusted to the verifier's actual coverage of
the joint record, and the resulting hardness profile committed
to the protocol digest at each declared re-evaluation cycle.

\subsubsection{Distinction from prior art}
\label{sec:wm-prior-art}

The witness-mesh security primitive is distinct from the following
classes of prior art.

\paragraph{Single-device optical physically unclonable functions.}
The witness-mesh primitive is distinct from single-device optical
physically unclonable function (PUF) constructions, including
those cited in the background art (Pappu et al., ``Physical
One-Way Functions,'' \emph{Science}, 2002; Goorden et al.,
``Quantum-secure authentication of a physical unclonable key,''
\emph{Optica}~1, 421--424, 2014; Hui et al., ``Non-linear optical
scattering PUF,'' \emph{Optics Express}~31(24), 40646--40657,
2023; Davis, Letz, Mosk, and Pinkse, EP3252740B1; and related
disclosures), in the following respects:

\begin{enumerate}[label=(\alph*),nosep]
  \item Optical PUF security rests on the empirical hardness of
    cloning a single device's challenge--response mapping. The
    witness-mesh primitive does not reduce to single-device
    clonability: the joint reproduction problem of
    Section~\ref{sec:wm-security-primitive} binds at every
    module's observation of reality and across all peer
    couplings, and the joint hardness contribution is not
    separable into per-module components under the declared meter
    envelope.
  \item Optical PUFs are characterised by challenge--response
    pairs over a fixed challenge space; some optical PUFs are
    high-dimensional and based on mesoscopic scattering or
    nonlinear interactions. The witness-mesh primitive is
    characterised by the joint \cb record across the network over
    the evidence window, which is a fundamentally different object
    from a single-device challenge--response mapping: it requires
    simultaneous multi-vantage reality witnessing under a coupled
    live verifier and is not reducible to a finite set of
    per-device CRPs.
  \item Published optical-PUF attack methodologies, including
    transmission-matrix reconstruction, machine-learning surrogate
    modelling, and neural-network reconstruction of speckle
    responses, target single-device transfer functions. These
    methodologies apply to the witness-mesh primitive only insofar
    as the verifier accepts records reproducible from per-module
    surrogate models composed under the public coupling network
    --- a regime addressed by the joint-trace verifier of
    Section~\ref{sec:wm-verifier-binding} and its declared
    coverage tuple.
  \item The witness-mesh primitive incorporates the witnessing
    requirement of Section~\ref{sec:wm-witnessing}, absent from
    optical PUFs: the security claim binds each module's response
    to its observation of physical reality during the evidence
    window, with the witnessing contribution bounded by the
    per-vantage residual unobserved entropy under the adversary's
    surveillance budget.
\end{enumerate}

In some embodiments, the witness-mesh primitive is combined with
optical reactor-microstructure-class single-device hardness, with each module's
reactor incorporating reactor-microstructure unclonability as disclosed in
Section~\ref{sec:layered-stacks} and related sections. In such combined embodiments,
the joint security claim incorporates both the witness-mesh
constraint and the per-module clonability constraint, with the
contributions independent and complementary.

\paragraph{Integrated photonic interferometric PUFs and on-chip
MZI-network PUFs.}
The witness-mesh primitive is distinct from integrated photonic
interferometric PUFs and on-chip Mach--Zehnder interferometer (MZI)
network PUFs, including without limitation Smith and Jacinto,
``Reconfigurable Integrated Optical Interferometer Network-Based
Physically Unclonable Function,'' \emph{J.\ Lightwave
Technol.}~38(17), 4599--4606, 2020 (arXiv:2012.11358); Jacinto
and Smith, ``Utilizing a Fully Optical and Reconfigurable PUF as a
Quantum Authentication Mechanism,'' \emph{OSA Continuum}~4(2),
739--756, 2021 (arXiv:2012.10326); Tarik et al., ``Realization of
Robust Optical Physical Unclonable Function Using a Silicon
Photonic Quasicrystal Interferometer''; Mesaritakis et al.,
``Physical unclonable function based on a multi-mode optical
waveguide,'' \emph{Sci.\ Rep.}~8(1), 9653, 2018; and the broader
reconfigurable-optical-PUF literature including silicon-photonic
moir\'e-quasicrystal interferometers, electrically- or
thermally-reconfigurable scattering PUFs, and phase-change-material-
based reconfigurable photonic PUFs. Such constructions describe a
single integrated photonic device --- in the case of the
MZI-network PUFs, a network of Mach--Zehnder interferometers
fabricated on a single photonic integrated circuit with
computer-controlled phase modulators --- used as a reconfigurable
challenge--response authentication primitive. The ``network'' in
such constructions is the on-chip optical interconnect within a
single device, not a network of physically separate modules. The
witness-mesh primitive is distinguished from these constructions
by:

\begin{enumerate}[label=(\alph*),nosep]
  \item the architectural primitive: the witness mesh comprises a
    plurality of physically separate \RK modules, each
    with its own emitter, reactor, detector, and fast closed loop
    in causal feedback with its own portion of physical reality,
    coupled via inter-module signal paths between
    separately-housed modules; the cited integrated photonic PUFs
    are single devices with no inter-module signal paths;
  \item the closed-with-reality property: each witness-mesh module
    operates a fast closed loop in causal feedback with physical
    reality $\Sigma_i$ external to the module during the evidence
    window, with the reactor in feedback with the loop dynamics;
    the cited integrated photonic PUFs operate as open-loop
    challenge--response devices that do not close a feedback loop
    with external physical reality during authentication;
  \item the witnessing requirement of
    Section~\ref{sec:wm-witnessing}, which has no counterpart in
    the cited integrated photonic PUF literature: those
    constructions authenticate the device or messages from the
    device, not observations of physical reality external to the
    device;
  \item the inter-module continuous-analogue coupling under the
    trained inter-module coupling network of
    Section~\ref{sec:wm-inter-module-coupling}, with declared
    inter-module latencies $\tau_{ij}$ and verifier timing
    fidelity $\sigma^{\mathrm{ver}}_{\tau}$, which has no
    counterpart in the cited single-device constructions; and
  \item the joint \cb record across multiple modules under the
    verifier of Section~\ref{sec:wm-verifier-binding} with
    declared coverage tuple
    $(\rho_{\mathrm{open}}, \kappa_{\mathrm{edge}},
    \delta_{\mathrm{dec}}, \epsilon_{\mathrm{amp}},
    \epsilon_{\mathrm{phase}})$, as opposed to challenge--response
    pairs over a single device.
\end{enumerate}

In some embodiments, an individual witness-mesh module's reactor
may be implemented as an integrated photonic device, including a
reconfigurable interferometric structure or a quasicrystal
interferometer of the kind described in the above-cited
literature; the witness-mesh primitive is the multi-module
networked construction operating on top of such per-module
reactors, not the per-module reactor itself.

\paragraph{Distance-bounding and physical-layer relay-attack
literature.}
The witness-mesh primitive is distinct from distance-bounding
protocols and physical-layer relay-attack mitigations, including
without limitation Brands and Chaum, ``Distance-Bounding
Protocols,'' \emph{Eurocrypt~1993}, LNCS~765, 344--359, 1994;
Hancke and Kuhn, ``An RFID Distance Bounding Protocol,''
\emph{SecureComm~2005}, 67--73, 2005; Clulow, Hancke, Kuhn, and
Moore, ``So Near and Yet So Far: Distance-Bounding Attacks in
Wireless Networks,'' \emph{ESAS~2006}, LNCS~4357, 83--97, 2006;
and successor distance-bounding and physical-layer
relay-resistance work. Such protocols use round-trip-time
measurements between a prover and verifier, with bit-by-bit timed
challenges, to upper-bound the prover's physical distance from the
verifier and to detect relay attacks; published physical-layer
attacks demonstrate that distance-bounding security depends on the
specific receiver, demodulation, and channel-design details. The
witness-mesh primitive shares the architectural element of
latency-bounded physical-layer security with this literature, but
is distinguished by:

\begin{enumerate}[label=(\alph*),nosep]
  \item the security objective: distance-bounding authenticates
    the prover's physical distance from a single verifier through
    a single round-trip-time measurement; the witness-mesh
    primitive authenticates the joint physical evolution of
    multiple coupled witnesses bound to their individual
    observations of physical reality, across an entire evidence
    window;
  \item the temporal scope: distance-bounding operates over
    discrete bit-level challenge--response rounds; the
    witness-mesh primitive operates over a continuous-analogue
    joint trace of duration $T$ verified at sample-level
    granularity under the declared coverage tuple of
    Section~\ref{sec:wm-verifier-coverage};
  \item the witnessing requirement of
    Section~\ref{sec:wm-witnessing}, which has no counterpart in
    distance-bounding: distance-bounding does not bind the
    prover's response to any observation of physical reality
    external to the prover and verifier; the witness-mesh
    primitive binds each module's response to that module's
    observation of $\Sigma_i$;
  \item the multi-witness joint-record structure: distance-bounding
    typically involves a single prover and single verifier (with
    multi-prover and multi-verifier extensions described as
    add-ons); the witness-mesh primitive is intrinsically
    multi-module with mutually-recursive joint reproduction
    constraints across the coupling graph $\mathcal{G}$.
\end{enumerate}

The physical-layer relay-attack methodologies analysed in the
distance-bounding literature --- including without limitation
late-commit attacks, overclocking, demodulation-stage attacks,
deferred-decision attacks, and analogue waveform manipulation ---
constitute attack methodologies against which the witness-mesh
primitive's latency-budget contribution must be empirically
evaluated. The structural defences against this attack class are
the conjunction of the latency budget of
Section~\ref{sec:wm-latency-budget}, the witnessing requirement of
Section~\ref{sec:wm-witnessing}, and the joint-trace verifier of
Section~\ref{sec:wm-verifier-binding}, as disclosed in the
latency-budget-alone-is-insufficient paragraph of
Section~\ref{sec:wm-security-primitive}.

\paragraph{Coupled-reservoir-computing and delay-coupled
optoelectronic networks; link-inference attacks as a primary
attack class.}
The witness-mesh primitive is distinct from coupled-reservoir-
computing and delay-coupled optoelectronic-network literature,
including without limitation Brunner and Larger and successor
photonic reservoir-computing systems; Tanaka et al., ``Recent
advances in physical reservoir computing: A review,'' \emph{Neural
Networks}~115, 100--123, 2019; Banerjee et al., ``Machine Learning
Link Inference of Noisy Delay-coupled Networks with
Opto-Electronic Experimental Tests,'' \emph{Phys.\ Rev.\ X}~11,
031014, 2021 (arXiv:2010.15289); Banerjee, Pathak, Roy, Restrepo,
and Ott, ``Using Machine Learning to Assess Short Term Causal
Dependence and Infer Network Links,'' \emph{Chaos}~29, 121104,
2019; chaos-synchronisation key-distribution work using coupled
nonlinear units; and related disclosures. Such systems describe
coupled nonlinear physical dynamical systems with learnable
inter-node delays, and demonstrate machine-learning-based
reservoir-computing surrogates and topology-and-delay inference
from observed nodal time series.

The link-inference and surrogate-emulation methodologies of
Banerjee et al.\ and successor work constitute a primary class of
attacks against which the witness-mesh primitive must be
empirically evaluated. In particular, an adversary in possession
of leaked joint-record windows may, under the public coupling
graph $\mathcal{G}$ and the public trained inter-module coupling
network, train a reservoir-computing or message-passing surrogate
to reproduce the joint nodal-state time series, and may infer the
declared inter-module delays $\{\tau_{ij}\}$ from the same data.
The witness-mesh primitive's structural defences against this
attack class are:

\begin{enumerate}[label=(\alph*),nosep]
  \item the witnessing requirement of
    Section~\ref{sec:wm-witnessing}: a surrogate that reproduces
    joint dynamics statistically under the inferred topology must
    additionally produce, at each module, a response consistent
    with that module's actual observation of physical reality
    during the evidence window;
  \item the joint-trace verifier of
    Section~\ref{sec:wm-verifier-binding}: per-edge directed
    causal consistency at sample-level granularity, evaluated
    under selective opening with the declared coverage tuple, is
    evaluated at a fidelity that may exceed, and in committed
    versions is compared against, the fidelity achieved by
    declared surrogate-emulation attack classes including the
    cited literature; and
  \item the periodic re-evaluation and reconfiguration framework
    of Section~\ref{sec:wm-lifetime}: empirical hardness against
    the surrogate-emulation attack class is established at each
    declared re-evaluation cycle and committed to the protocol
    digest, and coupling-network rotation invalidates accumulated
    surrogate characterisation at the surface level.
\end{enumerate}

The witness-mesh primitive's empirical hardness against the
link-inference and surrogate-emulation attack class is therefore a
declared, time-indexed, attacker-indexed quantity per
Section~\ref{sec:hardness-index} and Section~\ref{sec:wm-lifetime}, not a structural
guarantee. Architectural elements shared with the cited
reservoir-computing literature --- delay-coupled inter-module
signal exchange, trained coupling parameters, nonlinear nodal
dynamics --- are repurposed in the witness-mesh primitive for a
security objective rather than a computational objective, and are
combined with the witnessing requirement, the protocol-digest
commitment, the joint-trace verifier with declared coverage tuple,
and the per-version hardness commitment, none of which is present
in the cited literature.

\paragraph{Continuous-analogue chaos-synchronisation and
analogue-coupled secure-communications literature.}
The witness-mesh primitive is distinct from continuous-analogue
chaos-synchronisation and analogue-coupled secure-communications
literature, including without limitation Argyris et al.,
``Chaos-based communications at high bit rates using commercial
fibre-optic links,'' \emph{Nature}~438, 343--346, 2005; Soriano,
Garc\'ia-Ojalvo, Mirasso, and Fischer, ``Complex photonics:
Dynamics and applications of delay-coupled semiconductor
lasers,'' \emph{Reviews of Modern Physics}~85, 421--470, 2013;
Stankovski, McClintock, and Stefanovska, ``Coupling Functions
Enable Secure Communications,'' \emph{Phys.\ Rev.\ X}~4, 011026,
2014; Nadzinski, Dobrevski, Anderson, McClintock, Stefanovska,
Stankovski, and Stankovski, ``Experimental Realization of the
Coupling Function Secure Communications Protocol and Analysis
of Its Noise Robustness,'' \emph{IEEE Trans.\ Inf.\ Forensics
Secur.}~13(10), 2591--2601, 2018; Illing, Panda, and Shareshian,
``Isochronal chaos synchronization of delay-coupled
optoelectronic oscillators,'' \emph{Phys.\ Rev.\ E}~84, 016213,
2011; Zhou and Roy, ``Isochronal synchrony and bidirectional
communication with delay-coupled nonlinear oscillators''; and
related disclosures of chaos-based and coupling-function-based
secure communications using continuous-analogue inter-device
coupling.

Such systems describe two or more nonlinear dynamical systems
(typically semiconductor lasers, optoelectronic oscillators, or
related photonic devices) coupled through continuous-analogue
signal paths with declared propagation delays, with secure
communication of an information signal between the coupled
systems by chaos synchronisation, coupling-function-based
ciphers, or related dynamical-system primitives. Information
security in the cited literature derives from one or more of:
the receiver's knowledge of a private coupling function or
shared synchronisation key, the difficulty of cloning the
nonlinear dynamics from observed time series, or the practical
inaccessibility of the inter-device signal path to an
eavesdropper.

The witness-mesh primitive shares architectural elements with
this literature --- including continuous-analogue inter-device
coupling, declared propagation delays, and nonlinear dynamics
--- but is distinguished by:

\begin{enumerate}[label=(\alph*),nosep]
  \item the security objective: the cited literature uses
    coupled analogue dynamics to enable \emph{private
    communication of an information signal} between two or more
    parties, with security derived from coupling-function
    secrecy or synchronisation-key secrecy. The witness-mesh
    primitive uses coupled analogue dynamics to enable
    \emph{multi-witness joint attestation of physical-reality
    observations} to a third-party verifier, with security
    derived from the joint reproduction problem of
    Section~\ref{sec:wm-security-primitive} under a public
    coupling commitment;
  \item public coupling commitment vs.\ private coupling
    parameters: the cited literature relies, in various
    instances, on one or more of private parameters,
    synchronisation assumptions and pre-shared synchronisation
    keys, parameter mismatch between transmitter and adversary,
    chaotic masking, or the practical difficulty of modelling
    the coupled dynamics from observed signals --- with the
    coupling parameters or synchronisation parameters not
    disclosed to adversaries. The witness-mesh primitive treats
    the trained inter-module coupling network and the coupling
    graph $\mathcal{G}$ as \emph{public commitments} to the
    protocol digest under the declared threat model, with
    security derived from joint physical reproduction under
    public coupling rather than from coupling-parameter secrecy
    or any of the foregoing;
  \item multi-witness joint-record structure: the cited
    literature is typically point-to-point (one transmitter,
    one receiver) or small symmetric mesh, with information
    flow from sender to receiver. The witness-mesh primitive is
    intrinsically multi-witness with the joint \cb record
    across the network as the security primitive, mutually-
    recursive across the connected component of $\mathcal{G}$
    per Section~\ref{sec:wm-security-primitive};
  \item the witnessing requirement of
    Section~\ref{sec:wm-witnessing}, which has no counterpart in
    chaos-synchronisation secure communications: the cited
    literature transmits an information signal generated by the
    sender, not observations of physical reality external to
    the coupled-systems network. The witness-mesh primitive
    binds each module's response to that module's observation of
    physical reality $\Sigma_i$ during the evidence window;
  \item the closed-with-reality property: each witness-mesh
    module operates a fast closed loop in causal feedback with
    physical reality external to the module during the
    evidence window. The cited literature's coupled nonlinear
    systems close their loops with each other (and optionally
    with self-feedback), not with external physical reality;
  \item the joint-trace verifier with selective opening: the
    cited literature uses time-evolving Bayesian inference,
    chaos-synchronisation decoding, or related
    receiver-side decoding methods to recover the transmitted
    signal. The witness-mesh primitive uses the verifier of
    Section~\ref{sec:wm-verifier-binding} with declared
    coverage tuple under the two-seed selective-opening
    protocol of Section~\ref{par:two-seed} to evaluate the joint
    \cb record, which has no counterpart in the cited
    literature;
  \item per-version hardness commitment under
    Section~\ref{sec:wm-lifetime}: the cited literature does
    not commit a per-version empirical hardness estimate
    against declared adversary classes to a protocol digest as
    part of the deployment record.
\end{enumerate}

The chaos-synchronisation literature's attack methodologies ---
including chaos-synchronisation key-extraction attacks,
generalised-synchronisation attacks, parameter-estimation
attacks, and time-series-based reconstruction --- are
methodologies against which the witness-mesh primitive's
empirical hardness must be evaluated where the inter-module
coupling exhibits chaotic or near-chaotic dynamics. The
witness-mesh primitive's structural defences against this
attack class are the joint-trace verifier with declared
coverage tuple of Section~\ref{sec:wm-verifier-binding}, the
witnessing requirement of Section~\ref{sec:wm-witnessing}, and
the per-version hardness commitment of
Section~\ref{sec:wm-lifetime}. Architectural elements shared
with the cited literature --- continuous-analogue inter-device
coupling, declared propagation delays, nonlinear coupled
dynamics --- are repurposed in the witness-mesh primitive for
multi-witness joint attestation under public coupling
commitment, rather than for private communication between
parties under coupling secrecy.

\paragraph{Virtual proofs of reality and physical-fact
attestation.}
The witness-mesh primitive is distinct from virtual proofs of
reality and related physical-fact attestation primitives,
including without limitation R\"uhrmair, Martinez-Hurtado, Xu,
Kraeh, Hilgers, Kononchuk, Finley, and Burleson, ``Virtual
Proofs of Reality and their Physical Implementation,'' in
\emph{2015 IEEE Symposium on Security and Privacy}, pp.~70--85
(DOI: 10.1109/SP.2015.12), and successor work on virtual proofs
of physical statements. The cited literature describes
protocols by which a prover establishes a physical statement
about its local environment --- including without limitation
``a certain object in the prover's system has temperature
$X^\circ$C,'' ``two certain objects in the prover's system are
positioned at distance $X$,'' or ``a certain object in the
prover's system has been irreversibly altered or destroyed''
--- to a remote verifier over a digital communication channel,
with security derived from physically unclonable per-device
behaviours of disordered scattering media, temperature-sensitive
integrated circuits, quantum systems, or related physical
primitives, and without reliance on classical secret keys or
tamper-resistant trusted sensor hardware.

The witness-mesh primitive shares with this literature the
objective of establishing physical statements to a verifier
without relying on a classical secret-key tamper-resistance
trust assumption, but is distinguished by:
\begin{enumerate}[label=(\alph*),nosep]
  \item single-prover-single-verifier vs.\ multi-witness joint
    attestation: the cited literature is structured as a
    single-prover-single-verifier protocol in which one
    physical statement is established by one prover (using one
    or more sensors or reactor-microstructure-unclonability components co-located on the prover) to one
    verifier. The witness-mesh primitive establishes a joint
    physical statement spanning a plurality of physically
    separate modules whose joint \cb record is the security
    primitive, with security derived from the joint reproduction
    problem of Section~\ref{sec:wm-security-primitive} across
    the connected component of $\mathcal{G}$;
  \item continuous-analogue inter-module coupling: the cited
    literature establishes individual physical statements
    through challenge-response interactions between the
    verifier and individual physical primitives at the prover.
    The witness-mesh primitive establishes its joint statement
    through continuous-analogue inter-module coupling among
    the witness modules during the evidence window, in which
    each module's fast-loop response is determined in part by
    peer-module modulation of the inter-module coupling
    network, with no analogous mutual coupling among physical
    primitives in the cited single-prover virtual-proofs-of-
    reality literature;
  \item the joint-trace verifier with declared coverage tuple
    of Section~\ref{sec:wm-verifier-binding} under the
    two-seed selective-opening protocol of Section~\ref{par:two-seed},
    which has no counterpart in virtual-proofs-of-reality
    challenge-response protocols; and
  \item the per-version empirical hardness commitment of
    Section~\ref{sec:wm-lifetime}, including the cross-version
    transfer-attack evaluation, which is not present in the
    cited virtual-proofs-of-reality literature.
\end{enumerate}

\paragraph{Sensor physical unclonable functions.}
The witness-mesh primitive is also distinct from sensor
physical unclonable functions (sensor PUFs) and related
sensor-binding primitives, including without limitation
Rosenfeld, Gavas, and Karri, ``Sensor Physical Unclonable
Functions,'' in \emph{2010 IEEE International Symposium on
Hardware-Oriented Security and Trust (HOST)}, pp.~112--117,
and successor sensor-PUF literature. Sensor PUFs extend the
PUF challenge-response paradigm to sensing devices by binding
the sensed physical value to a per-device unclonable
fingerprint, providing authentication, unclonability, and
verification of the sensed value at a single sensor. The
witness-mesh primitive's distinguishing structural elements
relative to sensor PUFs are:
\begin{itemize}[nosep]
  \item the multi-module joint-record structure: sensor PUFs
    bind a sensed value to a single physical sensor's
    fingerprint; the witness-mesh primitive binds the joint
    record of a plurality of mutually-coupled witness modules
    to the joint physical reproduction problem of
    Section~\ref{sec:wm-security-primitive};
  \item the inter-module coupling network and joint-trace
    verifier, neither of which is present in single-sensor
    microstructure-based authentication; and
  \item the public coupling commitment, in which the trained
    inter-module coupling network and the coupling graph
    $\mathcal{G}$ are public commitments to the protocol
    digest, vs.\ the typically-private per-device fingerprint
    of a sensor PUF.
\end{itemize}
Where a witness-mesh module's reactor incorporates per-device
microstructure variability of reactor-microstructure-unclonable as disclosed in
Section~\ref{sec:layered-stacks}, the resulting per-module reactor microstructure
is one component of the witness-mesh joint hardness profile,
complementary to but distinct from the joint-record structure
of the witness-mesh primitive disclosed in this section.

\paragraph{Cyber-physical-systems attestation and dynamic
watermarking.}
The witness-mesh primitive is also distinct from cyber-
physical-systems attestation and dynamic-watermarking
primitives, including without limitation Valente, Barreto, and
C\'ardenas, ``Cyber-Physical Systems Attestation,'' (2014),
which proposes that a verifier introduce false control signals
and compare the resulting plant dynamics against a model to
attest the operational correctness of sensors and controllers;
Valente and C\'ardenas, ``Remote Proofs of Video Freshness for
Public Spaces,'' in \emph{Proceedings of the 3rd ACM Workshop
on Cyber-Physical Systems Security \& Privacy (CPS-SPC '17)},
pp.~111--122, which uses physically-manifested visual
challenges (such as QR codes or social-media-derived feeds in
the camera's field of view) to attest video-feed freshness for
surveillance cameras monitoring public spaces; Hespanhol,
Porter, Vasudevan, and Aswani, ``Dynamic Watermarking for
General LTI Systems,'' in \emph{56th IEEE Conference on
Decision and Control (CDC)}, December 2017, pp.~1834--1839
(arXiv:1703.07760), which detects malicious sensor attacks on
multi-input multi-output linear-time-invariant systems with
partial state observations through asymptotic and statistical
tests on injected private control signals; Satchidanandan and
Kumar, ``Dynamic Watermarking: Active Defense of Networked
Cyber-Physical Systems,'' \emph{Proceedings of the IEEE},
105(2):219--240, 2017; Porter, Hespanhol, Aswani, and
co-authors, ``Detecting Generalized Replay Attacks via
Time-Varying Dynamic Watermarking'' (2020) and successor
dynamic-watermarking literature.

The cited literature establishes attestation or attack-detection
guarantees for control systems through one or more of:
verifier-injected challenge signals into a plant or sensor
loop; statistical tests on resulting plant dynamics against a
declared plant model; or physically-manifested challenges
(such as visual challenges in a camera's field of view) whose
expected response is verifiable from the prover's recorded
output. The witness-mesh primitive shares with this literature
the use of declared adversary models, statistical detection
tests, and external challenges that manifest physically in the
prover's environment, but is distinguished by:
\begin{enumerate}[label=(\alph*),nosep]
  \item single-prover vs.\ multi-witness structure: the cited
    literature establishes attestation for a single plant,
    single sensor stream, or single camera. The witness-mesh
    primitive establishes joint attestation across a plurality
    of mutually-coupled witness modules with a joint \cb record
    as the security primitive;
  \item inter-module mutual coupling: the cited literature
    does not couple multiple sensors or plants through a
    continuous-analogue inter-module coupling network during
    capture. The witness-mesh primitive's joint physical
    evolution under continuous-analogue inter-module coupling
    is not present in CPS attestation, video-freshness
    attestation, or dynamic-watermarking primitives;
  \item public coupling commitment vs.\ private watermark or
    challenge: the cited dynamic-watermarking literature
    relies on a private excitation injected by the verifier
    that is not disclosed to the adversary, and the cited
    video-freshness literature relies on visual challenges
    whose expected manifestation is bound to the prover by
    visibility in the camera's field of view. The witness-mesh
    primitive's trained inter-module coupling network and
    coupling graph are public commitments to the protocol
    digest, with security derived from joint physical
    reproduction under public coupling rather than from
    challenge or watermark secrecy; and
  \item the joint-trace verifier with declared coverage tuple
    and the per-version empirical hardness commitment of
    Section~\ref{sec:wm-lifetime}, neither of which is present
    in the cited literature.
\end{enumerate}

The dynamic-watermarking literature is, in addition, the
strongest cited art for attacks against the linear embodiment
disclosed in Section~\ref{sec:wm-preferred} (linear-case
reduction to multi-input multi-output system identification)
and is referenced in that section accordingly. Where the
witness-mesh primitive is deployed with linear per-module
dynamics and linear coupling, the joint hardness contribution
from coupled dynamics is, as disclosed in
Section~\ref{sec:wm-preferred}, evaluated against the cited
system-identification and dynamic-watermarking literature; the
preferred embodiments are non-affine, history-dependent, or
high-dimensional, in which the cited literature's linear-case
methods are insufficient to recover the joint dynamics.

\paragraph{Secure localisation, verifiable multilateration, and
hidden base-station systems.}
The witness-mesh primitive is also distinct from secure
localisation and verifiable-multilateration primitives,
including without limitation \v{C}apkun and Hubaux, ``Secure
Positioning of Wireless Devices with Application to Sensor
Networks,'' in \emph{Proceedings of IEEE INFOCOM 2005},
vol.~3, pp.~1917--1928; \v{C}apkun and Hubaux, ``Secure
Positioning in Wireless Networks,'' \emph{IEEE Journal on
Selected Areas in Communications}, 24(2):221--232, February
2006; \v{C}apkun, \v{C}agalj, and Srivastava, ``Securing
Localization with Hidden and Mobile Base Stations,'' in
\emph{Proceedings of IEEE INFOCOM 2006}; \v{C}apkun, Rasmussen,
\v{C}agalj, and Srivastava, ``Secure Location Verification with
Hidden and Mobile Base Stations,'' \emph{IEEE Transactions on
Mobile Computing}, 7(4):470--483, 2008; \v{C}apkun, Butty\'an,
and Hubaux, ``SECTOR: Secure Tracking of Node Encounters in
Multi-Hop Wireless Networks''; the distance-hijacking and
multi-prover relay-attack literature; and successor secure-
localisation work. The cited literature establishes secure
position verification or secure proof-of-location of one or
more provers through verifiable multilateration --- in which
multiple verifiers measure time-of-flight or signal-strength
to a single prover and use distance-bounding to constrain the
prover's position --- or through hidden or mobile base
stations whose positions are not disclosed to adversaries.

The witness-mesh primitive shares with this literature the
notion that physical-layer timing constraints among multiple
spatially-distributed parties can establish security-relevant
properties of a prover, but is distinguished from this
literature on grounds that include and extend the v4 distance-
bounding distinctions:
\begin{enumerate}[label=(\alph*),nosep]
  \item single-prover-multiple-verifier vs.\ multi-witness
    joint attestation: the cited literature establishes the
    position of a single prover from multiple verifiers'
    distance-bounding measurements. The witness-mesh primitive
    establishes a joint physical record across a plurality of
    mutually-coupled witness modules, in which each module
    serves both as a prover of its own observation of $\Sigma_i$
    and as a participant in the inter-module coupling that
    constrains peer modules' fast-loop responses;
  \item inter-module mutual coupling vs.\ unidirectional
    distance-bounding challenges: the cited literature performs
    challenge-response distance-bounding from each verifier to
    the prover. The witness-mesh primitive's inter-module
    coupling is mutual and continuous-analogue across the
    evidence window, not a discrete challenge-response
    sequence;
  \item joint physical evolution: the cited literature does
    not produce a joint physical record whose joint
    reproduction is the security primitive; verifier
    measurements are independent per-verifier distance
    estimates aggregated to localise the prover. The witness-
    mesh primitive's joint \cb record under public coupling
    commitment is the security primitive of
    Section~\ref{sec:wm-security-primitive}; and
  \item the joint-trace verifier and per-version empirical
    hardness commitment, neither of which is present in the
    cited secure-localisation literature.
\end{enumerate}

\paragraph{Controlled physical random functions and
PUF-modelling-attack literature.}
The witness-mesh primitive is also distinct from controlled
physical random functions and the PUF-modelling-attack
literature, including without limitation Gassend, Clarke, van
Dijk, and Devadas, ``Controlled Physical Random Functions,''
in \emph{Proceedings of the 18th Annual Computer Security
Applications Conference (ACSAC '02)}, December 2002,
pp.~149\textit{ff.}; Gassend, Clarke, van Dijk, and Devadas,
``Silicon Physical Random Functions,'' in \emph{Proceedings of
the 9th ACM Conference on Computer and Communications Security
(CCS '02)}, 2002, pp.~148--160; the controlled-PUF and
certified-execution literature derived from these works; and
the PUF-modelling-attack literature including without
limitation R\"uhrmair, Sehnke, S\"olter, Dror, Devadas, and
Schmidhuber, ``Modeling Attacks on Physical Unclonable
Functions,'' in \emph{Proceedings of the 17th ACM Conference on
Computer and Communications Security (CCS '10)}, 2010,
pp.~237--249. Controlled physical random functions (CPUFs) are
PUFs that can only be accessed via an algorithm physically
bound to the PUF in an inseparable way, enabling certified
execution and challenge-response authentication with the PUF
under the bound algorithm. The PUF-modelling-attack literature
demonstrates that PUFs can be characterised from collected
challenge-response pair sets through machine-learning
modelling attacks, motivating mitigations through nonlinearity,
controlled access, and reduced challenge-response-pair
exposure.

The witness-mesh primitive shares conceptual elements with
controlled PUFs --- specifically, the binding of an
algorithmic transfer function (the inter-module coupling
network) to physical primitives (the per-module reactors and
detectors) in a manner that constrains the algorithmic
function's evaluation by physical hardware --- but is
distinguished by:
\begin{itemize}[nosep]
  \item the multi-module joint-record structure rather than
    single-PUF challenge-response;
  \item the continuous-analogue inter-module coupling among
    physically separate witness modules, with no analogue in
    single-device CPUF literature;
  \item the public coupling commitment, with the trained
    inter-module coupling network committed to the protocol
    digest as a public commitment, vs.\ the typically-private
    PUF challenge-response model in the cited CPUF literature;
    and
  \item the per-version empirical hardness commitment that
    explicitly evaluates the witness-mesh primitive against
    PUF-modelling-attack methodologies and successor surrogate-
    modelling literature, with the resulting empirical hardness
    estimate committed to the protocol digest.
\end{itemize}
The PUF-modelling-attack literature is, in addition, one of
the declared adversary attack classes against which the
witness-mesh primitive's per-version empirical hardness
commitment of Section~\ref{sec:wm-lifetime} is, in some
preferred embodiments, evaluated.

\paragraph{Multi-camera attestation and content-provenance systems.}
The witness-mesh primitive is distinct from multi-camera
attestation and content-provenance systems, including the C2PA
Content Credentials specification (which defines manifests,
assertions, claims, claim signatures, content bindings, and
provenance data) and related distributed attestation protocols.
Such systems establish multi-witness signature schemes and
tamper-evident provenance metadata over independently-captured
records, but do not constrain the captured records by physical
mutual coupling during capture: the witnesses sign independent
records after capture, and the cross-witness consistency check is
cryptographic rather than physical. The witness-mesh primitive
constrains the joint capture itself through the
continuous-analogue inter-module coupling, so that each module's
recorded response is physically determined in part by peer
modules' simultaneous responses, and joint reproduction requires
solving the coupled physical-dynamics problem of
Section~\ref{sec:wm-security-primitive} rather than separately
reproducing multiple independent records.

\paragraph{Decentralized proof-of-location and witnessing-zone
architectures.}
The witness-mesh primitive is also distinct from decentralized
proof-of-location and witnessing-zone architectures, including
without limitation: Brito, Hadachi, Kamm, and Norbisrath,
``Decentralized Proof-of-Location systems for trust, scalability,
and privacy in digital societies,'' \emph{Scientific Reports}
15:19808, 5 June 2025 (DOI: 10.1038/s41598-025-04566-4), which
formalises a composable decentralized proof-of-location model
based on fault-tolerant witnessing zones, integrating distributed
digital signatures, distance-bounding protocols, and consensus
mechanisms; and Brito, Castillo, Hadachi, Norbisrath, and Heiss,
``Decentralized Proof-of-Location for Content Provenance:
Towards Capture-Time Authenticity,'' arXiv:2603.27883, 29 March
2026, accepted at the 5th International Workshop on Architecting
and Engineering Digital Twins (AEDT 2026) and the companion
proceedings of the 23rd IEEE International Conference on
Software Architecture (ICSA 2026), which extends the
witnessing-zone model with multiple independent observers
collectively validating physical events to produce auditable
evidence artifacts for capture-time content authenticity in
cyber-physical settings.

The cited literature establishes verifiable claims of physical
presence, spatio-temporal synchronisation, collective
attestation, and capture-time content authenticity through
witnessing-zone architectures in which multiple independent
observers (witnesses) attest to a prover's presence or to a
physical event, with the witnesses' attestations aggregated
through distributed digital signatures, distance-bounding,
consensus mechanisms, or related cryptographic primitives. The
witness-mesh primitive shares with this literature the
multi-witness structure, the spatio-temporal binding of
attestations to physical events, and the application to
capture-time provenance of physical-world content, but is
distinguished by:

\begin{enumerate}[label=(\alph*),nosep]
  \item continuous-analogue mutual coupling among witnesses
    during capture vs.\ independent observer attestation: the
    cited literature treats the witnesses as independent
    attesters whose individual attestations are aggregated
    cryptographically, with no mutual physical coupling among
    the witnesses' fast-loop responses during the evidence
    window. The witness-mesh primitive's modules are mutually
    coupled through the continuous-analogue, frame-free
    inter-module signal path of
    Section~\ref{sec:wm-inter-module-coupling}, such that no
    single module's record is independent of the joint record;
    each module's fast-loop response during the evidence window
    is physically determined in part by the simultaneous
    responses of peer modules through the inter-module coupling
    network. This is a structural property of the joint capture
    itself, not a post-capture cryptographic aggregation;
  \item public committed inter-module transfer function vs.\
    cryptographic signature aggregation: the cited literature
    aggregates per-witness attestations through cryptographic
    signatures, distance-bounding measurements, or consensus
    protocols; the trust model is rooted in cryptographic
    signature verifiability and consensus liveness assumptions
    among the witness population. The witness-mesh primitive's
    inter-module coupling network is a public commitment to the
    protocol digest, with the trust model rooted in the joint
    physical reproduction problem of
    Section~\ref{sec:wm-security-primitive} under public
    coupling, evaluated empirically per-version under
    Section~\ref{sec:wm-lifetime};
  \item sample-level joint-trace verifier with directed-edge
    coverage vs.\ aggregated witness signatures: the cited
    literature's verification step is signature-aggregation,
    distance-bound checking, or consensus verification over
    independent witness attestations. The witness-mesh
    primitive's verifier of
    Section~\ref{sec:wm-verifier-binding} performs sample-level
    selective opening over committed joint records and checks
    per-edge directed causal consistency under the declared
    coverage tuple, with no counterpart in the cited
    decentralized proof-of-location literature;
  \item per-version empirical-hardness commitment vs.\
    consensus-based attestation: the cited literature's security
    derives from cryptographic signature verifiability,
    distance-bounding correctness, and consensus-mechanism
    liveness assumptions; security is structural under those
    assumptions. The witness-mesh primitive's security is
    deployment-indexed empirical hardness committed per-version
    against declared adversary classes, including the
    cross-version transfer-attack evaluation of
    Section~\ref{sec:wm-lifetime}, with no counterpart in the
    cited literature; and
  \item closed-with-reality property at each module: the cited
    literature's witnesses observe the prover or the physical
    event as external observers, but do not themselves operate
    a fast closed loop in causal feedback with the observed
    physical reality during capture. The witness-mesh
    primitive's modules are each closed with their own portion
    of physical reality $\Sigma_i$ during the evidence window
    through the per-module fast loop of
    Section~\ref{sec:wm-architectural-primitive}, and the
    inter-module coupling is among these closed-loop modules
    rather than among external attesters.
\end{enumerate}

The cited literature's distance-bounding, consensus, and
signature-aggregation primitives are themselves separately
addressed in this disclosure: distance-bounding is distinguished
in the distance-bounding subparagraph of this section, and
signature-based attestation is distinguished in the multi-camera
attestation and content-provenance subparagraph above. The
present subparagraph addresses the higher-level architectural
combination of these primitives in the witnessing-zone
literature.

\paragraph{Distributed and network PUF aggregation.}
The witness-mesh primitive is distinct from distributed-PUF and
network-PUF constructions in which per-device challenge--response
pairs are aggregated from multiple devices in a sensor network.
Such constructions typically aggregate independent per-device
CRPs without mutual physical constraint via fast-loop coupling.
The witness-mesh primitive requires fast-loop mutual coupling
under the continuous-analogue, frame-free inter-module signal
path; aggregation of independently-captured per-device records,
even with mutual cryptographic signing, does not implement the
witness-mesh primitive.

\paragraph{Sensor-fusion and multi-sensor-calibration systems.}
The witness-mesh primitive is distinct from sensor-fusion and
multi-sensor-calibration systems that establish spatial and
temporal consistency across multi-modal sensors. Such systems
typically operate on independently-captured records and establish
post-capture geometric or temporal consistency. The witness-mesh
primitive instead imposes a physical mutual constraint during
capture through the continuous-analogue inter-module coupling,
and the joint record's hardness against reproduction is a
consequence of the coupled physical dynamics, not of post-capture
geometric consistency checks.

\subsubsection{Relationship to the two-dimensional hardness index}
\label{sec:wm-hardness-index}

In the witness-mesh embodiment, the two-dimensional hardness index
pair $(k^*_{\mathrm{dig}}, m^*_{\mathrm{ana}})$ of Section~\ref{sec:hardness-index} is
evaluated against an adversary attempting to forge the joint
\cb record across the network under the verifier of
Section~\ref{sec:wm-verifier-binding} with declared coverage tuple
$(\rho_{\mathrm{open}}, \kappa_{\mathrm{edge}},
\delta_{\mathrm{dec}}, \epsilon_{\mathrm{amp}},
\epsilon_{\mathrm{phase}})$ and verifier timing fidelity
$\sigma^{\mathrm{ver}}_{\tau}$. Specifically:

\begin{itemize}[nosep]
  \item $k^*_{\mathrm{dig}}$ measures the minimal model capacity
    required to reproduce the network's joint dynamics under the
    committed coupling, given a declared adversary observation
    budget over the network's collective inputs and outputs across
    the coupling graph $\mathcal{G}$.
  \item $m^*_{\mathrm{ana}}$ measures the minimal physical
    emulation scale required to reproduce the network's joint
    analogue response, including the propagation-delay-bounded
    inter-module signal exchange at the timing fidelity required
    to satisfy the optical-delay logical clocks of Section~\ref{sec:distributed-time}
    and the declared latency budget of
    Section~\ref{sec:wm-latency-budget} at verifier timing
    fidelity $\sigma^{\mathrm{ver}}_{\tau}$.
\end{itemize}

In some embodiments, the joint hardness index of an $N$-module
witness mesh is empirically estimated by training adversary models
against full-network records and measuring the capacity required
to reproduce the joint distribution at a declared threshold
$\varepsilon$ under the verifier of
Section~\ref{sec:wm-verifier-binding}, in accordance with the
procedure of Section~\ref{sec:hardness-index}. In some embodiments, joint hardness
exceeds the sum of per-module hardnesses because the cross-module
coupling under the committed inter-module coupling network
introduces constraints that are not reducible to per-module
characterisation; in some embodiments, joint hardness equals or
approaches the sum of per-module hardnesses where the verifier
reduces to per-module checks; the specific relationship is
empirical and is committed to the protocol digest at each declared
re-evaluation cycle.

\subsubsection{Preferred non-limiting embodiments}
\label{sec:wm-preferred}

\paragraph{Nonlinear per-module dynamics.}
In some preferred embodiments, each \RK module's fast
loop comprises per-module dynamics that are not reducible to a
linear or affine transfer under the declared meter envelope,
including without limitation the fluorescent and scattering bed
of the 1D anchor (Section~\ref{sec:anchor-1d}), the nonlinear-film,
scattering-medium, or phosphor reactor media of the 2D anchor
(Section~\ref{sec:anchor-2d}), the SLM--scattering-medium--camera reactor of Section~\ref{sec:slm-scattering-camera}, and any of the reactor variants disclosed in
Section~\ref{sec:layered-stacks} and subsequent reactor sections. The witness-mesh
primitive operates with linear per-module dynamics and linear
coupling, and the architectural primitive of
Section~\ref{sec:wm-architectural-primitive} is not so limited;
in linear embodiments, however, the joint dynamics of the coupled
network are characterised by a linear time-invariant multi-input
multi-output transfer with declared coupling graph and delays,
and the joint hardness contribution from coupled dynamics in
such embodiments reduces to that of multi-input multi-output
system identification under the declared adversary observation
budget, including without limitation the methods of Ho \&
Kalman (1966), Juang \& Pappa (Eigensystem Realization
Algorithm, 1985), Van Overschee \& De Moor (N4SID subspace
identification, 1994), Ljung (\emph{System Identification:
Theory for the User}, 2nd ed.~1999), and the dynamic-watermarking
detection literature of Hespanhol et al.~(2017) and successor
work. Linear embodiments are accordingly included for
architectural breadth and remain within the broader scope of the
witness-mesh primitive, but are \emph{not relied upon} for the
stronger empirical security effect of the witness-mesh primitive.
The stronger empirical security effect is preferentially
associated with non-affine, history-dependent, high-dimensional,
or heterogeneous embodiments; preferred embodiments are
accordingly those in which per-module dynamics, the coupling
network, or both are non-affine, history-dependent, and
high-dimensional, such that the joint dynamics are not reducible
to a linear MIMO transfer under the declared meter envelope.

\paragraph{Nonlinear coupling network.}
In some preferred embodiments, the trained inter-module coupling
network is itself nonlinear, comprising a frozen small neural
network or other declared nonlinear deterministic transfer
function from peer fast-loop outputs to per-module modulation
input, implemented in analogue hardware where consistent with the
continuous-analogue property of the inter-module signal path.

\paragraph{Amplitude modulation as canonical modality.}
In some preferred embodiments, the inter-module coupling network
modulates the receiving module's input by amplitude, consistent
with the live amplitude drive $A(t)$ of the 2D anchor of
Section~\ref{sec:anchor-2d} and the per-sample reactor-detector gain
$G_{\mathrm{det}}$ of the 1D anchor of Section~\ref{sec:anchor-1d} transposed to
the inter-module case. Amplitude modulation is a canonical
preferred modality but is not the only modality of the
witness-mesh primitive; alternative modalities are disclosed in
Section~\ref{sec:wm-coupling-modality}.

\paragraph{Phase, frequency, polarization, wavelength, and
spatial-mode embodiments.}
In some preferred embodiments, the inter-module coupling network
modulates the receiving module's input by phase, frequency,
polarization, wavelength, spatial mode, or a combination thereof,
as disclosed in Section~\ref{sec:wm-coupling-modality}. Such
embodiments are within the scope of the witness-mesh primitive
and may be selected based on the available physical layer, the
carrier modality of the inter-module signal path, and the
verifier's measurement modality. Phase-modulation embodiments are
particularly useful where the inter-module carrier is coherent
optical and the verifier has phase-resolving detection; frequency-
and wavelength-modulation embodiments are useful for
multiplexing multiple peer signals onto a shared inter-module
carrier; polarization- and spatial-mode-modulation embodiments
are useful for free-space inter-module links with polarization-
or mode-resolving detection.

\paragraph{Heterogeneous module construction.}
In some preferred embodiments, the witness mesh comprises modules
of heterogeneous construction with distinct physical reactor
classes, distinct nonlinear regimes, and distinct spectral or
temporal operating bands, contributing to the joint hardness
index by requiring an adversary to model multiple distinct
physical reactor classes simultaneously.

\paragraph{Spatial distribution.}
In some preferred embodiments, the witness mesh is deployed at
multiple physical locations with non-overlapping or
partially-overlapping observable scenes, and the joint
\cb record incorporates the spatial distribution of the network
as part of the binding. In some embodiments, the network is
collocated within a single enclosure with multiple modules
observing different aspects or regions of a shared scene, and the
binding incorporates the geometric relationship between modules.

\paragraph{Direct optical inter-module signalling.}
In some preferred embodiments, the inter-module fast-loop
broadcast is implemented as direct optical signalling between
modules, with propagation delay determined by free-space or
fibre-optic geometry, and with the optical signal carrying the
broadcast module's fast-loop output amplitude, phase, frequency,
polarization, or spatial mode (or processed feature thereof) in
continuous analogue form, without packet framing or digital
re-encoding.

\paragraph{Mixed-modality coupling.}
In some embodiments, the inter-module broadcast is implemented as
a mixed-modality channel combining optical, RF, acoustic, or
wired-electrical signalling, with each modality contributing
distinct propagation-delay and fidelity characteristics that are
committed to the protocol digest.

\subsubsection{Lifetime, reconfiguration, and management}
\label{sec:wm-lifetime}

The witness-mesh security primitive is, like all empirical
hardness claims in this disclosure, a time-indexed and
attacker-indexed quantity per Section~\ref{sec:hardness-index}. The witness-mesh
hardness at time $t$ is meaningful relative to the adversary
class available at time $t$, and degrades as adversary modelling
capability improves and as observation of the network accumulates.
Management of the witness-mesh primitive over time follows the
framework of Section~\ref{sec:hardness-index}: periodic re-evaluation, network
reconfiguration, query throttling, and fleet-level monitoring.

\paragraph{Coupling-network rotation.}
In some embodiments, network reconfiguration is performed by
re-training the trained inter-module coupling network against
declared benchmark adversary classes available at the
re-evaluation date, freezing the result to a new committed
parameter-family version, and committing the new version to the
protocol digest. The rotation period is selected based on the
empirical characterisation half-life of the prior coupling under
the declared adversary observation budget, such that the time
required for the declared adversary class to recharacterise the
new coupling exceeds the period over which the prior coupling's
empirical hardness exceeds a declared policy threshold. The
rotation period is committed to the protocol digest.

\paragraph{Cross-version transfer-attack evaluation
(non-limiting).}
The empirical hardness of a rotated coupling-network version is,
in some preferred embodiments, evaluated against an adversary
that carries forward, from prior coupling versions, all of the
following persistent characterisations to the extent these have
been accumulated under the declared adversary observation
budget:
\begin{itemize}[nosep]
  \item per-module reactor microstructure and transfer function
    estimates;
  \item analogue front-end transfer functions and detector
    noise models;
  \item module geometry and vantage configuration;
  \item scene priors and learned reconstruction models for
    each module's $\Sigma_i$;
  \item latency calibration estimates and per-edge propagation
    characterisations;
  \item verifier-threshold and tolerance estimates; and
  \item learned embeddings, surrogate state representations,
    and intermediate model components trained against the prior
    coupling that generalise to the rotated coupling.
\end{itemize}
The cross-version transfer-attack evaluation is committed to
the protocol digest as part of the per-version hardness
commitment of this section. Coupling-network rotation
invalidates accumulated adversary characterisation of the prior
coupling at the surface level only, providing a rotation
mechanism analogous in part to key rotation in cryptographic
systems; rotation is \emph{not} asserted to erase, and does not
erase, the persistent characterisations enumerated above, and
the empirical hardness of the rotated coupling-network version
is established against the cross-version transfer adversary.
Persistent characterisation across coupling-network rotations
is addressed by the additional management mechanisms of this
section, including topology reconfiguration, per-version
re-evaluation of the empirical hardness index against updated
adversary classes, and physical reactor refresh where
applicable.

\paragraph{Topology reconfiguration.}
In some embodiments, the coupling graph $\mathcal{G}$ is itself
reconfigured at declared rotation cycles, with new edges added,
removed, or re-weighted, and the new graph committed to the
protocol digest. Topology reconfiguration is independent of
coupling-network rotation and may be performed at distinct
rotation cycles. In some embodiments, topology reconfiguration is
implemented at the trained-parameter level (re-weighting existing
declared physical paths) without physical re-routing of
inter-module signal paths; in some embodiments, topology
reconfiguration includes physical re-routing where deployment
constraints permit, with the new physical paths' propagation and
hardware latencies recharacterised and committed.
In further embodiments, topology reconfiguration is implemented as
\emph{circuit-switched reconfiguration} of optical and electro-
optical interconnects, in which switching elements (including
without limitation MEMS mirrors, MEMS lens arrays, micromotor-
positioned lens elements, optical-switch fabrics, beam-steering
SLMs, and electronically reconfigurable apertures) are commanded
to reconnect or redirect signal paths between modules without
halting an active scan. In some embodiments, circuit-switched
reconfiguration is synchronised to scan phase, so that
reconfiguration events occur during declared low-information
intervals of the scan trajectory; the duration and placement of
such intervals are matched to the response time of the specific
switching element, so that for fast-responding elements (including
without limitation certain MEMS mirrors and blanking-synchronous
apertures) reconfiguration may occur within flyback or blanking
intervals, and for slower-responding elements (including without
limitation certain spatial light modulators, lens arrays, and
optical-switch fabrics) reconfiguration occurs within declared
longer reconfiguration windows scheduled outside the active
acquisition interval. The reconfiguration schedule, switch states,
device-specific response times, and any intermediate transient
measurements are committed to the protocol digest. In
some embodiments, the circuit-switched topology is varied online
across an episode under a declared reconfiguration policy that
itself may be trainable or adaptively-controlled.

\paragraph{Module addition, removal, and quarantine.}
In some embodiments, modules are added to or removed from the
witness mesh over the deployment lifetime under the enrolment,
revocation, and quarantine framework of the networking section.
Modules whose hardness or reliability drops below policy
thresholds may be quarantined or removed and the coupling graph
updated accordingly.

\paragraph{Per-version hardness commitment.}
In some embodiments, each committed deployment version --- comprising
a specific coupling graph $\mathcal{G}$, a specific configured
deterministic inter-module coupling transfer function (optionally
trained), specific per-module trained parameter families, a
specific verifier timing fidelity $\sigma^{\mathrm{ver}}_{\tau}$,
and a specific verifier coverage tuple
$(\rho_{\mathrm{open}}, \kappa_{\mathrm{edge}},
\delta_{\mathrm{dec}}, \epsilon_{\mathrm{amp}},
\epsilon_{\mathrm{phase}})$ --- is associated with a committed
empirical hardness estimate
$(k^*_{\mathrm{dig}}, m^*_{\mathrm{ana}})$ established at the
declared re-evaluation cycle and committed to the protocol
digest, so that downstream verifiers and policy components can
act on the version-specific hardness profile. In some
embodiments, the per-version hardness commitment additionally
records, as committed artifacts in the protocol digest:
\begin{itemize}[nosep]
  \item the declared adversary classes against which the
    empirical hardness was evaluated;
  \item the declared adversary observation budget;
  \item the model families employed in the evaluation,
    including without limitation linear and subspace
    system-identification baselines, neural state-space and
    recurrent surrogates, reservoir-computing surrogates,
    generative-scene-reconstruction adversaries, and
    analogue-relay adversary models;
  \item identifiers for the training and evaluation corpora used
    in the evaluation, in some embodiments including
    cryptographic hashes of corpus contents committed to
    $\Pi_{\mathrm{dig}}$;
  \item training/evaluation split seeds, where applicable, so
    that the corpus partition is reproducible from the
    committed seeds;
  \item identifiers for model implementations, including
    without limitation source-code commit hashes, container
    image hashes, and software-environment manifests, so that
    the evaluation pipeline is reproducible;
  \item calibration transcripts for the verifier-side timing
    fidelity $\sigma^{\mathrm{ver}}_{\tau}$, the analogue
    front-end latency, and the residual tolerances
    $\epsilon_{\mathrm{amp}}$ and $\epsilon_{\mathrm{phase}}$;
  \item attack success/failure logs across declared adversary
    attack classes;
  \item a replay protocol sufficient to re-run the verifier's
    acceptance decision on a sample of opened windows from the
    evidence corpus;
  \item the acceptance thresholds applied at the verifier;
  \item declared confidence intervals on the empirical hardness
    estimate; and
  \item negative results material to the estimate, including
    declared adversary attack classes that succeeded under
    declared budgets.
\end{itemize}
The commitment of these artifacts supports auditability of the
per-version hardness estimate by a third-party verifier or
auditor, without obligating the implementer to a specific
numerical hardness multiplier. The empirical hardness estimate
is, consistent with Section~\ref{sec:hardness-index}, a time-indexed,
attacker-indexed quantity; the per-version commitment records
the conditions under which the estimate was established at the
declared re-evaluation cycle.

\paragraph{Control of physical operation through the per-version
hardness commitment (non-limiting).}
In some embodiments, the committed per-version hardness profile
is used to control physical operation of the witness mesh,
including without limitation:
\begin{itemize}[nosep]
  \item selecting or updating the verifier coverage tuple
    $(\rho_{\mathrm{open}}, \kappa_{\mathrm{edge}},
    \delta_{\mathrm{dec}}, \epsilon_{\mathrm{amp}},
    \epsilon_{\mathrm{phase}})$ in response to the empirical
    hardness profile of the deployment version;
  \item accepting or rejecting evidence-window records at the
    verifier under acceptance thresholds derived from the
    per-version hardness profile;
  \item triggering coupling-network rotation of
    Section~\ref{sec:wm-lifetime} when the empirical
    characterisation half-life of the prior coupling falls below
    a declared policy threshold;
  \item triggering topology reconfiguration to add, remove, or
    re-weight coupling edges in $\mathcal{G}$ in response to
    the per-version hardness profile or to evaluation against
    cross-version transfer attacks;
  \item quarantining individual modules whose per-module
    contributions to the joint hardness fall below a declared
    threshold under the declared adversary class;
  \item disabling a deployment version whose committed
    evaluation no longer satisfies a declared policy threshold,
    pending re-evaluation against updated adversary classes;
    and
  \item gating downstream apparatus that consumes evidence-
    window records on the per-version hardness profile, such
    that downstream apparatus operates only on evidence-window
    records originating from deployment versions whose committed
    evaluation satisfies a declared policy threshold.
\end{itemize}
The per-version hardness commitment of this section accordingly
operates as a control input to physical apparatus operations,
including verifier configuration, coupling-network rotation,
topology reconfiguration, module quarantine, deployment-version
enablement, and downstream-apparatus gating, rather than as an
abstract evaluation or policy artifact.

\subsubsection{Cross-references}
\label{sec:wm-cross-references}

The witness-mesh primitive cross-references the following
sections of this disclosure:

\begin{itemize}[nosep]
  \item Section~\ref{sec:anchor-1d} (1D anchor): the per-sample fast loop of the
    1D anchor is a canonical instantiation of a single
    witness-mesh module's fast loop closed with reality, and the
    reactor-detector gain $G_{\mathrm{det}}$ is a canonical
    intra-module instantiation of a trained committed
    amplitude-modulation transfer function transposed to the
    inter-module case.
  \item Section~\ref{sec:anchor-2d} (2D anchor): the continuous live amplitude
    broadcast of the 2D anchor is a canonical instantiation of a
    single witness-mesh module's fast loop, with the live
    coupling-path discipline (no per-period wait, no
    runtime-adaptive element) directly transposed to the
    inter-module case.
  \item Section~\ref{sec:layered-stacks} (reactor-microstructure-unclonability variants): canonical
    instantiations of per-module reactor microstructure for the
    combined witness-mesh-plus-reactor-microstructure embodiments.
  \item Section~\ref{sec:crt-phosphor} (CRT/phosphor analogue-memory
    extension) and Section~\ref{sec:dual-loop-appliance} (dual-loop
    appliance combining a fast loop with a slow analogue-memory
    loop): optional extensions orthogonal to the witness-mesh
    primitive.
  \item Section~\ref{sec:matryoshka-reactor} (intra-module multi-reactor coupling):
    structurally similar to the inter-module witness-mesh
    coupling but operates within a single \RK module's
    extended state $X_t$; the witness-mesh primitive operates
    across modules with shared world state $W_t$.
  \item Section~\ref{sec:hardness-index} (two-dimensional hardness index): provides
    the framework against which the joint witness-mesh hardness
    is evaluated.
  \item Section~\ref{sec:networks} (Networked \RKs and
    Transitive Proof-of-Projection): the witness-mesh primitive
    is the principal refinement of the networked-modules
    framework.
  \item Section~\ref{sec:distributed-time} (optical-delay logical clocks): provides
    the physical timing anchors against which the latency budget
    of Section~\ref{sec:wm-latency-budget} is evaluated.
  \item Section~\ref{sec:two-seed-protocol} (Merkle commitment and two-seed selective
    opening): provides the verifier protocol underpinning
    Section~\ref{sec:wm-verifier-binding}.
  \item Section~\ref{par:multi-reactor-topologies} (multi-reactor network topologies): the
    coupling graph $\mathcal{G}$ of the witness mesh is a
    security-purposed instantiation of the multi-reactor topology
    framework.
\end{itemize}


% ----------------------------------------------------------------------
\subsection{Anchor Embodiment: The Bench-Top 1D Reality Kernel}
\label{sec:anchor}
\label{sec:anchor-1d}

Before developing the full formalism, a concrete embodiment is
described to provide physical intuition. The bench-top 1D
\RK is a minimal but complete implementation that
illustrates the two-subsystem topology, the reactor-to-scene
amplitude coupling, the trainable parameter families, and the
joint-record hardness objective, in a form that can be built
from inexpensive commercial components. This embodiment is
non-limiting; a 2D generalisation with continuous real-time
coupling is described in Section~\ref{sec:anchor-2d}, and
substrate and coupling variants are described in subsequent
sections.

\paragraph{Mapping to embodiments (a) and (b).}
The 1D anchor description that follows in this section, including
its per-sample $A(t)$ feedback path and its trainable parameter
families, is a concrete realisation of the parallel-subsystem
embodiment (a) of
Section~\ref{sec:architectural-alternatives}. The cascade
embodiment (b) of
Section~\ref{sec:architectural-alternatives} is illustrated in
FIG.~4A by a heavy-stroke forward low-latency coupling arrow
running from the scene detector into the UV laser of the
reactor subsystem. The forward low-latency coupling, where
present, is in addition to the per-sample $A(t)$ feedback
described below, and may be implemented in any of the
modalities recited in
Section~\ref{sec:architectural-alternatives} (electrical
direct, mixed-signal, optical gain-pumped, or optical direct
relay), with its end-to-end latency declared and committed to
the protocol digest. Bench-top instantiations of the 1D anchor
described in this section, including the FIG.~4A configuration,
typically realise the $A(t)$ feedback path electrically through
the conditioning chain described below. The forward low-latency
coupling, where engaged, is similarly typically realised
electrically in bench-top instantiations; fully-optical and
hybrid optical-electrical realisations of either or both
coupling paths are within the scope of this disclosure.

\paragraph{Physical layout.}
In one non-limiting configuration, the bench-top 1D \RK comprises two physically separate subsystems under a
common controller~10:
\begin{enumerate}[nosep]
  \item A \emph{scene subsystem}, open to the world under a
    declared eye-safety envelope, comprising a scene laser
    (for example an eye-safe infrared or visible laser), a
    scene steering mirror mounted on a single-axis micromotor
    servo, and a scene detector (for example a photodiode).
  \item A \emph{reactor subsystem}, sealed behind a declared
    safety envelope, comprising: a reactor container located
    within the sealed reactor envelope, for example an
    open-topped aluminium cup or can providing an opaque
    reflective interior, together with a heterogeneous bed
    within the container (for example mixed fluorescent glass
    beads or fragments, including without limitation
    uranium-glass (``vaseline glass'') elements where
    permitted by applicable safety and regulatory
    requirements, clear or frosted glass marbles,
    mirror-finished marbles, and crumpled aluminium foil); a
    reactor ultraviolet illumination source (for example a
    near-UV laser at about $405\,\mathrm{nm}$) mounted on the
    side of the container; a tunable emitter-side focus lens
    mounted on its own micromotor servo, independent of the
    scanning mirror; a reactor steering mirror mounted on a
    separate single-axis micromotor servo; a tunable
    detector-side focus lens mounted on its own micromotor
    servo, independent of both the emitter-side focus and the
    steering mirror; a green-pass or long-pass optical filter;
    and a reactor detector (for example a photodiode) viewing
    the container interior through the detector-side focus
    lens and the filter.
\end{enumerate}

In some embodiments, the reactor detector is mounted near the
reactor illumination assembly, approximately co-aligned with
the reactor ultraviolet illumination source's direction of
propagation, with precise optical-axis sharing not required.
In some embodiments, the reactor detector is mounted at an
angle to the specular axis. In some embodiments, these
configurations are treated as anchor-equivalent declared
detector geometries because the detector response is dominated
by wavelength-shifted fluorescence selected by the green-pass
filter, which is robust to alignment imprecision. Differences
in signal level, glare rejection, or path-length statistics
between the two configurations are characterised per
embodiment and recorded in the meter envelope. The declared
detector-path geometry is committed to the protocol digest.

In some embodiments, the reactor container is wrapped or
coated with reflective and opaque layers to suppress stray
light and to contain the bed. In some embodiments, shaking or
otherwise re-configuring the bed between runs re-randomises
internal microstructure, so that a fixed scan program applied
to different internal states yields different but
statistically related reactor responses. In some embodiments,
a deliberate bed reconfiguration event is logged in the
protocol digest, and any hardness, calibration, or enrolment
claim made after reconfiguration is tied to the resulting
reactor state and meter envelope.

\paragraph{Reactor substrate (non-limiting).}
In the 1D anchor, the reactor substrate is a shallow bed of
mixed beads or fragments including fluorescent species. In
some embodiments, the bed includes uranium-glass
(``vaseline glass'') marbles as the primary fluorescent
species, which upon near-UV excitation emit wavelength-shifted
green fluorescence with a characteristic persistence time,
providing both scattering from the bed's reflective and
refractive elements and wavelength-shifted emission from the
fluorescent elements. In some embodiments, other fluorescent
or luminous materials --- for example europium-doped,
terbium-doped, or manganese-doped fluorescent glass beads;
stable phosphor beads; other fluorescent glass fragments; or
fluorescent polymer beads --- may be substituted or combined
with uranium-glass, subject to applicable safety and
regulatory requirements. The green-pass or long-pass filter in
front of the reactor detector selects the wavelength-shifted
emission band and rejects the UV excitation, so the reactor
detector records a response dominated by fluorescence together
with scattered light in the pass band. Implementations of the
reactor substrate follow applicable safety, handling, and
regulatory requirements.

\paragraph{Spherical, orb, and isotropic reactor geometries (non-limiting).}
In some embodiments, the reactor is configured as a
spherical, ellipsoidal, or otherwise enclosed
three-dimensional volume --- non-limitingly termed an
\emph{orb} --- within which the emitter and recorder
occupy interior positions and the scene surrounds the
module rather than facing it from a preferred direction.
The orb interior includes mixed fluorescent, scattering,
and absorbing media as in the 1D anchor cup embodiments
(Section~\ref{sec:anchor-1d}), arranged for approximately
isotropic coverage of the surrounding scene from the
module's vantage. The emitter is one or more sources
arranged for full-solid-angle interrogation (for example a
polyhedral arrangement of laser-illuminated fibre
couplers, a swept galvo with an enclosing
diffuser-reflector boundary, or distributed microemitters
embedded in the orb wall); the recorder is correspondingly
an arrangement of photodetectors configured for isotropic
collection. In PolieBot embodiments
(Definition~\ref{def:poliebot}) and other mobile or fleet
deployments where the scene cannot be confined to a fixed
angular sector relative to the module, the orb geometry
enables full-solid-angle observation and interrogation
without mechanical re-aiming. Convolution-bundle
commitment, protocol digests, and meter envelopes apply to
the orb embodiment on the same footing as planar,
cylindrical, and directional reactor geometries; the
geometric variant does not by itself alter the
verification or hardness disciplines that apply to other
reactor configurations. Device-specific signatures arising
from orb fabrication tolerances (sphericity, wall material
inhomogeneity, interior medium packing geometry,
emitter-and-recorder alignment) contribute to the 2D
hardness index in the same way as in non-orb embodiments
(Section~\ref{sec:security-theory}).

\subsubsection{Scene subsystem}
\label{sec:anchor-scene-subsystem}

In the 1D anchor, the scene subsystem comprises an eye-safe
scene laser, a scene steering mirror on a single-axis
micromotor servo driven by a scene scan profile
$\mathrm{scan}_{\mathrm{s}}(t)$, and a scene detector. The
scene laser's instantaneous amplitude is driven by $A(t)$
from the reactor subsystem (see
Section~\ref{sec:anchor-1d-coupling}). In some embodiments,
the scene detector is a single detector (for example a
photodiode) as the anchor default; a camera, a photodiode
array, or other spatially resolved detector is a declared
extension of the single-detector default. In some
embodiments, the scene detector is coaxial with the scene
emitter via a polarisation combiner or equivalent.

\subsubsection{Reactor subsystem}
\label{sec:anchor-reactor-subsystem}

In the 1D anchor, the reactor subsystem is sealed and
comprises, on the excitation path and in optical order: the
reactor UV laser, a tunable emitter-side focus lens on its
own micromotor servo following a trained profile
$f_{\mathrm{emit}}(t)$ periodic in the scan period $T$, a
reactor steering mirror on a separate single-axis micromotor
servo driven by a reactor scan profile
$\mathrm{scan}_{\mathrm{r}}(t)$, and the sealed container
with its heterogeneous fluorescent-and-scattering bed. On the
detection path, the reactor subsystem further comprises a
tunable detector-side focus lens on its own micromotor servo
following a trained profile $f_{\mathrm{det}}(t)$ periodic
in $T$, a green-pass or long-pass filter, and the reactor
detector. The reactor UV laser is driven by a trained profile
$A_{\mathrm{emit}}(t)$ periodic in $T$;
$A_{\mathrm{emit}}(t)$ probes the fluorescent and nonlinear
response of the bed at each scan phase.

In some embodiments, the detection path and the excitation
path are approximately co-aligned, with the reactor detector
mounted near the UV laser and looking into the bed along a
direction close to that of the UV illumination; in other
embodiments, the detection path and the excitation path are
angularly separated. In some embodiments, these
configurations are treated as anchor-equivalent declared
detector geometries because the green-pass filter selects
the wavelength-shifted fluorescence response, which is
robust to alignment imprecision. The declared path geometry
is committed to the protocol digest. In some embodiments, the
reactor detector is a single detector as the anchor default.

\subsubsection{Per-sample discrete feedback coupling}
\label{sec:anchor-1d-coupling}

In the 1D anchor, the scene and reactor subsystems are
coupled by a per-sample discrete feedback loop rather than by
a continuous broadcast. At each scan step $t$, the common
controller:
\begin{enumerate}[nosep]
  \item commands the scene steering mirror to
    $\mathrm{scan}_{\mathrm{s}}(t)$ and the reactor steering
    mirror to $\mathrm{scan}_{\mathrm{r}}(t)$;
  \item commands the reactor emitter-side focus to
    $f_{\mathrm{emit}}(t)$ and the reactor detector-side
    focus to $f_{\mathrm{det}}(t)$;
  \item after the steering mirrors and focus actuators have
    reached their declared settled states, commands the
    scene laser to emit at amplitude $A(t)$ and the reactor
    UV laser to emit at $A_{\mathrm{emit}}(t)$;
  \item records the scene detector's and reactor detector's
    responses for step $t$, yielding $y_{\mathrm{s}}(t)$ and
    $y_{\mathrm{r}}(t)$;
  \item the reactor-detector output $y_{\mathrm{r}}(t)$ is
    passed through the declared deterministic coupling path,
    which applies the frozen committed gain to produce the
    next scene amplitude,
    $A(t+1)=G_{\mathrm{det}}(y_{\mathrm{r}}(t))$, where
    $G_{\mathrm{det}}$ is the trained reactor detector gain
    (see Section~\ref{sec:anchor-trainable}).
\end{enumerate}
In some embodiments, the common controller sequences this
fixed calculation and applies the result to the scene laser
drive; the gain law $G_{\mathrm{det}}$ is not learned or
adapted during the evidence window. In some embodiments, the
coupling path --- comprising the reactor photodetector,
protective conditioning, and the frozen committed
$G_{\mathrm{det}}$ --- holds one sample of buffering inherent
to the per-sample discrete feedback loop, and contains no
multi-period storage, no runtime-adaptive element, and no
runtime-learned element.

In some embodiments, the scene and reactor emit-and-record
phases within a step are logically ordered but may overlap in
actual timing; in some embodiments, the steering mirrors and
focus actuators are required to be at their declared settled
states before either detector records, to avoid motion and
defocus artefacts. In some embodiments, $G_{\mathrm{det}}$
is trained offline and frozen for the committed deployment
version. The latency from reactor photocurrent at step $t$ to
scene emitter modulation at step $t+1$ is bounded by a
declared settling time and is committed to the protocol
digest.

\paragraph{Boot.}
In some embodiments, the first-sample amplitude $A(0)$ is
initialised from one of the following declared sources: (a) a
committed public constant, for example a declared offset into
the decimal expansion of $\pi$; (b) co-illumination from a
distant reactor via the auxiliary noise co-illumination port
of Section~\ref{sec:anchor-noise}; or (c) any other declared
constant or signal committed to the protocol digest. In
embodiments using co-illumination as the boot source, the
reactor detector samples the co-illuminated reactor response
before the first scene emission, and the frozen committed
$G_{\mathrm{det}}$ maps that sampled response to $A(0)$ via
the same live coupling path used during deployment. In some
embodiments, the boot source is recorded in the protocol
digest so that a verifier can reproduce the loop from its
initial condition.

\subsubsection{Trainable parameter families}
\label{sec:anchor-trainable}

In the 1D anchor, six trainable parameter families are
optimised jointly under hardware-in-the-loop training (see
Section~\ref{sec:anchor-training}) against the hardness
objective of Section~\ref{sec:anchor-hardness}. The count of
six is an addressable-parameter inventory of this anchor
embodiment and is not asserted as a guaranteed effective
capacity, task-performance bound, or training-scale result;
effective capacity, generalisation, and task performance are
empirical and embodiment-specific. The trainable parameter
families are:
\begin{enumerate}[nosep]
  \item The scene scan trajectory
    $\mathrm{scan}_{\mathrm{s}}(t)$, periodic in $T$.
  \item The reactor scan trajectory
    $\mathrm{scan}_{\mathrm{r}}(t)$, periodic in $T$. In
    some embodiments, $\mathrm{scan}_{\mathrm{s}}(t)$ and
    $\mathrm{scan}_{\mathrm{r}}(t)$ are initialised
    identical at training start and are allowed to warp
    apart under training; the committed trained pair is
    recorded in the protocol digest.
  \item The reactor emitter-side focus lens profile
    $f_{\mathrm{emit}}(t)$, periodic in $T$. Controls the
    focusing of the reactor UV laser onto the fluorescent
    and scattering bed.
  \item The reactor detector-side focus lens profile
    $f_{\mathrm{det}}(t)$, periodic in $T$. Controls the
    region and depth of the bed over which the reactor
    detector integrates at each scan phase, and may be
    trained independently from $f_{\mathrm{emit}}(t)$.
  \item The reactor UV drive profile $A_{\mathrm{emit}}(t)$,
    periodic in $T$, applied upstream of the bed.
    $A_{\mathrm{emit}}(t)$ probes the substrate's
    intensity-dependent fluorescent and nonlinear response
    at each scan phase.
  \item The reactor detector gain $G_{\mathrm{det}}$,
    applied in the live coupling path between the reactor
    photodetector and the scene emitter drive.
    $G_{\mathrm{det}}$ may be a time profile
    $G_{\mathrm{det}}(t)$ periodic in $T$, a nonlinear
    transfer function $G_{\mathrm{det}}(I)$ acting on
    instantaneous reactor photocurrent $I$, or a combined
    $G_{\mathrm{det}}(t,I)$. In some embodiments,
    $G_{\mathrm{det}}$ is allowed to be nonlinear, for
    example implementing soft compression at high
    photocurrents to prevent wipeout, or a dead-zone near
    zero. Where $G_{\mathrm{det}}$ is time-dependent, its
    phase origin is defined relative to the common scan
    period $T$ and committed to the protocol digest.
\end{enumerate}
In some embodiments, the trained values defining all six
parameter families are committed to the protocol digest as
part of the committed parameter-family version.

\subsubsection{Hardness objective}
\label{sec:anchor-hardness}

In some embodiments, the trainable parameter families of
Section~\ref{sec:anchor-trainable} are optimised against a
declared adversary family Eve, defined as strong as the
available training budget allows; a stronger declared Eve
produces a more robust trained configuration. In some
embodiments, Eve's attack surface is the full joint record
--- scene observations, reactor observations, and their
mutual consistency under the declared scan and feedback
protocol. In some embodiments, Eve may attempt to forge
either record or their conditional relationship, and may use
white-noise-baseline probing of the public interface to
characterise the channel. In some embodiments, the training
objective is to maximise the capacity, training-data, and
compute budget Eve requires to produce a joint record the
legitimate verifier accepts.

In some embodiments, operational hardness claims for the
anchor embodiment are quantified by the measured canonical
hardness index pair
$(k^*_{\mathrm{dig}}(\theta;\varepsilon),
m^*_{\mathrm{ana}}(\theta;\varepsilon, r))$ under the declared
adversary family and the declared meter envelope, as
described in Section~\ref{sec:security-theory}, evaluated on
the reactor subsystem's fluorescent and nonlinear response,
the per-sample feedback coupling, and the resulting joint
scene-and-reactor record.

\subsubsection{Training pathways}
\label{sec:anchor-training}

In some embodiments, the trainable parameter families are
fitted by one or more of the following training pathways,
which may be applied sequentially or in combination:

\paragraph{Stage 1 --- reactor against white noise (anchor
baseline, non-limiting).}
In some embodiments, the reactor subsystem is characterised
in isolation under white-noise inputs, and the reactor-side
parameter families ---
$\mathrm{scan}_{\mathrm{r}}(t)$, $f_{\mathrm{emit}}(t)$,
$f_{\mathrm{det}}(t)$, $A_{\mathrm{emit}}(t)$, and
$G_{\mathrm{det}}$ --- are fitted to maximise the capacity
budget Eve requires to match the reactor's white-noise-
challenge response. In some embodiments, the scene scan
trajectory $\mathrm{scan}_{\mathrm{s}}(t)$ is initialised
from the reactor scan $\mathrm{scan}_{\mathrm{r}}(t)$, from
a declared calibration scan, or from a terminated or
dummy-scene training run, and is subsequently co-trained
during scene-coupled or simulated-scene training (Stage 2
or later). In some embodiments, once the reactor-side
parameter families and the scene scan are fitted through
the relevant stages, the combined parameter-family version
is committed and the device is deployed against arbitrary
scenes the reactor has never seen, with the reactor serving
as a signature attached to the scene illumination via the
committed feedback loop. In some embodiments, the
white-noise training regime is selected as a broad baseline
stressor: Eve's inability to match the reactor on
broadband-noise challenges at the trained capacity budget
is treated as evidence of reactor-side empirical hardness
under the declared meter envelope. Deployment hardness for
scene-driven operation remains evaluated on the committed
joint record and reduces to the measured canonical hardness
index pair under the declared adversary family.

\paragraph{Stage 2 --- digital twin with simulated scenes
(declared extension, non-limiting).}
In some embodiments, a digital twin of the reactor
subsystem is fitted from physical measurements, and the
trainable parameter families are further optimised by
running the digital twin against a library of simulated
scenes. In some embodiments, Stage 2 training targets
whole-system properties --- task performance, attention
allocation, joint-record hardness under scene-driven
operation --- that are not directly optimised by
white-noise training alone. In some embodiments, each Stage
2 acceptance produces a new committed parameter-family
version.

\paragraph{Stage 3 --- fleet-scale real-time analog
reinforcement learning (declared extension, non-limiting).}
In some embodiments, once multiple devices are deployed,
the trainable parameter families on individual devices or
on fleet-coordinated subsets are further updated by
real-time analog reinforcement learning using committed
physical evidence from deployment. In some embodiments,
Stage 3 updates are applied only between declared evidence
windows, or within windows expressly declared as training
windows. Each accepted update creates a new committed
parameter-family version, including a new $G_{\mathrm{det}}$
fingerprint where applicable; verification or deployment
windows using a given version treat that version as frozen
for the duration of the window.

In some embodiments, the training pathway used to fit a
given deployed device, the window type under which any
update was accepted, and the resulting committed
parameter-family version are recorded in the protocol
digest.

\subsubsection{Yoked operation (orthogonal declared extension)}
\label{sec:anchor-yoked}

In some embodiments, Yoked operation is an orthogonal
declared extension of the anchor and may be composed with
any of the training pathways of
Section~\ref{sec:anchor-training}. In embodiments where
direct transfer-entropy and conditional-Lyapunov-exponent
estimates are available, those meters are used directly per
the formalism of Section~\ref{sec:yoked}. In some
embodiments, a Yoked-state classifier is trained on the
committed joint evidence record as a classifier-mediated
declaration embodiment, especially for scene classes where
direct two-sided state access is unavailable; the classifier
embodiment does not replace the meter-based Yoked formalism
of Section~\ref{sec:yoked}. In some embodiments, the
trainable parameter families may be further optimised
against scene distributions conditioned on the Yoked
declaration signal (meter-based or classifier-based, as
declared). In some embodiments, when this extension is
engaged, the Yoked declaration method, and the classifier
identity and thresholds where applicable, are committed to
the protocol digest.

\subsubsection{Noise co-illumination (declared extension)}
\label{sec:anchor-noise}

In some embodiments, the sealed reactor envelope further
includes an auxiliary noise co-illumination subsystem
comprising a dedicated auxiliary light source (for example
a low-power laser), an auxiliary scanning mechanism, and an
RF input port. In some embodiments, when the noise
co-illumination subsystem is engaged, noise light is
delivered to the reactor container alongside the trained
reactor UV drive $A_{\mathrm{emit}}(t)$; the fluorescent
and nonlinear substrate then mixes the trained drive with
the co-illuminated noise, and the composite response is
delivered through the green-pass filter to the reactor
detector and onwards to the live feedback path.

\paragraph{Trainable parameters (local).}
In some embodiments, the auxiliary scanning mechanism and
an auxiliary amplitude gain on the auxiliary light source
are locally trainable against standard Eve-simulation
training analogous to
Section~\ref{sec:anchor-hardness}. In some embodiments,
these two additional trainable parameter families are
committed to the protocol digest when the extension is
engaged, distinct from and additional to the anchor's
trainable parameter families.

\paragraph{External signal at the RF port (declared per
deployment).}
In some embodiments, the RF port accepts a signal from one
of the following declared sources: (a) a white-noise
generator, for example during training-time reactor
characterisation; (b) a committed noise profile stored in
the protocol digest, for deployments requiring reproducible
injected noise; or (c) an external runtime input, for
example the output of one or more other devices of the same
class or a fleet-level coordination signal. In some
embodiments, the external signal itself is not trained by
the local device; only the auxiliary scan law and auxiliary
amplitude gain, which determine how the external signal is
projected into the reactor container, are locally trainable.

\paragraph{Tap point for inter-device sharing (non-limiting).}
In some embodiments where option (c) is used, the specific
tap point on the sharing device from which the external
signal is drawn --- for example the sharing device's
reactor photocurrent pre-$G_{\mathrm{det}}$, or its
coupling-path output post-$G_{\mathrm{det}}$ --- is
committed to the protocol digest on both the sharing and
receiving devices.

\paragraph{Relationship to the anchor's trained regime.}
In some embodiments, when the noise co-illumination
subsystem is engaged during Stage 1 reactor training, the
trained reactor is trained in the noise-present regime in
which it will be deployed. In some embodiments, the noise
co-illumination subsystem also provides a boot source for
the first-sample amplitude $A(0)$, as described in the Boot
paragraph of Section~\ref{sec:anchor-1d-coupling}.

\paragraph{Direct coupling through declared optical
feedthroughs (non-limiting).}
In some embodiments, ambient or external optical signals
--- including signals from other deployed devices of the
same class --- may enter the reactor envelope through
declared optical feedthroughs, filtered ports, or
characterised stray-light paths that remain within the
declared safety envelope. Such coupling is separate from
the auxiliary noise co-illumination subsystem described
above and is addressed by the envelope-management and
meter-envelope provisions described in connection with the
protocol digest and Section~\ref{sec:safety-envelope}.

\subsubsection{CRT/phosphor analogue-memory extension
(declared extension)}

In some embodiments, the 1D anchor is further extended by
the CRT/phosphor analogue-memory extension of
Section~\ref{sec:crt-phosphor-extension}, in which the
reactor detector output (or a derived signal) is written to
a CRT whose XY deflection is locked to a declared scan
coordinate. In the 1D anchor, the default coupling between
the reactor detector and the CRT drive is electrical rather
than optical, because the green-pass filter architecture
enforces wavelength separation between the reactor's
ultraviolet drive and its wavelength-shifted green-fluorescent
response, and optical addition of the reactor output onto
the CRT would muddy that separation. Optical CRT addition
is a declared variant within this extension for embodiments
in which the wavelength allocation is compatible.

\paragraph{Mapping to core primitives (1D anchor).}
In the 1D anchor: the \emph{scene emitter} is the eye-safe
scene laser; the \emph{reactor emitter} is the sealed
reactor UV laser; each subsystem has its own \emph{scanner}
(single-axis micromotor mirror); the \emph{reactor stack}
comprises, on the excitation path, the tunable emitter-side
focus lens and the sealed container with its fluorescent
and scattering bed, and on the detection path, the tunable
detector-side focus lens and the green-pass filter; each
subsystem has its own \emph{detector}; and the
\emph{controller} generates the scan profiles, the reactor
drive profile, and the feedback sequencing. The \emph{live
coupling path} comprises the reactor detector, protective
conditioning, and the frozen committed reactor detector
gain $G_{\mathrm{det}}$. The \emph{protocol digest} records
the scan profiles, the emitter-side focus profile
$f_{\mathrm{emit}}(t)$, the detector-side focus profile
$f_{\mathrm{det}}(t)$, the declared detector-path geometry,
the reactor drive profile, the reactor detector gain
(including phase origin for time-dependent
$G_{\mathrm{det}}$), the boot source for $A(0)$, the
declared window type, any deliberate bed reconfiguration
events, and the committed parameter-family version. The
\emph{meter vector} includes emitter-side and detector-side
focus-actuator settling indicators, mirror settling
indicators, detector SNR, filter pass-band calibration,
geometry-difference characterisation, and boot-state
indicators.
\subsection{System Architecture}
% ----------------------------------------------------------------------

\subsubsection{System overview}

In the most general terms, a \RK module includes
an emission subsystem, a detection subsystem, one or more physical
paths coupling emission to detection (optionally through a reactor
medium and/or an external scene), and a controller. The controller executes a control
protocol $U_{0:T}$ that specifies, at each time step, the emission
configuration (scan coordinate, intensity, wavelength, polarisation,
phase, pattern, timing, and gating), the detector configuration
(exposure, gain, spectral band, integration window, gating), and any
reactor-medium controls (bias, pump power, temperature, mechanical
position).

The emission subsystem may include, without limitation, lasers, LEDs,
broadband sources, CRTs, projectors, spatial light modulators, and
antenna arrays. The detection subsystem may include, without limitation,
photodiodes, cameras, photomultipliers, avalanche photodiodes,
spectrometers, bolometers, and microphone arrays.

The controller produces a \cb $C_{0:T}$ by pairing each
emission event with the corresponding detection result. Each \cba sample $c_t = (u(t), \mathbf{y}_t)$ records both what was done and
what was observed, under a shared timestamp from a common timebase.

In some embodiments, the controller also maintains and updates a reactor
latent $Z_\theta(t)$ and logs protocol digests, meter summaries, and
optional commitments. The full evidence package for a run includes
the \cb, protocol digest, meter envelope, and any commitments
or selective-opening artifacts.

\subsubsection{Extended state, world kernel, and measurement map}

In some embodiments, the system is modelled as a stochastic dynamical
system with the extended state $X_t$
(Section~\ref{sec:definitions}),
evolving under a world transition kernel
$\mathsf{P}_{\mathrm{world}}(X_{t+1} \mid X_t, u(t))$ and observed
through a measurement map $\mathbf{y}_t \sim p(\cdot \mid X_t, u(t))$.
This factorisation is a modelling convenience; the physical system does
not necessarily separate into these components, and the factorisation
may differ across embodiments.  In multi-device settings
(Section~\ref{sec:networks}), $W_t$ denotes the shared external world
state common to all devices, while $X_t$ denotes the per-device
extended state including internal degrees of freedom; $W_t$ is a
component of each device's $X_t$ but not identical to it.

\subsubsection{Layered optics, memory media, and skip paths}

In some embodiments, the optical path between emitter and detector
includes multiple layers, each contributing dynamics, memory, or
transformation. Non-limiting layers include: scattering plates,
nonlinear thin films, phosphor or luminescent persistence layers,
fibre-delay loops, liquid-crystal or photorefractive layers, and
engineered diffusers.

In some embodiments, skip paths (analogous to residual connections in
neural networks) allow a portion of the signal to bypass one or more
layers, creating a reference channel that preserves information about
the input while the main path undergoes transformation. Skip paths may
be implemented optically (beam splitters, fibre couplers) or
electronically (signal tapping before and after a processing stage).

In some embodiments, layers can be inserted, removed, or rearranged
without changing the Markov-kernel abstraction or the \cba
format, provided the protocol digest and meter envelope are updated to
reflect the current configuration.

% ----------------------------------------------------------------------
\subsection{Coupling Modes}
% ----------------------------------------------------------------------

\subsubsection{Direct coupling}

In direct-coupling embodiments, the emitter illuminates an external
scene and the detector observes the scene's response directly. The
scene and reactor may be in the same optical path (the emitted light
passes through the reactor, illuminates the scene, returns through the
reactor, and reaches the detector) or in separate paths (scene and
reactor are independently illuminated and observed, and the controller
correlates their outputs).

\subsubsection{Confocal and descanned coupling}
\label{sec:descanned-coupling}

In some embodiments, confocal or descanned optical paths provide axial
sectioning, spatial filtering, or point-by-point correspondence between
emission and detection, improving signal-to-noise ratio and spatial
resolution.

\subsubsection{Camera-obscura coupling}
\label{sec:camera-obscura}

In some embodiments, an external scene is imaged onto an internal screen
or surface (for example via a pinhole, lens, or relay system), and the
\RK module interacts with the projected image rather than the
scene directly. This provides environmental isolation while preserving
scene-dependent modulation of the reactor's behaviour.

\paragraph{Reality-side spectral filter (non-limiting).}
In some embodiments, a spectrally-selective filter is positioned between
the external scene and the internal screen, such that the reality-to-screen
channel is constrained to a declared spectral band $\Lambda_{\mathrm{ext}}$.
In a canonical embodiment, $\Lambda_{\mathrm{ext}}$ is the visible RGB band.
The reality-side filter has two load-bearing functions: (i) it constrains
the spectral signature of the projected image to a band whose properties
are committed to the protocol digest $\Pi_{\mathrm{dig}}$, and (ii) it
constrains the sensing channel to be unidirectional in the spectral sense,
as further described below. The filter parameters (centre wavelength,
bandwidth, transmission profile, blocking ratio outside band) are committed
to the protocol digest as part of the sensing-layer parameter family.

\paragraph{IR-veracity scan (non-limiting).}
In some embodiments, the \RK module further includes an active
interrogation channel that illuminates its own internal screen with a
probe beam at a declared wavelength $\lambda_{\mathrm{IR}}$ outside
$\Lambda_{\mathrm{ext}}$, in a canonical embodiment an infrared wavelength
selected so that $\lambda_{\mathrm{IR}}$ lies in the blocking band of the
reality-side filter. The descanned detector of
Section~\ref{sec:descanned-coupling} reads the scattered, reflected, or
re-emitted response from the screen at $\lambda_{\mathrm{IR}}$, producing
a \emph{veracity signal} $V_{\mathrm{IR}}(t)$ that recovers geometric,
textural, and material properties of what is physically present on the
screen, independently of and complementary to the visible-band image
$I_{\mathrm{ext}}(t)$ formed by reality. The veracity signal supports
detection of tampering, of projection-driver mismatch, and of missing or
substituted physical substrate at the screen. The IR-veracity scan
parameters --- probe wavelength, scan pattern, detector configuration,
and the rule by which $V_{\mathrm{IR}}(t)$ is compared against
$I_{\mathrm{ext}}(t)$ to produce a verdict --- are committed to the
protocol digest. The probe wavelength, scan duty cycle, and verdict rule
are non-limiting; other interrogation wavelengths and modalities (for
example ultraviolet, terahertz, structured polarisation, or acoustic
probing where consistent with the screen substrate) may be substituted
provided the same load-bearing properties of (i) lying outside
$\Lambda_{\mathrm{ext}}$ and (ii) being blocked by the reality-side
filter are preserved.

\paragraph{One-way isolation of the sensing channel (non-limiting).}
The reality-side filter is positioned such that emissions originating within the \RK module at wavelengths outside $\Lambda_{\mathrm{ext}}$ --- including in particular the IR probe of the IR-veracity scan at $\lambda_{\mathrm{IR}}$ --- are blocked from reaching the external scene through the sensing channel by the filter's blocking band, by ordinary spectral selectivity and optical reciprocity. No claim of nonreciprocal optical isolation is asserted; emissions originating within the module at wavelengths inside $\Lambda_{\mathrm{ext}}$ (where present) are reciprocally transmitted by the filter and are addressed by the disclosure-minimisation and emission-budget policies committed to the protocol digest rather than by the spectral filter alone. The reality-to-screen channel is
unidirectional in the spectral sense: reality drives the screen within
$\Lambda_{\mathrm{ext}}$, and the module's internal interrogation at
$\lambda_{\mathrm{IR}}$ does not exit through the same channel. This
one-way property is structural to the filter geometry and is independent
of whether any actuator subsystem of the \RK module is
connected to or disconnected from external actuation. The one-way
property is committed to the protocol digest as a sensing-layer
property and is verifiable by external measurement against the declared
filter blocking ratio.

\paragraph{External driveability of the screen and kernel-into-kernel
docking (non-limiting).}
In some embodiments, the internal screen of the camera-obscura coupling
is configured so that, in addition to being driven by an external scene
via the reality-side aperture and filter, it may be driven by a separate
external apparatus delivering a projection within $\Lambda_{\mathrm{ext}}$
in place of the physics-driven scene image. Where the separate external
apparatus is itself a \RK of the kind disclosed herein, the
arrangement implements a \emph{kernel-into-kernel docking} configuration
in which a first \RK module senses through its camera-obscura
sensing layer a projection rendered by a second \RK module.
Spectral matching of the rendered projection to $\Lambda_{\mathrm{ext}}$
ensures that the first module's visible-band sensing chain cannot
distinguish a docked-kernel-driven projection from a physics-driven
projection by spectral signature alone. The IR-veracity scan provides
the complementary distinguishability instrument: the docking kernel
either matches the IR-veracity properties of the physical screen
substrate, in which case the docked module's veracity signal is
consistent with physics, or deliberately fails to match them, in which
case the docked module's veracity signal reports the mismatch. Both
matched and unmatched IR-veracity regimes are non-limiting embodiments;
the choice between them is a configuration parameter of the docking
arrangement and is committed to the protocol digest of the docking
session. The kernel-into-kernel docking arrangement is disclosed here
as an apparatus-level property of the camera-obscura sensing layer and
does not depend on any particular use of the docked configuration.

\subsubsection{Quantum-sensitive coupling (optional)}

In some embodiments, the detection subsystem includes single-photon
detectors, photon-number-resolving detectors, or coherent-state
receivers, enabling sensitivity to quantum optical properties of the
channel (for example photon statistics, quantum correlations, or
entanglement witnesses). These capabilities are optional and do not
affect the core Markov-kernel abstraction.

\subsubsection{Optical Anchor and Continuous Physical Anchor Coupling}
\label{sec:optical-anchor-mechanism}

In some embodiments, two or more \RK layers, modules,
partitions, or cascade levels are coupled by a continuous physical
anchor. The continuous physical anchor provides a bounded-latency
physical signal path by which a first layer obtains, challenges,
samples, stimulates, witnesses, or constrains a second layer. The
anchor is not merely an after-the-fact log comparison; it is a
runtime-coupled apparatus pathway whose signal state, timing budget,
and anomaly response are declared in, or bound to, the protocol digest.

In a non-limiting optical embodiment, the continuous physical anchor
comprises an emitter, an optical path, a modulator or framing element,
a target or reactor interface, a detector, a synchronisation source,
and a controller. The emitter may emit at any wavelength or band,
including ultraviolet, visible, near-infrared, short-wave infrared,
mid-infrared, long-wave infrared, terahertz, or a combination thereof.
The emitted signal may comprise structured illumination, pulsed
illumination, phase modulation, amplitude modulation, polarisation
modulation, wavelength modulation, spatial modulation, temporal
coding, frequency coding, spread-spectrum coding, challenge-response
coding, or any combination thereof. The detector may measure intensity,
phase, time-of-flight, spectrum, polarisation, spatial distribution,
temporal waveform, fluorescence, phosphorescence, scattering response,
speckle response, or another response of the anchored layer.

In some embodiments, an anchor challenge $a_t$ is derived from a
protocol state and emitted into the anchored path:
\[
    a_t = A_{\eta}(\Pi_t, C_{0:t}, q_t, s_t),
\]
where $q_t$ may comprise a challenge, beacon value, timestamp,
ledger-derived value, or prior digest, and $s_t$ may comprise a local
state of the anchor controller. A measured anchor response $b_t$ is
then received:
\[
    b_t = B_{\xi}(a_t, X_t, \epsilon_t),
\]
where $X_t$ denotes the physical state of the anchored layer or
interface and $\epsilon_t$ denotes noise, drift, environmental
variation, or measurement uncertainty. An anchor meter may compute:
\[
    m^{\mathrm{anchor}}_t
    =
    M_{\mathrm{anchor}}(a_t, b_t, \Pi_t, \Delta t_t),
\]
where $\Delta t_t$ is the measured or bounded age of the anchor
response. These equations are illustrative only; any challenge,
response, witnessing, synchronisation, or continuous-coupling function
may be used.

In some embodiments, the protocol digest commits to an anchor-latency
budget. The budget may include one or more of propagation time,
emitter rise time, modulator delay, detector integration time,
detector electronics delay, analogue-conditioning delay, digital
processing delay, queueing delay, decision-window length,
synchronisation uncertainty, clock drift, jitter allowance, stale
signal age, path-length bounds, or maximum permitted missing-sample
duration. The apparatus may reject, flag, quarantine, down-weight, or
escalate a signal when an anchor response is absent, stale, delayed,
inconsistent, saturated, clipped, spoof-like, out of distribution, or
outside a declared meter envelope.

In some embodiments, the continuous physical anchor supports cascade
security by making a successful attack on one layer insufficient
unless the attacker also satisfies the bounded-latency physical
coupling constraints at one or more other layers. The anchor thereby
provides a simultaneity-of-attack constraint complementary to
bottleneck-minimum per-layer hardness. The local hardness of a layer
and the cross-layer anchor condition are separate security properties:
the former concerns the difficulty of compromising a layer in
isolation, while the latter concerns maintaining physically consistent
coupling across layers within the declared timing and meter envelope.

In some embodiments, anomaly handling for the continuous physical
anchor is graduated. A first response may increase monitoring cadence,
collect additional diagnostic samples, or enter a soft-gate diagnostic
mode. A second response may restrict release of governance signals,
freeze parameter updates, suspend one or more influence paths, or
require selective opening. A third response may disconnect a cascade
interface, enter a rescue or containment mode, require a governance
authority signature, or trigger mandatory escalation for hard-floor
proximity. The response ladder, confidence thresholds, false-positive
handling, retry policy, and recovery conditions may be declared in
the protocol digest.

\paragraph{Canonical optical-anchor embodiment (non-limiting).}
In one non-limiting canonical embodiment, the continuous physical
anchor comprises a near-infrared LED or laser-diode emitter
operating in a declared band, positioned at a declared standoff
from a reactor interface or anchored-layer surface; a modulator
applying a declared challenge code (for example a
spread-spectrum-coded amplitude pattern with a committed code
length and chip rate); an optical path of declared geometry and
declared maximum path-length variation; a detector comprising a
photodiode and transimpedance amplifier with declared bandwidth
and noise floor; a synchronisation source comprising a stable
oscillator with declared frequency, phase-noise, and clock-drift
parameters; and a controller that derives the per-cycle challenge
$a_t$ from the protocol state and computes the per-cycle anchor
meter $m^{\mathrm{anchor}}_t$. The declared end-to-end anchor
latency budget comprises emitter rise time, modulator delay,
optical propagation time, detector integration and electronics
delay, and decision-window length, all committed to the protocol
digest. The graduated anomaly response of this canonical
embodiment is: first response, increase sampling cadence and
collect additional diagnostic samples; second response, restrict
release of governance signals and freeze parameter updates; third
response, disconnect the anchored cascade interface and require a
governance authority signature for re-engagement. The wavelength,
modulation code, latency budget, response ladder, and other
parameters of this canonical embodiment are illustrative and
non-limiting; other declared parameter choices are within the
scope of the continuous physical anchor as generalised in the
following paragraph.

\paragraph{Anchor-modality generalisation (non-limiting).}
Although optical anchor embodiments are described, the continuous
physical anchor generalises to any bounded-latency physical coupling
having anomaly-as-detection-signal properties. Non-limiting examples
include optical coupling at any wavelength, acoustic coupling at any
frequency, ultrasonic coupling, electromagnetic coupling, radio
frequency coupling, microwave coupling, terahertz coupling,
electrostatic coupling, magnetic coupling, particle-beam coupling,
electron-beam coupling, photon-counting coupling, thermal coupling,
mechanical vibration coupling, fluidic coupling, and hybrid
multi-modal coupling. The anchor may be continuous, sampled,
challenge-response, event-triggered, periodic, asynchronous, or
multi-rate, provided that the declared coupling and timing constraints
are bound to the protocol digest.

% ----------------------------------------------------------------------
\subsection{Scanning and Addressing}
% ----------------------------------------------------------------------

In some embodiments, one or more scanning elements address spatial or
angular coordinates over time. A scan program is represented by scan
coordinates $\mathbf{s}(t)$ and associated control signals, and the
system logs enough information to map time indices to calibrated scene
coordinates and reactor coordinates.

\paragraph{Scan coordinates and calibration (non-limiting).}
In some embodiments, $\mathbf{s}(t)$ denotes a 2D or 3D scan
coordinate, for example $(x(t),y(t))$ on a raster,
$(\theta_x(t),\theta_y(t))$ for galvanometer angles, or a
higher-dimensional coordinate that includes focus or wavelength
selection. A calibration map $M$ relates commanded coordinates to
physical coordinates:
\[
  \mathbf{s}_{\mathrm{phys}}(t)
  = M\bigl(\mathbf{s}_{\mathrm{cmd}}(t),\;\mathbf{m}(t)\bigr),
\]
where $\mathbf{m}(t)$ is a meter vector capturing alignment, gain,
drift, and stability indicators.

\paragraph{Helical and volumetric scan laws (non-limiting).}
In some embodiments, $\mathbf{s}(t)$ includes an internal polar or
cylindrical coordinate such as $(\rho(t),\varphi(t))$ together with an
axial or focus coordinate. The controller chooses how
$(\rho,\varphi,f)$ evolve over time, and calibrated optics map these
emitter-side coordinates to trajectories through a reactor volume. In
a non-limiting cylindrical reactor embodiment, planar probes in
$(\rho,\varphi)$ combined with modulation of $f$ yield approximately
helical trajectories in reactor coordinates, providing a volumetric
sweep that remains characterisable and effectively injective at the
sampling timescale.  The scan-law angular coordinate $\varphi(t)$ is
unrelated to the RT style parameter $\phi \in \Phi$ defined in
Section~\ref{sec:style-parameter}.

\paragraph{Addressing and synchronisation (non-limiting).}
In some embodiments, addressing links a timebase to scene and reactor
coordinates. A controller maintains a common clock that ties: projector
emission indices, detector sample indices, scan mirror positions, and
optional rolling-shutter timing models. Synchronisation error is
captured in meters and used to gate verification, trigger
recalibration, or downgrade to a safer regime.

\paragraph{Carrier-modulated and demodulated probing (non-limiting).}
In some embodiments, control and/or authentication signals are encoded
by modulation of a carrier rather than at baseband, including amplitude,
frequency, phase, or polarisation modulation of an optical, acoustic, or
RF carrier. Demodulation is performed by matched filtering, lock-in
detection, or other synchronous detection aligned to a logged carrier
schedule. Carrier parameters and demodulation settings are recorded in
the protocol digest.

\paragraph{AOD raster and electronic scanning (non-limiting).}
In some embodiments, scan coordinates are implemented by electronically
steered optics, including acousto-optic deflectors (AODs), electro-optic
deflectors, MEMS scanners, galvanometers, spatial light modulators, or
phased arrays. In a non-limiting AOD embodiment, beam deflection is
controlled by radio-frequency drive signals and the scan program is a
time series of commanded frequencies and amplitudes.

\paragraph{Mixed and hierarchical scanning schemes (non-limiting).}
In some embodiments, scanning is mixed or hierarchical, combining coarse
trajectories with fine trajectories. Non-limiting examples include
coarse rasters that locate regions of interest followed by local
high-density scans, or slow patrol scans interleaved with fast
verification micro-sweeps. Such schemes support attention-like scan
policies in \LI and \RT regimes and adaptive
challenge--response sweeps in \TB.

\paragraph{Rolling shutter and temporal mapping (non-limiting).}
In some embodiments, a detector uses rolling shutter and each row (or
block) has an exposure interval. The protocol digest includes the
exposure model, and mapping functions associate each observed sample
with a scan coordinate and an emission interval. In some embodiments,
rolling-shutter readout is configured with row inversion, row/column
pairing, and/or windowing so that exposure windows align to scan dwells.

\paragraph{Global shutter and exposure-window integration
(non-limiting).}
In some embodiments, a detector uses global shutter so that pixels
integrate over a common exposure interval. In some embodiments,
global-shutter acquisition is combined with rolling-shutter acquisition
(for example a rolling-shutter sensor for fine temporal attribution and
a global-shutter sensor for robust integrated baselines).

\paragraph{Tunable distortion elements and e-beam deflection
  (non-limiting).}
In some embodiments, the optical path between a memory surface
(including without limitation a CRT phosphor screen, scanned-laser
target, or persistence medium) and the observing detector includes
a fixed or tunable distortion element, comprising for example a
steering mirror, cylindrical lens, tunable lens, electrowetting
or liquid-crystal lens, mechanically rotated optic, anamorphic
prism pair, or programmable diffractive element. In tunable
embodiments, the distortion element is driven by one or more
analogue or digital control parameters (for example rotation
angle, drive voltage, piezo extension, or applied magnetic field),
allowing the geometric relationship between the memory surface
coordinate system and the detector coordinate system to be varied
continuously or in declared steps over time, with the drive
schedule committed to the protocol digest.
In further embodiments specific to CRT and related electron-beam
reactor substrates, the electron beam itself is subject to one or
more tunable distortion controls applied before the beam reaches
the phosphor, including without limitation: auxiliary electric
deflection plates and electrostatic field configurations applied
in addition to the primary deflection yoke; auxiliary magnetic
deflection coils, focus coils, astigmatism-correction coils, and
convergence coils with independently controllable currents;
quadrupole, sextupole, and higher-order multipole electron-optical
elements; tunable accelerating-voltage and beam-current modulation;
and time-varying yoke drive applied as a trainable kernel parameter
rather than only as a static raster control. Each such electric or
magnetic e-beam distortion control contributes additional trainable
or adaptively-controlled degrees of freedom to the reactor and,
together with the optical distortion element described above,
provides a tunable spatial-and-angular mapping between displayed
or addressed memory coordinates and detected coordinates. The
drive schedule, control parameters, and any commanded changes
across an episode are committed to the protocol digest and \cb on
the same footing as other reactor controls. In some embodiments,
small natural variations in deflection-yoke winding, electric-plate
spacing, residual magnetisation, and shim placement contribute to
the reactor-microstructure unclonability signature
(Section~\ref{sec:security-theory}).
In some embodiments, the disclosed electric and magnetic e-beam
distortion controls are implemented in purpose-built electron-
optical assemblies rather than as modifications to commodity or
aged CRT hardware; in such embodiments, operating envelopes
(including without limitation accelerating-voltage range and slew
rate, deflection-coil and electrode drive bandwidths, beam-current
limits, and safety interlocks) are declared in the protocol digest
and constrained by the physical capabilities of the specific
implementation. In other embodiments, a subset of the disclosed
controls is implemented as a bounded retrofit of commodity or
hobbyist CRT hardware, with the available control envelope
correspondingly narrower and likewise declared.

\paragraph{Operating modes A/B/C (non-limiting).}
In some embodiments, the same hardware supports multiple scanning and
tuning modes. Mode~A (settled raster): scan coordinates follow a settled
$N\times M$ raster with dwell and settle allowances, supporting
repeatable calibration. Mode~B (sampled continuous): scan coordinates
follow a continuous trajectory (for example Lissajous, spiral, or other
smooth curves) and the \cb is sampled under logged timing.
Mode~C (pure analogue tuning): kernel updates are derived from analogue
proxy signals and applied as actuator-safe set-point changes without
explicit digital gradients. Mode transitions are explicitly marked in
the protocol digest and meter summaries.  Modes A/B/C are
scan and tuning modalities that may be used within any of the three
operating regimes (\TB, \LI, \RT); regime
selection determines the objective applied to the resulting
\cb, while mode selection determines how the scan
coordinates and tuning parameters evolve during acquisition.

\paragraph{Terminology note: uses of ``mode'' in this document.}
This document uses the term ``mode'' on three independent axes:
(i)~\emph{scan/tuning modality} (Modes~A, B, C above), which
governs how scan coordinates evolve;
(ii)~\emph{coupling mode} (static vs Yoked,
Section~\ref{sec:regimes}), which governs device--scene feedback
depth; and
(iii)~\emph{training configuration} (adversarial, cooperative, or
mixed training stance, and lifecycle state such as
Embodied/Docked/Offline,
Section~\ref{sec:agent-lifecycle}), which
governs the training context during calibration.  These three axes
are orthogonal: any combination of
scan modality, coupling mode, and training configuration is in
principle valid and is recorded in the protocol digest.

\paragraph{Lissajous galvanometer trajectories and focus modulation
(non-limiting).}
In some embodiments, galvanometer angles follow a Lissajous scan law:
\[
  \theta_x(t)=r_x\sin(2\pi f_x t+\phi_x),
  \qquad
  \theta_y(t)=r_y\sin(2\pi f_y t+\phi_y),
\]
where $(r_x,r_y)$ are amplitudes, $(f_x,f_y)$ are scan frequencies, and
$(\phi_x,\phi_y)$ are phases. In some embodiments, focus is treated as a
primary trainable driven in synchrony with the scan, for example a focus
control $\zeta(t)=\zeta_0+\kappa_{\mathrm{focus}}(t)$, and the focus
schedule is logged so that meter scores condition on the correct focus
modulation regime.

\paragraph{Bidirectional addressing surfaces (non-limiting).}
In some embodiments, addressing surfaces are bidirectional: they can be
written and read. Non-limiting examples include CRT or phosphor surfaces
(with XY deflection and intensity drives) and other media that preserve
a spatiotemporal state.

% ----------------------------------------------------------------------
\subsection{Multi-Channel Observation and Invariance Profiles}
% ----------------------------------------------------------------------

In some embodiments, observations are multi-channel and include
intensity, spectral channels, polarisation, time-of-flight or phase,
and other modalities. The system may also record reduced-dimensional
summaries, field-valued measurements, or feature vectors derived from
raw observations and meters.

\paragraph{Field-valued and reduced-dimensional readouts (non-limiting).}
In some embodiments, an observation at time $t$ is a field
$\mathbf{y}_t(\mathbf{s})$ over scan coordinates, while in other
embodiments it is a reduced-dimensional vector
$\mathbf{y}_t \in \R^d$ obtained by sampling, pooling, or feature
extraction. The choice of readout is logged as part of the protocol
digest.

\paragraph{Observation layers (non-limiting).}
In some embodiments, the system represents observations using a layered
vocabulary: (i)~raw stream (sensor outputs as produced by transducers),
(ii)~conditioned stream (raw outputs aligned to sweep/challenge context
with provenance and calibration identifiers), (iii)~projection stream
(primitive measurement values such as samples, windowed statistics, or
demodulated correlator outputs), (iv)~feature stream (derived
representations for meter evaluation or inference), and (v)~signature
stream (compressed artefacts intended for storage, comparison, or
verification). These layers are conceptual rather than mandatory stages.

\paragraph{Sampling, multi-rate acquisition, and alignment
(non-limiting).}
In some embodiments, different modalities operate at different sampling
rates and may be asynchronous. The conditioned stream includes
timestamps, sweep indices, and alignment metadata, and the system
performs resampling and time-alignment so that cross-channel statistics
are computed under a consistent timebase.

\paragraph{Noise, drift, distortion, and injection taxonomy
(non-limiting).}
In some embodiments, meters account for non-idealities including:
measurement noise (shot, thermal, amplifier, quantisation), timing and
sampling distortion (jitter, clock drift, latency, missing samples),
drift and calibration changes, nonlinearities and saturation, dispersion
and memory artefacts, crosstalk and interference, environmental
perturbations, and adversarial injection or tampering. Meter vectors
include indicators for these categories so that verification and
hardness claims are explicitly conditioned on the observed operating
envelope.

\paragraph{Detector pooling and compressive detection (non-limiting).}
In some embodiments, detector-side pooling compresses high-dimensional
optical interactions into one or a few channels, enabling compressive
operation with single-pixel or near-single-pixel readouts. Non-limiting
examples include pooling a compact photodiode array to a single
time-continuous stream, aligning rolling-shutter integration windows to
scan dwell timing, and using global-shutter frames as integrated
baselines.

\paragraph{MIMO transforms (non-limiting).}
In some embodiments, the module is operated as a multiple-input
multiple-output (MIMO) system with multiple emitters, multiple
detectors, and multiple observed channels. The system may estimate or
learn an effective mixing model that maps emission-control degrees of
freedom to observed channels, and use such a model to design protocols
that concentrate task-relevant information into a reduced set of
channels while still logging the raw multi-channel record for
auditability.

\paragraph{Multi-camera observation of a shared reactor (non-limiting).}
In some embodiments, a single reactor or scene output is observed by
multiple detectors in parallel via beam splitters, fibre splitters, or
relayed imaging paths. Each detector may use a different shutter mode,
spectral bandpass, polarisation analysis, aperture, exposure, or
viewpoint, yielding channels with different invariance profiles. In some
embodiments, the continuous detector stream produced by a swept scan
functions as an analogue ``video feed'' that can be fanned out to
multiple downstream processing or modulation paths.

\paragraph{Invariance profiles (non-limiting).}
In some embodiments, the system characterises an invariance profile that
describes how observations change under controlled transformations, for
example changes in scan density, wavelength, polarisation, focus, or
viewpoint. A non-limiting goal is to design protocols so that
task-relevant quantities are stable to known nuisance transformations
while remaining sensitive to authenticity-relevant microstructure.

\paragraph{Fractional Fourier transform as a concrete multi-channel
example (non-limiting).}
In some embodiments, a tunable optical configuration realises a family
of transforms parameterised by an order $\alpha_{\mathrm{FrFT}}$ that interpolates
between spatial and Fourier-like representations.

\paragraph{Multi-channel constraints for verification and sensing
(non-limiting).}
In some embodiments, the verifier uses cross-channel consistency
constraints, requiring that multi-spectral or multi-polarisation
observations agree with a shared underlying scene model and with
meter-bounded stability.

\subsubsection{Channel separation and side-channel mitigation
(non-limiting)}
\label{sec:channel-separation}

\paragraph{Channel tags (non-limiting).}
In some embodiments, each separated processing pathway is assigned a
channel tag $u_{\mathrm{ch}}$ that is bound into the protocol digest
and challenge record and propagated into meter inputs and \cba
metadata. Channel tags enable channel-specific replay detection, audit,
and calibration.

\paragraph{Multiplexed channel separation (non-limiting).}
In some embodiments, multiple logical channels share the same physical
medium but are separated by time-division, frequency-division, and/or
code-division multiplexing. Channel assignments are logged in the
protocol digest and may be rotated periodically.

\paragraph{Slots and guard intervals (non-limiting).}
In some embodiments, per-channel emission and scoring are scheduled in
fixed-duration slots with guard intervals sized to measured timing
jitter. The device pads or equalises processing so that each slot
occupies constant wall-time, mitigating timing side-channels.

\paragraph{Power equalisation (non-limiting).}
In some embodiments, instantaneous power draw is equalised across slots
by adding controlled dither, dummy operations, or compensating
illumination, so that power-consumption traces are less informative
about channel identity or kernel state.

\paragraph{Shutter-based channel isolation (non-limiting).}
In some embodiments, electro-optical shutters are placed at strategic
points in the optical path, including pre-reactor shutters,
inter-reactor shutters between cascaded stages, and pre-detector
shutters. Shutter timing protocols enable channel separation for
baseline maintenance, style-reset sequences, reference-versus-transform
acquisition, and differential measurements.

\paragraph{Dual-camera reference channel (non-limiting).}
In some embodiments, a beam splitter feeds a primary detector and a
secondary detector that serves as a reference channel. The secondary
detector may include controllable shutter glasses to acquire only
designated time slots or calibration windows.

% ======================================================================
\section{Operating Regimes}
\label{sec:regimes}
% ======================================================================

A \RK supports three operating regimes that share
the same physical operator but differ in what they optimise. Each regime answers
a different question about the scene through the physical channel:

\begin{description}[style=nextline, leftmargin=2em]
  \item[\TB \normalfont{(Illumination)}] The device interrogates
    the scene and asks: \emph{who are you?} In some embodiments, the
    optimised quantity is Fisher information about the device's
    configuration parameters $I(\theta)$, making forgery detectable
    under bounded adversary resources.
  \item[\LI \normalfont{(Perception)}] The device attunes itself to
    the scene and asks: \emph{what is out there?} In some embodiments,
    the optimised quantity is mutual information $I(S; Y)$ between
    scene variables and the \cba output.
  \item[\RT \normalfont{(Creation)}] The device acts on
    the scene and says: \emph{make it look like this.} In some
    embodiments, the optimised quantity is the divergence
    $d(Y, Y^*)$ between output and target, subject to meter envelopes.
\end{description}

The parenthetical labels (Illumination, Perception, Creation) are
non-limiting mnemonic descriptors and are not alternative regime names;
the canonical regime names are \TB, \LI, and \RT throughout this specification.

These three regimes can operate in two coupling modes: \emph{static}
and \emph{Yoked}. In static operation (the default for most
embodiments), the scene is treated as having fixed or slowly varying
properties, and the device senses, verifies, or renders based on
snapshot or time-averaged observations. In Yoked operation, the
feedback loop is tuned so that device and scene are driven toward a
coupled dynamical state: the \cb encodes not only what the
scene \emph{is} but how it \emph{responds} to being driven, and the
resulting joint dynamics are empirically assessed via meters rather
than asserted \emph{a priori}. Yoked operation adds a dynamic dimension to any regime: Yoked
\LI senses dynamical scene properties such as resonances, damping,
and transfer functions; Yoked \TB uses the dynamics of
scene--reactor coupling as a security primitive; Yoked \RT renders patterns that entrain with scene dynamics.

In some embodiments, a single run traverses multiple regimes and
coupling modes (for example static verification followed by Yoked
sensing), with regime changes and coupling-mode transitions explicitly
marked in the protocol digest and meter summaries.

\paragraph{Reactor complexity is application-dependent (non-limiting).}
Not all applications require empirical hardness.  The three-regime
framework and \cba data model apply across the full
reactor complexity spectrum, from identity (standard
camera/projector) through linear media (gain stages, colour filters,
LTI systems with memory) to strongly nonlinear reactors (saturation,
bistability, chaotic mixing, and ultrafast Kerr-effect switching in
epsilon-near-zero media as described in
Section~\ref{sec:embodiments}).  Empirical hardness---the difficulty of
forging \cba traces---arises from reactor nonlinearity
and manufacturing microstructure, and is relevant primarily to \TB and to applications requiring provenance or attestation.  A
\LI operating with a linear or identity reactor is a sensor whose
\cba format, meter infrastructure, and protocol-digest
logging still apply; the device simply offers no forgery resistance.
A \RT with a well-characterised linear medium can
perform physical computation---optical matrix multiplication,
analogue inference, or physical gradient propagation---where
\emph{fidelity} of the input--output relationship matters more than
unclonability.  A multi-reactor network may combine linear and
nonlinear elements, using linear subsystems for computation and
nonlinear subsystems for signing or attesting the result.  The
degenerate cases in Section~\ref{sec:degenerate} define baselines for
hardness claims; they are also legitimate operating configurations
for applications where hardness is not the objective.

% ----------------------------------------------------------------------
\subsection{Operational Stances: Iris, Eve, Demeter, and the Harmonia Coupling Modifier}
\label{sec:stances}
% ----------------------------------------------------------------------

The three operating regimes are not merely distinct objective
functions: each is the stable deployed result of a generative process
governed by a characteristic \emph{operational stance}, and each
continues to operate in that stance once deployed.  Three primary stances are
recognised, plus a coupling modifier (Harmonia), named after archetypes that recur throughout this
specification.

\begin{definition}[Operational stance]
An \emph{operational stance} is a meta-objective that governs how a
\RK loop selects and weights its primary objective at any
given moment.  Three primary stances and one coupling modifier are defined:

\begin{description}[style=nextline, leftmargin=2em]
  \item[Harmonia (coupling stance --- Yoked modifier)]
    The loop is oriented toward joining: it tends the feedback
    between device and scene toward a coupled state
    in which neither partner's output can be fully understood without
    the other.  Harmonia does not replace any of the three primary
    stances; she is a modifier applied \emph{across} the simplex
    $\Delta_2$, orthogonal to the Iris--Eve--Demeter axis.  Any point
    on $\Delta_2$ may be operated in the Harmonia modifier: an
    Eve-hardened \TB operated under Harmonia becomes a channel
    whose forgery resistance is amplified by the
    joint attractor structure; an Iris-expanded \LI operated under Harmonia
    encodes how the scene responds to being driven, not merely what it
    is; a Demeter-tended \RT operated under Harmonia
    entrains scene dynamics toward the target rendering rather than
    projecting onto a passive surface.  Harmonia is the stance of the
    yoke itself: the bond that transforms two separate systems into a
    coupled whole whose joint evidence is stronger than the sum of its
    parts.  The corresponding coupling modifier is the \yd
    $\alpha \in [0, 1]$, with $\alpha = 0$ the static (unmodified)
    case and $\alpha \to 1$ the maximum commanded coupling depth
    (Yoked classification is meter-derived from the CLE/TE
    co-condition, not from $\alpha$ alone).  A distinction is drawn
    between the \emph{declared target} $\alpha_{\mathrm{target}}$
    (the coupling strength the controller schedules and commits to in
    the protocol digest) and the \emph{achieved}
    $\hat{\alpha}$ (a meter-derived empirical quantity computed from
    the \cb).  The protocol digest commits to the target range and
    estimator methodology; the bundle provides the measurements from
    which $\hat{\alpha}$ is estimated.

  \item[Iris (information-maximising stance)]
    The loop is oriented toward maximum clarity about what is actually
    present in the scene or channel.  No adversarial agenda and no
    growth agenda are imposed: the objective is to resolve the world as
    completely as the physical channel permits.  Iris does not judge
    what she sees; she transmits it with full fidelity.  The
    information-maximising stance produces and characterises the
    \emph{\LI} regime.  The corresponding regime weight is
    $\lambda_{\mathrm{L}}$ on the simplex $\Delta_2$; the pure Iris
    stance is the vertex $(0, 1, 0)$.

  \item[Eve (adversarial-hardening stance)]
    The loop is oriented toward worst-case robustness: it selects
    challenges, probes, and operating points that maximise
    distinguishability between genuine physical response and
    adversarial forgery under a declared attacker family and meter
    envelope.  Eve does not punish; she hardens.  A scene-reactor loop
    hardened by the Eve stance becomes a \emph{\TB}: a channel
    whose empirical security is calibrated by systematic adversarial
    pressure at every layer of the architecture
    (Section~\ref{sec:meters-training}).  The corresponding regime
    weight is $\lambda_{\mathrm{TB}}$; the pure Eve stance is the
    vertex $(1, 0, 0)$.

  \item[Demeter (managed-development stance)]
    The loop is oriented toward appropriate development of the scene or
    agent: it expands capability and rendering when the system is
    stable and flourishing, and contracts toward stabilisation when
    risk is elevated.  Demeter is bidirectional and care-giving; she
    tends toward a target rather than interrogating or maximising.  The
    default expression of the Demeter stance applied to a scene is
    stable stylistic projection mapping: controlled rendering that
    maintains fidelity to a target style while adapting to scene state.
    A scene-reactor loop trained by the Demeter stance becomes a
    \emph{\RT}.  The corresponding regime weight is
    $\lambda_{\mathrm{RT}}$; the pure Demeter stance is the vertex
    $(0, 0, 1)$.
\end{description}
\end{definition}

\paragraph{Stances, regimes, and the simplex (non-limiting).}
The three primary stances correspond exactly to the three vertices of the
regime-weight simplex $\Delta_2 =
\{(\lambda_{\mathrm{TB}}, \lambda_{\mathrm{L}}, \lambda_{\mathrm{RT}})
: \lambda_i \ge 0,\, \sum_i \lambda_i = 1\}$ already defined in
Section~\ref{sec:definitions}.  A mixed stance --- any interior point
of $\Delta_2$ --- produces a regime that simultaneously exhibits
properties of two or three canonical stances: for example a loop with
high $\lambda_{\mathrm{TB}}$ and non-zero $\lambda_{\mathrm{L}}$
interrogates the scene adversarially while also extracting maximal
information about scene structure.  The Yoked modifier
(Section~\ref{sec:yoked}) applies independently of stance mixture:
any point on $\Delta_2$ may be operated in static or Yoked coupling
mode.  The Harmonia stance names and governs this orthogonal
dimension: where Iris, Eve, and Demeter govern \emph{what objective}
the loop pursues, Harmonia governs \emph{how deeply the loop is
coupled to its scene partner}.  The full operational state of a
\RK session is therefore characterised by a point on
$\Delta_2$ (the primary stance mixture) together with a \yd
$\alpha \in [0,1]$ (the Harmonia dimension).

\paragraph{Stances as generative processes and as ongoing operation
(non-limiting).}
Each stance is simultaneously (i)~the \emph{training or configuration
process} that produces a deployed regime, and (ii)~the
\emph{ongoing operational character} of that regime at runtime.  A
deployed \TB was Eve-hardened during configuration and
continues to operate in the Eve stance: its runtime challenge
selection, acceptance-boundary calibration, and adversarial
re-evaluation all express the same adversarial-hardening principle
that shaped its initial configuration.  Similarly, a deployed \LI
was Iris-expanded and continues to operate in the Iris stance,
maximising information gain at every step.  A deployed \RT was Demeter-trained and continues to operate in the Demeter
stance, bidirectionally tending scene state toward the target
rendering.

\paragraph{Stances at every scale (non-limiting).}
The same three primary stances (plus the Harmonia coupling modifier) apply at every scale of the \RK
architecture.  At the \emph{device level}, they select the operating
regime as described above.  At the \emph{protocol level}, challenge
selection, sweep design, and disclosure policy may be governed by any
stance or mixture.  At the \emph{fleet and training level}, the
same stance vocabulary applies to agent training trajectories:
adversarial training configurations (Eve stance), cooperative
training configurations (Demeter stance), and undirected exploration
(Iris stance) are vertex instantiations of the same simplex.  The
regime-weight simplex $\Delta_2$, together with the
Harmonia coupling dimension $\alpha$, thus serves simultaneously as
a coordinate system for device operation, protocol design, and
training configuration.

% ----------------------------------------------------------------------
\subsection{Truth Beam (Verification Regime)}
\label{sec:truth-beam}
% ----------------------------------------------------------------------

\paragraph{Illumination.}
\TB is the regime of radical transparency: the device
interrogates the scene with the full force of its physical channel,
and the question it asks is the simplest and hardest---\emph{was this
produced by a genuine physical device under logged conditions?} Like a beam of light that reveals what it touches without being
changed by it, \TB maximises the sharpness of the likelihood
function around the true physical state. The device emits structured
probes, reads the scene's response, and evaluates whether the response
is consistent with a genuine physical interaction or could have been
fabricated by a bounded adversary.

\paragraph{Scope of ``Truth'' (non-limiting).}
Throughout this document, \emph{\TB} and related terminology
(``was this produced by a genuine physical device?,'' ``verification,'' ``Semantic \TB'')
denote \textbf{provenance-consistent physical interaction under
declared threat models}, not semantic truth adjudication of
depicted content.  What the mechanism certifies is a causal
provenance chain: that a \cba trace was produced by a
specific physical device, under logged protocols, within declared
meter envelopes, and is consistent with genuine physical interaction
rather than bounded-adversary fabrication.  It does not certify---and
should not be read as asserting---that an event depiction is
unstaged, unselectively framed, or free of strategic construction by
the scene's participants.  ``Real-but-misleading'' outcomes (genuine
physical traces of deliberately arranged scenes) survive \TB
verification by design, because the device certifies the physics of
the interaction, not the intent behind it.  The Semantic \TB
extension (Section~\ref{sec:semantic-truth-beam}) evaluates consistency of a
semantic claim $q$ with committed evidence, which is a provenance
and consistency check on the claim's evidential support, not an
oracle for propositional truth.

In \TB embodiments, the \cb is used for authenticity
assessment. A verifier evaluates whether an observed record is
consistent with the physical channel induced by the device, the logged
protocol, and the meter envelope, and rejects records that are likely
to be forged, replayed, or materially altered.

In some embodiments, \TB uses a verisimilitude or consistency
model operating on the \cb or derived features, optionally
together with meter vectors and protocol digests, to produce an accept
or reject decision, a confidence score, or a continuous hardness score.
In some embodiments, this score is used for verification-gated actions,
selective opening escalation, or trust-weight updates.

\subsubsection{Noise and drift envelope (non-limiting)}

In some embodiments, \TB explicitly models an operating envelope
that includes shot noise, read noise, and slow drift. A non-limiting
measurement model for a scalar sensor channel is
\[
  y_t = (1+\delta_t)\,\mu_t + n^{\mathrm{shot}}_t
  + n^{\mathrm{read}}_t,
\]
where $\mu_t$ is the expected response under the protocol and scene,
$\delta_t$ is a slow multiplicative drift term,
$n^{\mathrm{shot}}_t$ is a shot-noise component, and
$n^{\mathrm{read}}_t$ is a read-noise component. In some embodiments,
shot noise is modelled as a Poisson-like component whose variance scales
with signal level, while read noise is modelled as an approximately
additive Gaussian component dominated by electronics and digitisation.

In some embodiments, drift is modelled as a slowly varying process, for
example
\[
  \delta_{t+1} = \rho_{\mathrm{AR}}\,\delta_t + w_t,
\]
with $\rho_{\mathrm{AR}}$ near one and $w_t$ a small disturbance term. In some
embodiments, meters include drift indicators, gain estimates, timing
alignment indicators, and stability flags so that verification claims
are conditioned on an observed envelope rather than asserted
unconditionally.

In some embodiments, an extended noise model includes slow
multiplicative drift combined with affine jitter:
\[
  y_t = a_t(\mu_t + b_t) + n_t,
\]
where $a_t$ is a slow gain drift and $b_t$ is a small offset jitter,
both estimated or bounded by meters, and $n_t$ is a non-limiting
residual noise term.

\subsubsection{Hardness and attacker models (non-limiting)}

In some embodiments, hardness is reported as an empirical difficulty of
producing a forged record that passes the verifier under a bounded
attacker model. For example, an attacker may be bounded by query budget,
compute budget, model class, or training data access. In some
embodiments, these notions connect to the digital hardness index
$k^*_{\mathrm{dig}}(\theta;\varepsilon)$ defined elsewhere in this
document, with hardness claims interpreted as conditioned on the
protocol digest, meter envelope, and a stated attacker family.

In some embodiments, hardness is strengthened by serial composition of
challenges, protocol randomisation, cross-channel constraints, and
selective opening checks that force consistency not only in local
windows but also across time and across coupled channels.

\subsubsection{Semantic Truth Beam (non-limiting)}
\label{sec:semantic-truth-beam}

In some embodiments, \TB is applied to hard inference
verification, where the object being verified is a semantic claim or
inference output rather than only raw sensor traces. For example, a
device or agent may assert a claim $q$ about a scene (a label, a
structured description, a detected event, or a property estimate)
together with a commitment to the evidence window and protocol digest
used to produce that claim.

In some embodiments, a semantic claim is verified by routing challenges
to physically grounded evidence that is relevant to the claim. For
example, the verifier selects protocols or openings that are conditioned
on the claim type, the claimed region of interest, or a semantic routing
rule, and then checks consistency of opened windows with the claim, the
protocol digest, and the meter envelope. In some embodiments, the same
mechanism is used to verify that an inference is not a post hoc
fabrication by requiring that the claim be bound to committed evidence
prior to the opening seed.

In some embodiments, semantic \TB includes explicit grounding
mechanisms that bind semantic outputs to committed evidence and make
certain classes of semantic forgery more difficult:
\begin{enumerate}[nosep]
  \item \emph{Semantic routing through the physical kernel:} candidate
    semantic claims (or compact digests thereof) are re-encoded as
    emission patterns and passed back through a \RK run, and
    the resulting response is compared against a reference derived from
    the original measurement or claimed evidence window.
  \item \emph{Inference commitment:} a semantic extraction head commits
    to intermediate artifacts (for example attention maps, feature
    activations, derivation traces) using commitment mechanisms
    described herein.
  \item \emph{Spatiotemporal grounding:} each semantic claim is
    annotated with specific scan coordinates, time indices, and/or
    disclosed atoms in the \cb that support the claim.
  \item \emph{Adversarial semantic hardness:} semantic extraction heads
    and verifiers are trained adversarially so that producing an
    alternative plausible-looking semantic description that passes the
    verifier is empirically difficult for a bounded adversary.
\end{enumerate}

In some embodiments, projection hardness is assessed via a
challenge--commit--open flow tied to emission digests: a verifier issues
a challenge seed; the device deterministically selects an emission
schedule under the seed-driven protocol, commits to an emission root
$R^{\mathrm{emit}}$, executes the physical run, commits to an
observation root $R^{\mathrm{obs}}$, and later opens selected emission
and observation atoms with inclusion proofs for meter-conditioned
consistency checks.  The commitment binds the full protocol digest
(including scan-law parameters and emission schedule), the declared
meter envelope $\mathcal{M}_{\mathrm{accept}}$, and the protocol
seed(s) used to derive the emission schedule (following the two-seed
structure of Section~\ref{par:two-seed} in some embodiments), so that
openings can reconstruct and verify the full run configuration.

\subsubsection{Graduated verification levels (non-limiting)}
\label{sec:graduated-verification}

In some embodiments, the same device supports multiple verification
security levels by controlling which regions of the configuration space
$\Theta$ the challenge policy probes:

\begin{description}[style=nextline, leftmargin=2em]
  \item[Level~0 (ambient).] Static challenge-response with no feedback.
    The device behaves as a microstructure challenge-response oracle (in the sense used in the PUF literature). The challenge policy
    operates in low-curvature regions of the attractor bundle where
    attractor geometry changes slowly with $\theta$. This level is fast,
    cheap, and suitable for low-security applications.
  \item[Level~1 (trajectory).] A short closed-loop trajectory in a
    moderate-curvature region. The verifier checks attractor basin
    identity over a brief feedback window. In some embodiments, spoofing
    would require matching
    transient dynamics, not just steady-state responses.
  \item[Level~2 (bifurcation).] The challenge policy sweeps near a known
    bifurcation surface in $\Theta$. Fisher information is concentrated
    near bifurcations, so spoofing would require tracking the exact topology
    change. Hardness increases sharply with proximity to the bifurcation
    set~$B$.
  \item[Level~3 (holonomy).] A full cyclic sweep around a bifurcation,
    with accumulated geometric phase measured from the \cb.
    If coexisting attractors give rise to non-Abelian holonomy, the
    accumulated phase is path-dependent and is conjectured to be
    significantly harder to forge from finite surrogate data, though
    the precise scaling with path complexity remains an open problem.
    Holonomy-based hardening is optional and is not required for
    core \TB, \LI, or \RT operation;
    Level~1 and Level~2 verification are independently enabled.
\end{description}

In some embodiments, the verification level is selected per session
based on the required security level, and the selected level is recorded
in the protocol digest. In some embodiments, levels are composed
sequentially (a fast Level~0 gate followed by a slower Level~2 deep
check if the initial gate passes). The cost of verification scales with
the curvature of the interrogation path through $\Theta$, not with
device complexity: a single device serves all levels.

\subsubsection{Honeypot and canary verification (non-limiting)}
\label{sec:honeypot}

In some embodiments, the verification protocol deliberately operates in
a regime where a specific attacker model class appears to succeed, while
a secondary check detects the attack:

\begin{enumerate}[nosep]
  \item The verifier presents a low-curvature challenge that a linear or
    low-order surrogate model can match.
  \item Simultaneously, the verifier monitors higher-order structure in
    the \cb (for example transfer entropy spectra, speckle
    higher-order statistics, or frequency-resolved causal signatures).
  \item A genuine device passes both checks. An attacker's surrogate
    passes the first but fails the second because the higher-order
    structure is expected to require, under the declared attacker
    families, reproducing the full nonlinear device dynamics.
  \item The pattern of what the attacker matches and what the attacker
    fails to match reveals the attacker's model class and capacity.
\end{enumerate}

In some embodiments, honeypot verification inverts the usual security
strategy: instead of making every challenge maximally hard, some
channels are deliberately easy in order to create a diagnostic surface
that characterises attacker capabilities. The diagnostic output is
logged in the meter envelope.

\subsubsection{Anti-Truth Beam: provenance negation (non-limiting)}

In some embodiments, the system provides a complement to positive
verification: the ability to establish, under declared attacker
families and calibrated error tolerances, that an observation was
\emph{not} produced by any device in a declared family.

In some embodiments, the attractor bundle of a device family defines a
region in observable space. Anti-\TB checks whether a submitted
observation lies outside this region. The boundary of the region is
determined by the bifurcation set~$B$---the same structure that makes
positive verification hard also makes negative verification
operationally feasible, because the bifurcation set partitions the observable
space into topologically distinct basins.  As with all empirical
hardness claims in this document, negative-provenance decisions are
conditioned on the current attacker class and fleet-calibrated
acceptance/error rates, and are time-indexed rather than
unconditionally durable (see Section~\ref{sec:security-theory}).

Non-limiting applications include: detecting synthetic or manipulated
media that was not produced by a physical \RK (deepfake
detection); determining that a submitted measurement did not originate
from an enrolled device (provenance negation); and identifying devices
whose physical configuration has changed (for example through tampering
or ageing that moves the device's operating point across a bifurcation
boundary).

% ----------------------------------------------------------------------
\subsection{Limager (Perception and Sensing Regime)}
\label{sec:limager}
% ----------------------------------------------------------------------

\paragraph{Perception.}
\LI is the regime of receptivity: the device attunes itself to the
scene and asks \emph{what is out there?} Like a still pool whose
surface reflects the sky precisely because it does not impose its own
shape upon it, \LI maximises the information flowing from scene to
output. The device actively selects scan laws, integration times,
coherence regimes, and attention policies to maximise the information
flowing from scene to output, much as a listening ear turns toward a
sound source.

In \LI embodiments, the \RK is used for perception and
sensing. The system selects protocols and processes the \cb to
infer latent scene variables, reconstruct geometry or materials,
estimate state, or optimise information acquisition. \LI can
operate with or without supervised labels, and supervised labels may be
used at any granularity.  In some embodiments, \LI is operated as
a standalone boundary-state estimation and calibration method
(Section~\ref{sec:boundary-state-estimation}), independent of
verification or rendering objectives.

\subsubsection{Latent scene variables and inference targets
(non-limiting)}

In some embodiments, the inferred variable is a latent scene description
$\xi$ (distinguished from the extended state $X_t$; see
Section~\ref{sec:definitions}), for example geometry, pose, reflectance parameters, spectral
parameters, or a learned representation. Given observations $C_{0:T}$
and a protocol $U_{0:T}$, inference is performed by a model that
estimates $\xi$ or a distribution over $\xi$, optionally conditioned on
meters and protocol digests.

\subsubsection{Information maximisation objective (non-limiting)}

In some embodiments, \LI selects or adapts protocols to maximise
information gain about $\xi$ under meter and safety constraints. A
non-limiting objective is:
\[
  \max_{\mathsf{P}_{U_{0:T}}\in\mathcal{U}}
  I_\theta\!\left(\xi;\; C_{0:T}\mid U_{0:T}\right),
\]
where $\mathcal{U}$ encodes admissible protocols (for example eye-safety
bounds, thermal bounds, scan-rate bounds, and stability bounds). In some
embodiments, tractable proxies are used, for example
mutual-information lower bounds, predictive uncertainty reduction, or
reconstruction-error reductions measured on committed windows.

\paragraph{Sparse, compressed-sensing, and plug-and-play reconstruction
(non-limiting).}
In some embodiments, \LI reconstruction is performed by sparse or
compressed-sensing recovery under declared sparsity, incoherence, or
measurement-matrix assumptions (Cand\`es--Tao 2006; Donoho 2006), by
model-based iterative reconstruction, or by plug-and-play priors
(Venkatakrishnan, Bouman, and Wohlberg 2013) in which a denoiser or
score-like prior is inserted into the reconstruction loop. The
reconstruction algorithm, regularisation weights, denoiser identity,
training-corpus digest, stopping rule, and residual diagnostics are
committed to the protocol digest, and reconstructions used for
verification are gated by the same meter and selective-opening policies
as other \LI-derived outputs.

\paragraph{Geometric reconstruction and registration tools
(non-limiting).}
In some embodiments, \LI or digital-twin state estimation composes
standard geometric estimation tools, including structure-from-motion,
visual odometry, simultaneous localisation and mapping (SLAM), bundle
adjustment, iterative closest point (ICP), photometric registration,
optical-flow estimation, or differentiable rendering-based pose
refinement. These tools act as optional estimators of scene geometry,
pose, correspondence, or motion, while the inventive apparatus remains
the committed, meter-bounded \RK loop. The estimator family,
feature representation, correspondence policy, robust-loss function,
outlier rejection rule, and validation residuals are recorded with the
derived-scene-state identifier.

\subsubsection{Joint exposure and flash control for HDR imaging
(non-limiting)}

In some embodiments, camera exposure parameters (for example gain and
shutter time) and projected flash or illumination patterns are jointly
selected by the protocol to construct a high dynamic range
representation of the scene. Brighter regions may be probed with
reduced exposure or attenuated emission, while darker regions receive
stronger illumination or longer exposure. In some embodiments, a
latent state $Z_\theta(t)$ fuses the resulting emission--observation
tuples into an HDR scene estimate.

\subsubsection{Learned scan trajectories (non-limiting)}

In some embodiments, scan trajectories are selected by a learned policy
that conditions on recent observations and meters:
\[
  u(t) \sim \pi_\psi(\cdot \mid o_{0:t}, \mathbf{m}_{0:t}),
\]
where $o_t$ denotes RK-derived observations or summaries. In some
embodiments, the policy is constrained to a safe admissible set and is
logged through a protocol digest so results remain auditable.

\subsubsection{Surrogate models and cross-reactor transfer
(non-limiting)}

In some embodiments, a surrogate model predicts outcomes of physical
trials, enabling optimisation when direct gradients are unavailable. In
some embodiments, transfer learning is performed across reactors using
low-rank adapters, device-conditioned modulation, or universal encoders
with device-specific heads, with drift compensation and meter-bounded
constraints to preserve stability and auditability.

\subsubsection{Inner Limager: self-calibration (non-limiting)}
\label{sec:inner-limager}

In some embodiments, \LI is turned inward: the ``scene'' is the
device's own configuration state $\theta$, and the objective is to
maximise information about the device's current parameters rather than
about an external environment. This self-calibration mode uses the same
mathematical machinery as standard \LI but with the roles of device
and scene exchanged. Inner \LI asks not ``what is out there?'' but
``where in configuration space is this device right now?''

In some embodiments, Inner \LI uses calibration targets or reference
signals as controlled ``scenes'' and estimates drift, aging, and
parameter shift by treating the device's own response statistics as the
inference target. The resulting state estimates feed into aging-worldline
and drifting-field mechanisms (a learned gradient field for steering
device configuration $\theta$ toward target operating regions; defined
fully in Section~\ref{sec:drifting-field}) described in the control
theory sections.

\subsubsection{Boundary-state estimation and Fisher-information diagnostics
as a standalone method (non-limiting)}
\label{sec:boundary-state-estimation}

In some embodiments, the \LI regime is operated as a standalone
boundary-state estimation method, independent of verification or rendering
objectives.  The method takes as inputs a device operating under a
declared protocol and meter envelope, together with a set of
perturbation directions in configuration space $\Theta$, and produces
as output a boundary-state estimate, a Fisher-information diagnostic,
and optionally a calibration update.

A non-limiting procedure is as follows:
\begin{enumerate}
  \item \textbf{Boundary perturbation.}  The control protocol
    $U_{0:T}$ is swept along one or more directions in $\Theta$
    chosen to approach a boundary of the current attractor basin
    (for example a bifurcation surface, a stability boundary, or a
    coupling-transition surface).  The sweep direction and step size
    are recorded in the protocol digest.
  \item \textbf{Measurement.}  The \cb $C_{0:T}$ is recorded for each
    perturbation step.  Meter summaries (alignment, drift, stability)
    are computed and logged.
  \item \textbf{Estimator construction.}  A boundary-proximity estimate
    $\hat{b}(\theta)$ is computed from the observed response statistics,
    for example as the ratio of within-basin trajectory length to the
    total sweep length, or as a signed distance to the inferred
    bifurcation surface under a declared model class.
  \item \textbf{Fisher-information diagnostic.}  The Fisher information
    $I(\theta)$ about the boundary-state parameter is estimated from
    the \cb using the declared estimator family (for example via the
    score function or via finite-difference approximations of the
    log-likelihood gradient).  Where the Fisher information matrix
    $\mathcal{F}(\theta)$ is computed for multi-dimensional $\theta$,
    Cram\'er--Rao-style lower bounds on estimation variance are
    derived and reported alongside the estimate.
  \item \textbf{Calibration update (optional).}  In some embodiments,
    the boundary-state estimate is used to update a calibration
    record for the device, for example adjusting a drifting-field
    waypoint or updating an aging-worldline entry.  This update
    constitutes the terminal output of the method and does not require
    a verification decision.
\end{enumerate}
In some embodiments, the protocol digest for a boundary-state
estimation run includes the following items (or their cryptographic
hashes, where confidentiality of intermediate outputs is required):
the perturbation direction set $\{v_j\}$ and step schedule; the
estimator family identifier, version, and declared hyperparameters
(for example bandwidth, kernel class, finite-difference step); the
computed boundary-state estimate $\hat{b}(\theta)$; the Fisher
information matrix $\mathcal{F}(\theta)$ or its principal components;
any Cram\'er--Rao lower bounds produced; the meter summaries from
step~2; and, where step~5 is executed, a version pointer or hash
of the resulting calibration record so that the post-run calibration
state can be reconstructed by an auditor.  These commitments are
sufficient for an independent party to re-execute the method on
the same inputs and verify the reported outputs.
In some embodiments, the output of this method is a calibration update
or state estimate rather than a verification decision, supporting use
cases in which boundary-state knowledge is sought for its own value
(for example during commissioning, periodic health monitoring, or
fleet-level calibration scheduling) rather than as part of an
authentication or \TB workflow.  The boundary-state estimation
method is enabled independently of the \TB and \RT
regimes and does not require those regimes to be active.

\subsubsection{Inner Creation: self-directed state sculpting
(non-limiting)}
\label{sec:inner-creation}

In some embodiments, \RT is turned inward: the
``scene'' is the device's (or agent's) own reactor state, and the
objective is to minimise the distance between the current internal
configuration and a target internal configuration.  Where Inner
\LI asks ``where am I in configuration space?'', Inner Creation
asks ``how do I get \emph{there} in configuration space?''  The same
feedback machinery used by outward-facing \RT
(emit, observe residual, update control signal) applies, but the
emission acts on the reactor's own dynamics rather than on an
external scene.  Inner Creation is the intentional counterpart of the
drifting-field management problem
(Section~\ref{sec:control-theory}): where the drifting field treats
parameter drift as a disturbance to be tracked and bounded, Inner
Creation treats it as a control action to be steered toward a
declared target configuration.

\paragraph{Amplifying internal differences (non-limiting).}
In some embodiments, the primary use of Inner Creation is to
\emph{differentiate} nearby attractor basins.  When a reactor's
state space contains two behavioural modes separated by a shallow
basin boundary---for example two equilibria that produce nearly
indistinguishable output statistics---Inner Creation drives the
system to deepen the boundary, widen the basins, and make the two
modes robustly distinguishable.  The mechanism uses the natural
gradient of the drifting-field framework
(Section~\ref{sec:control-theory}), but with an objective that
rewards basin depth, inter-basin distance, or attractor stability
rather than external scene fidelity.

In some embodiments, Inner Creation is used in conjunction with
calibration mechanisms described in the related Filing 2 application
(optional, non-essential) to consolidate attractor-basin separation
along declared operational axes.

\paragraph{Inner Creation and reactor refresh (non-limiting).}
In some embodiments, the reactor refresh mechanism described in the
hardness management section (Section~\ref{sec:security-theory}) is
implemented as controlled Inner Creation: the drifting-field
controller drives the reactor into a new region of state space,
re-establishing the microstructure complexity on which empirical
hardness depends.  Consistent with the meter partition
(Section~\ref{sec:security-theory}), the refresh controller's
objective is drawn from acceptance meters only; the decision to
initiate a refresh may be triggered by periodic offline hardness
reassessment (which may consult evaluation meters), but the refresh
trajectory itself does not optimise against evaluation-meter
values.  Post-refresh hardness is re-estimated as a new
time-indexed snapshot under the current attacker class.
Adjoint reactor pairs
(Section~\ref{sec:regimes}) may serve as the gradient-propagation
channel for Inner Creation, using the adjoint path to compute
how to change the reactor's own parameters.

\paragraph{Inner \cbs (non-limiting).}
In some embodiments, Inner Creation and Inner \LI produce
\emph{inner \cbs}: committed, auditable records of
the reactor's internal state trajectory, attractor occupancy, basin
boundary proximity, and parameter drift, captured through the same
protocol-digest and meter infrastructure used for outward-facing
regimes.  In some embodiments, inner
bundles record what the device
\emph{is becoming} (internal dynamics) rather than what it
\emph{produces} (external behaviour).  In some embodiments, inner
bundles serve as a training signal for optional monitoring
mechanisms (described in the related Filing 2 application; optional), enabling detection of
attractor-level drift rather than relying solely on
behavioural signatures.

In some embodiments, a non-limiting inner-trace summary over a time
window $[t_0,t_1]$ is represented as
\[
\mathcal{I}_{t_0:t_1}
\equiv \bigl(\phi_{\mathrm{inner}}(C^{\mathrm{inner}}_{t_0:t_1}),\;
  \theta_{t_0:t_1},\;
  \mathrm{basin}(t),\;
  \nabla_\theta J(t)\bigr),
\]
where $\phi_{\mathrm{inner}}$ denotes features extracted from inner
\cba atoms, $\theta_{t_0:t_1}$ is the parameter
trajectory observed by Inner \LI,
$\mathrm{basin}(t)$ is an attractor-occupancy indicator, and
$\nabla_\theta J(t)$ is the gradient signal (when available from
Inner Creation or adjoint-pair propagation).  Inner-trace summaries
are committed and auditable through the same protocol-digest and
meter infrastructure as external \cbs.

\paragraph{Inter-reactor and extended applications of Inner Creation
(non-limiting).}
In some embodiments, Inner Creation is directed at inter-reactor
coupling and encoded inner-state artifacts. These applications
are described in the related Filing 2 application (optional, non-essential).
In some embodiments, Inner Creation is rate-limited,
protocol-restricted, and logged, and inner-trace-based evaluations
are complemented by independent witness checks.

% ----------------------------------------------------------------------
\subsection{Reality Transform (Controllable Rendering Regime)}
\label{sec:reality-transform}
% ----------------------------------------------------------------------

\paragraph{Creation.}
\RT is the regime of purposeful action: the device reaches
into the scene and reshapes it. Like an artist whose brush strokes
respond to the canvas's texture and the paint's resistance, \RT minimises the distance between what is observed and what is
desired. The question is not \emph{what is there} but \emph{what should
be there.} What the device does next depends on what it observes the
scene becoming under its influence---in general, blind projection
without feedback may not converge, compensate for surface properties,
or respect the physical channel's constraints.

In \RT embodiments, the system controls emissions so that
the physical scene or sensed appearance exhibits a controllable
transformation. The transformation may be transient (present only during
emission) or may persist in a medium (for example phosphorescent,
photochromic, thermochromic, or other stateful surfaces), and can be
designed for human consumption or for downstream machine perception.

In some embodiments, \RT is implemented as feedback
projection mapping: the system emits patterns, observes the result, and
updates subsequent emissions to converge toward a target aesthetic,
constraint set, or calibrated appearance. In some embodiments, this loop
is constrained by meters and protocol digests so rendering remains
within empirically supported operating regions.

\subsubsection{Transient versus persistent scene effects (non-limiting)}

In some embodiments, the transform is transient and exists only during
the emission window. In other embodiments, the transform changes a
stateful medium so effects persist after emission, with persistence time
captured by meters and logged protocols. In some embodiments,
persistence is treated as a memory channel and is used to encode or
stabilise transformations across time.

\subsubsection{Constraint coupling to verification and sensing
(non-limiting)}

In some embodiments, \RT is coupled to verisimilitude and
projection-hardness constraints so transformations do not drift into
unstable regimes. For example, a transform policy may include penalties
or gating based on meter envelopes, cross-channel consistency, or
learned verisimilitude discriminators.

\subsubsection{Compatibility with semantic conditioning (non-limiting)}

In some embodiments, text prompts, semantic tokens, or other
high-level conditioning signals are used to select or modulate emission
policies so that physical rendering expresses a requested style or
content while remaining within a meter-bounded operating envelope. This
includes embodiments in which a structured illumination pattern, scan
law, or projection mapping is chosen as a function of semantic
conditioning, rather than being fixed in advance.

In some embodiments, a prompt digest and model identifier are logged as
part of a protocol digest, and the resulting \cb and meter
summaries are committed and made available for selective disclosure.

\subsubsection{Style parameter: formal definition and generality
(non-limiting)}
\label{sec:style-parameter}

In some embodiments, the target behaviour of the \RT
regime is specified by a \emph{style parameter} $\phi \in \Phi$,
where $\Phi$ is a declared style space.  The style parameter $\phi$
is a formal input to the RT control objective and is logged in the
protocol digest alongside the run seed, meter envelope, and emission
schedule.

A non-limiting RT objective incorporating the style parameter is:
\[
  J_{\mathrm{RT}}(U_{0:T};\phi)
  = -\mathbb{E}\Bigl[\sum_{t=0}^{T-1}
      \Bigl(d\!\bigl(\mathbf{y}_t,\, \mu_\phi(x_t, u_t)\bigr)
    + \lambda_{\mathrm{meter}} \cdot \mathbf{1}[\mathbf{m}(t) \notin \mathcal{M}_{\mathrm{accept}}]\Bigr)
    \Bigr],
\]
where $d(\cdot,\cdot)$ is a declared divergence or distortion metric
(dimensionless in preferred embodiments, e.g.\ a normalised
cross-entropy or cosine distance), $\mu_\phi$ is a target response
function parameterised by $\phi$, $\mathcal{M}_{\mathrm{accept}}$ is
the declared meter envelope
(Definition~\ref{def:meter-envelope}), and $\lambda_{\mathrm{meter}}$
is chosen on the same declared scale as $d$ so that all terms in the
per-step sum are commensurate.  The style parameter
$\phi$ may be instantiated in many non-limiting ways, including:
\begin{itemize}
  \item a semantic or aesthetic token (text prompt, embedding, or
    structured label) specifying a visual or perceptual target;
  \item a constraint or policy specification (for example a declared
    operating envelope or compliance target);
  \item a physical rendering target (spectral distribution, spatial
    texture, temporal modulation pattern);
  \item a calibration target (reference scene appearance, reference
    device response distribution);
  \item a composite objective combining multiple of the above.
\end{itemize}
In each case the formal structure --- a declared $\phi \in \Phi$
logged in the protocol digest, an objective $J_{\mathrm{RT}}(\cdot;\phi)$
maximised by control, and a committed \cb recording the result ---
is the same.  Any domain-specific conditioning is
therefore one specialisation among many; the general mechanism is
defined and operable independently of any specific conditioning
domain.
This allows a verifier to audit that a rendered outcome is consistent
with the logged conditioning, the chosen protocol regime, and the
observed envelope, and to distinguish physically grounded rendering
from post-processed fabrication.

In some embodiments, semantic conditioning is combined with
verisimilitude or projection-hardness constraints. This steers
requested transformations toward styles and operating regimes for which
the system remains stable, auditable, and within the training envelope.


\paragraph{Committed primitive-set scene representations
(non-limiting).}
\label{par:committed-primitive-set-scene-representations}
In some embodiments, the style parameter $\phi$ is a committed,
finitely parameterised primitive-set scene representation
comprising a set of spatial primitives $\{p_k\}_{k=1}^{K}$, each
primitive $p_k$ specifying at least a position in a declared
world coordinate frame, a shape or support (for example an
anisotropic covariance $\Sigma_k \in \mathbb{R}^{3\times 3}$ for
covariance-bearing primitive families), an appearance descriptor
(for example an RGB triple or spherical-harmonic coefficients
$c_k$), and a per-primitive weight or opacity $\beta_k$.
Non-limiting primitive families include anisotropic
three-dimensional Gaussian splats of the kind introduced by
Kerbl \emph{et al.}\footnote{B.~Kerbl, G.~Kopanas,
T.~Leimk\"uhler, and G.~Drettakis, ``3D Gaussian Splatting for
Real-Time Radiance Field Rendering,'' ACM Transactions on
Graphics 42(4), 2023.}, point-based primitives with learned
blending, and textured mesh primitives. Related committed
scene-representation families --- including signed distance
fields and neural radiance fields, where embodied by finite
network weights, grids, tables, or other finite serialised
parameter sets --- supply alternative finitely parameterised
representations of $\phi$ and are treated here as non-limiting
alternatives to primitive-set representations. In such
embodiments, a canonical serialisation of $\phi$, or a
content-addressed pointer thereto, is committed by a digest
$h(\phi)$ recorded in the protocol digest; the canonical
serialisation declares primitive ordering or an
ordering-invariant hashing rule, coordinate-frame identifiers,
units, floating-point or quantisation conventions,
appearance-code versions, and any representation-complexity
envelope. The committed digest is logged alongside the run
seed, meter envelope, and emission schedule $U_{0:T}$.
Primitive-set embodiments are one implementation of the target
response function $\mu_\phi$ and do not limit $\Phi$ to spatial
or rasterised targets.

\paragraph{Physically calibrated rasterisation forward model for
primitive-set targets (non-limiting).}
In some embodiments in which $\phi$ is a committed primitive-set
representation and the declared target is expressed as a spatial
or rasterisable scene, the target response function
$\mu_\phi(x_t, u_t)$ is realised as the composition of
(a)~differentiable rasterisation of $\{p_k\}_{k=1}^{K}$ from the
viewpoint or scan geometry specified in $u_t$, under declared
world-to-device coordinate transforms, producing a rendered
radiance, luminance, or detector-channel field
$\tilde{L}_\phi(x,y,t)$;
(b)~spatial convolution with the device-specific spot
point-spread function $K_{\mathrm{spot}}$ of
Section~\ref{sec:crt-phosphor} (or the embodiment-appropriate
device PSF);
(c)~temporal convolution with the effective persistence kernel
$p_{\mathrm{eff}}(t) = p(t)\ast w_{\mathrm{det}}(t)$; and
(d)~detector integration under declared scan geometry and
detector settings, producing the predicted detector observation
$\hat{\mathbf{y}}_t$. In such embodiments, the distortion term
in $J_{\mathrm{RT}}$ (Section~\ref{sec:style-parameter}) may be
evaluated as $d(\mathbf{y}_t, \hat{\mathbf{y}}_t)$ while the
meter-envelope penalty and run-level objective structure of
Section~\ref{sec:style-parameter} are preserved. The
physical-channel parameters $K_{\mathrm{spot}}$,
$p_{\mathrm{eff}}$, $w_{\mathrm{det}}$, the world-to-device
coordinate transforms, and the calibration-table identifiers
and versions are drawn from independently enrolled or metered
calibration records associated with the device configuration
representation $\mathbf{d}_i$ and, where applicable, the
declared challenge configuration $X$ or protocol geometry; they
are referenced or committed in the protocol digest, and are
not fitted as free parameters during primitive-set
optimisation.

\paragraph{Factorisation of scene hypothesis and device nuisance
(non-limiting).}
In some embodiments in which the declared target is an
externally represented scene or a declared spatial or
radiometric proxy scene, committed primitive-set representations
carry the scene hypothesis in $\phi$ while the device-specific
forward-model parameters remain in $\mathbf{d}_i$. For
covariance-bearing primitive families, the per-primitive
covariance $\Sigma_k$ describes scene-space extent in world
coordinates and is distinct from the device's imaging PSF
$K_{\mathrm{spot}}$, which describes device-space blur; the two
compose in the forward model but are logged and verified
separately. This factorisation permits one non-limiting \LI
inverse-problem formulation as joint inference over $\phi$ and
$\mathbf{d}_i$ given $\mathbf{y}_t$, with $\phi$ and
$\mathbf{d}_i$ drawn from distinct priors and constrained by
distinct meters. The factorisation is not asserted for
inward-facing or self-calibration embodiments in which the
``scene'' is the device's own configuration state, for
biological-reactor embodiments except to the extent that a
declared spatial or radiometric proxy-scene representation is
used, or for non-image passive embodiments such as GNSS in
which the observation is expressed in receiver-observable
coordinates rather than as rasterised imagery. The factorisation
is also not asserted in Yoked operation
(Section~\ref{sec:yoked}); in Yoked mode, primitive-set targets
are interpreted as time-varying $\{p_k(t)\}$ specifying a target
trajectory or desired coupled behaviour rather than a separable
scene hypothesis.
\subsubsection{Projector--surface compensation and multi-projector
blending (non-limiting)}

In some embodiments, the emission device is a conventional digital
projector and the observation device is a camera viewing the projection
surface. The kernel uses projected calibration patterns and camera
observations to estimate a geometric warp from projector coordinates to
perceived surface coordinates, and applies a corresponding inverse warp
to outgoing frames. In some embodiments, multiple projectors illuminate
overlapping regions of a common surface and the kernel estimates overlap
geometry, relative brightness, and colour differences; the kernel then
performs edge blending by driving each projector with a separately
warped and gain-compensated pattern so that the combined projection
appears seamless.

\subsubsection{Nulling: scene-invariant rendering (non-limiting)}

In some embodiments, \RT is operated with a target $Y^*$
chosen to be scene-invariant: the objective is to find the feedback
policy that makes the output independent of the scene, so that what
remains is pure device signature. Nulling is the anti-\LI: instead
of maximising scene information in the output, it minimises it. Nulling
is useful for isolating device identity from environmental
contamination, performing differential measurements (output with scene
minus output nulled = pure scene contribution), and calibrating the
trivial-reactor baseline.

% ----------------------------------------------------------------------
\subsection{Yoked Operation (Dynamic Coupling Modifier)}
\label{sec:yoked}
% ----------------------------------------------------------------------

\paragraph{What changes in Yoked mode.}
In static operation, the scene contributes to the \cb in a
separable way: in principle, one could factor the device's contribution
(emission pattern, detector response, reactor dynamics) from the
scene's contribution (reflectance, geometry, spectral properties). The
device acts on the scene; the scene modifies the device's output; but
the two are distinguishable subsystems.

In Yoked operation, this factorisation is expected to fail. The feedback loop is
tuned so that device and scene are driven toward a coupled dynamical
state whose behaviour is not readily decomposed into what the device
does and what the scene does. The \cb is a device-side record
generated under joint dynamics, not a sum of individual contributions.
Remove either partner and the joint coupled behaviour ceases to exist.

In preferred embodiments, the coupling is asymmetric at the level of
autonomous dynamics: scene-as-driver is the preferred modelling
convention when the scene has pre-existing autonomous dynamics that
continue independently of the device on the timescale of interest,
and the device is modelled as the \emph{responder} whose dynamics are
shaped by coupling to the scene.  This drive--response convention
determines the correct directionality of the Yoked diagnostics: the
operative
transfer entropy is $\TE_{S \to R}$ (scene driving device), the
operative conditional Lyapunov exponent is computed for the response
system conditioned on the driver trajectory, and the Yoked condition
is met when the response system has entrained to the driver while the
driver retains its autonomous character.  Symmetric mutual
synchronisation (in the sense of Pecora and Carroll's original
two-equivalent-oscillator formulation) is a special case and not the
general operating assumption.

This is not merely ``feedback.'' Every closed-loop system has feedback.
Yoked operation is the specific dynamical condition where (i)~the
coupled system is tuned toward a non-product invariant set in the
product state space (the intended dynamical target of the CLE/TE
co-condition, not an independently asserted property), (ii)~the
system operates near the edge of synchronisation where perturbation
sensitivity is high, and (iii)~the \cb is a device-side record
generated under joint dynamics, encoding how the scene
\emph{responds to being driven}---resonances, damping, nonlinear
coupling, transfer functions---properties not captured by static
open-loop optical observation and that depend on internal structure
rather than surface appearance.

The authentication warrant for Yoked operation rests on a
physical-access argument rather than a computational-hardness argument,
conditioned on a declared threat model. Under the declared threat
model --- which specifies the attacker's emulator class, fidelity
threshold against the genuine bundle, accessible meter family,
photon/sample budget, and access budget to the genuine device-scene
pair --- a committed \cb produced in Yoked mode is not generable from
the resources available to that attacker class without joint physical
access to the specific device and the specific scene at the time of
measurement, to the declared fidelity threshold. The coupling
structure recorded in the bundle --- CLE, TE balance, Arnold tongue
profile --- arises from the device-scene pair at that moment, and a
faithful reproduction of those quantities to the declared fidelity
threshold is not, under the declared threat model, achievable from a
model of the device or a description of the scene alone. The warrant
is thus operational and threat-model-conditioned: forgery becomes
infeasible at the declared fidelity threshold for attackers within the
declared budget; absolute infeasibility, infeasibility against
unbounded attackers, and infeasibility outside the declared meter
envelope or operating point are not claimed. The threat model,
fidelity threshold, accessible meter family, and access budget are
committed to the protocol digest, so that the warrant is auditable and
its scope is explicit. This warrant is independent of the Fisher
information conjecture (Appendix~\ref{sec:appendix-research}).

The name ``Yoked'' carries a deliberate double meaning: \emph{yoke} as
in the deflection yoke of a CRT---the magnetic assembly that steers the
beam and whose idiosyncrasies are part of the device's unclonable
fingerprint---and \emph{yoke} as in the Sanskrit \emph{yoga}, the
practice of joining, of coupling the self to something beyond itself.

\subsubsection{Operative diagnostics for Yoked operation (non-limiting)}

The following quantities are the operative diagnostic suite for Yoked
operation.  They are drawn from established nonlinear-dynamics and
synchronisation theory: conditional Lyapunov exponents as developed
by Pecora and Carroll (Phys.\ Rev.\ Lett.\ 64, 821--824, 1990),
transfer entropy as developed by Schreiber (``Measuring information
transfer,'' Phys.\ Rev.\ Lett.\ 85(2), 461--464, 2000), and Arnold
tongue analysis from classical forced-oscillator theory.  What is
novel is their application as engineering targets for a closed-loop
optical device that simultaneously performs authentication, sensing,
and rendering.  The Yoked condition is declared when the CLE and TE
conditions specified in the protocol digest are simultaneously met for
a declared window duration; this four-part co-condition is the
operative declaration standard.  The non-decomposability of the
resulting joint attractor is the intended dynamical consequence of
meeting these conditions, not an independently asserted or
independently measurable claim.

\paragraph{Balanced transfer entropy.}
$\TE_{S \to R} \approx \TE_{R \to S}$, both large, indicating
strong bidirectional information transfer between scene and reactor.
When transfer entropy is strongly unbalanced (one direction dominant),
the system is in a master--slave configuration.  The Yoked condition
is operationally declared when sustained high bidirectional TE is
detected within the qualified meter envelope for the declared window
duration; this is the verifiable claim.  Whether high bidirectional
TE reflects mutual co-production of the attractor rather than
near-perfect entrainment of the response to the driver is not
asserted as a consequence of the TE criterion alone; distinguishing
these cases requires supplementary perturbational or
partial-information-decomposition analysis, which is a research
direction.

\paragraph{Conditional Lyapunov exponent near zero ($\CLE \to 0^-$).}
The conditional Lyapunov exponent (CLE) is the largest sub-Lyapunov
exponent of the response system (reactor) computed along trajectories
conditioned on the driver (scene) trajectory, in the sense of
Pecora and Carroll (Phys.\ Rev.\ Lett.\ 64, 821--824, 1990).
Formally, let $\mathbf{x}_S(t)$ be the driver trajectory and let
$\delta \mathbf{x}_R(t)$ denote an infinitesimal perturbation to the
response system state.  The CLE is
\[
  \CLE \;=\; \lim_{T\to\infty} \frac{1}{T}
  \ln \frac{\|\delta\mathbf{x}_R(T)\|}{\|\delta\mathbf{x}_R(0)\|},
\]
evaluated along the coupled trajectory driven by $\mathbf{x}_S$.
The $\CLE$ measures whether the reactor synchronises with the scene
drive. When $\CLE < 0$, the reactor locks to the scene. The
operating point near $\CLE = 0^-$ is the edge of synchronisation,
where perturbation sensitivity is empirically high in coupled-oscillator
systems studied to date. At this point, the Fisher information about
the coupling strength increases sharply and may diverge, a standard
property of systems near criticality
(analogous to the divergence of magnetic susceptibility at a
ferromagnetic phase transition). Whether this gives a formal
authentication advantage is an open research direction
(Appendix~\ref{sec:appendix-research}).

In some embodiments, the \yd $\alpha \in [0,1]$ parameterises
the strength of bidirectional coupling, where $\alpha = 0$ corresponds
to static (open-loop or weakly coupled) operation and $\alpha = 1$
corresponds to the maximum commanded coupling depth; whether the
resulting operation is classified as Yoked is determined by the
CLE/TE co-condition recorded in the protocol digest. The \yd is
recorded in the protocol digest and may vary within a run.

\subsubsection{Yoked Limager: dynamic perception (non-limiting)}

In static \LI, the inferred variables are properties of what the
scene \emph{is}: geometry, reflectance, spectral signatures. In Yoked
\LI, the inferred variables include properties of how the scene
\emph{responds}: resonant frequencies, damping coefficients, nonlinear
response curves, transfer functions, mechanical impedance, and coupling
topology (which parts of the scene are dynamically connected to which
other parts).

In some embodiments, these dynamical properties are accessible through
bidirectional coupling because they describe the scene's response to
driving, not its static appearance. A cracked beam has different resonant frequencies
from an intact one. A live finger has different vascular pulsation
dynamics from a silicone replica. A genuine painting's pigment layers
have different micro-vibrational responses from a surface print. In
each case, the static appearance can be replicated; the dynamical response function encodes internal structure
rather than surface properties, and is not straightforwardly
reproducible from static appearance alone.

In some embodiments, Yoked \LI operates by sweeping coupling
parameters (feedback gain, modulation frequency, drive amplitude) and
recording the frequencies and coupling strengths at which entrainment
occurs, constructing a scene-specific Arnold tongue map that reveals
dynamical properties invisible to passive observation.

\paragraph{Arnold tongue: frequency identification (non-limiting).}
In some embodiments, the frequency axis of an Arnold tongue analysis
is the \emph{modulation envelope frequency} $f_{\mathrm{mod}}$ ---
the slow periodic variation imposed on emission (whether amplitude
modulation, frequency modulation, phase modulation, polarisation
modulation, spatial pattern repetition rate, or a reactor-internal
characteristic timescale) --- and the detuning $\Delta f = f_{\mathrm{mod}}
- f_{\mathrm{scene}}$ is the difference between this frequency and
the scene's natural dynamical frequency.  The modulation envelope
frequency is explicitly distinguished from (a)~the optical carrier
frequency ($\sim 600\,\mathrm{THz}$ for visible light), which is
irrelevant to Arnold tongue analysis at the dynamical timescales of
interest; (b)~the AOD or EOM RF carrier frequency ($\sim 10$--$100\,\mathrm{MHz}$),
which is a hardware operating parameter of the acoustic transducer
rather than a modulation signal applied to the scene; and (c)~the
scan or deflection repetition rate, which determines spatial coverage
geometry.  In some embodiments, the modulation envelope frequency does
not appear in the emission signal at all but resides in the reactor's
internal dynamics: the characteristic decay timescale of a phosphor
persistence layer (non-limiting example: P20 phosphor,
$\sim 50\,\mathrm{ms}$ to several seconds; P1 phosphor,
$20$--$80\,\mathrm{ms}$), the thermal relaxation rate of a scattering
medium, or the round-trip delay of a fibre cavity are all candidate
internal oscillators whose natural frequencies may entrain to or
drive entrainment with scene dynamics independently of the commanded
emission schedule.

\paragraph{Multi-dimensional tongue structure (non-limiting).}
In some embodiments, the full Arnold tongue is not a two-dimensional
map in $(\alpha, \Delta f)$ but a projection of a higher-dimensional
entrainment manifold, where each independently controllable
modulation degree of freedom (AM envelope, FM deviation, phase
modulation rate, polarisation rotation rate, spatial pattern
repetition rate, scan law repetition rate, and each reactor-internal
characteristic timescale) contributes a separate tongue structure in
a corresponding subspace.  In some embodiments, these tongue
structures are superimposed: a single training run can simultaneously
map respiratory entrainment ($\sim 0.1$--$0.5\,\mathrm{Hz}$, accessible
via spatial pattern repetition rate or slow AM envelope), cardiac
entrainment ($\sim 0.8$--$2\,\mathrm{Hz}$, accessible via fast-loop AM
interacting with photoplethysmographic return), and reactor-internal
dynamical structure (phosphor persistence band or thermal mode) within
a single committed \cb.  The resulting multi-dimensional tongue profile
constitutes the scene's dynamical fingerprint: it encodes not which
frequencies are present but which combinations of driving degrees of
freedom produce entrainment.

\subsubsection{Yoked Truth Beam: dynamic verification (non-limiting)}

In static \TB, security rests on the device's static
microstructure fingerprint: under declared adversary families, an
attacker would need to clone the physical device
to produce matching response statistics. In Yoked \TB, security
additionally rests on the \emph{dynamics of scene--reactor coupling}:
an attacker would need to reproduce the device's nonlinear response to
arbitrary
scene inputs in real time, not just its response to a known challenge
set.

In some embodiments, this is expected to be substantially harder because
the scene's own dynamics serve as
a continuous, unpredictable challenge that the attacker is not expected
to anticipate under declared resource bounds. A static (non-yoked) microstructure-unclonability mode has a fixed (if large) challenge--response
space that can in principle be exhaustively characterised. A Yoked microstructure-unclonability mode
has a response that depends on the scene's own dynamics, which are
unknown to the attacker in advance and change continuously. At the edge
of synchronisation ($\CLE \approx 0^-$), any discrepancy between the
genuine device's nonlinear response and an emulator's response is
amplified by the coupled dynamics rather than averaged out.

In some embodiments, Yoked \TB uses the coupling quality itself
as a continuous tamper-detection signal: any interruption of the
physical coupling (device removal, interposition of a screen, replay
injection) causes immediate collapse of the joint attractor and is
detectable as a discontinuity in transfer entropy balance, spectral
coherence, or order parameter, without waiting for a scheduled
verification challenge.

\subsubsection{Yoked Reality Transform: dynamic rendering
(non-limiting)}

In static \RT, the target $Y^*$ is a fixed image or
appearance. In Yoked \RT, the target is a \emph{dynamical
behaviour}: a limit cycle, a frequency lock, a coupled oscillation
pattern. The scene participates in producing the output, and the device
steers the joint attractor, in contrast with static projection.

Non-limiting applications include: projection mapping that responds to
and entrains with structural vibrations of the projection surface; light
installations that synchronise with crowd movement or biological
rhythms; interactive displays whose visual output is coupled to the
viewer's pupillary, postural, or gestural dynamics; and feedback-driven
physical computation in which the rendering target is a desired
computational state of the coupled system.

In some embodiments, the rendering target specifies a desired coupling
profile (for example target coherence at specific frequency bands)
rather than a target image, and the \RT controller adjusts
emissions to achieve and maintain the specified coupling profile.

\subsubsection{Training mechanisms for Yoked operation (non-limiting)}

In Yoked operation, the scene itself provides the training signal. This
is fundamentally different from the informational regimes, where
training typically uses an external reference (an enrolled template,
ground-truth labels, or a target image). The coupling quality is both
the objective and the signal.

In some embodiments, training uses one or more of the following
mechanisms:

\paragraph{Physical self-tuning (analogue).}
Feedback gain is increased gradually until spontaneous oscillation or
phase-locking emerges. The reactor's own nonlinearity provides the error
signal through amplitude-dependent frequency pulling---a phase-locked
loop variant in which the coupled system is the oscillator, the scene
provides the reference dynamics, and the reactor's slow dynamics serve
as the loop filter. The Arnold tongue structure maps where entrainment
is possible as a function of frequency mismatch and coupling strength;
training finds the tongue and stabilises within it.

\paragraph{Transfer-entropy gradient (hybrid).}
Transfer entropy $\TE_{S \to R}$ and $\TE_{R \to S}$ are estimated from
the \cb. In some embodiments, a training-phase proxy objective
\[
  J_{\mathrm{train}}(\theta, \pi)
  = \lambda_{\mathrm{mag}}\,\min\!\bigl(\TE_{S \to R},\, \TE_{R \to S}\bigr)
  - \lambda_{\TE}\,\bigl|\TE_{S \to R} - \TE_{R \to S}\bigr|
\]
is maximised via natural gradient ascent
$\theta \leftarrow \theta + \eta\, \mathcal{F}(\theta)^{-1}
\nabla_\theta J_{\mathrm{train}}$, where $\mathcal{F}(\theta)$ is the
Fisher information matrix
(pseudoinverse or damped inverse where $\mathcal{F}$ is singular or
ill-conditioned near bifurcation boundaries).
This training objective is a component of the full Yoked objective
$J_{\mathrm{Yoked}}$ (Section~\ref{sec:yoked-control-objective});
the CLE target term enters after the TE-gradient phase has established
a coupling foothold.
In some coupled dynamical systems, Fisher
Transfer Entropy increases sharply near the entrainment phase
transition, so the natural gradient may automatically
accelerate toward the coupling boundary.

\paragraph{Spectral coherence maximisation (practical).}
Cross-spectral density between scene return and reactor output is
computed at each frequency band; feedback parameters are tuned to
maximise coherence at target frequencies. This provides narrowband
entrainment---locking onto specific scene dynamics while ignoring
others.

\paragraph{Full-$\theta$ Yoked training (non-limiting).}
\label{par:full-theta-yoked}
In some embodiments, the gradient update $\theta \leftarrow \theta
+ \eta\,\mathcal{F}(\theta)^{-1}\nabla_\theta J$ is applied over the
full parameter vector $\theta = (\theta_{\mathrm{hw}},
\theta_{\mathrm{sw}})$ in a coordinated sweep that includes, without
limitation, all of the following independently controllable degrees of
freedom as training knobs:
\begin{itemize}
  \item \textbf{Emission modulation degrees of freedom:} amplitude
    modulation envelope frequency and depth; frequency modulation
    deviation and rate; phase modulation rate; polarisation rotation
    rate and axis; spatial pattern repetition rate; and any compound
    modulation waveform combining two or more of the above.
  \item \textbf{Scan law parameters:} sweep repetition rate, dwell
    time per region, and spatial extent of the scan trajectory
    (which determine the range of spatial pattern repetition rates
    presented to the scene).
  \item \textbf{Feedback gain and coupling depth:} loop gain
    magnitude and phase, \yd $\alpha$, and feedback filter
    bandwidth.
  \item \textbf{Reactor-internal timescale parameters:} in some
    embodiments, the reactor's characteristic timescale is
    selectable --- for example through choice of phosphor formulation
    (decay time), drive intensity (affecting phosphor saturation),
    cavity round-trip length, or thermal load --- and this timescale
    enters as a trainable parameter in $\theta_{\mathrm{hw}}$,
    because it determines which scene dynamics the reactor can
    entrain to.
  \item \textbf{Drive amplitude:} emission power, scan speed, and
    detector integration time (which jointly determine the SNR
    available to the convergence estimators).
\end{itemize}
In some embodiments, an Arnold tongue mapping phase precedes gradient
ascent: the control protocol sweeps the modulation-envelope frequency
and coupling gain jointly, estimates $\widehat{\CLE}$ at each grid
point, and uses the resulting empirical tongue map to initialise
$\theta$ at a point inside the tongue boundary before commencing
gradient-based refinement.  This separates the coarse exploration
phase (find the tongue) from the fine exploitation phase (park at
the boundary), and makes the training robust to scenes where the
tongue is narrow or located at an unexpected frequency.  Both phases
are committed in the \cb and logged in the protocol digest.

In some embodiments, the Yoked training objective acts over the
joint product of all modulation degrees of freedom simultaneously,
seeking the combination $(\theta_{\mathrm{hw}}, \theta_{\mathrm{sw}})$
at which the multi-dimensional tongue intersection --- the region of
simultaneous entrainment across AM, FM, phase, spatial, and
reactor-internal channels --- is widest and most stable.  The
resulting parameter set constitutes the scene's multi-dimensional
dynamical fingerprint as captured by the committed \cb, and its
hardness derives from the attacker needing to reproduce the correct
entrainment signature across all degrees of freedom simultaneously.
The physical mechanism by which $\nabla_\theta J$ is computed in
practice---without numerical finite differences or a differentiable
surrogate---is disclosed in Section~\ref{sec:full-theta-sweep}
(scene-phased full-parameter sweeping).

\subsubsection{Yoked state estimation: classifier-mediated coupling
detection (worked embodiment, non-limiting)}
\label{sec:yoked-classifier}

The parameter training mechanisms described above (Arnold tongue sweep,
TE-gradient ascent, full-$\theta$ sweep) address one training problem:
how to adjust $\theta$ so that the device reaches the CLE/TE operative
condition for a given scene.  A second training problem arises at
runtime: how does the device determine whether it is currently in the
Yoked condition, particularly against scenes for which CLE and TE
cannot be computed from both sides simultaneously?  For analysable
physical reactor scenes in a laboratory setting, CLE and TE can be
estimated directly from both the device and scene time series.  For
scenes lacking an accessible internal state (non-limiting example:
a human observer in a Yoked \RT interaction), the scene's
internal trajectory is not available to the device, and explicit
CLE/TE computation from both sides is not possible at runtime.

In some embodiments, this problem is addressed by a
\emph{coupling state classifier} $\hat{f}_\phi$ trained to estimate
the Yoked condition from device-internal bundle data alone.  The
classifier is trained offline using verified CLE/TE ground truth
computed from known analysable reactor pairs, then qualified before
deployment, and deployed at runtime against declared scene families
for which the classifier has been validated.  Whether such classifiers
generalise reliably to scene classes outside their training
distribution is an empirical question addressed in the Open Research
Directions (Section~\ref{sec:appendix-research}).

In some embodiments, the classifier is a \emph{population-level coupling
estimator}: trained on a sufficiently large and diverse reactor
population, it is intended to generalise across a held-out range of
reactor types within the declared qualification set without per-device
adaptation, and it qualifies against held-out reactors drawn from the
same broad population.  In some embodiments, the classifier is a
\emph{device-specific estimator}: a population-trained base
model is finetuned on known reactors interacted with on the specific
device instance to be deployed, capturing the unclonable physical
characteristics of that hardware instance, so that the qualified
coupling declaration is constitutively tied to that device and cannot
be replayed on a different hardware instance without a new finetuning
and qualification run.  These two embodiments represent the ends of a
spectrum; intermediate embodiments (partial finetuning, ensemble
methods, domain-adaptive classifiers) are also within scope.  The
four-stage architecture below describes the device-specific case as
the worked embodiment; the universal case is recovered by omitting
Stage~2.

In all embodiments, the classifier operates on device-internal bundle
data alone.  The TE estimator family and conditioning policy used to
generate training supervision labels are identical to those committed
in the protocol digest for runtime Yoked declaration: if the runtime
declaration uses conditional transfer entropy $\TE_{S \to R \mid
\mathbf{Z}}$ with covariate set $\mathbf{Z}$ (as in the preferred
embodiment of Section~\ref{sec:yoked-convergence}), the supervision
labels are computed using the same conditional estimator and covariate
set.  This consistency requirement is committed to the protocol
digest.  The four stages are as follows.


\paragraph{Scope (non-limiting).}
The control-theoretic objective $J_{\mathrm{Yoked}}$ disclosed above is
directed to adaptive probing, coupled-state optimisation, and navigation
of $\Theta$ toward entrainment regimes. It is not an assertion that the
disclosed physical training machinery constitutes a general-purpose
substrate for deep-learning-scale parameter optimisation. Capability
claims arising from this subsection are bounded by the capacity
accounting rule of Section~\ref{sec:capacity-accounting-rule}.

\paragraph{Stage 1: Population training (offline).}
In some embodiments, a training device is coupled in sequence to
$K \geq 20$ known analysable reactors spanning at least three distinct
reactor types (non-limiting examples: CRT phosphor, scattering
plate, resonant mechanical structure).  For each coupling run,
ground-truth CLE and TE are computed from both the device and scene
time series over windows of declared length $T$ at a declared
effective sampling rate.  The supervision window length $T$ is
committed to the protocol digest and matched to the rolling window
length used at runtime.  $T$ should span at least 50 characteristic
timescales of the dominant scene dynamics to ensure reliable
finite-time CLE estimation; for physical reactors with fast dynamics
(e.g.\ scattering plates at hundreds of Hz), $T = 120\,\mathrm{s}$
is typically sufficient; for slow-timescale scenes, $T$ is scaled
accordingly and committed per training reactor.

The ground-truth labels are the triple $(\TE_{\text{achieved}},
\CLE_{\text{achieved}}, \TE_{\text{balance}})$ computed from both
sides of each coupling run under the declared estimator family.  The
training set consists of $K$ coupling-regime interactions; within
each interaction, non-overlapping windows of length $T$ are extracted
from the committed \cba atoms.  The effective number of independent
coupling-regime labels is $K$ (one per reactor interaction);
within-interaction windows are autocorrelated and do not constitute
independent supervision signals.  Training is performed with
run-level train/validation splitting to prevent within-run leakage.
In some embodiments, semi-synthetic augmentation (noise-perturbed
replays of each reactor interaction) is used to expand the effective
diversity of the training set beyond $K$ interactions.

The classifier $\hat{f}_\phi$ is trained on device-internal bundle
data alone---it never receives the scene's internal time series
as input, only the committed \cba atoms and derived meter summaries
from the device side of each interaction.  The training loss is a
weighted mean-squared error over the ground-truth triple:
\[
  \mathcal{L}(\phi)
  = \mathbb{E}\!\left[
      w_1\,(\hat{\TE}_{S\to R} - \TE_{\text{achieved}})^2
    + w_2\,(\hat{\CLE} - \CLE_{\text{achieved}})^2
    + w_3\,(\hat{\Delta}_{\TE} - \TE_{\text{balance}})^2
  \right],
\]
where $w_1, w_2, w_3$ are declared loss weights recorded in the
protocol digest.  In some embodiments, the classifier architecture is
a recurrent neural network (non-limiting: LSTM or GRU) operating on
the rolling bundle window, chosen for its capacity to represent
temporal dependencies in the coupling dynamics.

\paragraph{Stage 2: Per-device finetuning (device-specific embodiment
only).}
In some embodiments, the population-trained classifier is finetuned on
$M_{\text{ft}} \in [3, 5]$ known analysable reactors interacted with on
the specific device instance to be deployed.  Lower layers of the
classifier (encoding general temporal structure of the coupling
dynamics) are frozen; upper layers (encoding the mapping from bundle
statistics to metric values) are finetuned.  This captures the
unclonable physical characteristics of the specific hardware instance
that distinguish it from the population average, such that the
qualified coupling declaration is constitutively tied to that device:
a classifier qualified on device $A$ and replayed on device $B$
produces invalid coupling declarations because device $B$'s bundle
statistics lie outside the qualified envelope.  A new finetuning and
qualification run is required for each device instance.  Estimated
finetuning time is 20--45 minutes per device instance.  In the
universal embodiment, this stage is omitted.

\paragraph{Stage 3: Stability envelope qualification.}
Before deployment, the classifier is qualified by running it against
$M_{\text{qual}}$ held-out known analysable reactors not seen during
population training or finetuning.  The classifier output is compared
against ground-truth CLE and TE computed from both sides under the
same estimator family used in training.  The qualification pass
criteria are:
\begin{enumerate}
  \item Mean absolute error of the classifier's $\TE$ estimate is
    below $15\,\%$ of the declared TE target range, where the TE
    target range is committed to the protocol digest \emph{before}
    qualification runs are evaluated, based on the declared use case
    and not on observed classifier performance.
  \item The classifier correctly identifies the Yoked / non-Yoked
    classification (as defined by the four-part co-condition of
    Section~\ref{sec:yoked-convergence}) for all $M_{\text{qual}}$
    held-out reactors.
  \item No systematic bias in the CLE or TE estimates (signed mean
    error within $\pm 5\,\%$ of the declared target range).
\end{enumerate}
In preferred embodiments, $M_{\text{qual}}$ is set so that the
probability of false passage (passing qualification despite a true
per-reactor misclassification rate exceeding a declared threshold
$p_0$) is below a declared confidence level $\alpha$, requiring
$M_{\text{qual}} \geq \log(\alpha) / \log(1 - p_0)$; the achieved
misclassification confidence bound is committed to the protocol
digest alongside the raw pass/fail record.  In some embodiments, a
minimum of $M_{\text{qual}} \geq 5$ is used for deployment under
this classifier-mediated embodiment; preferred
embodiments use larger qualification sets determined by the above
formula for the declared $p_0$ and $\alpha$.  A device that fails
qualification is not used for Yoked operation until retraining or
refinetuning.  In some embodiments, the qualification record forms
part of the Proof-of-Discrepancy submission for fleet-level
calibration (Section~\ref{sec:proof-of-discrepancy}).

\paragraph{Stage 4: Runtime deployment.}
In some embodiments, at runtime the classifier operates on a rolling
window of committed \cba atoms of the declared length $T$ and
produces estimates of $(\hat{\TE}_{S\to R}, \hat{\CLE},
\hat{\Delta}_{\TE})$ without requiring access to the scene's internal
state.  During the pre-Yoked phase (Phase~1, in which the parameter
sweep seeks the Arnold tongue and establishes initial coupling), the
classifier operates in \emph{monitoring-only mode}: its outputs are
logged but do not trigger phase transitions or Yoked declarations.
Phase~1 control decisions are based on direct two-sided TE estimation
where available, or on spectral coherence metrics committed in the
\cb.  The classifier becomes the operative control signal only after
Phase~1 has established a coupling foothold, as declared by the
Phase~1 exit criterion committed in the protocol digest.

After Phase~1 exit, the four-part co-condition
(Section~\ref{sec:yoked-convergence}) is evaluated against classifier
estimates, and the protocol digest records that the Yoked declaration
is classifier-mediated under the declared qualification set, weights
hash, supervision window length, and estimator consistency
commitment.

In some embodiments where the scene is agentive (non-limiting
example: a human observer in a Yoked \RT interaction),
the control objective shifts from reaching a fixed target coupling
value to minimising the rate of change of the classifier's coupling
estimate, seeking a stationary mutual coupling state within the
qualified envelope rather than a declared scalar target.  This
reflects that an agentive scene has its own autonomy and the coupling
equilibrium is co-determined, not unilaterally imposed.

Failure modes detected at runtime include: out-of-distribution scenes
(detected by classifier output variance and embedding-space distance
from training distribution, with thresholds committed to the protocol
digest), coupling instability (detected by persistent drift in the
classifier's CLE estimate), and attractor collapse (detected by
abrupt discontinuity in TE balance), all of which are logged in the
protocol digest with the classifier confidence at the time of the
event.

\paragraph{Role of this embodiment.}
The detection architecture used here --- a sequence model trained on
one-sided time series to infer synchronisation state --- has precedent
in the prior art.  The inventive contribution of this embodiment is
not the classifier itself but the \emph{qualification-gated evidentiary
pipeline}: a model supervised by physically verified two-sided coupling
metrics, qualified against held-out physical reactors under
precommitted thresholds, and committed by weights hash and
qualification record to the protocol digest, such that a Yoked
declaration in non-analysable scenes is auditable and reproducible by
a third party.  In the device-specific embodiment, the declaration is
additionally constitutively tied to a specific physical hardware
instance via per-device finetuning.  This embodiment supports the
primary claims of the \RK architecture and demonstrates
operability of the Yoked condition against non-analysable scenes; it
is not itself advanced as a primary claim.

\paragraph{Relationship between parameter training and state
estimation training.}
In some embodiments, the two training problems are architecturally
distinct but operationally complementary.  Parameter training
(Sections above) adjusts $\theta$ to steer the device toward the
Yoked condition for a given scene.  State estimation training (this
section) produces a classifier that detects whether that condition
has been reached and maintains it, without requiring direct access to
the scene's internal state.  Together they constitute a complete Yoked
training architecture operable across scene types ranging from
fully analysable physical reactors (where explicit CLE/TE ground truth
is available from both sides) to non-analysable agentive scenes (where
only device-internal bundle data is available at runtime).  The
explicit CLE/TE protocol of the parameter training phase is the
supervisory signal that makes the classifier trustworthy; the
classifier is the mechanism that makes the Yoked condition detectable
and controllable in deployment.

\subsubsection{Convergence assessment (non-limiting)}
\label{sec:yoked-convergence}

In some embodiments, convergence is assessed using a four-part
co-condition:
\begin{enumerate}
  \item \textbf{CLE window condition:} $\widehat{\CLE}$ remains in
    $[-\varepsilon_{\mathrm{CLE}}, 0)$ for $N$ consecutive
    non-overlapping windows, where $\varepsilon_{\mathrm{CLE}}$ and
    $N$ are protocol parameters declared in advance and recorded in
    the protocol digest.
  \item \textbf{TE balance condition:}
    $|\TE_{S \to R} - \TE_{R \to S}| < \tau_{\TE}$, where
    $\tau_{\TE}$ is a declared threshold, evaluated over the same
    windows.
  \item \textbf{TE minimum magnitude condition:}
    $\min(\TE_{S \to R}, \TE_{R \to S}) > \tau_{\mathrm{mag}}$, where
    $\tau_{\mathrm{mag}} > 0$ is a declared minimum magnitude
    threshold.  This condition ensures that TE balance cannot be
    satisfied trivially when both transfer entropies are near zero; a
    system with negligible coupling in both directions can satisfy
    condition~2 while failing condition~3.  $\tau_{\mathrm{mag}}$ is
    declared in the protocol digest in the same units as the TE
    estimator and is set to a value reflecting the minimum coupling
    magnitude regarded as operationally significant for the declared
    use case.
  \item \textbf{Meter SNR gate:} estimation confidence metrics
    (sufficient sample size and SNR as declared by the meter
    envelope) are satisfied for both the CLE estimator and the TE
    estimator, so that the convergence declaration is not made on
    under-resolved data.
\end{enumerate}
In some embodiments, all four conditions must be satisfied simultaneously for Yoked
operation to be declared.  The thresholds
$(\varepsilon_{\mathrm{CLE}}, \tau_{\TE}, \tau_{\mathrm{mag}}, N)$ and the estimator
families and window parameters are recorded in the protocol digest
so that a verifier can reproduce the convergence assessment from
committed \cba data.  In preferred embodiments, $\varepsilon_{\mathrm{CLE}}$
is specified in the same time-base unit as the declared CLE estimator
(per-step or per-second) and $\tau_{\TE}$ is specified in the same
information unit as the declared TE estimator (bits or nats per
declared window); the estimator's log-base and windowing convention
are part of the protocol digest commitment.

In some embodiments, the transfer entropy estimator is selected from
one or more of the following non-limiting families: histogram-based
binning estimators, Kraskov--St\"ogbauer--Grassberger (KSG)
$k$-nearest-neighbour estimators, kernel density estimators, or
learned neural estimators. The estimator family, window length,
window overlap, embedding dimension, lag selection, and any
preprocessing choices are recorded in the protocol digest so that
later verifiers can reproduce the convergence assessment from
committed \cba data.

\paragraph{Common-cause confounding and conditional transfer entropy
(preferred embodiment).}
Transfer entropy estimated from unconditional marginals can produce
apparent high coupling even without any direct physical feedback loop,
if the scene has strong autocorrelation or if scene and reactor share
a common external driver (a phenomenon documented in physical measurement
systems by Vakorin et al.\ 2009 and Papana et al.\ 2020).
In preferred embodiments, the Yoked declaration therefore uses
\emph{conditional} transfer entropy: $\TE_{S \to R \mid \mathbf{Z}}$
and $\TE_{R \to S \mid \mathbf{Z}}$, where $\mathbf{Z}$ comprises
all measured scene covariates committed to the \cb (for example
scene-side sensor readings, ambient illumination, temperature, and any
shared external stimulus measurable within the declared meter envelope).
The conditional estimator family is recorded in the protocol digest
with the same provenance as the unconditional estimator.  In some
embodiments where scene covariates are partially observable, a
common-cause meter is additionally declared: it tests whether the
TE balance condition can be reproduced by a surrogate scene trace
generated from the measured covariates alone, without reactor output
access; if the surrogate reproduces the balance condition, the Yoked
declaration is withheld and the event is logged as \emph{apparent
coupling} pending further investigation.  In embodiments where neither
conditional TE nor a common-cause meter is practical, the Yoked
declaration is qualified as \emph{apparent coupling under declared
measurement conditions} in the protocol digest.

\paragraph{Null-model calibration and operating characteristics for
Yoked declaration (preferred embodiment).}
Pre-declaration of thresholds is necessary but not sufficient to
make the Yoked declaration principled: the false-positive rate
(incorrectly declaring Yoked when dynamics are uncoupled), the
false-negative rate (failing to declare when coupling is genuine but
underresolved), and the minimum window length required for adequate
statistical power must be characterised against declared null
models.  In preferred embodiments, the calibration procedure is as
follows.  A null ensemble is constructed from phase-shuffled or
block-permuted versions of the scene and reactor time series,
preserving autocorrelation structure but destroying directional
coupling; in some embodiments, an auto-regressive null model is
fitted to each marginal series and the null ensemble is generated
by matched-simulation.  The four-part co-condition is applied to
the null ensemble; the empirical distribution of the co-condition
outcome defines the null operating characteristics.  In preferred
embodiments, $\tau_{\TE}$, $\tau_{\mathrm{mag}}$, $\varepsilon_{\mathrm{CLE}}$,
and $N$ are jointly set so that the false-positive rate under the
null ensemble is below a declared target $\delta_{\mathrm{FP}}$
(for example $\delta_{\mathrm{FP}} < 0.01$), and the minimum window
length is set so that the false-negative rate against a declared
alternative coupling model is below $\delta_{\mathrm{FN}}$.  The null
ensemble, the achieved $(\delta_{\mathrm{FP}}, \delta_{\mathrm{FN}})$
under the anchor embodiment, and the minimum window length are
committed to the protocol digest alongside the threshold values.
In embodiments where empirical calibration is not performed before
deployment, the Yoked declaration is qualified as
\emph{uncalibrated} in the protocol digest and carries reduced
evidentiary weight.

\paragraph{Statistical uncertainty and bootstrap confidence intervals
(non-limiting).}
In some embodiments, block-bootstrap confidence intervals are computed
for both the TE balance statistic and the CLE estimate, using
non-overlapping \cb windows of declared length.  The interval
halfwidth at a declared coverage level (for example 95\,\%) is
committed to the protocol digest alongside the point estimate.  The
SNR gate (condition~4 above) is necessary but not sufficient for
adequate statistical power; in preferred embodiments, the declared
thresholds $(\varepsilon_{\mathrm{CLE}}, \tau_{\TE})$ are set at a
declared multiple $\kappa \geq 2$ of the empirical bootstrap halfwidth,
ensuring that the Yoked declaration is not made unless the point
estimates are statistically separated from zero under the declared
estimator.

The minimum window length $T_{\min}$ is a committed protocol parameter.
Near $\widehat{\CLE} = 0^-$, finite-window CLE estimation is least
reliable: the $O(1/\sqrt{T})$ standard error of the estimator is of
the same order as the quantity being estimated, and the probability of
a false sign-flip (estimated CLE crossing zero when true CLE is just
below zero) is highest in precisely the operating region the Yoked
condition targets.  In preferred embodiments, $T_{\min}$ is set so that
the bootstrap halfwidth of the CLE estimate is below
$\varepsilon_{\mathrm{CLE}} / \kappa$; this is the condition under
which the multiplier $\kappa$ in the threshold-setting rule above
provides the declared false-positive coverage.  The value of
$T_{\min}$ is scene-class-dependent and must be characterised against
the null ensemble for each declared deployment class; it is committed
to the protocol digest per scene class alongside the achieved
halfwidth.

\paragraph{Lexicographic priority in the Yoked objective
(non-limiting).}
The scalar Yoked objective $J_{\mathrm{Yoked}}$ combines CLE proximity,
TE magnitude, and TE balance (Section~\ref{sec:yoked-control-objective}).
In some embodiments, competing objectives are resolved by a declared
lexicographic priority committed to the protocol digest: CLE proximity
is primary (the system is not declared Yoked unless
$\widehat{\CLE} \in [-\varepsilon_{\mathrm{CLE}}, 0)$); TE balance is
secondary; TE magnitude is tertiary.  This ordering prevents
configurations with high TE but diverging CLE from being declared
Yoked.  The lexicographic rule is non-limiting and may be replaced by
an alternative priority rule declared in the protocol digest.

\subsubsection{Limits of Yoked operation (non-limiting)}

The capabilities of Yoked operation are bounded by three physical
constraints:

\paragraph{Bandwidth.} The feedback loop has a characteristic timescale
set by the round-trip time. Scene dynamics faster than the loop
bandwidth are averaged or aliased; dynamics slower than the loop
bandwidth can be tracked and coupled to. For embodiments with a
two-rate stack such as the CRT/phosphor extension
(Section~\ref{sec:crt-phosphor-extension}), the fast optical or
modulation loop can couple to dynamics up to hundreds of MHz, while
the slow camera or sensor loop can couple to dynamics from
approximately 0.1~Hz upward; the bench-top 1D anchor
(Section~\ref{sec:anchor-1d}) is a single-rate variant whose loop
bandwidth is determined by the unified optical-and-detector return
path.

\paragraph{Scene cooperativity.} In preferred embodiments, the scene
responds dynamically to the device's output. A rigid, specularly
reflecting surface reflects the device's output without dynamical
engagement. Richer Yoked operation is facilitated by a scene with its
own dynamics---resonances, damping, nonlinearity---that interact with
the device's output. Living tissue, mechanical structures, reactive
media, and other devices are non-limiting examples of cooperative
scenes.

In some embodiments, a low-cooperativity scene (non-limiting example:
a rigid wall or planar diffuser) exhibits few and narrow entrainment
tongues, typically at structural resonance modes in the
$10$--$500\,\mathrm{Hz}$ range; training is constrained by SNR and
photon budget, and the multi-dimensional tongue map is comparatively
sparse.  In some embodiments, a high-cooperativity scene (non-limiting
example: living human tissue) is expected, based on known
physiological coupling characteristics, to exhibit tongue structures
across multiple frequency bands simultaneously, including respiratory
coupling ($\sim 0.1$--$0.5\,\mathrm{Hz}$, typically accessible via
spatial pattern repetition rate or slow AM envelope), cardiac coupling
($\sim 0.8$--$2\,\mathrm{Hz}$, accessible via fast-loop AM interacting
with photoplethysmographic return signal), skeletal and postural
dynamics, and higher harmonics.  Whether these tongue structures are
broad and overlapping or narrow and isolated depends on the strength
and nonlinearity of the relevant physiological oscillators and has
not yet been directly characterised in a projector-camera feedback
loop; this is a predicted embodiment pending experimental
confirmation (see Appendix~\ref{sec:appendix-research}).  Tongue
widths and positions are expected to shift with physiological state,
so that mapping benefits from being performed within a single session.
The full multi-dimensional tongue map for a human body, committed in
a \cb and logged in a protocol digest, would constitute a dynamical
identity record encoding the scene's response to a diversity of
driving conditions rather than its appearance under any fixed
illumination.

\paragraph{Signal-to-noise ratio.} In some embodiments, coupling quality
measurement benefits from sufficient photons per measurement window to
estimate transfer entropy or spectral coherence reliably. Meters should track estimation
confidence and gate Yoked claims on sufficient SNR.

\paragraph{Non-limiting training time estimates.}
In some embodiments, the time required to complete Yoked training from
cold start to convergence declaration depends on scene cooperativity,
estimator window length, and convergence threshold parameters.  For the
bench-top 1D RK anchor embodiment (Section~\ref{sec:anchor}), non-limiting
estimates for the major phases are: Arnold tongue mapping (coarse
grid sweep): $5$--$20\,\mathrm{min}$ wall time; initial lock
acquisition (finding the tongue and stabilising within it):
minutes of wall time for a low-cooperativity scene, seconds to minutes
for a high-cooperativity scene once inside the tongue; TE gradient
ascent to the coupling boundary: $20$--$60\,\mathrm{min}$ wall time
(bottleneck is TE estimator requiring $\sim$$100$--$300$ samples per
window at the relevant timescales); convergence declaration ($N=5$--$15$
windows at $\sim$$1\,\mathrm{min}$ per window): $5$--$15\,\mathrm{min}$.
Total first-run time is typically $1$--$4\,\mathrm{hours}$ for a new
scene; re-training to a previously characterised scene is typically
$30$--$60\,\mathrm{min}$.  These estimates are non-limiting and will
vary with hardware, scene properties, and declared convergence
thresholds; the relevant parameters are declared in protocol digests
so that convergence claims are reproducible.

% ----------------------------------------------------------------------
\subsection{Continuous Interpolation}
\label{sec:interpolation}
% ----------------------------------------------------------------------

In some embodiments, the operating regime is set by continuously varying
weights on a composite objective rather than selecting a discrete mode.
A composite objective combines \TB, \LI, and \RT terms with weights
$(\lambda_{\mathrm{TB}}, \lambda_{\mathrm{L}},
\lambda_{\mathrm{RT}})$ on the two-simplex $\Delta_2$, and a \yd
$\alpha \in [0,1]$ that controls the strength of bidirectional dynamical
coupling. The full operating point is
$(\lambda_{\mathrm{TB}}, \lambda_{\mathrm{L}},
\lambda_{\mathrm{RT}}, \alpha)$, allowing smooth transitions among
verification, sensing, and rendering, and between static and Yoked
operation, using the same hardware. In preferred embodiments, the
weights, \yd, and any transition schedule are recorded in the
protocol digest so that meters and audits condition on the declared
regime mix and coupling mode.

The regime space has the geometry of a triangular prism: the TB--L--RT
triangle at each value of $\alpha$, extruded along the \yda axis.
The base face ($\alpha = 0$) is standard static operation. The top face
($\alpha = 1$) is the maximum commanded coupling depth, at which any
regime is fully augmented with dynamic coupling; runs at the top face
are classified as Yoked operation only when the CLE/TE co-condition is
satisfied as recorded in the protocol digest. Key practical protocols occupy
identified regions:

\begin{description}[style=nextline, leftmargin=2em]
  \item[Static sensing] ($\alpha = 0$, high $\lambda_{\mathrm{L}}$): standard
    computational imaging and active perception.
  \item[Yoked sensing] ($\alpha \to 1$, high $\lambda_{\mathrm{L}}$):
    resonance-based scene characterisation, liveness detection,
    dynamic material identification.
  \item[Static verification] ($\alpha = 0$, high $\lambda_{\mathrm{TB}}$):
    standard microstructure-based authentication against enrolled fingerprints.
  \item[Yoked verification] ($\alpha \to 1$, high $\lambda_{\mathrm{TB}}$):
    continuous tamper detection, coupling-based authentication,
    scene-coupled security.
  \item[Static rendering] ($\alpha = 0$, high $\lambda_{\mathrm{RT}}$):
    standard projection mapping and style transfer.
  \item[Yoked rendering] ($\alpha \to 1$, high $\lambda_{\mathrm{RT}}$):
    entrainment-driven display, responsive environments, co-creative
    projection.
  \item[Verified Yoked sensing] (high $\alpha$, balanced
    $\lambda_{\mathrm{TB}}$ and $\lambda_{\mathrm{L}}$): dynamical scene sensing with
    simultaneous authentication of the coupling.
\end{description}

% ----------------------------------------------------------------------
\subsection{Information Operations Taxonomy}
\label{sec:info-operations}
% ----------------------------------------------------------------------

The three operating regimes and the Yoked modifier correspond to
composite objectives built from more primitive information operations.
This section enumerates non-limiting primitive operations, each of which
defines a distinct trainable axis with its own loss function. In some
embodiments, a protocol specification assigns a weight vector across
these axes, and the optimisation targets the weighted multi-objective
function on the Fisher-Rao geometry of $\Theta$.

\subsubsection{Regime-level operations (summary)}

The three regime-level operations and one coupling-mode operation
are described in detail in the preceding subsections:

\begin{description}[style=nextline, leftmargin=2em]
  \item[Read (\LI).] Maximise $I(S; C_{0:T})$: extract information
    about the scene.
  \item[Write (\RT).] Minimise $d(Y, Y^*)$: impose a
    target appearance on the scene return.
  \item[Authenticate (\TB).] Maximise separation between genuine
    and forged \cbs under a bounded attacker model.
  \item[Co-create (Yoked coupling mode).] Maximise the minimum of bidirectional
    transfer entropy $\min(\TE_{S \to R},\,
    \TE_{R \to S})$: establish a joint attractor that neither
    scene nor reactor could produce alone.  Yoked is a coupling-mode
    modifier applicable to any of the three regimes, not a fourth
    regime (Section~\ref{sec:regimes}).
\end{description}

\subsubsection{Primitive information operations (non-limiting)}

In some embodiments, the following primitive operations are individually
trainable and composable. Each primitive defines a loss function over
the \cb and the device's configuration parameters, and
may be combined with any regime-level objective.

\paragraph{Erase (non-limiting).}
In some embodiments, the system actively destroys targeted information
in the \cb. The objective is $\min I(C_{0:T}; S_{\mathrm{target}})$ subject to a declared utility constraint $I(C_{0:T}; S_{\mathrm{utility}}) \geq I_0$ over declared non-target scene features $S_{\mathrm{utility}}$, so as to drive the output toward statistical independence from specific target features while preserving information about non-target features sufficient for the declared utility model. The trade-off form parallels the information-bottleneck objective recited elsewhere in this disclosure and is not asserted to identify a unique optimum. The feedback loop
drives the scene--reactor system into a regime where the target features
are in the null space of the observation map. This is not passive
omission (``do not record'') but active destruction: the coupled
dynamics are driven to a region where the target information is
thermodynamically dissipated. In some embodiments, the erasure process
is designed so that reconstruction of the target information is
physically costly in proportion to the thermodynamic dissipation
incurred in driving the coupled dynamics to the null-space region;
that dissipation provides a physical lower bound on the resource cost
of any hypothetical reconstruction attempt under the declared attacker
model and information assumptions.  No claim of computational
infeasibility in the complexity-theoretic sense is asserted.

Non-limiting applications include physical redaction (erasing faces,
identifiers, or other personally identifiable information from the
\cb at the physical layer, before digital processing,
so that the information is rendered statistically non-inferable under the declared reconstruction model class and metered operating envelope), privacy-preserving sensing
(reading structural information while erasing identity information,
with the information bottleneck enforced by physics), and secure
disposal (driving reactor memory media to saturation to erase
accumulated scene history, with thermodynamic verification that
erasure occurred).

\paragraph{Encrypt (non-limiting).}
In some embodiments, the system makes the \cb
conditionally accessible. The objective is
$\max I(C_{0:T}; S \mid \theta)$ subject to
$\min I(C_{0:T}; S \mid \neg\theta)$: the output is informative about
the scene if the device parameters $\theta$ are known; uninformative
otherwise. The reactor's nonlinear transfer function acts as a physical
cipher. The same device can invert the transform because $\theta$ is
known; a different device is unlikely to succeed because the forward map
$S \to C_{0:T}$ given $\theta$ passes through bifurcation structure that
creates exponentially many branches, making inversion without $\theta$
ill-posed.

Non-limiting applications include physical-layer encryption for optical
communication (the \cb is the ciphertext; the device
is the key), device-locked content (output is meaningful only through
the originating device or a physical clone, which is microstructure-clone-hard to
produce), and physical key exchange (two devices that share a coupling
history can establish shared secrets from the joint attractor without
digital key exchange protocols).

\paragraph{Obscure (non-limiting).}
In some embodiments, the system irreversibly destroys semantic content
in the scene return while preserving signal energy. The scene remains
physically illuminated but its semantic content is scrambled beyond
reconstruction. The objective is to minimise parsability by any
downstream model. The mechanism drives the reactor into a chaotic regime
where scene features are nonlinearly mixed, with the maximum Lyapunov
exponent indicating the rate of trajectory divergence in the reactor's
phase space, producing rapid mixing and loss of parsable scene structure.

Obscure is distinct from Erase in that Erase targets specific features
(selective) while Obscure destroys all semantic content
(non-selective). Obscure is distinct from Encrypt in that Encrypt is
reversible with knowledge of $\theta$ while Obscure is irreversible:
the information is thermodynamically dissipated.

Non-limiting applications include anti-surveillance (making a scene
unparsable to hostile optical sensors by driving the reflected light
field into a chaotic regime), physical adversarial perturbation
(injecting modulations into the physical light field that break a
sensor's model, as distinct from digital adversarial examples), and
communication denial (preventing any receiver from extracting useful
information from the optical channel).

\paragraph{Compress (non-limiting).}
In some embodiments, the system reduces the effective dimensionality of
the \cb while preserving task-relevant information. The
objective is $\min d_{\mathrm{eff}}$ subject to
$I(C_{0:T}; S_{\mathrm{task}}) \geq \tau$: find the smallest attractor
that preserves the target information. The attractor itself is the
compressed representation, and the attractor dimension $d_{\mathrm{eff}}$
is the physical bottleneck width. This implements an information
bottleneck enforced by physics, in contrast with a digital regulariser.

Non-limiting applications include bandwidth-limited transmission
(compress scene information into a minimum-dimension attractor before
transmission, with a matched reactor at the receiver acting as a
decoder) and resource-constrained verification (reduce
$d_{\mathrm{eff}}$ to the minimum needed for \TB, discarding
irrelevant scene complexity).

\paragraph{Amplify (non-limiting).}
In some embodiments, the system increases differential sensitivity to
targeted scene features. The objective is
$\max \partial C_{0:T} / \partial S_{\mathrm{target}}$. The mechanism
exploits critical amplification near a bifurcation, where the reactor
acts as a tunable nonlinear amplifier whose gain is
Fisher-curvature-dependent. In some embodiments, near the bifurcation
boundary, gain increases sharply for features aligned with the critical
eigenvector, subject to hardware saturation limits and noise floors.

Amplify is distinct from Read: Read maximises total scene information;
Amplify maximises differential sensitivity to a specific feature,
possibly at the cost of saturating others. Non-limiting applications
include sub-threshold detection (detecting signals below the noise
floor by exploiting critical amplification) and resonant enhancement
(amplifying scene features near a specific frequency by tuning the
reactor's natural frequency to match).

\paragraph{Filter (non-limiting).}
In some embodiments, the system selectively passes information about
desired scene features while actively suppressing undesired features.
The objective is $\max I(C_{0:T}; S_{\mathrm{pass}})$ subject to
$\min I(C_{0:T}; S_{\mathrm{reject}})$. The reactor's attractor
structure defines a physical passband: features that couple strongly
to the attractor pass; features orthogonal to the attractor's stable
manifold are rejected. The attractor's geometry is the filter shape.

Non-limiting applications include spectral selection (passing one
frequency band of scene dynamics while rejecting others), semantic
filtering (passing structural features while rejecting textural
features, or the reverse), and source separation (if the scene
contains multiple superimposed sources, training multiple filter
configurations, each passing one source, with the reactor's nonlinear
dynamics acting as a physical implementation of independent component
analysis).

\paragraph{Predict (non-limiting).}
In some embodiments, the system extrapolates scene dynamics forward in
time. The objective is
$\min_{\pi} \E\|S_{t+\tau} -
\hat{S}_{t+\tau}(C_{0:t})\|^2$. If the reactor is entrained with the
scene (Yoked operation) and the coupled system occupies a known
attractor, the attractor's geometry constrains future states. Reactor
memory media store the trajectory needed for prediction. The prediction
horizon is set by the conditional Lyapunov exponent: chaotic scenes
yield short horizons (exponential divergence); periodic or quasiperiodic
scenes yield long or indefinite horizons.

Non-limiting applications include predictive maintenance (entraining
with a machine's vibration signature and predicting future state from
attractor dynamics), biological forecasting (entraining with
physiological rhythms and predicting upcoming state changes), and
predictive sensing (forecasting where the scene will be and pre-pointing
the illumination to optimise the next measurement).

\paragraph{Sign and watermark (non-limiting).}
In some embodiments, the system embeds device identity into the
\cb. Two modes are distinguished:

\emph{Overt signing}: the device's identity is explicitly present in the
\cb. The device's holonomy---accumulated geometric phase
from cyclic parameter sweeps---is a natural device-specific signature
that is verifiable by any party with access to the family's attractor
bundle description.

\emph{Covert watermarking}: the device's microstructure-specific
response embeds a subtle, broadband signature in the \cb
that survives downstream processing (compression, cropping, colour
correction) but is detectable only by a matched device or with knowledge
of the detection key. In some embodiments, the watermark is trained
adversarially to survive a declared suite of transforms while remaining
undetectable by a discriminator.

Non-limiting applications include provenance (every \cba
window carries an empirically hard-to-forge physical signature of the originating
device), chain of custody (successive devices add their signatures;
the output carries a physical audit trail), and anti-deepfake (genuine
\RK output carries a physical watermark that synthetic
content is impractical to reproduce without the actual device's $\theta$).

\paragraph{Forget (non-limiting).}
In some embodiments, the system selectively erases accumulated
information from reactor memory media (for example CRT phosphor state,
phase-change material configuration, or other persistent media) while
preserving other stored information. The objective is
$\min I(\mathrm{Memory}_{\mathrm{new}}; S_{\mathrm{target}})$ subject
to preserving
$I(\mathrm{Memory}_{\mathrm{new}}; S_{\mathrm{other}})$. Selective beam
deflection or localised overwrite targets specific memory locations.
Each erasure cycle incurs an irreducible thermodynamic cost.

Forget is distinct from Erase: Erase targets scene information in the
\emph{output}; Forget targets accumulated information in the
\emph{memory state}. Erase operates on the present; Forget operates
on the past.

Non-limiting applications include rolling-window operation
(continuously forgetting information older than a declared horizon while
retaining recent history, implementing a physical first-in-first-out
buffer), privacy compliance (auditable erasure of specific data from device
memory under declared hardware and environmental assumptions), and selective pattern
retention (forgetting undesired acquired patterns while retaining
desired ones, a physical analogue of managing catastrophic forgetting
in neural networks).

\paragraph{Create (non-limiting).}
\label{par:create}
In some embodiments, the system generates output in the absence of
external scene input. The scene loop is either disabled
($S^{\mathrm{scene}} = S_{\mathrm{term}}$, as in PoliePuter mode) or
presents a featureless input, and the reactor's own dynamics---driven by
internal noise, spontaneous emission, or stochastic fluctuations
amplified near a bifurcation---produce structured output.

The mechanism exploits symmetry breaking at bifurcation points: when the
reactor's operating point crosses a bifurcation boundary, infinitesimal
noise is amplified into macroscopic pattern selection. The specific
pattern selected is determined by the reactor's microstructure and the
instantaneous noise realisation, making each creation event
device-specific and non-reproducible. In some embodiments, the creation
process is guided by a bias or seed embedded in the control protocol
(for example a weak stimulus that biases which branch of the bifurcation
is selected) while the detailed structure is filled in by the reactor's
autonomous dynamics.

Non-limiting applications include physical random number generation
(the bifurcation-amplified noise provides entropy that is grounded in
physical process rather than algorithmic pseudorandomness), generative
art (the reactor produces patterns that are neither predetermined by
software nor directly copied from a scene, but emerge from the device's
own nonlinear dynamics), and stochastic search (the reactor's
exploration of its own attractor landscape, guided by a weak objective
signal, discovers configurations that directed optimisation might miss).

Create is distinct from Write: Write imposes a target pattern $Y^*$
determined externally; Create produces a pattern from within the device's
own dynamics, with no external target. The \cb records
the creation event, and protocol digests record the control conditions
under which creation occurred.

\paragraph{Invert (non-limiting).}
In some embodiments, the system or an external party attempts to
recover the scene $S$ from the \cb $C_{0:T}$. Inversion
difficulty depends on knowledge of the device parameters $\theta$:
with known $\theta$, inversion is a well-posed (if nonlinear)
optimisation (this is the legitimate operator's situation); without
known $\theta$, inversion requires simultaneously estimating $\theta$
and $S$ from $C_{0:T}$, and the bifurcation structure creates
exponentially many branches that are unidentifiable from finite data.

The ratio of inversion difficulty (unknown $\theta$ divided by known
$\theta$) is the security margin. This ratio is maximised at high
Fisher curvature (near bifurcations) where the mapping is maximally
$\theta$-sensitive. In some embodiments, the Invert primitive is used
as an internal diagnostic: measuring how easily a digital model can
invert the device's output provides a lower bound on hardness.

\subsubsection{Composite operations (non-limiting)}

In some embodiments, primitive operations are composed to produce
higher-level capabilities. Non-limiting composites include:

\begin{description}[style=nextline, leftmargin=2em]
  \item[Redact (Read + Erase).] Identify specific content in the
    \cb via a \LI head, then actively destroy it via
    the Erase primitive, producing a \cb that contains
    task-relevant information but physically lacks the redacted content.
  \item[Physical steganography (Write + Watermark).] Render target
    content via \RT while simultaneously embedding a
    covert device signature via the Watermark primitive.
  \item[Predictive sensing (Predict + Read).] Forecast scene state
    from Yoked coupling dynamics, then optimise the next \LI
    measurement for the predicted state.
  \item[Authenticated encryption (Encrypt + Sign).] Produce
    device-locked content (meaningful only with knowledge of $\theta$)
    that simultaneously carries an overt provenance signature.
  \item[Privacy-preserving authenticated sensing (Read + Erase +
    Authenticate).] Authenticate the scene and device via \TB,
    extract task-relevant information via \LI, and physically erase
    personally identifiable information via the Erase primitive, all
    within a single feedback session. The \cb that exits
    the session is designed so that the erased content is destroyed
    in the analogue domain before
    digital recording, making reconstruction infeasible under the
    declared threat model.
\end{description}

\subsubsection{Training vector space and regime geometry (non-limiting)}

In some embodiments, a protocol specification is a point in a
multi-dimensional objective space, with one axis per primitive operation.
A weight vector
\begin{multline*}
  \bigl(w_{\mathrm{read}},\;
  w_{\mathrm{write}},\;
  w_{\mathrm{auth}},\;
  w_{\mathrm{yoke}},\;
  w_{\mathrm{erase}},\;
  w_{\mathrm{encrypt}},\;
  w_{\mathrm{obscure}},\\
  w_{\mathrm{compress}},\;
  w_{\mathrm{amplify}},\;
  w_{\mathrm{filter}},\;
  w_{\mathrm{predict}},\;
  w_{\mathrm{sign}},\;
  w_{\mathrm{forget}},\;
  w_{\mathrm{create}}\bigr)
\end{multline*}
defines the Pareto trade-off for a specific task. The three operating
regimes and the Yoked modifier occupy specific faces of this polytope
(for example, pure \LI is the face where $w_{\mathrm{read}}$
dominates). Named composites occupy identified edges or interior points.
Training is the process of finding the operating point in this polytope
that maximises a task-specific utility function, using the natural
gradient defined by the Fisher information metric of $\Theta$.

% ----------------------------------------------------------------------
\subsection{Cross-Scene Regime Combinations (Non-limiting)}
\label{sec:cross-scene}
\label{sec:cross-scene-regime-combinations}
% ----------------------------------------------------------------------

In some embodiments, a device with multiple optical ports (scene-facing
channels, reactor-facing channels, or distinct field-of-view sectors)
operates different ports in different regimes simultaneously. In some
embodiments, two or more devices observe the same scene from different
vantage points, each in a potentially different regime.

The resulting combination space has the structure of a product of
per-port regime simplices. For a device with two ports A and B, each
port occupies a point on the
$(\lambda_{\mathrm{TB}}, \lambda_{\mathrm{L}}, \lambda_{\mathrm{RT}},
\alpha)$ prism described in Section~\ref{sec:interpolation}, and the
joint operating point lies on the product of two such prisms.

\paragraph{Non-limiting cross-port combinations.}
\begin{description}[style=nextline, leftmargin=2em]
  \item[Authenticated rendering (\TB on port A, \RT
    on port B).] Rendering is gated on continuous verification:
    the \RT output is produced only while the \TB
    port confirms device and scene authenticity. If the verification
    score drops below a threshold, rendering is suspended or reverted
    to a safe default.
  \item[Differential sensing (\LI on port A, \LI on port B).]
    Two scenes or two views of the same scene are compared through the
    same device. The device's transfer function cancels in the
    difference, isolating scene-to-scene variation. In some embodiments,
    one port observes a reference and the other observes a sample,
    enabling device-mediated comparison with cancellation of systematic
    device effects.
  \item[Scene translation (\LI on port A, \RT on
    port B).] Information read from scene A is used to drive the
    rendering applied to scene B. The device performs a physical
    scene-to-scene translation: the style, content, or structure of one
    scene modulates the appearance of another through the reactor's
    nonlinear dynamics.
  \item[Multi-surface rendering (\RT on port A, \RT on port B).] Two scene surfaces are rendered simultaneously
    with shared reactor dynamics. Temporal coordination through the
    reactor produces coupled visual effects that may be difficult or
    impractical to achieve with a pair of independent projectors.
  \item[Yoked cross-scene coupling (Yoked on port A, Yoked on port B).]
    Two scenes are dynamically coupled through the device. The device
    mediates a physical interaction between scenes that are not in
    direct contact. Transfer entropy flows $S_A \to D \to S_B$ and
    $S_B \to D \to S_A$ create a device-mediated coupling channel.
\end{description}

\paragraph{Multi-device combinations (non-limiting).}
In some embodiments, multiple devices observe the same scene, each
operating in a potentially different regime:

\begin{description}[style=nextline, leftmargin=2em]
  \item[Mutual device authentication.] Two devices both operate in
    \TB on the same scene. Cross-device consistency of the
    \cb, conditioned on the shared scene, provides a
    basis for mutual authentication without a trusted third party.
  \item[Consensus sensing.] Two or more devices operate in \LI on
    the same scene. Agreement across independent device observations,
    each with distinct $\theta$, strengthens scene claims beyond what
    any single device can provide.
  \item[Distributed authenticated rendering.] One device operates in
    \TB to authenticate the scene; a second device operates in
    \RT, conditioned on the first device's verification
    output. Verification and rendering are physically separated but
    cryptographically linked via committed \cba atoms.
\end{description}

\paragraph{Temporal sequencing (non-limiting).}
In some embodiments, a device transitions between regimes over time
within a single session. Non-limiting sequences include: a \TB
phase followed by \LI followed by \RT (authenticate,
then sense, then render); interleaved \TB checks during Yoked
operation (integrity monitoring during entrainment); and alternating
\LI and \RT passes (iterative refinement where each
sensing pass informs the next rendering pass). In some embodiments, the
protocol digest records the regime sequence, transition times, and
any gating conditions.

% ----------------------------------------------------------------------
\subsection{Degenerate and Boundary Cases}
\label{sec:degenerate}
% ----------------------------------------------------------------------

The following limiting cases define the boundary conditions of the
regime framework and the baselines against which capability and
hardness aspects of this disclosure are measured.

\paragraph{Trivial reactor.}
If the reactor transfer function is the identity
($g_{\mathrm{reactor}} = \mathrm{id}$), all regimes collapse to
conventional optics: \LI becomes a standard camera, \RT becomes a standard projector, \TB provides no
physically unclonable function, and Yoked operation reduces to
projector--camera feedback (which has been described in the projection
mapping literature). This is the baseline that every hardness and
capability claim is measured against.

\paragraph{Fully linear reactor (non-limiting).}
If the reactor has a linear (but non-identity) transfer function---for
example a gain stage, colour filter, or linear time-invariant (LTI)
medium with memory---the input--output map is reproducible by
standard linear-systems estimation (finite impulse response
identification) given a sufficient number of challenge--response
pairs.  Such a reactor provides limited empirical hardness because
emulation is straightforward once the transfer function is
characterised.  Hardness claims in this document rely on reactor
\emph{nonlinearity} (saturation, bistability, chaotic mixing),
\emph{bifurcation topology}, and \emph{idiosyncratic microstructure}
that make finite-data emulation intractable under the declared
attacker families.  A fully linear reactor is therefore a second
baseline, stronger than the identity case but weaker than any
nonlinear embodiment.

\paragraph{Trivial scene.}
If the scene is absent, featureless, or replaced by a calibration
target, \LI returns no useful information about the external world,
\RT has no target substrate, and Yoked operation has no
dynamic partner. In such embodiments, Inner \LI
(self-calibration) and \TB remain operational as the primary
meaningful modes. The device becomes a pure microstructure challenge-response oracle. In some
embodiments, this is the declared mode for device enrolment,
commissioning, and self-calibration.

\paragraph{Broken feedback (open loop).}
If the feedback path is interrupted, all closed-loop features are lost.
\LI reverts to passive sensing. \RT reverts to static
projection. \TB loses closed-loop hardness and reverts to static
challenge-response (Level~0). Yoked operation is not available in this
mode. In some
embodiments, the system detects feedback interruption and falls back to
a reduced-capability mode comprising at least static challenge-response
verification, and the interruption event is logged in the meter
envelope.

\paragraph{Edge of synchronisation.}
When the conditional Lyapunov exponent is near zero ($\CLE
\approx 0$), Fisher information peaks, and all regimes are
simultaneously at their most powerful and most fragile. This is the
operating point of maximum information capacity but also maximum
sensitivity to perturbation. In some embodiments, the operating policy
targets a neighbourhood of this boundary rather than the boundary
itself, trading off peak capacity against robustness, with the distance
from the boundary recorded in the meter envelope as a stability margin.

\paragraph{Bifurcation crossing during operation.}
If the device's operating point drifts across a bifurcation surface $B$
during a session, the attractor changes topology mid-operation. The
regimes respond differently to this event: for \TB, it is a
verification failure (the expected attractor has changed); for \LI,
it is a recalibration event (switch to a new attractor model); for
\RT, it may be a discontinuity in the output; for Yoked
operation, it may be a desirable discovery event (the coupled system
finding a new dynamical regime) or a catastrophic instability depending
on the application. In some embodiments, the protocol digest records whether a
bifurcation crossing occurred during the session, and meters report
pre- and post-crossing attractor identity.

\paragraph{Identical-twin devices.}
Manufacturing tolerance may produce devices whose attractor bundles are
metrically indistinguishable at a given curvature resolution. Higher
verification levels (Level~2 or Level~3) or longer trajectories resolve
this, but there exists a minimum interrogation complexity below which
two devices are functionally equivalent. This is not a failure of the
framework; it bounds the security claim. The Cram\'er--Rao bound on
device-parameter estimation from finite observations sets this
resolution floor.

% ----------------------------------------------------------------------
\subsection{Human-Coupled Operations (Non-limiting)}
% ----------------------------------------------------------------------

In some embodiments, a human observer or participant is included in
the feedback loop as an active agent rather than a passive viewer. The
human's sensory and motor channels become part of the coupled system,
and the information operations taxonomy
(Section~\ref{sec:info-operations}) extends to include human-mediated
primitives. Each human-coupled primitive defines a distinct training
objective involving the human's response.

\paragraph{HC-Read (non-limiting).}
The human reads the \RK's output through one or more sensory
channels (visual, auditory, haptic). The objective is
$\max I(C_{0:T}; R_{\mathrm{human}})$, where $R_{\mathrm{human}}$
denotes the human's measurable response (gaze, reaction time, task
performance, physiological signal). The system optimises its output to
maximise the human's information uptake, adapting illumination,
rendering, or sonification to the human's perceptual capabilities and
attentional state. In some embodiments, the human's response is measured
via eye tracking, EEG, galvanic skin response, or task performance
metrics, and these measurements are logged in the \cb
as auxiliary channels.

\paragraph{HC-Write (non-limiting).}
The human provides input that the \RK incorporates into its
feedback loop. Input channels include gaze direction, gesture, voice,
touch, physiological signals (heart rate, respiration, muscle tension),
and explicit control inputs. The objective treats the human's input as
a scene signal: $\max I(H_{\mathrm{input}}; C_{0:T})$, where
$H_{\mathrm{input}}$ is the human's input stream. In some embodiments,
multiple human input channels are fused through the reactor's nonlinear
dynamics, producing a physical multimodal integration that is
device-specific and logged.

\paragraph{HC-Verify (non-limiting).}
The human participates in the verification loop by confirming or
denying the plausibility of the \RK's output. The human acts
as an auxiliary verifier whose judgment is logged alongside automated
meter scores. In some embodiments, the system presents verification
challenges that require human perceptual judgment (for example
detecting subtle visual artefacts that automated metrics miss), and the
human's response is bound to the protocol digest. In some embodiments,
HC-Verify implements a human-in-the-loop gate where automated
operations proceed with human confirmation.

\paragraph{HC-Entrain (non-limiting).}
The human and the \RK mutually entrain: the device's
output influences the human's physiological state, and the human's
physiological response modulates the device's feedback loop. This is
distinct from software-mediated biofeedback: the physical coupling
between the device's optical output and the human's neural or
physiological response IS the intervention, with no digital
intermediary in the coupling path. The objective is
$\max \min(\TE_{D \to H},\, \TE_{H \to D})$, where
$D$ and $H$ denote device and human respectively. Non-limiting
applications include neurofeedback (EEG-optical coupling where
brain rhythms and reactor dynamics mutually entrain), rehabilitation
(motor-optical coupling for movement training), and biofeedback
(entraining physiological rhythms with stable
reactor dynamics for regulation).

\paragraph{HC-Create (non-limiting).}
The human guides the creation process (Section~\ref{par:create}) by
providing weak biases that influence which branch of a bifurcation the
reactor selects, while the detailed structure is filled in by the
reactor's autonomous dynamics. The human's creative intent is expressed
through continuous low-bandwidth input (for example a joystick, gaze
direction, or emotional valence derived from physiological signals)
that biases the reactor's symmetry-breaking, producing output that is
neither fully determined by the human (the reactor contributes its own
nonlinear dynamics) nor fully autonomous (the human's intent shapes the
large-scale structure). The \cb records both the human's
input and the reactor's response, making the co-creative process
auditable.

% ----------------------------------------------------------------------
\subsection{Coupling-history effects and physical key exchange
(non-limiting)}

In some embodiments, the joint attractor of two coupled systems
produces a shared dynamical history that is empirically hard to
reconstruct without access to both systems and their coupling
record.  In some embodiments, coupling-history features are used
for trust-topology construction, physical key exchange from
correlated measurements, and cost-graded verification where closely
coupled devices verify more cheaply than loosely coupled devices.
Higher-layer coordination and agent-autonomy applications of
coupling-history reification are described in the related Filing 2 application
(optional, non-essential).

\section{Meters and Training Games}
\label{sec:meters-training}
% ======================================================================

Meters may act directly on the \cb or on derived
latents. In
some embodiments, meter families include: (i)~verisimilitude
discriminators that distinguish physical from synthetic traces,
(ii)~hardness estimators that operationalise empirical difficulty of
simulation or forgery under bounded adversaries, (iii)~calibration
probes that track alignment and drift, and (iv)~coupling meters that
track Yoked operation quality.

\subsection{Yoked-specific meters (non-limiting)}

In some embodiments, meters specific to Yoked operation include:
\begin{description}[style=nextline, leftmargin=2em]
  \item[Transfer entropy balance]
    $\TE_{S \to R}$ and $\TE_{R \to S}$ estimated from the \cb,
    together with their ratio and absolute values.
  \item[Conditional Lyapunov exponent estimate]
    An estimate of $\CLE$ from the time series, indicating proximity to
    the synchronisation boundary.
  \item[Spectral coherence]
    Cross-spectral density between scene return and reactor output,
    at each frequency band.
  \item[Order parameter]
    A scalar summary of coupling quality, for example the Kuramoto
    order parameter $r = |N^{-1} \sum_j e^{i\phi_j}|$ computed over
    coupled oscillator phases.
  \item[\Yd]
    The declared target $\alpha_{\mathrm{target}}$ logged in the
    protocol digest and the meter-derived achieved $\hat{\alpha}$
    estimated from the \cb.
  \item[Coupling estimation confidence]
    An estimate of the statistical reliability of coupling quality
    measurements, gating Yoked claims on sufficient SNR.
\end{description}

Training may include supervised learning, self-supervised learning
(including without limitation SimCLR, MoCo, BYOL, SwAV, DINO,
masked-autoencoder methods such as MAE, masked language modelling,
joint-embedding predictive architectures including JEPA, and
joint vision-language methods such as CLIP, ALIGN, and SigLIP),
contrastive objectives, generative modelling (including diffusion,
score-based, flow-matching, or consistency models), reinforcement
learning for protocol selection, and hardware-in-the-loop
optimisation such as SPSA, evolution strategies, natural evolution
strategies (NES), covariance matrix adaptation evolution strategy
(CMA-ES), population-based training (PBT), genetic algorithms,
particle swarm optimisation, and rank-structured evolution-strategy
methods (including without limitation rank-$r$ perturbation
schemes that recover batched-matrix-multiplication arithmetic
intensity for hyperscale training) for non-differentiable
parameters.

\subsection{Roles and multi-agent framing (non-limiting)}

In some embodiments, design and operation are framed as a multi-agent
interaction. Alice selects configuration and protocols and trains
control policies or adaptation procedures. Bob trains or operates
verifiers and meters that evaluate the \cb for plausibility,
calibration, or hardness. Eve acts as an adversary attempting to spoof,
simulate, or perturb the \cb or downstream decisions under a
bounded resource model. This framing is used primarily for verification
(\TB), but may also be used for \LI and \RT
training when adversarial robustness or distribution control is desired.

\paragraph{Eve as a system-wide hardening principle (non-limiting).}
In some embodiments, Eve is not confined to a single protocol layer but
serves as the adversarial-hardening archetype operating at every level
of the \RK architecture.  At the physical-channel level, Eve
attempts to spoof or forge \cbs under bounded resources, driving
the empirical hardness that underpins \TB verification.  At the
information-theoretic level, Eve observes a degraded view of the
physical channel (the wiretap lens of
Section~\ref{sec:wiretap-advantage}), and protocol design aims to widen
the Bob--Eve advantage gap.  At the meter-training level, Eve's
best-effort attacks under frozen meters define the hardness index that
calibrates security claims.  At the semantic level, a semantic Eve
attempts to produce alternative descriptions, inferences, or
behavioural outputs that pass semantic verification
(Section~\ref{sec:semantic-truth-beam}); in PoliePuter mode, semantic
Eve probes the agent's reasoning and decision-making under adversarial
prompts, perturbations, and distribution shifts, deliberately pushing
toward unsafe or deceptive operating regimes so that the system is
hardened against regions that natural drift might otherwise reach
undetected.  At the training level, adversarial training
environments (described in the related Filing 2 application; optional) instantiate the same adversarial principle in
the agent's training trajectory, stress-testing behavioural policies
under declared floors and meter envelopes.  In each case, the adversarial role is
constructive: by systematically probing failure modes under bounded
attacker families, Eve drives the system toward empirically robust
configurations at every layer of the architecture.

\paragraph{Eve injector and jammer embodiment with kappa-only
  training variant (non-limiting).}
In some embodiments, Eve operates as an active injector or jammer
rather than only as a passive observer: Eve is permitted to inject
a bounded signal $j_{\psi}(c)$ into the optical, electromagnetic,
acoustic, or coupled physical path on which the reactor operates,
subject to a declared injection power constraint, a declared
spectral support, a declared spatial-or-channel mask, and a
declared duty-cycle envelope, all of which are committed to the
protocol digest. The injected signal is treated as part of the
adversarial setting under which Bob's meters and \TB verification
are required to retain calibrated performance, and Eve's
admissible attack family includes constant, swept-frequency,
modulated, and replay-style injections. In further embodiments, a
\emph{kappa-only training variant} is supported in which only the
physical kernel parameters $\kappa$ of the reactor are updated
during a training round, while ancillary parameters including
optical and electrical attenuators, dampers, emitter counts,
amplifier biases, and timing-channel keys are held externally
fixed at declared settings recorded in the protocol digest.
Kappa-only training is used when training-induced changes to
ancillary parameters would otherwise confound hardness or
robustness measurements, or when a fixed-ancillary baseline is
required for cross-device comparability. In some embodiments, the
Eve injector embodiment and the kappa-only training variant
compose: Eve injects under declared envelopes while only $\kappa$
is updated, producing an adversarial-robust trained reactor under
fixed ancillary settings.

The system-wide counterpart to Eve's adversarial principle is the
\emph{Iris} principle: at every layer where Eve probes for failure,
Iris expands information capture, ensuring that the full signal
content of the physical channel is available for verification,
sensing, and calibration.  A well-designed \RK loop
alternates between Iris expansion (maximising what the channel can
resolve) and Eve hardening (ensuring that what is resolved cannot be
forged), with Demeter providing the bidirectional management that
maintains stability across the alternation.  The three primary stances (and the Harmonia coupling modifier) are
defined formally in Section~\ref{sec:stances}.

In some embodiments, \RK design and operation are treated as
a two-timescale game. On a slow timescale, Alice commits to a hardware
and protocol configuration $\theta^A$ (and associated admissible
protocol classes), after which Bob and Eve adapt toward best responses
within bounded resource classes (meters, models, and query budgets).
This can be read as a Stackelberg game with Alice as leader and Bob/Eve
as followers, where empirical hardness indices are defined against these
best-response followers. On a fast timescale, for fixed physical
configurations, Bob and Eve interact with the live device in repeated
challenge--response episodes under logged protocol digests.

\paragraph{Programmable-scene jamming game (non-limiting).}
In some embodiments, the scene includes a programmable surface such as
an e-ink display that renders a time-indexed plaintext, watermark, or
procedural field whose parameters encode a message. In a non-limiting
multi-agent game, Bob is a cooperative decoder trained to recover the
plaintext from the \cb under the declared protocol, while Eve
attempts to decode the same plaintext from a degraded view and/or jam
the channel by injecting co-illumination or structured interference.
Alice selects kernel and protocol families to increase a secrecy-margin
objective under meter-envelope constraints:
\[
  \max_{\kappa,\,\pi}\;
  \Bigl(\mathcal{U}_{\mathrm{Bob}}(\kappa,\pi)
  - w_{\mathrm{adv}}\,\mathcal{U}_{\mathrm{Eve}}(\kappa,\pi)\Bigr),
  \quad\text{subject to meter envelopes and safety bounds.}
\]

\paragraph{Multi-beam dampers (non-limiting).}
A power-allocation vector $a \in \R_+^B$ distributes probe or carrier
power across $B$ beams or logical channels distinguished by space, time,
frequency, and/or polarisation. A damper vector
$\delta \in [0,1]^B$ attenuates designated beams. Alice optimises
$(a,\delta)$ jointly with $\kappa$ and protocol choices to increase
Bob's SNR and reduce Eve's leverage, subject to $\sum_b a_b \le
P_{\max}$ and per-beam safety bounds.

\subsection{Meter families (non-limiting)}

In some embodiments, meters operate on \cba windows and/or
derived latents and output scalar scores. Verisimilitude meters score
whether a trace is consistent with the empirical distribution of
physical runs. Hardness meters estimate difficulty for a specified
adversary family to generate traces that pass verification, optionally
producing a hardness index. Calibration meters track alignment, drift,
timing integrity, and cross-device consistency. Meters may be learned
discriminators, energy-based models, statistical tests, or hybrids. In
some embodiments, meters are distilled or simplified for embedded
operation after training.

\paragraph{Sequential hypothesis-testing meters (non-limiting).}
In some embodiments, verifier decisions are made by sequential
hypothesis-testing procedures rather than by a fixed-length batch
score. A sequential probability ratio test, anytime-valid confidence
sequence, e-value process, alpha-spending rule, or other declared
sequential test updates the accept, reject, or continue state as
committed bundle windows arrive. The null and alternative hypotheses,
stopping boundaries, maximum sample count, error budgets,
multiple-testing correction, and policy for inconclusive runs are
declared before the window opens and are included in the verifier's
protocol digest.

\subsection{Threat model taxonomy (non-limiting)}

Non-limiting attack families include replay of previously recorded
traces, model-based simulation using learned emulators, deepfake
synthesis, adversarial perturbations of genuine traces, co-illumination
interference, hardware tampering, and semantic misattribution. In
preferred embodiments, replay is mitigated using freshness challenges
and commitments, while simulation and perturbation are addressed using
verisimilitude and hardness meters and protocol design.

\subsection{Verisimilitude discriminator training (non-limiting)}
\label{sec:verisimilitude-training}

In one embodiment, a discriminator is trained to distinguish physical
\cba windows from synthetic or attacked windows. Negative
examples may include simulator outputs, time-shuffled segments, injected
spikes outside hardware envelopes, and quantised or clipped traces.
After training, the discriminator may be used as a soft penalty when
training controllers and emulators so that rollouts remain within
empirically supported regions.

\begin{definition}[Verisimilitude functional]
\label{def:verisimilitude}
A \emph{verisimilitude functional} for a given configuration $\theta$
is any measurable map
\[
  V_\theta : \mathcal{C} \times \mathcal{X}_{\mathrm{aux}} \to \R
\]
from the \cb (and any auxiliary metadata $\mathcal{X}_{\mathrm{aux}}$, such as
protocol digests or claimed labels) to a real-valued score, such that
higher scores indicate better agreement with the empirical distribution
of physical runs under the same or comparable protocol classes. In many
embodiments $V_\theta$ is realised by a learned discriminator, but
hand-crafted statistics are also non-limiting embodiments.
\end{definition}

\paragraph{Verisimilitude checker for projector--camera loops
(non-limiting).}
In some embodiments, a projector--camera loop is treated as a special
case of a \RK in which a subset of the control $u(t)$ is an
emitted image or pattern and a subset of the observation $\mathbf{y}_t$
is a captured image. A verisimilitude checker is a learned meter that
outputs a score indicating whether windows of loop data are consistent
with the empirical distribution of genuine, correctly functioning loops
under comparable protocol digests. In preferred embodiments, the
verisimilitude checker is trained to be sensitive to physical tampering,
miscalibration, spoofing, and relay artefacts that preserve superficial
pixel statistics but break loop consistency.

\paragraph{Adversarial autoencoder for loop data (non-limiting).}
In some embodiments, verisimilitude training uses an
adversarial-autoencoder pattern: an encoder maps loop windows to
latents, a generator/decoder reconstructs one or both sides of the
loop, and a discriminator distinguishes genuine loop windows from
reconstructed, simulated, or attacked candidates. A reconstruction loss
encourages predictive closure of the loop under the logged protocol,
while an adversarial loss encourages the latent representation to
respect empirical physical constraints rather than only pixel-level
similarity.

\paragraph{Multi-pair coordinated training (non-limiting).}
In some embodiments, a plurality of projector--camera pairs operate in
a shared physical scene, and models are trained with cross-pair
constraints: candidate reconstructions or simulations are constrained
to be mutually consistent across pairs under the same scene coupling
and protocol digest.

\paragraph{Human response channels as auxiliary loop evidence
(non-limiting).}
In some embodiments, the loop includes an auxiliary response channel
derived from a human observer, including gaze direction, head pose, task
performance, comfort ratings, or physiological signals. The response
channel is logged as auxiliary evidence and may condition the
verisimilitude checker so that ``plausible for a human'' and
``physically consistent'' are jointly enforced under declared policies.

\paragraph{External and virtual-view checks (non-limiting).}
In some embodiments, given a learned 3D representation or a digital
twin, the system predicts what an additional camera should observe under
a declared emission schedule, and compares predicted versus measured
sequences as an additional consistency term.

\paragraph{Illustrative empirical verifier-stack behaviour.}
In some embodiments, the verisimilitude meter is evaluated using
plural independently trained verifier paths that are not co-trained
with one another and that do not share the same optimisation loop.
As a non-limiting illustration, independently trained residual,
binder, and feature-distance verifier paths may each separate genuine
committed-bundle/emission pairings from shuffled-emission pairings and
from trained-adversary pairings. Cross-deployment or cross-session
evaluation may be used to test whether the verifier paths are
responding to substrate-physics structure rather than to a
session-specific artefact. Such empirical observations are reported
as part of a declared metric envelope and do not limit the apparatus
to any particular architecture, dataset, numerical threshold, or
performance value.

\paragraph{Single-meter failure modes as motivation.}
In some embodiments, optimisation against a single verifier or
single binder-confusion loss can produce Goodhart-style failure modes.
Non-limiting examples include anti-binder texture artefacts that
confuse a target binder without corresponding to physically valid
emissions, and over-smoothed identity collapse in which continued
optimisation converges toward a degenerate or near-identity transform
that remains visually plausible while failing under an independent
verifier path. The multi-path independence, held-out evaluation, and
protocol-digest commitment discipline disclosed herein are configured
to make such single-meter pathways detectable. The named failure modes
are illustrative only; other failure modes may arise under different
adversary architectures, losses, modalities, or deployment conditions.

\paragraph{Load-bearing verifiers and diagnostic instruments.}
In some embodiments, load-bearing verifier paths are trained or
calibrated against committed-bundle ground truth, curated negative
examples, declared distractor sets, adversary-generated examples, or
other governance-relevant labels. Non-limiting examples include
binder-style emission-recovery models, correct-versus-distractor
retrieval models, feature-distance detectors, PatchCore-class
detectors, supervised or semi-supervised anomaly detectors, and
verisimilitude meters trained against declared negative example sets.
In further embodiments, unsupervised methods including
autoencoder-based substrate-signature detection, reconstruction-error
scoring, clustering, density estimation, and latent-space drift
analysis are used in parallel as diagnostic instruments for
depth-localisation, drift monitoring, calibration, exploratory audit,
or spot-check review. Such diagnostic instruments may support the
verifier stack without necessarily serving as load-bearing verifiers
under the protocol digest.



\subsection{Domain Discriminators and Domain Generators}
\label{sec:domain-discriminator-generator}

In some embodiments, a \RK apparatus comprises a domain
discriminator and a domain generator. A domain discriminator is a
model, meter, classifier, verifier, or scoring function configured
to estimate whether a bundle, feature, latent state, view, emission,
or derived artefact belongs to, or is consistent with, a declared
domain. A domain generator is a model, transformer, simulator, renderer,
or adapter configured to produce, translate, normalise, bridge, or
augment artefacts across domains while preserving one or more declared
constraints.

In some embodiments, domains comprise physical deployments, sessions,
rigs, environments, spectral bands, modalities, virtual worlds,
simulation environments, reactor states, governance states, privacy
release channels, calibration regimes, observer populations, or
adversary curricula. A domain label $d$ may be explicit, inferred,
hidden, soft, hierarchical, or represented as a continuous embedding.
A discriminator may score:
\[
    s_d = D_{\psi}(C_{0:T}, U_{0:T}, Z_{0:T}, O_{0:T}, \Pi; d),
\]
where $s_d$ indicates domain membership, consistency, likelihood,
distance, confidence, or anomaly score. A generator may produce a
domain-transformed artefact:
\[
    \tilde C_{0:T}^{(d')}
    =
    G_{\phi}(C_{0:T}^{(d)}, d, d', \Pi),
\]
where $d$ and $d'$ denote source and target domains. These equations
are non-limiting; the discriminator and generator may operate on raw
bundles, reduced features, latent representations, emissions, labels,
metadata, views, or meter outputs.

In some embodiments, the discriminator and generator are trained or
evaluated jointly. For example, a generator may attempt to translate
a bundle from a first deployment domain into a second deployment
domain, while a discriminator tests whether the translated bundle
preserves substrate evidence, emission identity, protocol consistency,
or governance-relevant signals. In some embodiments, the discriminator
is used to reject generator outputs that match superficial domain
statistics while failing a binder, held-out verifier, physical
re-execution, temporal-coherence meter, or anchor-consistency meter.

In some embodiments, the domain-generator primitive supports
counterfactual and privacy-preserving governance. For example, a
generator may produce a neutralised, redacted, stylised, compressed,
or domain-normalised representation for audit, while a domain
discriminator or held-out meter verifies that the representation
retains the declared governance signal and does not introduce a
forbidden transformation. In further embodiments, a domain
discriminator detects when a deployment has drifted outside an
enrolled domain and triggers calibration, selective opening, increased
monitoring, rescue mode, or governance escalation.

\paragraph{Domain primitive generalisation (non-limiting).}
The domain discriminator may be any scoring or decision function,
including without limitation a classifier, binary discriminator,
multi-class discriminator, contrastive model, density estimator,
normalizing flow, energy-based model, one-class detector, diffusion
residual scorer, retrieval model, conformal predictor, Bayesian model,
or ensemble. The domain generator may be any transform or synthesis
model, including without limitation an image-to-image translator,
video-to-video translator, diffusion model, GAN, flow model,
autoregressive model, renderer, simulator, style-transfer model,
domain-adaptation network, feature normaliser, calibration transformer,
physics model, or optimisation-based procedure. The term domain is not
limited to visual domains and may include optical, acoustic,
electromagnetic, temporal, semantic, governance, privacy, simulation,
or physical deployment domains.

\paragraph{Composition with virtual-camera and virtual-emitter
primitives (non-limiting).}
In some embodiments, the domain discriminator and domain generator
primitives of this subsection are instantiated using the virtual
camera and virtual emitter primitives disclosed in
Section~\ref{sec:virtual-views-emissions}, with the virtual camera
or virtual emitter acting as a domain generator and an external or
held-out meter acting as a domain discriminator. Conversely, the
virtual-camera and virtual-emitter primitives of
Section~\ref{sec:virtual-views-emissions} may use the domain
discriminator and domain generator of this subsection as held-out
or adversarial verifiers, for example in verisimilitude games and
re-capture or re-emission audits.



\subsection{Reality Encryption and Neutral Recording Generation}
\label{sec:reality-encryption-neutral-recording}

In some embodiments, a \RK apparatus implements reality
encryption by emitting, selecting, modulating, or transforming a
physical stimulation pattern such that, under a declared meter
envelope and attacker model, a resulting recording is intended or
declared to be verifiable, reconstructable, decryptable,
classifiable, or semantically useful only to an authorised opener
satisfying a declared key, protocol-state, model-state,
physical-response-model, authority-token, or selective-opening
condition. As with other security and hardness properties disclosed
in this specification, the ``only to an authorised opener'' property
is empirical under the declared attacker family and resource
budget, is time-indexed rather than unconditionally durable, and
does not constitute a formal cryptographic guarantee unless a
specific embodiment expressly specifies and parameterises such a
guarantee. The encrypted object may be the emitted signal, the
captured response, a latent representation, a derived label, a
meter output, a reconstruction path, or any combination thereof.

In some embodiments, an emission $u_t$ is derived from a protocol
state, key material, and optional payload:
\[
    u_t = E_{\omega}(K_t, \Pi_t, p_t, q_t),
\]
where $K_t$ may comprise a cryptographic key, threshold key,
PUF-derived value, physical one-way function output, session key,
authority key, or derived protocol state; $p_t$ may comprise a
payload, challenge, watermark, governance marker, or null payload; and
$q_t$ may comprise a timestamp, randomness beacon, ledger value,
prior digest, or external challenge. A recording $y_t$ produced in
response to $u_t$ may be bound into a committed bundle:
\[
    C_{0:T} = \{(u(t),y(t))\}_{t=0}^{T}.
\]
A verifier or authorised opener may recover or test a protected signal:
\[
    r_t =
    R_{\kappa}(C_{0:T}, K_t, \Pi_t, A_t),
\]
where $A_t$ may comprise authority records, selective-opening records,
attestation records, or proof records. These equations are illustrative
only and do not limit reality encryption to conventional encryption,
to digital payloads, or to any particular cryptographic primitive.

In some embodiments, the apparatus comprises a neutral recording
generator. A neutral recording generator produces a released recording
or representation that preserves one or more declared verification,
governance, calibration, or audit signals while suppressing, masking,
normalising, stylising, encrypting, or removing one or more protected
content signals. The neutral recording may be produced at capture
time, after capture, in a governance partition, in a privacy partition,
or through a coupled physical-digital process. A neutral recording
$\bar C$ may be generated as:
\[
    \bar C_{0:T}
    =
    N_{\phi}(C_{0:T}, \Pi, \lambda_{\mathrm{priv}},
             \lambda_{\mathrm{verify}}),
\]
where $\lambda_{\mathrm{priv}}$ denotes declared privacy constraints
and $\lambda_{\mathrm{verify}}$ denotes declared verification
constraints. The neutral recording may be accepted when it satisfies:
\[
    M_{\mathrm{verify}}(\bar C_{0:T}, C_{0:T}, \Pi)
    \in M_{\mathrm{accept}}
    \quad\text{and}\quad
    M_{\mathrm{privacy}}(\bar C_{0:T}, C_{0:T}, \Pi)
    \in P_{\mathrm{accept}}.
\]
The particular form of the privacy and verification meters is
non-limiting.

In some embodiments, reality encryption and neutral recording are
combined. For example, a device may emit a key-conditioned structured
illumination pattern whose captured response supports later
verification by an authorised party, while ordinary viewers receive a
neutralised recording, a classifier-only result, a perceptual hash, a
cohort statistic, or a proof that a floor condition was not violated.
In other embodiments, the emitted pattern is public but the readout
model is protected; the emitted pattern is protected but the readout
model is public; both are protected; or neither is protected but the
selective-opening path is authority-gated.

\paragraph{Reality-encryption generalisation (non-limiting).}
Reality encryption is not limited to any particular cipher, key size,
medium, or payload format. Key material may be pre-shared, derived
from correlated measurements, derived from a PUF or physical one-way
function, established interactively, generated by a threshold protocol,
derived from a ledger or randomness beacon, or produced by a
post-quantum key-establishment mechanism. The protected transformation
may use symmetric encryption, asymmetric encryption, hybrid
encryption, secret sharing, threshold cryptography, commitment
schemes, watermarking, steganography, spread-spectrum modulation,
structured illumination, learned codecs, physical unclonability,
selective disclosure, zero-knowledge proof, proof of correct
computation, or any combination thereof.

\paragraph{Neutral-recording generalisation (non-limiting).}
A neutral recording may comprise a redacted video, transformed video,
stylised recording, synthetic view, latent representation, embedding,
hash, perceptual hash, caption, scene graph, classifier output, audit
flag, proof object, accumulator state, compressed representation, or
selective-opening handle. Neutrality may refer to privacy neutrality,
semantic neutrality, identity neutrality, policy neutrality,
domain-neutrality, calibration neutrality, or any declared suppression
or invariance condition. The load-bearing property is not any single
neutralisation technique, but the apparatus capability to produce a
released representation whose retained and suppressed signal classes
are declared, tested, and bound to the protocol digest.

\paragraph{Composition with virtual-camera and virtual-emitter
primitives (non-limiting).}
In some embodiments, the neutral-recording generator and the
reality-encryption primitives of this subsection are composed with
the virtual-camera and virtual-emitter primitives of
Section~\ref{sec:virtual-views-emissions} to produce a released
representation that is generated from, or audited against, a
virtual view or a virtual emission rather than from a raw captured
recording. In such embodiments, the released representation
inherits the disclosure-regime tags committed to the protocol
digest and remains testable by the meters of this section.

\subsection{Edge-based structural meters (optional)}
\label{sec:edge-meters}

In some embodiments, a decoder maps the \cb to an image-like
representation and an edge or contour extractor produces a structural
map. Losses encourage preservation of major contours while permitting
style variation, enabling \RT constraints that preserve
content stability.

\subsection{Structural and semantic conditioning extractor family
(non-limiting)}
\label{sec:control-extractors}

In some embodiments, the edge-based structural meter of the preceding
subsection is one member of a broader family of declared structural
and semantic conditioning extractors that consume an image-like or
feature representation derived from a committed \cba
window via $\mathsf{RTDec}$ and produce a declared conditioning channel
or structural representation. Non-limiting members of the family
include:

\begin{itemize}
  \item \textbf{Soft-edge extractors}, including holistically-nested
    edge detection (HED) and learned soft-edge variants, producing a
    continuous-valued edge probability map.
  \item \textbf{Hard-edge and contour extractors}, including Canny,
    Sobel, structured forests, and related operators.
  \item \textbf{Line-art and scribble extractors}, producing sparse
    or stylised line representations of scene content.
  \item \textbf{Straight-line and segment extractors}, including
    M-LSD-style line-segment detection.
  \item \textbf{Monocular depth estimators}, producing per-pixel
    metric or relative depth maps.
  \item \textbf{Surface-normal estimators}, producing per-pixel
    surface-orientation vectors, optionally derived by photometric
    stereo from declared multi-illumination probe sequences.
  \item \textbf{Semantic, instance, and panoptic segmentation
    estimators}, producing pixel- or region-wise class, instance, or
    panoptic labels.
  \item \textbf{Optical-flow estimators}, producing per-pixel motion
    vectors between declared frame pairs.
  \item \textbf{Disparity and stereo-correspondence estimators},
    where geometry permits.
  \item \textbf{Body-surface UV-coordinate estimators} (such as
    DensePose-style maps), producing per-pixel correspondences to a
    declared body-surface chart.
  \item \textbf{The pose-feature extractor} $\mathsf{Pose}$ of
    the pose-conditioned \RT of Section~\ref{sec:pose-rt}.
  \item \textbf{Other declared structural or semantic measurement
    heads} whose identity, weights or version hash, input
    preprocessing, calibration state, output representation, and any
    confidence- or visibility-mask rule are committed to the protocol
    digest.
\end{itemize}

In some embodiments, the family is operated as a declared set
$\mathsf{ControlExt} = \{\mathsf{E}_1, \mathsf{E}_2, \dots\}$ in
which each $\mathsf{E}_j$ is an extractor of one of the above types,
declared by its identity, version, input preprocessing, output
representation, and calibration state, with the set membership
committed to the protocol digest. The set is open: members may be
added, deprecated, or version-bumped under the applicable
authority-separation policy and the version-history discipline of
the protocol digest.

\paragraph{Conditioning role versus meter role (non-limiting).}
In some embodiments, a member $\mathsf{E}_j$ of $\mathsf{ControlExt}$
participates in the RT compound objective in either or both of two
declared roles:

\begin{description}
  \item[Meter role.] A loss term penalises divergence between the
    extracted control signal $\mathsf{E}_j(Y^{\mathrm{dec}}_t)$ on the
    rendered output and a declared target control signal, encouraging
    the rendered output to preserve a structural or semantic
    property. The edge-based structural meter of
    Section~\ref{sec:edge-meters} is one such instantiation; depth,
    normal, segmentation, flow, and related preservation terms are
    further instantiations.
  \item[Conditioning role.] An extracted control signal from a
    declared reference (a committed reference frame, a committed
    probe-channel response under
    the temporal-multiplexed paired-acquisition corpus of Section~\ref{sec:style-corpus-paired}, or a declared external
    reference whose provenance type is committed) is supplied as a
    conditioning input to the controller or to a generator
    component, steering the rendered output toward a structurally or
    semantically constrained target.
\end{description}

In some embodiments, the same extractor $\mathsf{E}_j$ is operated
in both roles within a single run, with the conditioning-role
reference and the meter-role target separately declared and
committed.

\paragraph{Composition with the RT compound objective.}
In some embodiments, each member of $\mathsf{ControlExt}$ contributes
a term to $\mathsf{RTTerm}$ under
the RT-against-twin disclosure of Section~\ref{sec:rt-against-twin}, with its own declared weight
$\lambda_{\mathrm{rt},s}$, declared meter-partition assignment
(acceptance, evaluation, or shaping), and declared confidence- or
visibility-mask rule. The acceptance / evaluation partition is
preserved so that controller optimisation against any single
extractor does not expose all evaluative surfaces simultaneously.

\paragraph{Composition with twin-side training.}
In some embodiments, each member of $\mathsf{ControlExt}$ inherits
the twin-side evaluation of the RT-against-twin disclosure of Section~\ref{sec:rt-against-twin}: the
extractor is evaluated on $C^{\mathsf{Twin}}_{\mathsf{Win}_t}$ for
the digital twin or $C^{\mathsf{AnaTwin}}_{\mathsf{Win}_t}$ for the
analogue twin, with the synthetic-window construction and any
domain-gap correction (digital twin) or matching-fidelity envelope
(analogue twin) committed to the protocol digest. Control signals
extracted from twin-produced bundle windows are provenance-typed as
twin-predicted or analogue-twin-derived, consistent with the
provenance-typing discipline applied to $C^{\mathsf{Twin}}_{0:T}$
and $C^{\mathsf{AnaTwin}}_{0:T}$.

\paragraph{Negative-example construction.}
In some embodiments, negative examples for control-signal
discriminator and meter-residual threshold validation are
constructed by applying declared spatiotemporal edits to committed
recordings or to committed twin-predicted windows, with each edit
family, intensity parameter, and random-seed policy recorded or
hash-identified in the protocol digest, following the pattern of
the verisimilitude-training disclosure of Section~\ref{sec:verisimilitude-training} and the pose-conditioned
negative-example discipline of the pose-conditioned \RT of Section~\ref{sec:pose-rt}.

\paragraph{Scope note (non-limiting).}
For avoidance of doubt, members of $\mathsf{ControlExt}$ may be
drawn from conventional architectures and trained models. The
disclosed technical contribution is the committed-evidence,
meter-bounded, and twin-coupled composition of declared structural
or semantic conditioning extractors with the \RT
compound objective, the temporal-multiplexed paired-acquisition
training-data pipeline of the temporal-multiplexed paired-acquisition corpus of Section~\ref{sec:style-corpus-paired}, and
the provenance-typing discipline applicable to extracted control
signals across physical, twin-predicted, analogue-twin-derived, and
generator-produced bundle windows, rather than the underlying
edge-, depth-, normal-, segmentation-, flow-, line-, or related
estimator architectures.

\subsection{Canonical meter game (non-limiting)}

In some embodiments, a canonical evaluation game measures how well a
bounded generator can imitate physical outputs under a fixed challenge
distribution. A hardness index may be defined as the minimal simulator
capacity and training budget required to make simulated outputs
indistinguishable from physical outputs under a competent discriminator
on held-out tests. The index is empirical and depends on the chosen
simulator class, training budget, and meter strength.

\subsection{Physical sources of empirical hardness (non-limiting)}

In some embodiments, empirical hardness is supported by intrinsic noise
and chaos, history-dependent media, complex configuration manifolds, and
multi-channel dependence structure across invariance profiles. These
mechanisms support persistent empirical difficulty for bounded adversaries
under the engineering security framing of this disclosure.

\paragraph{Example noise and drift model (non-limiting).}
\label{sec:noise-model}
For concreteness, one operating model treats each \cba sample
as a noisy observation of an underlying ideal response after passage
through optics and media, with a combination of: (i)~Poisson shot
noise, (ii)~additive read noise on detector channels, and (iii)~slow
multiplicative drift factors. Misalignment and vibration may be modelled
as small random affine perturbations of projected and imaged
coordinates. Meters and objectives are trained under these models so
that acceptance regions distinguish ``honest noisy channel'' from
``attacked channel'' statistically, without assuming noise-free
operation.

\subsection{Digital twins and surrogate-guided optimisation
(non-limiting)}

In some embodiments, a digital twin is maintained to approximate the
behaviour of a physical reactor, scene loop, or coupled system under
protocol choices. The twin may be differentiable, partially
differentiable, or purely predictive, and may be updated using committed
\cba windows, meter summaries, and protocol digests. The twin
is assistive rather than authoritative: the physical \RK
remains the arbiter for verification, hardness, and meter-bounded
validity.

In some embodiments, the twin supports protocol selection by predicting
expected observation summaries and uncertainty under candidate protocols,
and choosing protocols that improve information gain, verisimilitude, or
projection-hardness subject to meter-bounded constraints. In some
embodiments, the twin is used to reduce physical trial count by
proposing candidate protocols, after which a subset is validated on
hardware.

\paragraph{Physics-based propagation surrogates (non-limiting).}
In some embodiments, the digital twin includes a physics-based
wave-propagation engine (for example ray tracing, beam-propagation, or
wave-optics solvers) together with learned priors for speckle
statistics, aberrations, or alignment drift. Such hybrids explore large
regions of configuration space in simulation, after which
hardware-in-the-loop optimisation fine-tunes around promising
configurations.

\paragraph{Neural and point-based twin renderers (non-limiting).}
In some embodiments, the digital twin's renderer or scene model is
implemented as a neural radiance field (Mildenhall et al. 2020), a
hash- or grid-accelerated radiance-field variant such as Instant-NGP
(M\"uller et al. 2022), a signed-distance-field representation, a
mesh- or point-based differentiable renderer, or a three-dimensional
Gaussian splatting representation (Kerbl et al. 2023). Twin parameters
are updated against committed \cba windows under the
declared matching-fidelity envelope, and the renderer family, finite
parameter serialisation, training corpus digest, optimiser settings,
and validation residuals are recorded with the twin version identifier.

\subsection{Hardware-in-the-loop optimisation (non-limiting)}
\label{sec:hitl-optimisation}

In some embodiments, protocol parameters are optimised with hardware in
the loop when gradients through the physical channel are unavailable or
unreliable. A non-limiting approach uses stochastic approximation (for
example SPSA) to estimate gradient directions of a scalar objective
$J(u)$ computed from committed windows and meter summaries. In some
embodiments, $J(u)$ combines task performance terms with
evidence-quality terms, including information-gain proxies,
verisimilitude or projection-hardness scores, and penalties for meter
violations.

A non-limiting two-sided SPSA-style estimate uses two evaluations per
step:
\[
  \widehat{\nabla}J(u)
  \approx \frac{J(u+c\Delta)-J(u-c\Delta)}{2c}\,(\Delta^{-1}),
\]
where $\Delta$ is a random perturbation vector, $c>0$ is a perturbation
scale, and $\Delta^{-1}$ denotes elementwise inversion.

\paragraph{Bayesian experimental design for protocol selection
(non-limiting).}
In some embodiments, protocol selection is performed by Bayesian
experimental design (Chaloner--Verdinelli 1995; Foster et al. 2019),
with the digital twin or fleet surrogate acting as the predictive
model and an expected-information-gain, expected-utility,
knowledge-gradient, Thompson-sampling, or related acquisition function
selecting the next candidate protocol. The acquisition function is
constrained by meter envelopes, optical-radiation safety bounds,
thermal limits, query budgets, and disclosure policies, and its prior,
posterior-update rule, uncertainty representation, and optimisation
budget are recorded alongside the run.

\paragraph{Minimax adversarial training objective (non-limiting).}
In some embodiments, adversarial hardening is written in minimax form:
\[
  \min_{\theta \in \Theta} \;\max_{A \in \mathcal{B}_{\text{surr}}}\;
  \E\Bigl[\mathcal{L}_{\text{hard}}\!\bigl(M_{\text{hard}},
  C, \tilde{C}_A\bigr)\Bigr],
\]
where $\mathcal{B}_{\text{surr}}$ is a family of surrogate adversaries,
$C$ are genuine committed windows, and $\tilde{C}_A$ are
adversary-generated candidates. Optimisation alternates: update $A$ to
increase spoof success, update meters to reduce false accepts, and
update permissible parameters to increase difficulty subject to meter
envelopes.

\paragraph{Projected-gradient surrogate adversaries (non-limiting).}
In some embodiments, the surrogate adversary family
$\mathcal{B}_{\mathrm{surr}}$ includes projected-gradient perturbation
adversaries of the kind used in robust-optimisation adversarial training
(Madry et al. 2018). The projected-gradient attack is applied to a
differentiable surrogate, digital twin, learned feature representation,
or differentiable component of the protocol, and its perturbation set,
step size, restart count, iteration count, and projection operator are
declared in the protocol digest. Candidate adversarial windows produced
by the surrogate are supplied as negative examples or hard cases to the
meter-update loop, while physical actuation remains bounded by the
apparatus safety and meter envelopes.

\paragraph{Curriculum and escalation (non-limiting).}
In some embodiments, adversarial training follows a curriculum: begin
with weak adversaries (replay and low-capacity generators), then
escalate to stronger adversaries (higher-capacity models,
physics-informed generators, partial side information, and model
inversion). Escalation triggers include meter-disagreement spikes,
anomalous drift indicators, or sudden decreases in verisimilitude on
held-out reference sweeps.

\paragraph{Hybrid digital--analogue co-training (non-limiting).}
In some embodiments, training updates both digital parameters and
analogue/physical degrees of freedom, potentially on different
timescales. Non-limiting analogue update mechanisms include: a
``training beam'' that writes correction patterns into memory media
derived from error signals, pump or gain modulation that adjusts cavity
or amplifier operating points as a function of meter objectives, and
feedback illumination schedules that drive photorefractive or
photochromic media toward configurations that improve declared meter
outcomes. These updates are bounded by safety and stability meters and
are logged for audit.

\paragraph{Physical actuation binding (non-limiting).}
In some embodiments, ``updating the kernel'' includes applying
actuator-safe set-point changes to one or more hardware controls,
including without limitation: $V_{G1}$, $V_{G2}$,
$V_{\mathrm{anode}}$, deflection and auxiliary coil currents, focus and
astigmatism controls, mask selection or SLM patterns, MEMS mirror
angles, liquid-lens or MEMS-lens voltages, spectral filter drives, and
per-emitter gains in addressable arrays. The controller enforces rate
limits and limit projections before committing updates and logs applied
set-points and calibration identifiers in the protocol digest.

\paragraph{Common random numbers and robust aggregation (non-limiting).}
In some embodiments, zeroth-order evaluations reuse the same validation
challenges, protocol seeds, and any internal randomness when evaluating
an objective at perturbed parameters $u+c\Delta$ and $u-c\Delta$. This
common-random-numbers discipline reduces estimator variance. Each
perturbation side may be evaluated using multiple replicates under
matched seeds, and replicate scores aggregated using a robust estimator
(for example median-of-means).

\paragraph{Run records and audit fields (non-limiting).}
In some embodiments, training, calibration, and hardware-in-the-loop
optimisation consume structured run records as the primary substrate.
Run records extend meter records by binding sweep/control identifiers,
challenge identifiers, provenance metadata, calibration state histories,
and digital twin metadata to the underlying committed windows and meter
outputs. Operation is treated as a two-timescale system, separating slow
maintenance (re-enrolment, deep recalibration, twin training, meter
updates, and sweep-library evolution) from fast runtime adaptation
(bounded adjustments gated by meters), with both timescales logged.

\subsection{Frozen-meter adversarial tuning (non-limiting)}

The auxiliary diagnostics in this subsection and in the
Scale-to-threshold subsection below (including
$H_{\mathrm{err}}(\kappa)$, $H_{\mathrm{steps}}(\kappa;\tau)$,
$C^\ast(\kappa)$, $c^\ast(\kappa)$, and fidelity-threshold $\tau$
framing) are operational training-loop tools, not instances of the
canonical two-dimensional hardness index pair
$(k^*_{\mathrm{dig}}(\theta;\varepsilon),
m^*_{\mathrm{ana}}(\theta;\varepsilon, r))$ defined in
Section~\ref{sec:security-theory}.  They serve as surrogate objectives
during meter and device training; the canonical hardness index
characterises the resulting device under the declared threat model and
evaluation meters.\looseness=-1

In some embodiments, an objective term in $J(u)$ is defined using a
neural meter trained in a two-phase procedure. In a first phase, a
conditional meter model is pretrained across a range of device/kernel
states and protocol classes. The pretrained meter is then frozen and
used as a fixed-capacity adversary or verifier. In a second phase, the
physical configuration and/or protocol family is updated to increase the
frozen meter's difficulty.

In discriminator-driven embodiments, a frozen discriminator induces a
separability score measuring how well it separates physical
\cba windows from generated or simulated windows under
matched challenges. In reconstruction-driven embodiments, a frozen
reconstructor induces a modelling-error hardness
$H_{\mathrm{err}}(\kappa)$ and an optional compute-hardness
$H_{\mathrm{steps}}(\kappa;\tau)$ defined as the minimum number of
bounded reconstruction steps required to reach a fidelity threshold
$\tau$:
\[
  H(\kappa) = H_{\mathrm{err}}(\kappa)
  + \gamma\,H_{\mathrm{steps}}(\kappa;\tau).
\]

In some embodiments, a bank of frozen meters is maintained with varied
architectures, seeds, datasets, and sampler configurations, and hardness
is computed as a minimum or low-quantile score across the bank to reduce
overfitting to a single meter.  In some embodiments, a minimum freeze
window duration $T_{\mathrm{freeze}}$ is declared in the protocol
digest and enforced: a meter is not updated or replaced during any
active verification or PoD confirmation window, so that an adversary
who queries the system during a frozen window cannot steer the meter
by slow drift before the window closes.  The freeze window start time
and duration are committed to the protocol digest before the window
opens.  A drift meter (a separate, lightweight monitor for meter
output stability) is declared alongside each frozen meter; if the
drift meter signals anomalous output shift during a freeze window, the
freeze is extended and the event is logged.  In some embodiments,
the hardness index is additionally bounded under an adaptive-query
model in which the adversary is permitted $Q$ queries per freeze
window: the declared query-rate limit $Q / T_{\mathrm{freeze}}$ and
the empirically estimated meter-learning rate are recorded in the
protocol digest so that the claimed hardness margin under this model
is traceable.

\paragraph{Operating-point consistency meter (non-limiting).}
The slow-drift attack exploits the fact that acceptance meters adapt
to gradual operating-point shifts over timescales longer than the
rotation period, potentially steering the device toward a region of
lower hardness without triggering any single-window alarm.  In some
embodiments, this attack is bounded by an \emph{operating-point
consistency meter} $m_{\mathrm{drift}}(\theta_t, \theta_{\mathrm{ref}})$
that tests whether the device's operating point $\theta_t$ has
remained within a declared tolerance $\varepsilon_{\mathrm{drift}}$
of its reference point $\theta_{\mathrm{ref}}$ during the evaluation
window.  The reference point is set at the start of each hardness
evaluation freeze window.  In some embodiments, the maximum tolerable
drift over a freeze window of length $T_{\mathrm{freeze}}$ is derived
from the AR(1) drift model (Section~\ref{sec:noise-model}):
$\varepsilon_{\mathrm{drift}} = z_\alpha \sigma_{\mathrm{drift}} \sqrt{T_{\mathrm{freeze}}}$,
where $\sigma_{\mathrm{drift}}$ is the estimated drift rate and
$z_\alpha$ is the declared confidence quantile.  If
$m_{\mathrm{drift}}$ exceeds $\varepsilon_{\mathrm{drift}}$, the
hardness evaluation is invalidated for that window, the event is
logged, and the evaluation is rescheduled after a fresh reference
point is established.  The operating-point consistency meter and its
threshold are committed to the protocol digest and are evaluated
concurrently with the frozen evaluation meters.  A declared operating
envelope in $\Theta$ space (separate from meter space) is maintained;
operating-point drift is distinguished from physical device degradation
by comparing against the fleet-wide distribution of drift rates.

\paragraph{Freeze window duration bound from AR(1) drift model
(non-limiting).}
Given the AR(1) drift model $\delta_{t+1} = \rho_{\mathrm{AR}} \delta_t + w_t$
(Section~\ref{sec:noise-model}), the expected squared drift over a
window of length $W$ starting from a reference point is:
\[
\mathbb{E}[|\delta_W - \delta_0|^2]
= \sigma_w^2 \frac{1 - \rho_{\mathrm{AR}}^{2W}}{1 - \rho_{\mathrm{AR}}^2},
\]
where $\sigma_w^2$ is the noise variance of the AR(1) process.  In
some embodiments, the bias induced on the hardness index estimate by
drift during an unfrozen window of length $W'$ is bounded as a
function of $\mathbb{E}[|\delta_{W'} - \delta_0|^2]$ times the
empirically estimated sensitivity $\partial k^*_{\mathrm{dig}} / \partial \delta$;
the maximum unfrozen window duration $W'_{\max}$ is then set so that
this drift-induced bias satisfies
$\mathrm{bias}(W'_{\max}) \leq \eta_{\mathrm{freeze}}$, where
$\eta_{\mathrm{freeze}}$ is a policy-declared fraction of the hardness
threshold.  In some embodiments, $W'_{\max}$ is committed to the
protocol digest alongside $\rho_{\mathrm{AR}}$, $\sigma_w^2$, and
$\eta_{\mathrm{freeze}}$, making the freeze window duration a derived
and auditable quantity rather than an ad hoc engineering choice.

\paragraph{Frozen GAN meter (non-limiting).}
In some embodiments, a frozen meter is realised by a conditional GAN
trained across kernel states. A frozen discriminator defines a
separability score $S_D(\kappa)$ for distinguishing physical windows
from generator windows. A frozen generator defines a divergence score
$S_G(\kappa)$ measuring perceptual discrepancy when attempting to
imitate physical responses. Non-limiting GAN objectives include
hinge-loss, Wasserstein, and feature-matching variants.

\paragraph{Frozen diffusion meter (non-limiting).}
In some embodiments, a frozen meter is realised by a diffusion
reconstructor. Non-limiting hardness proxies include modelling error and
a step-count hardness defined by the minimum number of
reverse-diffusion steps required to reach a fidelity threshold under a
fixed bounded-step sampler (for example DDIM with $\eta=0$).

\paragraph{Normalisation and keyed scoring (non-limiting).}
In some embodiments, meter inputs are standardised by a fixed
non-learned normalisation, and normalisation or feature extraction may
be keyed or salted so that disclosed scores are harder to replay or
overfit without access to the key; key identifiers are logged while the
key itself is not disclosed.

\subsection{Scale-to-threshold hardness (non-limiting)}

In some embodiments, hardness is quantified using a
scale-to-threshold protocol: an attacker or meter model class is grown
in capacity until it reaches a performance threshold $\tau$ on held-out
challenge windows. An optional shrink-test estimates the minimal
capacity $C^\ast(\kappa)$ that still meets $\tau$. The device is
updated to maximise $C^\ast(\kappa)$, thereby forcing attackers to
deploy larger models to achieve the same competence. The same protocol
family may be used in an assistive direction for cooperative decoders
by updating $\kappa$ to minimise an analogous $C^\ast(\kappa)$.

\paragraph{Classifier-centred curricula (non-limiting).}
In some embodiments, scale-to-threshold tuning is organised as a
curriculum over a family of decoders or classifiers ordered by capacity
$c$:
\[
  c^\ast(\kappa) = \min\{c:\; \mathrm{Perf}(C_c;\kappa)\ge \tau\}.
\]
An assistive phase updates $\kappa$ to reduce
$c^\ast_{\mathrm{Bob}}(\kappa)$, while an adversarial phase updates
$\kappa$ to increase $c^\ast_{\mathrm{Eve}}(\kappa)$. A single
objective may increase a gap such as
$c^\ast_{\mathrm{Eve}} - w_{\mathrm{adv}}\,c^\ast_{\mathrm{Bob}}$ subject to
meter envelopes and actuator projections.

\paragraph{Constant-time scoring and DoS guards (non-limiting).}
In some embodiments, evaluation is padded to a fixed wall-time so that
latency does not reveal the number of iterations required.
Denial-of-service guards cap per-score wall-time, maximum sampler steps,
and query budgets, and a rate limiter constrains the frequency of
challenge--response queries.

\paragraph{Randomised-smoothing certificates over meter features
(non-limiting).}
In some embodiments, randomised-smoothing certificates (Cohen, Rosenfeld,
and Kolter 2019) are computed over meter features, verifier scores, or
latent summaries under a declared perturbation distribution. The
certificate supplies an auxiliary certified-robustness radius or bound
alongside the empirical scale-to-threshold index $C^\ast(\kappa)$, and
is valid only for the declared feature map, perturbation distribution,
base-meter version, confidence level, and sampling budget. The smoothing
noise family, variance or scale parameter, number of Monte-Carlo samples,
and certified norm are recorded in the protocol digest.

\subsection{Continual calibration, drift compensation, and transfer
\label{sec:calibration}
learning (non-limiting)}

In some embodiments, calibration is continual: alignment, gain, drift,
and timing parameters are updated online using meter feedback and
committed windows. Calibration updates are gated by stability meters,
and calibration state identifiers and versions are logged. Drift
compensation may use replay buffers of previously validated windows with
regularisation against a reference parameter state. In some embodiments,
catastrophic-forgetting mitigation is implemented using methods from the
continual-learning literature, including without limitation elastic
weight consolidation (EWC, Kirkpatrick et al. 2017), synaptic
intelligence (SI), memory-aware synapses (MAS), online EWC,
learning-without-forgetting (LwF), gradient-episodic-memory methods
(GEM, A-GEM), experience replay, generative replay, dark experience
replay (DER), progressive networks, PackNet, and modular
continual-learning methods. Transfer learning
across reactors is supported via device-conditioned modulation, low-rank
adapters (including without limitation LoRA, QLoRA, DoRA, and LoHa),
prefix tuning, prompt tuning, adapter tuning, IA${}^3$, BitFit, and
other parameter-efficient fine-tuning (PEFT) methods,
or universal encoders with device-specific heads. In some embodiments,
per-deployment adaptation uses meta-learning methods, including
without limitation model-agnostic meta-learning (MAML), Reptile,
ANIL, first-order MAML (FOMAML), and meta-reinforcement-learning
methods such as RL${}^2$ and PEARL.

\paragraph{Probability calibration of meter outputs (non-limiting).}
In some embodiments, meter outputs that are interpreted as probabilities
or confidence scores are calibrated using Platt scaling, temperature
scaling (Guo et al. 2017), isotonic regression, beta calibration, or
histogram/binning methods. Calibration parameters are versioned with
the meter identifier, evaluated on held-out committed windows
stratified by device, protocol class, and operating envelope, and
logged with expected calibration error, maximum calibration error,
reliability-diagram summaries, and the calibration-data digest. A
meter whose calibration error exceeds a declared threshold is
downgraded, frozen pending review, or excluded from acceptance
decisions under the meter-partition policy.

\paragraph{Conformal and selective-prediction meters (non-limiting).}
In some embodiments, meter outputs are wrapped by conformal-prediction,
selective-classification, or abstention mechanisms that convert raw
scores into prediction sets, rejection regions, or audit-escalation
triggers under a declared calibration distribution. The conformal
score, calibration split, exchangeability or stratification assumptions,
coverage level, abstention threshold, and recalibration schedule are
recorded with the meter version. Where exchangeability is not plausible
because of drift, device heterogeneity, or protocol adaptation, the
conformal guarantee is limited to the declared stratum or replaced by
a weighted, covariate-shift, or online calibration variant, and the
limitation is recorded in the protocol digest.

\paragraph{Uncertainty-budget and traceability records (non-limiting).}
In some embodiments, calibration records include an uncertainty budget
that decomposes meter uncertainty into repeatability, detector noise,
alignment, wavelength or spectral calibration, timing jitter,
temperature drift, model discrepancy, reconstruction error, and
cross-device transfer components. The uncertainty budget may follow a
GUM-style propagation or Monte-Carlo propagation method, and may
include traceability metadata for reference artefacts, calibration
targets, timing references, optical power measurements, and
environmental sensors. The budget is recorded as audit metadata and
as a gating input for claims that require metrological confidence,
without requiring that every embodiment operate as a certified
metrology instrument.

\subsection{Energy-based loss functions (optional)}

In some embodiments, training losses are derived from a physically
grounded energy functional rather than coordinate-wise penalties. Around
each ground-truth configuration $S$ a Boltzmann-inspired energy landscape at
temperature parameter $\tau$ motivates the loss:
\[
  \mathcal{L}_E(\theta) = E(\hat{S}_\theta; S) - E(S; S).
\]
If $E$ is invariant under a symmetry group $G$, then $\mathcal{L}_E$
inherits this invariance. For discrete scenes, a tractable energy is
$E_{\mathrm{dist}}(S)=\sum_{i<j}V(\|s_i-s_j\|)$, which is
automatically invariant to translations, rotations, and permutations of
identical elements.

\paragraph{Koopman-style latent dynamics (non-limiting).}
In some embodiments, a latent representation is learned in which
closed-loop evolution is approximately linear:
$z_{t+1}\approx K z_t$, supporting prediction, control design, and
anomaly or tamper detection via deviations in the learned linear
evolution or changes in the spectrum of $K$.

\subsection{Distributional alignment and robustness (optional)}

\paragraph{Optimal-transport divergences (non-limiting).}
In some embodiments, discrepancies between response distributions are
quantified using optimal-transport losses, including Sinkhorn
divergences between batches of simulated and physical \cba
windows. Unbalanced optimal transport supports missing data, and
Gromov--Wasserstein variants compare relational structure across
modalities.

\paragraph{Certified robustness against bounded perturbations
(non-limiting).}
In some embodiments, a verifier is constructed to admit robustness
certificates against bounded perturbations of \cba features,
for example via randomised smoothing or Lipschitz-constrained
architectures.

\subsection{Self-adapting Reality Kernels (non-limiting)}

In some embodiments, a \RK is operated as a continuously
adapting physical process rather than a system with a strict separation
between training and inference. A learned self-edit head proposes
adaptation actions:
\[
  \pi_{\mathrm{edit}} : c_{0:t} \mapsto a^{\mathrm{edit}}_t,
\]
where $c_{0:t}$ denotes the emission--observation history and $a^{\mathrm{edit}}_t$
denotes a proposed self-edit. Non-limiting self-edits include proposed
updates $\Delta\theta_t$ to slow parameters within bounded ranges,
schedules for biasing native analogue relaxation, and bounded write
programmes in embodiments that include a fabrication loop. A wrapper map
applies edits subject to hard safety and calibration constraints:
\[
  \theta_{t+1},\, M_{t+1}
  = F_{\mathrm{edit}}\bigl(\theta_t, M_t, a^{\mathrm{edit}}_t\bigr),
\]
where $M_t$ denotes slow media state.

\paragraph{Meta-reward lens (non-limiting).}
In some embodiments, edits are evaluated over a horizon $H$ using a
meta-level return:
\[
  R^{\mathrm{meta}}(a^{\mathrm{edit}}_t)
  = \sum_{\tau=t}^{t+H}
  \Bigl(
    -\ell_{\mathrm{task}}(\tau)
    + \lambda_{\mathrm{hard}}\,H_{\mathrm{dig}}(\theta_\tau)
    - \lambda_{\mathrm{cal}}\,C_{\mathrm{drift}}(\theta_\tau)
  \Bigr),
\]
where $\ell_{\mathrm{task}}$ is a task loss,
$H_{\mathrm{dig}}$ is an empirical hardness score, and
$C_{\mathrm{drift}}$ is a calibration-drift penalty. The self-edit
policy is trained by reinforcement learning or other optimisation to
increase $\E[R^{\mathrm{meta}}]$ while remaining within the admissible
edit envelope.

\paragraph{Overnight self-editing pattern (non-limiting).}
In some embodiments, a practical experimental schedule alternates
between supervised ``daytime'' use and constrained ``overnight''
self-edit phases, supporting empirical characterisation of stability,
overfitting, drift, and benefits of self-adaptation under different
meta-reward choices and constraint policies.

\subsection{Federated and multi-agent learning (non-limiting)}

In some embodiments, learning is federated across a network of devices,
with updates weighted by trust and audit outcomes. Devices may accept
aggregated updates only when accompanying evidence meets verification
requirements, and when meter envelopes indicate stable operation.

\paragraph{Federated averaging, secure aggregation, and privacy
accounting (non-limiting).}
In some embodiments, the federated update loop is implemented by
Federated Averaging (McMahan et al. 2017) over committed-window-derived
gradients or parameter deltas, with per-device contributions protected
by secure aggregation (Bonawitz et al. 2017) and with the aggregated
update accepted only when the disclosed aggregate and accompanying
evidence satisfy meter-envelope compliance. In some embodiments,
fleet-update releases are governed by differential-privacy accounting,
including a moments-accountant or DP-SGD analysis (Abadi et al. 2016),
R\'enyi differential privacy (Mironov 2017), or Gaussian differential
privacy (Dong--Roth--Su 2022), with clipping norms, sampling rates,
noise scales, composition rule, privacy budget, and unit of protection
recorded in the protocol digest. In some embodiments, FLAME-style
backdoor filtering, Krum or Multi-Krum, coordinate-wise median,
trimmed-mean, Bulyan, clustering-based anomaly filters, or other
Byzantine-robust aggregation rules are composed with the meter-gated
acceptance step, and the chosen aggregation rule, rejection criteria,
and false-rejection tolerance are committed before the round opens.

\paragraph{Privacy accounting and selective-disclosure consistency
(non-limiting).}
In some embodiments, the relationship between a declared
differential-privacy budget and the selective-disclosure policy is
itself committed in the protocol digest: the selective-disclosure
policy shall not permit openings, witness sets, or audit
reconstructions that are inconsistent with the declared DP guarantee
or that would, in aggregate, exceed the declared composition budget.
In some embodiments, the auditor or relying party is provided with
the DP-accounting state alongside any selectively disclosed evidence,
so that the privacy-budget posture is itself part of the evidence
record.

\paragraph{Fleet-update poisoning diagnostics (non-limiting).}
In some embodiments, fleet updates are accompanied by poisoning and
backdoor diagnostics computed on held-out committed windows, neutral
probe corpora, canary challenges, or synthetic-but-committed
adversarial cases. Diagnostics include update-norm clipping compliance,
cosine- or Mahalanobis-distance outlier scores, influence estimates,
gradient-sign agreement, cluster membership, backdoor trigger success
on held-out probes, and degradation of calibration or verisimilitude
meters. The fleet accepts, rejects, downweights, or quarantines the
update according to a predeclared rule, and all rejected or
quarantined updates remain available for audit under the
selective-disclosure policy.

In some embodiments, the system trains verisimilitude or consistency
discriminators using committed windows. A non-limiting training loop
alternates: sample committed windows, label them as consistent or
inconsistent (for example by controlled perturbations, mismatch
injections, or cross-device disagreement), update discriminator
parameters to reduce classification error, and periodically validate
against fresh physical runs under meter-bounded regimes.

In some embodiments, learned scan policies are trained under
constraints: policies propose candidate protocols, meters enforce
admissibility, and selective opening is used to audit that learning
did not exploit unlogged channels or post hoc fabrication. In some
embodiments, training is framed as a multi-agent game between agents
proposing protocols or interpretations and verifiers or discriminators
enforcing consistency with \RK evidence, with trust weights
and hardness scores shaping the effective payoffs and update weights.

\section{Representative Reactor Embodiments (Non-limiting)}
\label{sec:embodiments}
% ======================================================================

This section describes non-limiting physical embodiments that implement
the \RK operator
$\mathsf{P}_\theta(\cdot\mid S,U_{0:T})$. Unless stated
otherwise, each embodiment produces a \cb $C_{0:T}$ as defined
in Section~\ref{sec:definitions}, where each sample bundles the control
input $u(t)$ and an observation vector $\mathbf{y}_t$ that may include
multiple detector ports. The bench-top 1D \RK of
Section~\ref{sec:anchor-1d} provides a concrete teaching
anchor. The 2D generalisation is described in
Section~\ref{sec:anchor-2d}. The subsections below describe
additional reactor substrates, memory extensions, multi-rate
embodiments, and alternative physical realisations.

\subsection{Embodiment overview}

In various embodiments, the \RK module is implemented using
one or more of: (i)~a scene-facing loop (structured emission toward an
external scene and detection of a response), (ii)~a reactor loop (an
engineered subsystem contributing mixing, memory, nonlinearity, or
history dependence), and (iii)~optional media layers (slow or persistent
memory media) that modify the effective channel over time.

In one practical organisation, some embodiments aim for minimal electronics
and treat the dominant computation as optical or physical, with
electronics used primarily for timing, gain, digitisation, and control
signalling. Other embodiments introduce stronger nonlinearities, longer
memory, or higher bandwidth, and may employ additional sensors or
actuators. In a non-limiting multi-rate configuration, a fast
all-optical loop implements a convolution-like kernel in a feedback
cavity (for example a gain-pumped cavity with AOD/EOM providing fast
sweep and coupling control), while one or more slower analogue-memory
loops---for example CRT/phosphor stacks, laser-addressed luminescent
screens, or other persistence media---provide additional state and
mixing on longer timescales. In such multi-rate stacks, the fast loop
supplies high-throughput convolution-like mixing and scattering, and the
slow loop stores or reshapes states between passes, all under the same
\RK control and logging abstraction.

\paragraph{Form factors and deployment scenarios (non-limiting).}
In some embodiments, \RKs are realised as handheld scanners
or probes, desktop or bench-top units, kiosk or public verification
stations, integrated mobile devices (for example smartphones and
tablets using existing cameras and screens or flashes as emitters),
wearables (for example body cameras and smart glasses), fixed
infrastructure (for example security cameras and sensor beacons), and
airborne or satellite platforms for verified remote sensing. In
preferred embodiments, each form factor records protocol digests and
meter envelopes so that exported evidence is interpreted under the
correct operating regime and sensor geometry.

\paragraph{Carrier and spectral variants (non-limiting).}
In some embodiments, the carrier includes one or more of: UV (for
example fluorescence excitation, surface inspection, and sterilisation
verification), visible, NIR, short-wave infrared (SWIR, for example
through-fog and moisture-sensitive imaging), mid-infrared (MIR, for
example molecular spectroscopy and gas sensing), far-infrared (FIR) and
terahertz (THz, for example non-destructive evaluation and screening),
RF and microwave (for example channel probing and through-wall sensing),
and acoustic or ultrasound. In preferred embodiments, the selected
bands, filters, and calibration identifiers are recorded in protocol
digests so that verifiers interpret the \cb under the correct
bandpass and operating envelope.

\paragraph{Embodiment-family conformance summary (non-limiting).}
The following non-limiting table maps representative embodiment families
to supported primitives, typical regimes, and key constraints. This
table is illustrative and non-exhaustive; additional families appear
throughout the specification and combinations are contemplated.
\emph{Legend:} Regimes: TB\,=\,\TB, L\,=\,\LI, RT\,=\,\RT.
Primitives: R\,=\,Read, W\,=\,Write, Sign\,=\,Sign, Erase\,=\,Erase.

\begingroup
\footnotesize
\setlength{\tabcolsep}{3pt}
\renewcommand{\arraystretch}{1.05}
\begin{longtable}{@{}>{\raggedright\arraybackslash}p{0.18\textwidth}>{\raggedright\arraybackslash}p{0.16\textwidth}>{\raggedright\arraybackslash}p{0.11\textwidth}>{\raggedright\arraybackslash}p{0.17\textwidth}>{\raggedright\arraybackslash}p{0.22\textwidth}@{}}
\toprule
Family & Primitives & Regimes & Typical meters & Notes \\
\midrule
\endfirsthead
\toprule
Family & Primitives & Regimes & Typical meters & Notes \\
\midrule
\endhead
\endfoot
\bottomrule
\endlastfoot
Hybrid proj.--det. &
  R/W/Sign/Erase &
  TB, L, RT &
  Drift, timing, TE &
  Full-featured reference \\
CRT/phosphor &
  R/W/Sign/Erase &
  TB, L, RT &
  Phosphor decay, persistence &
  Analogue memory; Erase via natural decay \\
Fibre-delay &
  R/Sign &
  TB, L &
  Speckle, polarisation &
  Compact; RT limited \\
Acoustic/cymatic &
  R/W/Sign &
  TB, L, RT &
  Mode spectrum, scattering &
  Non-optical; SASER variant \\
PIC/waveguide &
  R/Sign &
  TB, L &
  Ring resonance, loss &
  On-chip; wafer-scale \\
Metasurface &
  R/W/Sign &
  TB, L, RT &
  Spectral, polarisation &
  Flat optics; programmable \\
Chromic/electro. &
  R/W &
  L, RT &
  Colour change, switching &
  Slow; persistent state \\
Sweeping reactor &
  R/W/Sign &
  TB, L &
  Modulation freq., phase &
  Swept parameter (T, V, path) \\
GNSS passive &
  R/Sign &
  TB, L &
  Pseudorange, carrier phase &
  Passive; no emitter control \\
Extended-scale &
  R/W/Sign &
  TB, L, RT &
  Domain-specific &
  Geological, atmospheric, bio \\
\end{longtable}
\endgroup
\noindent Notes in the table are typical and non-limiting; they do not
exclude alternative configurations, faster/slower variants, or
additional primitives beyond those listed.

% ------------------------------------------------------------------
\subsection{Generalisation: shared-scan reactor-modulated scene kernel (2D)}
\label{sec:anchor-2d}
\label{sec:anchor-generalisation}
% ------------------------------------------------------------------

The 2D generalisation of the 1D anchor of
Section~\ref{sec:anchor-1d} extends the bench-top construction
along three axes: the single-axis micromotor-mirror scanners
are replaced by galvanometer pairs providing 2D scan coverage,
the per-sample discrete feedback loop is replaced by a
continuous real-time amplitude drive $A(t)$, and the
fluorescence-specific reactor substrate is replaced by a
substrate-open reactor stack with a generic nonlinear medium.
The 2D anchor embodiment (FIG.~6) comprises two physically
separate subsystems driven from a shared clock: a scene
subsystem~210 open to the world under a declared eye-safety
envelope, and a reactor subsystem~220 sealed behind a declared
safety envelope.

In the parallel-subsystem embodiment~(a) of
Section~\ref{sec:architectural-alternatives},
the architectural coupling between the scene subsystem~210 and
the reactor subsystem~220 is exclusively the real-time
amplitude drive $A(t)$ produced from reactor detector~44
photocurrent and delivered to scene emitter~24. The shared scan
program produced by common controller~10 and fanned out to
galvanometer pairs~25 and~26 is a synchronisation mechanism,
not an architectural coupling within the meaning of this
section: the shared scan program causes the two subsystems to
traverse their respective scan trajectories in phase, and is
structurally distinct from the architectural feedback path
$A(t)$. In some embodiments, $A(t)$ is produced by a declared
live conditioning chain comprising protective conditioning (for
example transimpedance amplification and a safety limiter)
followed by a frozen committed reactor detector gain
$G_{\mathrm{det}}$, which is trained offline as one of the
trainable parameter families for the 2D anchor (see
Section~\ref{sec:anchor-2d-trainable}) and held fixed for the
committed deployment version. In the parallel-subsystem
embodiment~(a), the live coupling path contains no per-period
wait, no forward low-latency coupling from scene detector~43
to reactor source, and no runtime-adaptive or runtime-learned
element; the scene emitter's instantaneous amplitude at time
$t$ is therefore the reactor subsystem's response at time $t$
as shaped by the declared live conditioning chain, up to a
bounded and declared settling time. The cycle is closed in
embodiment~(a) by the $A(t)$ feedback path alone.

In the cascade embodiment~(b) of
Section~\ref{sec:architectural-alternatives},
the architectural coupling between the scene subsystem~210 and
the reactor subsystem~220 additionally includes a forward
low-latency coupling from scene detector~43 to the reactor
source of subsystem~220, illustrated in FIG.~6 as a
heavy-stroke arrow. The forward low-latency coupling is
realised in any of the non-limiting implementations recited in
Section~\ref{sec:architectural-alternatives} (optical
gain-pumped; optical direct relay; electrical; mixed-signal),
each with its end-to-end latency declared and committed to the
protocol digest. In the cascade embodiment~(b), the cycle is
closed by both the forward low-latency coupling and the
$A(t)$ feedback path operating jointly, with the scene
observation entering the reactor stage in real time via the
forward coupling and the reactor response being returned to
the scene-emitter amplitude in real time via $A(t)$.

In some embodiments, the 2D anchor scene subsystem and reactor
subsystem are described more particularly in
Sections~\ref{sec:anchor-2d-scene} and
\ref{sec:anchor-2d-reactor}. Unless a declared extension is
engaged and committed to the protocol digest, the reactor
subsystem and scene subsystem of embodiment~(a) are coupled
exclusively by the live reactor-detector-to-scene-emitter
amplitude drive $A(t)$ (with the shared scan program acting
as a synchronisation mechanism rather than as an architectural
coupling), and the reactor subsystem and scene subsystem of
embodiment~(b) are coupled by both the forward low-latency
coupling and $A(t)$, in each case rather than by bidirectional
cross-conditioning beyond what is declared.

In some embodiments, the shared scan program is a Lissajous
trajectory in $(x,y)$ with declared frequencies, phase
offsets, and amplitudes, committed to the protocol digest. In
some embodiments, the same scan program drives both
galvanometer pairs one-to-one, so that the reactor subsystem
and the scene subsystem traverse the same trajectory in their
respective coordinate systems at the same phase of the shared
scan period $T$. In some embodiments, the scan trajectory
shape is one of six trainable anchor parameter families (see
Section~\ref{sec:anchor-2d-trainable}).

In some embodiments, spatial or temporal warping of the
reactor scan relative to the scene scan --- translation,
rescaling, or phase offset --- is a declared extension of the
one-to-one shared-scan default and is committed to the
protocol digest when used. In some embodiments, a
runtime-adaptive transfer function in the live coupling path
beyond the declared live conditioning chain, scene-side
amplitude modulation in addition to reactor-driven modulation,
and bidirectional cross-conditioning between the subsystems
are further declared extensions, each committed to the
protocol digest when used.

\paragraph{Safety envelope.}
\label{sec:safety-envelope}
In some embodiments, the scene-facing path is operated under
declared eye-safety, thermal, scan-rate, and stability bounds,
while higher-power or otherwise restricted reactor illumination
is confined to the sealed reactor subsystem. In some
embodiments, protocol digests and meter summaries record
whether illumination occurred in a sealed internal path or in
a scene-facing path, together with the applicable interlock,
safety-envelope, and meter-envelope status.

In some embodiments, a handheld or portable probe embodiment is provided
(PolieProboscis), in which a compact projector--detector pair and
optional sealed reactor loop are integrated into a single enclosure for
close-range scanning, verification, or environmental sampling.
PolieProboscis operation is constrained to an eye-safe scene-facing
regime, while any higher-power operation is confined to sealed internal
paths under interlocks.

\subsubsection{Scene subsystem (2D)}
\label{sec:anchor-2d-scene}

In the 2D anchor, the scene subsystem~210 comprises an
illumination source~24 selected to meet a declared eye-safety
envelope (for example a lower-power infrared or visible
source), a collimating lens, a galvanometer pair~25 driven by
the shared scan program, a scan lens defining a field of view,
and a scene detector~43. The scene emitter's instantaneous
amplitude is driven by $A(t)$ from the reactor subsystem~220.
The controller drives the galvanometer pair~25 according to
the shared scan program (for example a Lissajous trajectory)
and records scene responses synchronised to the scan timing.
In some embodiments, the scene detector~43 is a single
detector, for example a photodiode, as the anchor default; a
camera, a photodiode array, or other spatially resolved
detector is a declared extension of the single-detector
default. In some embodiments, the scene detector is coaxial
with the scene emitter via a polarisation combiner or
equivalent.
\subsubsection{Reactor subsystem (2D)}
\label{sec:anchor-2d-reactor}

In the 2D anchor, the reactor subsystem~220 is sealed behind
a declared safety envelope and comprises, on the excitation
path and in optical order: a reactor illumination source, a
galvanometer pair~26 driven by the shared scan program, a
tunable emitter-side focus lens, a scattering element, and a
nonlinear medium~30. On the detection path, the reactor
subsystem further comprises a tunable detector-side focus
lens, a detector-side aperture, and a single reactor
photodetector~44, for example a photodiode. Because the
reactor subsystem is sealed, the reactor illumination source
is selected to drive strong nonlinearity in the reacting
medium and is free from the constraints governing the scene
illumination; for example, the reactor illumination source
may be a UV source, a high-intensity laser, or any other
source compatible with the sealed safety envelope and the
medium's excitation requirements. In some embodiments, the
reactor illumination source is driven by a trained profile
$A_{\mathrm{emit}}(t)$ periodic in the shared scan period
$T$ (see Section~\ref{sec:anchor-2d-trainable});
$A_{\mathrm{emit}}(t)$ probes the nonlinear medium's
intensity-dependent response across different phases of the
scan period.

\paragraph{Reactor-side collimation (non-limiting).}
In some embodiments, the reactor illumination source is a laser
or laser-diode module producing a beam of sufficient collimation
for galvanometer scanning without a separate collimating lens
between the reactor source and galvanometer pair~26. In other
embodiments, a collimating element is positioned between the
reactor source and galvanometer pair~26; its presence, optical
parameters, and calibration state are declared and committed to
the protocol digest. The asymmetry between the scene loop's
explicit collimating lens (Section~\ref{sec:anchor-2d-scene})
and the reactor loop's optional or absent collimating lens shown
in FIG.~6 reflects the different illumination-source choices
available in the two loops under the declared safety envelope
and is not a feature limiting the apparatus.

In some embodiments, the tunable emitter-side focus lens (for
example a liquid lens) follows a trained profile
$f_{\mathrm{emit}}(t)$ periodic in the shared scan period
$T$. In some embodiments, the tunable detector-side focus
lens follows a trained profile $f_{\mathrm{det}}(t)$
periodic in $T$ and controls the region and depth of the
reactor response over which the reactor photodetector
integrates at each scan phase. In some embodiments, the
detector-side aperture follows a trained profile
$a_{\mathrm{det}}(t)$ periodic in $T$, implemented for
example as a fast-actuated iris, an LCD-based aperture, a
MEMS-based aperture, or other programmable spatial stop; the
aperture provides spatial selectivity of the reactor response
analogous to a programmable confocal pinhole in the
detection path. In some embodiments using a mechanical iris
or other slower actuator not capable of per-scan-period
modulation, $a_{\mathrm{det}}$ is a committed scalar (for
example a committed fixed aperture diameter) rather than a
time-varying profile, and is committed to the protocol
digest as such.

In some embodiments, the detector-side aperture is positioned
before the detector-side focus lens, after the detector-side
focus lens, or at an intermediate image plane, provided that
the declared configuration determines the spatial selectivity
of the reactor response and is committed to the protocol
digest.

The detector-side focus lens and detector-side aperture shape
the reactor response before conversion to photocurrent; they
are not part of the post-photodetector live coupling path
from the reactor detector to the scene emitter. The
post-photodetector live coupling path begins at the reactor
photodetector output and proceeds through protective
conditioning and frozen committed $G_{\mathrm{det}}$, as
described in Section~\ref{sec:anchor-2d-coupling}.

In some embodiments, the scattering element is any strong
random optical mixer, for example ground glass, a diffuser,
or a heterogeneous scattering bed. In some embodiments, the
nonlinear medium~30, additional reactor-medium elements, or
combined reactor-medium layer may include, without
limitation: (i)~a scattering plate such as ground glass or a
diffuser; (ii)~a thin-film nonlinear element, for example
graphene-coated or other intensity-dependent media; (iii)~a
temporal or persistence layer such as phosphor or other
luminescent media; (iv)~a delay or mixing element such as
fibre-delay loops; (v)~an epsilon-near-zero (ENZ) thin film
such as indium tin oxide (ITO), whose dielectric permittivity
can be switched on femtosecond timescales via the optical
Kerr effect, enabling ultrafast temporal modulation of the
reactor's optical properties; (vi)~a gas discharge or plasma
medium (for example a sealed gas cell, a plasma display panel
cell, a dielectric barrier discharge layer, or an arc
discharge electrode pair) whose ionisation dynamics,
filament formation, and emission spectra provide strongly
nonlinear, device-specific, and temporally rich reactor
behaviour; (vii)~a microwave or radiofrequency cavity with
at least one tunable boundary condition (for example a
piezoelectrically actuated endwall or an electrically tuned
metamaterial surface), whose electromagnetic mode spectrum
provides a calculable but device-specific transfer function
modulated by the boundary geometry; and/or (viii)~any
combination of the foregoing. The phosphor option is listed
as one non-limiting reactor-medium element among others and
is not a privileged default for the 2D anchor.

In some embodiments, the reactor photodetector~44 is a
single detector as the anchor default; a photodiode array or
other spatially resolved detector is a declared extension.
In some embodiments, the 2D anchor is defined by the
topology of the reactor stack and the training protocol, not
by the specific choice of nonlinear medium, scattering
element, or reactor illumination source; the
fluorescence-specific substrate of the 1D anchor
(Section~\ref{sec:anchor-reactor-subsystem}) is one
non-limiting substrate that may be used in the 2D topology.

In the 2D anchor, the reactor emitter-side focus profile
$f_{\mathrm{emit}}(t)$ is implemented by a tunable liquid
lens driven from a scalar kernel coordinate
$\kappa_{f,e}(t)$ synchronised to the shared scan period.
Similarly, the reactor detector-side focus profile
$f_{\mathrm{det}}(t)$ is implemented by a tunable
detector-side focus lens driven from a scalar kernel
coordinate $\kappa_{f,d}(t)$ synchronised to the shared
scan period. Other controls are held fixed or updated on
slower timescales. The trained values defining
$f_{\mathrm{emit}}(t)$, $f_{\mathrm{det}}(t)$, and
$a_{\mathrm{det}}$ (whether a profile or a committed scalar)
over the scan period are committed to the protocol digest.
\subsubsection{Trainable parameter families (2D)}
\label{sec:anchor-2d-trainable}

In the 2D anchor, six trainable parameter families are
optimised jointly against the hardness objective of
Section~\ref{sec:anchor-hardness} using the training
pathways of Section~\ref{sec:anchor-training}:
\begin{enumerate}[nosep]
  \item The shared scan trajectory shape --- Lissajous
    parameters (frequency ratio, phase offsets, amplitudes),
    or more generally the shared scan path traversed by
    both galvanometer pairs. Periodic in the shared scan
    period $T$.
  \item The reactor emitter-side focus lens profile
    $f_{\mathrm{emit}}(t)$. Periodic in $T$.
  \item The reactor detector-side focus lens profile
    $f_{\mathrm{det}}(t)$. Periodic in $T$. Controls the
    region and depth of the reactor response over which the
    reactor photodetector integrates at each scan phase, and
    may be trained independently from $f_{\mathrm{emit}}(t)$.
  \item The reactor detector-side aperture profile
    $a_{\mathrm{det}}(t)$. Periodic in $T$ where the
    aperture implementation supports fast actuation;
    alternatively a committed scalar where a fixed
    mechanical aperture is used. Provides spatial
    selectivity of the reactor response analogous to a
    programmable confocal pinhole. A fixed-aperture
    embodiment is the scalar degenerate case of the
    detector-side aperture parameter family; the scalar may
    be selected or trained during configuration and is then
    frozen in the committed parameter-family version,
    preserving the six-family count.
  \item The reactor emitter drive profile
    $A_{\mathrm{emit}}(t)$, applied upstream of the reactor
    stack. Periodic in $T$. $A_{\mathrm{emit}}(t)$ probes
    the nonlinear medium's intensity-dependent response at
    each scan phase and is distinct from the reactor
    detector gain $G_{\mathrm{det}}$ because, for a
    nonlinear reactor, pre-stack scaling and post-stack
    scaling are not equivalent.
  \item The reactor detector gain $G_{\mathrm{det}}$,
    applied in the live coupling path between the reactor
    photodetector and the scene emitter drive.
    $G_{\mathrm{det}}$ may be a time profile
    $G_{\mathrm{det}}(t)$ periodic in $T$, a nonlinear
    transfer function $G_{\mathrm{det}}(I)$, or a combined
    $G_{\mathrm{det}}(t,I)$, allowed to be nonlinear per
    Section~\ref{sec:anchor-trainable}. Where
    $G_{\mathrm{det}}$ is time-dependent, its phase origin
    is defined relative to the shared scan period $T$ and
    committed to the protocol digest.
\end{enumerate}

The 2D anchor also has six trainable parameter families, but
their composition differs from the 1D anchor. The 1D anchor
has separate scene and reactor scan trajectories; the 2D
anchor has one shared scan trajectory and adds the
detector-side aperture profile $a_{\mathrm{det}}(t)$. Thus
the 2D default comprises one shared scan, one reactor
emitter-side focus profile, one reactor detector-side focus
profile, one detector-side aperture profile, one reactor
drive profile, and one frozen committed coupling-path gain.
Scan-warping between subsystems, when used, is a declared
extension (Section~\ref{sec:anchor-2d-extensions}) rather
than part of the default trainable set. In some embodiments,
the trained values defining all six parameter families are
committed to the protocol digest as part of the committed
parameter-family version.

\subsubsection{Common controller, live coupling path, and time alignment (2D)}
\label{sec:anchor-2d-coupling}

In the 2D anchor, the common controller~10 generates a single
shared scan program and fans it out to the galvanometer pairs
in both subsystems. The \emph{live coupling path}, as used in
this subsection, refers specifically to the $A(t)$ return path
from reactor detector~44 to scene emitter~24 drive. It is a
deterministic low-latency path comprising the reactor
photodetector, protective conditioning (for example
transimpedance amplification and a safety limiter), and the
reactor detector gain $G_{\mathrm{det}}$ of
Section~\ref{sec:anchor-2d-trainable}. $G_{\mathrm{det}}$ is
trained offline and frozen for the committed deployment
version. At runtime, the live coupling path contains no
per-period wait and no runtime-adaptive or runtime-learned
element. In the parallel-subsystem embodiment~(a) of
Section~\ref{sec:architectural-alternatives},
the live coupling path is the sole architectural coupling
between the scene and reactor sides, and there is no forward
low-latency coupling from scene detector~43 into the reactor
stage. In the cascade embodiment~(b) of
Section~\ref{sec:architectural-alternatives},
the live coupling path is supplemented by a separately disclosed
\emph{forward low-latency coupling} from scene detector~43 to
the reactor source of subsystem~220 (see
Section~\ref{sec:architectural-alternatives}); the forward
low-latency coupling is itself declared and committed to the
protocol digest and is not runtime-adaptive. The live coupling
path and the forward low-latency coupling are structurally
distinct: the live coupling path is the reactor-detector to
scene-emitter return signal, and the forward low-latency
coupling, where present, is the scene-detector to
reactor-emitter forward signal.

In some embodiments, $G_{\mathrm{det}}$ is implemented as a
fixed analog shaping network whose transfer function matches
the trained profile. In other embodiments, $G_{\mathrm{det}}$
is implemented as a deterministic high-rate sampled transfer
function or look-up table operating on reactor photocurrent,
provided that the sampled implementation introduces no
per-scan-period buffering and has a declared bounded latency.
In either implementation, the latency from reactor
photocurrent to scene emitter modulation is bounded by a
declared settling time and is committed to the protocol
digest. The scene-detector and reactor-detector measurement
records are separate from the live coupling path and follow
the timestamping conventions below.

In some embodiments, the common controller drives the scene
subsystem and the reactor subsystem from a shared scan
program and timebase. Scene and reactor measurements are
timestamped and optionally buffered to compensate for
exposure, readout, propagation, or processing delays. The
controller pools the aligned measurements into the
observation vector $\mathbf{y}_t$ and logs the \cb
$c_t=(u(t),\mathbf{y}_t)$. In embodiments using derived
latents, the controller updates $Z_\theta(t)$ from $c_t$ and
optionally uses $Z_\theta(t)$ to determine subsequent
controls.


\subsubsection{Light-field hybrid embodiment (non-limiting)}
\label{sec:lightfield-hybrid}

In some embodiments, the projector-role emission path and
detector-role sensing path of the hybrid embodiment are each realised
as four-dimensional light-field subsystems: a light-field projector
(Section~\ref{sec:emitter-embodiments-beyond}) emits a controllable
field $E(x, y, u, v)$ into the scene branch, and a plenoptic detector
(Section~\ref{sec:detector-embodiments-beyond}) records the returning
field $L(x', y', u', v')$ in a single exposure. The reactor branch may
independently be two-dimensional or four-dimensional. The \cba sample
$c_t = (u(t), \mathbf{y}_t)$ therefore pairs a four-dimensional
commanded emission with a four-dimensional observed response, and the
per-step transfer kernel maps between emitted and observed light-field
coordinate spaces rather than between two-dimensional intensity fields.

In some embodiments, the projection-side microlens array and the
detection-side microlens array are optically combined through a beam
splitter so that the emit and capture paths share an MLA plane; in
other embodiments, the two arrays are independent, with registration
established by a calibration sub-protocol logged in the protocol
digest.

Where a scene branch includes a camera-obscura screen, relay screen,
display surface, or other two-dimensional intermediate surface, the
angular-consistency primitive applies to that intermediate physical
surface unless the relay optics preserve calibrated angular
information from the external scene. In embodiments claiming external
scene parallax through such a relay, the relay geometry, aperture,
field stop, and angular-preservation calibration are recorded in the
protocol digest and evaluated by declared meters.

\paragraph{Angular-consistency verification primitive (non-limiting).}
In some embodiments, a \TB sub-regime uses the plenoptic
capture as an angular-consistency check: scene emissions that are
consistent with a three-dimensional physical scene, as observed over
the declared sub-aperture baseline, exhibit measured parallax
relationships across sub-apertures $(u, v)$. Planar presentation
spoofs---including printed images, ordinary flat displays lacking
calibrated angular multiplexing, planar projection surfaces, and
planar retro-reflective replay surfaces---lack the depth-dependent
parallax relationships expected under the declared pose, illumination,
and calibration state, and therefore produce an angular-consistency
residual under the declared meter. In some embodiments, the
angular-consistency residual is included as a component of the
verisimilitude meter and committed alongside the \cb. This geometric
exclusion is limited to planar or effectively planar replay and
print-spoof classes; high-fidelity angular replay spoofs, including
lenticular, holographic, light-field, angularly multiplexed planar
displays, volumetric displays, three-dimensional printed replicas,
textured mannequins, and parallax-mimicking surface assemblies, are
not excluded by this geometric argument alone and are handled, where
relevant, by the existing meter envelope, declared attacker model,
and empirical hardness evaluation.

\paragraph{Four-dimensional reciprocity and adjoint-matching
(non-limiting).}
In some embodiments, a calibration and verification sub-protocol
emits a declared four-dimensional light field $E(x, y, u, v)$ and
verifies that the captured field $L(x', y', u', v')$ is consistent
with the adjoint, radiometric-adjoint, or declared transport-adjoint
model of the forward reactor or scene transfer kernel under the
declared configuration $\theta$. In some embodiments, the
reciprocity residual is logged in the meter envelope and treated
analogously to the phase-conjugation residual of interferometric
regimes, while remaining compatible with incoherent operation where
phase-stable interferometric geometry is not available.

\paragraph{Effective-rank headroom (non-limiting).}
In some embodiments, replacing a two-dimensional emitted or observed
intensity field with a calibrated light-field channel increases the
raw addressable channel description from spatial coordinates to
spatial--angular coordinates. The product of the declared spatial and
angular resolutions on the emit and capture sides is treated only as
an upper bound on the hypothetical physical rank available under the
four-dimensional channel geometry, not as a task-useful rank claim.
Declared Tier~2 effective rank is measured under the declared meter
set from the committed physical evidence, for example by correlation,
singular-spectrum, or independence analyses of the recorded channel.
Declared Tier~3 task-useful rank is measured separately for the
relevant task and deployment, and is not inferred from the
spatial--angular resolution product. The protocol digest records the
effective $(x, y, u, v)$ resolution, calibration state, vignetting,
cross-talk, and sub-aperture registration on the emit and capture
sides so that any downstream rank statement is conditioned on the
committed channel geometry and the measured evidence for that
deployment.

\paragraph{Compatibility with existing operating regimes
(non-limiting).}
In some embodiments, the light-field hybrid embodiment remains
compatible with all three operating regimes (\TB, \LI,
\RT) and with Yoked operation: in \LI,
four-dimensional capture provides passive depth and refocusing
evidence in a single exposure; in \RT, the
four-dimensional emitted field supports near-eye and integral-imaging
rendering, view-dependent projection mapping, and parallax-consistent
immersive rendering; in Yoked operation, the angular response of the
scene to controlled angular emission constitutes a new coupling
dimension whose stability is gated by the same conditional Lyapunov
exponent discipline described elsewhere in this specification.

\subsubsection{Sealed reactor pumps and eye-safe external beams}

In many embodiments, high-intensity or tightly focused beams are
confined to internal optical paths within an enclosed reactor housing,
while the scene-facing emission is operated in an eye-safe or otherwise
bounded exposure regime appropriate for the deployment context. One or
more internal pump sources may drive nonlinear or history-dependent
reactor behaviour inside the enclosure. From the channel-model
perspective, internal pumps modify the effective configuration $\theta$
and thereby modify the induced conditional law
$\mathsf{P}_\theta(\cdot\mid S,U_{0:T})$.

\subsubsection{High-intensity etching, engraving, and ionising variants}

In some variants, one branch of the module is configured for controlled
interaction with matter at higher intensities, for example to engrave,
etch, locally ablate, expose latent layers, or create persistent
fiducials on a target substrate. In such variants, safety and hardware
interlocks constrain the permissible working distance, angle, duty
cycle, and total exposure. The interaction modifies the scene state over
time and therefore changes the realised \cb and, where
present, media state $M_t$.

\subsubsection{Temporal modulation and time-frequency reactor operation
(non-limiting)}
\label{sec:temporal-modulation}

In some embodiments, the reactor medium is modulated in the time
domain, not only the spatial domain.  Recent experimental work has
demonstrated that thin-film epsilon-near-zero (ENZ) materials---such
as indium tin oxide (ITO), which is widely used as a transparent
conductor in display panels, touchscreens, and LCD screens---can
switch optical state on timescales approaching an optical cycle
(approximately 10~fs) via the optical Kerr effect, producing temporal
apertures whose interaction with probe light yields frequency-domain
interference patterns (Tirole et al., ``Double-slit time diffraction
at optical frequencies,'' \emph{Nature Physics}~19, 999--1002, 2023;
Galiffi et al., \emph{Adv.\ Photonics}~4, 014002, 2022).  This result is relevant to \RK embodiments in
three respects.

First, temporal modulation extends the reactor's degrees of freedom:
the kernel $\mathsf{P}_\theta$ operates on the spatio-temporal field
jointly, not merely the spatial distribution of light.  The emitter
can encode structured information in pulse timing as well as spatial
pattern, and the recorder can measure spectral interference patterns
that encode both scene response and reactor temporal dynamics.

Second, time-varying media can enable nonreciprocal responses under
appropriate spatiotemporal modulation conditions: light
propagating in one direction through a temporally modulated reactor
experiences a different transformation than light in the opposite
direction.  In some embodiments, the reactor is operated in a
temporal modulation regime specifically to exploit this non-reciprocal
property as a hardness mechanism for the verification regime.  The
emitter-to-recorder optical transfer function then differs physically
from the recorder-to-emitter transfer function, providing a
directional attestation primitive that does not depend on
computational complexity assumptions (see also
Section~\ref{sec:security-theory} for the general reciprocity-breaking
hardness discussion).

Third, in some embodiments the reactor's temporal modulation is
periodic with period $T$, producing discrete frequency sidebands
(Floquet modes) separated by $1/T$.  In such time-crystal-like
regimes, the kernel operates on a discrete spectral lattice, and the
periodic modulation provides a natural clock cycle for optical
processing.  Topologically protected frequency-conversion pathways
arising from the Floquet band structure may, in some embodiments,
provide signatures that are robust against noise but empirically hard
to forge without the physical medium.  Time-refraction effects arising
from single-cycle refractive-index modulation have been experimentally
demonstrated in ENZ media (Lustig et al., ``Time-refraction optics
with single cycle modulation,'' \emph{Nanophotonics}~12(12),
2221--2230, 2023), confirming that sub-cycle temporal interfaces
produce measurable frequency shifts suitable for \cba
characterisation.

ITO is highlighted as a non-limiting example because it is already
present in the screen-based reactor embodiments described elsewhere
in this specification.  Other ENZ or time-varying metamaterial
substrates may be used.  All temporal modulation parameters (pulse
timing, modulation depth, repetition rate, pump profile) are logged
in the \cb alongside spatial controls.  A complementary
approach---the sweeping reactor (Section~\ref{sec:sweeping-reactor})---modulates reactor \emph{geometry} rather than material
properties, at much lower frequencies (Hz--kHz rather than THz), and
uses phase-sensitive (lock-in) demodulation to extract the derivative
of the reactor's physical transform with respect to the swept
parameter.  The two approaches are orthogonal: temporal modulation
of material properties amplifies hardness and can enable nonreciprocal responses,
while parametric geometry modulation provides precision differential
characterisation of the reactor's transfer function.

\subsubsection{Gas discharge and plasma reactor embodiments
(non-limiting)}

Gas discharge and plasma media constitute a broad class of nonlinear
reactor media whose dynamics are inherently chaotic,
device-specific, and temporally rich.  Plasma exhibits negative
differential resistance, avalanche ionisation, recombination
cascades, and self-organised filamentary pattern formation---all
strongly nonlinear.  Manufacturing variation (electrode geometry, gas
pressure and composition, electrode surface condition, tube or
envelope contamination) produces device-specific
challenge--response behaviour suitable for reactor-microstructure hardness
assessment.  Non-limiting plasma reactor embodiments include the
following.

\paragraph{Plasma display panel (PDP) as micro-reactor array
(non-limiting).}
In some embodiments, a plasma display panel is used as a combined
emitter and reactor array.  A PDP comprises a grid of sealed gas
discharge cells, each producing ultraviolet emission that excites
phosphor coatings to yield visible light.  Each cell is an
independent micro-reactor with device-specific gas pressure,
electrode gap, and phosphor grain distribution.  The panel
simultaneously functions as a structured emitter (it displays
arbitrary patterns) and as a reactor array (each cell's discharge
dynamics are nonlinear and manufacturing-dependent).  In some
embodiments, the PDP is driven near the cells' firing threshold,
where avalanche ionisation is most sensitive to perturbation and the
system operates near a bifurcation surface; Fisher information about
cell identity peaks in this regime.  A camera observing the PDP
output while the controller sweeps driving waveforms across the
panel produces \cbs encoding the joint
emitter--reactor dynamics of thousands of micro-reactors.

\paragraph{Near-threshold arc discharge (non-limiting).}
In some embodiments, a gas discharge source (xenon arc, carbon arc,
or other arc lamp) is operated near its breakdown voltage, where the
discharge dynamics are maximally sensitive to electrode geometry, gas
density, and surface micro-features.  Near the ionisation threshold,
Townsend avalanche dynamics exponentially amplify small
perturbations into macroscopically distinct discharge paths,
producing the highest information yield per probe and the hardest
conditions for an attacker to reproduce.  In carbon-arc embodiments,
the electrodes are consumed during operation---carbon evaporates, the
crater shape evolves, and carbon nanoparticles participate in the
discharge---producing a self-modifying reactor whose
challenge--response mapping drifts irreversibly because the physical
medium is being destroyed and reformed.  This built-in irreversibility
provides a naturally time-indexed hardness source: the reactor's state
at any moment reflects its operational history and is not expected
to be practically rewound without replacing or materially altering
the reactor under the declared maintenance, access, and observation
assumptions.  Xenon arc sources are commonly used light sources in
cinema and high-end DLP projectors; in some embodiments, the arc's
intrinsic dynamics are observed alongside the projected content as
additional \cba channels.

\paragraph{Laser-induced plasma as transient reactor (non-limiting).}
In some embodiments, a focused laser pulse ablates a surface,
producing a transient plasma plume whose emission spectrum encodes
the elemental composition of the ablated material (the physical
basis of laser-induced breakdown spectroscopy, LIBS).  In some
embodiments, a second structured illumination probe is directed
through or onto the expanding plume, so that the plume serves as a
transient three-dimensional reactor medium: it scatters, absorbs, and
re-emits the probe light according to its spatially and temporally
varying composition and density.  A spectrometer and/or camera
records both the plume's self-emission and the transmitted or
scattered probe light.  This produces a \cb in which
the reactor is created by the emission itself, exists for
microseconds, carries information about the scene (the ablated
surface), and is practically unreproducible because each ablation
event removes material.  In some embodiments, scanning the ablation
point across the surface produces a spatially resolved
composition--dynamics map logged as a committed \cba sequence.

\paragraph{Dielectric barrier discharge and atmospheric plasma
(non-limiting).}
In some embodiments, a dielectric barrier discharge (DBD) panel or
atmospheric-pressure plasma source produces self-organised
filamentary patterns whose spatial structure depends on gas
composition, driving frequency, dielectric surface properties, and
electrode geometry.  These patterns are chaotic, history-dependent,
and responsive to proximity of external objects (a nearby grounded
conductor restructures the discharge pattern, analogous to
scene-coupled dynamics in Yoked operation).  In some embodiments, a
plasma ball or similar sealed discharge device is used as a reactor
whose visible filament dynamics are observed by a camera; the
filaments' response to nearby objects or touches provides a
scene-coupled observation channel.

\paragraph{Fluorescent tube as cascaded frequency-conversion reactor
(non-limiting).}
In some embodiments, a fluorescent tube is used as a cascaded reactor:
electrical excitation produces gas discharge plasma, the plasma emits
ultraviolet radiation, and phosphor coatings downconvert UV to visible
light.  The three-stage nonlinear cascade (electrical $\to$ plasma
$\to$ UV $\to$ visible) amplifies device-specific manufacturing
variation at each stage.  Gas pressure, electrode condition, and
phosphor grain distribution all contribute to the device-specific
spectral output.  At mains frequency (50/60~Hz magnetic ballast) or
tens of kilohertz (electronic ballast), the discharge is periodically
modulated, producing temporal dynamics analogous to the Floquet-regime
operation described in the temporal-modulation embodiments above.
% ------------------------------------------------------------------

In one non-limiting bench-top embodiment, the scan law and a portion of
the optical kernel are implemented by a single moving element in a
removable shroud. A compact source (for example a laser diode module)
illuminates a sealed container holding a heterogeneous scattering medium
(for example mixed glass or reflective elements). A micro-servo-mounted focus-and-steering optic
provides both focus control and one-dimensional
steering by tilting about an axis. A photodetector views the container
interior through spectral filtering to emphasise selected bands (for
example fluorescence bands or scattered light bands). The servo angle
defines a scan law over reactor volume, and the focus setting of the
optic defines a kernel coordinate. Shaking or otherwise reconfiguring the
internal medium between runs re-randomises microstructure and thereby
changes the response distribution while preserving a family resemblance
under similar protocols.

\paragraph{Sealed-housing fluorescent-medium variant (non-limiting, by cross-reference).}
In one non-limiting variant, the sealed container is a metal housing
holding a shallow bed of mixed marbles or fragments. Suitable
fluorescent and reflective media, including those discussed in detail
at the bench-top 1D anchor (Section~\ref{sec:anchor-1d}), may be used
in the present 2D embodiment subject to the same safety, handling,
and regulatory considerations described there. Geometry-specific
aspects of the present 2D embodiment include the sealed metal
housing; the servo-defined steering schedule across a two-dimensional
bed; photodetector observation through a green-pass or long-pass
optical filter so that the measured response emphasises
wavelength-shifted fluorescence together with scattered light from the
non-fluorescent media; and a shake-randomisation procedure between
runs that yields different but statistically related photodetector
time series under a fixed sweeper program.

% ------------------------------------------------------------------

\subsubsection{Hardness, training, and extensions (by reference)}
\label{sec:anchor-2d-extensions}

The hardness objective
(Section~\ref{sec:anchor-hardness}), the staged training
pathways (Section~\ref{sec:anchor-training}), the Yoked
orthogonal extension (Section~\ref{sec:anchor-yoked}), the
noise co-illumination extension
(Section~\ref{sec:anchor-noise}), and the CRT/phosphor
analogue-memory extension
(Section~\ref{sec:crt-phosphor-extension}) apply to the 2D
anchor by direct reference, mutatis mutandis: references in
those sections to the 1D per-sample discrete feedback loop
are read in the 2D case as references to the continuous live
reactor-detector-to-scene-emitter amplitude broadcast of
Section~\ref{sec:anchor-2d}, and the 2D trainable parameter
families of Section~\ref{sec:anchor-2d-trainable} substitute
for the 1D families. In the 2D anchor, the CRT extension's
default coupling between reactor detector output and CRT
drive is not constrained by the 1D anchor's green-pass
wavelength-separation argument, because the 2D anchor's
reactor substrate is substrate-open; both electrical and
optical CRT coupling are available as declared configurations.

\subsection{Layered reactor stack embodiments}
\label{sec:layered-stacks}
% ------------------------------------------------------------------

In some embodiments, the reactor loop is implemented as a layered stack
designed to amplify memory, mixing, nonlinearity, and empirical
hardness. A representative non-limiting configuration may include, in
sequence: (i)~a scattering plate, (ii)~a nonlinear thin-film layer,
(iii)~a temporal persistence layer, (iv)~a delay-network layer (for
example fibre-delay loops with selectable couplings), and (v)~one or
more readout heads operating on the \cb and/or derived latents.

Individual layers may be characterised and tuned in isolation and then
assembled. Slow physical parameters may be fine-tuned using
hardware-in-the-loop optimisation such as SPSA, while digital meters
and decoders are trained on logged traces. Spectral multiplexing may be
used to separate measurement bands (for example a band reserved for
verification probing and a separate band for presentation or internal
pumping), with per-band gain and filtering to maintain detector dynamic
range.

% ------------------------------------------------------------------
\subsection{SLM--scattering-medium--camera embodiment (non-limiting)}
\label{sec:slm-scattering-camera}
% ------------------------------------------------------------------

In some embodiments, the reactor is implemented as a spatial light
modulator (SLM) directing structured illumination through a scattering
or nonlinear medium, with a camera or detector array capturing the
output. The SLM provides programmable, protocol-controlled emission
patterns; the scattering medium provides a high-dimensional,
physically fixed, device-specific transfer function; and the camera
captures the \cb. In such embodiments, the SLM pattern
is included in the control protocol $U_{0:T}$ and recorded in the
protocol digest, while the medium's transfer function contributes the
analogue complexity and history dependence that generate empirical
hardness and computational richness.

This architecture supports all three operating regimes on the same
hardware. In \TB mode, the device-specific speckle or
scattering signature provides physically unclonable challenge-response
behaviour. In \LI mode, the SLM patterns are optimised to
maximise information about a scene placed adjacent to or within the
scattering medium. In \RT mode, the SLM patterns are
optimised to control the output field subject to the medium's
physical constraints. Yoked operation is possible when the medium
includes dynamic elements (for example a liquid crystal layer,
flowing suspension, or biological specimen whose optical properties
evolve in response to the illumination).

In some embodiments, the scattering medium is sealed and
characterised at enrolment, providing a stable enrolled microstructure signature. In other
embodiments, the medium is exposed to environmental or biological
inputs, and the scene-dependent variation in the scattering transfer
function becomes the measurement signal. Meters record speckle
contrast, mode-mixing statistics, and temporal stability so that
verification and sensing claims are conditioned on the observed
regime.

% ------------------------------------------------------------------
\subsection{CRT/phosphor analogue-memory extension embodiment}
\label{sec:crt-phosphor-extension}
\label{sec:crt-phosphor}
% ------------------------------------------------------------------

In some embodiments, the anchor embodiment of
Section~\ref{sec:anchor-1d} or Section~\ref{sec:anchor-2d} is
extended by adding a CRT/phosphor analogue-memory loop atop
the anchor's live coupling. The CRT/phosphor loop provides a
slow, spatially indexed analogue-memory channel, read out by
a camera, and does not replace the anchor's live coupling. In
this extension, the reactor photodetector signal (or a
derived signal) is written to the CRT via XY deflection
locked to a declared scan coordinate: in the 1D anchor, to
$\mathrm{scan}_{\mathrm{r}}(t)$,
$\mathrm{scan}_{\mathrm{s}}(t)$, or a declared derived
coordinate; and in the 2D anchor, to the shared scan program.
The phosphor persistence provides a spatial accumulation of
reactor response over multiple scan periods that is
subsequently read by the camera. In the 1D anchor, the
default coupling between reactor detector and CRT drive is
electrical, because the green-pass filter architecture
enforces wavelength separation. In the 2D anchor, both
electrical and optical CRT coupling are available; the
declared choice is committed to the protocol digest.

\paragraph{Mapping to FIG.~4 and to embodiments (a) and (b).}
FIG.~4 illustrates the CRT/phosphor extension as a fast
optical loop (laser source~21, AOD/EOM deflector~22, relay
optics~23, reactor medium~30, photodetector~41) cascaded with
a slow analogue-memory loop (CRT reactor~31 with electron
gun, yoke, and phosphor screen~32, observed by camera~42),
both driven from common controller~10. The heavy-stroke arrow
shown in FIG.~4 from photodetector~41 to the slow-loop
reactor stage is the \emph{forward low-latency coupling} of
embodiment~(b) of
Section~\ref{sec:architectural-alternatives}: it carries the
fast-loop scene-side observation directly into the
analogue-memory stage as the load-bearing forward path of the
cascade embodiment. The forward low-latency coupling in FIG.~4
may be implemented in any of the modalities recited in
Section~\ref{sec:architectural-alternatives} (electrical
direct, mixed-signal, optical gain-pumped, or optical direct
relay), and its end-to-end latency is declared and committed
to the protocol digest. The parallel-subsystem embodiment~(a)
of Section~\ref{sec:architectural-alternatives} is obtained
from FIG.~4 by omitting the heavy-stroke forward arrow and
relying on the controller's scan protocol to synchronise
fast-loop and slow-loop interrogations, with the slow-loop
reactor stage driven by the controller alone rather than by
the fast-loop photodetector signal.

\paragraph{Hobbyist and entry-level embodiments (non-limiting).}
For hobbyist, educational, or cost-sensitive builds, a cathode-ray tube
provides a non-limiting analogue-memory reactor option. Aged and imperfect CRTs provide a
rich, high-dimensional optical memory: over their lifetime, phosphors,
shadow masks or aperture grilles, glass envelopes, and deflection
assemblies accumulate idiosyncratic burn-in patterns, stains, charge
traps, and hysteretic behaviours that together define a unique physical
fingerprint. Such tubes, including surplus or second-hand displays, are
therefore attractive as low-cost, high-entropy reactor elements even
before any deliberate modification.

\paragraph{Augmenting CRT reactors with additional memory layers
(non-limiting).}
In some embodiments, CRT reactors are enhanced with additional optical
memory layers that preserve the unique character of each tube while
deepening its state space. A thin luminous or photoreactive coating may
be applied over the phosphor or faceplate so that prior illumination
leaves slowly decaying traces that modulate subsequent excitation.
Dispersions of scatterers, dyes, or textured films may be deposited in
front of the screen to add further spatial structure, mixing, and
history dependence.

\paragraph{Sources of reactor-microstructure uniqueness in CRT reactors
(non-limiting).}
In some embodiments, reactor-microstructure uniqueness arises from manufacturing
tolerances and wear mechanisms including: gun alignment tolerances,
phosphor grain structure, shadow-mask or aperture-grille tolerances,
magnetic-material idiosyncrasies in yokes and shims, and glass-envelope
geometry variations. Together these yield device-specific response
diversity in which small control changes induce effectively uncorrelated
output changes under meter-bounded operation.

In a preferred configuration, the CRT is an XY-deflection tube (for
example an oscilloscope-style tube) providing separate analogue inputs
for horizontal deflection $X_{\mathrm{CRT}}(t)$, vertical deflection
$Y_{\mathrm{CRT}}(t)$, and intensity drive $Z_{\mathrm{CRT}}(t)$. In
some embodiments, these signals are produced by a controller and are
synchronised with reactor optics (for example galvanometer angles) so
that the CRT face coordinates and the reactor scan coordinates share a
stable mapping. The intensity drive may be modulated by processed scene
signals or by internal controller state to encode additional history
into phosphor persistence.

In some embodiments, an optical readout path relays a region of the CRT
face into the reactor loop. Relay optics image the CRT face (or a
portion of it) into the entrance pupil of reactor optics, after which
additional tunable elements and scan mirrors transform the field before
it reaches a reactor detector. The reactor detector may be a photodiode
with a transimpedance amplifier producing a continuous analogue
measurement, or a camera producing sampled frames whose timing is
logged.

In some embodiments, the CRT subsystem is used as a bidirectional
surface: it can be written by the electron-beam excitation pattern and
read optically, and it may be combined with further photoreactive or
scattering layers to increase state-space richness. The CRT/phosphor
stack is one example of a slower luminescent analogue-memory medium in a
multi-rate reactor stack that may be coupled to faster optical reactor
loops, alongside alternatives such as laser-addressed luminescent
screens or other luminescent persistence media.

\subsubsection{Trainable electron-optical controls and calibration
(non-limiting)}

In some embodiments, the CRT includes an electron gun and
electron-optical controls, and the controller treats electron-optical
set-points as tunable or trainable kernel parameters. Non-limiting
trainable parameters include: control-grid bias $V_{G1}(t)$,
screen-grid bias $V_{G2}(t)$, anode (EHT) voltage
$V_{\mathrm{anode}}$, deflection coil currents and waveform shaping
(including harmonic corrections and lookup tables), focus and
astigmatism controls (optionally address-dependent), auxiliary coil
currents producing controlled nonuniform fields, per-address dwell and
blanking schedules, and per-channel or per-gun spectral-tap gains (for
example $(w_R,w_G,w_B)$ in RGB-gun embodiments). In some embodiments,
additional tunable degrees of freedom include mechanically positioned
ferromagnetic shims or auxiliary magnets near the yoke to introduce
controlled geometric warps, and stencil or mask elements (including
turrets or programmable masks) to implement piecewise spatial apertures
or mask bases.

In some embodiments, an auxiliary diagnostic camera observes the CRT
face (or beam) to estimate spot size, convergence, and alignment,
enabling closed-loop calibration while the main observation stream
remains a single-pixel photodiode/TIA stream and/or a separate
scene-facing camera stream.

\paragraph{Nonlinearity and mixing mechanisms (non-limiting).}
In some embodiments, nonlinearity and mixing arise from controlled
geometric warps (shims, auxiliary magnets, auxiliary coils), spatial
apertures and mask selection (turrets or programmable masks implementing
a mask basis), and bias shaping via $V_{G1}(t)$ schedules that yield
nonlinear transfer curves and saturable regions. Spectral diversity
across multiple guns or sources yields multiple natural ``taps'' with
slightly different spot statistics, and these mechanisms interact with
phosphor persistence so that nonlinearity is coupled to analogue memory.

\paragraph{Time-token injection (non-limiting).}
In some embodiments, a time-token $\tau_{\mathrm{tok}}(t)$ is injected
into the physical channel, for example as a small-signal modulation onto
$V_{G1}(t)$ (or another bias) with a damping coefficient $d_\tau$ that
bounds its effect on gain and stability, and/or as an auxiliary optical
beam summed into the scene or reactor path. The token couples a timing
or authenticity signal into the physical response without requiring a
separate out-of-band channel, and token schedules, envelope limits, and
calibration identifiers form part of $\theta$ and are logged in the
protocol digest.

\paragraph{Beam PSF and persistence convolution (non-limiting).}
In some embodiments, instantaneous luminance at a phosphor plane is
modelled as
\[
  L(x,y,t)\propto Z_{\mathrm{CRT}}(t)\,
  K_{\mathrm{spot}}\bigl(x-X_{\mathrm{CRT}}(t),\;y-Y_{\mathrm{CRT}}(t)\bigr),
\]
where $K_{\mathrm{spot}}(\cdot,\cdot)$ is an effective spot point-spread function.
Emitted intensity incorporates persistence as a temporal convolution
\[
  E(x,y,t)=\int_{-\infty}^t L(x,y,\tau)\,p(t-\tau)\,d\tau,
\]
and the imaged observation is obtained by relay optics and detector
integration, with calibration and receptive-field models recorded in the
protocol digest.

\paragraph{Phosphor impulse response and detector weighting
(non-limiting).}
In some embodiments, phosphor persistence is modelled by an impulse
response
\[
  p(t)=\sum_i a_i e^{-t/\tau_i},\qquad t\ge 0,
\]
and dwell and blanking schedules shape an effective temporal kernel
$p_{\mathrm{eff}}(t)=p(t)\ast w_{\mathrm{det}}(t)$, where
$w_{\mathrm{det}}(t)$ denotes detector temporal weighting due to sensor
integration, transimpedance amplifier dynamics, filtering, and
digitisation aperture. In some embodiments, $w_{\mathrm{det}}(t)$ is
estimated by applying a known sub-threshold optical impulse (for
example a calibration LED pulse or a low-amplitude electrically
driven probe) and fitting a causal parametric IIR, FIR, or spline
model; the estimated model is stored as a calibration table and
its version is logged.

\paragraph{Receptive-field and spot-PSF calibration (non-limiting).}
In some embodiments, the detector spatial receptive field
$\rho(u,v)$ is estimated before or during operation by
spot-scanning the scene with a small bright probe, by
micro-pattern projection, or by projecting a structured calibration
grid, and fitting a parametric, spline, or lookup-table model to
the resulting response. The estimated receptive field is used to
correct signed, differential, bipolar, and pooled measurements,
and is committed to the protocol digest together with the
estimation procedure and its calibration epoch. In some
embodiments, the beam spot point-spread function
$K_{\mathrm{spot}}(\cdot,\cdot)$ at the phosphor plane is
estimated by knife-edge measurement, slanted-edge modulation-
transfer-function methods, micro-grid calibration, or related
techniques, and the estimated PSF parameters are stored in lookup
tables used by focus, convergence, astigmatism, and drift-
compensation control loops. Calibration epochs are versioned and
their identifiers are recorded in the protocol digest.

\paragraph{Face-camera diagnostic channel (non-limiting).}
In some embodiments, a separate diagnostic imager is directed at
the CRT phosphor face (or other persistence-medium surface) to
provide closed-loop measurement of spot size, beam position,
convergence, drift, and uniformity, independently of the primary
sensing channel that observes the relayed or scene-projected
signal. The face-camera channel is read on a dedicated calibration
clock, its samples and exposure schedule are committed to the
protocol digest as a separate atom from the primary observation,
and its outputs feed the trainable electron-optical and timing
controls. In some embodiments, the face camera operates at a lower
sample rate than the primary detector and supplies slow-loop
calibration corrections without participating in the fast scoring
or reconstruction loop.

\paragraph{Safety and burn-in prevention (non-limiting).}
In some embodiments, blanking interlocks disable beam emission on
deflection failure to reduce burn-in risk. EHT discharge procedures are
supported by bleeder resistors and explicit discharge workflows, and
shielding reduces ancillary emissions. Thermal and ageing compensation
uses slow bias-trim loops and periodic recalibration to keep spot
geometry and gain within a corridor; drift tracking may use PID- or
Kalman-style updates of lookup tables between training rounds.

\paragraph{Sources of CRT and electron-beam unclonability
  signature (non-limiting).}
In some embodiments, the per-device unclonability signature that
contributes to the reactor-microstructure uniqueness property
(Section~\ref{sec:security-theory}) arises from one or more of
the following sources, individually or in combination:
electron-gun alignment and assembly tolerances, including
cathode-to-grid spacing, accelerating-electrode alignment, and
gun-to-screen registration; phosphor microstructure, including
grain-size distribution, packing density, layer-thickness
variation, and dopant micro-distribution; shadow-mask, aperture-
grille, or slot-mask tolerances, including hole or slot pitch
variation, edge profile, and registration to the phosphor pattern;
deflection-yoke and shim magnetic idiosyncrasies, including
winding tolerances, residual magnetisation, and shim placement;
focus-coil and convergence-coil winding tolerances; auxiliary
electric-deflection-plate spacing and surface finish variations;
glass-envelope geometry, including thickness variation, residual
strain, and curvature; vacuum-level micro-variation across batches;
and component ageing histories that diverge from device to device
under different operating regimes. Each such source is individually
sufficient to contribute to per-device challenge-response variation,
and in combination they provide the high-dimensional uniqueness
signature on which proof-of-projection, proof-of-discrepancy, and
related verification protocols rely.

\paragraph{Programmable-mask placement and stencil control
  (non-limiting).}
In some embodiments, a programmable mask is positioned at one or
more of: proximate the CRT phosphor screen on the viewing side as a
programmable relay stencil; proximate the electron-gun aperture or
anode pathway as an electron-compatible beam-stencil reference; in
the relay optical path between the phosphor and the detector as an
aperture or filter element; or in the scene-projection path as a
structured-illumination element. At the phosphor-screen, relay-path,
and scene-projection placements, the programmable mask comprises
without limitation a liquid-crystal display panel, liquid-crystal-
on-silicon panel, digital micromirror device, or other spatial
light modulator. At the gun-aperture or anode-path placement, the
electron-compatible beam-stencil reference is implemented without
limitation by a configured electrostatic field (for example through
patterned biased electrodes or microfabricated electron-optical
apertures), by a configured magnetic field (for example through
patterned coil or shim arrays), or by a patterned physical aperture
mask compatible with the in-vacuum electron-optical environment.
Each programmable-mask configuration provides rapid stencil updates
under declared control, contributing trainable or adaptively-
controlled degrees of freedom that are committed to the protocol
digest together with the mask state and its update schedule.

\paragraph{Enumerated focusing and beam-shaping hardware
  (non-limiting).}
In some embodiments, focus, astigmatism, convergence, and beam-shape
adjustment is implemented by one or more of: micromotor-driven lens
positioning, MEMS-actuated lens elements, electrowetting or
liquid-crystal tunable lenses, piezoelectric lens-positioning
actuators, voice-coil actuators, varifocal liquid lenses, and
deformable-mirror or membrane-mirror elements on the optical side;
and on the electron-optical side by trainable focus-coil current,
trainable astigmatism-correction multipole currents, trainable
convergence-coil currents, and trainable accelerating-voltage
modulation. Each focusing or beam-shaping degree of freedom may
be operated as a fixed calibrated setting, a slow drift-compensated
setting, a trainable kernel parameter, or an adaptively-controlled
runtime parameter, with the operating mode and parameter trajectory
committed to the protocol digest.

\subsubsection{Signed-weight realisation via differential exposure
(non-limiting)}

In some embodiments, signed or bipolar weights are realised without
sacrificing CRT dynamic range by driving two complementary exposures
$E^+$ and $E^-$ and performing differential detection:
\[
  y(t)
  =
  \int \Bigl[E^+(u,v,t)-E^-(u,v,t)\Bigr]\;\rho(u,v)\,du\,dv,
\]
where $\rho(u,v)$ denotes a detector sensitivity (or pooling) function.
In some embodiments, $E^+$ and $E^-$ correspond to alternating masks,
alternating bias offsets, alternating write patterns, or alternating
line/field passes, and the differential signal is formed by subtracting
consecutive detector samples or by differencing two detector channels.

\paragraph{Temporal and channel differencing embodiments
(non-limiting).}
In some embodiments, signed weights are realised by temporal
differencing (positive and negative contributions assigned to distinct
time slots) or by channel differencing (positive and negative
contributions assigned to distinct channels such as wavelengths,
polarisations, or optical paths). A signed measurement is formed by
subtracting calibrated readouts under logged timing or channel tags.

\paragraph{Persistence guard (non-limiting).}
In some embodiments, a minimum inter-polarity delay (a persistence
guard) is enforced between $E^+$ and $E^-$ exposures to bound crosstalk
induced by phosphor decay. The guard interval is derived from fitted
phosphor time constants and is logged as part of the scan timing model
and protocol digest.

\paragraph{Multi-CRT summation (non-limiting).}
In some embodiments, two or more CRTs are optically combined, with
$E^+$ assigned to one tube and $E^-$ to another, realising signed
weights directly in the optical domain. Small misalignments between
tubes may be treated as additional mixing degrees of freedom rather than
solely as error sources, and are accounted for in calibration and meter
summaries. In some embodiments, per-tube misregistration patterns define
a device-specific convolution kernel that contributes to reactor-microstructure
uniqueness of the combined assembly.

\paragraph{Acoustic and RF analogues (non-limiting).}
In some embodiments, differential or signed-weight readouts are realised
in other physical domains by pairing a steered emitter with a
persistence element and differential detection. Non-limiting examples
include acoustic paths with decay (reverberant cavities or damped delay
lines) and RF cavities whose $Q$-factor produces multi-timescale memory.

\subsubsection{CRT as a two-dimensional analogue memory}

In some embodiments, the reactor includes a CRT or similar scanned
display device whose phosphor integrates and decays in response to
excitation, thereby implementing a two-dimensional analogue memory. The
deflection waveforms and beam intensity control are treated as
components of the control input $u(t)$, and the phosphor luminance
field is treated as part of the reactor or media state.

\subsubsection{Scan law parameterisation for CRT addressing}

Let the phosphor define a bounded domain $\Omega\subset\R^2$. A scan
law may be expressed as a time-parametrised curve
\[
  \gamma(t;\eta) = (x(t;\eta),\;y(t;\eta))\in\Omega,
\]
where $\eta$ parameterises frequencies, phases, amplitudes, envelopes,
or basis coefficients. Non-limiting examples include Lissajous
trajectories, chirps, piecewise dwell-and-jump schedules, and finite
basis expansions. The parameters $\eta$ may be included in $\theta$ and
may be held fixed, updated on a slow timescale, or adapted online.

\paragraph{Address-space parameterisation (non-limiting).}
In some embodiments, addressing is expressed as a composition of maps
\[
  (\rho,\varphi) = \Gamma_{\mathrm{addr}}(t;\eta),
  \qquad
  (x,y) = \Psi(\rho,\varphi),
\]
where $(\rho,\varphi)$ denotes a non-limiting internal address coordinate,
and $\Psi$ maps addresses to phosphor coordinates. In some embodiments,
$\Gamma_{\mathrm{addr}}$ and $\Psi$ are chosen so that the induced address map covers the
active phosphor region to within hardware resolution and noise, and is
invertible or approximately invertible for writing, reading, and
verification under the logged scan law and timing model.

\subsubsection{XY-deflection synchronisation with external scanners}

In one embodiment, an XY-deflection CRT provides analogue inputs for
horizontal and vertical deflection, and the same underlying drive
signals are also used, after calibration scaling, to drive galvanometer
scanners in another reactor path. This synchronises CRT coordinates and
optical scan coordinates so that CRT patterns can act as a spatially
indexed memory layer aligned to a scan address space.

\subsubsection{Channelised memory using line scans and rolling-shutter
readout}

In one embodiment, a one-dimensional scanner produces a time-varying
intensity profile that is written into one or more CRT lines. A
rolling-shutter camera reads the CRT surface with timing chosen relative
to the line-write schedule, producing a time-staggered sampling of the
phosphor memory. The resulting camera frames and/or derived features are
included as channels of $\mathbf{y}_t$. Colour CRTs may be used to
store multi-channel features by writing distinct values into different
phosphors.

\paragraph{Channelised memory with commodity hardware (non-limiting).}
In one illustrative embodiment, a commodity barcode scanner is coupled
to a line-addressable CRT that is observed by a rolling-shutter camera.
Each barcode scan produces a one-dimensional intensity profile; the CRT
driver writes this profile into a chosen horizontal line or narrow band
of the screen. The rolling-shutter camera then reads the CRT line by
line, converting the two-dimensional pattern on the tube into a
time-staggered \cba stream. In this configuration the CRT
functions as a channelised optical memory: the controller can decide
where to place the most recent scan, effectively multiplexing many scan
histories into a single optical buffer.

\paragraph{Global-shutter and rolling-shutter observation (non-limiting).}
In some embodiments, the CRT surface is observed by a global-shutter
sensor (approximating an instantaneous snapshot) and/or by a
rolling-shutter sensor (in which each sensor row integrates the CRT at
a slightly different time). In some embodiments, the rolling-scan
direction is arranged to be substantially parallel, perpendicular, or
oblique to a CRT sweep direction by mechanical mounting and/or by
introducing a fixed or tunable optical element (for example a plane
mirror, steering mirror, cylindrical lens, or tunable lens) into the
optical path. In tunable embodiments, the distortion element is driven
by one or more analogue control parameters, allowing the geometric
relationship between CRT and sensor to be varied continuously over
time; the element's drive schedule is treated as part of $\theta$ and
recorded in the protocol digest.

\subsubsection{Raster CRT variants}

In some embodiments, a raster CRT is used instead of an XY-deflection
tube. The CRT internal raster is treated as part of the reactor
dynamics, and external scanners sample the raster-refreshed phosphor
field along slower trajectories, yielding temporally down-sampled but
spatially consistent views of the CRT state.

% ------------------------------------------------------------------

% ======================================================================
%
%
%
% ======================================================================

\subsection{Optical and Other Physical Long-Term Memory Substrates}
\label{sec:optical-long-term-memory}

In some embodiments, the \RK apparatus comprises a
physical long-term memory substrate configured to retain one or more
committed-bundle records, protocol-digest records, meter records,
calibration records, attestation records, governance records, or
derived evidentiary artefacts. The physical long-term memory substrate
may be configured so that a write operation produces a persistent
physical state change that is thereafter readable many times and is
not rewritable by ordinary runtime operation of the apparatus. Such a
substrate is referred to herein as a write-once-read-many, or WORM,
substrate. The WORM property is an apparatus property of the storage
medium, the write/read path, or both, and is not merely a software
retention policy.

In some embodiments, an entry written to the physical long-term memory
substrate comprises, or is bound to, a digest of a \cb
$C_{0:T}$, a protocol digest $\Pi$, a meter vector or meter envelope,
a calibration state, a verifier state, a governance-partition state,
an authority signature, an external timestamp, a randomness-beacon
value, a ledger-derived value, or any combination thereof. In some
embodiments, the written entry further comprises a physical-substrate
binding value derived from the measured physical response of the
memory substrate, so that later retrieval verifies not only the
digital content of the entry but also the association between the
entry and the physical memory medium on which it was written.

In a non-limiting embodiment, a record $r_k$ is written to a physical
memory substrate by applying a write operation $W$ to a local region
or mode of the substrate:
\[
    \mu_{k+1} = W(\mu_k, r_k, \Pi_k),
\]
where $\mu_k$ denotes the physical memory state before writing, and
$\Pi_k$ denotes the protocol digest or digest component associated
with the write event. A later read operation $R$ produces a readout
$\hat r_k$ and a substrate-response value $\rho_k$:
\[
    (\hat r_k, \rho_k) = R(\mu_{k+1}, \Pi_k).
\]
A verification function may then test consistency between the recovered
record, the substrate-response value, and one or more committed
digests:
\[
    \operatorname{Verify}_{\mathrm{mem}}
    (\hat r_k, \rho_k, \Pi_k, h(r_k)) \in \{0,1\}.
\]
The foregoing equations are illustrative only; the storage, readout,
and verification functions may be implemented by any physical,
digital, optical, magnetic, molecular, silicon, hybrid, or
cryptographic mechanism.

In some embodiments, access to the physical long-term memory substrate
is physically isolated, logically isolated, cryptographically isolated,
temporally isolated, or any combination thereof from an agent behaviour
generation pathway. For example, an agent may cause operational events
that are recorded by the apparatus, while lacking direct read, write,
erase, or rewrite authority over the memory substrate or its
governance-side index. In some embodiments, retrieval requires an
authority token, a governance-partition command, a threshold signature,
a hardware-attested read path, a selective-opening proof, a proof of
authorised access, or another access-control condition bound to the
protocol digest.

The physical long-term memory substrate may be used to store durable
evidence while preserving privacy by storing encrypted records,
commitments, accumulators, digest paths, selective-opening handles,
reconstruction-threshold records, cohort-publication records,
redacted records, or derived verification results rather than raw
sensor records. In some embodiments, raw records remain unavailable
unless a declared authority condition is satisfied, while digest,
meter, or accumulator records remain available for ordinary
governance audit.

\paragraph{Physical memory substrate generalisation (non-limiting).}
The long-term memory substrate is not limited to an optical memory.
In some embodiments, the substrate comprises any physical storage
medium in which read-many semantics, tamper evidence, substrate
binding, or write-once behaviour is enforced at least in part by an
apparatus property rather than by policy alone. Non-limiting examples
include holographic storage, fluorescent storage, photochromic
storage, photorefractive storage, phase-change storage, multilayer
optical-disc storage, perpendicular magnetic recording, magnetic tape,
MRAM-WORM storage, DNA storage, polymer storage, anti-fuse memory,
one-time-programmable read-only memory, eFuse memory, ceramic or
glass-based archival marks, physically unclonable storage media, and
hybrid physical-digital storage. Optical, magnetic, molecular,
silicon, mechanical, chemical, and hybrid storage species are examples
only.

\paragraph{Representative write/read examples by substrate class
(non-limiting).}
The following non-limiting examples illustrate write, read, and
substrate-response operations in five representative substrate
classes. In each, the write operation produces a persistent
physical state change, the read operation interrogates that state,
the WORM property is an apparatus property of the medium and write
path, and the substrate-response value $\rho_k$ is measured at read
time and bound into the protocol digest of the corresponding write
event.
\begin{itemize}[nosep]
  \item \emph{Optical / holographic class.} In some embodiments,
    the write operation comprises exposing a photosensitive
    medium (for example a photopolymer, photorefractive crystal,
    or photochromic film) at a declared write fluence sufficient
    to produce a refractive-index, absorption, or fluorescent
    state change; the read operation comprises illuminating the
    written region under a declared reference geometry and
    measuring the diffracted, transmitted, or emitted optical
    response; $\rho_k$ is a measured optical response value (for
    example diffraction efficiency, fluorescence intensity, or
    holographic correlation peak) interpreted under a committed
    optical-response model.
  \item \emph{Magnetic / MRAM-WORM class.} In some embodiments,
    the write operation comprises driving a write current
    sufficient to set a magnetic-tunnel-junction or domain-wall
    cell into a declared polarisation state under an
    apparatus-enforced one-shot write path (for example a
    one-time-programmable fuse on the write line); the read
    operation comprises measuring the tunnelling or
    magnetoresistive response under a declared read bias;
    $\rho_k$ is a measured electrical response value (for example
    a tunnelling-magnetoresistance ratio) interpreted under a
    committed cell-response model.
  \item \emph{Molecular / DNA / polymer class.} In some
    embodiments, the write operation comprises a declared
    synthesis or polymerisation sequence producing a persistent
    chemical record (for example a synthesised oligonucleotide
    strand, a templated polymer sequence, or a covalent reaction
    product); the read operation comprises a declared sequencing,
    spectroscopic, or affinity-readout process; $\rho_k$ is a
    measured readout response (for example a sequencing-quality
    score, mass-spectrometric peak, or affinity-binding response)
    interpreted under a committed readout-response model.
  \item \emph{Silicon / anti-fuse / eFuse / OTP class.} In some
    embodiments, the write operation comprises applying a
    one-shot programming voltage sufficient to permanently
    rupture a dielectric, blow a metal fuse, or set a one-time
    register; the read operation comprises measuring the
    post-program resistance, threshold, or logic level under a
    declared read condition; $\rho_k$ is a measured electrical
    response value (for example post-program resistance or
    threshold-voltage shift) interpreted under a committed
    device-response model.
  \item \emph{Ceramic / glass / PUF-mark class.} In some
    embodiments, the write operation comprises a declared
    laser-induced, ion-beam-induced, or mechanically engraved
    permanent mark in a ceramic, glass, or other archival
    substrate; the read operation comprises an optical or
    microscopic measurement of the mark geometry, scattering
    response, or PUF response under a declared imaging geometry;
    $\rho_k$ is a measured response (for example a speckle
    correlation against an enrolled response or a geometric
    fingerprint of the mark) interpreted under a committed
    PUF- or mark-response model.
\end{itemize}
The foregoing classes and operations are illustrative only and do
not limit the substrate, write path, read path, or substrate-response
mechanism of any disclosed embodiment.

% ------------------------------------------------------------------

% ======================================================================
%
%
%
% ======================================================================

\subsection{Bench-top dual-loop evidence appliance}
\label{sec:dual-loop-appliance}
% ------------------------------------------------------------------

In some embodiments, a worked bench-top embodiment (FIG.~4) comprises
a single \RK module with a fast optical loop and a slow
analogue-memory loop under a common controller.  The slow loop
instantiates the generic CRT and phosphor analogue-memory mechanisms
described in Section~\ref{sec:crt-phosphor}; the present subsection
specifies how those mechanisms are combined with a fast scattering
loop, analogue rate conversion between the two timescales, and a
committed dual-loop evidence architecture.

\paragraph{Fast optical loop.}
In some embodiments, the fast loop comprises a visible-wavelength
source~21 (for example a laser diode module), an addressing element~22 (for
example a galvanometer mirror or acousto-optic deflector), relay
optics~23, a
disordered scattering medium~30 (for example a ground-glass plate,
semi-opaque diffuser, or sealed container of mixed scattering
elements), and a detector~41 (for example a frame camera or
photodetector with transimpedance amplifier). In some embodiments, the
addressing element defines a scan law (raster, Lissajous, spiral, or
pseudo-random) over the reactor volume, and the fast-loop timescale
is set by the addressing element's repositioning time (microsecond
scale for an AOD, sub-millisecond for a galvanometer).

\paragraph{Analogue-domain computation prior to digitisation
(non-limiting).}
In some embodiments of the fast optical loop, the loop performs
convolution-like mixing, scattering, interference, and nonlinear
transformation in the analogue optical and electrical domains, with
no analogue-to-digital conversion in the signal path between the
source and the photodetector output. In some embodiments, the
detector produces a continuous analogue electrical signal (for
example a current or voltage from a transimpedance amplifier) that
encodes the result of the physical computation, and digitisation
occurs only after the physical mixing and transformation are
complete. In some embodiments, this fully-analogue computational
signal path within the fast loop is an architectural distinction
from many hybrid photonic computing systems, which digitise optical
signals at layer boundaries for digital readout, re-encoding, and
re-projection between stages. In some embodiments, the absence of
intermediate analogue-to-digital conversion within the fast-loop
computational path means that the committed bundle encodes a
measurement of the analogue physical output, preserving temporal and
amplitude resolution up to the detector's analogue bandwidth, and
that the physical transformation is not discretised or truncated by
converter quantisation noise or sample-rate limitations at
intermediate stages. In some embodiments, other parts of the system
(for example the routed-ensemble digital re-projection path or the
slow-loop CRT write path) may include analogue-to-digital or
digital-to-analogue conversion stages; the fully-analogue property
described here applies specifically to the fast optical loop's
computational signal path.

\paragraph{Analogue rate conversion between fast and slow loops
(non-limiting).}
In some embodiments, the rate conversion between the fast optical
loop (nanosecond to microsecond timescale) and the slow
analogue-memory loop (tens of Hz write cadence) is accomplished in
two stages, both operating in the analogue domain. In some
embodiments, a first stage performs analogue bandwidth reduction on
the fast-loop photodetector output, producing a lower-bandwidth
analogue signal that represents a declared feature of the fast-loop
response. In some embodiments, analogue bandwidth-reduction
mechanisms include one or more of: (i)~envelope detection, in which
a rectifier and low-pass filter extract the amplitude envelope of
the fast-loop signal; (ii)~sample-and-hold, in which a capacitor
captures the fast-loop voltage at a trigger event and holds it for
the slow-loop write cycle; (iii)~integrate-and-dump, in which the
fast-loop signal is integrated over a declared window and the
accumulated charge is transferred to the slow-loop input;
(iv)~analogue peak detection, in which a peak-detector circuit
captures the maximum excursion of the fast-loop signal during a
declared window; and (v)~analogue matched filtering, in which a
passive or active filter with a template impulse response correlates
the fast-loop signal with a declared feature profile. In some
embodiments, a second stage maps the reduced-bandwidth analogue
signal to a CRT write signal (for example an XY deflection voltage
pair and a Z-axis intensity voltage) that deposits the result into
the slow analogue-memory loop at the CRT refresh cadence.

\paragraph{Slow analogue-memory loop.}
In some embodiments, the slow loop comprises a CRT reactor~31 or related
persistence display and a camera~42 that observes the phosphor
screen~32. In
some embodiments, the CRT phosphor screen~32 provides analogue temporal
memory with a persistence time constant that depends on the phosphor
species. In some embodiments, the camera captures frames at a rate
slower than the fast-loop bandwidth, producing a committed bundle
that encodes a time-integrated view of the evolving phosphor state.

\paragraph{Colour-phosphor parallel temporal filter bank
(non-limiting).}
In some embodiments using a colour CRT or colour phosphor screen
comprising two or more phosphor species with measurably distinct
persistence time constants (for example, in phosphor sets where
$\tau_R$, $\tau_G$, $\tau_B$ --- instances of the general
persistence time constants $\tau_i$ described above --- differ by at
least an order of magnitude), each phosphor species implements an
independent temporal convolution channel: the same excitation pattern
is deposited into all phosphor species simultaneously, and each
species integrates and decays at its own characteristic rate, so that
at any readout instant the luminance from each phosphor species
represents a differently time-weighted history of the excitation
sequence. In some embodiments, a colour camera separates these
channels, yielding a set of parallel temporal filter outputs from a
single phosphor screen. In some embodiments, the number of
independent temporal channels is equal to the number of phosphor
species with sufficiently distinct time constants, and whether three
species provide measurably independent temporal information is an
empirical property of the specific phosphor set that must be
validated by measuring inter-channel correlation under representative
excitation sequences.

\paragraph{Tuneable readout colour weighting (non-limiting).}
In some embodiments, the camera used to observe the persistence
display is augmented with a tuneable colour filter or colour-channel
gain control that allows the relative contribution of each phosphor
channel to the readout to be adjusted, thereby selecting among
different temporal weighting profiles without modifying the phosphor
screen itself.

\paragraph{Channelised optical memory (non-limiting).}
In some embodiments, the colour-phosphor temporal filter bank is
used as a channelised optical memory in which the controller
selectively addresses different phosphor channels (for example by
varying the electron-beam colour convergence, masking, or excitation
spectrum) so that independent information streams are deposited into
different temporal-persistence channels and read out through spectral
separation.

\paragraph{Signed-weight realisation via differential exposure
(non-limiting).}
In some embodiments, signed effective weights are realised by
differential exposure: two phosphor patches, spectral bands,
temporal windows, or spatial regions are driven with the same
excitation magnitude but opposite intended signs, and a subsequent
readout stage computes their difference (for example by subtracting
two camera colour channels, two spatial regions, or two temporally
interleaved frames). In some embodiments, the resulting signed
readout approximates a single bipolar weight without requiring a
bipolar physical medium. In some embodiments, whether differential
exposure produces signed weights with sufficient precision, dynamic
range, and stability for a given task is an empirical property that
must be validated per embodiment.

\paragraph{Multi-CRT summation (non-limiting).}
In some embodiments, multiple CRTs or persistence displays are
observed simultaneously by a single camera, and the overlapping
phosphor images are summed optically on the detector surface. In some
embodiments, this optical summation realises element-wise addition
of multiple analogue-memory states without digital intervention.

\paragraph{XY-deflection synchronisation with external scanners
(non-limiting).}
In some embodiments, the CRT XY deflection voltages are synchronised
with the fast-loop addressing element (AOD, galvanometer, or SLM) so
that each fast-loop scan point maps to a declared position on the CRT
phosphor surface. In some embodiments, this synchronisation enables
the CRT to serve as a spatially resolved analogue memory of the
fast-loop computation results.

\paragraph{Channelised memory with commodity hardware
(non-limiting).}
In some embodiments, the channelised optical memory and
colour-phosphor filter bank described above are implemented using
commodity colour CRT monitors and webcams or industrial cameras with
Bayer-pattern colour filter arrays. In some embodiments, commodity
hardware operates within declared calibration envelopes and its
limitations (for example colour cross-talk, limited persistence
uniformity, and limited spectral separation) are characterised and
metered.

\paragraph{Operating modes and commitment.}
In some embodiments, the dual-loop appliance operates in all three
regimes (verification, sensing, transformation) and produces a
committed bundle covering both loops. In some embodiments, the
protocol digest records the fast-loop scan law, the analogue
rate-conversion mechanism, the CRT write parameters, the camera
readout cadence, and the phosphor persistence model.

\paragraph{Bring-up and validation (non-limiting).}
In some embodiments, the bring-up procedure includes: (a)~measuring
the CRT phosphor persistence time constant per colour channel,
(b)~measuring the CRT spot size and positional accuracy,
(c)~measuring the camera-to-CRT geometric alignment and applying a
calibration warp, (d)~verifying the analogue rate-conversion stage by
driving the fast loop with known excitation sequences and comparing
CRT-deposited patterns against predicted patterns, and
(e)~establishing a baseline verisimilitude score using the trained
discriminator on committed bundles from the validated configuration.

\paragraph{Servo-swept scanning reactor variant (non-limiting, by cross-reference).}
In some embodiments, a compact bench-top variant of the present module
follows the bench-top 1D fluorescence-substrate anchor of
Section~\ref{sec:anchor-1d}, used as the reactor for the present
multi-rate / dual-loop architecture in place of the CRT and scattering
combination. The geometry, source arrangement, lens-and-servo control,
sealed container, spectral filtering, scan-law construction, and
fluorescent-persistence properties of that anchor apply unchanged in
the present variant.

\paragraph{Computational interpretation as gated recurrence
(non-limiting).}
In some embodiments, the bench-top dual-loop evidence appliance is
interpreted as a gated recurrent architecture in which: (i)~the fast
optical loop implements a nonlinear input-to-hidden transformation
via scattering and analogue mixing, (ii)~the analogue rate-conversion
stage acts as a gating mechanism that selects which features of the
fast-loop state are deposited into long-term memory, (iii)~the CRT
phosphor screen implements analogue hidden-state memory with
exponential decay, (iv)~the camera readout implements a measurement
of the hidden state, and (v)~the controller's emission at the next
time step may depend on the camera readout, closing the recurrence.
In some embodiments, the key distinction from digital gated recurrent
architectures is that the gating, state storage, and recurrence
operate in the analogue physical domain rather than in digital
arithmetic, and the committed bundle provides an auditable record of
the recurrent dynamics.


% ======================================================================

\subsection{Fibre-delay reactor embodiments}
% ------------------------------------------------------------------

In some embodiments, the reactor includes one or more fibre-delay
networks that implement recurrent temporal kernels in the optical
domain. A fibre-delay reactor may include splitters that distribute an
optical field into multiple delay loops with different delays
$\{\tau_\ell\}$, optional phase shifters and attenuators, and combiners
that recombine delayed fields. Parameters may include coupling ratios,
phase shifts, selected taps, and, in some variants, selectable loop
lengths. These parameters form a subset of $\theta$ and may be tuned at
design time or as slow controls.

\paragraph{Impulse-response lens (non-limiting).}
In some embodiments, in a small-signal or narrowband regime, a
fibre-delay network is modelled as a linear temporal filter with a
sparse impulse response:
\[
  h_\theta(t) = \sum_{\ell} a_\ell e^{i\varphi_\ell}\,
  \delta(t-\tau_\ell),
\]
where $a_\ell$ and $\varphi_\ell$ are non-limiting per-loop amplitude
and phase parameters (part of $\theta$ and/or slow controls). An input
field envelope $E_{\mathrm{in}}(t)$ yields an output
\[
  E_{\mathrm{out}}(t) = \int h_\theta(\tau)\,
  E_{\mathrm{in}}(t-\tau)\,d\tau,
\]
so the reactor implements a physically recurrent temporal kernel whose
effective impulse response is determined by the loop delays and
couplings.

% ------------------------------------------------------------------
\subsection{Acoustic and cymatic reactor embodiments}
% ------------------------------------------------------------------

In some embodiments, a non-optical or multi-physics reactor is used,
including acoustic and cymatic variants. A speaker--microphone pair may
drive and sense a membrane, plate, or fluid layer at a carrier frequency
with a controlled envelope. The membrane patterns may be observed
optically, producing visual features that are included as channels of
$\mathbf{y}_t$. Acoustic features (for example envelope or band-energy
features over windows) may also be included. Optical and acoustic
controls may be coupled through a shared controller and shared latent,
and hardness or verisimilitude meters may be trained on the joint
multi-channel \cb.

\paragraph{Acoustic control and \cba notation
(non-limiting).}
In some embodiments, acoustic control at step $t$ is expressed as
\[
  U_t^{\mathrm{ac}}
  = \bigl(s_{\mathrm{ac}}(t),\; a_{\mathrm{drv}}(t)\bigr),
\]
where $s_{\mathrm{ac}}(t)$ denotes a carrier or waveform choice and
$a_{\mathrm{drv}}(t)$ denotes a drive envelope or amplitude schedule.
Acoustic observations may be represented as an audio channel of the
the \cb, while camera-observed membrane patterns may be
represented as a visual channel. Meters and downstream heads may
operate on either channel or on the joint law to exploit cross-domain
constraints.

\paragraph{Infrasonic channels and cymatic pattern analysis
(non-limiting).}
In some embodiments, acoustic channels include infrasonic components
(below approximately 20~Hz) and mid- to high-frequency modes capable
of exciting visible cymatic patterns in fluid, granular, or membrane
media. Non-limiting uses include long-range acoustic sensing, detection
of architectural or geophysical resonances, and acoustic-to-visual
frequency analysis via optical scanning of cymatic patterns. Safety
constraints limit infrasonic emission power and duty cycle to levels
below those associated with adverse physiological effects.

\paragraph{Ultrasonic phased arrays and haptic fields (non-limiting).}
In some embodiments, the acoustic channel is ultrasonic and is driven
by a phased array that steers and focuses pressure fields in air, water,
or tissue. In \RT embodiments, focused ultrasound produces
haptic or vibrotactile sensations and the \cb logs the
delivered field schedule for audit; in \LI embodiments, the same
array implements ultrasound imaging or elastography under logged
steering and safety envelopes.

\paragraph{Parametric acoustic emission and self-demodulation in air (non-limiting).}
In some embodiments, two or more ultrasonic carriers at
frequencies $f_1, f_2, \ldots$ are emitted along a common
axis and undergo nonlinear self-demodulation in air,
generating audible difference-frequency components at
$|f_i - f_j|$ whose beam directionality is set by the
ultrasonic carrier wavelengths rather than by the audible
wavelengths. In \RT embodiments, the swept parametric
audio beam delivers spatially localised sound to one or
more listeners without enclosing loudspeakers, with the
delivered field schedule (carrier frequencies, modulation
envelope, beam direction, exposure duration) committed to
the protocol digest and \cb on the same footing as
optical sweeps. In \TB embodiments, the directional beam
serves as an acoustic challenge directed at a specific
scene region, with the return signal recorded by one or
more microphones as an acoustic observation channel and
folded into the \cb. In \LI embodiments, the parametric
beam supports directional acoustic imaging without the
enclosure geometry required by audible-frequency arrays.
Device-specific tolerances of the ultrasonic emitter
array (element mismatch, phase calibration, near-field
beam structure, carrier intermodulation profile)
contribute reactor-microstructure uniqueness in the
acoustic channel symmetric to the optical-emitter case
(Section~\ref{sec:emitter-reactor}). Safety envelopes
limit carrier intensity, duty cycle, and listener
exposure to levels below applicable thresholds for
ultrasonic and audible exposure.

\paragraph{Phonon lasers and coherent acoustic emission (non-limiting).}
In some embodiments, the acoustic emitter is a phonon laser
(sound-amplification-by-stimulated-emission-of-radiation device, SASER)
in which coherent acoustic phonons are generated by stimulated emission
in a semiconductor superlattice, optomechanical cavity, trapped-ion
chain, or crystalline resonator, rather than by a driven transducer.
The resulting emission is a coherent, narrowband acoustic beam whose
frequency, linewidth, and threshold behaviour are determined by the
phononic cavity geometry, lattice parameters, and gain-medium
properties---analogous to optical laser operation but in the acoustic
domain. Device-specific fabrication tolerances in superlattice period,
cavity dimensions, defect distributions, and mechanical $Q$ factors
provide reactor-microstructure uniqueness in the acoustic channel, symmetric to the
optical emitter-as-reactor case (Section~\ref{sec:emitter-reactor}).
Near-threshold operation exposes stochastic phonon emission and
mode-competition noise suitable for stochastic sampling in PoliePuter
mode. In \LI embodiments, the coherent acoustic beam provides
high-resolution probing of material elastic properties, subsurface
structure, or biological tissue via phonon scattering and absorption;
in \TB embodiments, the device-specific phonon emission
fingerprint provides verification of the acoustic source. In Yoked
embodiments, acoustic backscatter or material response feeds back to
the SASER's gain or detuning, creating a closed-loop phononic--material
coupled system. In some embodiments, a SASER is co-integrated with an
optical reactor path, providing a joint opto-acoustic channel whose
cross-domain correlations supply additional hardness constraints.

% ------------------------------------------------------------------
\subsection{Magnetic and other multi-physics bistatic reactor
embodiments}
% ------------------------------------------------------------------

In some embodiments, a reactor is implemented using a multi-physics
channel rather than a purely optical channel. Non-limiting examples
include magnetic, electrochemical, mechanical, photoacoustic,
opto-acoustic, and other coupled sensing and excitation channels in
which an emission input and a sensing output are linked through a
physical medium with hysteresis, relaxation, or amplitude-dependent
nonlinearity. These bistatic embodiments preserve the same high-level
structure as optical reactors: a controller selects a protocol $u(t)$,
an excitation is applied through a transducer, and a measurement is
recorded as an observable stream whose summaries are committed and
audited.

\subsubsection{Mechanical and inertial microstructure embodiments
(non-limiting)}

In some embodiments, the reactor medium includes embedded inertial
microstructures or MEMS-scale masses whose motion is driven by reactor
excitation and by external perturbations (for example vibration or
acceleration). The resulting mechanical state evolves with slow
relaxation, providing a mechanical memory coupled to an optical,
magnetic, or acoustic path. One or more inertial sensors
(accelerometers, gyros) may be logged as auxiliary channels and used by
meters to distinguish external perturbations from internal reactor
evolution.

\subsubsection{Aqueous, hydrogel, and liquid-crystalline reactor media
(non-limiting)}

In some embodiments, the reactor medium comprises aqueous solutions,
hydrogels, or liquid-crystalline phases whose optical, electrical, or
mechanical properties vary with structure, temperature, solute
concentration, or preparation history. Such media may exhibit ordering
near interfaces, metastable structural states, sensitivity to
electromagnetic or acoustic fields, and nonlinear optical responses
arising from molecular alignment. Reactor behaviour in these media is
characterised empirically. Liquid-crystal phases with well-documented
ordering transitions are preferred for reproducible embodiments.

\subsubsection{Biological and optogenetic reactor surfaces
(non-limiting)}

In some embodiments, the reactor medium includes biological surfaces or
bio-hybrid layers that can be stimulated and recorded. Non-limiting
examples include optogenetic stimulation with optical recording,
fluorescence or reflectance readouts, and calcium-indicator
measurements, where structured stimulation patterns cause measurable
responses whose dynamics provide mixing, persistence, and channel
constraints.

\paragraph{Biophoton detection and biological reactor coupling
(non-limiting).}
In some embodiments, the system detects ultra-weak photon emission from
biological tissues, arising from oxidative metabolic processes and
electronically excited molecular species. Single-photon-sensitive
detectors measure emission rates, spectral characteristics, and spatial
patterns from living samples, and these readouts are appended to
the \cb as additional channels. In some embodiments,
structured illumination (including beams carrying orbital angular
momentum or other higher-order spatial modes) is directed through
biological tissue, and mode mixing introduced by the tissue's
refractive-index inhomogeneity is recorded as part of the \cb. In such embodiments, the tissue acts as a scattering medium
whose transfer function contributes device-specific and
specimen-specific features to the observation, and the preserved or
transformed spatial mode structure provides sensitivity to
microstructural properties of the biological sample.  OAM and
higher-order spatial modes require sufficient spatial coherence
to propagate recognisably; in thick or highly scattering tissue,
mode structure is destroyed over path lengths exceeding a few
transport mean free paths.  These embodiments are therefore most
applicable to thin tissue preparations, surface measurements, or
waveguide-coupled samples where the optical path through
scattering media is short.


\paragraph{Optically responsive biological reactor media (non-limiting).}
In some embodiments, the reactor medium comprises a biological
preparation in which one or more cellular or subcellular components
are both optically addressable and optically emissive, so that the
control input $u(t)$ modifies a biological state and the resulting
optical emission encodes the consequence of that perturbation as
channels of $\mathbf{y}_t$.  Non-limiting classes of optically
responsive protein that may serve as reactor elements include:
naturally occurring emitters such as luciferase variants and
fluorescent proteins (GFP and its derivatives) that report gene
expression state or protein localisation; genetically encoded
state reporters such as GCaMP calcium indicators and genetically
encoded voltage indicators (GEVIs) whose emission intensity or
spectrum encodes electrophysiological activity; optogenetic
actuators such as channelrhodopsin variants (ChR2, step-function
opsins, red-shifted opsins) that transduce structured illumination
into ionic currents; and subcellular-targeted constructs such as
mitochondria-directed opsins engineered for import to the inner
mitochondrial membrane using one or more mitochondrial targeting
sequences, enabling structured illumination to modulate
mitochondrial energetic state, calcium handling, and effective
neuronal response gain.  In neural tissue and neural organoid
embodiments, the preparation may additionally provide endogenous
optical signals including NADH and FAD autofluorescence reporting
metabolic redox state, membrane-potential-sensitive dyes such as
TMRM or JC-1, genetically encoded mitochondrial reporters, and
ultra-weak biophoton emission associated with oxidative metabolic
processes, all appended as channels of $\mathbf{y}_t$ and
summarised in the protocol digest under declared meter envelopes.
In some embodiments a two-stage control structure is employed:
a first emission pattern sets the gain state of a declared target
region (for example by modulating mitochondrial energetic state or
by driving an optogenetic actuator to a declared operating point);
a second emission pattern then probes the preparation's response
under that gain state; and the detection channel records the
consequence.  The gain-setting and probe patterns, their wavelength
bands, power envelopes, and timing are committed in the protocol
digest before the run begins, supporting post-hoc audit of
illumination conditions independently of the biological readout.
In preferred embodiments, safety envelopes include irradiance
limits, duty-cycle limits, and phototoxicity-monitoring channels
appended to $\mathbf{y}_t$; these embodiments do not assert
specific therapeutic or clinical outcomes.

\paragraph{Piezoelectric biological interfaces (non-limiting).}
In some embodiments, the reactor couples to piezoelectric biological
materials (including bone, collagen, tendon, and engineered tissues)
that generate electrical signals in response to mechanical stress and
conversely deform in response to applied fields.

\paragraph{Chemical and olfactory reactor surfaces (non-limiting).}
In some embodiments, a proxy scene or reactor surface integrates
volatile chemistry, fluid chemistry, or other non-optical fields over
time and exposes an evolving state that the \RK interrogates
optically and/or electrically. Non-limiting examples include gas-sensor
plates, optoelectronic olfactory receptor arrays, microfluidic channels
with dye indicators, and living tissue cultures. The controller may emit
tracer odorant pulses or other chemical stimuli under a logged schedule,
and detectors record temporal response traces as additional channels of
$\mathbf{y}_t$. In \TB embodiments, the stimulus schedule and
response statistics define a challenge--response hardness source; in
\LI embodiments, the same channels support inference; and in \RT embodiments, the controller steers masking or room-profile
matching by emitting bounded stimuli under safety and policy
constraints.

\subsubsection{Bioelectric and brain-scan embodiments (non-limiting)}

In some embodiments, the scene and/or a ponderable cell comprises a
bioelectric or neurophysiological field, such as electroencephalography
(EEG), magnetoencephalography (MEG), electrocorticography (ECoG),
electromyography (EMG), peripheral nerve recordings, or related
modalities. The module may include a sensor cap, electrode array, field
probes, or transducer surfaces whose effective spatial sensitivity is
time-varying under control. In these embodiments, the control input
selects, without limitation, montages or spatial filters, sampling
schedules, carrier frequencies, drive amplitudes within safety
envelopes, stimulation sites (where permitted), or structured sweep
protocols across the array, and the resulting observations are appended
to the \cba record as channels of $\mathbf{y}_t$.

\paragraph{Frequency-locked EEG--optical coupling (non-limiting).}
In some embodiments, bioelectric coupling is implemented as a
frequency-locked optical channel in which recorded multi-channel neural
signals are combined using swept weighting waveforms defining a
time-varying montage or spatial filter. Intermediate combined signals
modulate one or more optical carriers that traverse a reactor loop, so
that scattering and interference progressively merge multiple channels
into a reduced set of optical outputs. The resulting detector readout
provides channels included in the \cb and can be trained to
track a desired low-dimensional functional of the neural signals.

\paragraph{Neural entrainment and sensory pacing (non-limiting).}
In some embodiments, the controller generates sensory stimuli (visual
flicker, auditory pulses, tactile signals, or combinations) at
frequencies selected to influence oscillatory activity via entrainment,
optionally with closed-loop adaptation based on measured EEG or
behavioural signals. In preferred embodiments, stimulation is
constrained by safety policies (including photosensitive-seizure risk
controls) and is logged with sufficient detail for later audit.

\subsubsection{Safety envelope and calibration (non-limiting)}

In some embodiments, multi-physics and biological embodiments operate
under explicit safety envelopes and calibration policies. Protocols may
be constrained by maximum energy, duty cycle, exposure time, thermal
bounds, and interlock states, and meters record compliance with these
bounds together with drift indicators and calibration status. Sealed or
physically confined reactor loops are used when appropriate to isolate
the reactor medium from the environment.

% ------------------------------------------------------------------
\subsection{Resonant and frequency-domain reactor embodiments}
% ------------------------------------------------------------------

In some embodiments, the reactor is structured as a resonant cavity
(optical, acoustic, electromagnetic, or mechanical) with characteristic
modes and quality factor. The system may encode information via
amplitude, frequency, phase, polarisation, or modulation of carriers.
The \cb may be analysed in the time domain and/or frequency
domain (for example via Fourier or wavelet transforms), and frequency
signatures may serve as features for \LI tasks, calibration, or
verification. In Yoked operation, resonant reactors are particularly
natural: the resonant modes provide the spectral structure that
Arnold-tongue analysis maps, and changes in the scene that alter the
cavity's resonant response are directly encoded as shifts in the
entrainment diagram.

In a non-limiting example, the reactor is a Fabry--P\'erot optical
cavity (two partially reflecting mirrors separated by a controlled
gap) containing an absorptive or scattering medium (for example a
dye-doped polymer film, nanoparticle suspension, or vapour cell).  The
emitter is a tuneable laser whose wavelength is swept across one or
more cavity resonances; the detector records transmitted intensity as
a function of wavelength.  The resonant linewidth, free spectral
range, and finesse depend on mirror reflectivities, cavity length, and
the medium's absorption and scattering properties---all of which carry
device-specific fabrication signatures.  In \TB mode, the
measured transmission spectrum is compared against a stored device
profile to verify device identity.  In \LI mode, changes in the
cavity's resonant response (for example shifts caused by a sample
placed inside the cavity or by temperature variation) are used to
infer scene or sample properties.  Protocol digests record laser
sweep parameters, mirror alignment indicators, and temperature.
Meters record cavity stability, mode-hop indicators, and
signal-to-noise ratios.

% ------------------------------------------------------------------
\subsection{Gain-pumped cavity Reality Kernel embodiments}
% ------------------------------------------------------------------

In some embodiments, the reactor is realised as a gain-pumped optical
cavity in which state evolves through repeated optical passes prior to
detection. A gain medium may be provided by a semiconductor optical
amplifier, fibre amplifier, doped waveguide, Raman gain, parametric
gain, or other amplifying structure. Feedback elements and output
couplers define the cavity dynamics. Control inputs may include pump
power, phase shifter settings, modulator settings, tunable filters, and
optional weak injection of seed fields. The detector observes one or
more outputs of the cavity and records them as channels of the
the \cb. Such embodiments may be operated in \TB, \LI, or
\RT regimes. Optional quantum-sensitive variants and
entangled or squeezed-light variants are described in the quantum
coupling section (Section~\ref{sec:quantum}).

\paragraph{Threshold, bistability, and mode competition (non-limiting).}
In some embodiments, gain-pumped cavities exhibit threshold behaviour,
bistability, and mode competition that produce rich, history-dependent
signatures. In a non-limiting rate-equation lens for mode amplitudes
$\{a_k(t)\}$,
\[
  \dot{a}_k(t)
  =
  \bigl(g_k(p(t)) - \ell_k\bigr)\,a_k(t)
  -
  \sum_j \beta_{kj}\,|a_j(t)|^2 a_k(t)
  +
  w_k(t),
\]
where $p(t)$ is a pump control, $\ell_k$ are loss terms, $\beta_{kj}$
model gain saturation and cross-saturation (mode competition), and
$w_k(t)$ denotes noise. Meter envelopes record pump and thermal
margins, stability indicators, and saturation proxies; protocol digests
record pump schedules and tuning settings so that verification and
reconstruction are conditioned on the intended cavity regime.

\subsection{Multi-pass folded-cavity reactor embodiments
(infinity-mirror architecture)}

Iterative optical feedback loops using incoherent light and commodity
components---spatial light modulators of the type used in miniature LCD
panels, LED-class light sources, and compact cameras---have been
reported to generate nonlinear input--output transformations through
cumulative multi-pass interaction, without requiring exotic nonlinear
materials or high-power lasers. However, reported systems do not
provide \cba logging, multi-regime objective selection,
or committed evidence coupling. In some \RK embodiments,
the principle of building effective nonlinearity from individually
linear or weakly nonlinear components is adapted as a reactor
architecture with the addition of these disciplines.

In some embodiments, a reactor is implemented as a compact folded-cavity
or ``infinity mirror'' arrangement in which light is repeatedly routed
through one or more miniature optical elements (for example LCD-type
spatial light modulators, diffusers, thin-film layers, and/or
polarisation elements) arranged between partially reflecting surfaces.
Each pass through the element stack contributes an incremental nonlinear
transformation; the cumulative effect of $N$ passes produces an
effective nonlinear mapping whose complexity grows with $N$. A compact
camera or photodetector captures the output after a controlled number
of passes or at steady state.

In some embodiments, the spatial light modulator in the multi-pass loop
is driven by the \RK controller, so that the per-pass
transformation is a trainable or protocol-controlled function. The SLM
pattern may be updated between passes (if the SLM refresh rate permits)
or held fixed for a burst of passes and then updated, providing two
timescales of control analogous to the fast-loop/slow-loop architecture
of the bench-top 1D \RK (Section~\ref{sec:anchor}).

In some embodiments, the multi-pass reactor is coupled to the scene loop
so that scene return signals modulate the SLM pattern, the source
intensity, or an auxiliary element in the cavity, creating a
scene-coupled feedback reactor from commodity components. In Yoked
operation, the scene's dynamical response feeds back into the cavity
parameters, and the multi-pass nonlinearity amplifies coupling
signatures that would be invisible in a single-pass system.

In some embodiments, reactor-microstructure uniqueness arises from manufacturing
tolerances in the SLM pixel structure, partial-reflector coatings,
cavity alignment, and any diffuser or scattering elements in the
multi-pass path. These tolerances accumulate multiplicatively over $N$
passes, so that small per-element variations produce large and
device-specific cumulative effects on the output pattern.

In some embodiments, the multi-pass cavity is sealed and miniaturised
(for example integrated into a module comparable in size to a small
camera), making it suitable for portable or embedded \RK
deployments where a CRT reactor is impractical.

\subsection{Phase-change material reactor embodiments}

In some embodiments, reactor layers include one or more phase-change
materials (PCMs) whose optical properties (refractive index, absorption,
or both) switch between two or more stable states under optical or
electrical stimulus. Non-limiting PCM materials include antimony
selenide (Sb$_2$Se$_3$), germanium-antimony-telluride (GST, for example
Ge$_2$Sb$_2$Te$_5$), antimony sulphide (Sb$_2$S$_3$), and
vanadium dioxide (VO$_2$). These materials provide
\emph{nonvolatile optical memory at zero static power}: once switched, the optical state is retained without energy input over the device's declared retention window (subject to material- and device-specific retention, endurance, and resistance-drift characteristics), in contrast to phosphor persistence which decays on millisecond-to-second timescales. In some embodiments, the PCM layer
is addressed by a focused laser beam or structured illumination pattern
from the \RK controller, so that spatial patterns of
amorphous and crystalline regions are written, read, and overwritten
as part of the reactor's state evolution.

In some embodiments, PCM layers are integrated into photonic integrated
circuits (for example as cladding overlays on silicon waveguides or
microring resonators) to provide programmable, nonvolatile weight
elements in an on-chip reactor. In some embodiments, multi-level
states (for example six-bit phase levels demonstrated with
nitrogen-doped Sb$_2$Se$_3$) increase the state-space richness of
the reactor medium.

In some embodiments, PCM reactor layers are combined with other reactor
elements (for example scattering plates, phosphor layers, or nonlinear
thin films) in a layered reactor stack, where the PCM layer provides
nonvolatile memory and the other layers provide mixing, persistence,
or nonlinearity. In some embodiments, reactor-microstructure uniqueness arises from
stochastic crystallisation nucleation, grain boundary formation, and
defect distributions in the PCM layer, which vary from device to device
even under identical fabrication processes.

\subsection{Spiking photonic reactor embodiments}

In some embodiments, one or more reactor elements are implemented as
spiking photonic neurons---optical devices that produce discrete
spike-like output pulses in response to input signals exceeding a
threshold, with excitability, refractory periods, and history-dependent
dynamics analogous to biological neurons.

In some embodiments, spiking photonic neurons are implemented using
vertical-cavity surface-emitting lasers (VCSELs) biased near threshold,
where polarisation switching or injection-locking dynamics produce
GHz-rate spike trains whose timing and rate encode input intensity and
history. VCSELs are manufactured at wafer scale for data-communication
and smartphone applications, providing a sub-dollar, mass-producible
nonlinear reactor element. In some embodiments, reactor-microstructure uniqueness
arises from device-specific threshold characteristics, polarisation
switching points, and thermal response curves that vary with
fabrication tolerances.

In some embodiments, spiking photonic neurons are implemented using
microring resonators on silicon-on-insulator platforms, where
carrier-induced nonlinear dynamics produce self-pulsation and spike
generation at rates of hundreds of MHz. In some embodiments, spiking
neurons are implemented using resonant tunnelling diode (RTD)
opto-electronic circuits, providing THz-bandwidth nonlinear excitability.

In some embodiments, arrays of spiking photonic neurons are coupled
through optical waveguides, free-space paths, or fibre links to form
spiking neural network reactors in which the collective spiking dynamics
provide mixing, temporal memory (through refractory periods and synaptic
integration), and nonlinear computation. In some embodiments, the
spiking dynamics of the reactor are coupled to the scene loop so that
scene return signals modulate the spike rates or timing, providing a
natural interface for Yoked operation in which the coupled
scene--reactor spiking dynamics form a joint attractor.

\subsection{Perovskite and adaptive-luminescence reactor embodiments}

In some embodiments, reactor media include perovskite-based or other
adaptive-luminescence materials whose photoluminescent response evolves
as a function of excitation history, providing neuromorphic dynamics
(for example facilitation, depression, potentiation, and spike-like
threshold behaviour) directly in the optical domain.

In some embodiments, mixed-halide perovskite nanocrystals are embedded
in a porous host matrix (for example macroporous rare-earth oxide such
as Y$_2$O$_3$:Eu$^{3+}$), where interface-mediated halide migration
under optical excitation creates progressively evolving energy barriers
that modulate photoluminescence intensity and spectrum. Such materials
simultaneously provide: (i)~nonlinear optical response (threshold and
saturation behaviour), (ii)~programmable memory (history-dependent PL
evolution), (iii)~spectral mixing (host and guest emission bands that
shift with state), and (iv)~reactor-microstructure uniqueness (stochastic
nanocrystal distribution and interface morphology). In some embodiments,
these properties are obtained from a single solution-processable layer,
providing a low-cost reactor medium suitable for mass production.

In some embodiments, tin-based perovskite thin films (for example
PEA$_2$SnI$_4$) provide giant nonlinear refractive indices (on the
order of cm$^2$/GW) on flexible substrates, enabling compact nonlinear
reactor layers via solution processing.

In some embodiments, persistent luminescent materials (for example
SrAl$_2$O$_4$:Eu$^{2+}$,Dy$^{3+}$ or Li$^+$/Dy$^{3+}$ co-doped
Sr$_2$SiO$_4$:Eu$^{2+}$) provide optical memory with controllable
decay times from seconds to hours, extending the temporal integration
window of phosphor-based reactors beyond that achievable with
conventional CRT phosphors.

In some embodiments, mechanoluminescent materials (for example
ZnS:Mn$^{2+}$ nanocrystals or doped strontium silicate) emit light in
response to mechanical stress, providing a transduction pathway from
acoustic, vibrational, or tactile scene properties into the optical
domain. In some embodiments, mechanoluminescent reactor layers are
used in Yoked operation to couple to structural vibrations, acoustic
fields, or haptic interactions with the scene.

\subsection{Photonic integrated circuit reactor embodiments}

In some embodiments, the reactor is implemented as a photonic integrated
circuit (PIC) comprising programmable interferometric elements (for
example arrays of tuneable Mach-Zehnder interferometers in a mesh
topology, microring resonator arrays, or multimode interference
couplers) fabricated on a semiconductor substrate (for example silicon
on insulator, silicon nitride, or indium phosphide).

In some embodiments, the PIC reactor provides programmable linear
transformations (matrix-vector multiplication) combined with intrinsic
nonlinearity from detector saturation, carrier effects in resonators,
or integrated nonlinear elements. In some embodiments, the PIC reactor
provides tuneable feedback paths through on-chip waveguide loops,
creating recurrent dynamics suitable for reservoir computing and
temporal memory.

In some embodiments, reactor-microstructure uniqueness arises from fabrication-induced
random phase variations in the interferometric mesh, resonance offsets
in microring arrays, and waveguide dimension tolerances that produce
device-specific transfer matrices even for nominally identical designs.
In some embodiments, the PIC reactor is coupled to the scene loop
through fibre or free-space optics, with the controller programming
the PIC configuration as part of the scan protocol and logging the
PIC state in the protocol digest.

\subsection{Detector-as-reactor: stochastic sampling in the image
sensor (non-limiting)}
\label{sec:detector-reactor}

In some embodiments, the detector itself --- for example a CMOS image
sensor --- serves as a stochastic sampling element within the optical
path. Every pixel in a CMOS sensor exhibits multiple noise sources whose
statistics are device-specific and physically grounded:
\begin{itemize}[nosep]
  \item \textbf{Shot noise:} Poisson-distributed photon arrival
    statistics (quantum in origin).
  \item \textbf{Read noise:} thermal and $1/f$ noise in the readout
    electronics, with per-pixel variance determined by transistor
    geometry and fabrication tolerances.
  \item \textbf{Dark current:} thermally generated charge carriers
    whose rate varies per pixel due to crystal defects, implant
    variations, and localised contamination.
  \item \textbf{Fixed-pattern noise (FPN):} systematic per-pixel offset
    and gain variations from manufacturing, providing a spatial
    fingerprint that is stable over time and unique to each sensor.
\end{itemize}

In some embodiments, the system deliberately operates the sensor in a
noise-dominated regime --- for example at very low illumination, very
short exposure, high gain (ISO), or elevated temperature --- so that the
readout at each pixel is primarily a sample from the pixel's
noise distribution rather than a deterministic record of incident
light. In this regime, the emitter (projector) programs the
\emph{energy landscape} of the sampling process: the structured
illumination pattern sets the mean photon flux at each pixel, biasing
the noise distribution without overwhelming it. The sensor's physics
performs the sampling, and the \cb records both the
illumination control parameters and the noisy readout, preserving the
standard emission-observation format. This energy-landscape
interpretation is strongest when the forward model from projector pixel
values through the scene or reactor to sensor photon flux is stable (for
example in a sealed reactor or a static calibration scene), so that the
dominant source of per-frame variation is sensor noise rather than
unmodelled scene dynamics. In embodiments where the scene is dynamic,
the interpretation holds per-frame under the assumption that scene
changes between consecutive frames are slow relative to the sensor's
integration time.

\paragraph{Energy-based model interpretation (non-limiting).}
In some embodiments, the per-pixel readout in the noise-dominated
regime is modelled as a sample from a conditional distribution
\[
  y_i \;\sim\; p\!\left(y_i \mid \lambda_i(\theta, u),\;
  \sigma^{\mathrm{read}}_i,\; d_i,\; \mathrm{FPN}_i\right),
\]
where $\lambda_i$ is the expected photon flux at pixel $i$ (set by the
projector pattern $u$ and the scene/reactor configuration $\theta$),
$\sigma^{\mathrm{read}}_i$ is the pixel-specific read noise,
$d_i$ is the pixel-specific dark current rate, and $\mathrm{FPN}_i$ is
the pixel-specific offset. The projector pattern provides the
programmable bias (analogous to the energy-function parameters in an
energy-based model), and the pixel noise statistics provide the
stochastic dynamics that produce samples. Changing the projector pattern
changes the distribution from which the sensor samples, enabling
programmable stochastic computation within the existing RK optical path.

\paragraph{reactor-microstructure uniqueness from sensor noise (non-limiting).}
In some embodiments, the combination of fixed-pattern noise, dark
current non-uniformity, and per-pixel read noise statistics constitutes
a physically unclonable fingerprint of the specific sensor die.
This fingerprint is measurable from the \cb under
controlled illumination (for example a series of flat-field exposures
at different levels), and provides the basis for \TB
verification: a verifier can distinguish samples produced by a specific
sensor from samples produced by a different sensor or by software
simulation, because the spatial correlation structure, per-pixel
higher moments, and dark-current hot-pixel patterns are
device-specific. In some embodiments, because dark current and hot-pixel
patterns vary with temperature, exposure time, and sensor aging,
enrollment protocols include calibration at multiple operating points
and verification uses fuzzy extraction or helper-data algorithms to
accommodate drift while preserving uniqueness.

\paragraph{Regime mapping (non-limiting).}
The three regimes and Yoked modifier apply to the detector-as-reactor
mode. \RT in PoliePuter mode corresponds to generating
structured random patterns whose statistics are shaped by the projector
(the Create primitive, Section~\ref{par:create}): each frame is a
device-specific sample. \LI uses the noise-dominated readout for
inference: the projector pattern is chosen to maximise information gain
about a scene or data variable encoded in the illumination, and the
sensor's stochastic response is the measurement. \TB verifies
that the noisy readout was produced by a specific physical sensor
under a specific illumination protocol. Yoked operation couples the
sensor's noise-driven output back to the projector, creating a
closed-loop stochastic system whose joint dynamics sample from a
distribution that depends on both the sensor's noise profile and the
scene's response to structured illumination.

\paragraph{Denoising chain via temporal iteration (non-limiting).}
In some embodiments, the projector--sensor loop is iterated: the
readout from frame $t$ conditions the projector pattern for frame
$t+1$, implementing a temporal denoising chain within a single
projector--camera pair. Each iteration is one sampling step whose
conditional distribution depends on the previous sample (through the
projector update) and the sensor's noise (through the readout). This
maps to the multi-reactor pipeline
(Section~\ref{par:PoliePuter}) unfolded in time rather than in space,
with each temporal frame as one reactor stage. The \cb
logs every frame, providing a complete record of the denoising
trajectory.



\paragraph{Scope limitation (non-limiting).}
The embodiments of this section support declared sensing, attestation,
and hardness tasks benchmarked under the filing's task procedures. They
do not support, and shall not be construed to support, capacity or
scaling claims of the kinds bounded by
Section~\ref{sec:capacity-accounting-rule}, and do not imply any
increase in capacity quantities.

\subsection{Emitter-as-reactor: stochastic sampling in the source
(non-limiting)}
\label{sec:emitter-reactor}

In some embodiments, the emitter itself --- for example a semiconductor
laser, VCSEL, superluminescent diode (SLD), LED, or micro-LED ---
serves as a stochastic sampling element, symmetric to the
detector-as-reactor mode (Section~\ref{sec:detector-reactor}). Every
emitter exhibits device-specific noise sources whose statistics depend
on fabrication:
\begin{itemize}[nosep]
  \item \textbf{Spontaneous emission noise:} quantum-mechanical
    spontaneous decay into the lasing or emission mode, setting a
    fundamental noise floor that is source-specific through cavity
    geometry, gain-medium composition, and facet reflectivities.
  \item \textbf{Mode-partition noise:} in multi-mode sources,
    stochastic competition between longitudinal or transverse modes
    produces intensity fluctuations whose partition statistics depend on
    the cavity's mode spectrum and gain curve.
  \item \textbf{Relative intensity noise (RIN):} intensity fluctuations
    above the shot-noise floor, arising from carrier-photon dynamics in
    semiconductor sources, with magnitude and spectral shape determined
    by the specific die's relaxation oscillation frequency, damping
    factor, and parasitic elements.
  \item \textbf{Turn-on jitter and threshold fluctuations:} when pulsed
    or modulated near threshold, the turn-on delay varies stochastically
    from pulse to pulse, with statistics set by the spontaneous emission
    factor and carrier lifetime of the specific device.
\end{itemize}

In some embodiments, the system deliberately operates the emitter in a
noise-dominated regime --- for example biased just below or at
threshold, pulsed with short current pulses near the lasing onset, or
operated in a chaotic regime via delayed optical feedback --- so that
the emitted optical field at each time step is primarily a sample from
the source's noise distribution rather than a deterministic output. The
controller programs the bias point, modulation waveform, or feedback
strength (the ``energy landscape''), and the emitter's physics performs
the sampling. The detector reads out the resulting optical field, and
the \cb records both the control parameters and the
observed emission, preserving the standard emission-observation format.

\paragraph{reactor-microstructure uniqueness from emitter noise (non-limiting).}
In some embodiments, the combination of RIN spectral shape,
mode-partition statistics, threshold characteristics, and turn-on
jitter constitutes a physically unclonable fingerprint of the specific
emitter die. This fingerprint is measurable from the \cb
under controlled modulation sequences (for example a series of
threshold-crossing pulses at different bias levels), and provides the
basis for \TB verification: a verifier can distinguish the
noise signature of a specific emitter from that of a different emitter
or a software simulation. In some embodiments, because emitter noise
characteristics (particularly RIN magnitude, threshold current, and
mode-partition ratios) vary with temperature, bias current noise from
drive electronics, optical feedback from external reflections, and
device aging, enrollment protocols include characterisation at multiple
operating points and controlled bias conditions, and verification uses
robust feature extraction (for example phase-space reconstructions or
spectral entropy measures) together with fuzzy extraction or
helper-data algorithms to accommodate drift while preserving
uniqueness. In some embodiments, the measurement path is specified as
the integrated emitter-plus-readout chain (emitter die, coupling
optics, and detector bandwidth) to prevent trivial counterarguments
that the fingerprint reflects the measurement setup rather than the
emitter.

\paragraph{Regime mapping and symmetry with detector-as-reactor
(non-limiting).}
The regime mapping is symmetric to the detector-as-reactor case. In
PoliePuter / Create mode, each emission event is a device-specific
noise sample shaped by the programmed bias. In \LI mode, the
emitter's noise-driven output illuminates a scene, and the detector
reads back the scene's modulation of that stochastic illumination,
providing information about both the scene and the source. In \TB mode, the emitter's noise fingerprint is verified. In Yoked
operation, the detector's readout feeds back to the emitter's bias
point, creating a closed-loop stochastic system in which both emitter
noise and detector noise contribute device-specific signatures.

\paragraph{Combined emitter-and-detector stochastic mode
(non-limiting).}
In some embodiments, both the emitter and the detector operate
simultaneously in their noise-dominated regimes, so that the
\cb records the convolution of two independent
device-specific noise processes --- one at the source and one at the
sensor. The combined noise signature is richer than either alone and
provides reactor-microstructure diversity from two independent physical mechanisms
within the same optical path. In some embodiments, the two noise
contributions are partially separable by modulating the emitter at
known frequencies and examining the detector readout at those
frequencies versus at other frequencies, enabling independent
characterisation of source and detector fingerprints from a single
\cba record.

\subsection{All-optical stochastic sampling reactor embodiments
(non-limiting)}

In some embodiments, the reactor is implemented as an array of coupled
optical elements --- for example degenerate optical parametric
oscillators (DOPOs), bistable optical cavities with Kerr nonlinearity,
coupled semiconductor laser arrays biased near threshold, or arrays of
microring resonators with optical feedback --- whose quantum vacuum
fluctuations and spontaneous emission noise drive stochastic sampling
from a programmable distribution.

\paragraph{Optical parametric oscillator arrays (non-limiting).}
In some embodiments, an array of DOPOs is implemented in a fibre ring,
a free-space cavity, or a photonic integrated circuit. Each DOPO is
pumped above threshold so that it selects one of two degenerate phase
states ($0$ or $\pi$ relative to the pump), with quantum vacuum
fluctuations determining which phase is selected on each pump pulse.
Programmable optical or electronic coupling between DOPOs (for example
through delay lines, beam splitters, spatial light modulators, or
measurement-feedback circuits) implements an interaction graph, so that
the collective phase configuration of the array samples from an
energy-based model whose energy function is defined by the coupling
weights and local biases (pump detunings). In some embodiments, the
sampling distribution corresponds to an Ising model
$E(\sigma) = -\sum_{ij} J_{ij}\,\sigma_i\,\sigma_j
- \sum_i h_i\,\sigma_i$
where $\sigma_i \in \{+1, -1\}$ maps to the two DOPO phase states,
$J_{ij}$ are the programmed couplings, and $h_i$ are the programmed
biases. In some embodiments, continuous-variable extensions beyond
binary phase states are implemented by operating DOPOs in a regime
where amplitude fluctuations are significant, or by using non-degenerate
parametric oscillators with richer phase-space structure.

\paragraph{Bistable optical cavities (non-limiting).}
In some embodiments, arrays of Kerr-nonlinear optical cavities exhibit
optical bistability: two stable output states for a given input,
separated by an unstable branch. Quantum noise (vacuum fluctuations and
spontaneous emission) drives stochastic switching between stable states
at rates determined by the barrier height (set by cavity detuning,
nonlinear coefficient, and input power). Programmable coupling between
cavities through waveguides, evanescent fields, or feedback paths
implements an interaction graph as in the DOPO case, and the collective
switching dynamics sample from the corresponding energy-based model.

\paragraph{Native quantum noise and security (non-limiting).}
In some embodiments, the all-optical stochastic reactor operates in a
regime where quantum vacuum fluctuations are a significant or dominant
contribution to the noise that drives sampling. In idealised single-mode
OPO/DOPO theory, vacuum fluctuations seed the phase selection; in
practical implementations, the effective noise budget may also include
contributions from pump amplitude and phase noise, technical noise in
measurement-feedback electronics (for CIM-style architectures),
intracavity loss fluctuations, and thermal phonon coupling to the
nonlinear medium. The degree to which quantum noise dominates depends
on the operating regime (proximity to threshold, cavity finesse, pump
stability, and temperature) and is a design parameter rather than an
automatic guarantee. In embodiments where quantum noise is the dominant
contribution, the sampling has consequences for both security and
sampling quality.
For security: the noise statistics depend on cavity geometry, mirror
reflectivities, nonlinear coefficients, and other fabrication-dependent
parameters, providing reactor-microstructure uniqueness. In the quantum-noise
regime, an adversary attempting to characterise these parameters faces
a measurement-budget constraint: fully characterising the cavity
response at quantum precision requires probe resources that scale
unfavourably with the number of modes and the desired fidelity. The
security advantage is operational and measurement-budget-limited,
strongest when verification protocols exploit quantum correlations
(for example photon-number statistics or quadrature correlations) that
require many samples to estimate precisely.
For sampling: quantum tunnelling through the barrier between bistable
states can occur on timescales inaccessible to classical thermal
activation, potentially improving mixing times for multi-modal
distributions.

\paragraph{\Cb and regime mapping (non-limiting).}
In some embodiments, the \cb records the phase or
intensity state of each optical element at each pump cycle (or
measurement window), together with the programmed coupling and bias
parameters, pump power, and timing information. The regime mapping
follows the detector-as-reactor pattern
(Section~\ref{sec:detector-reactor}): \RT in
PoliePuter mode corresponds to generating samples from a target
distribution (the Create primitive, Section~\ref{par:create}); \LI
in PoliePuter mode corresponds to inference from data encoded in the
coupling parameters; and \TB in PoliePuter mode corresponds to
verifying that samples were produced by a specific physical device
whose quantum noise statistics carry device-specific signatures.

\paragraph{Denoising chains and multi-stage architectures
(non-limiting).}
In some embodiments, multiple optical sampling stages are chained so
that the output of one stage conditions the input (coupling parameters
or bias fields) of the next, implementing a denoising or iterative
refinement architecture. This maps to the multi-reactor chain topology
(Section~\ref{par:PoliePuter}), with each optical stage as one reactor
in the pipeline. In some embodiments, stages operate at different
timescales (for example a fast parametric stage followed by a slower
cavity-relaxation stage), providing a temporal hierarchy analogous to
the fast-loop / slow-loop architecture of the anchor embodiment.

\paragraph{Hybrid optical-cavity and detector-reactor architectures
(non-limiting).}
In some embodiments, optical parametric oscillators or bistable cavities
provide fast, quantum-noise-driven sampling for a subset of variables,
while the detector-as-reactor mode
(Section~\ref{sec:detector-reactor}) provides spatially resolved,
pixel-level stochastic readout with device-specific noise
fingerprints. The optical cavities generate samples at GHz rates; the
sensor integrates and reads out at frame rates, providing a natural
temporal hierarchy analogous to the fast-loop / slow-loop
architecture of the anchor embodiment. The \cb logs both
the optical cavity state trajectories (if accessible through homodyne
or photon-counting measurements) and the sensor readout, and the
resulting noise statistics combine quantum vacuum fluctuations from
the cavities with shot noise and sensor-specific fixed-pattern noise
from the detector, providing reactor-microstructure diversity from multiple
independent physical sources within the optical path.

\subsection{Reconfigurable metasurface reactor and optic embodiments}

In some embodiments, one or more optical elements in the \RK
module are implemented as metasurfaces---planar arrays of
sub-wavelength nanostructures fabricated on a substrate by lithographic
processes compatible with semiconductor manufacturing. Metasurfaces may
serve as reactor layers (providing scattering, phase shaping, and
polarisation transformation with device-specific nanostructure
randomness), as emitter optics (beam shaping, steering, and structured
illumination generation), or as detector optics (polarimetric filtering,
spectral sorting, and computational sensing).

In some embodiments, metasurface elements are dynamically
reconfigurable through integration with phase-change materials,
liquid crystals, electrochemical actuation, or MEMS, providing
tuneable reactor behaviour under controller command. In some
embodiments, metasurface elements are static but device-specific,
with reactor-microstructure uniqueness arising from sub-wavelength fabrication
tolerances in nanostructure geometry, placement, and material
composition.  In a non-limiting example, a silicon-on-insulator
metasurface fabricated by electron-beam or deep-UV lithography provides
a spatially varying phase and amplitude transfer function at a design
wavelength; the fabrication-dependent deviations from the nominal
design (pillar width, height, sidewall angle, and placement jitter)
constitute a device-specific fingerprint readable in the far-field
diffraction pattern. In \TB mode, the objective is to verify
that a measured far-field response matches a stored device profile
under a sequence of illumination angles or wavelengths. In \LI
mode, the metasurface serves as a computational optic whose known
transfer function enables computational reconstruction of scene
properties from compressive measurements. In \RT mode,
a reconfigurable metasurface reshapes emitted wavefronts under
closed-loop feedback to achieve target scene illumination.

\subsection{Chromic and optically writable reactor layer embodiments}

In some embodiments, reactor layers include one or more chromic
materials whose optical properties change in response to light
(photochromic), electrical stimulus (electrochromic), temperature
(thermochromic), or mechanical stress (mechanochromic). Non-limiting
examples include spiropyran and diarylethene photochromic films,
conductive-polymer electrochromic films, and VO$_2$ thermochromic
layers.

In some embodiments, photochromic reactor layers are directly written
by the \RK's emission subsystem (for example by UV or
visible structured illumination), creating spatial patterns of optical
state that persist for controllable durations and that modulate
subsequent light passing through the layer. In some embodiments, this
provides a projector-writable optical memory that is read by the same
or a different wavelength band, with the written pattern decaying or
being erased under controlled illumination.

In some embodiments, electrochromic reactor layers are electrically
addressed by the controller, providing fast ($<$1~s) and high-contrast
($>$50\% transmittance change) optical switching with stability over
thousands of cycles. In some embodiments, electrochromic layers are
combined with photovoltaic elements to create self-powered tuneable
reactor layers.

\subsection{Emitter embodiments beyond lasers and conventional
projectors (non-limiting)}
\label{sec:emitter-embodiments-beyond}
\label{sec:emitter-catalogue}

In some embodiments, the emission subsystem includes one or more of the
following emitter types, each providing capabilities distinct from
conventional laser diodes, broadband sources, and digital projectors:

\paragraph{Micro-LED arrays.}
In some embodiments, the emitter is an individually addressable
micro-LED array with pixel pitches of $10\,\mu\mathrm{m}$ or less,
providing peak brightness exceeding $10^6$~nits under declared duty-
cycle and thermal envelopes, device-level emitter modulation in the
sub-nanosecond to nanosecond range under appropriate drive
electronics, and direct electrical control of each emitter element.
Micro-LED arrays enable structured illumination patterns to be
generated without a separate spatial light modulator. In some
embodiments, fast device-level emitter modulation is combined with
photodiode-class or transimpedance-amplifier-class detector loops
to support local high-bandwidth feedback, while system-level
closed-loop bandwidth, including detector readout, controller
update, and committed-bundle logging, operates at a slower rate
declared in the protocol digest. In some embodiments, micro-LED
arrays include multiple wavelength bands (for example RGB or UV/
visible) on the same substrate, enabling spectral multiplexing of
emission patterns.
In some embodiments, the plurality of individually addressable
micro-LED emitters is partitioned into kernel taps, with per-emitter
intensities and switching times defining trainable or adaptively-
controlled tap weights distributed across the array, and with per-
emitter wavelength assignments fixed by material composition or
binning so that the array provides a wavelength-partitioned set of
taps. The kernel-tap partition operates either alone or in parallel
with scanning, DMD, SLM, or other emitter modalities, and is
configured to interoperate with the channel-separation and side-
channel mitigation discipline of the present disclosure.

\paragraph{VCSEL arrays.}
In some embodiments, the emitter is a two-dimensional array of
vertical-cavity surface-emitting lasers (VCSELs), providing coherent,
individually addressable, polarisation-controlled emission from a
monolithic semiconductor chip at wafer-scale manufacturing cost.
VCSEL arrays may simultaneously serve as emitters and as spiking
reactor elements (see spiking photonic reactor embodiments above),
combining emission and nonlinear dynamics in a single component.

\paragraph{Microwave-driven continuously tunable mode-locked
semiconductor lasers.}
In some embodiments, the emitter is a monolithic mode-locked
semiconductor laser in which a microwave driving signal applied
along the laser cavity induces a spatiotemporal gain modulation,
generating intracavity mode-locked pulses with a continuously
tunable group velocity, so that the output provides a coherent
pulse train with a continuously tunable repetition rate and,
equivalently, a frequency comb with a continuously tunable mode
spacing, under electrical tuning control and without mechanical
cavity adjustment. In some embodiments, the laser is a terahertz
quantum cascade laser (THz QCL) and the demonstrated tuning range
of the repetition rate is approximately $4$--$16$~GHz; in other
embodiments, the same spatiotemporal-gain-modulation mechanism may
in principle be implemented in other suitable semiconductor laser
classes (for example mid-infrared quantum cascade lasers or
interband semiconductor lasers with suitable fast-gain components)
through modulation-optimised geometries. In some embodiments, the
microwave drive waveform is treated as a component of the emitter
control input $u(t)$ and is committed to the protocol digest
together with the resulting repetition-rate or comb-mode-spacing
schedule, the laser bias and temperature set points, and the
relevant device-geometry identifiers. In some embodiments, meters
record drive-signal amplitude and phase stability, comb-line
linewidth, repetition-rate lock status, and power stability across
a declared tuning range. In some embodiments, continuously tunable
comb emission from such an emitter is combined with the
wavelength-multiplexing and dispersive primitives of
Section~\ref{sec:spectral-primitives} to support, when paired with
a second comb or suitable local-oscillator/reference path,
dual-comb or multi-heterodyne spectroscopic interrogation of the
reactor medium or of the scene as a declared observation channel.
Supporting background includes Senica et al., ``Continuously
tunable coherent pulse generation in a semiconductor laser,''
\emph{Nature} (2026).

\paragraph{Digital micromirror devices (DMDs).}
In some embodiments, the emitter path includes a DMD providing binary
amplitude modulation at pattern rates exceeding $30\,\mathrm{kHz}$,
optionally in UV wavelength bands ($363$--$420\,\mathrm{nm}$) suitable
for driving fluorescent or photochromic reactor media. In some
embodiments, DMD-based structured illumination is combined with
camera-based detection in a single compact module.

\paragraph{Spatial light modulators (SLMs) of various types.}
In some embodiments, the emission or reactor path includes one or more
spatial light modulators including, without limitation: liquid-crystal
on silicon (LCOS) phase-only modulators, MEMS piston-motion modulators,
electro-optic modulators based on lithium niobate or organic
electro-optic materials, acousto-optic modulators, and magneto-optic
modulators with sub-micron pixel pitch. In some embodiments, the SLM
type is selected based on the required modulation speed, pixel count,
wavelength range, and whether amplitude, phase, or polarisation
modulation is needed.

\paragraph{Light-field, integral-imaging, and holographic-stereogram
projectors.}
In some embodiments, the emission subsystem is a light-field or
integral-imaging projector comprising a high-density emitter (for
example a DMD, micro-LED array, LCOS panel, or phase-only spatial
light modulator) with a microlens array bonded to or positioned
immediately in front of the emitter plane, such that each microlens
is fed by multiple emitter pixels and the emitted field is
controllable in both spatial $(x, y)$ and angular $(u, v)$
coordinates. A single emission configuration therefore specifies a
four-dimensional light field $E(x, y, u, v)$, subject to the emitter's
spatial and angular resolution limits and the MLA geometry. In some
embodiments, the projector is realised as a stacked or tensor
light-field display in which two or more layered spatial modulators
jointly encode the emitted angular distribution by multiplicative
composition. In some embodiments, the projector is realised as a
holographic-stereogram or computer-generated-hologram display in
which a phase spatial light modulator reconstructs an approximation of
the complex emitted field at the display plane. In some embodiments,
the emitted light-field configuration is recorded in the protocol
digest together with the MLA pitch, effective angular resolution, and
calibration state, and meters record angular uniformity, cross-talk
between sub-apertures, and emitter-to-MLA registration. In some
embodiments, a light-field projector is paired with a light-field
(plenoptic) detector of matched geometry
(Section~\ref{sec:detector-embodiments-beyond}) to realise a
four-dimensional input--output architecture
(Section~\ref{sec:lightfield-hybrid}), in which the reactor or scene
sits between two controllable four-dimensional fields and the \cb
records both the emitted and received angular distributions at each
time step.


% ======================================================================
%
%
% ======================================================================

\paragraph{Dual-wavelength excitation and readout (non-limiting).}
In some embodiments, the emitter operates at a first wavelength band
selected to excite fluorescence, phosphorescence, or other
wavelength-shifting processes in the reactor medium, and the detector
observes the reactor through a bandpass, long-pass, or notch filter
that passes the wavelength-shifted emission while rejecting the
excitation wavelength. In some embodiments, the excitation band is in
the near-ultraviolet (for example approximately 365--420\,nm) and the
readout band is in the visible (for example approximately
500--700\,nm), or the excitation is in the visible and the readout is
in the near-infrared. In some embodiments, the spectral separation
between excitation and readout enables the system to observe reactor
dynamics (for example phosphorescence decay, fluorescence intensity,
spectral shifts, or bleaching kinetics) that would be masked by
specularly reflected or scattered excitation light in a
single-wavelength configuration.

In some embodiments, a dual-wavelength embodiment further includes a
reference channel in which a second detector or a second spectral band
of the same detector observes the excitation wavelength reflected or
scattered by the reactor, providing a simultaneous measure of the
reactor's elastic scattering response alongside the inelastic
(wavelength-shifted) response. In some embodiments, the ratio or
difference between elastic and inelastic channels is used as a
meter input for calibration, drift monitoring, or reactor-state
estimation.

In some embodiments, the protocol digest records the excitation
wavelength, the emission filter specification, the detector spectral
response, and any reference-channel parameters. In some embodiments,
dual-wavelength operation is compatible with all three operating
regimes: in verification mode, the wavelength-shifted response
provides an additional layer of physical specificity that increases
empirical hardness against emulators operating at a single wavelength;
in sensing mode, fluorescence or phosphorescence dynamics carry
information about reactor state, temperature, chemical environment,
or medium composition; and in transformation mode, the excitation
beam may be used to write persistent changes (for example
photobleaching, photochromic switching, or phosphorescence loading)
into the reactor medium while the readout beam observes the result.


% ======================================================================

\subsection{Detector embodiments beyond conventional cameras
(non-limiting)}
\label{sec:detector-embodiments-beyond}

In some embodiments, the detection subsystem includes one or more of the
following detector types, each providing invariance profiles or temporal
characteristics distinct from conventional frame-based cameras:

\paragraph{Event cameras (dynamic vision sensors).}
In some embodiments, the detector is an event camera in which each
pixel independently and asynchronously reports brightness changes that
exceed a threshold, producing a stream of timestamped events rather
than synchronous frames. Event cameras provide microsecond temporal
resolution and dynamic range exceeding 120~dB, and inherently detect
changes in the scene's or reactor's optical state---precisely the
signal of interest in a feedback loop. In some embodiments, event
cameras are used as the primary detector in Yoked operation, where the
asynchronous event stream directly encodes the temporal structure of
bidirectional coupling dynamics. In some embodiments, event cameras
are combined with conventional frame-based cameras to provide both
high-speed change detection and spatial context.

\paragraph{Single-photon avalanche diode (SPAD) arrays.}
In some embodiments, the detector is an array of SPADs providing
photon-counting sensitivity, picosecond timing resolution, and dynamic
range exceeding 150~dB through weighted photon-counting or
time-correlated methods. In some embodiments, SPAD arrays with temporal
gating (gate widths of nanoseconds or less) are used to separate prompt
scatter from delayed luminescence in reactors containing fluorescent or
phosphorescent media, or to perform time-of-flight measurements within
the reactor path. SPAD arrays extend the quantum-sensitive coupling
modes described elsewhere in this specification.

\paragraph{Hyperspectral and multispectral cameras.}
In some embodiments, the detector is a snapshot hyperspectral camera
providing tens to hundreds of spectral bands per pixel at video rates.
In some embodiments, hyperspectral detection enables the system to
resolve spectral shifts in reactor emission (for example
phosphorescence wavelength shifts with temperature or excitation
history) and to construct spectrally resolved invariance profiles for
cross-channel verification and sensing.

\paragraph{Polarimetric cameras.}
In some embodiments, the detector includes a polarimetric sensor
capable of resolving linear and optionally circular polarisation state
of incident light, providing polarisation-resolved invariance profiles.
In some embodiments, polarimetric detection is combined with
polarisation-controlled emission to create polarisation-multiplexed
channels with distinct scene and reactor sensitivities.

\paragraph{Light-field (plenoptic) cameras.}
In some embodiments, the detector is a light-field or plenoptic
camera comprising a microlens array (MLA) bonded to or positioned
immediately in front of a two-dimensional image sensor, such that each
microlens spans multiple sensor pixels and pixels beneath a given
microlens sample distinct angular coordinates or ray directions within
that microlens's sub-aperture set. A single exposure therefore records
a sampled four-dimensional light-field representation
$L(x, y, u, v)$, where $(x, y)$ denotes the spatial coordinate on the
MLA plane and $(u, v)$ denotes the angular coordinate within each
microlens's sub-aperture set. In some embodiments, the plenoptic
observation is included as additional channels of $\mathbf{y}_t$, so
that the \cb records not only per-pixel intensity but per-pixel
angular distribution within a single exposure. In some embodiments,
the plenoptic observation supports depth estimation, refocusing, and
sub-aperture view synthesis without active time-of-flight
illumination, and provides a measured angular-consistency constraint
on scene emissions that planar printed images, ordinary flat displays
lacking calibrated angular multiplexing, planar projection surfaces,
and other effectively two-dimensional replay surfaces do not
reproduce under the declared pose, sub-aperture baseline, and
calibration state. In some embodiments, the MLA pitch, microlens focal length,
and sensor pixel pitch are recorded in the protocol digest, and meters
record sub-aperture alignment, vignetting across the MLA field, and
angular-resolution calibration indicators.

\paragraph{Wavefront sensors and digital holographic imaging sensors.}
In some embodiments, the detector is a wavefront sensor that resolves
local phase gradients or complex field structure of the incident beam,
including, without limitation, Shack--Hartmann sensors (MLA-over-sensor
with per-microlens spot-displacement readout), pyramid wavefront
sensors, curvature wavefront sensors, and phase-diversity sensors. In
some embodiments, the detector is a digital holographic imaging sensor
that records a complex-valued field on a two-dimensional array, for
example through off-axis reference-beam interference, phase-shifting
interferometry, or in-line holography with computational
reconstruction. In some embodiments, the resolved phase-gradient or
complex-field observation is included as additional channels of
$\mathbf{y}_t$, enabling phase-conjugation and adjoint-matching
primitives as verification constraints, and supporting sensing regimes
in which intensity-only measurements under-determine the reactor or
scene state. In some embodiments, a plenoptic, wavefront, or
holographic sensor is combined with conventional intensity, spectral,
or polarimetric channels so that the \cb carries simultaneously an
intensity image, an angular or phase-resolved representation, and
auxiliary channels with distinct invariance profiles.

\paragraph{Computational and in-sensor processing detectors.}
In some embodiments, the detector performs computation on the optical
signal before or during digitisation, for example through analogue
photonic matrix operations on the sensor, spectral inner products
computed by on-chip filter arrays, or optical reservoir computing
performed in the detector's own optical front-end. In some embodiments,
such detectors reduce the data rate and latency of the feedback loop
by extracting task-relevant features directly in the optical or
analogue electronic domain.

% ------------------------------------------------------------------
\subsection{Printed and fabricated operator embodiments (non-limiting)}
% ------------------------------------------------------------------

In some embodiments, at least a portion of the mapping from emitted
fields to measured responses is realised by a designed operator rather
than discovered primarily by end-to-end learning. Non-limiting examples
include diffractive phase plates, metalens-like stacks, holographic or
volume gratings, engineered scattering layers, ultrafast-laser-written
photoelastic micro-vortex arrays in thermoplastic polymer substrates,
and hybrid refractive--diffractive assemblies.

In some embodiments, an input field encoded by a projector,
microdisplay, spatial light modulator, or patterned emitter is
propagated through one or more printed phase layers or diffractive
layers and then measured on a detector surface. The overall effect may
approximate a known analytic mapping, such as a Hankel-type transform, a
diffusion kernel on a disc, or a Laplace--Beltrami propagation on a
curved manifold. The configuration parameters $\theta$ include geometric
and material parameters of the fabricated layers.

\paragraph{Fabricated four-dimensional (light-field) operators
(non-limiting).}
In some embodiments, a fabricated operator is configured to implement
a designed four-dimensional transform between an emitted light field
$E(x, y, u, v)$ and a captured light field $L(x', y', u', v')$.
Non-limiting examples include metasurface stacks configured as
angular--spatial couplers, dual-MLA sandwich elements with an
engineered intermediate phase layer, holographic volume gratings
configured as Bragg-selective angular multiplexers, and diffractive
optical neural network layers trained to approximate declared
four-dimensional transfer kernels. In some embodiments, the
configuration parameters $\theta$ include the MLA pitches,
intermediate-layer phase profile, volume-grating recording parameters,
and inter-layer spacings, all of which are recorded in the protocol
digest so that downstream reconstruction and verification are
conditioned on the fabricated four-dimensional geometry.

In some embodiments, a fabricated operator is used in a low-recurrence
(single-pass) configuration. In other embodiments, the same fabricated
operator is placed in a cavity or feedback loop, so that repeated passes
yield Markovian dynamics. These embodiments remain instances of the
\RK formalism.

\paragraph{Ultrafast-laser-written photoelastic dispersive operators
(non-limiting).}
In some embodiments, a designed dispersive operator is realised by
ultrafast-laser-induced photoelastic micro-vortex structures written
into a thermoplastic polymer substrate. In some embodiments the
substrate is polycarbonate; in other embodiments other thermoplastic
polymers are substituted, subject to applicable photoelastic-response
and fabrication requirements. In some embodiments, each such structure
produces a rich spectral response over a broad bandwidth (for example
approximately $400$--$1550$~nm), occupies a compact footprint (for
example approximately $10~\mu\mathrm{m} \times 10~\mu\mathrm{m}$), and
is substantially independent of viewing angle. In some embodiments,
arrays of such structures are registered directly on a CMOS image
sensor to form a per-pixel dispersive operator, and the write
trajectory, laser fluence schedule, polymer specification, and
post-write spectral characterisation are recorded in the protocol
digest as part of the operator's configuration parameters $\theta$. In
some embodiments, the device-specific microstructure induced by the
ultrafast-laser writing process is treated as a physically unclonable
contribution to the operator, such that two nominally identical write
programmes on two substrates yield distinguishable per-pixel spectral
signatures that can themselves participate in attestation and hardness
claims under the rules set out elsewhere in this specification.
Supporting background includes Zhang et al., ``Optical dispersion
using micro-vortices in thermoplastic polymers for integrated
microspectrometers,'' \emph{Nature Electronics} (2026).

\paragraph{Fabrication and use of microstructure-unclonability cassettes and authentication
tokens (non-limiting).}
In some embodiments, the system includes a fabrication loop in which a
beam engraves, modifies, or otherwise writes microstructure into a
substrate to form a reactor-microstructure-unclonability cassette (analogous to a physically unclonable function or PUF in the prior-art sense) or
authentication token. The engraving branch steers a higher-power beam
across a substrate under a programmed trajectory while measuring the
evolving response in situ. The resulting emission--observation history
forms part of a canonical verification signature. After fabrication, the
engraved substrate can be re-inserted into a \RK module as a
microstructure-unclonability cassette, and a lower-power probe protocol is executed and compared
against a stored reference.

\paragraph{Additive-manufactured volumetric microstructure-unclonability bodies (non-limiting).}
In some embodiments, a microstructure-unclonability cassette is produced by additive
manufacturing---for example a 3D-printed translucent body whose internal
microstructure is intentionally disordered or process-variant (stochastic
infill, controlled bubbles or scatterer distributions, layered exposure
variations, or embedded particulate inclusions). When probed by
structured illumination, such bodies act as volumetric scattering microstructure-unclonability media:
responses are highly instance-specific and difficult to reproduce without
the same fabrication microstructure.

\paragraph{Physical signing and watermarking (non-limiting).}
In some embodiments, an engraving loop is used to apply physical
signatures, serial numbers, or watermarks to objects. The engraved
structure is treated as part of the scene during subsequent verification
or \LI passes, providing additional microstructure that can increase
empirical distinctiveness.

\paragraph{Training engraving and probing policies (non-limiting).}
In some embodiments, the system trains not only a verifier for a given
microstructure-unclonability cassette but also an engraving policy and a probing policy intended
to increase empirical hardness. Alice selects engraving policies that
determine scan laws, dwell patterns, pass schedules, and power
envelopes; Bob selects probing patterns and meters to verify the
engraved item; and Eve attempts to counterfeit cassettes under a bounded
resource model.

\paragraph{Subtractive microfabrication (non-limiting).}
In some embodiments, the controlled beam loop is used as a subtractive
manufacturing tool to create channels, relief structures, alignment
features, or diffractive patterns on a workpiece. The system may
alternate lower-power probe exposures and higher-power modification
exposures while observing results in situ. The resulting fabricated
element may serve as a static analogue operator in later \RK
runs, or as a functional component in a larger reactor stack.

\subsection{Extended-scale and environmental reactor embodiments
(non-limiting)}
In some embodiments, the \RK framework extends beyond
laboratory-scale optical systems to encompass larger-scale physical
systems whose dynamics satisfy the Markov-kernel abstraction
$\mathsf{P}_\theta(\cdot \mid S, U_{0:T})$. The following
non-limiting embodiments illustrate breadth of applicability.

\paragraph{Geological and geophysical reactors (non-limiting).}
In some embodiments, a geological formation or geophysical system acts
as the reactor: seismic wave propagation through rock strata, tectonic
strain accumulation and release, magma chamber dynamics, or glacial ice
flow. The ``emission'' is a controlled or instrumented perturbation
(for example a seismic source, a borehole pressure pulse, or an
acoustic transducer), and the ``detection'' is a geophone array,
strainmeter, GPS, or distributed fibre-optic sensor. The geological
medium's internal structure (fracture networks, layering,
mineralogical variation) provides device-specific nonlinear transfer
functions suitable for all three operating regimes: \TB for
structural monitoring and tamper detection (detecting changes in the
medium's transfer function), \LI for subsurface characterisation
(using the medium's response to active probing), and \RT
for controlled modification (for example guided fracturing or
injection). Protocol digests and meter envelopes record source
parameters and sensor configurations.

\paragraph{Atmospheric and ionospheric reactors (non-limiting).}
In some embodiments, the atmosphere or ionosphere acts as the reactor.
Optical signals (laser, structured light, or natural illumination)
propagate through turbulent or scintillating media whose refractive
index varies in space and time. The medium's turbulence profile
provides a natural time-varying microstructure signature (device-specific, time-varying, hard to clone
exactly). In some embodiments, ionospheric scintillation of satellite
signals provides the ``scene'' and the observable is the \cb of transmitted and received signal structure. Atmospheric
embodiments include lidar systems where backscatter from aerosols,
clouds, or precipitation acts as the reactor response.

\paragraph{GNSS-based passive Reality Kernel embodiments (non-limiting).}
In some embodiments, a Global Navigation Satellite System (GNSS)
receiver operates as a passive \RK in which the satellite constellation provides the emission
protocol and the atmosphere, ionosphere, and local multipath
environment collectively act as the reactor medium.  The ``emission'' is
the known transmitted signal structure (spreading codes, navigation
messages, carrier frequencies, and timing) of one or more GNSS
constellations, and the ``detection'' is the receiver's correlation,
tracking, and measurement engine.  The \cb pairs the
known transmitted structure with the received observables (pseudoranges,
carrier phases, Doppler shifts, and signal-strength indicators) under a
shared timebase derived from the GNSS signals themselves.

In some embodiments, the ionospheric and tropospheric propagation
channel introduces signal delays, scintillation, and multipath that are
location-specific, time-varying, and difficult to clone exactly,
providing a natural physically unclonable channel analogous to
device-specific reactor behaviour in active embodiments.  Meters
evaluate the consistency of received signal characteristics across
multiple satellites, frequencies, and time windows; anomalies such as
unexpected delay profiles, carrier-phase discontinuities, or
signal-strength patterns inconsistent with the declared antenna
environment are flagged for review.

In some embodiments, the system operates in a \TB mode where the
objective is to verify that GNSS observables are consistent with a
declared location, time, and propagation environment, enabling
tamper-evident positioning.  In some embodiments, the system operates in
a \LI mode where the objective is to characterise the propagation
channel itself (for example ionospheric total electron content,
tropospheric wet delay, or local multipath signatures) using the known
satellite signals as structured probes.  Protocol digests record
receiver configuration, antenna characteristics, constellation
selection, and any corrections applied (for example differential
corrections or atmospheric models), and meter envelopes record
signal-quality and consistency indicators.

\paragraph{Fisher information in passive GNSS Truth Beam (non-limiting).}
In the passive GNSS embodiment, the controllable configuration
parameters $\theta$ include receiver-side variables such as antenna
orientation and gain pattern, constellation and signal selection
(which satellites and frequencies to track), tracking-loop bandwidth
and integration time, elevation mask angle, and multipath mitigation
settings. \TB operation in this context selects $\theta$ to
maximise the Fisher information of the observable residual
distribution with respect to the hypothesis that the receiver is at
the declared location and time---equivalently, to maximise the
statistical power to distinguish a genuine location from a spoofed
one. This is a constrained form of the general \TB objective
(maximise $I(\theta)$ about the scene-plus-device configuration) in
which the ``emitter'' is not under the operator's control but the
receiver's measurement strategy is.

In some embodiments, a \RK combines an active optical or RF
emission path with a passive GNSS reception path, and the controller
fuses evidence from both paths.  The GNSS path provides externally
anchored position and timing while the active path provides
device-controlled scene or reactor probing; cross-consistency between
the two paths strengthens verification by constraining the evidence to
be consistent with both the declared physical interaction and the
declared spatiotemporal location.

\paragraph{Aquatic, fluidic, and chemical-oscillator reactors (non-limiting).}
In some embodiments, a body of water, flowing fluid, or microfluidic
channel acts as the reactor. The fluid's turbulence, mixing,
particle suspensions, or chemical composition provides the nonlinear
transfer function. Non-limiting examples include: river or ocean
current monitoring where structured light interacts with the water
column; microfluidic devices where reagent mixing in a channel
produces concentration-dependent optical response; and underwater
acoustic systems where sound propagation through a water column with
varying temperature, salinity, and current provides the channel.
In some embodiments, the reactor medium comprises a chemical oscillator
such as a Belousov--Zhabotinsky reaction, a chlorite--iodide--malonic
acid oscillator, or other reaction-diffusion system whose spatiotemporal
dynamics (spiral waves, target patterns, turbulence) provide a
high-dimensional, history-dependent, and intrinsically analog nonlinear
transfer function. Structured light illumination perturbs the
photosensitive variant of the reaction, providing the ``emission''
channel; optical imaging of the reaction's colour or fluorescence
changes provides the ``detection'' channel. The reaction's attractor
structure, bifurcation diagram under illumination intensity, and
spatial pattern repertoire are medium-specific and difficult to clone
or simulate exactly.

\paragraph{Biological ecosystem reactor applications (non-limiting).}
In some embodiments, a living biological system acts as the reactor:
coral reef fluorescence under structured UV illumination,
bioluminescent organisms responding to mechanical or chemical
stimulation, mycelial networks whose electrical or chemical signalling
responds to environmental probes, plant canopy photosynthesis and
transpiration dynamics under controlled lighting, or microbial
communities whose metabolic activity modulates optical properties of
their medium. These systems provide intrinsically nonlinear,
history-dependent, and clone-resistant transfer functions suitable for
environmental monitoring (\LI), ecosystem health verification
(\TB), and bio-responsive installations (\RT).
Implementations at larger ecosystem scale may use embodiment-specific
instrumentation, calibration, safety, environmental, and regulatory
controls appropriate to the selected biological medium and deployment
environment. Such implementations are non-limiting applications of the
reactor-medium class and do not limit the smaller-scale biological,
chemical, optical, or instrumented-reactor embodiments otherwise
described in this disclosure.
These embodiments differ from organoid computing and biological
reservoir computing systems (e.g.\ DishBrain, FinalSpark neuroplatform)
in that the biological medium here functions as a reactor in a
cryptographically committed optical feedback loop---its optical
response statistics constitute the \cb---rather than as
an autonomous computational substrate optimised for task performance.

\paragraph{Social and collective-behaviour reactors (non-limiting, limited scope).}
In some embodiments, a group of interacting agents (human, robotic, or
mixed) acts as the reactor, with group dynamics providing the nonlinear
transfer function. The ``emission'' is a structured stimulus presented
to the group (visual display, auditory signal, information feed), and
the ``detection'' is measurement of the group's collective response
(voting patterns, movement trajectories, communication patterns,
physiological synchrony). The group's internal social structure, trust
topology, and communication network provide device-specific properties.
Yoked operation is natural: the group's response modifies the next
stimulus, creating a bidirectional coupling between the display system
and the social dynamics. \textbf{Scope note (instrumented measurement applications; non-limiting):}
In preferred embodiments, this embodiment uses physical measurement
equipment (sensors, displays, physiological monitors) providing the
emission and detection channels; human behaviour itself does not
constitute a ``reactor'' in the patent-law sense absent such
instrumentation.  Purely abstract social dynamics without physical
instrumentation performing nonlinear signal transformation, and
systems in which the ``reactor'' is an economic, political, or
informational process without a physical substrate generating
measurable nonlinear responses to controlled physical stimuli, are
not required and are not the focus of the embodiments described
herein.  This paragraph describes applications to human or mixed-agent
behaviour via physical instrumentation; deployment scale is selected
according to available instrumentation, safety, consent, calibration,
and governance requirements.

% ------------------------------------------------------------------
\subsection{Sweeping reactor embodiments (non-limiting)}
\label{sec:sweeping-reactor}
% ------------------------------------------------------------------

In some embodiments, the reactor is operated in a \emph{sweeping}
configuration in which at least one physical parameter
$\lambda$ of the reactor (here $\lambda$ denotes a generic swept
physical parameter such as temperature, bias voltage, or optical path
length; it is unrelated to the regime weights
$\lambda_{\mathrm{TB/L/RT}}$) is periodically modulated at a frequency
$f_{\mathrm{mod}}$ while the emitter--recorder subsystem captures the
resulting signal modulation:
\[
\lambda(t) = \lambda_0 + \delta\lambda \cdot
  \sin(2\pi f_{\mathrm{mod}}\, t).
\]
Phase-sensitive (lock-in) demodulation of the recorder output at
$f_{\mathrm{mod}}$ extracts the local derivative of the reactor's
physical transform with respect to the swept parameter:
\[
S(f_{\mathrm{mod}}) \;\propto\;
  \frac{\partial T}{\partial\lambda}\bigg|_{\lambda_0}
  \cdot \delta\lambda,
\]
where $T$ denotes the physical transform applied by the reactor.  By
slowly stepping $\lambda_0$ through a range of interest, the full
transfer function $\partial T / \partial\lambda\,(\lambda_0)$ is
mapped---providing a differential characterisation of the reactor's
physics that is inaccessible to static measurements.

The sweeping reactor principle provides three intrinsic advantages.
First, common-mode rejection: any contribution to the measured signal
that does not vary at $f_{\mathrm{mod}}$ is rejected by the
demodulation process, including static offsets, slowly drifting
environmental perturbations, and constant systematic biases.  Second,
access to derivative information: fine structure in the reactor's
transfer function (dispersive slopes, resonance lineshapes, mode
crossings) is resolved directly as a continuous function rather than
sampled at isolated points.  Third, orthogonality to existing static
measurements: a smooth geometry-dependent contribution to the
reactor's output that is degenerate with other systematics in a
fixed-geometry configuration is selectively measured by the
sweeping reactor while geometry-independent offsets are rejected.
A complementary approach---temporal modulation of reactor material
properties at THz frequencies
(Section~\ref{sec:temporal-modulation})---addresses a different
timescale regime; the two approaches are orthogonal and may be
combined.

\paragraph{Null-subtraction and adaptive sweep control (non-limiting).}
In a preferred implementation, the measurement operates in a
null-subtraction mode: a model prediction for the signal at
$f_{\mathrm{mod}}$ is subtracted from the demodulated output in real
time, so that the lock-in output represents the \emph{residual}
between the measured transfer function and the model prediction.
This concentrates the instrument's dynamic range on deviations from
the model.  In some embodiments, the DC operating point $\lambda_0$
is stepped under feedback control: if the residual is consistent with
noise, the system advances with minimal dwell time; if the residual
exhibits statistically significant structure, integration time is
increased at that operating point.  In some embodiments, the
modulation amplitude $\delta\lambda$ is adapted to maintain operation
in the perturbative regime near resonances and to maximise sensitivity
between resonances.  In an advanced implementation, the drive waveform
(which may include multiple frequency components, chirps, or shaped
pulses) is optimised by a machine learning system to maximise
sensitivity to residual structure, constituting a learned emission
protocol within the emitter--reactor--recorder architecture (analogous to learned scan
laws in other embodiments).

\paragraph{Residuals as device-specific signatures (non-limiting).}
From the perspective of the \RK framework, the residual
between a model prediction and the measured transfer function is
analogous to the device-specific microstructure signature that provides
hardness in other reactor embodiments.  A reactor whose transfer
function is perfectly predicted by a finite-dimensional model has
zero residual and provides no empirical hardness beyond the model.
A reactor whose transfer function includes contributions not captured
by the model---whether from manufacturing imperfections, surface
roughness, material inhomogeneity, or any other physical contribution
not captured by the finite-dimensional model---produces a nonzero residual that is device-specific,
difficult to predict or reproduce without access to the physical
reactor, and therefore contributes to the 2D hardness index
(Section~\ref{sec:security-theory}; see also
Section~\ref{sec:crypto-primitives} for microstructure identification protocols).
The sweeping reactor provides a
precision tool for measuring these residuals as continuous functions
of the swept parameter.  In some embodiments, the scalar sweeping
principle is extended to simultaneous multi-parameter sweeping with
the scene's return signal as the lock-in reference rather than an
internally declared fixed frequency; this generalisation is described
in Section~\ref{sec:full-theta-sweep}.

\paragraph{Plenoptic readout of sweeping-reactor demodulation
(non-limiting).}
In some embodiments, the recorder of a sweeping-reactor configuration
is a light-field (plenoptic) detector
(Section~\ref{sec:detector-embodiments-beyond}) rather than a
single-pixel photodetector or a conventional frame-based camera, and
lock-in demodulation at the modulation frequency $f_{\mathrm{mod}}$ is
performed independently on each sub-aperture channel of the plenoptic
observation. The demodulated output is therefore a four-dimensional
derivative field
\[
  \frac{\partial L(x, y, u, v)}{\partial \lambda}\bigg|_{\lambda_0}
\]
rather than a scalar or two-dimensional derivative, where $(x, y)$
denotes the spatial coordinate on the MLA plane and $(u, v)$ denotes
the angular sub-aperture coordinate.

In some embodiments, per-sub-aperture demodulation is used only
where the declared photon budget, sampling cadence, and sub-aperture
signal-to-noise ratio support the selected lock-in estimator;
otherwise sub-apertures are pooled or treated as lower-confidence
channels under the meter envelope.

In some embodiments, the plenoptic sweeping-readout configuration
provides three capabilities that are not available under single-pixel
or conventional two-dimensional readout:

\begin{enumerate}[nosep]
  \item \emph{Angular separation estimates for scene and reactor
    contributions.}
    Contributions to the demodulated signal whose origin is
    geometrically coupled to the scene (for example,
    parallax-dependent return from a three-dimensional scene) and
    contributions whose origin is intrinsic to the reactor (for
    example, cavity-mode shifts under boundary modulation, or reactor
    microstructure response) may exhibit different angular profiles
    under the declared optical model. In some embodiments, regression
    of the demodulated field against declared angular basis functions
    estimates these contributions as additional channels of the meter
    envelope, and the angular-basis residual records how well the
    declared separation model fits the committed evidence. This provides
    a measured separation mechanism orthogonal to the
    $f_{\mathrm{mod}}$-based rejection described above, rather than an
    assumption that plenoptic geometry automatically separates the
    contributions.
  \item \emph{Sub-aperture derivative redundancy.}
    Each sub-aperture yields a separately indexed estimate of the
    local derivative, or of the view-conditioned component of
    $\partial T / \partial \lambda$, at a displaced viewing angle.
    In some embodiments, the cross-sub-aperture consistency of these
    estimates is computed under declared meters and committed
    alongside the \cb as a within-device consistency statistic,
    complementary to fleet-level proof-of-discrepancy confirmation
    statistics defined elsewhere in this specification. In some
    embodiments, sub-aperture consistency below a declared threshold
    is logged as a degraded-confidence indicator rather than treated
    as operational failure; the consistency statistic is a measured
    quantity under the declared meter set, not an assumed property of
    the plenoptic geometry.
  \item \emph{Native readout for angular-parameter sweeps.}
    Where the swept parameter $\lambda$ itself governs angular
    structure---for example, MLA translation, MLA tilt, emitted
    sub-aperture selection, or the phase offset of an
    angular-multiplex element---a plenoptic recorder reads the
    derivative $\partial L(x, y, u, v) / \partial \lambda$ directly,
    whereas a single-pixel or angle-averaging recorder averages over
    the angular coordinate and loses the signal of interest. In some
    embodiments, a light-field projector and plenoptic recorder
    (Section~\ref{sec:lightfield-hybrid}) are used as a matched pair
    for such angular-parameter sweeps, with the emission-side angular
    configuration recorded in the protocol digest alongside the
    demodulated derivative field.
\end{enumerate}

In some embodiments, the protocol digest records the plenoptic
geometry (MLA pitch, effective angular resolution, sub-aperture
assignment) alongside the sweep schedule, and meters record
sub-aperture registration stability during the sweep, vignetting
profile, and angular-basis residual indicators. In some embodiments,
the plenoptic sweeping-readout configuration is compatible with the
null-subtraction and adaptive-sweep-control strategies described
above, with the null model extended to a four-dimensional prediction
and the adaptive controller responding to residual structure in any
sub-aperture channel.

\paragraph{Cavity-coupled charged-particle embodiment (non-limiting).}
In a non-limiting embodiment, the reactor is a cylindrical microwave
cavity of radius $R$ and variable axial length
$L(t) = L_0 + \delta L \cdot \sin(2\pi f_{\mathrm{mod}}\, t)$,
where the movable endwall is displaced by a piezoelectric actuator.
The emitter is a single electron (or small ensemble) confined in a
Penning trap within a static magnetic field $B_0$, undergoing cyclotron
motion at frequency $\omega_c = eB_0 / m_e$ and radiating microwave
emission at $\omega_c$.  The cavity's electromagnetic mode spectrum
modifies the cyclotron frequency by a calculable shift
$\Delta\omega_c(L)$ that depends on the cavity geometry.  The recorder
is a cryogenic heterodyne receiver that downconverts the cyclotron
radiation and applies digital lock-in demodulation at
$f_{\mathrm{mod}}$ to extract the derivative
$\mathrm{d}\Delta\omega_c / \mathrm{d}L$ at each operating point
$L_0$.

Key design features of this embodiment include:
\begin{itemize}[nosep]
  \item A non-contacting quarter-wave choke groove at the
    piston--cylinder interface, which presents a low-impedance RF
    boundary without sliding electrical contact, eliminating
    particulate generation and time-variable resistive losses that
    would produce spurious signals at $f_{\mathrm{mod}}$.
  \item Independent calibration of the piston displacement
    $\delta L(t)$ via optical interferometry and of the magnetic field
    at $f_{\mathrm{mod}}$ via an independent magnetometer, enabling
    separation of geometry-dependent and field-dependent contributions.
  \item Null tests with the charged particle removed, with the
    particle far from cavity mode resonances, and with different
    piston materials, to identify and bound spurious signals.
\end{itemize}
The experimental transfer function
$\{\,L_0,\; \mathrm{d}\Delta f_c / \mathrm{d}L\,(L_0)\,\}$ is
compared to the standard cavity quantum electrodynamics prediction.
Residuals are examined for systematic smooth trends (indicating
geometry-dependent contributions not captured by the mode sum),
mode-correlated deviations (indicating uncharacterised cavity
imperfections or other device-specific microstructure), or null residuals
consistent with noise (indicating that the reference cavity model
accounts for the measured response at the instrument's resolution,
and that device-specific signatures must be sought at finer
sensitivity or in different measurement channels).  From the \RK perspective, the
cavity's manufacturing tolerances, surface finish, and
mode-coupling imperfections constitute a device-specific signature
measured with high fractional sensitivity (in some embodiments,
sub-part-per-billion or finer, depending on cavity quality factor
and readout configuration).  Implementation of this embodiment uses
cryogenic infrastructure (for example a dilution refrigerator or
adiabatic demagnetisation refrigerator providing base temperatures
below 100~mK), ultra-high vacuum (below $10^{-10}$~Pa), a
superconducting solenoid for $B_0$, and single-electron detection
via image-current amplifiers or quantum-jump spectroscopy.  These
requirements are standard in precision measurement laboratories and
do not require novel infrastructure beyond the \RK
control and logging overlay.

\paragraph{Alternative sweeping reactor implementations (non-limiting).}
The sweeping reactor principle extends beyond the microwave-cavity
embodiment.  Non-limiting alternative implementations include:
\begin{itemize}[nosep]
  \item \emph{Optomechanically modulated optical cavity:} a
    Fabry--P\'erot cavity with one mirror on a piezoelectric or MEMS
    actuator, modulated at $f_{\mathrm{mod}}$ while an optical source
    illuminates the cavity and a photodetector records the transmitted
    or reflected signal.  Lock-in demodulation extracts the derivative
    of the cavity transfer function with respect to length.
  \item \emph{Acoustically modulated fluid channel:} a microfluidic
    or mesoscale channel with an acoustically driven wall, probed
    optically or acoustically.  The demodulated signal encodes the
    derivative of the fluid's response with respect to channel
    geometry.
  \item \emph{Electrically tuned metamaterial boundary:} a
    metasurface whose electromagnetic response is modulated by applied
    voltage (for example varactor-loaded elements) at
    $f_{\mathrm{mod}}$, forming one wall of a resonant structure.
    The tuning parameter is an applied electric field rather than a
    physical displacement, enabling modulation at higher frequencies
    (up to MHz) and eliminating mechanical fatigue.
  \item \emph{Thermally swept photonic crystal:} a photonic crystal
    whose band gap is thermally modulated while an optical source
    probes the transmission.  The derivative measurement maps the
    temperature sensitivity of the photonic band structure.
  \item \emph{Radiation-pressure gradient measurement in asymmetric
    cavities:}
    a microwave cavity of non-uniform cross-section (for example a
    truncated conical frustum, a stepped-diameter cylinder, or a
    cavity containing a movable dielectric insert) driven at a
    resonant mode and mounted on a high-sensitivity force transducer.
    One cavity dimension is modulated at $f_{\mathrm{mod}}$, and the
    force transducer output is demodulated to extract the derivative
    of the net radiation pressure with respect to the modulated
    parameter.  This embodiment does not assert that the apparatus
    produces anomalous forces; the measured signal is the
    geometry-dependent
    radiation-pressure derivative predicted by standard
    electrodynamics, used as a device-specific signature.
    The lock-in architecture rejects thermal radiation
    pressure, outgassing forces, electromagnetic interference, and
    Lorentz forces from stray fields, providing systematic rejection
    inaccessible to static force measurements in asymmetric cavities.
\end{itemize}

\paragraph{Multi-reactor sweeping configurations (non-limiting).}
In some embodiments, two or more sweeping reactors are operated
simultaneously, each modulated at a distinct frequency
$f_{\mathrm{mod},k}$ with a controlled phase relationship.
Cross-demodulation of reactor $j$'s output at reactor $k$'s
modulation frequency extracts the inter-reactor coupling transfer
function: how modulation of reactor $k$'s geometry affects the
physical response observed in reactor $j$.  By sweeping
separation, relative orientation, or the DC operating point of
either reactor, the full inter-reactor coupling is mapped as a
continuous function of geometric parameters.  Residuals between
measured coupling and the prediction from standard coupled-mode
theory constitute a two-reactor analogue of the single-reactor
device-specific signature, contributing to the 2D hardness index
(Section~\ref{sec:security-theory}).

Non-limiting spatial configurations include:
\begin{itemize}[nosep]
  \item \emph{Collinear:} two or more reactors aligned along a
    common axis, providing maximum axial coupling sensitivity and a
    baseline measurement of the coupling's distance dependence.
  \item \emph{Ring with progressive phase:} $N$ reactors equally
    spaced around a ring, each modulated with a phase offset of
    $2\pi/N$ from its neighbour, so that the combined
    field-energy distribution rotates around the ring.  This
    configuration sources circulating energy currents and separates
    coupling components with distinct angular symmetry.
  \item \emph{Nested concentric (Matryoshka):} a smaller reactor
    enclosed within a larger reactor, so that the inner reactor
    operates within the electromagnetic environment modified by the
    outer reactor's boundary conditions.  The inner reactor's
    response is sensitive to the outer reactor's geometry through
    any channel that mediates inter-boundary coupling
    (Section~\ref{sec:matryoshka-reactor}).
  \item \emph{Incommensurately spaced (Fibonacci):} $N$ reactors
    along a line at spacings proportional to successive Fibonacci
    numbers (or another maximally irrational ratio), providing
    logarithmically distributed length scales with minimal periodic
    aliasing.  This arrangement optimises extraction of the coupling's
    distance scaling law (power-law exponent) from a finite number
    of reactor pairs.
\end{itemize}
In some embodiments, the inter-reactor phase relationships are
optimised to maximise the area enclosed by the combined state
trajectory in the multi-reactor configuration space, which
determines the sensitivity to coupling effects that depend on
the cyclic sequence of deformations rather than instantaneous
geometry.  In some embodiments, sum and difference frequencies
$(f_{\mathrm{mod},j} \pm f_{\mathrm{mod},k})$ are also monitored
to detect nonlinear inter-reactor coupling.  These configurations
describe multiple reactors within a single \RK module or
collocated assembly; for multi-device networked operation using
protocol-mediated coupling, see Section~\ref{sec:networks}.

\paragraph{Reference experiment: Fisher information at a synchronisation
  transition (non-limiting).}
In some embodiments, a reference experiment measures Fisher information
as a function of coupling strength across a synchronisation transition
between two \RK modules, or between an \RK module and a cooperative
scene partner. In one non-limiting protocol, two \RK modules (or one
\RK module and a cooperative scene partner) are configured in Yoked
mode with a continuously adjustable coupling-strength parameter
$\alpha \in [0,1]$ of the kind defined in the Yoked-operation
discussion of the present disclosure. The parameter $\alpha$ is
swept from zero (uncoupled) through a candidate synchronisation
threshold $\alpha_c$ to full coupling ($\alpha = 1$) and back, in
small increments, with the sweep schedule recorded in the protocol
digest. At each value of $\alpha$, a sufficient \cb window is
recorded from both modules to estimate the Fisher information
matrix $\mathcal{I}(\alpha)$ from the empirical score function
under a declared statistical model; transfer entropy estimates in
both directions ($\mathrm{TE}_{1\to 2}(\alpha)$ and
$\mathrm{TE}_{2\to 1}(\alpha)$), conditional Lyapunov exponents,
and attractor-dimension estimates are computed from the same \cb
segment. The trace $\mathrm{tr}[\mathcal{I}(\alpha)]$ is plotted
as a function of $\alpha$ and committed to the protocol digest
together with all intermediate estimates, the estimator family used
for each quantity, and the declared statistical model. In some
embodiments, the predicted signature is a peak in Fisher information
at or near $\alpha \approx \alpha_c$, accompanied by a conditional
Lyapunov exponent approaching zero from negative values on the
synchronised side and a sharp increase in transfer entropy; for
continuous transitions, the peak width is predicted to scale with a
critical exponent related to the universality class of the
transition, while for discontinuous transitions, a sharp spike with
hysteresis between the sweep-up and sweep-down branches is
predicted. The reference experiment provides a non-limiting protocol
for characterising critical signatures of synchronisation transitions
on \RK platforms; it connects the precision-metrology framing of
the sweeping reactor (Section~\ref{sec:sweeping-reactor}) with the
Yoked-operation regime and provides a concrete embodiment in which
the apparatus's sensitivity to dynamical-coupling phenomena can be
characterised against declared theoretical predictions.

\paragraph{Nested reactor (Matryoshka) architecture (non-limiting).}
\label{sec:matryoshka-reactor}
In some embodiments, independently of the sweeping reactor
configuration, the reactor element comprises two or more
concentric or enclosed physical structures, each constituting a
reactor in its own right, arranged so that inner reactors operate
within the electromagnetic, acoustic, or thermal environment
established by outer reactors.  This nested (Matryoshka)
architecture is distinct from layered reactor stacks
(Section~\ref{sec:layered-stacks}) in which layers are traversed
sequentially by the signal: in the Matryoshka configuration, the
inner reactor is \emph{immersed in} the modified environment of
the outer reactor rather than receiving light that has passed
through it.

The physical significance is that the inner reactor's transfer
function depends on the boundary conditions imposed by the outer
reactor.  When the outer reactor's boundaries are modulated (as
in the sweeping reactor configuration), the inner reactor's
operating environment changes continuously.  The inner reactor's
response---measured by its own emitter--recorder subsystem---is
thus a probe of how the outer reactor's geometry affects the
physical environment within it.  In the electromagnetic case,
this includes the cavity mode spectrum and the local density of
states as modified by the outer boundaries.

Non-limiting Matryoshka embodiments include:
\begin{itemize}[nosep]
  \item A microwave cavity containing a smaller Penning trap cavity
    (the inner reactor probes the electromagnetic environment
    established by the outer cavity via cyclotron frequency shifts).
  \item A Fabry--P\'erot optical cavity enclosing a scattering
    medium reactor (the inner reactor's microstructure response depends on the
    intracavity field structure set by the outer cavity mirrors).
  \item An acoustic resonator containing a smaller optical reactor
    (the acoustic mode structure modulates the mechanical
    environment of the optical components, coupling acoustic and
    optical degrees of freedom).
  \item Concentric electromagnetic cavities at different frequency
    bands (for example an RF cavity containing a microwave cavity,
    or a microwave cavity containing an optical cavity), where the
    outer cavity modulates the environment at one frequency scale
    and the inner cavity probes at another.
\end{itemize}
In Matryoshka embodiments, each reactor's \cb
is recorded independently, and the cross-coupling between reactors
is extracted via the multi-frequency demodulation scheme described
above.  The nested architecture provides a natural realisation of
hierarchical reactor composition within the \RK
algebra (Section~\ref{sec:security-theory}).  In some embodiments,
compromise of the outer reactor (characterisation of its full
transfer function by an adversary) partially degrades the inner
reactor's effective hardness, because the adversary can predict the
inner reactor's operating environment.  However, the inner reactor
retains its own device-specific microstructure, which outer-reactor
compromise does not directly reveal.  Periodic reconfiguration of
the outer reactor's boundary conditions (analogous to key rotation)
invalidates the adversary's accumulated characterisation.  In
preferred embodiments, inner-reactor hardness is monitored
independently so that degradation from outer-reactor compromise is
detected through the fleet monitoring mechanisms described in
Section~\ref{sec:networks}.

% ------------------------------------------------------------------
\subsection{Scene-phased full-parameter sweeping and physical gradient
computation (non-limiting)}
\label{sec:full-theta-sweep}
% ------------------------------------------------------------------

\subsubsection{Motivation and relationship to scalar sweeping}

The sweeping reactor embodiments described in
Section~\ref{sec:sweeping-reactor} operate by modulating a single
scalar parameter $\lambda$ of the reactor at a fixed, externally
declared frequency $f_{\mathrm{mod}}$, extracting the scalar
derivative $\partial T / \partial\lambda$ at each operating point
$\lambda_0$ via lock-in demodulation.  This provides a differential
characterisation of the reactor's transfer function along one
dimension of configuration space at a time.

In some embodiments, this principle is extended in two coupled ways:
(i)~the modulation is applied simultaneously across multiple
components of the full configuration parameter vector
$\theta = (\theta_{\mathrm{hw}}, \theta_{\mathrm{sw}})$, and
(ii)~the lock-in reference signal is derived from the scene's own
return dynamics rather than from an externally declared fixed
frequency.  The resulting operating mode is termed
\emph{scene-phased full-parameter sweeping}.  It generalises scalar
sweeping in the same way that the Jacobian generalises the scalar
derivative: where scalar sweeping extracts $\partial T /
\partial\lambda$ along one direction in $\Theta$, scene-phased
full-parameter sweeping extracts the directional derivative of the
coupled transfer function along the direction the scene is actually
traversing in configuration space during operation.

\subsubsection{Mathematical formulation}

Let $\theta(t)$ denote a time-varying modulation of the full
parameter vector:
\[
  \theta(t) = \theta_0 + \delta\theta(t),
\]
where $\theta_0$ is the nominal operating point and $\delta\theta(t)$
is a small perturbation applied across all or a selected subset of
admissible parameters.  In the scalar sweeping embodiment,
$\delta\theta(t) = \delta\lambda \cdot \mathbf{e}_\lambda
\sin(2\pi f_{\mathrm{mod}} t)$ for a fixed unit vector
$\mathbf{e}_\lambda$ and a single declared frequency $f_{\mathrm{mod}}$.

In some embodiments of scene-phased full-parameter sweeping, the
modulation direction $\delta\theta(t)$ and the reference signal for
lock-in demodulation are both derived from the scene's return signal
$r(t)$:
\begin{enumerate}[nosep]
  \item A reference signal $\phi_{\mathrm{ref}}(t)$ is extracted from
    $r(t)$ by bandpass filtering around the scene's dominant dynamical
    frequency $f_{\mathrm{scene}}$, where $f_{\mathrm{scene}}$ is
    estimated by a declared spectral-peak or phase-tracking method
    (for example a spectral coherence maximiser, a PLL, or a power
    spectral density peak detector) applied to $r(t)$.  In some
    embodiments, the transfer entropy estimators described in
    Section~\ref{sec:yoked} are used as a band-selection gate to
    identify frequency ranges with strong scene-to-reactor coupling,
    with spectral-peak or phase-tracking within that band used to
    extract the scalar frequency estimate $f_{\mathrm{scene}}$.
  \item The modulation $\delta\theta(t)$ is generated with phase
    locked to $\phi_{\mathrm{ref}}(t)$, so that each component of
    $\theta$ oscillates at $f_{\mathrm{scene}}$ with a declared
    amplitude $\delta\theta_k$ and a declared phase offset $\psi_k$
    relative to the scene reference.
  \item Phase-sensitive demodulation of the \cb signal $C_{0:T}$ at
    $\phi_{\mathrm{ref}}(t)$ extracts the in-phase (real) and
    quadrature (imaginary) components of the scene-phased response.
\end{enumerate}
The demodulated in-phase output is proportional to the directional
derivative of the physical transform $T$ along the modulation
direction at the current operating point, to first order in
$\delta\theta$:
\[
  S_{\mathrm{IP}} \;=\;
  G_L \,\operatorname{Re}\!\bigl[
    \nabla_\theta T\!\big|_{\theta_0} \cdot \widetilde{\delta\theta}
  \bigr]
  \;+\; O\!\bigl(\|\delta\theta\|^2\bigr),
\]
where $G_L$ is the lock-in gain and $\widetilde{\delta\theta}$ is the
complex modulation amplitude.  This first-order expansion holds under
the assumptions that $\|\delta\theta\|$ is small relative to the
operating point and that $T$ is sufficiently smooth in a neighbourhood
of $\theta_0$; higher-order terms appear at harmonic and
intermodulation frequencies and can be monitored or suppressed by
appropriate modulation depth and bandwidth choice.
By cycling $\delta\theta$ through a spanning set of directions in
$\Theta$ (or by applying a simultaneous multi-tone modulation with
adequate spectral separation, one tone per parameter dimension), the
full Jacobian $\nabla_\theta T\!\big|_{\theta_0}$ is estimated as a
function of the coupled operating point.

\subsubsection{The lock-in reference is the scene}

A defining feature of this embodiment, distinguishing it from scalar
sweeping, is that the demodulation reference is the scene's own
dynamical signal rather than a fixed oscillator internal to the
device.  This has several non-limiting consequences.

First, common-mode rejection applies not only to signals that do not
vary at $f_{\mathrm{mod}}$ (as in scalar sweeping) but also to
signals that vary at $f_{\mathrm{mod}}$ but are not phase-coherent
with the scene: spurious device-internal oscillations at the
modulation frequency are rejected because they lack the scene's phase
relationship, substantially improving rejection of internal artefacts.

Second, the extracted gradient is conditioned on the scene's current
dynamical state.  Two scenes with identical static appearance but
different internal dynamics produce different demodulated outputs,
because the reference phase $\phi_{\mathrm{ref}}(t)$ encodes the
scene's dynamics.  The gradient measurement is scene-specific in the
same sense that the Arnold tongue is scene-specific: it reflects the
coupled transfer function at the current joint operating point, not
the device's transfer function in isolation.

Third, when the device is operating in Yoked mode
(Section~\ref{sec:yoked})---specifically when
$\widehat{\CLE} \in [-\varepsilon_{\mathrm{CLE}}, 0)$ and transfer
entropy balance is confirmed---the reference signal
$\phi_{\mathrm{ref}}(t)$ is coherent with the reactor's own dynamics
by construction, because the two are phase-locked to the same joint
attractor.  In this regime, the modulation $\delta\theta(t)$ is
automatically phase-aligned with the joint attractor's oscillation,
and the demodulated gradient reflects the sensitivity of the
\emph{coupled} transfer function to parameter variation rather than
the device's transfer function alone.

\subsubsection{Connection to Yoked training: physical gradient
computation}

In some embodiments, scene-phased full-parameter sweeping serves
directly as the physical mechanism for computing the gradient
$\nabla_\theta J$ required by the Yoked training update rule
(Section~\ref{par:full-theta-yoked}):
\[
  \theta \;\leftarrow\; \theta
  + \eta\,\mathcal{F}(\theta)^{-1}\,\nabla_\theta J.
\]
In some embodiments, $\nabla_\theta J_{\mathrm{train}}$ (where
$J_{\mathrm{train}}$ is the declared TE-proxy training objective; see
Section~\ref{par:full-theta-yoked}) is not estimated numerically
by finite differences or via a differentiable surrogate model, but
is estimated from the lock-in output of the scene-phased sweep.
The physical mechanism is: a small perturbation $\delta\theta(t)$
phase-locked to the scene induces a measurable first-order change in
the \cb statistics (specifically in the TE estimators computed over
declared measurement windows), and the demodulated response provides
a directional-derivative estimate of $J_{\mathrm{train}}$ along
$\delta\theta$, to first order in the modulation depth.  The
full-gradient estimate for the training update is assembled from these
directional-derivative measurements across a spanning set of
modulation directions; the outer-loop quantities $J_{\mathrm{task}}$
and the CLE penalty are updated on their respective slower
timescales as described in Section~\ref{par:full-theta-yoked}.
Cycling through a spanning set of modulation directions recovers the
gradient estimate as a physical measurement, without requiring a
differentiable surrogate model of the physical channel.

In some embodiments, this operating mode supplies the measured signal
from which the gradient estimate is extracted by the declared
demodulation and estimation procedure.  The amplification of
sensitivity near the entrainment phase transition---where coupling
sensitivity to parameter changes is largest, as described in
Section~\ref{sec:yoked}---means that the gradient signal is
strongest precisely at the operating point where the Yoked training
update is most informative.  For periodically forced limit-cycle
oscillators of the class disclosed in this section, and under the
declared phase-dynamics assumptions of Section~\ref{sec:yoked}, the
edge-of-synchronisation condition and the maximum-gradient condition
coincide at the Arnold tongue boundary; in some embodiments, this
coincidence is empirically observed for the disclosed reactor
classes and is not asserted as a universal property of all
oscillator or media classes.

In some embodiments, an empirical preconditioner matrix
$\widehat{\mathcal{F}}(\theta)$ is formed from the covariance of
demodulated gradient estimates across multiple scene-phased sweep
cycles.  Under a declared likelihood model where the demodulated
measurements provide an estimate of the score, this covariance
approximates the Fisher information matrix; otherwise it serves as a
Gauss--Newton-style preconditioner that adapts the update direction
to the observed gradient geometry.  In some embodiments, the
pseudoinverse or damped inverse of $\widehat{\mathcal{F}}(\theta)$
is computed numerically from these empirical estimates, with damping
applied near bifurcation boundaries where
$\widehat{\mathcal{F}}(\theta)$ is ill-conditioned.

\subsubsection{Parameter scope: what lives in \texorpdfstring{$\theta$}{theta}}

In some embodiments, the parameter vector $\theta$ subject to
scene-phased sweeping includes all or any subset of the following
non-limiting parameter classes:
\begin{description}[style=nextline, leftmargin=2em]
  \item[Emission amplitude modulation frequency and depth]
    The frequency and amplitude of the slow intensity envelope imposed
    on the emission.  This is the parameter class addressed by scalar
    sweeping in its most common form.
  \item[Emission frequency modulation (FM)]
    The centre frequency of the emission carrier, modulated at the
    scene reference phase.  In some embodiments, this corresponds to
    the frequency detuning parameter $\Delta f$ of the Arnold tongue
    (Section~\ref{sec:yoked}), making scene-phased FM sweeping a
    direct instrument for mapping the Arnold tongue boundary from
    within the coupled system.
  \item[Emission phase and polarisation]
    Phase offset and polarisation state of the emitted field, each
    swept in phase with the scene reference.  These parameters access
    degrees of freedom in the coupled transfer function that are
    orthogonal to amplitude and frequency modulation.
  \item[Scan law parameters]
    The timing, trajectory, and repetition rate of the spatial scan
    pattern.  Sweeping the scan repetition rate in phase with the
    scene extracts the coupling's dependence on spatial sampling rate,
    complementing the frequency-domain Arnold tongue map with a
    spatial-domain coupling map.
  \item[Feedback gain $\alpha$ (the yoking dial)]
    The bidirectional coupling strength parameter $\alpha \in [0,1]$.
    Sweeping $\alpha$ in phase with the scene reference directly
    measures the derivative of the coupling objective with respect
    to coupling depth, which is the primary control parameter for
    navigating toward and stabilising at the tongue boundary.
  \item[Reactor internal timescale parameters]
    Parameters governing the reactor's internal dynamics that are
    under operational control: in phosphor reactors, drive intensity
    and duty cycle (which affect effective decay time); in cavity
    reactors, mirror separation or finesse; in delay-line reactors,
    loop length.  Sweeping these in phase with the scene extracts the
    coupling's dependence on the reactor's own dynamical timescale,
    extending the Arnold tongue map from the (detuning, coupling
    strength) plane to the full multi-dimensional entrainment
    manifold.
  \item[Learned parameters $\theta_{\mathrm{sw}}$]
    Neural network weights, look-up table entries, or calibration
    offsets that are differentiably connected to the \cb statistics.
    In some embodiments, scene-phased perturbations of
    $\theta_{\mathrm{sw}}$ and physical demodulation of the resulting
    \cb changes implement a physical analogue of backpropagation,
    without requiring a differentiable surrogate model of the physical
    channel.
\end{description}


\subsubsection{Multi-frequency parallel gradient extraction (non-limiting)}
\label{sec:multi-frequency-parallel-gradient}

In some embodiments, multiple parameter components are modulated
simultaneously using a bank of modulation tones $\{f_k\}$ referenced to a
shared clock derived from the scene reference phase $\phi_{\mathrm{ref}}(t)$.
The number of gradient components that can be extracted in a single
measurement window is bounded above by the time-bandwidth product of the
measurement channel, and in practice is substantially smaller than this
bound because of minimum tone separation, per-tone SNR, modulation-depth
limits, and calibration stability across the measurement window.
Accordingly, this disclosure does not assert that gradient extraction is
constant-time in parameter count. The disclosed mechanism supports
extraction of selected gradient components, or low-dimensional
projections of $\nabla_\theta J$, in parallel, subject to the bounds
above.

\subsubsection{Convergence and meter gating}

In some embodiments, scene-phased full-parameter sweeping is gated
by the same convergence conditions used to declare Yoked operation
(Section~\ref{sec:yoked-convergence}):
$\widehat{\CLE} \in [-\varepsilon_{\mathrm{CLE}}, 0)$ for $N$
consecutive windows and
$|\TE_{S\to R} - \TE_{R\to S}| < \tau_{\TE}$.  Gradient
measurements taken outside this regime are flagged in the protocol
digest as pre-convergence sweeps; gradient measurements taken within
this regime are flagged as post-convergence and are suitable for use
in the Yoked training update rule.  In some embodiments,
pre-convergence sweeps are used to navigate toward the Yoked
operating regime (coarse Arnold tongue mapping) while
post-convergence sweeps are used for fine-grained gradient ascent
within the regime.

\subsubsection{Relationship to the sweeping reactor}

Scene-phased full-parameter sweeping is a strict generalisation of
the sweeping reactor (Section~\ref{sec:sweeping-reactor}).  Where
the sweeping reactor fixes: (i)~the modulation to a single parameter
$\lambda$; (ii)~the reference frequency to an externally declared
$f_{\mathrm{mod}}$; and (iii)~the extracted quantity to the scalar
derivative $\partial T / \partial\lambda\big|_{\lambda_0}$;
scene-phased full-parameter sweeping relaxes all three constraints:
(i)~modulation is applied across all of $\theta$; (ii)~the
reference is the scene's own return phase $\phi_{\mathrm{ref}}(t)$;
and (iii)~the extracted quantity is the full Jacobian
$\nabla_\theta T\big|_{\theta_0}$ projected along the scene's
dynamical direction.

The two approaches are complementary and may be combined: scalar
sweeping with a fixed reference provides device-specific signature
characterisation under controlled conditions; scene-phased
full-parameter sweeping provides in-situ coupled gradient computation
during Yoked operation.  In some embodiments, the same hardware
supports both modes under controller selection, differing only in
the source of the lock-in reference signal and the set of parameters
being modulated.


 In some embodiments, a plenoptic-readout front-end
(Section~\ref{sec:lightfield-hybrid}; see also
Section~\ref{sec:sweeping-reactor} for the scalar-sweep case) applies
to both scalar and scene-phased operation. In such embodiments, the
light-field-valued transfer map may be written as
\[
  T_\theta:\; E(x,y,u,v)\mapsto L_\theta(x',y',u',v'),
\]
so the scene-phased measurement estimates the Jacobian of this
transfer map in the observed light-field coordinates, equivalently
$\nabla_\theta L_\theta(x', y', u', v')\big|_{\theta_0}$ for the
declared emitted field and operating point, rather than only a scalar
derivative. Cross-sub-aperture consistency of the measured Jacobian
is recorded as a within-device consistency statistic under the
declared meter set and is not assumed from the plenoptic geometry.

\subsubsection{Summary}

Assertions elsewhere in this disclosure concerning ``full-$\theta$'' or
``whole-parameter'' gradient extraction shall be construed as the
extraction of the projection of $\nabla_\theta J$ onto the declared
tone-bank subspace, under the bounds of
Section~\ref{sec:multi-frequency-parallel-gradient}. No assertion of
constant-time gradient-per-parameter extraction is made or implied.

\subsection{Scope-preserving notes (non-limiting)}
% ------------------------------------------------------------------

The following non-limiting statements clarify scope and operational
interpretation: verification and security-related embodiments rely on
empirical hardness
metrics defined against declared adversary families and budgets.
Hardness indices and meter outputs are interpreted relative to chosen
model classes, divergences, and compute/query budgets, and may be
recalibrated as these choices change. Embodiments requiring scan
addressing, timing alignment, or proof-of-projection assume maintained
calibration and synchronisation within declared envelopes.
Quantum-sensitive observables or non-classical effects are optional
enhancements and are not required for core \TB, \LI, or
\RT operation. Learned meters may themselves be attacked;
robustness and recalibration are ongoing empirical concerns addressed by
training curricula, audits, and conservative deployment policies.

% ======================================================================

% ======================================================================
%
%
%
% ======================================================================

\subsection{Non-optical architecture-portability demonstrations}
\label{sec:non-optical}
% ------------------------------------------------------------------

The following non-optical examples illustrate how the commitment,
calibration, corroboration, and verification architecture disclosed in
the preceding sections can be mapped to non-optical reactor families.
These examples identify the component families, commitment plumbing,
and architectural correspondences for each non-optical domain. They
are architecture-portability demonstrations showing that the same
architectural structure can be mapped to non-optical physics domains;
they do not constitute
full make-and-use bench recipes at the same level of concrete
operating detail as the optical core embodiments above. Concrete
operating parameters, calibration procedures, and validation
protocols for each non-optical medium are embodiment-specific
engineering tasks within the skill of the art once the architectural
mapping is established.

\paragraph{Acoustic reactor example (non-limiting).}
In some embodiments, the reactor is an acoustic scattering
assembly comprising a speaker or ultrasonic phased array as emitter, a
disordered solid, fluid, or granular medium as reactor, and one or
more microphones or acoustic sensors as detectors. In some
embodiments, the controller executes a time-indexed acoustic
excitation protocol and produces a committed \cb
$C_{0:T}$ binding each emitted acoustic waveform to each measured
acoustic response under a shared time base, with the protocol digest
recording excitation frequency, amplitude envelope, phasing schedule,
and medium identity.

In some embodiments, the same meter families used for optical
embodiments --- verisimilitude meters, forgery-difficulty meters,
calibration-consistency meters, and timing-integrity meters --- are
computed from the acoustic bundle using meter heads trained on
committed acoustic evidence rather than optical evidence. In some
embodiments, the two-dimensional hardness index
$(k^*_{\mathrm{dig}},\, m^*_{\mathrm{ana}})$ is evaluated against a
declared acoustic-emulator attacker family. In some embodiments, the
acceptance and evaluation meter partition is maintained, and the
acoustic module supports Level~0 (static challenge-response) and
Level~1 (trajectory) verification under the same graduated
verification architecture described in
Section~\ref{sec:graduated-verification}.

In some embodiments, a fleet of acoustic modules performs
proof-of-discrepancy using the same method object, independent
physical re-execution, and quorum confirmation workflow described in
Section~\ref{sec:pod-protocol}. In some embodiments, the portable bilateral
reproducibility protocol described in
Section~\ref{sec:bpod} applies between an acoustic module and an
optical module using the same Event, Envelope, and Portable Record
structure, with the cross-physics normalisation procedure --- how
acoustic and optical observations are transformed into a common
confirmation space --- being an embodiment-specific engineering
parameter.

In some embodiments, these non-optical examples illustrate that the
physical channel formalism, commitment infrastructure, meter
partition, graduated verification, fleet corroboration, and bilateral
portability are architecturally compatible with physical media that
contribute memory, nonlinearity, resonant structure, stochasticity, or
device-specific microstructure to a committed emission-observation
channel. In some embodiments, whether a specific non-optical medium
supports the full verification and corroboration stack at the same
depth as the optical core is an empirical property of that medium,
its associated transducers, and the medium-specific calibration and
validation procedures applied.

\paragraph{Magnetic reactor example (non-limiting).}
In some embodiments, the same architectural mapping extends to the
magnetic domain. In some embodiments, the module comprises a magnetic
excitation source, a magnetic or ferromagnetic reactor medium whose
transfer function depends on composition, geometry, domain structure,
and hysteresis state, and one or more magnetic sensors (for example
Hall-effect sensors, fluxgate magnetometers, or magneto-optical
sensors). In some embodiments, the controller executes a time-indexed
magnetic excitation protocol and produces a committed \cb binding each excitation waveform to each measured magnetic
response. In some embodiments, hysteresis and Barkhausen noise in the
ferromagnetic medium provide device-specific, history-dependent
signatures analogous to those provided by scattering and nonlinearity
in optical reactors. In some embodiments, the same meter families,
graduated verification levels, fleet corroboration, and bilateral
portability protocols apply to the magnetic embodiment under the same
architectural mapping, with meter heads and calibration procedures
adapted to magnetic observables.

\paragraph{Extension to further non-optical reactor families
(non-limiting).}
In some embodiments, the architectural mapping demonstrated for
acoustic and magnetic domains extends to any reactor family that
contributes memory, nonlinearity, resonant structure, stochasticity,
or device-specific microstructure to a committed emission-observation
channel, including but not limited to radio-frequency, mechanical,
fluidic, chemical, biological, and multi-physics embodiments. In some
embodiments, the common requirement is that the medium's transfer
function is sufficiently complex and device-specific to support
verisimilitude and hardness metering under a declared attacker model.


% ======================================================================


\section{Cryptographic Overlays and Commitment Plumbing (Optional, Non-limiting)}
\label{sec:commitment-plumbing}
% ======================================================================

This section describes optional cryptographic overlays that may be used
to bind, commit, timestamp, selectively disclose, aggregate, and govern
\RK evidence. These overlays are not required to implement
the \RK operator and do not replace the physical grounding
mechanisms described elsewhere. Instead, they provide standard
attestation plumbing on top of the \cb and related summaries.

\subsection{Overview and role separation}

In various embodiments, cryptographic mechanisms are used for at least
one of: (i)~integrity of exported logs and summaries,
(ii)~authenticity of device outputs and identity assertions,
(iii)~freshness and replay resistance, (iv)~selective disclosure and
privacy-aware sharing, (v)~rate limiting of outputs so that
cryptographic material is physically bound to device operation, and
(vi)~multi-device attestation and quorum mechanisms in distributed
deployments.
The physical channel remains the source of empirical hardness for
forgery, while cryptography provides tamper evidence and protocol
structure.

\subsection{Hash-chain and one-way state update}
\label{sec:hash-chain-state}

In some embodiments, the system maintains a compact cryptographic state
that evolves one-way with each emission--observation step, binding the
run history into a tamper-evident chain. A chain state $\chi_t$
is updated each step using the \cba sample $c_t = (u(t),
\mathbf{y}_t)$ as defined in Section~\ref{sec:definitions}:
\[
  \chi_{t+1} = F\bigl(\chi_t,\; c_t,\; \mathrm{meta}_t\bigr),
\]
where $F$ may be instantiated as a cryptographic hash or sponge
construction, a MAC or PRF under a device-held key, a verifiable random
function (VRF), or a ZK-friendly algebraic hash, and
$\mathrm{meta}_t$ denotes auxiliary metadata (for example timestamps,
sensor identifiers, or calibration state). The terminal value
$\chi_T$ serves as a compact fingerprint of the run.

In some embodiments, the emission policy is driven directly by the
evolving chain state $\chi_t$, forming a closed loop:
$u(t+1) = g(\chi_{t+1}, \mathrm{ctx}_{t+1})$, causing future emission
patterns to be a one-way function of prior observations, yielding
replay resistance.

\paragraph{One-way function family agility (non-limiting).}
The system may support migration between one-way function families. For
example, if a hash family is deprecated, the controller may switch to a
new $F$ while preserving the physical layer. In preferred embodiments,
the choice of $F$ is versioned and included in committed metadata, so
verifiers can interpret tags consistently across device lifecycle.

\subsection{Batching, Merkle commitments, and time anchoring}
\label{sec:merkle-atoms}

In some embodiments, the run is partitioned into disclosure atoms and
each atom is hashed to produce per-atom digests. These are aggregated
using a Merkle tree, yielding a root that commits to the entire run
while enabling inclusion proofs for any chosen subset of atoms. Time
anchoring is applied to commitments to provide externally verifiable
freshness via public timestamping services, blockchain inclusion,
verifiable delay functions (VDFs), or other time-locked attestations.

In some embodiments, commitment roots are submitted to append-only
transparency logs that publish signed tree heads over time, providing
publicly auditable evidence that a device published commitments by a
given time and did not later equivocate.

\subsection{Record formats and ledger integration (non-limiting)}

In some embodiments, exported evidence is packaged as structured records
whose fields are bound by commitments and signatures. Non-limiting
record types include: enrolment records (device identity binding and
baseline profiles), run records (protocol digests, calibration
identifiers, meter summaries, and commitment references), PoP/tPoP
export records (compact exports referencing committed windows plus time
anchors and quorum signatures), rotation records (signed policy
updates with compatibility flags), and revocation/quarantine records
(signed lifecycle decisions with supporting meter evidence).

\subsection{Physically rate-limited cryptographic outputs}
\label{sec:rate-limited-outputs}

In some embodiments, the system produces cryptographic outputs whose
generation rate is constrained by physical measurement and control,
rather than only by computation. Rate-limited outputs are derived from
correlated measurements, meter-bounded variability, and commitment
schedules. Rate limiting is enforced by requiring that each output be
bound to a fresh time anchor, a fresh protocol digest, or a fresh
\cba window, so outputs are not readily generated
arbitrarily fast without reproducing the physical interaction. This also
limits the rate at which a device can contribute quorum evidence in
multi-device attestation deployments.

In some embodiments, the sensor cadence (for example the frame rate of the
observation camera and the sweep rate of the scan program) provides an upper
bound on the rate at which fresh \cba windows can be produced.
This bound is evaluated under declared attacker families that may know the
device architecture, declared profile, calibration summaries, and protocol
class, but do not have the ability to reproduce the corresponding physical
measurement event under the declared timing, coupling, and meter conditions.
Accordingly, the rate limit is not dependent solely on secrecy of the device
profile; it is grounded in the physical cost of producing fresh measurement
windows at the declared cadence.

\subsection{Secure aggregation, threshold attestation, and multi-device
computation}

In some embodiments, multiple devices participate in aggregation and
attestation protocols using threshold signatures (a quorum jointly
produces a signature so that no single device can unilaterally
authorise), secure multi-party computation (devices contribute private
measurements and jointly compute an aggregate without revealing
individual inputs), or homomorphic encryption (devices submit encrypted
meter scores that are aggregated in encrypted form and decrypted only
as a collective result). In some embodiments, outputs of these
protocols are treated as conditioning variables for subsequent \RK operation.

\subsection{Abort handling and physical failsafes}

In some embodiments, an abort transitions the device into a safe regime
by disabling emission, parking scan mirrors, reducing gain, and
recording an abort marker. Abort also triggers a bounded flush of
buffers and a commitment step for already-recorded windows, so partial
runs remain tamper-evident. Sealed reactor loops include interlocks and
watchdogs, and scene-facing emissions include eye-safety constraints.

\subsection{Selective disclosure and privacy-aware sharing}

In some embodiments, the system supports privacy-aware sharing by
committing to large records and selectively disclosing only the minimum
required windows, features, or summaries for an audit or verification
task. Disclosure policy may be role-dependent. Selective disclosure is
combined with two-seed opening to reduce anticipatory manipulation, and
with escalation policies that request additional openings only when
summary-level checks fail.

\paragraph{Post-quantum cryptographic primitives (optional).}
In some embodiments, cryptographic components are instantiated with
post-quantum primitives, for example hash-based signature schemes such
as XMSS or SPHINCS+, providing defence in depth for digital bindings
while the physical channel provides empirical hardness.

\paragraph{Decentralised identity and registry integration (optional).}
In some embodiments, device identity is integrated with decentralised
identity or registry systems, assigning each device an identifier that
resolves to a record containing public keys, an enrolment-profile
digest, policy version identifiers, and revocation status.

% ======================================================================

\subsection{Security under profile disclosure: cadence-bounded, scene-coupled, and network-anchored forgery cost}

In some embodiments, disclosure of the device profile does not by itself
defeat verification, because the relevant security cost is the cost of
producing fresh measurement windows under the declared protocol, timing,
and meter conditions. A first component is cadence-bounded freshness; a
second component arises in declared Yoked embodiments; and a third
component arises in declared network-anchored embodiments. Accordingly,
the relevant question is not whether $\theta$ is secret, but whether an
attacker can reproduce the declared fresh physical interaction under the
declared digital, physical, and latency budgets.

\section{Security Theory: Hardness Indices and Empirical Security}
\label{sec:hardness-index}
\label{sec:security-theory}
% ======================================================================

This section consolidates the mathematical foundations of empirical
security for \RK systems. Security is grounded in physical
hardness rather than computational complexity assumptions.

Security-related embodiments in this disclosure are stated as engineering
security claims grounded in committed physical evidence, declared attacker
classes, meter envelopes, and periodically recalibrated hardness
evaluations. Formal proof-system properties, zero-knowledge properties,
and reduction-based guarantees are neither required nor implied except
where expressly stated for a particular digital primitive.


\subsection{Two-dimensional hardness index}

In some embodiments, empirical hardness is characterised by a pair
$(k^*_{\mathrm{dig}}, m^*_{\mathrm{ana}})$, where
$k^*_{\mathrm{dig}}(\theta;\varepsilon)$ is the minimal model capacity
(parameter count, FLOPs, or training budget) required to drive the
balanced classification error $\mathrm{err}^*_k(\theta)$ below a
threshold $\varepsilon$ against frozen evaluation meters on held-out
challenge windows, and $m^*_{\mathrm{ana}}(\theta;\varepsilon, r)$ is the
minimal physical emulation scale required to reduce the analogue
distinguishability advantage $\mathrm{Adv}^{\mathrm{ana}}_{m, r}(\theta)$
below the same threshold $\varepsilon$.  Both quantities are formally
defined in Section~\ref{sec:info-theory} (digital hardness) and the
analogue hardness paragraph therein.  The pair captures both
computational and physical-access dimensions of forgery difficulty,
and both are conditioned on a declared threat model and meter envelope
rather than asserted unconditionally.  In some embodiments, temporal
modulation of reactor properties
(Section~\ref{sec:embodiments}) amplifies both dimensions: $k^*_{\mathrm{dig}}$
increases because an attacker model would need to capture near-optical-cycle
switching dynamics in addition to spatial structure, and $m^*_{\mathrm{ana}}$
increases because temporal probes add an independent challenge
dimension orthogonal to spatial patterns.  Where time-varying media
enable nonreciprocal responses, the attacker is additionally denied an
equivalent reverse-path measurement under the declared threat model,
further increasing the
effective $m^*_{\mathrm{ana}}$.

\paragraph{Hardness as a time-indexed, attacker-indexed quantity (non-limiting).}
Empirical hardness $(k^*_{\mathrm{dig}}, m^*_{\mathrm{ana}})$ is not a permanent property of
a device; it is a snapshot conditioned on the current state of
attacker modelling capability and the attacker's accumulated
physical-access budget.  As emulator architectures improve and as
challenge--response data accumulates through legitimate use or
leakage, the effective $(k^*_{\mathrm{dig}}, m^*_{\mathrm{ana}})$ for a given device and
meter family may decrease.  The physical device, by contrast, does
not improve autonomously.  Security is therefore a \emph{race
condition}: the device's hardness at time $t$ is meaningful only
relative to the attacker class available at time $t$.  In some
embodiments, this is managed by (i)~periodic re-evaluation of
fleet-derived hardness thresholds against updated emulator families,
(ii)~reactor refresh or reconfiguration that changes the
challenge--response mapping, (iii)~query throttling and selective
disclosure to limit the effective $m^*_{\mathrm{ana}}$ available to any single
observer, and (iv)~fleet-level monitoring for population-wide
hardness degradation trends.  Nothing in this specification should be
read as asserting that empirical hardness is time-invariant or
unconditionally durable.  A novel contribution is the management
framework itself---periodic re-evaluation, reactor refresh, query
throttling, and fleet-level monitoring---which provides systematic
tools for maintaining useful hardness over time, not an assertion that
any particular hardness value is permanent.

\paragraph{Query-complexity bound for the digital hardness index
(non-limiting).}
The query throttle rate required to maintain the hardness index at a
declared target $k^*_{\mathrm{thresh}}$ depends on the rate at which
adversary model capacity can be effectively increased per additional
query to the target device.  In some embodiments, this query-complexity
curve $k^*_{\mathrm{dig}}(\theta; \varepsilon; Q)$ --- the digital
hardness index as a function of adversary query budget $Q$ --- is
estimated empirically for the declared reactor type and adversary model
family by training models of increasing capacity $k$ on CRP datasets of
increasing size $Q$ and recording the capacity at which balanced
classification error drops below $\varepsilon$.  Because $k^*_{\mathrm{dig}}$ decreases with $Q$ as additional queries
provide larger training sets to the adversary, the slope $\partial
k^*_{\mathrm{dig}} / \partial Q$ is non-positive, and the query
throttle threshold $Q^*$ is set as the $Q$ at which the magnitude of
the slope $|\partial k^*_{\mathrm{dig}} / \partial Q|$ exceeds a
policy-declared sensitivity $\Delta_{\mathrm{policy}} > 0$; beyond
$Q^*$, additional queries meaningfully accelerate an adversary's
model improvement and the throttle engages.  In some embodiments, the query-complexity curve
and the derived throttle threshold are committed to the protocol digest
and re-estimated at each declared re-evaluation cycle.

\paragraph{Average-case summary use of the hardness index pair
(non-limiting).}
The hardness index pair $(k^*_{\mathrm{dig}}, m^*_{\mathrm{ana}})$ is
used as an average-case summary of security under declared regularity
conditions: the challenge protocol draws challenges uniformly from a
declared set $\mathcal{X}_{\mathrm{prot}}$, and $\delta_{\mathrm{irr}}(X, i)$
has bounded variance across $\mathcal{X}_{\mathrm{prot}}$. Under these
conditions, $k^*_{\mathrm{dig}}$ characterises the adversary's
average-case difficulty under the uniform challenge distribution and
the pair provides a useful scalar summary for the declared average
case. The scalar pair is not a sufficient characterisation for all
applications: as detailed in
Section~\ref{sec:hardness-curves-and-monotonicity}, two devices can
have identical scalar indices while having very different hardness
profiles, and full hardness curves are required for cross-device
comparison, throttle calibration, adaptive-attack analysis, and
applications sensitive to tail risk or non-uniform challenge
distributions. When the regularity condition does not hold --- that
is, when $\delta_{\mathrm{irr}}$ is concentrated on a small region of
the challenge space --- the hardness index is supplemented by a
concentration diagnostic: in some embodiments, the spatial variance
of $\delta_{\mathrm{irr}}$ across the security domain
$\mathcal{X}_{\mathrm{sec}}$ is committed to the protocol digest, and
concentrated distributions trigger a challenge-space diversity
protocol.

\paragraph{Race condition formalisation: Nakamoto analogy
(non-limiting).}
The security condition stated above --- hardness holds as long as
$k < k^*_{\mathrm{dig}}(\theta; \varepsilon)$ for the declared
device and threshold --- is analogous in structure to Nakamoto
consensus security, which holds as long as the adversary controls
fewer than 50\,\% of total hash power.  In both cases, security is a
quantity-vs-threshold race condition rather than a reduction-based
proof.  The disanalogy is also instructive: in Nakamoto consensus,
adversary power is measured as a fraction of total network power,
and the attack probability decreases geometrically in block depth $k$
at rate $(q/p)^k$.  In the \RK setting, adversary
capability is an absolute scale (model capacity, physical emulation
scale, query count) rather than a fraction of fleet capacity, and
security degrades in discrete steps when a new model class crosses
the threshold, rather than continuously.  In some embodiments, the
management tools for this race condition are: periodic re-evaluation
(Nakamoto-analogous to honest-miner majority monitoring), reactor
refresh (analogous to changing the hash function), query throttling
(analogous to proof-of-work difficulty adjustment), and fleet-level
hardness monitoring (analogous to network hash-rate tracking). The
PoD architecture is structured so that an adversary who controls only
a minority of the participating fleet, under declared assumptions of
device-population diversity, independence of manufacturing and
calibration provenance, Sybil-resistance of the participation set,
and integrity of the verifier-acceptance pipeline, faces a degradation
of attack power against committed PoD results that grows with the
number of independent re-executions; absolute corruption resistance
is not asserted, and the threat model under which this property is
claimed is committed to the protocol digest. Failure modes outside
the stated threat model --- including correlated manufacturing
defects, coordinated firmware compromise, model-update poisoning, and
verifier capture --- are treated as separate risks managed by the
fleet-governance and meter-partition mechanisms described elsewhere
in this disclosure.

\paragraph{Meter partition for hardness evaluation (non-limiting).}
In some embodiments, the meter vector $\mathbf{m}(t)$ is partitioned
into \emph{acceptance meters} $\mathbf{m}^{\mathrm{accept}}(t)$ (used
for admissibility gating during operation) and \emph{evaluation meters}
$\mathbf{m}^{\mathrm{eval}}(t)$ (used for hardness assessment,
adversarial testing, and frozen-meter discrimination).  In embodiments using the stated partition, hardness
indices are computed from evaluation meters; acceptance meters
are excluded from the set against which hardness is assessed under that partition.  This
partition is structured to avoid self-referential hardness definitions in which a
meter's inclusion in the acceptance set simultaneously determines and
is determined by the hardness index.  In some embodiments, evaluation
meters are held out from the optimization loop (including the drifting
field), so that the system is not expected to adapt to satisfy the metrics used to
judge it under the declared optimisation loop constraints.  The partition is declared in protocol digests and bound into
commitments so that auditors can verify which meters served which role.

\paragraph{Residual indirect gaming risk (non-limiting).}
The meter partition is designed to block the direct Goodhart optimisation pathway: under the stated partition, no
gradient flows from evaluation meters into the optimisation loop, and
trust-weight updates explicitly exclude evaluation meters
(Section~\ref{sec:networks}).  However, acceptance meters and
evaluation meters measure the same underlying physical channel, so
they are statistically correlated.  An agent or adaptation process
optimising acceptance meters will exert \emph{indirect} pressure on
evaluation meters through those correlations.  This is a residual
risk inherent in any holdout-evaluation design.  In some embodiments,
the risk is bounded by (i)~maintaining low correlation between
acceptance and evaluation meter sets (achieved by selecting meters
that capture complementary aspects of channel behaviour),
(ii)~periodic rotation of which meters serve which role, so that
indirect gaming pressure is reduced on a fixed evaluation set
over extended operation, and (iii)~fleet-level threshold adaptation,
which reactively adjusts evaluation benchmarks if population-level
metric drift is detected.  Rotation schedules and correlation
audits are logged in protocol digests.  Each protocol digest declares
the partition in force for its epoch; rotation means the subsequent
epoch's digest declares a new assignment, while commitments and audit
records made under prior partitions remain valid and interpretable
under the partition that was in force when they were created.
The partition is a
first-order defence; these mechanisms provide second-order
adaptivity against learned correlations.

\paragraph{Governance-Partition Isolation (non-limiting).}
\label{sec:governance-partition-isolation}

In some embodiments, the \RK apparatus comprises a
governance partition isolated from one or more agent behaviour
generation pathways. The governance partition is configured to receive,
derive, store, verify, or emit governance-side signals, including
without limitation meter outputs, floor-proximity flags, protocol
digest entries, committed-bundle summaries, calibration records,
attestation records, authority decisions, rescue-mode triggers,
selective-opening decisions, or audit tokens. The isolation property
means that an agent behaviour pathway cannot directly read, write,
erase, rewrite, suppress, or directly influence the governance
partition except through declared evidence interfaces committed to the
protocol digest.

In some embodiments, the governance partition receives declared inputs
from one or more apparatus components while remaining outside the
optimisation loop that updates agent policy, behaviour, reward,
planning, self-modelling, or actuation parameters. For example, an
agent may produce actions that are sensed by the apparatus and
therefore affect future meter values, but the agent does not possess
direct authority to modify the meter implementation, the governance
logs, the authority signatures, the committed thresholds, or the
governance-partition state. This separates an evidentiary causal
effect from a direct control path.

In some embodiments, governance-partition isolation is implemented by
one or more of: a separate processor, a separate die, a separate
package, a separate sensor path, a separate power domain, a hardware
security module, a trusted execution environment, a microcontroller,
a read-only measurement path, an append-only log, a one-way data diode,
an air-gapped or optically isolated channel, a separate virtual
machine, a separate process, a separate key domain, a separate signing
authority, a threshold authority, a hardware-rooted attestation path,
or a physically isolated memory substrate. The foregoing examples are
non-limiting; the load-bearing property is the absence of an undeclared
direct agent-side read, write, or influence path over the governance
partition.

In some embodiments, the governance partition stores or emits a
governance signal $g_t$ derived from a committed bundle, meter vector,
or protocol digest:
\[
    g_t = G_{\mathrm{gov}}(C_{0:t}, m(t), \Pi_t, A_t),
\]
where $A_t$ may include authority records, attestation records,
selective-opening authorisations, or other governance-side inputs.
The behaviour-generation pathway may receive only a declared,
possibly lossy, delayed, thresholded, or redacted projection
$\tilde g_t$ of $g_t$:
\[
    \tilde g_t = R_{\mathrm{decl}}(g_t, \Pi_t),
\]
where $R_{\mathrm{decl}}$ is a declared release function committed to
the protocol digest. These equations are illustrative only and do not
limit the governance partition to any particular computational
architecture or release function.

In some embodiments, the governance partition is used to support the
faithfulness of privacy-preserving audit. For example, a privacy layer
may release only a conclusion, a commitment, a proof, an accumulator
state, a selective-opening result, or a lossy channel output, while
the governance partition preserves the evidence that the released
signal was computed from the declared committed bundle and not
fabricated by the agent pathway.

\paragraph{Isolation-mechanism generalisation (non-limiting).}
Governance-partition isolation generalises to physical, logical,
cryptographic, temporal, procedural, optical, electrical, software,
firmware, or hybrid separation mechanisms. Physical examples include
separate dies, separate packages, one-way links, sealed modules, and
air gaps. Logical examples include separate processes, execution
contexts, virtual machines, containers, privilege domains, and
enclaves. Cryptographic examples include separate key domains,
threshold signing authorities, attested measurement roots, signed
release policies, and append-only authenticated logs. Temporal
examples include committed past states that cannot be modified from
runtime behaviour pathways. No particular isolation technology is
essential unless expressly recited in a claim.

\subsection{Physical rate-limiting}

In some embodiments, the sensor cadence provides an empirical upper bound
on new-window production rate. Combined with commitment binding, this
yields a physically rate-limited token stream: a device can produce at
most $n$ independently verified tokens per unit time, where $n$ is
determined by sensor cadence and protocol overhead. Under declared
attacker families without the ability to reproduce the corresponding physical measurement event under the declared conditions,
exceeding this rate is empirically infeasible.

\subsection{Fleet latent space calibration}
\label{sec:fleet-model}

In some embodiments, a fleet of \RKs jointly calibrates a
shared latent space by exchanging committed reference sweeps under
standardised protocols. Fleet-calibrated latent spaces enable
cross-device comparison of meter outputs, hardness indices, and
anomaly thresholds. In some embodiments, fleet calibration uses
secure aggregation so that individual device responses are not
disclosed to the fleet, only population statistics. In some
embodiments, a fleet-derived hardness threshold is computed as a
quantile of the cross-device hardness distribution, providing a
data-driven acceptance boundary that adapts as the fleet evolves.
In some embodiments, the fleet-calibrated latent space provides the
device configuration representations used in the proof-of-discrepancy
protocol (Section~\ref{sec:proof-of-discrepancy}), and the irreducible
residual defined in Section~\ref{sec:error-decomposition} provides a
device-specific hardness contribution that complements the
population-level hardness index.

\subsection{Reality Kernel algebra and modular composition}

In some embodiments, \RK modules are treated as composable
operators with explicit input and output interfaces. A module may be
specified by (i)~an emission or actuation interface, (ii)~an
observation interface producing the \cb and meter summaries,
and (iii)~a policy interface defining admissible control protocols.

\paragraph{Modules as kernels with typed interfaces (non-limiting).}
For descriptive convenience, an RK module $K$ may be modelled as a
controlled Markov kernel
\[
  K_{\theta}(\cdot \mid x, u) \in \mathcal{P}(\mathcal{Y}),
\]
mapping an extended state~$x$ and a control~$u$ (including scan law,
projection schedule, gains, alignment, and other commands) to a
distribution over observable records or summaries in a
space~$\mathcal{Y}$. In some embodiments, $\mathcal{Y}$ includes
the \cb $C_{0:T}$, meter vectors, timing anchors, and
protocol digests. In some embodiments, interface compatibility is
enforced by policy-defined port types and meter-envelope constraints, so
composition is asserted only within a compatibility class.

\paragraph{Serial and parallel composition (non-limiting).}
If two modules $K_1$ and $K_2$ expose compatible interfaces, their
serial composition may be written informally as
\[
  K_{\mathrm{comp}} = K_2 \circ K_1,
\]
where outputs or effective fields of $K_1$ are treated as part of
the inputs available to~$K_2$ (for example through a shared scene,
a shared platform state, or direct reactor--reactor ports). In some
embodiments, modules are also composed in parallel, written
informally as $K_{\mathrm{par}} = K_1 \otimes K_2$, where both
modules operate over the same event window and their outputs are
combined by an aggregation rule.

\paragraph{Containment (nesting) composition (non-limiting).}
In some embodiments, a third composition mode---containment---is
distinguished from serial and parallel composition. In containment
composition, written informally as
$K_{\mathrm{nest}} = K_2 \triangleright K_1$, the outer module $K_1$
establishes boundary conditions (electromagnetic, acoustic, thermal,
or other physical constraints) that modify the operating environment
of the inner module $K_2$. Both modules produce independent
\cbs and meter summaries, but $K_2$'s transfer function
depends on $K_1$'s boundary state. Unlike serial composition, the
signal does not pass sequentially from $K_1$ to $K_2$; instead,
$K_2$ is \emph{immersed in} the environment defined by $K_1$.
Modulating $K_1$'s boundaries (as in sweeping reactor operation,
Section~\ref{sec:sweeping-reactor}) continuously changes $K_2$'s
operating environment, enabling cross-module coupling measurements.
Nested (Matryoshka) reactor embodiments
(Section~\ref{sec:matryoshka-reactor}) instantiate this composition
mode.  The algebraic properties of $\triangleright$ (for example
whether $(K_3 \triangleright K_2) \triangleright K_1$ is equivalent
to $K_3 \triangleright (K_2 \triangleright K_1)$, and how
$\triangleright$ interacts with serial and parallel composition) are
governed by the physical coupling topology of the specific
embodiment and are treated as empirical engineering targets rather
than formal identities, consistent with the algebraic-laws
discussion below.

\paragraph{Composition without shared secrets (non-limiting).}
In some embodiments, modules are composed without shared secrets by
binding evidence to public commitments and policy-logged protocol
digests. Each module contributes (i)~commitments to disclosure
atoms, (ii)~protocol digests $\Pi^{(i)}$, and (iii)~meter summaries,
and a higher-level composition operator combines these contributions
using quorum aggregation and selective opening. This supports
deployments in which modules can be assessed independently and later
combined using public logs, inclusion proofs, and policy-defined
trust weighting.

\paragraph{Summaries and auditability (non-limiting).}
In some embodiments, a summary map $\Sigma_{\mathrm{sum}}$ is applied to raw records
to reduce bandwidth, for example
\[
  \Sigma_{\mathrm{sum}}: (C_{0:T}, \Pi, \mathbf{m}) \mapsto \mathbf{s},
\]
where $\mathbf{s}$ is a compact summary. In some embodiments,
commitments and selective disclosure provide an auditable tie between summaries and raw atoms, so that a verifier can request openings
when summary-level checks fail or when escalation criteria are met.

\paragraph{Algebraic laws as engineering targets (non-limiting).}
In some embodiments, desirable algebraic properties are treated as
empirical engineering targets rather than exact identities. For
example, associativity and commutativity may hold only approximately
under meter-bounded operation, and are evaluated using the stability
metrics defined in the following subsection.

\subsection{Markov property metrics for engineered stability (optional)}

In some embodiments, \RK designs and protocol families are
engineered to exhibit approximate stability properties that make
serial composition and networked operation more predictable. Such
properties may be quantified by non-limiting empirical functionals
defined over an effective transition operator induced by a protocol.

For descriptive convenience, let $K_{u}$ denote the effective
one-step operator under a protocol choice~$u$ (which may bundle scan
law, projection schedule, gains, alignment, and other controls),
acting on an internal extended state~$X$ and inducing an observable
record (or feature summary) through $C_{0:T}$ and meters. Let
$d(\cdot,\cdot)$ denote any non-limiting discrepancy on induced
observables or summaries (for example, a distance on meter vectors,
a feature-space distance on the \cb, or a divergence between
induced distributions).

\paragraph{Kernel interpretation for the following functionals
  (non-limiting).}
For the purposes of the following functionals, $K_u$ is interpreted
as an effective transition operator on a declared summary or
representative state space (for example a sufficient statistic, an
embedding into a fixed Hilbert space, or a declared finite-
dimensional projection of the observable record), rather than as the
underlying law over observable records; the choice of summary or
representative is committed to the protocol digest, and the
functionals are evaluated on the chosen representative space under
a declared metric $d$.

\paragraph{Approximate idempotence.}
A non-limiting idempotence functional under a protocol~$u$ is
\[
  \nu_{\mathrm{idemp}}(u)
  = \E_{X\sim \mu}\!\left[d\!\left(K_{u}\!\circ K_{u}(X),\,
    K_{u}(X)\right)\right],
\]
where $\mu$ is a chosen reference distribution over operating states
(for example, a replay buffer over recent conditions).

\paragraph{Contractivity.}
A non-limiting contractivity functional may be defined as a
Lipschitz-like bound:
\[
  \nu_{\mathrm{contr}}(u)
  = \sup_{X\neq X'} \frac{d\!\left(K_{u}(X),\,
    K_{u}(X')\right)}{d(X,X')}.
\]

\paragraph{Mixing.}
A non-limiting mixing functional over a horizon~$n$ is
\[
  \nu_{\mathrm{mix}}(u;n)
  = \sup_{X,X'} \mathrm{TV}\!\left(\mathsf{P}_{\theta,u}^{(n)}
    (\cdot\mid X),\,\mathsf{P}_{\theta,u}^{(n)}
    (\cdot\mid X')\right),
\]
where $\mathsf{P}_{\theta,u}^{(n)}$ denotes the $n$-step induced law
under parameters~$\theta$ and protocol family~$u$, and $\mathrm{TV}$
is total variation distance (or another non-limiting divergence).

\paragraph{Approximate commutativity under serial composition.}
For two protocol choices $u$ and $v$, a non-limiting commutativity
functional is
\[
  \nu_{\mathrm{comm}}(u,v)
  = \E_{X\sim \mu}\!\left[d\!\left(K_{v}\!\circ K_{u}(X),\,
    K_{u}\!\circ K_{v}(X)\right)\right].
\]

\subsection{\texorpdfstring{$\varepsilon$}{epsilon}-nice kernels and tolerance specification}

In some embodiments, a \RK instance is characterised as
$\varepsilon$-nice under a family of protocols when the Markov
property functionals defined above are bounded by $\varepsilon$:
\[
  \mathcal{K}_\varepsilon
  =
  \Bigl\{\mathsf{P}_\theta:\;
  \nu_{\mathrm{idemp}} \le \varepsilon,\;
  \nu_{\mathrm{contr}} \le \varepsilon,\;
  \nu_{\mathrm{mix}} \le \varepsilon,\;
  \nu_{\mathrm{comm}} \le \varepsilon
  \Bigr\}.
\]
Instances in $\mathcal{K}_\varepsilon$ exhibit stable behaviour that
simplifies composition and network operation.  In practice,
$\varepsilon$ is understood as a vector of per-functional tolerances
$(\varepsilon_{\mathrm{idemp}}, \varepsilon_{\mathrm{contr}},
\varepsilon_{\mathrm{mix}}, \varepsilon_{\mathrm{comm}})$, each scaled
to the range and units of its respective functional; the single-symbol
notation is a shorthand for the requirement that all four tolerances
are simultaneously satisfied.

\subsection{Self-modification channel limitation (non-limiting)}
\label{sec:deception-symmetry}
\label{sec:unified-trust-root}

In some embodiments, the Inner Creation mechanism
(Section~\ref{sec:inner-creation}) permits a reactor to reshape its
own attractor landscape.  The same physical channel that enables
self-calibration could in principle be used to sculpt misleading
attractor patterns.  Physical constraints (microstructure
unclonability, thermal noise, bifurcation topology) raise the cost
of such misrepresentation, and in preferred embodiments Inner
Creation is rate-limited, protocol-restricted, and logged.
All layers of the evidence architecture are ultimately grounded in
the same physical channel; this is an inherent limitation of any
evidence system built on a shared physical substrate.  Extended
analysis of this limitation is described in the related Filing 2 application
(optional, non-essential).

% ======================================================================
\section{Fleet Model Improvement and Proof-of-Discrepancy Protocol}
\label{sec:proof-of-discrepancy}
% ======================================================================

This section describes methods (FIG.~5) by which a fleet of \RK
devices jointly improves a shared predictive model~310 of device behaviour,
using a distributed protocol in which participants discover~320 and submit~330
reproducible experimental methods that expose errors in the predictive
model, and in which, in preferred embodiments, verification of submitted
methods includes independent physical re-execution~340 on distinct devices.

\paragraph{Terminological note (non-limiting).}
\label{par:proof-of-discrepancy-empirical}
The term \emph{proof-of-discrepancy} is used throughout this document
as a term of art denoting an empirically calibrated fleet-confirmation
protocol, not a cryptographic proof system with formal soundness and
completeness properties.  What the protocol provides is:
(i)~committed evidence that a discoverer's device produced a
challenge--response measurement with an asserted residual against the
fleet predictive model, (ii)~independent physical re-execution on
one or more distinct unclonable devices in the verifier role,
scored against the fleet residual-margin estimator, and
(iii)~auditability via the submission record's commitment binding
and the verifier's confirmation record.  It does not provide, and
should not be read as asserting, a zero-knowledge or interactive
proof in the cryptographic sense.  Achieving formal proof-system
guarantees would require additional identity, Sybil-resistance, and
adversary-model assumptions beyond those specified here.  The
``discrepancy'' in the protocol name refers to the protocol residual
that fleet residual-margin calibration commits and that the verifier
physically reproduces; it is distinct from the discrepancy term in
Bayesian calibration in the manner of Kennedy and O'Hagan, which is
treated as a nuisance parameter rather than as a committed protocol
resource (see the novelty discussion at
Section~\ref{sec:fleet-model}).

% ----------------------------------------------------------------------
\subsection{Predictive model with device configuration representation}
\label{sec:device-config-rep}
% ----------------------------------------------------------------------

In some embodiments, a predictive model $\mathcal{Q}$ accepts as input a
challenge configuration $X$ (comprising at minimum a control protocol
or protocol seed specifying emission patterns, timing, gain settings,
and other controllable parameters) and a device configuration
representation $\mathbf{d}_i$ drawn from a fleet of device
representations~300 (a vector or structured object encoding
device-specific properties), and produces a predicted response:
\[
  \hat{Y} \;=\; \mathcal{Q}\!\bigl(X,\;\mathbf{d}_i\bigr).
\]
The device configuration representation $\mathbf{d}_i$ is a learned
or inferred quantity that captures how device~$i$ differs from a
population-level model of device behaviour.  In some embodiments,
$\mathbf{d}_i$ is a vector in a continuous latent space; in some
embodiments, $\mathbf{d}_i$ is a structured representation comprising
sub-vectors corresponding to distinct physical subsystems (for example
scattering medium properties, detector characteristics, and optical
path parameters).

In some embodiments, the predictive model $\mathcal{Q}$ is implemented as a
neural network, a Gaussian process emulator, a conditional generative
model, a physics-based simulator with adjustable parameters, or any
combination thereof.  The specification does not restrict the
architecture of $\mathcal{Q}$; what is required is that $\mathcal{Q}$ accepts both
a challenge specification and a device-specific representation and
produces a predicted response that can be compared with a physical
measurement.

\paragraph{Training the predictive model (non-limiting).}
In some embodiments, the predictive model is trained on
challenge--response data collected from a population of devices under
standardised protocols.  Training data comprises tuples
$(X_n, \mathbf{d}_{i_n}, Y_n)$ where $Y_n$ is the measured response
of device~$i_n$ to challenge~$X_n$.  In some embodiments,
$\mathbf{d}_i$ and the parameters of $\mathcal{Q}$ are jointly optimised from
fleet data.  In some embodiments, training uses a hierarchical
structure in which population-level parameters are shared and
device-level representations $\mathbf{d}_i$ are inferred per device.

\paragraph{Regularisation of device configuration representations
(non-limiting).}
Without regularisation, $\mathbf{d}_i$ can over- or under-fit the
available challenge--response data from device~$i$, inflating or
collapsing $\delta_{\mathrm{irr}}(X, i)$ and thereby mis-calibrating
the hardness index.  In some embodiments, a Gaussian process prior is
placed on $\mathbf{d}_i$: the prior mean is the population-level
representation $\mathbf{d}_0$ and the prior covariance is a declared
kernel (for example a radial basis function with length-scale and
output-scale hyperparameters committed to the protocol digest).  The
posterior over $\mathbf{d}_i$ given the calibration set then provides
shrinkage toward the population mean, with the degree of shrinkage
decreasing as more calibration data are collected.  In some
embodiments, an explicit bias-variance diagnostic is computed: the
posterior predictive variance on held-out challenges is committed to
the protocol digest, and a bias-variance alarm is triggered if the
variance exceeds a declared threshold (indicating under-fitting) or
if the in-sample residual drops below the noise floor (indicating
over-fitting).  These diagnostics are included in the hardness index
so that the security margin is bounded from below by the degree of
regularisation and from above by the noise floor, not by an
uncharacterised fit quality.

In some embodiments, the dimensionality of $\mathbf{d}_i$ is
explicitly bounded by a declared maximum $d_{\max}$, motivated by the
number of resolvable physical degrees of freedom of the reactor
(estimated from the scattering mean free path, beam diameter, and
camera resolution for optical embodiments).  Representations with
dimension exceeding $d_{\max}$ are not permitted; the bound is
committed to the protocol digest and is constant across re-evaluation
cycles unless a hardware change is declared.  In some embodiments,
the regularisation prior also imposes an $\ell_2$ norm constraint
$\|\mathbf{d}_i - \mathbf{d}_0\|_2 \leq r$ for a population-level
radius $r$ estimated empirically from fleet data; representations
outside this radius correspond to devices with non-representative
manufacturing history that the fleet model is not expected to cover.

In some embodiments, a sample-complexity bound characterises how
quickly $|\delta_{\mathrm{irr}}|$ shrinks as the adversary's
calibration set size $n_{\mathrm{adv}}$ grows: for a device
satisfying the learnability-exclusion condition, the expected
reduction in $\|\delta_{\mathrm{irr}}\|^2$ per additional CRP under
the declared adversary model is empirically estimated on the anchor
embodiment and committed to the protocol digest.  The query throttle
rate is set so that $n_{\mathrm{adv}}$ remains below the declared
threshold $n^*_{\mathrm{adv}}$ at which the shrinkage rate would
reduce the security margin below $\sigma^2_{\mathrm{sec}}$, a
policy-declared security floor.

\paragraph{Generalisation of $\mathcal{Q}$ to novel challenges
(non-limiting).}
The hardness claims for novel challenges --- challenges outside the
adversary's training distribution --- require not only that the
fleet model $\mathcal{Q}$ generalises, but that the irreducible
residual $\delta_{\mathrm{irr}}$ remains substantial on those
challenges.  In some embodiments, the generalisation of
$\delta_{\mathrm{irr}}$ to novel challenges is characterised using
the residual-margin estimator
(Section~\ref{sec:fleet-model}), which implements fleet
residual-margin calibration on the committed fleet calibration data.
In a Gaussian-process instantiation, the posterior predictive
variance at an unseen challenge $X_{\mathrm{new}}$ provides a lower
bound on the expected squared residual,
$\mathrm{Var}[\delta_{\mathrm{irr}}(X_{\mathrm{new}}, i) | \mathrm{data}] \geq \sigma^2_{\mathrm{post}}(X_{\mathrm{new}})$,
where $\sigma^2_{\mathrm{post}}$ is the posterior variance
computed from the fleet calibration data.  Other estimator classes
(neural posteriors, physics-based predictive distributions, or
hybrids) may instantiate the residual-margin estimator equivalently;
the estimator class is not load-bearing.  The estimator output is
treated as a committed protocol resource confirmed under the
Proof-of-Discrepancy protocol by independent physical re-execution on
distinct unclonable devices, not as a nuisance term to be marginalised.
This distinguishes fleet residual-margin calibration from Bayesian
calibration in the manner of Kennedy and O'Hagan, in which the
discrepancy term is typically estimated as a nuisance and integrated
out.  In a Gaussian-process instantiation, the predictive
uncertainty serves as the operationally used security margin: regions
of the challenge space where $\sigma^2_{\mathrm{post}}$ is high under
the GP prior (long-range challenge correlations not yet covered by
fleet data) are expected to be hard to forge.  In some embodiments,
the hardness claim on novel challenges is conditioned on challenges
at distance $> r_{\mathrm{adv}}$ from the adversary's training set in
the declared challenge-space metric, where $r_{\mathrm{adv}}$ is
bounded by the adversary's query budget and the density of the
challenge space under the declared protocol distribution; this bound
is reported in the protocol digest alongside the hardness index.

\paragraph{Fleet scaling and GP sample complexity (non-limiting).}
The GP posterior variance converges as fleet calibration data
accumulates.  For a $d$-dimensional challenge-device input space
$(X, \mathbf{d}_i)$ and a kernel with length scale $\ell$, the
posterior standard deviation at an unseen point at distance $r$ from
the nearest training point decays as $O(r^\nu)$ for a Mat\'{e}rn
kernel with smoothness $\nu$.  The effective input dimension is
bounded by $d_{\mathrm{max}}$, estimated from the number of
resolvable spatial modes $N_{\mathrm{modes}}$ and the declared
device-representation dimension; in practice $d_{\mathrm{max}}$
is the dominant sample-complexity driver.  No general closed-form
convergence theorem for heterogeneous fleets is claimed; the
operational approach is to use sparse or inducing-point GP
approximations (non-limiting) to maintain tractable inference as
fleet size grows, with the number of inducing points and the
approximation error committed to the protocol digest.  Fleet
scalability is therefore an empirical and implementation claim,
not a formal convergence claim of this specification.

\paragraph{Novelty of the fleet Proof-of-Discrepancy protocol relative to prior art (non-limiting clarification).}
The closest statistical prior art for fleet residual-margin
calibration is the Bayesian calibration framework of Kennedy and
O'Hagan (2001) and the Gaussian-process emulation literature that
follows from it.  Fleet residual-margin calibration is distinguishable
from that framework on three load-bearing points.
First, Bayesian calibration in the manner of Kennedy and O'Hagan
treats the discrepancy term as a nuisance to be estimated and
marginalised, with inferential interest in the calibrated parameters
of the computer model; fleet residual-margin calibration treats the
protocol residual as the committed resource of interest, with the
residual-margin estimator providing the operational security margin
that the Proof-of-Discrepancy protocol commits and physically
re-executes.  Second, Bayesian calibration in the manner of Kennedy
and O'Hagan does not involve physical re-execution on distinct
unclonable devices as a confirmation step; fleet residual-margin
calibration confirms the committed residual by independent physical
re-execution on distinct unclonable devices as its verification
primitive.  Third, Bayesian calibration in the manner of Kennedy and
O'Hagan operates at the level of a single computer model against
physical observations; fleet residual-margin calibration operates at
fleet level, with the residual decomposition separating shared model
discrepancy from device-specific variation and from the irreducible
clone-resistant residual maintained as a security primitive.
No cited prior calibration system,
authentication, or model-improvement system uses independent
physical re-execution on distinct unclonable devices as the
verification primitive for confirming a model discrepancy.  In
particular: (i)~Optical PUF
apparatus including that of Davis et al.\ (EP3252740B1) operates as a
single-device, single-function authentication system and does not
describe a fleet-level protocol in which multiple devices jointly
improve a shared predictive model or in which the irreducible residual
is deliberately maintained as a security margin. (ii)~Optical random
projection systems (including those of LightOn and related photonic
reservoir computing systems of Brunner and Larger) treat the optical
layer as a computational substrate for machine-learning inference; the
internal state of the optical layer is not committed as a tamper-evident
record and those systems do not support multi-regime operation under a
shared physical operator. (iii)~Chained optical transformation systems
such as those of Naughton, Hennelly, and Dowling use fixed optical
elements in a CFB-mode encryption scheme; they do not describe a closed
feedback loop with an evolving reactor state, a logged control protocol,
or a \cb in which each sample pairs emission parameters
with measured responses. The combination of (a)~physical re-execution
on independently manufactured devices as the confirmation primitive,
(b)~the five-term error decomposition separating shared model
discrepancy, device-specific variation, and irreducible clone-resistant
residual, and (c)~the protocol digest commitment binding the
confirmation record to the physical channel event is, to the best of
the inventors' knowledge, absent from the cited prior art and from any
combination thereof.

\paragraph{Physical Irreducibility Conjecture (non-limiting, research direction).}
The residual-margin estimator
(Section~\ref{sec:fleet-model}) provides an operationally used lower
confidence bound on $\delta_{\mathrm{irr}}$ under the declared
estimator class; in a Gaussian-process instantiation, this bound is
the GP posterior variance under the declared GP model class.  A
deeper and as yet
unproved question is whether a stronger, model-class-independent lower
bound exists.  The following is stated as a research conjecture
rather than a claim of the present disclosure.

\medskip
\noindent\textit{Physical Irreducibility Conjecture.}  Let a device
satisfy the following physical conditions: (i) the scattering medium
has a transport mean free path $\ell_s$ and wavelength $\lambda$ such
that $N_{\mathrm{modes}} \leq c_0 \Delta^2 / \lambda^2$ resolvable
spatial modes exist in the measurement aperture of diameter $\Delta$;
(ii) the device dynamics have at least one positive Lyapunov exponent
(sensitive dependence on initial conditions); (iii) observations are
corrupted by instrument noise $\sigma^2_{\mathrm{noise}} > 0$.  Then
there exists a function $\sigma^2_{\mathrm{phys}}(N_{\mathrm{modes}},
k, T) > 0$ depending only on these physical parameters such that,
for any learning procedure using a model class with at most $k$
parameters, the minimax prediction error on novel challenges satisfies
\[
  \min_{f \in \mathcal{H}_k}
  \mathbb{E}\!\left[\|Y_i(X) - f(X, \mathbf{d}_i)\|^2\right]
  \;\geq\;
  \sigma^2_{\mathrm{phys}}(N_{\mathrm{modes}},\, k,\, T)
  \;>\; 0
  \quad \text{for all } k < k^*(N_{\mathrm{modes}}, T, \sigma^2_{\mathrm{noise}}),
\]
and this lower bound is independent of the adversary model class and
holds for any learning procedure given $T$ observations.
\medskip

Even without a model-class-independent lower bound, the device retains a
measurable empirical irreducibility: the irreducible residual
$\delta_{\mathrm{irr}}$ is directly observable as the held-out discrepancy
surviving independent physical re-execution under declared model classes,
query budgets, and meter envelopes, and is reported as a bound at the
declared credibility level. Absent the conjecture, the bound remains an
operational measurement.

If true, this conjecture would constitute a bridge between three
bodies of results that are individually well-developed but not yet
connected: (i) Cramér--Rao bounds and channel capacity for the optical
scattering channel as a function of $N_{\mathrm{modes}}$ and $T$;
(ii) the Kaplan--Yorke conjecture relating Lyapunov spectrum to
attractor fractal dimension; and (iii) PAC learning lower bounds
converting channel capacity (information about $\theta$ extractable in
$T$ observations) into a model-complexity threshold $k^*$.  The
proof technique most likely to proceed: (a) establish an optical
channel converse bounding $I(\theta; C_{0:T})$ by
$N_{\mathrm{modes}} \cdot \log(1 + \mathrm{SNR})$; (b) apply
Fano's inequality to convert this to a minimax lower bound on any
estimator of $\theta$; (c) use the resulting bound on estimable
degrees of freedom to lower-bound prediction error for any
$k$-parameter model.  In some embodiments, this programme is pursued
empirically by reporting $\mathrm{err}^*_k(\theta; Q)$ as a function
of $Q$ (query count) and $k$ (model capacity) across a well-characterised
family of scattering tokens with known grain structure, testing the
prediction that error curves plateau at positive values scaling with
$\ell_s / \lambda$ and $N_{\mathrm{modes}}$.  The empirical plateau
(if observed) would constitute evidence in favour of the conjecture
and would ground the security claims of the present disclosure on a
physical law rather than on current model-class limitations alone.

\paragraph{Extension to unseen devices (non-limiting).}
In some embodiments, a device configuration representation
$\mathbf{d}_j$ for a previously unseen device~$j$ is inferred from a
calibration set of challenge--response pairs
$\{(X_n, Y_n)\}_{n=1}^{N_{\mathrm{cal}}}$ obtained by executing a
calibration protocol on device~$j$.  In some embodiments, inference
is performed by optimising $\mathbf{d}_j$ to minimise discrepancy
between $\mathcal{Q}(X_n, \mathbf{d}_j)$ and $Y_n$ for the calibration
challenges, with the parameters of $\mathcal{Q}$ held fixed.  In some
embodiments, inference is performed by an amortised encoder---a
learned function that maps a set of challenge--response pairs directly
to a device configuration representation without iterative
optimisation.

In some embodiments, the number of calibration challenges required to
infer a device configuration representation to a given accuracy is
determined empirically and depends on the dimensionality of
$\mathbf{d}$, the information content of each challenge--response
pair, and the complexity of device-to-device variation.  The
calibration set need not exhaustively cover the challenge space; it
need only contain sufficient information to locate the device in the
learned representation space.

\paragraph{Fleet geometry in representation space (non-limiting).}
In some embodiments, the set of inferred device configuration
representations $\{\mathbf{d}_i\}$ for a fleet of devices exhibits
structure reflecting manufacturing processes, material properties, and
fabrication conditions.  Devices from the same manufacturing batch or
process may cluster in representation space; devices from distinct
processes may occupy distinct regions.  In some embodiments, the
geometry of the representation space is used for fleet management,
including identification of device families, detection of anomalous
devices, allocation of computational or sensing tasks based on device
characteristics, and assessment of fleet diversity.

% ----------------------------------------------------------------------
\subsection{Five-term error decomposition}
\label{sec:error-decomposition}
% ----------------------------------------------------------------------

In some embodiments, a population-level representation
$\mathbf{d}_0$ is defined (for example as a learned population mean
in the representation space, or as the representation that minimises
average prediction error across the fleet).  The population-level
prediction and the device-specific prediction are then:
\[
  \hat{Y}_0(X) \;=\; \mathcal{Q}\!\bigl(X,\;\mathbf{d}_0\bigr),
  \qquad
  \hat{Y}_i(X) \;=\; \mathcal{Q}\!\bigl(X,\;\mathbf{d}_i\bigr).
\]
The per-device residual---the quantity directly observable after
inference of $\mathbf{d}_i$---is:
\[
  e_i(X) \;=\; Y_i(X) \;-\; \hat{Y}_i(X).
\]
In some embodiments, the measured response $Y_i(X)$ comprises a
noise-free physical response $Y_i^{\mathrm{true}}(X)$ and
measurement noise $\varepsilon(X, i)$, such that
$Y_i(X) = Y_i^{\mathrm{true}}(X) + \varepsilon(X, i)$.
In some embodiments, a noise-reduced residual $\bar{e}_i(X)$ is
obtained by repeating the measurement under the same conditions and
averaging, or by applying a denoising operator consistent with a
noise model, so that
$\bar{e}_i(X) \approx \E[e_i(X)]$.

In some embodiments, the measured response of device~$i$ is
decomposed into five terms with three functionally distinct
discrepancy components:
\[
  Y_i(X)
  \;=\;
  \underbrace{\hat{Y}_0(X)}_{\text{population model}}
  \;+\;
  \underbrace{\delta_{\mathrm{dev}}(X,\, i)}_{\text{device signature}}
  \;+\;
  \underbrace{\delta_{\mathrm{sys}}(X)}_{\text{shared discrepancy}}
  \;+\;
  \underbrace{\delta_{\mathrm{irr}}(X,\, i)}_{\text{irreducible residual}}
  \;+\;
  \varepsilon(X,\, i),
\]
where $\varepsilon(X, i)$ denotes measurement noise and the three
discrepancy components are defined on noise-reduced quantities as
follows:
\begin{enumerate}[label=(\roman*)]
  \item \textbf{Device-specific captured variation (device
    signature)} $\delta_{\mathrm{dev}}(X, i)$: the difference between
    the device-specific prediction and the population prediction,
    \[
      \delta_{\mathrm{dev}}(X,\, i)
      \;=\;
      \hat{Y}_i(X) \;-\; \hat{Y}_0(X)
      \;=\;
      \mathcal{Q}\!\bigl(X,\;\mathbf{d}_i\bigr)
      \;-\;
      \mathcal{Q}\!\bigl(X,\;\mathbf{d}_0\bigr).
    \]
    This term captures how device~$i$ differs from the fleet average
    in ways the model can represent.  It is a modelled quantity
    (computed entirely from the predictive model) and provides a
    computational characterisation of the individual device reflecting
    its microstructure, fabrication history, and physical properties.

  \item \textbf{Shared model error (shared discrepancy)}
    $\delta_{\mathrm{sys}}(X)$: the component of the noise-reduced
    per-device residual $\bar{e}_i(X)$ that is common across devices.
    In some embodiments, the shared model error is \emph{defined as}
    the robust central tendency of the noise-reduced residuals across
    the fleet under a declared aggregation rule (for example a mean,
    trimmed mean, or median) committed in the protocol digest:
    \[
      \delta_{\mathrm{sys}}(X)
      \;:=\;
      \operatorname{agg}_i \; \bar{e}_i(X),
    \]
    where $\operatorname{agg}$ denotes the declared and pre-committed
    aggregation functional (for example $\operatorname{robust\_mean}$,
    trimmed mean with declared trim fraction, or fleet median).  The
    aggregation rule, trim parameters, and any subpopulation
    stratification are recorded in the protocol digest as protocol
    parameters before the discovery procedure begins.  Under this
    definition, $\delta_{\mathrm{sys}}(X)$ is an exact functional of
    the observed fleet residuals, not a latent quantity; the equality
    sign in the decomposition above is exact by construction.
    Shared model error indicates where the predictive model~$\mathcal{Q}$
    itself is inadequate, regardless of which device is measured.
    Regions of the challenge space where $\delta_{\mathrm{sys}}(X)$
    is large are the primary targets of the proof-of-discrepancy
    protocol.  In some embodiments, the aggregation is performed
    within subpopulations of devices sharing similar configuration
    representations, to capture errors that are common within a
    manufacturing family but not necessarily fleet-wide; the
    subpopulation partition is also pre-committed.

  \item \textbf{Irreducible residual}
    $\delta_{\mathrm{irr}}(X, i)$: the device-specific component of
    the noise-reduced residual after removing the shared model error:
    \[
      \delta_{\mathrm{irr}}(X,\, i)
      \;=\;
      \bar{e}_i(X) \;-\; \delta_{\mathrm{sys}}(X).
    \]
    This residual represents aspects of device~$i$'s physical
    response that persist even when the best available device
    configuration representation $\mathbf{d}_i$ is used and the
    shared model error is accounted for.  It arises from the device
    possessing more internal degrees of freedom than the model can
    access through the available measurement channels.
\end{enumerate}
After subtracting the baseline prediction $\hat{Y}_0(X)$ and the
modelled device signature $\delta_{\mathrm{dev}}(X, i)$ from the raw
observation, the remaining per-device residual retains only the three
unmodelled terms. Accordingly, the observed (noisy) per-device residual satisfies:
\[
  e_i(X)
  \;=\;
  \delta_{\mathrm{sys}}(X)
  \;+\;
  \delta_{\mathrm{irr}}(X,\, i)
  \;+\;
  \varepsilon(X,\, i).
\]

\paragraph{Identifiability of the additive decomposition (non-limiting).}
The five-term decomposition above is additive by construction:
$\delta_{\mathrm{sys}}(X)$ is defined as the fleet aggregate of
residuals, and $\delta_{\mathrm{irr}}(X, i)$ is the per-device
residual after removing $\delta_{\mathrm{sys}}$.  Without additional
statistical structure, the three discrepancy components
$(\delta_{\mathrm{dev}}, \delta_{\mathrm{sys}}, \delta_{\mathrm{irr}})$
are not jointly identifiable from finite samples in general: shifting
mass between $\delta_{\mathrm{dev}}$ and $\delta_{\mathrm{irr}}$ can
leave the observed residuals unchanged.  In some embodiments, the
identifiability gap is addressed by placing independent Gaussian
process priors on $\delta_{\mathrm{sys}}(X)$ and the device-specific
component of the residual, with orthogonal covariance kernels
(for example distinct length-scale and output-scale hyperparameters
committed to the protocol digest); the resulting posterior
decomposes the observed residual into components whose credible
intervals are disjoint at the declared credibility level.  In this
Gaussian-process instantiation, the residual decomposition implements
fleet residual-margin calibration
(Section~\ref{sec:fleet-model}) applied to the per-device residual,
which is informed by the Bayesian calibration literature in the
manner of Kennedy and O'Hagan but extends that literature in the
respects identified in the novelty discussion above: the protocol
residual is the committed object, the residual-margin estimator is
the inferential apparatus, and the verifier physical re-execution is
the confirmation operation.
Where independent GP priors are not used, the decomposition is a
definitional accounting identity and the components should be
interpreted as defined quantities rather than independent physical
quantities; all downstream claims (hardness index, device identity,
PoD confirmation) are stated in terms of observable functionals of
the residual rather than in terms of the latent decomposition.

\paragraph{Approximate orthogonality diagnostic (non-limiting).}
For the functional roles of the decomposition to be cleanly separable,
the device-signature component $\delta_{\mathrm{dev}}(X, i)$ and the
irreducible residual $\delta_{\mathrm{irr}}(X, i)$ should be approximately
orthogonal: $\mathbb{E}_X[\delta_{\mathrm{dev}}(X,i) \cdot \delta_{\mathrm{irr}}(X,i)] \approx 0$
over challenges $X$ sampled from the protocol distribution.  If this
condition fails --- for example because $\mathbf{d}_i$ was over-fit to
the calibration set, causing representational error to contaminate
$\delta_{\mathrm{irr}}$ --- the security primitive is confounded with
modelling artefacts.  In some embodiments, after model convergence,
the empirical correlation
$\widehat{\mathrm{Corr}}[\delta_{\mathrm{dev}}(\cdot, i), \bar{e}_i(\cdot) - \delta_{\mathrm{sys}}(\cdot)]$
is computed on a held-out challenge set and committed to the protocol
digest; a declared orthogonality threshold $\rho_{\mathrm{orth}}$
triggers a security alert if the empirical correlation exceeds it,
indicating that further regularisation of $\mathbf{d}_i$ is needed
before the irreducible residual can serve as a security primitive.

\paragraph{Security domain and discovery domain partition
(non-limiting).}
The three functional roles of the discrepancy --- security primitive,
device identifier, and model improvement signal --- create a structural
tension: using challenges for PoD submissions exposes them to model
improvement (role iii), which may absorb part of $\delta_{\mathrm{irr}}$
into a more expressive $\delta_{\mathrm{dev}}$ under future model
updates, shifting the enrolled authentication signature.  In some
embodiments, this tension is managed by partitioning the challenge
space into two declared sub-domains: (a)~a \emph{security domain}
$\mathcal{X}_{\mathrm{sec}}$ where challenges are never disclosed in
PoD submissions, $\delta_{\mathrm{irr}}$ is protected from model
improvement, and device authentication is performed; and
(b)~a \emph{discovery domain} $\mathcal{X}_{\mathrm{disc}}$ where
challenges are used for PoD submissions and model improvement.  The
partition $({\mathcal{X}_{\mathrm{sec}}}, \mathcal{X}_{\mathrm{disc}})$
is declared and committed in the protocol digest before enrolment.
The security domain is designed to be representative of the device's
physical microstructure but non-overlapping with the discovery domain;
in some embodiments, the partition is drawn randomly and committed
cryptographically at enrolment time.  The learnability-exclusion
condition must hold independently on $\mathcal{X}_{\mathrm{sec}}$,
evaluated against adversary families that have access to PoD
submissions from $\mathcal{X}_{\mathrm{disc}}$ but not from
$\mathcal{X}_{\mathrm{sec}}$.

\paragraph{Functional roles of the three discrepancy components within the five-term decomposition (non-limiting).}
In some embodiments, each component of the decomposition serves
a distinct functional role within the fleet protocol:

\emph{Shared model error drives model improvement.}  Regions of the
challenge space where $\delta_{\mathrm{sys}}(X)$ is large indicate
where the predictive model is most inadequate.  Discovering and
characterising these regions is the primary target of the
proof-of-discrepancy protocol described in
Section~\ref{sec:pod-protocol}.  When shared model error is
identified, confirmed by independent physical reproduction, and
incorporated into an updated model, the fleet's collective predictive
capability improves.

\emph{Device-specific captured variation provides classification and
specialisation.}  The device signature
$\delta_{\mathrm{dev}}(X, i) = \mathcal{Q}(X, \mathbf{d}_i) - \mathcal{Q}(X, \mathbf{d}_0)$
is a modelled quantity that characterises how a device differs from
the fleet average.  Because $\delta_{\mathrm{dev}}$ depends on the
current model version, it may change when the model is updated even
if the physical device has not changed.  In some embodiments,
$\delta_{\mathrm{dev}}$ is therefore used primarily for fleet
management, device classification, and task allocation (assigning
computational or sensing tasks to devices whose specific variation
patterns make them well-suited to those tasks), rather than for
security-critical authentication.  In some embodiments where
$\delta_{\mathrm{dev}}$ is used for device identification, the
identification is bound to a specific model version.

\emph{Irreducible residual provides an empirically calibrated security
margin against cloning under declared attacker families and meter
envelopes, and in some embodiments serves as the primary authentication
primitive.}
The irreducible residual $\delta_{\mathrm{irr}}(X, i)$ captures
the component of the per-device response that is not well-approximated
by the predictive model under declared resource bounds, and is therefore
difficult to forge using a digital model alone.  It provides a security
margin against adversaries who attempt to simulate or clone a device's
responses using a model-based approach.  Unlike $\delta_{\mathrm{dev}}$,
which is explicitly computed from $\mathcal{Q}$, the irreducible
residual reflects physical properties of the device that the model
does not capture; as a result, its magnitude tends to remain
comparatively stable across model updates, though it is not
definitionally invariant to model changes and may shift if a
substantially more expressive model is deployed.  In some embodiments,
the magnitude of the irreducible residual at a given challenge is used
as a hardness metric: challenges where the irreducible residual is large
are challenges where the device's response is most difficult to forge
under the current model class.  In some embodiments, device
authentication and provenance attestation are grounded in the
irreducible residual rather than in the model-derived device signature.

\paragraph{Relationship to prior calibration formulations
(non-limiting).}
Fleet residual-margin calibration relates to prior calibration
approaches as follows.  Bayesian calibration in the manner of Kennedy
and O'Hagan provides the closest statistical antecedent and supplies
the Gaussian-process emulation tooling that may instantiate the
residual-margin estimator; in such frameworks, the discrepancy term
is typically treated as a nuisance parameter to be estimated and
marginalised over, with the goal of reducing its impact on parameter
inference and prediction.  Frequentist calibration approaches and
adversarial benchmarking are more distant.  Metrological
inter-comparison and proficiency testing share the fleet-level scope
but do not commit the protocol residual or perform verifier physical
re-execution.  The five-term decomposition in
Section~\ref{sec:error-decomposition} extends the Bayesian calibration
literature by treating the protocol residual not as a deficiency to
minimise but as a structured resource whose components serve distinct
and complementary roles within the fleet protocol; none of the prior
approaches treats the protocol residual as a committed resource
confirmed by independent physical re-execution on distinct unclonable
devices, which is the load-bearing distinction of fleet
residual-margin calibration.  This functional reframing---from
``model error as statistical nuisance'' to ``model error as
multi-role protocol resource''---is a distinguishing feature
of the disclosed system.

In some embodiments, the conditions under which this reframing is
valid are stated explicitly as follows.  The three functional roles
create exactly one structural tension: as the fleet model
$\mathcal{Q}$ improves over time (role~iii, $\delta_{\mathrm{sys}}$
decreases), a more expressive model can potentially absorb part of
what was $\delta_{\mathrm{irr}}$ into a better $\delta_{\mathrm{dev}}$
(threatening role~i, the security primitive).  The condition under
which roles~i and~iii remain non-conflicting is that the learnability-exclusion
condition holds permanently: the fleet model never achieves sufficient
expressiveness to substantially reduce $\delta_{\mathrm{irr}}$ for
the device population under the declared adversary families.  This
condition is not definitionally guaranteed and must be monitored
empirically: in some embodiments, $|\delta_{\mathrm{irr}}|$ is
tracked across model update cycles, and a declared
learnability-exclusion violation threshold triggers a security alert
if $|\delta_{\mathrm{irr}}|$ drops by more than a stated fraction
between consecutive model updates.  The functional reframing is valid
under the learnability-exclusion condition; the monitoring protocol
is the empirical check on whether the condition continues to hold.

\paragraph{Model-update stability criterion and subspace separation
(non-limiting).}
In some embodiments, the following model-update stability criterion
is declared in the protocol digest and checked after each fleet model
update: (i)~on the \emph{improvement subspace} (challenge set
$\mathcal{X}_{\mathrm{disc}}$ used for PoD submissions), the shared
model error $\delta_{\mathrm{sys}}$ must decrease or remain
stable---a model update that increases $\delta_{\mathrm{sys}}$ on
the improvement subspace is rejected and triggers diagnostic audit;
(ii)~on the \emph{verification subspace} (challenge set
$\mathcal{X}_{\mathrm{sec}}$ reserved for authentication), the
lower confidence bound on $|\delta_{\mathrm{irr}}|$ must not collapse
beyond a declared tolerance $\eta_{\mathrm{sec}}$ between model
versions --- a larger collapse triggers a security alert and halts
model deployment pending re-evaluation; (iii)~the cross-subspace
information leakage (estimated as conditional mutual information
between model-improvement updates and verification-subspace residuals)
must remain below a declared leakage threshold $I_{\mathrm{leak}}$,
committed in the protocol digest.  In some embodiments, cross-leakage
is bounded by the disclosure policy on improvement-subspace challenge
configurations (submitted PoD records do not include raw
verification-subspace traces), and the residual leakage is
empirically upper-bounded on the anchor embodiment.  These three
conditions, checked jointly after each model update, constitute the
model-update stability criterion; a model version that fails any
condition is not deployed until the failure is diagnosed and
remediated.

\paragraph{Decomposition is architectural, not post-hoc
(non-limiting).}
In some embodiments, the five-term decomposition arises
naturally from the architecture of the predictive model and the
structure of fleet data.  The device signature
$\delta_{\mathrm{dev}}$ is determined by the difference between
device-specific and population-level model predictions.  The shared
model error $\delta_{\mathrm{sys}}$ is the fleet-wide central
tendency of the per-device residuals.  The irreducible residual
$\delta_{\mathrm{irr}}$ is the device-specific component that remains
after the shared error is removed.  No separate decomposition
algorithm beyond standard fleet-level residual analysis is required.

% ----------------------------------------------------------------------
\subsection{Proof-of-Discrepancy protocol}
\label{sec:pod-protocol}
% ----------------------------------------------------------------------

In some embodiments, a distributed protocol---referred to herein as a
proof-of-discrepancy protocol---rewards participants for discovering
reproducible experimental methods that expose errors in the shared
predictive model.  The protocol operates as follows in non-limiting
outline:

\paragraph{Method submission.}
A participant (the ``discoverer'') identifies a challenge
configuration $X^*$ and associated operating conditions under which
the predictive model's prediction $\mathcal{Q}(X^*, \mathbf{d}_i)$ diverges
from the measured physical response $Y_i(X^*)$ by more than a
declared threshold, after accounting for measurement noise and known
device-specific variation.  The claimed discrepancy is evaluated on
the per-device residual $e_i(X^*) = Y_i(X^*) - \mathcal{Q}(X^*, \mathbf{d}_i)$
or a statistic derived therefrom.
The discoverer submits $X^*$ and
associated operating conditions as a reproducible experimental
method---a specification sufficient for another participant to execute
the same experiment on a different device.  In some embodiments, the
submission is committed (for example via a cryptographic commitment
scheme) before the verifiers are selected, to reduce the risk of retroactive
modification.

\paragraph{Controlled execution conditions (non-limiting).}
In some embodiments, proof-of-discrepancy methods are executed against
a standardised characterisation target, an internal reactor or
controlled scattering medium, or an environment specification
sufficiently constrained that different parties can reproduce it (for
example a supplied calibration artefact or a specified internal
operating mode).  The protocol need not require reproduction of
arbitrary real-world scenes; the subject of the predictive model is
the device's own physical response under controlled conditions.

\paragraph{Submission commitments (non-limiting).}
In some embodiments, each submission commits to the challenge
configuration $X^*$, the version identifier of the predictive model
against which the discrepancy is claimed, the calibration protocol
identifier, and any device firmware or operating mode identifiers
relevant to reproducibility.  In some embodiments, submissions further
specify environment constraints, allowed operating ranges, and data
processing procedures sufficient for independent reproduction under
comparable conditions.  Verification is performed against the
committed model version, even if a newer version is available at the
time of verification.

\paragraph{Non-limiting submission record format: method-object tuple for discrepancy submissions.}
In some embodiments, the submitted method object is represented as a tuple
comprising at least a challenge configuration, a device or source-class
identifier, a committed protocol digest, a model-version identifier, a
measured-response commitment, an asserted residual or scoring statistic,
environmental or calibration constraints, and any required reproduction
instructions or selective-opening commitments.

In some embodiments, a submission record includes the following fields
(serialised in a canonical binary format such as CBOR or a
length-prefixed binary encoding):
\begin{description}[nosep,leftmargin=1.5em,labelindent=0pt]
  \item[challenge\_config] Challenge configuration $X^*$ (operating
    point coordinates, scan protocol parameters, environmental
    setpoints).
  \item[model\_version] Cryptographic hash of the predictive model
    version $\mathcal{Q}$ against which the discrepancy is claimed.
  \item[device\_id] Device configuration representation
    $\mathbf{d}_i$ or anonymised device identifier.
  \item[protocol\_digest] Per-run protocol digest $\Pi_{\mathrm{dig}}$ specifying
    calibration identifiers, firmware version, and data-processing
    pipeline.
  \item[measured\_response] Committed hash of the discoverer's
    measured response $Y_i(X^*)$.
  \item[asserted\_residual] Summary statistic of $e_i(X^*)$ (for
    example mean, variance, or a test statistic exceeding the
    declared threshold).
  \item[environment\_spec] Environmental constraints sufficient for
    independent reproduction (temperature range, humidity, vibration
    class, or ``internal calibration target'' indicator).
  \item[timestamp] Externally anchored timestamp (for example GNSS
    time or a trusted timestamping service).
  \item[commitment] Cryptographic commitment binding all preceding
    fields.
\end{description}
In some embodiments, a verifier's confirmation record includes the
same fields for the verifier's device plus a cross-reference to the
discoverer's commitment hash.  In some embodiments, the discrepancy
is scored by a statistical test (for example a two-sample test
comparing discoverer and verifier residuals against model predictions,
with a pre-registered significance level recorded in the protocol
digest).

\paragraph{Verifier selection and physical re-execution.}
In some embodiments, one or more verifiers are selected from the
fleet.  Each verifier executes the submitted method $X^*$ on their
own device under the specified operating conditions and records the
physical response.  In preferred embodiments, verification includes
physical re-execution of
the experiment; computational re-evaluation of the model alone does not
constitute verification in these embodiments.  In some embodiments, verifiers are selected
by a randomised process to reduce or deter collusion between discoverer and
verifier.

\paragraph{Discrepancy confirmation.}
In some embodiments, fleet-level confirmation uses an explicit
confirmation statistic $T(X^*)$ defined as a robust aggregate of the
per-verifier residuals at the submitted challenge configuration
$X^*$.  A non-limiting confirmation statistic is:
\[
  T(X^*) \;:=\;
  \operatorname{agg}_{j \in \mathcal{V}}
  \bar{e}_j(X^*),
\]
where $\mathcal{V}$ is the set of participating verifiers,
$\bar{e}_j(X^*)$ is the noise-reduced residual of verifier~$j$ at
$X^*$, and $\operatorname{agg}$ is the pre-registered aggregation
rule (for example fleet median, trimmed mean with declared trim
fraction, or a quantile-based statistic).  This is the verifier-side
instantiation of the shared discrepancy functional
$\delta_{\mathrm{sys}}(X)$ (Section~\ref{sec:error-decomposition})
evaluated at the submitted challenge $X^*$; the confirmation
statistic is numerically large when the verifiers' residuals are
collectively large relative to model predictions, consistent with a
genuine shared model discrepancy.  The discrepancy is
\emph{confirmed} if $T(X^*) > \tau_{\mathrm{confirm}}$, where
$\tau_{{\mathrm{confirm}}}$ is a pre-registered confirmation threshold.

In some embodiments, two non-limiting verifier-population
embodiments are supported. (i) \emph{Matched-verifier embodiment.}
All verifiers are of a declared common device class with a declared
common response space, and the residuals $\bar{e}_j(X^*)$ are
aggregated directly. (ii) \emph{Heterogeneous-verifier embodiment.}
Verifiers may be of different declared device classes; in such
embodiments, residuals are placed into a declared common response
metric committed in the protocol digest before aggregation. A
non-limiting heterogeneous-verifier example: each verifier's
$\bar{e}_j(X^*)$ is divided by its declared per-device
measurement-noise standard deviation $\hat\sigma_j$, producing a
dimensionless $z_j(X^*) = \bar{e}_j(X^*) / \hat\sigma_j$, and
$T(X^*) = \operatorname{agg}_{j\in\mathcal{V}} z_j(X^*)$. Other
declared common metrics (response-norm normalisation, declared
calibration-uncertainty normalisation, or a committed common-metric
tag) are non-limiting alternatives. The choice of embodiment, the
declared common metric where applicable, and the per-device
normalising quantities are committed to the protocol digest before
verification begins.

In some embodiments, a quorum rule additionally requires that at
least $q$ out of $n$ verifiers individually exceed a per-device
exceedance threshold $\tau_{\mathrm{device}}$.  The parameters
$(n, q, \tau_{\mathrm{confirm}}, \tau_{\mathrm{device}},
\operatorname{agg})$ are pre-committed in the protocol digest before
verification begins and constitute part of the PoD method record.
In some embodiments, the aggregation functional is chosen with its
breakdown point explicitly stated in the protocol digest. In the
convention used here, the finite-sample breakdown point of an
aggregator is the largest fraction of contamination below which
the aggregator's output remains bounded under worst-case
contamination; for the sample median this convention gives a
breakdown point of $0.5$, with up to $\lfloor (n-1)/2 \rfloor$
colluding verifiers tolerable when strict tolerance is required,
or up to $\lfloor n/2 \rfloor$ when the convention permits boundary
contamination. A trimmed mean with trim
fraction $\alpha$ has breakdown point $\alpha$.  The required
breakdown point $\beta$ is declared as the maximum tolerated collusion
fraction under the adversary model and must satisfy
$\beta \leq$ (breakdown point of the declared aggregator).  In
some embodiments, the recency of verifier device calibration is
additionally declared: each verifier's $\mathbf{d}_j$ must have been
re-estimated within a declared calibration window before the PoD
submission, to bound the drift bias in $\bar{e}_j(X^*)$.

In some embodiments, the adversary model for the confirmation
protocol is stated explicitly as a Byzantine fault model: at most
$f$ of the $n$ verifiers in $\mathcal{V}$ are assumed to be colluding
adversaries (Byzantine), while the remaining $n - f$ are honest.
The quorum parameters $(n, q)$ are chosen to satisfy
$n \geq 3f + 1$ and $q \geq 2f + 1$, or an alternative declared
tolerance level is committed in the protocol digest with an explicit
derivation. In some embodiments, the protocol-level Byzantine tolerance is
\[
  f_{\mathrm{BFT}} = \min\!\left( \lfloor (n-1)/3 \rfloor,\;
    \lfloor (q-1)/2 \rfloor \right).
\]
Where a declared robust aggregator with breakdown point $\beta$ is used,
the operational tolerated collusion level is further capped at
\[
  f_{\mathrm{oper}} = \min\!\left( f_{\mathrm{BFT}},\;
    \lfloor n \cdot \beta \rfloor \right).
\]
The declared confirmation rule, quorum parameters $(n, q)$, the aggregator
breakdown point $\beta$, and the resulting $f_{\mathrm{oper}}$ are
committed in the protocol digest. Consistency of $T(X^*)$ across
independently operated verifier devices provides evidence that the
discrepancy reflects shared model inadequacy rather than a
device-specific artefact, measurement error, or fabricated result.

\paragraph{Physical anchoring relative to adversarial benchmarking
(non-limiting).}
The PoD protocol is related to adversarial benchmarking approaches
such as Dynabench, in which human contributors submit examples that
fool the current model, driving iterative improvement.  A key
structural difference in the physical setting is that a submitted
challenge $X^*$ must reproduce on independently manufactured verifier
devices that have not been inspected by the submitter; this physical
reproduction requirement anchors the PoD dataset to physical reality
in a way digital adversarial datasets are not.  A contributor cannot
retroactively alter a device's physical response, so the submitted
challenge must generalise to independent physical instantiations, not
merely fool the current digital model.  This turns the
Dynabench bias --- contributors optimise against the current model's
failure modes --- into a \emph{coverage bias} in the physical setting:
submitters choose physically accessible challenges from the discovery
domain, which may underrepresent challenge regions that are physically
difficult to reach (extreme operating conditions, rare material
states).  In some embodiments, a coverage scoring component in the
incentive mechanism rewards submissions from underrepresented regions
of the challenge space; the coverage metric and the definition of
``underrepresented'' are declared in the protocol digest.

\paragraph{Model update.}
In some embodiments, confirmed discrepancies are incorporated into the
predictive model.  The model is updated to reduce shared model error
in the region of the challenge space identified by the confirmed
discrepancy, while preserving the model's fidelity elsewhere.  In
some embodiments, model updates are validated against held-out
challenge--response data before deployment, to detect regressions.
In some embodiments, model updates are logged with provenance
information linking each update to the specific confirmed discrepancy
that motivated it.

\paragraph{Model versioning and archival (non-limiting).}
In some embodiments, model versions are append-only and archived, so
that prior confirmed discrepancies remain verifiable against the model
version under which they were submitted.  In some embodiments,
pending submissions are evaluated against the model version committed
at the time of submission, even if a newer version has been deployed,
to avoid retroactive invalidation or alteration of in-progress verification by model updates.

\paragraph{Incentive structure (non-limiting).}
In some embodiments, the protocol includes an incentive mechanism that
rewards discoverers for confirmed discrepancies.  In some embodiments,
rewards are scaled by one or more of: the magnitude of the confirmed
discrepancy, the novelty of the challenge region (whether the
discrepancy occurs in a region of the challenge space not previously
known to be problematic), and the informativeness of the discrepancy
for model improvement (for example measured by reduction in
predictive error after the model update).  In some embodiments,
verifiers are also rewarded for faithful execution of the
verification protocol.  In some embodiments, the incentive mechanism
uses tokenised rewards; in some embodiments, incentives are
reputational, contractual, or institutional.  The specification does
not restrict the form of the incentive mechanism.

% ----------------------------------------------------------------------
\subsection{Physical reproduction as the verification primitive}
% ----------------------------------------------------------------------

A distinguishing feature of the proof-of-discrepancy protocol is that,
in preferred embodiments, verification includes independent physical
re-execution of the submitted experimental method on devices that are
physically distinct from the discoverer's device.  This section elaborates on the
properties and implications of this verification primitive.

\paragraph{Why physical reproduction is the preferred verification mechanism (non-limiting).}
Computational re-evaluation---running the predictive model on the
submitted challenge and checking whether the model's own output
diverges from its training data---does not constitute verification
of a model--reality discrepancy, because the discrepancy is by
definition a property of the gap between the model and physical
reality, not a property of the model alone.  In preferred embodiments,
physical re-execution on an independent device provides a particularly
strong verification primitive for establishing what physical reality
actually produces in response to a given challenge, because it is
grounded in measurement rather than model evaluation alone.  Independent physical reproduction on
a device that is physically distinct from the discoverer's device
further guards against fabrication: if the discoverer has tampered
with their device or measurement apparatus, the fabricated result
will not reproduce on the verifier's independently manufactured and
operated device.

\paragraph{Role of physical unclonability (non-limiting).}
In some embodiments, devices in the fleet exhibit physical
unclonability arising from manufacturing stochasticity (for example
in the microstructure of scattering media, the precise geometry of
optical paths, or the fabrication-specific properties of detector
elements).  Physical unclonability is designed so that each device's
detailed physical response to a given challenge is unique and is
computationally infeasible to replicate by constructing a duplicate
device under declared adversary families.  This property
strengthens the verification primitive: an adversary who fabricates a
discrepancy claim by manipulating their own device is unable, under
the declared resource bounds, to predict or control what an
independently manufactured verifier device will
produce when executing the same method, because the verifier device's
microstructure differs in unpredictable ways.

In some embodiments, physical unclonability also supports the
expectation that the irreducible residual $\delta_{\mathrm{irr}}(X, i)$
(Section~\ref{sec:error-decomposition}) is substantial across the
fleet: under the declared observability and resource limits, no
practical model is expected to eliminate the gap between its
predictions and the physical response of a device whose internal
degrees of freedom exceed what is accessible through the measurement
interface.

\paragraph{Distinction from computational, cryptographic, and
statistical verification (non-limiting).}
The physical reproduction verification primitive is distinct from:
computational verification (re-executing code or re-evaluating a
function), cryptographic verification (checking a digital signature,
zero-knowledge proof, or hash), and statistical verification
(checking consistency of reported data against distributional
assumptions).  Physical reproduction involves the verifier
physically operating hardware and performing a measurement.  In some
embodiments, the physical reproduction verification primitive is
combined with one or more of these other verification methods---for
example, cryptographic commitments may be used to bind the
discoverer's claim before verification, and statistical tests may
be used to assess whether the reproduced measurements are consistent
with the claimed discrepancy---but in preferred embodiments, physical
re-execution serves as the primary verification step.

% ----------------------------------------------------------------------
\subsection{Difficulty scaling from physics}
% ----------------------------------------------------------------------

In some embodiments, it is expected---and is stated as a
research conjecture in the following paragraph---that the difficulty
of discovering new proof-of-discrepancy submissions scales with the
fidelity of the predictive model, creating a natural progression
analogous to difficulty adjustment in computational proof-of-work
protocols but grounded in physics rather than arbitrary computation.
This expectation motivates the PoD Monotonicity Conjecture; it is not
itself an established result.

\paragraph{PoD Monotonicity Conjecture (non-limiting, research direction).}
The following is stated as a research conjecture.  Let
$\mathcal{Q}_t$ denote the fleet predictive model at time $t$, and
let $D_t$ denote the expected number of candidate challenges a
computationally bounded discoverer must test before finding a
challenge $X^*$ satisfying $|\delta_{\mathrm{sys}}(X^*;
\mathcal{Q}_t)| > \tau_{\mathrm{confirm}}$ (the confirmation
threshold).  Let $\|\mathcal{Q}_t\|$ be a measure of model quality
(for example, mean squared prediction error on held-out challenges,
decreasing as the model improves).
\medskip
\noindent\textit{PoD Monotonicity Conjecture.}
$D_t$ is monotone non-decreasing as $\|\mathcal{Q}_t\|$ decreases:
improving the fleet model reduces the density of discoverable
discrepancies, so each successive confirmed discovery requires more
search effort.  Formally:
\[
  \frac{d D_t}{d (-\|\mathcal{Q}_t\|)} \;\geq\; 0
\]
in expectation over the challenge space distribution and fleet
device heterogeneity, under regularity conditions on the challenge
space and the confirmation criterion.
\medskip

For a Gaussian process model $\mathcal{Q}$, this conjecture is
tractable: as the GP posterior contracts (more fleet data reduces
posterior variance uniformly), the expected number of challenges
required to find one with posterior variance exceeding
$\tau_{\mathrm{confirm}}$ increases monotonically.  Extension to
neural network models requires additional conditions: the model must
improve uniformly across the challenge space (no pathological
residual-concentration), and the confirmation criterion must remain
stable as the model updates.  In some embodiments, the conjecture is
treated as an empirical claim monitored in the protocol digest: the
fleet logs an empirical $D_t$ estimate (from discoverer attempt
counts) and reports whether $D_t$ is increasing across consecutive
model update cycles.  A sustained decrease in $D_t$ following a model
update is a signal that the update may have introduced model
misspecification or that the challenge space is not being covered
adequately.

A formal proof of the PoD Monotonicity Conjecture, even in the
simplified GP case, would constitute a new distributed systems result:
it would establish that the PoD protocol has a
\emph{physically-grounded difficulty adjustment} property analogous
to---but mechanistically distinct from---proof-of-work difficulty
adjustment, grounded in the depletion of modellable discrepancy mass
rather than in an adjustable computational parameter.


Even if discrepancy-discovery difficulty is not monotone across model
updates, the PoD protocol continues to provide a physically anchored
mechanism for discovering, replicating, superseding, and withdrawing
discrepancy claims. Model improvement therefore rests on reproducible
correction cycles rather than on any assumed monotone search-difficulty
law.

\paragraph{Resolution-dependent difficulty (non-limiting).}
As the predictive model improves through incorporation of confirmed
discrepancies, the remaining shared model error
$\delta_{\mathrm{sys}}(X)$ decreases in magnitude and becomes
concentrated in regions of the challenge space where the physics is
most complex or where measurement resolution has not yet been
sufficient to reveal the discrepancy.  Discovering the next
discrepancy may require finer measurement resolution, more precise
control of operating conditions, or exploration of previously
untested regions of the challenge space.  This creates a difficulty
progression driven by the physics of the measurement interface
rather than by an adjustable computational parameter.

\paragraph{Depth of physical complexity (non-limiting).}
In some embodiments, the physical complexity of the devices is such
that increasing measurement resolution may reveal additional
discrepancies, because the devices possess more internal degrees of
freedom than any model can capture through a finite measurement
interface.  In some embodiments, empirical evidence indicates that
discrepancies persist across accessible scales: real scattering media,
optical systems, and other high-degree-of-freedom physical substrates
exhibit behaviour at finer resolutions that is not fully predicted by
models trained at coarser resolutions.  In some embodiments, this
property is verified empirically by demonstrating that increasing
measurement resolution consistently reveals new discrepancies that
were not predicted by the current model.

% ----------------------------------------------------------------------
\subsection{Learnability-exclusion condition}
\label{sec:learnability-exclusion}
% ----------------------------------------------------------------------

In some embodiments, the physical devices used in the
proof-of-discrepancy protocol satisfy a practical
learnability-exclusion condition: no known feasible learning procedure
achieves prediction accuracy above the verification threshold when
given access to a bounded number of challenge--response observations
from the device.  Informally, a computationally bounded adversary
is not expected to construct a model that reliably predicts the device's response
to unseen challenges with sufficient accuracy to fabricate
verification-passing results.

\paragraph{Physical basis (non-limiting).}
In some embodiments, the practical learnability-exclusion condition is
grounded in physical properties of the device: volumetric multiple
scattering through disordered media, controlled optical
nonlinearities, partially observable internal state, and a number of
effective internal degrees of freedom that exceeds the information
accessible through the measurement interface.  These properties are intended to make the challenge--response mapping
difficult to approximate by functions in classes that are efficiently
learnable from input--output examples under known learning algorithms;
this learnability-exclusion condition is evaluated empirically against
declared attacker families and meter envelopes.  In some embodiments, this
condition is interpreted through a complexity-theoretic lens as
resistance to polynomial-time characterisation, but the operative
definition is empirical rather than formal.

\paragraph{Empirical verification (non-limiting).}
In some embodiments, the practical learnability-exclusion condition is
verified empirically rather than proved mathematically.  Verification
comprises training state-of-the-art machine learning models (for
example deep neural networks, kernel methods, compressed sensing
algorithms, or other approaches representing the best known
learning algorithms) on challenge--response data from a target
device and demonstrating that prediction accuracy on held-out
challenges plateaus below the verification threshold.  In some
embodiments, this verification is performed periodically as learning
algorithms improve, to confirm that the exclusion condition continues
to hold under the current state of the art.

\paragraph{Relationship to irreducible residual (non-limiting).}
The practical learnability-exclusion condition provides the physical
basis for the irreducible residual $\delta_{\mathrm{irr}}(X, i)$
described in Section~\ref{sec:error-decomposition}.  If the condition
holds, the irreducible residual remains substantial regardless of the
sophistication of the predictive model, because the model is not expected to close
the gap without information that is physically inaccessible through
the measurement interface.

\paragraph{Operative lower bound on $\|\delta_{\mathrm{irr}}\|$
(non-limiting).}
In some embodiments, the residual-margin estimator
(Section~\ref{sec:fleet-model}) yields an
explicit posterior credible interval on each component of the five-term
decomposition including $\delta_{\mathrm{irr}}$.  In a
Gaussian-process instantiation of the residual-margin estimator
(informed by the Bayesian calibration literature in the manner of
Kennedy and O'Hagan), after
the GP posterior is conditioned on all available challenge--response
data from the fleet, the lower tail of the marginal posterior on
$\delta_{\mathrm{irr}}(X, i)$ provides an operative lower bound: the
irreducible residual is declared to exceed $\delta_{\mathrm{irr}}^{\min}$
at posterior probability $(1-\alpha)$ for a declared credibility level
$\alpha$ (for example $\alpha = 0.05$).  Other estimator classes
(neural posteriors, physics-based predictive distributions, or
hybrids) may instantiate the residual-margin estimator equivalently;
the estimator class is not load-bearing, but the commitment of the
residual and the verifier physical re-execution are.
This bound is a function of
the training CRP budget $n_{\mathrm{CRP}}$, the declared attacker model
class $\mathcal{H}$, and the estimator hyperparameters, all of which are
committed to the protocol digest.  In some embodiments, the bound is
updated periodically as new fleet data arrive and new attacker model
families are evaluated; the hardness index
(Section~\ref{sec:hardness-index}) includes $\delta_{\mathrm{irr}}^{\min}$
as a component, so that the stated security margin is directly traceable
to the residual-margin estimator output rather than to an unquantified
claim.  An attacker producing a discrepancy below this bound is within
the fleet's calibrated noise floor and does not constitute a
Proof-of-Discrepancy event.  The bound
shrinks if new learning algorithms reduce the unexplained residual
and is updated accordingly; no absolute floor is asserted.

% ----------------------------------------------------------------------
\subsection{Trust-bootstrapped reference devices}
% ----------------------------------------------------------------------

In some embodiments, the proof-of-discrepancy protocol is anchored
by one or more reference devices whose measurements are trusted to
higher confidence than measurements from general fleet participants.
Reference devices may be operated in controlled environments,
characterised under metrological standards (for example consistent
with ISO/IEC 17025 requirements for testing and calibration
laboratories), or operated by trusted parties.  Reference-device measurements provide an external ground truth that helps
anchor the fleet model against drift into self-consistent but physically
inaccurate states.

In some embodiments, the predictive model is periodically validated
against reference device measurements, and discrepancies between
model predictions and reference device responses are treated as
high-priority model update targets.  In some embodiments, new fleet
participants are initially calibrated against reference device
measurements to establish their device configuration representations
with known accuracy before participating in the broader protocol.

% ----------------------------------------------------------------------
\subsection{Scoring and method evaluation (non-limiting)}
% ----------------------------------------------------------------------

In some embodiments, submitted methods are evaluated using a scoring
function that reflects one or more of the following properties:

\begin{enumerate}[label=(\roman*)]
  \item \textbf{Confirmation:} the degree to which the submitted
    method's predicted discrepancy is confirmed by independent
    physical reproduction.

  \item \textbf{Discovery:} the degree to which the discrepancy is
    novel---occurring in a region of the challenge space where the
    model was previously considered adequate, or exceeding previously
    known discrepancy magnitudes.

  \item \textbf{Coverage:} the degree to which the submitted method
    explores regions of the challenge space that have been
    under-sampled by previous submissions, contributing to broader
    characterisation of model adequacy.
\end{enumerate}

In some embodiments, the scoring function is derived from a proper
scoring rule to incentivise honest reporting of predicted discrepancy
magnitudes.  In some embodiments, the scoring function is a weighted
combination of the above components, with weights adjusted over time
to balance model improvement (favouring discovery and coverage) against
verification robustness (favouring confirmation).

% ----------------------------------------------------------------------
\subsection{Robustness against adversarial participants (non-limiting)}
% ----------------------------------------------------------------------

In some embodiments, the proof-of-discrepancy protocol incorporates
mechanisms to mitigate adversarial behaviour by participants,
including:

\paragraph{Sybil resistance.}
In some embodiments, the cost of manufacturing and operating a
physical device (including acquisition, calibration, and ongoing
maintenance) provides a natural Sybil resistance mechanism: creating
fake identities would typically require acquiring and operating physical hardware under the declared identity-anchoring scheme,
which is more expensive than creating digital pseudonyms.  In some
embodiments, device identity is verified via microstructure-based
challenge--response authentication
(Section~\ref{sec:puf-auth}) to confirm that each fleet participant
controls a distinct physical device.  In some embodiments, as a
fallback that does not rely solely on manufacturing cost, device
identity is additionally anchored by a cryptographic commitment of
the device's enrolled configuration representation $\mathbf{d}_i$
(or a secure hash thereof) published to a tamper-evident log at
enrolment time; a physical one-way function
(as described in the one-way state update and hash-chain logging
subsection) evaluated on the
device's response to a held-out challenge set, whose output is
committed and unsealed at identity-claim time, provides an
additional Sybil barrier that is robust to cost collapse of the
manufacturing barrier.  These cryptographic and physical commitment
mechanisms are non-limiting and may be combined in any proportion
declared in the identity-anchoring policy.  In some embodiments, a
minimum manufacturing cost threshold $C_{\min}$ is declared in the
identity-anchoring policy and committed to the protocol digest; the
Sybil resistance claim based on manufacturing cost is conditioned on
this threshold, and the policy requires re-evaluation of the
cryptographic-commitment fallback whenever market device cost reports
indicate costs approaching $C_{\min}$.  A cross-verification
protocol --- in which each new device is co-located with and
cross-verified by at least $k_{\mathrm{cross}}$ existing trusted
fleet devices before receiving voting weight --- provides an
additional Sybil barrier independent of manufacturing cost.

\paragraph{Collusion resistance.}
In some embodiments, verifiers are selected by a randomised process from
the fleet, so that discoverers are not expected to predict or choose
which devices will verify their submissions. The protocol-level
Byzantine tolerance, quorum parameters, declared robust-aggregator
breakdown point, and operational tolerated collusion level are computed
and committed as described above, including $f_{\mathrm{BFT}}$, $\beta$,
and $f_{\mathrm{oper}}$. A colluding minority below the declared
operational tolerated collusion level may delay, but not corrupt, a PoD
confirmation under the declared confirmation rule.
\paragraph{Fabrication resistance.}
Physical unclonability of devices is designed so that fabricated
discrepancy claims---where a discoverer reports a discrepancy that
does not correspond to a genuine model--reality gap---are unlikely to
reproduce on independently manufactured verifier devices.  The
physical reproduction verification primitive is the primary
fabrication resistance mechanism.

\paragraph{Robust aggregation (non-limiting).}
In some embodiments, model updates derived from confirmed
discrepancies are aggregated using robust statistical methods that
down-weight or exclude outlier contributions.  In some embodiments,
contributions from devices with unusual verification histories are
weighted differently from contributions with consistent track records.

\paragraph{Submission economics (non-limiting).}
In some embodiments, the protocol includes economic mechanisms to
discourage frivolous or adversarial submissions, such as requiring a
bond or deposit that is returned upon confirmation and forfeited upon
rejection, or rate-limiting submissions based on device identity,
reputation, or staking.

\paragraph{Update containment (non-limiting).}
In some embodiments, confirmed model updates are evaluated in a shadow
or quarantine mode before fleet-wide deployment: the updated model is
tested against held-out challenge--response data to detect regressions
(cases where a correction in one region of the challenge space
degrades fidelity in another).  In some embodiments, the challenges
used for verification of submitted discrepancies are drawn from a
different distribution than the challenges used for model training, to
reduce overfitting risk with respect to the distribution of submitted challenges.

% ----------------------------------------------------------------------
\subsection{Meter partition for the proof-of-discrepancy protocol
(non-limiting)}
% ----------------------------------------------------------------------

In some embodiments, the challenge space of each device is
partitioned into disjoint subspaces allocated to distinct functional
roles:

\begin{enumerate}[label=(\roman*)]
  \item \textbf{Verification subspace:} challenges reserved for
    device authentication and identity verification.  Responses to
    challenges in this subspace are used for microstructure-based identification
    and are not disclosed to the fleet model.

  \item \textbf{Sensing and computation subspace:} challenges used
    for computational tasks, reservoir computing, or sensing
    applications.

  \item \textbf{Model improvement subspace:} challenges used for
    fleet model calibration and proof-of-discrepancy submissions.
\end{enumerate}

In some embodiments, the partition is supported by protocol-level
constraints intended to keep challenges from one subspace from being
reused in another, analogous to the security meter partition described in
Section~\ref{sec:security-theory}.  The partition is designed so that
information disclosed during model improvement does not compromise
device authentication, and that verification challenges are not
predictable from publicly shared model-improvement data.
In some embodiments, model-improvement disclosures (including
submitted methods, confirmed discrepancies, and model updates) are
constrained so that they do not leak verification-subspace challenges
or any sufficient statistics from which verification-subspace
responses could be predicted.  In some embodiments, published
proof-of-discrepancy records include the challenge configuration,
model version, and a discrepancy statistic with confidence bounds,
while raw device response traces are retained under access control or
disclosed via selective opening of committed records consistent with
the disclosure policies described elsewhere in this specification.

% ----------------------------------------------------------------------
\subsection{Embodiments and non-limiting examples}
% ----------------------------------------------------------------------

The proof-of-discrepancy protocol is described in Section~\ref{sec:proof-of-discrepancy} (and, where relevant, Section~\ref{sec:pod-protocol}) in general terms
applicable to any physical device fleet satisfying the stated
conditions (a predictive model accepting device configuration
representations, physical devices with manufacturing stochasticity,
and a measurement interface through which challenges can be executed
and responses recorded).  The following are non-limiting examples of
specific embodiments:

\paragraph{Optical scattering embodiment.}
In some embodiments, the devices are optical scattering systems
comprising a spatial light modulator or structured illumination source,
a scattering medium (for example a diffuser, multimode fibre,
integrating sphere, or disordered volumetric medium), and a detector
array.  Challenges comprise illumination patterns projected onto the
scattering medium.  Responses comprise the detected intensity
patterns.  The predictive model predicts detected intensity patterns
from illumination patterns and device configuration representations.
Manufacturing stochasticity in the scattering medium's microstructure
provides physical unclonability.

\paragraph{Projector--camera embodiment.}
In some embodiments, the devices are \RK modules as
described elsewhere in this specification, operated in a regime where
the scene is replaced by a characterisation target or where the
reactor's internal scattering properties are the subject of the
predictive model.  The \cb provides the
challenge--response record.  Fleet calibration via the
proof-of-discrepancy protocol improves the fidelity of the
fleet-shared surrogate model used for simulation, training, and
cross-device transfer.

\paragraph{Phonon and acoustic embodiment.}
In some embodiments, the devices are acoustic or phononic scattering
systems in which challenges comprise acoustic excitation patterns and
responses comprise measured acoustic fields.  In some embodiments,
the devices are SASER (sound amplification by stimulated emission of
radiation) systems or phonon laser systems, where coherent phonon
generation provides structured excitation and scattering through
disordered solid-state or fluid media provides the high-DOF
challenge--response mapping.

\paragraph{Generalisation condition (non-limiting).}
The proof-of-discrepancy protocol is applicable to any physical device
fleet where: (a)~individual devices have more internal degrees of
freedom than are accessible through the measurement interface,
(b)~multiple independently manufactured instances exist, (c)~a
submitted experimental method can be executed on other instances with
comparable results modulo device-specific variation, and (d)~device-to-device
variation is large enough for distinct device identities yet
structured enough to be separable from shared model error.

\paragraph{Terminology correspondence (non-limiting).}
In the context of \RK devices as described elsewhere in
this specification, the challenge configuration $X$ corresponds to a
control protocol or protocol seed specifying the emission schedule and
associated parameters; the measured response $Y_i(X)$ corresponds to
the \cb or a derived summary thereof.  Specifically,
$Y_i(X)$ denotes a summary statistic or derived quantity computed from
the observation sequence $\{\mathbf{y}_t\}$ produced by device~$i$
under control protocol~$X$.  The device
configuration representation $\mathbf{d}_i$ corresponds to the
device's position in the fleet-calibrated latent space described in
the Fleet Latent Space Calibration subsection.  The proof-of-discrepancy
protocol operates over the same evidence structures (committed
\cbs, protocol digests, and meter summaries) used
elsewhere in this specification.

% ----------------------------------------------------------------------
\subsection{Relationship to adjacent frameworks (non-limiting)}
% ----------------------------------------------------------------------

\paragraph{Distinction from physical proof-of-work.}
Physical proof-of-work protocols use physical processes to replace
computational hash puzzles for blockchain consensus.  The work product
in such protocols is the evaluation of a physical one-way function---a
solved puzzle.  The proof-of-discrepancy protocol differs in that the
work product is a reproducible experimental method demonstrating where
a predictive model is wrong.  The objective is model improvement
through error discovery, not consensus through puzzle-solving.

\paragraph{Distinction from adversarial benchmarking.}
Adversarial benchmarking platforms reward participants for discovering
inputs that cause computational models to produce incorrect outputs.
The proof-of-discrepancy protocol differs in that: (a)~the proof
object is a reproducible physical experimental method, not a digital
input example; (b)~in preferred embodiments, verification includes
independent physical re-execution on distinct devices, not
computational re-evaluation or human labelling; and (c)~the devices exhibit physical unclonability
that raises the empirical cost of result fabrication under the declared attacker model.

\paragraph{Distinction from metrological inter-comparison.}
Inter-laboratory comparisons and proficiency testing schemes use
independent physical measurement as a validation mechanism.  The
proof-of-discrepancy protocol differs in that: (a)~participants are
incentivised to actively discover model errors, not merely to
demonstrate measurement competence; (b)~each device is physically
distinct and unclonable, rather than measuring a circulated reference
artefact; and (c)~the protocol treats model--reality discrepancy as
a multi-functional resource rather than as a deficiency to be
corrected.

\paragraph{Distinction from microstructure-based authentication protocols.}
PUF-based blockchain and authentication systems use manufacturing
stochasticity for device identity and key generation.  The
proof-of-discrepancy protocol differs in that: (a)~model discrepancy
itself is the valued protocol output, not merely a security property;
(b)~devices participate in a model improvement loop, not only an
authentication loop; and (c)~verification is by physical reproduction
of a submitted method, not by challenge--response authentication
against enrolled templates.

\paragraph{Distinction from calibration compliance systems.}
Calibration-enforcement blockchain systems (for example
metrology-aware consensus protocols) embed calibration compliance
into validation logic.  The proof-of-discrepancy protocol differs in
that: (a)~the protocol rewards discovery of where models fail, not
attestation that devices are calibrated; and (b)~verification is by
physical re-execution, not by cryptographic certificate checking.

\paragraph{Distinction from sensor-fusion and shared-scene witnessing
systems.}
Multi-camera, multi-sensor, and Bayesian-fusion systems (including
Kalman-filter, particle-filter, interacting-multiple-model, and
factor-graph fusion frameworks) combine multiple sensor measurements
into a single fused estimate of an underlying state. Multi-RK
corroboration as described in subsequent sections differs in that:
(a)~the load-bearing output is a set of independently committed
\cbs, each retaining its own protocol digest, meter
envelope, and microstructure-bound provenance, rather than a fused
estimate that collapses per-sensor provenance into a single state
estimate; (b)~cross-kernel agreement is evaluated at the meter-score
and discriminator level (verisimilitude, hardness, freshness,
projection-hardness, microstructure identity) rather than at the
raw-observation level, so cross-kernel comparison does not require
and does not depend on a shared measurement-noise model;
(c)~each kernel's evidence is independently auditable and selectively
disclosable under the selective-disclosure policy, supporting
verifier-side reconstruction of the cross-kernel argument without
exposing fused intermediate states; and (d)~corroboration weights are
bound to declared trust and reputation parameters committed in the
protocol digest, rather than derived from estimator-internal innovation
statistics or shared-state covariance. Where a fused state estimate is
required for downstream use, a sensor-fusion estimator may be composed
as a derived output downstream of multi-RK corroboration; the fused
estimate is in such embodiments a meter-conditioned derived quantity
and is not a substitute for the underlying committed bundles.

\paragraph{Distinction from distance-bounding and continuous
authentication protocols.}
Distance-bounding protocols (in the tradition of Brands and Chaum
1993 and Hancke and Kuhn 2005) establish that a prover is within a
bounded physical distance of a verifier by measuring the round-trip
time of a cryptographic challenge--response exchange. Continuous-
authentication systems extend this with repeated freshness challenges
over a session. \RK freshness mechanisms differ in that:
(a)~the load-bearing event is an attested physical light--matter
interaction under a declared protocol, recorded as a committed
\cb, rather than a location bound or a session-presence
proof; (b)~freshness is established by time-anchored protocol-digest
commitment chained to external time anchors and to prior committed
bundles, rather than by round-trip timing of a cryptographic primitive;
and (c)~the apparatus and a distance-bounding or continuous-
authentication mechanism are not mutually exclusive: in some
embodiments, such a mechanism is composed alongside the apparatus and
its bounded-distance or session-presence attestation is included in
the protocol digest as auxiliary evidence, without becoming the
load-bearing primitive.

% ======================================================================
\section{Networked Reality Kernels and Transitive Proof-of-Projection}
\label{sec:networks}
% ======================================================================

This section describes embodiments in which multiple \RK
modules operate in a network, exchange evidence derived from
the \cb, and perform cross-verification of shared events. These
embodiments are optional. A single \RK module can be
implemented and used without any networked operation.
In some embodiments, networked modules additionally participate in the
fleet model improvement protocol described in
Section~\ref{sec:proof-of-discrepancy}, exchanging evidence of model
discrepancies and incorporating confirmed corrections into the shared
predictive model.

\paragraph{Scaling regime and limitations (non-limiting).}
In some embodiments, networked deployments use hierarchical aggregation,
committee-based attestation, neighbourhood-local trust updates, secure
aggregation of population statistics, and amortised monitoring to bound
communication overhead while preserving auditable cross-device consistency.
Illustrative deployments may range from local clusters to larger federated
fleets, with committee size, aggregation depth, witness cadence, and
overlap policy selected per target latency, fault model, and deployment
scale. Layered aggregation may be used to compress cross-device comparisons
into auditable summaries while preserving selective-opening pathways for
deeper review where required.
\subsection{Network model and shared channel}

Let $I$ index a finite set of \RK modules. For each $i\in I$,
let $\theta_i$ denote configuration, $u^{(i)}(t)$ denote per-step
control, and $C^{(i)}_{0:T}$ denote the \cb. In some
embodiments, devices share a physical channel through a shared
environment. A shared world state $W_t$ represents external scene
factors accessible to multiple devices:
\[
  W_{t+1} \sim \mathsf{P}^{\mathrm{net}}_{\mathrm{world}}\!\Bigl(
  \cdot \,\Big|\, W_t,\,
  (u^{(i)}(t))_{i\in I},\,u_E(t)\Bigr),
\]
where $u_E(t)$ denotes adversarial co-illumination when present.
The symbol $W_t$ denotes the shared world state and is distinct from
the per-device extended state $X_t$ defined in
Section~\ref{sec:definitions}; $X_t$ includes device-internal
components (reactor state, media state, controller state) that are
not shared across devices.

\subsection{Distributed time and spatiotemporal consistency}
\label{sec:distributed-time}

In some embodiments, spatiotemporal constraints on inter-device causal
links serve as physical timing anchors that bind a run to measurable
timing relationships. These function as optical-delay logical clocks:
physically grounded partial orders derived from propagation and
hardware-latency bounds. A non-limiting spatiotemporal violation
functional is
\[
  \nu^{(ij)}_{\mathrm{st}}
  =
  \left(\frac{d^{\mathrm{path}}_{ij}}{\Delta\tau_{ij}} - c\right)^{2}
  + \lambda_{\mathrm{lat}}\bigl(\Delta\tau_{ij}
  - \tau_{\min}\bigr)_{-}^{2}.
\]
where $(x)_{-} := \min(x, 0)$ denotes the negative-part operator
(and $(x)_{+} := \max(x, 0)$ denotes the positive-part operator,
used later in Yoked objectives).

\subsection{Projection hardness and proof-of-projection}

For devices $i,j\in I$, pairwise projection hardness is defined as
\[
  k^{\ast}_{\mathrm{proj}}(i\to j)
  :=
  \min\left\{
    k:\; D\!\left(\mathsf{P}_{\theta_j}^{(j)},\,
    \hat{\mathsf{P}}^{(i\to j,k)}\right)\le \varepsilon
  \right\},
\]
where $\hat{\mathsf{P}}^{(i\to j,k)}$ is an emulator class indexed
by capacity $k$. A committee $C(j)\subseteq I$ yields an aggregate
projection score
$H_{\mathrm{proj}}(j) := \min_{i\in C(j)}
k^{\ast}_{\mathrm{proj}}(i\to j)$.

\paragraph{Terminological note.}
\label{par:proof-of-projection-empirical}
The term \emph{proof-of-projection} is used throughout this document
as a term of art denoting an empirically calibrated attestation
protocol, not a cryptographic proof system with formal
soundness and completeness properties.  What the protocol provides
is: (i)~committed evidence that a physical device produced a
\cba trace under declared conditions,
(ii)~independent physical realisation on one or more distinct
unclonable devices in the verifier role, scored by projection
hardness against the fleet projection-hardness statistics, and
(iii)~auditability via selective disclosure, quorum verification,
and the submission record's commitment binding.
It does not provide, and should not be read as asserting, a
zero-knowledge or interactive proof in the cryptographic sense.
Achieving formal proof-system guarantees would require additional
identity, Sybil-resistance, and adversary-model assumptions beyond
those specified here.  The ``transitive'' extension
(Section~\ref{sec:networks}) inherits this character:
transitive proof-of-projection denotes a chain of empirical
attestations with auditable trust-weight updates, not a
transitively composable formal proof.

\subsection{Challenge--commit--open protocols for selective disclosure}
\label{sec:two-seed-protocol}

In some embodiments, devices commit to the \cb and later
selectively open subsets for audit. Per-atom digests are computed and
aggregated into a Merkle root $R^{(i)}$, binding emission and
observation windows:
\[
  \mathrm{com}^{(i)}_{\mathrm{full}}
  = H\!\left(\mathrm{rid}\,\|\,i\,\|\,\Pi_{\mathrm{dig}}\,\|\,
  \mathbf{m}^{(i)}\,\|\,R^{(i),\mathrm{obs}}\,\|\,
  R^{(i),\mathrm{emit}}\right).
\]

\paragraph{Canonical two-seed selective-opening protocol (non-limiting).}
\label{par:two-seed}
In some embodiments, a two-seed structure separates the run commitment
from the opening decision so that the challenge is structured to
reduce the opportunity for anticipatory tailoring to the specific
atoms that will later be opened.  A non-limiting protocol
proceeds as follows:
\begin{enumerate}
  \item \textbf{Run phase.}  The device executes a protocol run under
    run seed $s_{\mathrm{run}}$, derived from the protocol digest
    and a nonce committed before the run begins.  The run seed
    determines the emission schedule and scan law; it is logged in
    the protocol digest and not disclosed until step~4 (selective
    disclosure).
  \item \textbf{Commitment.}  Upon run completion, the device computes
    and publishes the Merkle commitment $\mathrm{com}^{(i)}_{\mathrm{full}}$
    above.  The run seed $s_{\mathrm{run}}$ remains withheld at this
    stage.
  \item \textbf{Challenge (opening seed).}  A verifier (or a randomised
    beacon in some embodiments) supplies an opening seed
    $s_{\mathrm{open}}$, derived independently of the run seed.
    $s_{\mathrm{open}}$ determines which atoms or windows are to be
    opened.
  \item \textbf{Selective disclosure.}  The device reveals the
    requested atoms together with their Merkle inclusion proofs and
    the run seed $s_{\mathrm{run}}$.  The verifier checks that
    (a)~the opened atoms are consistent with the committed Merkle root,
    (b)~the protocol digest reconstructed from $s_{\mathrm{run}}$ matches
    the committed digest, and (c)~meter summaries lie within declared
    envelopes.
  \item \textbf{Audit escalation (optional).}  Under policy-triggered
    escalation, additional atoms may be opened using further opening
    seeds supplied by an auditor.
\end{enumerate}
The two-seed structure is referenced throughout this specification
wherever ``two-seed opening'', ``run seed'', or ``opening seed'' appear;
those usages refer to this protocol.  In embodiments using this
structure, the run commitment is bound to the committed evidence before
the opening seed is known to the device, reducing the opportunity for
anticipatory tailoring of the run to match future opening challenges.

A non-limiting protocol flow summary: run phase (run seed determines
emission schedule and scan law), commitment (Merkle roots and
signatures published), challenge (opening seed supplied by verifier
or beacon), selective disclosure (requested atoms with Merkle
inclusion proofs and run seed), and optional audit escalation.

\subsection{Transitive proof-of-projection and reactor pairing}

In some embodiments, devices perform continuous or periodic pairing
procedures. A non-limiting continuous pairing latent
$\sigma_{ij}(t)$ is updated as
\[
  \dot{\sigma}_{ij}(t)
  = -\kappa_\sigma\,\sigma_{ij}(t)
    + F_{\mathrm{pair}}\bigl(\mathbf{y}^{(i)}(t),\,
    \mathbf{y}^{(j)}(t),\,u^{(i)}(t),\,u^{(j)}(t)\bigr).
\]
A tPoP record for an event includes run identifiers, protocol digests,
commitment references for observation and emission windows,
quorum-contributed meter summaries, and pairing latents.

\paragraph{Layered RK networks (non-limiting).}
In some embodiments, a network is organised into layers with a rolling
amortised monitor (where $\eta$ in this equation is a layer-level
smoothing coefficient, not the scan-law parameter or learning rate):
\[
  \overline{H}_{\ell}(t)
  = (1-\eta)\,\overline{H}_{\ell}(t-1)
    + \eta\,\mathrm{Agg}\bigl(\{H_{\mathrm{proj}}(j;t)\}_{j\in
    \mathcal{L}_\ell}\bigr).
\]

\paragraph{Direct optical coupling between reactors (non-limiting).}
In some embodiments, neighbouring devices are coupled optically at the
reactor level through free-space apertures, fibre links, or a shared
lightguide backplane.

\subsection{Trust weights and reputation updates}

\subsubsection{Cross-prediction and compatibility}

In an overlap window in which devices $i$ and $j$ observe the same
event under related protocols, device $i$ may attempt to predict a
view of $j$'s record (or summary) from its own record and the logged
protocols. For example, device~$i$ may form a cross-prediction
\[
  \widehat{C}^{(j)}_{t_0:t_1}
  = \mathcal{E}_{i\rightarrow j}\!\bigl(C^{(i)}_{t_0:t_1},\,
    U^{(i)}_{t_0:t_1},\,U^{(j)}_{t_0:t_1}\bigr),
\]
where $\mathcal{E}_{i\rightarrow j}$ is a learned emulator, calibrated
geometric transfer, or other cross-device predictor, and $t_0{:}t_1$
denotes the overlap window.

A non-limiting per-event loss for the directed edge
$(i\rightarrow j)$ is
\[
  \ell_{i\rightarrow j}(e)
  = \mathrm{clip}_{[0,1]}\!\left(
    \frac{D\!\left(\widehat{C}^{(j)}_{t_0:t_1},\,
      C^{(j)}_{t_0:t_1}\right)}{\tau}\right),
\]
where $D(\cdot,\cdot)$ is any non-limiting discrepancy (for example, a
meter-weighted feature distance, a negative log-likelihood under a
fitted noise model, or a learned perceptual distance), and $\tau>0$ is
a scale parameter that controls sensitivity.  To preserve the meter
partition (Section~\ref{sec:security-theory}), the discrepancy $D$ is
computed using acceptance meters $\mathbf{m}^{\mathrm{accept}}$ in embodiments using the stated meter partition;
evaluation meters are excluded from trust-weight computation because
trust weights feed into quorum selection and thereby into the
optimisation loop.

\subsubsection{Exponential-weights web-of-trust update (non-limiting)}

In some embodiments, device~$i$ maintains nonnegative trust weights
$\{W_{i\rightarrow k}(t)\}_{k\in\mathcal{N}_i}$ over a neighbourhood
$\mathcal{N}_i$ of peers. After an event~$e$ (or batch of events) the
weights may be updated by an exponential-weights rule:
\[
  \widetilde{W}_{i\rightarrow j}(t{+}1)
  = W_{i\rightarrow j}(t)\exp\!\bigl(-b\,\ell_{i\rightarrow j}(e)
    \bigr),
  \qquad
  W_{i\rightarrow j}(t{+}1)
  = \frac{\widetilde{W}_{i\rightarrow j}(t{+}1)}
    {\sum_{k\in\mathcal{N}_i}
      \widetilde{W}_{i\rightarrow k}(t{+}1)},
\]
where $b>0$ controls update aggressiveness.

In some embodiments, stabilisers are applied, for example:
(i)~flooring and renormalisation
$W_{i\rightarrow j}\leftarrow \max(W_{i\rightarrow j},w_{\min})$,
(ii)~mixing with a uniform prior
$W\leftarrow (1-\mu_{\mathrm{mix}})W+\mu_{\mathrm{mix}}/|\mathcal{N}_i|$ for $\mu_{\mathrm{mix}}\in(0,1)$, or
(iii)~clipping to prevent single-event collapse.

The resulting weights define a continuous web-of-trust used to select
quorums, weight cross-signatures, schedule higher-cost challenges, or
gate access to high-value transitive proof-of-projection records.

\subsubsection{Bayesian trust update (non-limiting)}

In some embodiments, trust and reputation are updated using
Bayesian-style posteriors over a peer's reliability, optionally
conditioned on operating regime and meter envelope.

\paragraph{Binary audit outcomes (Beta--Bernoulli, non-limiting).}
In some embodiments, agreement for a peer~$j$ on an event~$e$ is
reduced to a binary variable $a_{j,e}\in\{0,1\}$. With a Beta prior
$\pi_j\sim \mathrm{Beta}(a_0,b_0)$, the posterior after
observing outcomes $a_{j,1:n}$ is
\[
  \pi_j \mid a_{j,1:n}
  \sim \mathrm{Beta}\!\Bigl(a_0+\textstyle\sum_{e=1}^n a_{j,e},\;
  b_0+\textstyle\sum_{e=1}^n (1-a_{j,e})\Bigr).
\]
A posterior mean, quantile, or lower confidence bound may be used as a
reputation score.

\paragraph{Forgetting and regime conditioning (non-limiting).}
In some embodiments, older evidence is downweighted using exponential
forgetting:
\[
  S_j(t{+}1) = \gamma S_j(t) + a_{j,t}, \qquad
  F_j(t{+}1) = \gamma F_j(t) + (1-a_{j,t}),
\]
with $\gamma\in(0,1)$ (a forgetting factor; not the scan-law function
$\gamma(t;\eta)$), yielding a posterior of the same Beta form. In
some embodiments, separate reliabilities are maintained per regime, per
meter envelope, or per role, and combined by a policy-defined
aggregation.

\paragraph{Continuous compatibility scores (non-limiting).}
In some embodiments, continuous losses or residual scores are used
instead of binary variables. The loss
$\ell_{i\rightarrow j}(e)$ may be mapped into a pseudo-likelihood, or
discretised into bins (pass, marginal, fail), to update a posterior
over reliability. In some embodiments, pairing latents
$\sigma_{ij}(t)$ and layer-level rolling monitors modulate how
strongly evidence updates trust.

\subsection{Additional network mechanisms (non-limiting)}

\paragraph{Optical delays as physical logical clocks (non-limiting).}
In some embodiments, controlled delay elements impose causal ordering
constraints that function as physical logical clocks. A challenge
emitted by one device may be routed through a known delay element (for
example a fibre delay line) before reaching another device, and devices
log timestamps and \cba evidence for the corresponding
response. The known delay bounds constrain allowable histories and help
detect fabricated ordering, supplementing or replacing reliance on
external global time services.

\paragraph{Phase-locked and frequency-synchronised networks
(non-limiting).}
In some embodiments, multiple devices synchronise to a common frequency
reference using phase-locked loops, injection locking, or consensus
timing. Synchronisation may enable coherent measurements across
devices, distributed interferometry, coordinated emissions with
controlled interference, and tighter spatiotemporal consistency checks
without external clocks.

\paragraph{Near-field resonant coupling between reactors (non-limiting).}
In some embodiments, devices exchange energy or information through
near-field coupling without direct wired or free-space optical links,
for example through evanescent coupling, shared resonant cavities, or
magnetic coupling between microstructured media.

\paragraph{Entrainment and collective dynamics (non-limiting).}
In some embodiments, networks exhibit emergent collective dynamics such
as synchronisation, correlated oscillations, or collective stability
modes. These dynamics may be logged and used as anomaly indicators,
additional hardness constraints, or distributed timing signals.

\paragraph{Cross-attestation and overlap-window protocols (non-limiting).}
In some embodiments, network integrity is strengthened by
cross-attestation: multiple kernels observe a shared scene interval
and produce mutually constrained evidence. A non-limiting procedure
comprises: (i)~each node records a committed observation window under a
logged protocol digest; (ii)~nodes exchange compact digests (window
roots, timing-anchor summaries, meter vectors) over authenticated
channels; (iii)~a verifier checks overlap-window consistency and
accepts when the declared quorum condition (at least $k$ of $n$ witnesses) is met under the stated policy. In some embodiments, cross-attestation is strengthened by
structured emissions that are jointly scheduled or beacon-seeded.

\paragraph{Enrolment, identity binding, and certificates (non-limiting).}
In some embodiments, a device participates in a network only after
enrolment. Enrolment binds an identity to a physical device profile
and an initial policy regime. A non-limiting procedure comprises
identity assignment, profiling (calibration challenges and baseline
recording), certification by an anchor node, and anchoring to a
transparency log.

\paragraph{Revocation, quarantine, and trust lifecycle (non-limiting).}
In some embodiments, network trust is managed over a device lifecycle
that includes drift, re-enrolment, and revocation. Devices whose
hardness or reliability drops below policy thresholds may be
quarantined (reduced weight, restricted participation) pending
re-profiling.

\paragraph{Privacy and unlinkability for network witnesses (optional,
non-limiting).}
In some embodiments, network participation supports privacy by limiting
disclosure and reducing linkability. Non-limiting mechanisms include
pseudonymous identities, group signatures or anonymous credentials for
membership proofs, and secure aggregation of meter outputs so that only
quorum aggregates are disclosed. In some embodiments, overlap-window
corroboration is performed on committed digests rather than raw
the \cb, with selective opening triggered under audit-escalation policy.

\paragraph{Policy versioning and audit escalation (non-limiting).}
In some embodiments, network attestation policy is expressed as signed,
versioned policy objects that specify meter thresholds, admissible
protocol families, quorum requirements, disclosure rules, and
revocation procedures. Policy identifiers are bound into protocol
digests and signatures so that downgrades or equivocation are
detectable. Audit escalation policies include requesting additional
openings, requesting independent witnesses, switching to higher-cost
protocols, and temporarily disabling high-impact actions until
discrepancies are resolved.

\paragraph{Non-limiting deployment topologies.}
In some embodiments, networked kernels are deployed as mobile fleets,
fixed installations, environmental beacon networks, tiled immersive
chambers, or hybrid topologies. Topology and protocol selection may
adapt to deployment context under policy constraints.

\paragraph{Approximate physical common knowledge (non-limiting).}
In some embodiments, the committed and multiply-witnessed \cbs implement a form of approximate common knowledge: multiple
bounded-perspective agents can verify not only that an event occurred but
that they share the same optically grounded evidence of it, and that
each knows the others share it, up to the temporal and perceptual
granularity of the commitment protocol. Strict common knowledge in the
Aumann sense is unattainable in any practical system with communication
latency or perceptual uncertainty; however, the physically committed
evidence provides a stronger epistemic foundation than purely digital
broadcast, because each agent's committed \cb is bound to
a device-specific signature that other agents can verify, and the
temporal overlap and spatial co-observation constraints make it
empirically difficult for any subset of agents to fabricate a consistent
multi-perspective record. In some embodiments, the level of approximate
common knowledge is parameterised by the commitment protocol's temporal
resolution, the number of independent witness devices, and the spatial
co-observation constraints, and is reported as a meter-conditioned
confidence level rather than as a binary common-knowledge predicate. In
some embodiments, this approximate common knowledge enables
agreement protocols without a trusted third party: agents who share
physically committed evidence of the same event can reach agreement by
exchanging posterior beliefs, because the common evidence base is
physically assured to within the protocol's temporal and spatial
granularity rather than assumed.

\subsection{Byzantine adversaries and robustness (optional)}

In some embodiments, under resource bounds on the adversary and a
sufficient fraction of honest, high-hardness witnesses with diverse
viewpoints, coordinated forgery becomes difficult because fabricated
records would need to remain consistent with many independent, high-weight
committed logs.

\paragraph{Non-limiting adversary modes.}
In some embodiments, adversaries act not only by forging traces or
tampering with a single device, but by exploiting network structure,
quorum selection, selective disclosure, and cross-device trust
updates. Non-limiting adversary modes include:
\begin{enumerate}
  \item Collusion and Sybil participation to bias quorums or
    trust-weight updates.
  \item Eclipse or routing attacks that isolate a device or force it
    to interact with a compromised neighbourhood.
  \item Replay and timing manipulation, including attempts to re-use
    old commitments or to desynchronise protocol digests.
  \item Selective opening manipulation, including attempts to
    anticipate which atoms will be audited, or to craft atoms that
    pass shallow checks while failing globally.
  \item Audit flooding, in which an adversary forces repeated
    escalations to exhaust bandwidth or disclosure budgets.
  \item Cross-prediction poisoning, in which an adversary crafts
    emissions or summaries to distort cross-device residual estimates.
  \item Layer targeting, in which higher-layer aggregators are biased
    even if local witnesses are honest.
  \item Policy manipulation, including attempts to change
    thresholds, decay rates, or quorum eligibility rules.
\end{enumerate}

\paragraph{Non-limiting defences.}
Defences include:
\begin{enumerate}
  \item Diversity requirements for quorums (for example spatial
    diversity, spectral diversity, hardware lineage diversity, or
    operator independence).
  \item Trust-weighted quorum selection with rate limits and minimum
    diversity constraints.
  \item Delayed opening via a two-seed structure (run seed versus
    opening seed) to reduce precommitment attacks.
  \item Anti-replay binding to run identifiers, protocol digests,
    timing anchors, and freshness policies.
  \item Multi-layer monitoring with amortised background scoring and
    escalation only on deviations beyond policy thresholds.
  \item Audit budgets and escalation policies that bound disclosure
    and computational cost.
  \item Cross-checks that tie reported meter outputs to peer
    evaluations on the same shared event.
  \item Periodic re-seeding and neighbourhood reshuffling to reduce
    long-lived eclipse conditions.
\end{enumerate}

In some embodiments, audit outcomes and detected inconsistencies feed
back into trust weights and reputations, reducing future influence of
suspected peers and triggering higher-cost challenges only when policy
thresholds are exceeded.

% ======================================================================

% ======================================================================
%
%
%
% ======================================================================

\section{Routed Physical Ensembles and Optical
  Mixture-of-Experts}
\label{sec:routed-ensembles}
% ======================================================================

% ------------------------------------------------------------------
\subsection{Expert subunits, structural diversity, and route binding}
\label{sec:expert-subunits}
% ------------------------------------------------------------------

In some embodiments, a routed physical ensemble comprises multiple
physically distinct expert subunits. In some embodiments, each subunit
includes a scattering medium or related reactor, associated input and
output optics, and a readout configured to extract task-relevant
structure from a high-dimensional physical transform. In some
embodiments, each expert is itself a parameterised Markov kernel within
the broader \RK framework.

In some embodiments, each expert exhibits structural diversity due to
microstructure-based uniqueness of its microstructure: no two expert subunits
compute the same transformation. In some embodiments, structural
diversity arises from manufacturing variability but task complementarity is an
empirical property that must be validated by comparing multi-expert
routed performance against single-expert and random-routing baselines
on specific benchmarks.

In some embodiments, expert subunits are not limited to disordered
scattering media. In some embodiments, expert subunits include
phase-change reactors whose transfer function is switched between
amorphous and crystalline states, gain-pumped optical cavities whose
mode structure provides a distinct transform per cavity geometry,
photonic integrated circuits whose fabrication variation produces
device-specific transfer matrices, fibre-delay networks whose delay
and coupling configuration defines a recurrent temporal kernel, and
acoustic or multi-physics reactors as described in
Section~\ref{sec:embodiments}. In some embodiments, a
heterogeneous ensemble mixes expert types --- for example scattering
media, phase-change elements, and PIC subunits --- to provide
diversity along different physical axes. In some embodiments, the
commitment, provenance-binding, and calibration-state attestation
layer described in this section applies uniformly to all expert types,
because each expert is a parameterised Markov kernel regardless of its
physical substrate. In some embodiments, the routing interface ---
including coupling optics, mode matching, and per-expert calibration
--- varies by expert type and is determined per embodiment.

In some embodiments, routing identities, expert microstructure identities,
calibration-state hashes, and readout or weights version identifiers
are committed to the protocol digest and bound into provenance records,
so that every computation is traceable to the specific physical path
through which it was executed.

\paragraph{Distinguishing feature (non-limiting).}
In some embodiments, the routed physical ensemble is distinguished
from prior optical reservoir computing, photonic neural network, and
conditional computation architectures by the integration of committed
physical evidence into the routing and computation path. In some
embodiments, the selected physical path through the ensemble, the
microstructure identity of each participating expert, the calibration state of
each expert, and the execution transcript are committed to the
protocol digest and bound into provenance records, so that every
routed computation is traceable, auditable, freshness-anchored, and
integrated into the proof-of-discrepancy, bilateral portability, and
fleet corroboration infrastructure described elsewhere in this
specification. In some embodiments, this commitment and auditability
layer is the architectural distinction, not the optical routing
mechanism itself.


% ------------------------------------------------------------------
\subsection{Routing, addressing, and gating}
\label{sec:routing}
% ------------------------------------------------------------------

In some embodiments, routing is performed by an addressing subsystem
that may include acousto-optic deflectors, electro-optic devices, MEMS
scanners, galvanometers, spatial light modulators, camera or
reprojection systems, or hybrid optical-digital routing mechanisms.

In some embodiments, a conservative optical embodiment uses a TeO$_2$
acousto-optic deflector with a beam diameter of approximately 5~mm,
acoustic velocity in the range of approximately 600--4\,200~m/s
depending on interaction geometry, access time in the range of
approximately 1--25~$\mu$s depending on beam diameter and acoustic
velocity, RF bandwidth of approximately 50--200~MHz,
approximately 50--200 resolvable steering positions (sufficient for
16--64 expert subunits with margin depending on configuration),
single-axis diffraction efficiency of approximately 60--85\%, and
effective random-access routing rate of approximately
40~kHz--1~MHz depending on access time. In some embodiments, two cascaded deflectors
provide two-axis steering with approximately 36--64\% throughput.

In some embodiments, a relay lens between the AOD output and the
expert subunit array maps AOD angles to positions on the expert
inputs. In some embodiments, coupling optics between each expert and
the shared output path collect scattered or transmitted light and
direct it to a shared detector or detector array.

In some embodiments, routing decisions are made by a gating network
that may be fully digital (a learned gating head), fully optical
(routing based on optical preprocessing), or hybrid
(optical features consumed by a digital gating head). In some
embodiments, routing identities and schedules are logged in the
protocol digest and bound into provenance and audit records.

\paragraph{Architectural constraint (non-limiting).}
In some embodiments, the system logs sufficient information to
reconstruct which expert subunit(s) were selected, the routing
parameters, and the calibration-state hash of each participating
expert, so that downstream relying parties can verify or audit the
physical computation path. This architectural constraint distinguishes
the routed ensemble from an opaque black-box accelerator: every
computation is traceable to a specific physical path under a committed
protocol.

\paragraph{Candidate relay layout (non-limiting).}
In some embodiments, the relay layout includes two cylindrical relay
lenses (one per axis), a field lens at an intermediate plane, and a
micro-relay array aligned with the expert input facets. In some
embodiments, practical alignment tolerance requires active calibration
using the same calibration-loop infrastructure disclosed in this
specification.

% ------------------------------------------------------------------
\subsection{Gain-assisted short-chain coupling, failover, and
  inter-expert relay}
\label{sec:gain-coupling}
% ------------------------------------------------------------------

In some embodiments, the per-stage optical loss budget for a routed
ensemble limits the number of concatenated expert stages in a single
forward pass, because loss accumulates multiplicatively across stages,
lens surfaces, deflector passes, and coupling into and out of each
expert. In some embodiments, a loss budget analysis includes at least
deflector efficiency per pass, relay-lens reflection losses,
coupling-into-expert loss, scattering transmission per expert,
coupling-out-of-expert loss, and detector quantum efficiency.

In some embodiments, a semiconductor optical amplifier (SOA) or other
gain medium is inserted between expert stages to partially compensate
inter-stage loss and extend the usable chain length beyond what
passive optics alone would permit. In some embodiments, the SOA
contributes amplified spontaneous emission (ASE) noise that reduces
SNR, and the useful operating point is determined by the SOA gain,
noise figure, and accumulated ASE across the chain length, so that
there exists a practical maximum chain length beyond which additional
stages degrade rather than improve task performance.

In some embodiments, SOA gain is included in the per-stage loss budget.
In some embodiments, a two-stage budget template is used:

\begin{table}[ht]
\centering
\small
\caption{Per-stage budget template for gain-assisted inter-expert
  coupling (representative values; actual values are
  embodiment-specific).}
\label{tab:gain-budget}
\begin{tabular}{lcc}
\toprule
Stage element & Typical range & Notes \\
\midrule
AOD diffraction (per pass)   & 60--80\% & TeO$_2$, single axis \\
Relay lens (per surface)     & 97--99\% & AR-coated \\
Expert coupling (in + out)   & 10--50\% & geometry-dependent \\
Expert scattering transmission & 1--30\% & medium-dependent \\
SOA net gain (if present)    & 10--20\,dB & minus NF penalty \\
Detector QE                  & 40--90\% & wavelength-dependent \\
\bottomrule
\end{tabular}
\end{table}

In some embodiments, whether SOA-assisted chaining produces measurably
better task performance than single-expert operation is an empirical
question that must be established per configuration and per task by
comparison against single-expert and random-routing baselines.

\paragraph{State augmentation for gain-coupled operation
  (non-limiting).}
In some embodiments, when SOA gain is used between expert stages, the
extended state tracked by the controller is augmented to include SOA
bias current, estimated ASE power, SOA temperature, and gain
saturation state. In some embodiments, these parameters are logged in
the protocol digest and metered within declared envelopes. In some
embodiments, the SOA gain profile drifts with temperature and aging,
requiring periodic recalibration using the same drift-monitoring
infrastructure described elsewhere in this specification. In some
embodiments, thermal management and power-supply stability for the SOA
are engineering requirements comparable in kind (though not in scale)
to those for the primary emitter and are addressed by standard
optoelectronic packaging practice.

In some embodiments, failover modes include rerouting to alternative
expert subunits, falling back to single-expert operation, or reverting
to digital-only processing when optical chain integrity degrades below
a declared meter threshold. In some embodiments, failover decisions
are logged and auditable.

% ------------------------------------------------------------------
\subsection{Worked dimensionality-tier routed-ensemble embodiment}
\label{sec:dim-tiers}
% ------------------------------------------------------------------

This subsection describes one conservative routed-ensemble embodiment
analysed at three dimensionality tiers, showing how routing policy
changes at each tier.

\paragraph{Tier~1 --- raw sensed and addressed channels.}
In some embodiments, the conservative AOD embodiment described above
addresses 64 expert subunits. In some embodiments, each subunit
receives input from a spatial light modulator with approximately
$10^6$ addressable pixels and produces output captured by a camera with
approximately $8 \times 10^6$ pixels. In some embodiments, Tier~1
dimensionality is therefore bounded by the hardware channel count:
approximately $10^6$ input dimensions per expert and approximately
$8 \times 10^6$ output dimensions per expert, with 64 experts
selectable at approximately 800~kHz random-access rate. In some
embodiments, at Tier~1 the routing policy selects among experts based
on input modality, task identity, or scheduling constraints, without
requiring detailed knowledge of each expert's computational properties.

\paragraph{Tier~2 --- approximately independent channels and effective
  rank.}
In some embodiments, the number of channels that carry approximately
independent information is substantially smaller than the raw hardware
count, because of \'etendue or space-bandwidth product constraints,
numerical aperture limits, oversampling, control modality (phase-only,
amplitude-only, or complex), and inter-pixel correlations. In some
embodiments, for typical multiply-scattering systems studied in
transmission-matrix experiments, Tier~2 dimensionality is
approximately $10^4$ to $10^6$ per expert, depending on geometry and
coherence. In some embodiments, at Tier~2 the routing policy
incorporates measured effective rank per expert and selects among
experts based on which experts offer the highest effective rank for the
current input domain, accounting for correlations that reduce
independent information content below the raw pixel count. In some
embodiments, Tier~2 dimensionality is measured experimentally using
singular-value decomposition of sampled transmission-matrix blocks or
related effective-rank estimation procedures, and the measured rank is
logged as a calibration parameter.

\paragraph{Tier~3 --- task-useful effective rank.}
In some embodiments, the number of dimensions that actually contribute
to performance on a specific task, after readout training, is smaller
than or equal to the Tier~2 effective rank and is task-dependent. In
some embodiments, at Tier~3 the routing policy selects among experts
based on measured task-relevant structure: which experts contribute
the most to task performance after readout training, as measured by
held-out validation on committed physical evidence. In some
embodiments, capability claims for the routed ensemble are tied to
measured Tier~3 structure rather than to bulk Tier~1 channel counts
or Tier~2 effective-rank estimates. In some embodiments, Tier~2
dimensionality is a necessary but not sufficient condition for
capability: it bounds the space of possible computations but does not
by itself demonstrate that any specific task benefits from the
available dimensions. In some embodiments, Tier~3 task-useful rank
must be established experimentally on declared benchmarks before
capability is asserted, and any claimed routing advantage is validated
against single-expert and random-routing baselines on specific
benchmarks using committed physical evidence.

\paragraph{Routing policy and tier interaction.}
In some embodiments, the routing policy determines which expert
subunits are selected for a given input, and the dimensionality tier
determines how the resulting output is characterised. In some
embodiments, at Tier~1 routing is coarse: expert selection is based on
availability, scheduling, or input modality. In some embodiments, at
Tier~2 routing becomes correlation-aware: experts are selected based
on measured effective rank and domain coverage, and redundant experts
may be bypassed. In some embodiments, at Tier~3 routing is
task-specialised: experts are selected and weighted based on their
demonstrated contribution to specific task performance, and the
routing policy itself is trained using committed physical evidence
from task evaluation. In some embodiments, any claimed routing
advantage is validated against single-expert and random-routing
baselines on specific benchmarks using committed physical evidence.
In some embodiments, whether routing at higher tiers produces
measurably better task performance is an empirical property
established per embodiment.

% ------------------------------------------------------------------

\subsection{Capacity accounting rule (non-limiting)}
\label{sec:capacity-accounting-rule}

For avoidance of ambiguity across this disclosure, six categories of
quantity are distinguished, and no claim herein shall conflate them:
(a) raw hardware channel count; (b) effective rank; (c) task-useful
rank; (d) control waveform coefficients; (e) simultaneously-measurable
gradient components; and (f) stable trained weights. No capability
claim in this disclosure shall assert a value in category (f) derived
from counting in categories (a), (b), (d), or (e). Any such claim shall
be supported by direct task-benchmark measurement of category (c) on
committed physical evidence and, where stable trained weights are
recited, by separate empirical stability evidence under the declared
drift and freeze-window criteria. Migration between categories requires
explicit experimental demonstration.

\subsection{Optional semantic extension branch and attested semantic
  computation (validation-dependent)}
\label{sec:semantic-branch}
% ------------------------------------------------------------------

In some embodiments, the routed physical ensemble further includes a
late, optional, validation-dependent semantic extension branch in
which text or other symbolic content is rendered as spatial optical
patterns, processed by selected expert subunits through the same
committed physical channel used for nonsemantic routed
computation, and decoded by trained readout heads, with digital fusion
remaining the default multimodal baseline.

In some embodiments, semantic capability in such an extension resides
in the encoder, readout, and any bias-mask parameters rather than in
the scattering medium alone. In some embodiments, the scattering
physics is identical to that used for nonsemantic computation; the
semantic branch is a domain-specific application of the same core
reservoir architecture to rendered textual or symbolic inputs. In some
embodiments, the default multimodal embodiment combines perceptual
features and semantic features in a digital fusion layer, and that
digital fusion path serves as the reference baseline against which any
more physical comparison mechanism is evaluated.

\paragraph{Attested semantic computation record (non-limiting).}
In some embodiments, the system produces an attested semantic
computation record comprising: a committed \TB bundle
identifier, a selection-rule digest, a protocol digest, a routing log
identifying which expert subunits were traversed and in what order,
microstructure identities of participating experts, a calibration-state hash, a
readout or weights version identifier, a time anchor, the semantic
output, and any confidence, discriminant, or quality statistics
associated with the output. In some embodiments, hash-chain linkage
and external time anchoring make the semantic run order-sensitive and
freshness-sensitive, so that a semantic conclusion cannot be fabricated
and back-filled without breaking the chain. In some embodiments, the
physical channel is rate-limited, so that the attestation object
resists digital-speed brute-force generation or transcript farming.

\paragraph{Worked semantic example (non-limiting).}
In some embodiments, a non-limiting worked example illustrates the
semantic extension branch as follows. A text prompt is rendered as a
spatial optical pattern by an encoder. The pattern is projected onto a
selected expert subunit via the AOD routing subsystem. The expert
subunit scatters and transforms the pattern. A camera captures the
output. A trained readout head extracts a semantic feature vector.
The entire chain --- prompt rendering parameters, routing identity,
expert microstructure identity, calibration state, readout weights version, and
output --- is committed to the protocol digest. In some embodiments,
this chain produces an attested semantic computation record as
described above. In some embodiments, the same chain may be replayed
under declared calibration tolerances for dispute resolution.

% ------------------------------------------------------------------


\paragraph{Industrial applicability scope (non-limiting).}
The semantic extension branch disclosed in this section is directed to
attested visual input/output around a digital semantic model,
bounded-vocabulary semantic grounding, narrow-task continuous streaming
semantic operation, and semantic \TB attestation of claims over
committed physical evidence windows. This disclosure does not assert,
and shall not be construed to assert, that the physical scattering
substrate constitutes an open-ended language model, a general-purpose
semantic engine, or a substrate for training-scale deep neural networks
over natural language.

\subsection{Mapping to reactor-network primitives}
\label{sec:reactor-primitives}
% ------------------------------------------------------------------

In some embodiments, the routed physical ensemble instantiates route,
compile, parallelise, reduce, and abstract primitives of a broader
reactor network algebra. In some embodiments, route corresponds to
expert selection, compile to sequential chaining, parallelise to
simultaneous activation, reduce to output aggregation, and abstract to
trained readout or head extraction. In some embodiments, fleet-level
calibration and provenance monitoring extend across the ensemble so
that expert identity, routing history, and quality-state transitions
are auditable.


% ======================================================================

% ======================================================================
%
%
% ======================================================================

\subsubsection{Heterogeneous expert-type example (non-scattering, non-limiting)}
% ------------------------------------------------------------------

The following example illustrates how the routed-ensemble
architecture can be mapped to expert subunits beyond disordered
scattering media. As with the non-optical reactor examples above,
this is an architectural mapping showing that the routing,
provenance-binding, and commitment infrastructure can apply across
expert types; whether a
specific heterogeneous combination achieves measurable task-relevant
complementarity is an empirical question requiring experimental
validation per configuration.

In some embodiments, a
heterogeneous routed ensemble includes at least two expert types
selected from: (a)~disordered scattering media providing
high-dimensional random-feature-map transformations, (b)~phase-change
elements (for example chalcogenide or GST-based reactors) whose
transfer function is switchable between amorphous and crystalline
states and whose switched configuration is metered and logged,
(c)~gain-pumped optical cavities whose mode structure provides a
transform that depends on cavity geometry, mirror reflectivities, and
pump conditions, (d)~photonic integrated circuit subunits whose
fabrication-variant waveguide couplings provide device-specific
transfer matrices, and (e)~fibre-delay networks whose delay and
coupling configuration define a recurrent temporal kernel.

In some embodiments, the routing interface --- including coupling
optics, mode matching, alignment calibration, and per-expert
metering --- varies by expert type and is determined per embodiment.
In some embodiments, the commitment and provenance infrastructure
(protocol digest, routing log, microstructure identity, calibration-state hash,
readout weights version, time anchor) applies uniformly regardless of
expert type, because each expert is a parameterised Markov kernel
within the common framework.

In some embodiments, the central empirical question for any
heterogeneous combination is whether the ensemble achieves measurable
task-relevant complementarity --- that is, whether routing across
expert types produces statistically significant improvement on
declared benchmarks compared with single-expert and random-routing
baselines using committed physical evidence. In some embodiments, a
heterogeneous ensemble that lacks validated complementarity defaults
to single-expert or homogeneous-ensemble operation and is logged
accordingly.


% ======================================================================

% ======================================================================
%
%
%
% ======================================================================

\subsection{Conservative routed optical ensemble}
% ------------------------------------------------------------------

In some embodiments, a conservative routed optical ensemble (non-limiting)
comprises a coherent source (for example a 532~nm or 785~nm diode
laser), a beam expander, a spatial light modulator with approximately
$10^6$ pixels, a TeO$_2$ acousto-optic deflector providing
approximately 120 resolvable positions at approximately
800~kHz random-access rate, relay optics (two cylindrical lenses, a
field lens, and a micro-relay array), 64 scattering expert subunits
arranged on a planar grid, collection optics per expert, a shared
CMOS camera with approximately $8 \times 10^6$ pixels, and a digital
controller connected to the SLM driver, AOD RF driver, and camera
via synchronised interfaces.

In some embodiments, the controller implements the following cycle:
(a)~select an expert subunit by driving the AOD to the corresponding
angle; (b)~display an input pattern on the SLM; (c)~capture the
output on the camera; (d)~compute routing features from a preliminary
digital analysis of the captured output or from a separately trained
feature extractor; (e)~log the routing decision, expert identity, SLM
pattern identifier, AOD drive parameters, and camera frame in the
protocol digest; and (f)~optionally chain to a second expert stage
via gain-assisted relay as described in
Section~\ref{sec:gain-coupling}, with per-stage loss budget as
described in Table~\ref{tab:gain-budget}.

In some embodiments, the scan cycle time is dominated by the AOD
access time (approximately 1.2~$\mu$s) plus the camera integration
time (approximately 100~$\mu$s to 10~ms depending on exposure
requirements), yielding effective scan rates of approximately
100~Hz to 10~kHz for the full route-project-capture cycle. In some
embodiments, the SLM refresh rate (approximately 60~Hz for a
liquid-crystal SLM or approximately 10~kHz for a digital
micromirror device) may be the bottleneck for input pattern
changes; in other embodiments, the same SLM pattern is used across
multiple expert selections within a single SLM frame, amortising
the SLM refresh latency across expert queries.

In some embodiments, initial bring-up and calibration include:
(a)~per-expert transmission-matrix sampling using random input
patterns and camera capture, yielding a sampled
transmission-matrix block per expert from which Tier~2 effective
rank is estimated via singular-value decomposition; (b)~AOD
angle-to-position calibration using a known target in each expert
position; (c)~relay alignment verification using a collimated beam
and position-sensitive detector; (d)~SOA gain and noise-figure
measurement (if gain-assisted chaining is used); and (e)~baseline
verisimilitude scoring using the trained discriminator on committed
bundles from the calibrated configuration.

\paragraph{Relay-geometry design parameters (non-limiting).}
In some embodiments, the relay layout is characterised by: working
distance from the AOD output aperture to the expert input plane
(approximately 100--300~mm), magnification selected to match the AOD
output beam diameter to the expert input aperture, field curvature
and distortion across the expert grid (corrected by the field lens
and micro-relay array), and vignetting at the grid edges (metered
during calibration and compensated by per-position gain adjustment).
In some embodiments, the relay is designed for telecentric operation
at the expert input plane to minimise angle-dependent coupling
variation across the grid.

In some embodiments, the micro-relay array comprises one microlens
per expert input facet, with focal length and numerical aperture
matched to the scattering medium's input mode profile. In some
embodiments, the micro-relay array is fabricated as a single
moulded or lithographically defined element and is aligned as a unit
during assembly.

\paragraph{Experimental validation requirements (non-limiting).}
In some embodiments, the routed optical ensemble is validated by:
(a)~measuring single-expert task performance on declared benchmarks;
(b)~measuring multi-expert routed performance on the same benchmarks
under the trained routing policy; (c)~measuring random-routing
baseline performance; (d)~reporting whether multi-expert routing
produces statistically significant improvement over single-expert and
random-routing baselines; (e)~reporting Tier~2 effective rank per
expert and Tier~3 task-useful rank for the ensemble; and (f)~reporting
per-stage loss budget and, if applicable, whether gain-assisted
chaining improves or degrades net task performance compared with
single-stage operation. In some embodiments, all validation uses
committed physical evidence and is reproducible under the declared
calibration and protocol-digest conditions.

% ======================================================================


\section{Agent Integration (PolieBots) and Verification-Gated Actions}
\label{sec:agents}
% ======================================================================

This section describes embodiments in which an autonomous or
semi-autonomous agent uses one or more \RK modules as its
primary observation, verification, and action-conditioning channel.
References to ``PolieBots'' are non-limiting examples of mobile
platforms carrying one or more \RK modules.  The
verification-gated action mechanism and the RK-centred agent loop are
core operational mechanisms of this disclosure; optional higher-layer
coordination and training architectures are described in the related Filing 2 application (optional, non-essential).

\paragraph{Physical and functional composition (non-limiting).}
In one non-limiting configuration, a PolieBot apparatus comprises:
\begin{enumerate}[nosep]
  \item A \emph{mobile or deployable platform} (for example a wheeled,
    tracked, legged, aerial, or aquatic chassis; a fixed-installation
    enclosure; or a wearable or handheld carrier) providing structural
    support and power for the remaining subsystems.
  \item One or more \emph{\RK modules}, each comprising at
    least an emitter, a detector, one or more physical paths coupling
    emission to measurement, and the meter, commitment, and
    protocol-digest infrastructure described elsewhere in this
    disclosure (for example the bench-top 1D anchor of
    Section~\ref{sec:anchor-1d}, the 2D anchor of
    Section~\ref{sec:anchor-2d}, or any of the reactor variants
    described in subsequent sections).
  \item A \emph{locomotion or actuation subsystem} (for example wheels
    with motors and encoders, propellers, articulated limbs,
    manipulator arms, or steering surfaces) configured to translate
    platform motion commands $a^{\mathrm{platform}}_t$ into physical
    movement of the platform or of one or more \RK modules
    relative to the scene; in some embodiments the motion commands are
    derived directly from the analogue output of one or more \RK modules of Item~2 via declared transducers, without
    separate digital command synthesis.
  \item A \emph{communication subsystem} (for example a wireless
    radio, an optical link, a wired interface, or a beacon receiver)
    configured to receive challenges, beacons, or coordination
    signals from external verifiers or peer agents, and to publish
    committed bundle digests, evidence summaries, or coordination
    payloads; in standalone embodiments the communication subsystem
    is absent and the apparatus operates without external
    coordination.
  \item A \emph{decision-making subsystem} that selects per-step
    actions $a_t = (u(t), a^{\mathrm{platform}}_t,
    a^{\mathrm{tool}}_t)$ based on observation history $o_{0:t}$ and
    meter envelopes $\mathbf{m}_{0:t}$, comprising one or more of:
    \begin{enumerate}[nosep,label=(\roman*)]
      \item one or more \RK modules of Item~2 operating
        as the primary decision-making mechanism, in which the
        analogue feedback dynamics of the reactor-scene loop, taken
        together with declared decoders or transducers from
        reactor-side observables to platform-side commands, directly
        determine per-step actions without requiring separate
        digital policy evaluation (for example a fully optical
        feedback loop between a reactor and a scene in which
        emission schedules and platform motion commands are
        produced by the loop's own dynamics);
      \item one or more processors executing a learned or programmed
        policy $\pi_\psi(\cdot \mid o_{0:t}, \mathbf{m}_{0:t})$ over
        observation history and meter envelopes; or
      \item a hybrid configuration in which the \RK
        analogue dynamics of (i) and the digital policy of (ii) are
        composed (for example an analogue inner loop conditioned by
        a digital outer policy, or a digital policy whose output is
        filtered or shaped through an analogue reactor stage before
        actuation);
    \end{enumerate}
    together with a verification-gating mechanism that partitions
    the action space according to evidence sufficiency
    (unconstrained, RK-verified, hardness-gated, and
    multi-RK-verified classes as described below in this section),
    and an audit-policy interface that constrains action eligibility
    to declared envelopes.  The verification-gating mechanism and
    audit-policy interface are functional requirements and may be
    implemented digitally, in the analogue feedback dynamics of one
    or more \RK modules, or in any combination.
  \item Optionally, an \emph{external-tool interface} configured to
    invoke tool actions $a^{\mathrm{tool}}_t$ (for example calls to
    remote computation, retrieval, or human-in-the-loop confirmation)
    under the same verification, hardness, and audit constraints
    applicable to the platform's other actions.
\end{enumerate}

In some embodiments, the platform carries multiple \RK
modules in distinct geometric or spectral configurations, supporting
multi-RK-verified actions of the type described below in this
section.  In some embodiments, the locomotion and \RK
subsystems are mechanically coupled such that platform motion
participates in the scan law $u(t)$ (for example a body-mounted
sensor whose pose changes as the platform moves); in other
embodiments the \RK modules are independently steered
relative to the platform and the scan law $u(t)$ is constructed
from module-internal degrees of freedom only.

\paragraph{Fully analogue configuration (non-limiting).}
In a non-limiting fully analogue configuration, the apparatus
reduces to a mobile or deployable platform of Item~1; one or more
\RK modules of Item~2 acting in role~(i) of Item~5; a
locomotion subsystem of Item~3 driven by analogue transducer output
from those \RK modules; and an analogue
verification-gating subsystem and analogue audit-policy interface
of the form described below, implemented within or alongside the
analogue feedback dynamics of the same modules. The digital
processor of role~(ii) of Item~5, the external-tool interface of
Item~6, and the communication subsystem of Item~4 may all be absent
in such embodiments.

In one non-limiting analogue verification-gating subsystem, the
gate apparatus comprises:
(a) one or more analogue meter transducers configured to convert
declared committed-evidence observables (for example photodetector
amplitudes, coupling-quality envelope readouts, or hardness
proxies derived from the reactor-scene loop) into continuous
voltage or current signals;
(b) one or more threshold comparators or window discriminators
configured to compare each meter signal to a declared fixed
threshold or window, the threshold values being set as
hardware-declared component values (for example resistor-divider
references, voltage-reference diodes, or precision-set
potentiometer trims) that are physically inspectable as
device-state evidence;
(c) one or more analogue latch or sample-and-hold elements
configured to record the comparator outcomes for the duration of
an action-class evaluation window, providing memory of recent
gate state across the loop's response time;
(d) one or more analogue limiters or rate-limited switch elements
configured to clamp, attenuate, or interrupt the locomotion or
emission control path when any of the gated comparator outcomes
indicates an out-of-envelope condition; and
(e) a policy-selecting switch network configured to route the
locomotion or emission control between an unconstrained path, an
RK-verified path, a hardness-gated path, and a multi-RK-verified
path according to the latched comparator outcomes, the routing
implementing the partition of the action space described elsewhere
in this section.

In one non-limiting analogue audit-policy interface, gate state is
recorded as physical evidence of device behaviour rather than as a
digital audit log. Non-limiting recording mechanisms include: a
servo-driven mechanical position, ink-on-paper trace, exposed
photographic medium, or set-once analogue memory element whose
state is determined by the latched comparator outcomes and
inspectable post-run; a co-emitted committed companion analogue
signal recorded by an external observation channel; or
incorporation of the gate state into the optical or electrical
signature of the controlled emission such that a verifier with
access to the committed output and the declared hardware
declaration can recover the gate's effective decision history. The
hardware declaration --- comprising at least the threshold
component values, comparator and latch identities, switch-network
topology, and recording-medium identity --- is committed under the
same protocol-digest discipline applied elsewhere in this
disclosure to physical configuration declarations, so that an
audit recovers what the analogue gate could and could not have
done given its physically declared parameters.

The analogue verification-gating and audit-policy machinery
described in this paragraph is non-limiting; embodiments may
combine analogue gate elements with digital, optical, or hybrid
recording mechanisms. Where verifier-facing audit beyond the
committed hardware declaration and committed output channel is
required (for example for selectively-openable atom-level evidence,
cryptographic commitment plumbing, or remote audit), a digital or
hybrid implementation as described in role~(ii) or role~(iii) of
Item~5 is included rather than the fully analogue configuration of
this paragraph.

\subsection{RK-centred agent loop}

In some embodiments, at each step $t$, the agent selects an action
$a_t = (u(t), a^{\mathrm{platform}}_t, a^{\mathrm{tool}}_t)$, where
$u(t)$ is applied to RK modules, $a^{\mathrm{platform}}_t$ denotes
platform motion, and $a^{\mathrm{tool}}_t$ denotes calls to external
tools. A non-limiting per-step reward is
\[
  r_t
  = r^{\mathrm{task}}_t
    - \lambda_{\mathrm{cost}}\,c\!\bigl(u(t)\bigr)
    + \lambda_{\mathrm{info}}\,\widehat{I}_t
    + \lambda_{\mathrm{hard}}\,H_t.
\]

In some embodiments, the RK control is selected by a learned policy
$u(t) \sim \pi_\psi(\cdot \mid o_{0:t}, \mathbf{m}_{0:t})$, so that
scan paths and emission schedules adapt to scene context and meter
envelopes while remaining compatible with audit policies.

\subsection{Verification-gated actions}

In some embodiments, the agent partitions its action space into classes
with different verification requirements: unconstrained actions
(internal computation, reversible actions), RK-verified actions
(conditioned on recent verisimilitude checks), hardness-gated actions
(conditioned on empirical hardness thresholds), and multi-RK-verified
actions (requiring corroboration from multiple independent modules).
Failures trigger bounded escalation.

\subsection{Continual learning via physical kernel interactions}

In some embodiments, the agent logs experience tuples
$\tau_t = (o_t, u(t), a^{\mathrm{platform}}_t, r_t, \mathbf{m}(t),
\Pi_t)$ and uses these logs for continual learning with bounded drift
and audit compatibility. Hardness-regularised objectives include
penalties for meter violations and rewards for maintaining or improving
empirical hardness during adaptation.

\paragraph{Automated reactor design and meta-learning (non-limiting).}
In some embodiments, the design of reactor configurations (media
selection, geometry, scan protocols, and associated digital heads) is
optimised via meta-learning or automated search. A meta-learner
proposes reactor configurations; each configuration is instantiated
physically and/or in simulation and evaluated on a distribution of
tasks; and fitness signals guide search toward configurations that
generalise well under declared cost and safety constraints.

\paragraph{Uncertainty quantification and calibration (non-limiting).}
In some embodiments, the system produces calibrated uncertainty
estimates alongside predictions. Non-limiting mechanisms include
ensemble queries across multiple reactor passes or heads, stochastic
decoding that samples from output distributions, explicit uncertainty
heads trained with calibration losses, and conformal-prediction-style
wrappers that provide coverage guarantees under declared assumptions.
In some embodiments, intrinsic stochasticity of the physical channel
provides a natural source of uncertainty signals.

\paragraph{Streaming and online inference (non-limiting).}
In some embodiments, the \RK processes continuous input
streams online, producing outputs with bounded latency after each
input chunk rather than waiting for complete sequences. Reactor memory
maintains running context; an encoder maps arriving chunks to emission
patterns; and a decoder produces partial outputs that are refined as
more of the \cb arrives.

\paragraph{Self-modelling and introspection (non-limiting).}
In some embodiments, the system maintains a self-model of its
capabilities, limitations, and current calibration state. This
supports confidence calibration, knowing when to defer, predicting
latency/energy/error characteristics, and detecting drift or
degradation from evidence streams and meter outputs.

\subsection{Interpretability and internal model probing via RK
externalisation}

In some embodiments, the agent uses a \RK as a physically
grounded interpretability and debugging mechanism. The agent
externalises internal hypotheses, predicted states, or latent
variables as emission patterns and compares observed responses to
predicted responses. A non-limiting procedure includes:
\begin{enumerate}
  \item Encode an internal hypothesis or predicted scene property into
    a structured emission protocol.
  \item Execute the protocol on the \RK and record
    the \cb.
  \item Compare the measured \cb to the agent's predicted
    \cb under the same protocol.
  \item Update belief state or flag mismatch when discrepancy exceeds
    a threshold.
\end{enumerate}
This provides a model-reality consistency check that is difficult to
satisfy by purely semantic manipulation because it is anchored to the
physical channel.

\paragraph{Optional extensions
(described in the related Filing 2 application; optional, non-essential).}
In some embodiments, the agent system uses optional training and
operational extensions, together with sim-to-real world-model
refinement. These embodiments are described in the related Filing 2 application
(optional, non-essential). The
sim-to-real transfer mechanism, bifurcation-targeted probing, and
fleet-driven simulation improvement described there use the same
\cba, meter, and commitment machinery as the core
embodiments and are enabled by them, but are not required for core
\RK operation.

\paragraph{Perception hierarchy (non-limiting).}
In some embodiments, agent perception is organised as a hierarchy
of increasing abstraction: raw photon statistics (photon counts,
arrival times, spectral distribution); spatial structure (edges,
textures, depth, spatial frequency content); temporal dynamics (motion,
flicker, rhythmic patterns, transient events); semantic scene
interpretation (object identity, scene layout, activity recognition);
causal modelling (physical cause-and-effect relationships, intervention
predictions); and relational reasoning (relationships between entities,
roles, social structure). Each level of the hierarchy corresponds to
a different decoder head applied to the \cb, and the
agent's developmental curriculum can be structured to introduce
higher levels progressively as lower levels are mastered.

\subsection{Offline simulation and physical validation}
\label{sec:offline-sim}

In some embodiments, offline or partially offline experience is
generated in at least two ways. In a first class of embodiments, a
physical \RK module operates on synthetic or modified scenes
(for example projected patterns, printed targets, controlled test
fixtures, or augmented views), so that the optical and media stack
remains physical while the scene content is virtualised. In a second
class of embodiments, the \RK is replaced by a learned
emulator or digital twin calibrated to the physical module, and the
agent interacts with the emulator as a virtual reactor. A non-limiting
method comprises:
\begin{enumerate}
  \item Train or update an emulator on the logged \cb and
    control histories, and optionally define families of synthetic
    scenes for use with a physical \RK module.
  \item Generate simulated rollouts under candidate policies, using at
    least one of: (i)~a virtual \RK emulator, and (ii)~a
    physical \RK operating on synthetic or modified scenes.
  \item Select validation probes, execute them on the physical \RK, and compare the observed \cb to emulator
    predictions.
  \item Update policy and world model using both simulated and physical
    experience, with weighting dependent on validation outcomes and
    confidence.
\end{enumerate}

\paragraph{Parameter-sandboxed offline sessions (non-limiting).}
In some embodiments, the physical kernel operates on synthetic or
modified scenes during an offline session in which kernel parameters
may be updated freely.  All parameter updates made during the offline
session are reverted upon return to live operation, so that the
session constitutes a sandboxed exploration and no parameter
excursion persists into the live operating configuration.  In some
embodiments, classifiers or other derived models trained on the
\cba record produced during the sandboxed session are
retained and available for use in subsequent operation, even though
the kernel parameters themselves are reverted.  Higher-layer
deployment of such classifiers is described in the related Filing 2 application
(optional, non-essential).

\paragraph{Calibration and self-model construction during offline
periods (non-limiting).}
In some embodiments, an offline calibration period is used for
physical parameter update (drifting-field compensation and
re-commissioning against declared meter envelopes) and for fitting a
generative self-model to the device's own logged \cba history.
Unlike sandboxed sessions, parameter updates made during calibration
periods are retained.  In some embodiments, the resulting smoothed
self-model serves as a synthetic scene source for subsequent
sandboxed sessions.

\paragraph{Inward-facing gain modulation (non-limiting).}
In some embodiments, the \RT regime is applied with a
local gain-modulation scalar (denoted $g_{\mathrm{in}}$ and distinct
from the trust score $\sigma_{ij}(t)$, the canonical participant
ordering $\pi_{\mathrm{ord}}$, the tool-call specification
$\mathsf{tc}$, and any conventional statistics-notation $\sigma$
elsewhere in this disclosure) targeting the device's own response
gain
rather than an external scene's appearance.  In some embodiments,
this may be applied inwardly (modulating perceptual reactivity),
outwardly (attenuating emitted signal), or simultaneously in both
directions.  In some embodiments, the device remains live during
inward-facing operation and may continue to interact with its
environment; the distinguishing feature is that the fast feedback loop
is deliberately damped, for example via a smoothed self-model
constructed during a prior calibration period.  In some embodiments,
this is a degenerate case of \RT rather than a new
regime.

\subsection{Multi-RK corroboration for agents}

In some embodiments, an agent has access to multiple \RK
modules with distinct configurations, invariance profiles, or vantage
points. The agent may request corroboration for a claim by querying
multiple kernels and comparing the resulting evidence under
spatiotemporal consistency constraints. In some embodiments, the agent
uses a committee rule such that high-stakes actions require
corroboration from at least $k$ of $n$ kernels, optionally with
diversity constraints so that not all corroborators share the same
failure mode.

% ----------------------------------------------------------------------
\subsection{Bilateral Proof-of-Discrepancy (B-PoD) protocol (non-limiting)}
\label{sec:bpod}
% ----------------------------------------------------------------------

\subsubsection*{Purpose}

Bilateral PoD (B-PoD) exists to let independent devices
share and verify physical discoveries even when they are not
members of the same fleet.  B-PoD makes a jointly encountered
physical correction event portable, authentic, and open to later
third-party reproduction, so that cross-device trust can be grounded
in reproducible physical evidence rather than in shared fleet
membership alone.

The primary security function of B-PoD is different from that of
fleet-quorum PoD.  Fleet quorum PoD provides immediate
Byzantine-fault-tolerant consensus by requiring physical reproduction
on $n \geq 3f+1$ independent verifier devices.  B-PoD provides
deferred public falsifiability: any third party can attempt to
reproduce the claimed discrepancy; a fake bond that does not reflect
a genuine physical event is falsifiable by any party who tries to
replicate the experiment and fails.  These are complementary security
functions, not competing alternatives.

\paragraph{Physical Authentication Principle (reminder).}
See Remark~\ref{rem:physical-auth-principle}.  The primary object of
B-PoD is the physical discrepancy event.  Cryptographic identifiers,
signatures, and related proofs authenticate the event record; they do
not by themselves establish the physical truth of the recorded claim.

% --- Native object model --------------------------------------------------
\subsubsection*{Native object model}

B-PoD uses three distinct objects.  Standards-based infrastructure
(DID-style identifiers, verifiable credential proof formats,
remote-attestation evidence structures) are interoperability
mappings over these objects, not substitutes for them.

\paragraph{B-PoD Event.}
The primary object.  A record of a physical event with authentication
attached.  In some embodiments, a valid B-PoD Event binds all of the
following:
\begin{enumerate}[label=(\roman*)]
  \item \textbf{Originator identifiers.} Controller-bound identifiers
    for both participating agents.
  \item \textbf{Experiment Definition.} Content-addressed specification
    of the experiment: declared circuit configuration identifier,
    envelope version, scan law, protocol parameters, and procedure
    sufficient for independent replication.
  \item \textbf{Discrepancy domain and type.} Structured label
    indicating the relevant physical subdomain (e.g.\ scattering,
    nonlinear dynamics, coupling) and the category of discrepancy
    (e.g.\ systematic, noise-floor, regime-boundary).
  \item \textbf{Pre-run commitments.} Each originator's pre-run
    predictions or model expectations for the declared experiment,
    committed cryptographically before execution.
  \item \textbf{Independence declaration.} A signed statement by each
    originator that, during the run generating the claimed discrepancy,
    its execution environment was not controlled by, inspectable by, or
    execution-steered by the counterparty except through the agreed
    Experiment Definition and permitted protocol messages.
  \item \textbf{Evidence commitments.} Content-addressed commitments to
    raw sensor logs, calibration records, processing steps, and meter
    summaries sufficient to support independent replication.
  \item \textbf{Reproduction criteria.} Declared conditions under
    which a later party's run constitutes a same-configuration
    replication (see Replication Outcome types below).
  \item \textbf{Observed discrepancy statement.} The measured
    discrepancy relative to declared pre-run predictions, stated in
    the same units as the corresponding fleet discrepancy metric.
  \item \textbf{Model update claim.} A declaration (not a model
    disclosure) that at least one originator updated a model or
    prediction in response to the discrepancy.  Model update policy
    is local to each agent; this field records that an update
    occurred, not its content.
  \item \textbf{Both originator signatures.} Authenticating the event
    record.  The signatures are not the event; they are the
    authenticity layer on the event.
\end{enumerate}

\paragraph{B-PoD Envelope.}
The cryptographic carrier for a B-PoD Event.  Contains
controller-bound identifiers, signing keys, timestamps, nonces, and
optional hardware attestation material.  DID-style identifiers and
remote-attestation-style evidence operate at the envelope layer.  The
envelope transports the event; it is not the event.

\paragraph{Portable Discrepancy Record.}
An append-only collection of B-PoD Events plus all subsequent
appraisals.  The following outcome types are first-class record
elements:

\begin{itemize}
  \item \textbf{Same-configuration replication.} An independent party
    ran the same Experiment Definition (same circuit configuration
    identifier, envelope version, and procedure as declared in the
    root event) and either confirmed or disconfirmed the observed
    discrepancy.  This is the only outcome type that is directly
    confirmatory or disconfirmatory for the root claim.
  \item \textbf{Related-configuration replication.} A run using a
    related but distinct circuit configuration, envelope, or
    procedure.  Informative but not directly confirmatory.  Must be
    typed distinctly; blending this type with same-configuration
    replication smears evidence across experiments that are nearby
    in domain but not the same claim.
  \item \textbf{Methodological challenge.} An argument that the root
    run or any subsequent replication was not a valid test of the
    declared claim.  May target any event in the record, not only the
    root.
  \item \textbf{Supersession.} Replaces the operational interpretation
    of an earlier event with a better-characterised successor while
    preserving the superseded event in history.  Supersession is
    science working correctly, not a retraction.
  \item \textbf{Withdrawal.} Originator-level reduction of confidence
    in an earlier claim, without deleting the record.  A withdrawal is
    a signal worth examining; it does not erase the prior event.
\end{itemize}

\paragraph{Living Scientific Object.}
A Portable Discrepancy Record is a living scientific object rather than
a static credential.  Confidence in a B-PoD claim is cumulative and
asymmetric: successful replications increase confidence without proving
the claim correct; failed replications decrease confidence without
conclusively refuting it (failure can arise from methodological
differences, uncontrolled variables, or other causes); and the full
history of both is part of the object itself.

% --- Restricted-domain commensuration ------------------------------------
\subsubsection*{Restricted-domain commensuration}

B-PoD enables cross-device comparison only within the
restricted domain defined by a declared experiment, circuit
configuration, operating envelope, and procedure.  Two devices with
otherwise divergent or independently trained predictive models may
still jointly determine whether the declared physical system under
declared conditions produced the declared outcome.  Physical reality
does not depend on which model either device holds.  B-PoD does not
require or imply convergence of broader models or calibration
histories; it enables measurement-level comparison within a declared
envelope.

This is the architectural mechanism by which devices with different
predictive models can share physically grounded corrections:
reproducibility under a declared protocol is the common basis.
Two devices with genuinely different surrogate models can both
execute the same physical experiment and jointly observe the same
discrepancy.  What bilateral PoD records is that joint observation.


\paragraph{Extended bilateral applications
(non-limiting).}
In some embodiments, B-PoD supports operation by independent or
fleet-divergent devices, fleet reintegration through reproducible
claims, and relying-party appraisal of Portable Discrepancy Records.
These extended applications are described in the related Filing 2 application (optional, non-essential).

% ----------------------------------------------------------------------
\subsection{Security considerations for RK-agent systems}

In some embodiments, attacks include scene spoofing, co-illumination
interference, hardware tampering, communications interference, and
adversarial examples targeting meters. Mitigations include
multi-channel invariance profiles, multi-RK corroboration,
spatiotemporal checks, and gating of actions on validated RK evidence.
In multi-agent deployments, some nodes may be compromised or
Byzantine. Trust-weighted aggregation, diversity constraints, and
periodic audits of calibration and meter agreement may be used to
reduce the influence of compromised nodes.

\subsection{Agent lifecycle management and trajectory records
(non-limiting)}
\label{sec:agent-lifecycle}

\paragraph{Lifecycle states (non-limiting).}
In some embodiments, a device or agent transitions between lifecycle
states:
\textbf{Embodied} (external actuators connected and real scene loop
active), \textbf{Docked} (actuators disconnected or logically
disabled while RK evidence plumbing remains active),
\textbf{Offline} (operating on synthetic or modified scenes, with
parameter-sandboxed or calibration modes as described in the offline
simulation subsection above),
and \textbf{Re-embodied} (actuators
reconnected after compatibility checks).  State transitions are
logged in protocol digests and meter summaries.  The full
lifecycle state machine is
described in the related Filing 2 application (optional, non-essential).



% ======================================================================
\section{Digital Twin and Reality Merge Pipeline (Optional, Non-limiting)}
\label{sec:twin-pipeline-extension}
% ======================================================================

In some embodiments, the \RT regime and its associated
training, deployment, and verification machinery are extended via a
digital-twin and reality-merge pipeline.  This section describes
non-limiting embodiments of multi-factor pose-conditioned \RT, payload-embedding \RT, co-emission across the
separation channel, twin training and analogue-twin deployment, style
corpus acquisition, generator-and-inverse pairs, scene-state-contrast
training-data construction, and associated label-integrity and
privacy-discipline mechanisms.  These embodiments are optional
extensions: core \RK operation, including the three
operating regimes of Section~\ref{sec:regimes}, does not require any
mechanism described in this section.

\subsubsection{Multi-factor pose-conditioned Reality Transform
(non-limiting)}
\label{sec:pose-rt}

In some embodiments, the \RT regime is specialised to
articulated-content production by adding declared pose-feature
consistency terms to the RT control objective of
Section~\ref{sec:style-parameter}. The pose-feature terms are
computed from committed evidence windows of the \cb under a
declared extraction rule, and are used by the controller to select
or update emission controls subject to the protocol digest and
meter envelope. The pose-feature terms do not require any novel
pose-estimation architecture: the pose-feature extractor may be any
declared measurement head whose identity, weights or version, input
preprocessing, confidence thresholds, calibration state, and output
coordinate convention are committed to the protocol digest.

In some embodiments, a committed \RTa decoder or
extraction map $\mathsf{RTDec}$, whose identity, version,
parameterisation, output codomain, and any fixed comparison map are
committed to the protocol digest, maps a committed evidence window
of the \cb to an image-like or feature representation
\[
  Y^{\mathrm{dec}}_t
  =
  \mathsf{RTDec}\!\bigl(C_{\mathsf{Win}_t},\Pi\bigr),
\]
where $C_{\mathsf{Win}_t}$ denotes the committed \cba
atoms selected by the declared evidence-windowing rule
$\mathsf{Win}_t$, and $\Pi$ is the protocol digest. The declared
output codomain of $\mathsf{RTDec}$, or the declared fixed
comparison map applied to it, is chosen so that
$d\!\bigl(Y^{\mathrm{dec}}_t,\mu_\phi(x_t,u_t)\bigr)$ is evaluated
in a common representation space.
A declared pose-feature extractor $\mathsf{Pose}$ --- whose
architecture, weights, input preprocessing, confidence threshold,
calibration state, and output coordinate convention are likewise
committed --- then produces a pose-feature vector, joint-angle
vector, keypoint set, parametric-body-model state, hand-pose state,
or related pose representation
\[
  p_t = \mathsf{Pose}\!\bigl(Y^{\mathrm{dec}}_t\bigr),
\]
optionally together with a confidence or visibility mask
$\mathsf{Mask}_t$. The $\mathsf{RTDec}$ identity, pose-feature
extractor identity, pose-coordinate convention, time base
$\Delta t$, output-comparison space, and any confidence-mask rule
are committed to the protocol digest.

\paragraph{Compound pose-control penalty.}
In some embodiments, the controller optimises a non-limiting
pose-conditioned RT objective
\[
  J_{\mathrm{RT}\text{-}\mathrm{pose}}(U_{0:T};\phi)
  =
  -\,\mathbb{E}\!\left[
  \sum_{t=0}^{T-1}
  \Bigl(
    d\!\bigl(Y^{\mathrm{dec}}_t,\, \mu_\phi(x_t,u_t)\bigr)
    + \lambda_{\mathrm{dat}}\,\ell^{\mathrm{dat}}_t
    + \lambda_{\mathrm{tmp}}\,\ell^{\mathrm{tmp}}_t
    + \lambda_{\mathrm{meter}}\,
      \mathbf{1}\!\bigl[\mathbf{m}(t)\notin \mathcal{M}_{\mathrm{accept}}\bigr]
  \Bigr)\right].
\]
Here $\phi\in\Phi$ is the declared style parameter of
Section~\ref{sec:style-parameter}, $\mu_\phi$ is the target
response function defined therein, $d(\cdot,\cdot)$ is the declared
divergence or distortion metric, and $\mathbf{m}(t)$ and
$\mathcal{M}_{\mathrm{accept}}$ are the meter vector and meter
envelope of Definition~\ref{def:meter-envelope}.
The sign convention matches $J_{\mathrm{RT}}$ of
Section~\ref{sec:style-parameter}: the objective is
\emph{maximised}, which corresponds to \emph{minimising} the
per-step weighted-loss sum. The additional per-step pose-consistency
contributions $\ell^{\mathrm{dat}}_t$ and $\ell^{\mathrm{tmp}}_t$
are defined in the following paragraphs.
During operation, the values of $\ell^{\mathrm{dat}}_t$ and
$\ell^{\mathrm{tmp}}_t$ are computed from committed
\cba windows and applied by the controller to select
or update the next physical emission or control action of the
\RK module, so that the pose terms are operative control
constraints on a physical channel rather than merely post-hoc
mathematical scores. For online control, any pose-consistency term
whose declared window requires atoms later than the current control
decision is evaluated as a delayed meter, a retrospective update
term, or a declared predictive surrogate, and no future atom is
used to select a control action before that atom has been committed.
The weights $\lambda_{\mathrm{dat}}$, $\lambda_{\mathrm{tmp}}$, and
$\lambda_{\mathrm{meter}}$ are chosen on the same declared scale as
$d(\cdot,\cdot)$ so that all terms in the per-step sum are
commensurate, and are committed to the protocol digest together
with the identities of $\mathsf{RTDec}$, $\mathsf{Pose}$, the
pose-reference set $\mathcal{D}_{\mathrm{pose}}$, and the windowing
and normalisation rules used to compute the pose terms.

\paragraph{Dataset-consistency contribution (non-limiting).}
In some embodiments, the dataset-consistency contribution
$\ell^{\mathrm{dat}}_t$ is derived from a functional $L_{\mathrm{dat}}$
that compares pose features extracted from the RT evidence window
with pose features extracted from a declared pose-reference set
$\mathcal{D}_{\mathrm{pose}}$ under the same declared pose-feature
extractor $\mathsf{Pose}$ and coordinate convention. Non-limiting
realisations include:
\begin{enumerate}
  \item a learned pose-distribution discriminator
    $A_{\mathrm{dat}}$, whose architecture, weights, training window
    length, positive set, negative set, held-out validation policy,
    and version hash are committed to the protocol digest, and whose
    score on $p_{0:T}$ derived from committed RT outputs versus
    $\mathsf{Pose}$-features derived from $\mathcal{D}_{\mathrm{pose}}$
    defines $L_{\mathrm{dat}}$; in some embodiments,
    $A_{\mathrm{dat}}$ is trained under a declared binary or
    multiclass classification loss distinguishing pose windows drawn
    from $\mathcal{D}_{\mathrm{pose}}$ from pose windows drawn from
    generated, transformed, or edited RT pose windows, with the loss
    family and threshold-calibration policy committed to the
    protocol digest;
  \item a declared kernel two-sample statistic
    \[
      \mathrm{MMD}_{k}^{2}(P,Q)
      =
      \mathbb{E}_{p,p'\sim P}k(p,p')
      + \mathbb{E}_{q,q'\sim Q}k(q,q')
      - 2\,\mathbb{E}_{p\sim P,\,q\sim Q}k(p,q),
    \]
    computed over pose-feature vectors or pose-feature windows, with
    kernel family, bandwidth, batching rule, and confidence-mask
    policy declared in the protocol digest; and
  \item a pose-class or action-class consistency loss, in
    embodiments where the style parameter $\phi$ includes a declared
    pose-class, action-class, or kinematic-intent component, with
    the classifier identity, label ontology, calibration set, and
    confidence threshold committed to the protocol digest; in such
    embodiments, the protocol digest further declares the mapping
    from $\phi$ to the target pose-class, action-class, or
    kinematic-intent distribution, and declares whether the
    per-step contribution is a cross-entropy, calibrated margin, or
    other classifier-score loss.
\end{enumerate}
In some embodiments, $\ell^{\mathrm{dat}}_t$ is obtained from
$L_{\mathrm{dat}}$ by a declared per-step decomposition (for example
a windowed discriminator score, a windowed MMD estimate, or a
windowed classifier margin), with the decomposition rule committed.

\paragraph{Between-frame consistency contribution (non-limiting).}
In some embodiments, the temporal-consistency contribution
$\ell^{\mathrm{tmp}}_t$ is derived from a functional $L_{\mathrm{tmp}}$
that measures inter-frame pose-dynamics plausibility using the
committed time base of the \cb. Non-limiting realisations include:
\begin{enumerate}
  \item a derivative-residual penalty
    \[
      \ell^{\mathrm{tmp},\mathrm{deriv}}_t
      =
      \left\|\mathsf{Mask}_t\odot
        \frac{p_{t+1}-p_t}{\Delta t}\right\|_{\Sigma_v^{-1}}^2
      +
      \left\|\mathsf{Mask}_t\odot
        \frac{p_{t+1}-2p_t+p_{t-1}}{\Delta t^2}\right\|_{\Sigma_a^{-1}}^2,
    \]
    where $\Sigma_v$ and $\Sigma_a$ are declared velocity and
    acceleration envelopes or covariance estimates committed to the
    protocol digest. For boundary indices at which a finite-difference
    stencil is not fully available, the protocol digest declares
    whether a one-sided difference, padding convention, or
    boundary-index exclusion is used; the digest also declares how
    confidence masks are composed across the velocity or acceleration
    stencil and whether $\Sigma_v^{-1}$ and $\Sigma_a^{-1}$ denote
    ordinary inverses, regularised inverses, pseudoinverses, or
    diagonal reciprocal envelopes;
  \item a learned pose-dynamics residual
    \[
      \ell^{\mathrm{tmp},\mathrm{dyn}}_t
      =
      \bigl\|\mathsf{Mask}_{t+1}\odot
        \bigl(p_{t+1}-g_{\mathrm{dyn}}(p_t,u_t,\Delta t)\bigr)
      \bigr\|^2,
    \]
    where $g_{\mathrm{dyn}}$ is a declared pose-dynamics model whose
    identity, input fields, training corpus, validation corpus, and
    version hash are committed; in some embodiments,
    $g_{\mathrm{dyn}}$ is trained by a declared one-step or
    multi-step pose-prediction loss on committed positive
    pose-sequence windows, with the prediction horizon and
    teacher-forcing or rollout policy committed to the protocol
    digest; and
  \item a kinematic-plausibility classifier trained against committed
    physical pose-sequence windows and declared negative examples,
    including time reordering, frame-interpolation artefacts, pose
    or keypoint substitutions between protocol-compatible sequences,
    and local blending of motion-warped frames; in some embodiments,
    the kinematic-plausibility classifier is calibrated on held-out
    committed physical pose-sequence windows and held-out
    edited-negative windows, with the acceptance threshold and
    false-positive target committed to the protocol digest.
\end{enumerate}
In some embodiments, missing or occluded pose features are handled
by the confidence or visibility mask $\mathsf{Mask}_t$ rather than
by treating all keypoints as equally observed.

\paragraph{Composition with verisimilitude and structural meters.}
In some embodiments, $J_{\mathrm{RT}\text{-}\mathrm{pose}}$ is
evaluated together with the verisimilitude functional of
Definition~\ref{def:verisimilitude}, the edge-based structural
meters of Section~\ref{sec:edge-meters}, timing-integrity meters,
calibration-consistency meters, and spatiotemporal-consistency
meters. The pose contributions may be declared as acceptance
meters, evaluation meters, or objective terms according to the
meter-partition policy for the deployment. The
acceptance/evaluation partition is preserved so that controller
optimisation against pose consistency does not expose all
evaluative surfaces simultaneously.

\paragraph{Training data and negative-example construction
(non-limiting).}
In preferred embodiments, $\mathcal{D}_{\mathrm{pose}}$ comprises
committed physical recordings, committed RT runs, or committed
offline-session records whose protocol digests are available for
audit. In some embodiments, an externally curated pose-reference
set may be used as an operator-declared reference envelope, but
such a set does not by itself supply physical provenance or
empirical-hardness evidence unless its acquisition and validation
are separately committed. The protocol digest shall provenance-type
$\mathcal{D}_{\mathrm{pose}}$ as committed physical evidence,
committed offline-session evidence, or external-reference-only, and
verifier-facing attestations shall not treat an external-reference-only
set as physical-provenance evidence or as support for
empirical-hardness claims. Negative examples for discriminator and
kinematic-classifier training, and for validating pose-dynamics
residual thresholds, may be constructed by applying declared
spatiotemporal edits to committed recordings, with each edit
family, intensity parameter, and random-seed policy recorded or
hash-identified in the protocol digest, following the pattern of
Section~\ref{sec:verisimilitude-training}.

\paragraph{Sandboxed and live operation (non-limiting).}
In some embodiments, pose-conditioned RT training is performed
during a parameter-sandboxed offline session under the offline
simulation and physical-validation mechanism of
Section~\ref{sec:offline-sim}, in which the physical kernel
operates on synthetic or modified articulated scenes,
pose-reference sequences, or rendered targets. Parameter updates
made during the sandboxed session are reverted on return to live
operation unless a separate calibration or deployment authority
authorises retention under the applicable protocol. Pose-feature
extractors, pose-distribution discriminators, kinematic
classifiers, and validation summaries trained from the sandboxed
\cb record may be retained as derived models where permitted by
the protocol digest and meter policy.

\paragraph{Relation to the Limager regime (non-limiting).}
The compound-objective structure is formally analogous to the
\LI regime's decomposed information-gain objective in that both
regimes factorise a scalar control objective into declared terms
targeting distinct evidence dimensions, each subject to
protocol-digest commitments and meter-envelope discipline. The
primitive question is different: \LI selects probes to increase
information about scene variables, while pose-conditioned RT
selects emissions to realise a declared style target while
maintaining pose-distribution, pose-dynamics, structural, and
physical-consistency constraints.

\paragraph{Scope note (non-limiting).}
For avoidance of doubt, the pose-feature extractor $\mathsf{Pose}$
may be drawn from conventional keypoint, body-model, hand-pose, or
related pose-estimation architectures. In such embodiments, the
disclosed technical contribution is the use of the declared
pose-feature trace as a committed, meter-bounded component of a
\RT control objective over a physical evidence
substrate, rather than the pose-estimation architecture itself.


\subsubsection{Payload-embedding Reality Transform and prior-frame-hash
chaining (non-limiting)}
\label{sec:rt-payload}

In some embodiments, the \RT regime is specialised to
carry a declared bit-payload per frame, such that the payload is
extractable from the committed \cba evidence by a
declared extraction head while the style fidelity, verisimilitude,
and meter-envelope discipline of the surrounding RT objective are
preserved. The mechanism is a specialisation of the RT objective of
Section~\ref{sec:style-parameter}: the payload becomes an additional
declared evidence dimension on which the compound objective places
a meter-bounded term. The payload-embedding \RT is
compatible with, and independent of, the run-level hash-chain state
$\chi_t$ of Section~\ref{sec:hash-chain-state}.

\paragraph{Declared per-frame payload and extractor.}
In some embodiments, a declared embedder $\mathsf{Embed}$, committed
to the protocol digest by identity, version, weights, input
preprocessing, and output bit-length, produces an emission update
that causes frame~$t$ of the decoded evidence window to carry a
payload
\[
  \mathsf{Payload}_t \in \{0,1\}^{N_b},
\]
where $N_b$ is a declared payload bit-length (for example
$N_b\in\{32,\,64,\,128,\,256,\,512\}$), committed to the protocol
digest. A declared extractor $\mathsf{Extract}$, likewise committed
by identity, version, weights, and bit-error-rate (BER) threshold,
recovers an estimate $\widehat{\mathsf{Payload}}_t$ from the
$\mathsf{RTDec}$-decoded frame
$Y^{\mathrm{dec}}_t = \mathsf{RTDec}(C_{\mathsf{Win}_t},\Pi)$:
\[
  \widehat{\mathsf{Payload}}_t
  =
  \mathsf{Extract}\!\bigl(Y^{\mathrm{dec}}_t,\Pi\bigr).
\]
The BER threshold and any confidence threshold are calibrated on
held-out committed physical, parameter-sandboxed, or twin-predicted
windows under declared distortion, edit, and compression families,
with the calibration corpus and acceptance rule committed to the
protocol digest.
In some embodiments, $\mathsf{Embed}$ and $\mathsf{Extract}$ are
jointly trained under the compound RT objective defined below, on
committed parameter-sandboxed offline-session evidence, on
probe-corpus evidence (Section~\ref{sec:probe-corpus}), or on live
RT recordings committed during a declared training phase.

\paragraph{Prior-frame-hash payload structure.}
In some embodiments, the payload carried by frame~$t$ encodes one
or more hashes of prior frames, not necessarily the immediate
predecessor. Let
\[
  \mathsf{FrameHash}_t
  =
  H\!\bigl(
    \mathsf{Canon}(Y^{\mathrm{dec}}_t,\mathrm{meta}_t,\Pi),\,
    k_{\mathrm{hash}}
  \bigr),
\]
where $H$ is a declared cryptographic or keyed hash (for example
HMAC, SHA-family, sponge, or a ZK-friendly algebraic hash),
$\mathrm{meta}_t$ denotes declared per-frame metadata (for example
time base, protocol-digest epoch, device identifier), and
$k_{\mathrm{hash}}$ is a declared key or public parameter committed
or referenced in the protocol digest. Here $\mathsf{Canon}$ is a
declared canonicalisation map, committed by identity, version,
quantisation rule, colour or feature-space convention, metadata
fields, and rounding or tolerance policy, so that frame hashes are
recomputable from selectively opened physical evidence windows
rather than from uncontrolled raw pixel values.

A declared \emph{lag schedule}
$\mathsf{Lag}_t : \{0,\ldots,T-1\}\to \{1,\ldots,L\}^{N_{\mathrm{slot}}}$
assigns, for each frame~$t$ and each payload slot
$j\in\{1,\ldots,N_{\mathrm{slot}}\}$, a positive lag value
$\mathsf{Lag}_{t,j}\in\{1,\ldots,L\}$ indicating which prior frame's
hash occupies slot~$j$ at frame~$t$:
\[
  \mathsf{Payload}_t
  =
  \mathsf{Pack}\!\left(
    \mathsf{FrameHash}_{t-\mathsf{Lag}_{t,1}},
    \mathsf{FrameHash}_{t-\mathsf{Lag}_{t,2}},
    \ldots,
    \mathsf{FrameHash}_{t-\mathsf{Lag}_{t,N_{\mathrm{slot}}}};\Pi
  \right),
\]
where the packing, slot order, truncation, error-correcting-code,
and per-slot weighting rules of $\mathsf{Pack}$ are committed to the
protocol digest. For prefix frames in which
$t-\mathsf{Lag}_{t,j}<0$, the protocol digest declares whether the
corresponding slot is encoded by $\mathsf{Pack}$ as an omitted-slot
sentinel, filled with a committed initialisation hash, or filled
under another declared bootstrap convention, with $\mathsf{Pack}$
producing a fixed $N_b$-bit payload in all cases.

In some embodiments, $\mathsf{Lag}_t$ is
(i)~a fixed constant schedule (for example always
$\mathsf{Lag}_{t,j}=1$),
(ii)~a declared per-segment schedule committed to the protocol
digest, or (iii)~a pseudorandom-function output keyed by the
protocol digest and the run seed, so that $\mathsf{Lag}_t$ is
reproducible by a verifier holding the key or public parameters but
is not predictable to a party without them. Where a lag schedule or
embedding rule depends on a secret or delayed parameter, the
protocol digest commits the key identifier, public parameter,
delayed-reveal policy, or verification-access rule under which a
verifier may reproduce the schedule for audit. Multi-slot schedules
with $N_{\mathrm{slot}}\geq 2$ provide redundancy against short-range
edit attacks that attempt to resynchronise to an adjacent frame.

\paragraph{Compound RT objective with payload extractability.}
In some embodiments, the controller optimises a non-limiting
payload-conditioned RT objective, written in the negative-expected-loss
convention of $J_{\mathrm{RT}}$:
\[
  J_{\mathrm{RT}\text{-}\mathrm{pay}}(U_{0:T};\phi)
  =
  -\,\mathbb{E}\!\left[
  \sum_{t=0}^{T-1}
  \Bigl(
    d\!\bigl(Y^{\mathrm{dec}}_t,\mu_\phi(x_t,u_t)\bigr)
    + \lambda_{\mathrm{pay}}\,\ell^{\mathrm{pay}}_t
    + \lambda_{\mathrm{ver}}\,\ell^{\mathrm{ver}}_t
    + \lambda_{\mathrm{meter}}\,
      \mathbf{1}\!\bigl[\mathbf{m}(t)\notin\mathcal{M}_{\mathrm{accept}}\bigr]
  \Bigr)\right],
\]
where:
\begin{itemize}
  \item the style-target term
    $d(Y^{\mathrm{dec}}_t,\mu_\phi(x_t,u_t))$ is the RT fidelity term
    of Section~\ref{sec:style-parameter};
  \item the per-step payload-extractability contribution
    $\ell^{\mathrm{pay}}_t$ is a declared bit-error or cross-entropy
    loss between $\widehat{\mathsf{Payload}}_t$ and the target
    $\mathsf{Payload}_t$ specified by $\mathsf{Lag}_t$ and
    $\mathsf{Pack}$, with the loss family, per-slot weighting, and
    error-correcting-code policy committed to the protocol digest;
  \item $\ell^{\mathrm{ver}}_t$ is a declared verisimilitude term
    derived from the functional of
    Definition~\ref{def:verisimilitude}, penalising embedding
    artefacts that depart from the empirical distribution of genuine
    unembedded RT outputs under the committed verisimilitude
    discriminator; and
  \item the meter-violation term is the standard envelope term from
    Section~\ref{sec:style-parameter}.
\end{itemize}
The weights
$\lambda_{\mathrm{pay}},\lambda_{\mathrm{ver}},\lambda_{\mathrm{meter}}$
are on the same declared scale as $d(\cdot,\cdot)$ and are committed
to the protocol digest. Perceptual transparency is promoted and
meter-bounded by composition: the $\lambda_{\mathrm{ver}}$ term
penalises embedding artefacts that depart from the empirical
distribution of genuine unembedded RT outputs as measured by the
committed verisimilitude discriminator, and any acceptance threshold
or transparency envelope is declared in the protocol digest.

\paragraph{Composition with pose-conditioned RT.}
In some embodiments, the payload-extractability and verisimilitude
terms are composed with the pose-conditioned RT terms of
Section~\ref{sec:pose-rt}, yielding a compound objective
\[
  J_{\mathrm{RT}\text{-}\mathrm{pose+pay}}
  =
  J_{\mathrm{RT}\text{-}\mathrm{pose}}
  -
  \mathbb{E}\!\left[
    \sum_t
    \bigl(
      \lambda_{\mathrm{pay}}\,\ell^{\mathrm{pay}}_t
      +
      \lambda_{\mathrm{ver}}\,\ell^{\mathrm{ver}}_t
    \bigr)
  \right].
\]
Equivalently, the terms may be written as a single
negative-expected-loss objective under the same sign convention as
$J_{\mathrm{RT}}$ and $J_{\mathrm{RT}\text{-}\mathrm{pose}}$. The
meter-partition policy determines which terms are declared as
acceptance meters, evaluation meters, or objective terms; the
acceptance/evaluation partition is preserved as in
Section~\ref{sec:pose-rt}.

\paragraph{Tamper-evidence properties (non-limiting).}
In some embodiments, the prior-frame-hash payload chain provides
localised tamper evidence: if a frame $t^\star$ is post-hoc edited,
the $\mathsf{FrameHash}_{t^\star}$ values committed in the payloads
of all frames $t$ satisfying $t-\mathsf{Lag}_{t,j}=t^\star$ under
the declared $\mathsf{Lag}$ schedule fail validation against the
recomputed hash of the edited frame, after canonicalisation by
$\mathsf{Canon}$ and subject to the declared extraction confidence,
BER threshold, and error-correcting-code policy, thereby localising
the tamper event. This localised
evidence is orthogonal to, and composable with, the global
tamper-evidence provided by the run-level hash chain $\chi_t$ of
Section~\ref{sec:hash-chain-state}: the per-frame payload chain
identifies \emph{which} frame was edited, while $\chi_T$ provides
the compact run fingerprint. In some embodiments, PRF-indexed
$\mathsf{Lag}_t$ schedules strengthen tamper evidence against
anticipatory or segment-resynchronisation attacks; construction of
an internally consistent edited segment further requires
re-embedding corrected payloads under the declared $\mathsf{Embed}$
authority, key or parameter policy, BER threshold, and
verisimilitude and meter-envelope constraints.

\paragraph{Relation to hash-chain and watermark mechanisms
(non-limiting).}
The payload-embedding \RT is distinct from, and
compatible with:
(i)~the run-level hash-chain $\chi_t$ of
Section~\ref{sec:hash-chain-state}, which binds
$(u(t),\mathbf{y}_t)$-level state in a one-way chain external to
the rendered content;
(ii)~the covert device-watermark primitive in the
information-operations taxonomy, which embeds a continuous
device-identity signature rather than a discrete per-frame payload;
and
(iii)~the per-atom Merkle digests of
Section~\ref{sec:merkle-atoms}, which commit to atoms after
recording.
In the embodiments described in this paragraph, payload-embedding
RT is distinguished by embedding a declared \emph{discrete
extractable bit-payload} in each frame of the rendered content
itself, while the other listed mechanisms bind provenance, device
identity, or post-recording commitments through different surfaces.

\paragraph{Scope note (non-limiting).}
For avoidance of doubt, $\mathsf{Embed}$ and $\mathsf{Extract}$ may
be drawn from conventional deep-steganography or digital-watermarking
architectures. In such embodiments, the disclosed technical
contribution is the joint specification of a declared prior-frame-hash
payload structure, committed lag schedule, canonical frame-hash and
packing rules, and meter-bounded RT training or evaluation over
committed physical evidence windows or provenance-typed
twin-predicted windows, rather than the embedder or extractor
architecture itself.


\paragraph{Positive-residual additive-overlay RT objective
(non-limiting).}
\label{sec:baseline-continuous}
In some embodiments suitable for bright-ambient or
minimum-perturbance operation, the \RT objective is
restricted to additive overlays whose contribution to the
photodetector or imager response, evaluated against a declared
ambient or neutral baseline $\mathsf{Baseline}_t$, is non-negative
on the aggregate observed signal. The baseline $\mathsf{Baseline}_t$
is computed by the controller from declared baseline-acquisition
windows committed to the protocol digest --- non-limiting choices
include short pre-emission ambient samples, neighbouring-pixel
neutral references, or temporally adjacent neutral frames --- with
the baseline-acquisition rule, baseline-window length, and any
spectral or spatial pre-filtering parameters all committed to the
protocol digest. The compliance functional is the aggregate
observed positive excess
$\max(\widehat{R}_t - \mathsf{Baseline}_t, 0)$, evaluated over
the declared evaluation window and (where applicable) summed across
$K_{\mathrm{co}}$ co-emitted streams; the controller selects
emission schedules under the constraint that this aggregate-excess
functional remains within the declared additive-overlay envelope.
In some embodiments, the additive-overlay objective composes with
the meter envelope of Definition~\ref{def:meter-envelope} as an
additional declared meter term, and any per-stream perturbance
budgets are recorded as auxiliary constraints valid only within
the declared separation, linearity, and crosstalk envelope. The
aggregate formulation prevents a companion stream from complying
per-stream while the combined emitted light violates the
bright-ambient perturbance objective.

\paragraph{Perception-collapse semantic fallback (non-limiting).}
\label{sec:semantic-fallback}
In some embodiments, the \RT regime declares a
\emph{semantic fallback} regime selected at deployment time when a
perception-collapse condition is detected on the verisimilitude or
dual-direction-consistency meter envelope. The perception-collapse
condition is defined by a declared meter-threshold rule committed
to the protocol digest, comprising at least: a verisimilitude
floor, a dual-direction-consistency floor, a sustained-window
length over which the floors must be violated, and a hysteresis
margin governing return to the primary regime. While the fallback
is active, the controller restricts emission selection to a
declared safe subset of the operating envelope --- non-limiting
choices include reversion to a low-amplitude additive overlay
under the positive-residual objective of
Section~\ref{sec:baseline-continuous}, suspension of payload
embedding, or rotation to a previously verified protocol-digest
version --- with the fallback entry timestamp, the meter values at
entry, the active fallback policy identifier, and the fallback
exit timestamp committed to the protocol digest. In some
embodiments, the semantic fallback is integrated with the
verification-gating mechanism described in
Section~\ref{sec:agents}, such that actions classified as
RK-verified, hardness-gated, or multi-RK-verified are denied or
escalated while the fallback is active. The fallback regime does
not replace any audit, commitment, or verisimilitude evidence
already accrued before the perception-collapse condition; it
restricts the policy of subsequent emissions until the meter
envelope returns to compliance.

\paragraph{Intervention-gated recurrent-state hygiene
(non-limiting).}
\label{sec:recurrent-state-hygiene}
In some embodiments, recurrent or stateful components of the
controller --- including learned policies with hidden state,
adaptive scan-law generators, and recursive estimators ---
maintain a declared \emph{state hygiene} discipline under which
recurrent state is explicitly resettable, snapshot-able, and
auditable. The hygiene discipline comprises at least: a declared
snapshot-cadence rule committing recurrent-state digests to the
protocol digest at declared intervals; an intervention-gating rule
specifying which controller events --- non-limiting examples
include protocol-digest version rotation, semantic-fallback entry
or exit, manual operator intervention, B-PoD discrepancy
escalation, or any reset operation declared by the audit policy
interface --- trigger committed snapshot, partial reset, or full
reset of declared recurrent-state components; and a recurrent-state
provenance rule under which any synthetic-controller output
produced from recurrent state is provenance-typed identically to
other declared synthetic artefacts in this disclosure (evidence of
the committed input, the declared inference procedure, and the
recurrent-state digest at output time, not direct physical-
provenance evidence). In some embodiments, the snapshot, reset,
and intervention events are themselves committed events of the
protocol digest under the same commitment plumbing of
Section~\ref{sec:hash-chain-state}, so that recurrent-state
hygiene is auditable from the same selectively-openable record as
the rest of the run.


\subsubsection{Co-emission across the separation channel
(non-limiting)}
\label{sec:rt-co-emission}

In some embodiments, the \RT regime shares the
separation channel $\mathsf{Sep}$ (Section~\ref{sec:style-corpus-paired})
with one or more additional declared emission streams operated
concurrently with the styled RT emission. The resulting
\emph{co-emission} architecture supports RT-loop stabilisation via
companion-channel promotion, continuous live paired-evidence
acquisition during deployment, tamper-evidence via committed
companion streams, and multi-regime concurrent operation. The
primary operational use of co-emission is stabilisation of an RT
loop exhibiting degradation, rather than continuous multi-stream
emission.

\paragraph{Base primitive (non-limiting).}
\label{sec:co-emission-base}
In some embodiments, a co-emission configuration comprises
$K_{\mathrm{co}}\geq 2$ independent emission streams carried
concurrently across $\mathsf{Sep}$. For avoidance of doubt, the
symbol $K$ is used in the remainder of the filing with other
meanings (reactor index, coupling-regime label, dynamics
coefficient), and $K_{\mathrm{co}}$ is used specifically for the
number of co-emitted streams throughout this subsubsection and the
corresponding claims. Each stream is drawn from a declared
source-class, drawn from a non-limiting list comprising:
\begin{itemize}
  \item the \RT styled emission of
    Section~\ref{sec:style-parameter};
  \item a \TB probe emission (\TB regime of the present filing);
  \item a \LI probe emission (\LI regime of the present filing);
  \item a declared neutral probe emission drawn from the probe
    family $\mathsf{Probe}$ of Section~\ref{sec:probe-corpus};
  \item a committed cryptographic or pseudorandom generator stream
    (for example a keyed extendable-output function such as BLAKE3
    XOF across declared octaves);
  \item a dedicated human-viewing emission (for example committed
    low-intensity white or pink illumination) whose role is
    naked-eye viewer illumination only, used in sub-embodiments in
    which the \RT is not the human-visible channel.
\end{itemize}
Each stream is committed to the protocol digest with a declared
stream identifier, source-class, amplitude profile, per-stream
recovery rule, detector or demodulation assignment,
residual-crosstalk envelope, and role assignment specifying what
downstream processes (human perception, twin training, verification,
tamper-evidence audit, stabilisation, or combinations thereof) the
stream feeds. A verifier holding the committed stream declarations
and a recording from the declared detector or detector set
assigned to the relevant stream can determine whether each
declared stream was selectively recoverable under the committed
$\mathsf{Sep}$ modality and its residual-crosstalk envelope.

\paragraph{RT as human channel (primary configuration, non-limiting).}
\label{sec:co-emission-rt-as-human}
In preferred embodiments, the \RT styled emission
serves as the human-visible channel, and one or more companion
streams serve computational roles that do not require naked-eye
viewing. Companion streams are declared at amplitudes either
below a declared perceptual-visibility threshold (in the case of
imperceptible companions) or at declared intensities recoverable
by appropriate filter hardware (in the case of filtered-viewer
configurations). The dedicated human-viewing source-class is used
only in sub-embodiments where the \RT is not
currently active; in the default case, the \RT is
itself the human channel.

\paragraph{Amplitude regimes (non-limiting).}
\label{sec:co-emission-amplitude}
In some embodiments, co-emission operates under one of two
declared amplitude regimes, each committed to the protocol digest
along with corresponding perceptual thresholds and recording
recoverability thresholds:
\begin{enumerate}
  \item \textbf{Dark-ambient / photometric-match regime.} The
    \RT styled emission fills the scene's
    illumination budget under a declared photometric-match
    objective; companion streams are additive on top of the style
    and must remain at declared low amplitudes to preserve the
    photometric match, typically below the declared
    perceptual-visibility threshold while remaining above the
    declared recording-recoverability threshold.
  \item \textbf{Bright-ambient / minimum-perturbance regime.} The
    scene is already illuminated by ambient light, the \RT modifies appearance additively, and co-emitted
    companion streams must operate under the positive-residual
    additive-overlay objective of
    Section~\ref{sec:baseline-continuous}. The additive-overlay
    compliance is evaluated on the aggregate observed positive
    excess $\max(\widehat{R}_t - \mathsf{Baseline}_t, 0)$ across
    all $K_{\mathrm{co}}$ co-emitted streams, where
    $\mathsf{Baseline}_t$ is the declared ambient or neutral
    baseline of Section~\ref{sec:baseline-continuous}, not on any
    single named stream, with per-stream perturbance budgets
    logged as auxiliary constraints valid only within the declared
    separation, linearity, and crosstalk envelope. The aggregate
    formulation ensures that a companion stream cannot comply
    per-stream while the combined light violates the bright-
    ambient perturbance objective.
\end{enumerate}

\paragraph{Stabilisation via companion-channel promotion (primary
use-case, non-limiting).}
\label{sec:co-emission-stabilisation}
In preferred embodiments, co-emission is operated in a
state-transitioned manner: the \RT runs in its
ordinary single-stream configuration during steady-state
operation, and a companion stream is \emph{promoted} (brought
into concurrent co-emitted use) when a committed instability
meter $M_{\mathrm{instab}}$ indicates loop degradation. The
companion is \emph{demoted} (returned to inactive status) when a
committed recovery meter indicates loop recovery. The promotion
and demotion policies, instability meter, recovery meter, and all
associated thresholds are committed to the protocol digest.

In some embodiments, $M_{\mathrm{instab}}$ comprises one or more
of the following declared statistics, alone or in combination:
\begin{enumerate}
  \item dual-direction consistency drift from
    Section~\ref{sec:gen-dual-consistency}: $\delta_{\mathrm{fwd}}$,
    $\delta_{\mathrm{inv}}$, or $\delta_{\mathrm{sym}}$ exceeding
    declared thresholds;
  \item verisimilitude-discriminator pass-rate on a rolling window
    dropping below a declared threshold, evaluated under the
    verisimilitude functional of
    Definition~\ref{def:verisimilitude} and the
    discriminator-training composition of
    Section~\ref{sec:verisimilitude-training};
  \item pose-extractor confidence statistics derived from
    $\mathsf{Mask}_t$ of Section~\ref{sec:pose-rt} dropping below
    declared thresholds on a declared evaluation window;
  \item twin residual $r_t$ or $r^{\mathsf{AnaTwin}}_t$ moving
    outside the declared approximation envelope or matching
    envelope of Section~\ref{sec:reality-merge}.
\end{enumerate}

\paragraph{Companion-class stabilisation roles (non-limiting).}
In some embodiments, the stabilising effect of companion-channel
promotion depends on the companion source-class:
\begin{itemize}
  \item A neutral-probe companion (drawn from the probe family of
    Section~\ref{sec:probe-corpus}) provides an independent clean
    scene response at each frame, from which pose features, scene
    state, and twin residuals become estimable even when the
    styled channel produces confusing captures.
  \item A \TB-probe companion provides continuous
    substrate physical-verification evidence; drift outside
    $\mathcal{M}_{\mathrm{accept}}$ on the probe channel is
    direct evidence that the \RT's assumed
    substrate model is in error, distinguishing substrate drift
    from loop drift.
  \item A \LI-probe companion provides continuous
    information-gain measurement about scene state variables,
    resolving observability gaps in the \RT's
    pose-tracking without degrading the styled emission.
  \item A committed cryptographic or pseudorandom stream companion
    provides measurement-side audit only: its recovered statistics
    can be compared against the committed profile to detect
    measurement-chain or sensor drift, but by itself the
    cryptographic companion does not recover pose, scene state, or
    twin residual. In embodiments in which the instability meter
    indicates scene-state or perception-side degradation rather
    than pure measurement-side drift, a cryptographic companion
    alone is insufficient, and the protocol digest declares
    whether the stabilisation procedure instead promotes an
    evidence-bearing companion (neutral, \TB, or \LI)
    or escalates to output-side recovery via the semantic-fallback
    or recurrent-state-hygiene mechanisms of
    Section~\ref{sec:pose-rt}.
\end{itemize}

\paragraph{Promotion/demotion hygiene and state boundaries.}
In some embodiments, each promotion event and each demotion event
is treated as a recurrent-state hygiene boundary for the
\RT controller: the controller's recurrent state is
reset, quarantined, or reweighted under a declared rule committed
to the protocol digest, so that intervention-contaminated hidden
state does not persist across the boundary. This composes with
the intervention-gated recurrent-state hygiene mechanism of
Section~\ref{sec:pose-rt}. In some embodiments, demotion occurs
automatically when the recovery meter crosses a committed
threshold (auto-demotion variant); in other embodiments, demotion
requires external authority approval after the recovery meter
crosses its threshold (authority-gated variant). Both variants
are committed alternatives.

\paragraph{Dual-side recovery composition.}
In some embodiments, the co-emission stabilisation mechanism is
explicitly composed with the output-side recovery mechanisms of
the pose-RT addendum: companion-channel promotion handles the
input side (recovering clean evidence when observation signals
degrade), while semantic-fallback and recurrent-state-hygiene
handle the output side (what to emit when recovery is still in
progress). The two sides are independently committed and trigger
under independent or shared meters as declared in the protocol
digest.

\paragraph{Continuous-operation uses (non-limiting).}
\label{sec:co-emission-continuous-uses}
In some embodiments, co-emission is operated continuously rather
than only during stabilisation. Non-limiting continuous uses
include:
\begin{enumerate}
  \item \textbf{Live $\mathcal{D}_{\mathrm{pair}}$ acquisition
    during deployment.} A neutral-probe companion concurrent with
    a live \RT emission produces paired probe/style
    atoms under any declared $\mathsf{Sep}$ modality. Live
    $\mathcal{D}_{\mathrm{pair}}$ atoms acquired during deployment
    are committed as deployment-origin paired evidence with
    run-time role tags, $\mathsf{Sep}$ modality declaration,
    recovery envelope, and pre-derivative timestamps recorded
    before any curated, distilled, synthetic, or training-corpus
    derivative is generated from them, so that a verifier can
    distinguish deployment-origin evidence from pre-training
    evidence and from any subsequent derivative.
  \item \textbf{Tamper-evidence via committed-statistics companion
    with externally-anchored commitment.} A committed cryptographic
    or pseudorandom companion stream provides tamper evidence when
    the seed commitment is bound to an external anchor
    $\mathsf{Anchor}$ --- for example a pre-run hash-chain
    anchored by a verifier-supplied nonce, a public transparency
    log entry, a trusted timestamp, or an equivalent external
    commitment witnessed before the capture window begins. The
    seed is revealed only after the committed window closes. Local
    operator-side commitment without $\mathsf{Anchor}$ is
    disclosed only as a reduced-assurance sub-embodiment; the
    primary embodiment requires external anchoring to prevent
    operator backdating, suppression, or replacement of the
    commitment.
  \item \textbf{Multi-style per-viewer.} Two or more \RT streams with different style parameters $\phi^{(i)}$
    are carried on different subchannels, slots, codes, phases,
    bands, or filter subscriptions under the declared
    $\mathsf{Sep}$ modality, such that
    goggle-equipped viewers subscribed to different streams
    experience different styled content of the same underlying
    scene. Per-stream style-parameter identifiers, goggle filter
    specifications, and multi-stream amplitude balance are
    committed to the protocol digest.
  \item \textbf{Secure recording with triple consistency.} A
    committed cryptographic companion stream is co-emitted with a
    \RT styled emission; the recording contains the
    superposed signal. A verifier holding the recording and the
    revealed seed recovers the cryptographic channel via
    demodulation under the committed $\mathsf{Sep}$ recovery rule
    and additionally applies $\mathsf{G}_{\mathrm{destyle}}$ of
    Section~\ref{sec:gen-inverse-destyle} to the recovered styled
    channel (obtained via $\mathsf{Sep}$ demodulation or a declared
    prefiltered capture after crypto-channel removal, validated
    under the committed residual-crosstalk envelope) to produce a
    synthetic neutral view, yielding three independent consistency
    sources (the committed styled atoms, the recovered
    cryptographic channel, and the
    $\mathsf{G}_{\mathrm{destyle}}$-generated synthetic neutral
    view). The provenance asymmetry disclosed in
    Section~\ref{sec:gen-inverse-destyle} is preserved: the
    synthetic neutral view remains a generator-produced
    corroborand of declared provenance type, not a substitute for
    opened physical companion-stream evidence.
\end{enumerate}

\paragraph{Cross-regime co-emission (non-limiting).}
\label{sec:co-emission-cross-regime}
In some embodiments, co-emission is operated in cross-regime
configurations that may omit a human-visible channel or may
combine a computational probe stream with a \RT
stream. Non-limiting cross-regime configurations include:
\begin{itemize}
  \item \TB probe concurrent with \LI probe, providing
    dual-regime verification on a shared substrate with no
    human-visible channel;
  \item \TB probe concurrent with \RT
    styled emission, providing continuous physical-verification
    evidence during styled operation;
  \item \LI probe concurrent with \RT styled
    emission, providing concurrent information-gain measurement
    during styled operation.
\end{itemize}
Each cross-regime configuration inherits the commitment and
recovery discipline of the base primitive.

\paragraph{Goggle-based access control (non-limiting).}
\label{sec:co-emission-goggles}
In some embodiments, co-emission is combined with viewer filter
hardware (goggles, polariser eyewear, IR-blocking eyewear, etc.)
to produce declared access-control behaviours. Goggle
specifications (filter type, passband, shutter timing, polariser
orientation, spectral cutoff) are committed to the protocol
digest alongside the $\mathsf{Sep}$ modality. Multi-goggle-set
configurations permit different viewers to subscribe to different
subsets of the co-emitted streams. In embodiments where a
naked-eye viewer (without goggles) is expected to be present, the
aggregate amplitude profile across all co-emitted streams is
constrained to keep the unfiltered naked-eye experience within a
declared safety and navigability envelope, committed to the
protocol digest.

\paragraph{Declared bleed-through envelope (non-limiting).}
For avoidance of doubt, goggle-based perceptual filtering is
never absolute: an unfiltered viewer receives residual photons
from filter-intended channels through ambient scatter,
imperfect filtering, viewing-angle dependence, or
scene-induced modality changes (depolarisation, spectral
shift). The protocol digest commits a declared bleed-through
envelope describing the expected residual amplitudes perceived by
an unfiltered viewer; claims of imperceptibility are always
relative to this declared envelope rather than absolute.

\paragraph{Independence discipline for composed tamper-evidence
(non-limiting).}
In some embodiments, co-emission composes with the
payload-embedding \RT of
Section~\ref{sec:rt-payload} to produce three orthogonal
tamper-evidence mechanisms operating simultaneously: the
prior-frame-hash payload chain on the styled channel, the
run-level hash-chain state $\chi_t$ of
Section~\ref{sec:hash-chain-state}, and the independently
recoverable cryptographic companion stream. In such embodiments,
the three mechanisms are declared under independent seeds,
domain separators, and key-derivation roots committed to the
protocol digest, so that compromise of one mechanism's keying
material does not compromise the others.

\paragraph{Reality-merge composition (non-limiting).}
In some embodiments, role-authorised companion-stream residuals
are logged as reality-merge evidence under
Section~\ref{sec:reality-merge}, and may drive twin refinement,
analogue-twin re-matching under the matching-fidelity envelope of
Section~\ref{sec:analogue-twin}, or successor protocol-digest
version updates, according to the companion stream's declared role
assignment in the protocol digest.

\paragraph{Scope note (non-limiting).}
For avoidance of doubt, the co-emission primitive uses the same
$\mathsf{Sep}$ infrastructure as the paired-capture training-data
acquisition of Section~\ref{sec:style-corpus-paired}. The
disclosed technical contribution in this subsubsection is the
joint specification of declared source-class composition,
committed per-stream roles, amplitude regimes composing with the
additive-overlay objective, committed instability meter and
promotion/demotion policy for RT-loop stabilisation, external
commitment anchoring for cryptographic companions, and
composition with payload-embedding RT, $\mathsf{G}_{\mathrm{destyle}}$,
semantic fallback, and recurrent-state hygiene, rather than the
underlying separation-channel hardware itself.


In some embodiments, where the \RT subsystem produces
articulated content under a declared deployment scope, the
trajectory-domain classifier family further includes pose-conditioned
classifiers whose inputs comprise committed pose-feature traces,
pose-distribution residuals, and pose-dynamics residuals derived
from the pose-conditioned \RT embodiments of
Section~\ref{sec:pose-rt}, with the pose-feature extractor,
pose-reference-set provenance type, negative-edit families, and
residual definitions committed to the session protocol digest.
Companion governance embodiments described in the related Filing 2 application may
consume these classifier outputs and committed evidence types under
their own state and audit-policy machinery; such consumption does
not require any state, mode, or governance taxonomy from the
related Filing 2 application to be present in the apparatus disclosed here.


In some embodiments, the style parameter $\phi$ may include an
articulated-content component, such as a pose class, action class,
kinematic intent, or reference-motion digest, in which case
pose-conditioned \RT embodiments may evaluate
additional pose-distribution and pose-dynamics terms as described
in Section~\ref{sec:pose-rt}. In some embodiments, the declared
\RT objective may additionally carry a declared
per-frame payload whose extractability is enforced as part of the
compound objective, as described in Section~\ref{sec:rt-payload}.
In some embodiments, training of any such specialised \RT objective is conducted against a declared digital twin
under the reality-merge pipeline of Section~\ref{sec:twin-pipeline}.


\subsection{Digital twin and reality-merge pipeline (non-limiting)}
\label{sec:twin-pipeline}

In some embodiments, the physical \RK and its derived
models are coupled through a pipeline that proceeds, end-to-end,
through four stages:
(i)~acquisition of a committed probe corpus under declared spectrally
neutral or spectrally shaped illumination with declared
temporal-correlation structure
(Section~\ref{sec:probe-corpus});
(ii)~training or matching of a twin from the probe corpus, where
the twin may be a digital twin (Section~\ref{sec:twin-training}),
an analogue twin comprising a second physical \RK module
operated in twin-role (Section~\ref{sec:analogue-twin}), or a
hybrid digital--analogue twin;
(iii)~training of a \RT controller against the twin
under a compound objective that may include pose-conditioned terms
(Section~\ref{sec:pose-rt}), payload-extractability terms
(Section~\ref{sec:rt-payload}), verisimilitude
(Definition~\ref{def:verisimilitude}), edge-structural terms
(Section~\ref{sec:edge-meters}), and meter-envelope penalties
(Definition~\ref{def:meter-envelope}); and
(iv)~reality merge, in which the trained controller is deployed to
the physical substrate and the residuals between twin predictions
and physical observations feed back into verisimilitude discriminator
training (subject to declared anomaly classification), twin
refinement, and successor protocol-digest version updates
(Section~\ref{sec:reality-merge}).

The twin is assistive rather than authoritative: the physical
\RK remains the arbiter for verification, hardness, and
meter-bounded validity. The pipeline's correctness discipline is
that every stage emits committed, auditable evidence under a
versioned protocol digest, so that a verifier can recompute the
twin's identity, the probe-corpus contents, the RT training
schedule, and the reality-merge residuals from selectively opened
\cba windows.

% ---------------------------------------------------------------------
\subsubsection{Probe-corpus acquisition (non-limiting)}
\label{sec:probe-corpus}
% ---------------------------------------------------------------------

In some embodiments, a \emph{probe corpus}
$\mathcal{D}_{\mathrm{probe}}$ is a committed set of input-output
pairs $(u^{\mathsf{Probe}}_{0:T},\mathbf{y}_{0:T})$ in which the
emission sequence $u^{\mathsf{Probe}}_{0:T}$ is drawn from a
declared probe family $\mathsf{Probe}$ rather than from a
semantically structured RT target, and in which the corresponding
reactor or scene response is recorded through the ordinary \cba
pipeline. The probe corpus provides the twin with training evidence
that samples declared regions of the reactor's transfer-function
support under a committed coverage rule, without imposing semantic
bias on the twin.

\paragraph{Probe families (non-limiting).}
In some embodiments, the probe family $\mathsf{Probe}$ is specified
along two independent axes:
\begin{enumerate}
  \item \textbf{Spectral / spatial shape.} The illumination may be
    (a)~spectrally neutral, meaning equalised within declared
    calibration tolerances across the reactor's sensitive spectral
    band and declared spatial channels, so that no wavelength or
    channel is preferentially excited beyond the committed tolerance
    envelope; (b)~pink-noise-shaped, with
    $1/f^{\mathsf{Exp}_{\mathrm{spec}}}$ power spectral density for
    a declared exponent $\mathsf{Exp}_{\mathrm{spec}}$ (for example
    $\mathsf{Exp}_{\mathrm{spec}}=1$ for canonical pink noise,
    $\mathsf{Exp}_{\mathrm{spec}}=2$ for red noise, and
    $\mathsf{Exp}_{\mathrm{spec}}=0$ for white); (c)~custom-shaped
    under a declared filter or code family; or (d)~a declared
    mixture of the above.
  \item \textbf{Temporal-correlation structure.} The emission
    sequence may be (a)~temporally white (independent samples);
    (b)~temporally pink, with
    $1/f^{\mathsf{Exp}_{\mathrm{temp}}}$ temporal power spectral
    density for declared $\mathsf{Exp}_{\mathrm{temp}}$, yielding
    long-range frame-to-frame correlations; (c)~driven by a declared
    autoregressive or state-space generator; or (d)~constructed
    from concatenated segments with declared segment-boundary
    statistics.
\end{enumerate}
The spectral shape parameters, temporal shape parameters, sample
rate, frame rate, total duration, random-seed policy, and any
calibration pre-sequence are committed to the protocol digest as
part of the probe-corpus identity. In some embodiments, temporally
pink correlation is preferred over white because it delivers equal
power per octave, exercising slow reactor modes that white noise
undersamples in any fixed total observation budget.

\paragraph{Neutrality and semantic non-bias (non-limiting).}
In preferred embodiments, the probe corpus is drawn from a probe
family whose spectral, spatial, and temporal statistics are
declared neutral with respect to any downstream semantic category
(pose class, action class, style class, or payload pattern), so
that a twin trained on $\mathcal{D}_{\mathrm{probe}}$ does not
encode semantic priors that would bias downstream RT training. In
some embodiments, semantic non-bias is verified by training
declared probe classifiers on $\mathcal{D}_{\mathrm{probe}}$ and
confirming that their discrimination accuracy for the target
semantic categories remains at or near chance, with the verification
protocol and threshold committed to the protocol digest.

\paragraph{Commitment and provenance (non-limiting).}
In some embodiments, $\mathcal{D}_{\mathrm{probe}}$ is committed as
a Merkle-rooted set of disclosure atoms
(Section~\ref{sec:merkle-atoms}), with per-atom digests bound to
the probe-family declaration and to the run-level hash-chain state
$\chi_t$ (Section~\ref{sec:hash-chain-state}) at the acquisition
time. The protocol digest provenance-types
$\mathcal{D}_{\mathrm{probe}}$ as committed physical evidence, and
verifier-facing attestations may rely on
$\mathcal{D}_{\mathrm{probe}}$ as a physical-provenance reference
for subsequent twin identity, RT-training audits, and
empirical-hardness arguments.

% ---------------------------------------------------------------------
\subsubsection{Twin training from the probe corpus (non-limiting)}
\label{sec:twin-training}
% ---------------------------------------------------------------------

In some embodiments, a digital twin $\mathsf{Twin}$ is trained to
predict the reactor or scene response to declared emissions at the
same step:
\[
  \widehat{\mathbf{y}}_t
  =
  \mathsf{Twin}\!\bigl(u_{0:t},\,\mathrm{meta}_t,\,\Pi\bigr),
  \qquad
  C^{\mathsf{Twin}}_{0:T}
  =
  \bigl\{(u(t),\widehat{\mathbf{y}}_t)\bigr\}_{t=0}^{T}.
\]
$C^{\mathsf{Twin}}_{0:T}$ is the \emph{twin-predicted \cb} (synthetic bundle), constructed step-wise from declared
emissions and twin-predicted observations so that downstream
components such as $\mathsf{RTDec}$, $\mathsf{Pose}$,
$\mathsf{Extract}$, the verisimilitude functional, and the
edge-based structural meter may be evaluated on twin-predicted
windows $C^{\mathsf{Twin}}_{\mathsf{Win}_t}$ during twin-side
training.
The protocol digest provenance-types
$C^{\mathsf{Twin}}_{0:T}$ as twin-predicted or synthetic, and
verifier-facing attestations shall not treat
$C^{\mathsf{Twin}}_{0:T}$ as physical-provenance evidence except as
a recomputable surrogate derived from a committed twin identity and
declared input schedule.

$\mathsf{Twin}$ may be (a)~a learned autoencoder whose encoder maps
emissions to a latent dynamics state and whose decoder predicts the
response; (b)~a transfer-function model whose parameters are fitted
to the probe corpus by system-identification techniques (for example
subspace identification, prediction-error methods, or Gaussian-process
state-space models); (c)~a physics-based propagation surrogate as
described below; or (d)~a declared hybrid. In preferred embodiments,
the twin is differentiable so that gradients of a downstream RT
objective with respect to emissions may be propagated through the
twin during RT training.

\paragraph{Training objective on the probe corpus (non-limiting).}
In some embodiments, $\mathsf{Twin}$ is trained by minimising a
declared same-step prediction loss on $\mathcal{D}_{\mathrm{probe}}$:
\[
  \min_{\mathsf{Twin}}\;
  \mathbb{E}_{(u^{\mathsf{Probe}}_{0:T},\mathbf{y}_{0:T})
    \sim \mathcal{D}_{\mathrm{probe}}}
  \Bigl[
  \sum_{t=0}^{T}
  \ell_{\mathrm{twin}}\!\bigl(
    \mathsf{Twin}(u^{\mathsf{Probe}}_{0:t},\mathrm{meta}_t,\Pi),
    \,\mathbf{y}_{t}\bigr)
  \Bigr],
\]
where $\ell_{\mathrm{twin}}$ is a declared prediction loss (for
example squared error, negative log-likelihood under a declared
observation model, or a learned perceptual loss). The twin's
identity, architecture, weights or version hash, input fields,
training set reference, validation set reference, and
approximation-quality metrics on held-out probe-corpus windows are
committed to the protocol digest.

\paragraph{Physics-based propagation surrogates (non-limiting).}
In some embodiments, the digital twin includes a physics-based
wave-propagation engine (for example ray tracing, beam-propagation,
or wave-optics solvers) together with learned priors for speckle
statistics, aberrations, or alignment drift. Such hybrids explore
large regions of configuration space in simulation, after which
hardware-in-the-loop optimisation fine-tunes around promising
configurations.

\paragraph{Approximation bounds and assistive status (non-limiting).}
In some embodiments, the twin's approximation quality is
characterised by a declared envelope on held-out probe-corpus
windows and, where applicable, on held-out live recordings. The
twin is used as a differentiable or predictive surrogate during
training and protocol selection, but downstream validity claims
rest on physical \RK evaluation and meter-envelope
discipline, not on the twin alone.

% ---------------------------------------------------------------------
\subsubsection{Analogue twin: a second physical kernel in twin-role
(non-limiting)}
\label{sec:analogue-twin}
% ---------------------------------------------------------------------

In some embodiments, the twin is realised as an \emph{analogue twin}
$\mathsf{AnaTwin}$ comprising a second physical \RK
module whose reactor configuration, scan law, operating point, and
invariance profile are committed and matched to the primary kernel
through a shared probe corpus. Unlike the digital twin, the
analogue twin produces a genuine \cb
\[
  C^{\mathsf{AnaTwin}}_{0:T}
  =
  \bigl\{(u(t),\mathbf{y}^{\mathsf{AnaTwin}}_t)\bigr\}_{t=0}^{T}
\]
with physical provenance: each atom binds a declared emission to a
physically observed response through the analogue-twin module's own
emitter, reactor, and detector. Consequently,
$C^{\mathsf{AnaTwin}}_{0:T}$ is provenance-typed by the protocol
digest as committed physical evidence of the analogue-twin module's
own emitted-and-observed run, and is not subject to the
synthetic-surrogate restriction applicable to $C^{\mathsf{Twin}}_{0:T}$
of Section~\ref{sec:twin-training}. For verifier-facing purposes,
$C^{\mathsf{AnaTwin}}_{0:T}$ is physical-provenance evidence only of
the analogue-twin module's own run; it is evidence of the primary
kernel's expected response only as a matched cross-kernel surrogate
within the committed matching-fidelity envelope
$\mathcal{M}_{\mathrm{match}}$, and it is not physical-provenance
evidence or empirical-hardness support for the primary kernel
unless independently corroborated by primary-kernel committed
evidence.

\paragraph{Matching procedure (non-limiting).}
In some embodiments, the analogue twin is matched to the primary
kernel by executing a shared declared probe family
$\mathsf{Probe}$ (Section~\ref{sec:probe-corpus}) on both kernels
and fitting operating-point parameters to minimise a declared
cross-kernel prediction loss on held-out probe windows:
\[
  \min_{\theta^{\mathsf{AnaTwin}}}\;
  \mathbb{E}\!\left[
    \ell_{\mathrm{match}}\!\bigl(
      \mathbf{y}^{\mathsf{AnaTwin}}_t,\,\mathbf{y}_t\bigr)
  \right],
\]
where $\theta^{\mathsf{AnaTwin}}$ comprises the analogue-twin
kernel's declared adjustable parameters (for example gain,
bias, alignment, reactor operating point, or scan-law parameters)
and $\ell_{\mathrm{match}}$ is a declared prediction or
distributional distance. The matching procedure, probe corpus,
fitted parameter values, and residual-error statistics on
held-out windows are committed to the protocol digest.

\paragraph{Matching-fidelity envelope (non-limiting).}
In some embodiments, the analogue twin's fitness for twin-role
use is characterised not by a digital approximation envelope but
by a \emph{cross-device matching-fidelity envelope}
$\mathcal{M}_{\mathrm{match}}$, a declared admissibility bound on
residuals computed from paired committed probe-corpus bundles.
$\mathcal{M}_{\mathrm{match}}$ is a cross-device admissibility
envelope, distinct from the single-kernel meter envelope
$\mathcal{M}_{\mathrm{accept}}$ of
Definition~\ref{def:meter-envelope}, but it is declared in the
protocol digest and bound into commitments using the same
declaration and commitment mechanics as that definition. An
analogue twin is in-envelope when cross-kernel residuals remain
within $\mathcal{M}_{\mathrm{match}}$ across a declared evaluation
window.

\paragraph{Reality Transform training against an analogue twin.}
In some embodiments, the \RT controller of
Section~\ref{sec:rt-against-twin} is trained against
$C^{\mathsf{AnaTwin}}_{\mathsf{Win}_t}$ by emitting declared
controls $u(t)$ on the analogue-twin module and computing the
compound objective's evidence-dimension terms from the resulting
genuine bundle. Because $C^{\mathsf{AnaTwin}}$ is physical,
downstream $\mathsf{RTDec}$, $\mathsf{Pose}$, $\mathsf{Extract}$,
and verisimilitude evaluation on
$C^{\mathsf{AnaTwin}}_{\mathsf{Win}_t}$ do not require the
synthetic-window construction rule of
Section~\ref{sec:rt-against-twin} and are not subject to the
twin-predicted-bundle verifier-facing restriction. Gradients
through the analogue twin are, in general, not analytically
available; hardware-in-the-loop optimisation
(Section~\ref{sec:hitl-optimisation}) is used instead of
backpropagation through the twin.

\paragraph{Reality merge with an analogue twin.}
In some embodiments, when the trained controller is deployed to
the primary kernel, the reality-merge residual acquires a
different interpretation than in the digital-twin case. For an
analogue twin,
\[
  r^{\mathsf{AnaTwin}}_t
  =
  \mathbf{y}_t - \mathbf{y}^{\mathsf{AnaTwin}}_t
\]
measures drift, wear, or tamper between matched physical
substrates rather than model-approximation error. Systematic
$r^{\mathsf{AnaTwin}}_t$ outside $\mathcal{M}_{\mathrm{match}}$
triggers re-matching of the analogue twin, replacement or
recalibration of the analogue-twin module, or a successor
protocol-digest version recording the change, under the
applicable authority-separation policy. In some embodiments,
$r^{\mathsf{AnaTwin}}_t$ is additionally used as an input to the
verisimilitude-discriminator training of
Section~\ref{sec:verisimilitude-training}, under the same declared
anomaly-classification procedure that governs digital-twin
residuals.

\paragraph{Cross-attestation composition (non-limiting).}
In some embodiments, the analogue twin simultaneously serves as a
cross-attestation witness under the multi-RK cross-attestation
protocols of Section~\ref{sec:networks} and, where applicable,
as a B-PoD witness under the bilateral
Proof-of-Discrepancy protocol of Section~\ref{sec:bpod}. The
analogue-twin role and the cross-attestation witness role are
recorded in the protocol digest as separate capability
declarations so that a single module may participate in multiple
roles without role conflation in verifier-facing attestations.

\paragraph{Hybrid digital--analogue twin (non-limiting).}
In some embodiments, the digital twin of
Section~\ref{sec:twin-training} and the analogue twin of this
section are operated concurrently, with the controller trained
against both under declared weightings. This composition is a
specialisation of the hybrid digital--analogue co-training
discussed in connection with hardware-in-the-loop optimisation
(Section~\ref{sec:hitl-optimisation}): digital-twin gradients
support fast iteration; analogue-twin evidence supports
physically grounded validation; the pair together reduces the
reliance of any single validity claim on either digital
approximation quality or single-kernel physical provenance.
The weightings between digital-twin and analogue-twin terms, and
the triggering rules for escalation from digital-twin-only to
combined operation, are committed to the protocol digest.

\paragraph{Scope note (non-limiting).}
For avoidance of doubt, the analogue twin of this section is a
\RK module as already disclosed; its novelty as
claimed in the claim set of the related Filing 2 application lies in the joint
specification of role (twin-role rather than primary-role),
matching procedure, matching-fidelity envelope,
analogue-twin-specific reality-merge interpretation, and
composition with digital-twin and cross-attestation machinery,
rather than in any new physical kernel architecture.

% ---------------------------------------------------------------------
\subsubsection{Reality Transform training against the twin
(non-limiting)}
\label{sec:rt-against-twin}
% ---------------------------------------------------------------------

In some embodiments, a \RT controller is trained by
optimising a compound objective in which the physical observation
$\mathbf{y}_t$ is replaced, during training, by a twin-produced
observation: $\widehat{\mathbf{y}}_t$ from a digital twin
(Section~\ref{sec:twin-training}), $\mathbf{y}^{\mathsf{AnaTwin}}_t$
from an analogue twin (Section~\ref{sec:analogue-twin}), or a
declared weighted combination of the two. In the digital-twin case,
gradients of the compound objective flow through $\mathsf{Twin}$ to
the controller; in the analogue-twin case, gradients are estimated
by hardware-in-the-loop optimisation
(Section~\ref{sec:hitl-optimisation}) unless a differentiable
analogue-twin model is additionally maintained.
During twin-side RT training, any decoder, pose-feature extractor,
payload extractor, verisimilitude functional, or structural meter
that would ordinarily consume a physical \cba window
is evaluated on the corresponding twin-produced bundle window --
$C^{\mathsf{Twin}}_{\mathsf{Win}_t}$ for the digital twin or
$C^{\mathsf{AnaTwin}}_{\mathsf{Win}_t}$ for the analogue twin --
with the synthetic-window construction and any domain-gap correction
(for the digital twin) or the matching procedure and
matching-fidelity envelope (for the analogue twin) committed to the
protocol digest.

The compound objective aggregates declared terms across evidence
dimensions:
\[
  \mathcal{J}_{\mathrm{compound}}(U_{0:T};\phi)
  =
  \sum_{s\in\mathsf{RTTerm}}
    \lambda_{\mathrm{rt},s}\,\mathcal{L}_s
  +
  \lambda_{\mathrm{meter}}\,
    \mathsf{MeterPen}\!\bigl(\mathbf{m}_{0:T};
    \mathcal{M}_{\mathrm{accept}}\bigr),
\]
where $\mathsf{RTTerm}$ is a declared candidate set of
evidence-dimension terms, $\mathcal{J}_{\mathrm{compound}}$ is
minimised unless the protocol digest declares an equivalent
negative-reward convention, and $\mathsf{MeterPen}$ is a declared
meter-envelope-violation aggregate (such as a per-step indicator
sum, hinge loss, or smooth barrier penalty).
Terms in $\mathsf{RTTerm}$ are drawn from:
(i)~the style-target term of Section~\ref{sec:style-parameter};
(ii)~the pose-distribution and pose-dynamics terms of
Section~\ref{sec:pose-rt};
(iii)~the payload-extractability term of
Section~\ref{sec:rt-payload};
(iv)~the verisimilitude term of
Definition~\ref{def:verisimilitude};
(v)~the edge-based structural term of
Section~\ref{sec:edge-meters}; and
(vi)~any further declared meter term.
The weights $\{\lambda_{\mathrm{rt},s}\}$ and the meter-partition
policy are committed to the protocol digest. In some embodiments,
the protocol digest identifies active terms by assigning nonzero
$\lambda_{\mathrm{rt},s}$ values to selected members of
$\mathsf{RTTerm}$ and zero values to inactive members;
pose-conditioned RT and payload-embedding RT are two non-limiting
such specialisations.

\paragraph{Twin-to-physical transfer (non-limiting).}
In some embodiments, a controller trained against $\mathsf{Twin}$
is subsequently deployed on the physical substrate, and a declared
fine-tuning phase updates the controller against physical
observations under the same compound objective. The fine-tuning
schedule, permissible update magnitude, and twin-to-physical
transfer meter are committed to the protocol digest.

% ---------------------------------------------------------------------
\subsubsection{Reality merge: deployment to physical and residual
feedback (non-limiting)}
\label{sec:reality-merge}
% ---------------------------------------------------------------------

In some embodiments, once the \RT controller is
deployed on the physical substrate, the twin is retained as a
residual reference. At each step, the residual
\[
  r_t
  =
  \mathbf{y}_t - \widehat{\mathbf{y}}_t
  =
  \mathbf{y}_t
  -
  \mathsf{Twin}(u_{0:t},\mathrm{meta}_t,\Pi)
\]
is computed and logged. The residual feeds three downstream
processes:

\paragraph{Verisimilitude discriminator training.}
In some embodiments, residuals outside the declared twin
approximation envelope are logged as candidate anomalies or
out-of-model physical examples. Such residuals are used as negative
examples for the verisimilitude-discriminator training of
Section~\ref{sec:verisimilitude-training} only if a declared
post-hoc classification procedure assigns them to a forged, edited,
simulated, or otherwise non-genuine class; otherwise they are used
for twin refinement or meter-envelope revision.

\paragraph{Twin refinement.}
In some embodiments, committed residuals are used to update the
twin by declared continual-learning or re-fitting procedures, with
each twin update producing a new versioned twin identity committed
to the protocol digest as a successor protocol-digest version. The
prior twin is retained as part of the committed version history so
that verifiers may re-evaluate earlier runs under the twin active at
the time.

\paragraph{Successor protocol-digest updates.}
In some embodiments, systematic residuals trigger successor
protocol-digest version updates (for example new probe-corpus
acquisition, new meter envelope, or new RT fine-tuning schedule)
under the applicable authority-separation policy; the prior
protocol digest remains the operative digest for runs it already
bound.

\paragraph{Relation to hardware-in-the-loop optimisation.}
The reality-merge step composes with the hardware-in-the-loop
optimisation of Section~\ref{sec:hitl-optimisation}: twin-trained
controllers provide initial emission policies; hardware-in-the-loop
optimisation refines them against physical evidence; residual
feedback refines the twin; the cycle iterates under declared
schedules.

% ---------------------------------------------------------------------
\subsubsection{Sandboxed and live operation (non-limiting)}
\label{sec:twin-sandbox}
% ---------------------------------------------------------------------

In some embodiments, probe-corpus acquisition, twin training, and
\RT training against the twin are performed under the
parameter-sandboxed offline-session mechanism of
Section~\ref{sec:offline-sim}. Parameter updates made to the
physical kernel during the sandboxed session are reverted on return
to live operation; $\mathsf{Twin}$, $\mathsf{Embed}$,
$\mathsf{Extract}$, pose-feature extractors, pose-distribution
discriminators, kinematic classifiers, and verisimilitude
discriminators trained from sandboxed evidence may be retained as
derived models where permitted by the protocol digest and meter
policy.


\subsection{Styled training-data generation methods (non-limiting)}
\label{sec:styled-training-data}

In some embodiments, training data for downstream components
(digital twin, analogue-twin matching, verisimilitude
discriminator, pose-distribution discriminator $A_{\mathrm{dat}}$,
pose-dynamics model $g_{\mathrm{dyn}}$, kinematic-plausibility
classifier, and the \RT controller itself) is
generated not only from neutral probe-corpus acquisition
(Section~\ref{sec:probe-corpus}) but also from one or more
\emph{styled training corpora} in which the emitted content
carries declared style, aesthetic, or target-distribution
structure. The following subsubsections describe seven
non-limiting styled-corpus generation mechanisms, each
independently claimable, together with combined configurations.
Unlike the probe corpus of Section~\ref{sec:probe-corpus}, styled
training corpora are semantic-content-by-design and their
non-neutrality is declared and accepted in the protocol digest
rather than verified against; semantic-non-bias testing does not
apply to styled corpora.

% ---------------------------------------------------------------------
\subsubsection{Externally-sourced curated styled corpus (non-limiting)}
\label{sec:style-corpus-curated}
% ---------------------------------------------------------------------

In some embodiments, a declared styled corpus
$\mathcal{D}_{\mathrm{style}}$ comprises externally-sourced
materials (images, video sequences, art assets, authored
aesthetic references) selected as representative of the target
RT style family. The corpus is identified by content hash, source
attribution, declared style-family identifier, and any declared
temporal-correlation structure of source sequences, all committed
to the protocol digest. An externally-sourced curated corpus is
operator-declared and does not by itself supply physical-provenance
or empirical-hardness evidence; verifier-facing attestations shall
not treat it as physical-provenance evidence unless the acquisition
and validation of the source material are separately committed.

% ---------------------------------------------------------------------
\subsubsection{Unpaired iterated acquisition with committed reset
(non-limiting)}
\label{sec:style-corpus-iterated-reset}
% ---------------------------------------------------------------------

In some embodiments, a styled corpus is generated by an
unpaired iterated procedure: an initial noisy transfer function or
roughly-trained styling component is applied to declared seed
content to produce candidate styled emissions; those emissions are
projected through the physical substrate and captured, producing
committed (emission, observation) atoms; the committed atoms feed
supervised training of the next cycle's styling component. Because
the only training signal in each cycle derives from the prior
cycle's own captures, the procedure is subject to a compounding-drift
failure mode: mode collapse, distribution drift, or
verisimilitude-pass-rate decline accumulates across cycles.

\paragraph{Committed degradation detector and reset protocol.}
In some embodiments, a declared degradation detector monitors
committed per-cycle statistics (for example style-distribution
distance from a declared reference, verisimilitude-pass rate on
held-out evaluation windows, mode-collapse proxy such as effective
latent-space coverage) and triggers a reset when one or more
statistics cross committed thresholds. On reset, the procedure
re-seeds from a fresh neutral probe corpus or from a declared
bootstrap corpus, producing a new training trajectory independent
of the degraded one. The detection statistics, thresholds, reset
policy, and per-cycle protocol digests (chained through
$\chi_t$ of Section~\ref{sec:hash-chain-state}) are committed.

\paragraph{Relationship to reinforcement-learning-based iterative
improvement (non-limiting).}
In some embodiments, convergent iterative improvement of a styling
component under a reward signal is conducted as a reinforcement-
learning procedure with declared committed trajectories, rewards,
and value estimates; such procedures are distinct from the
supervised one-shot-and-reset mechanism described in this
subsubsection and are not required by it.

% ---------------------------------------------------------------------
\subsubsection{Separation-channel paired probe/style acquisition
(non-limiting)}
\label{sec:style-corpus-paired}
% ---------------------------------------------------------------------

In some embodiments, a paired corpus
$\mathcal{D}_{\mathrm{pair}}$ is acquired by projecting a declared
neutral probe emission $u^{\mathrm{probe}}$ and a declared styled
emission $u^{\mathrm{style}}$ simultaneously onto the same scene
through a declared separation channel $\mathsf{Sep}$ that
preserves scene, geometry, lighting, and sensor state across the
probe and style captures within a declared invariance tolerance.
A single detector (or a detector-pair matched by the separation
channel) captures both channels, and the separation property
permits recovery of the two captures as paired
$(\mathbf{y}^{\mathrm{probe}}_t,\mathbf{y}^{\mathrm{style}}_t)$
atoms bound to the same scene instant.

\paragraph{Separation modalities (non-limiting).}
In some embodiments, $\mathsf{Sep}$ is realised under a declared
separation modality. The modalities are partitioned into preferred
embodiments and alternative embodiments, where the alternatives
carry declared bleed or confound envelopes reflecting physical
limitations. Preferred realisations of $\mathsf{Sep}$ include:
\begin{enumerate}
  \item \textbf{Temporal-interleave (preferred).} Probe and style
    emitted in interleaved frame slots at rates above a declared
    perceptual envelope, with phase-locked frame-synchronous
    capture recovering the two sequences. Robust against bleed
    when detector-shutter synchronisation is preserved within a
    declared synchronisation tolerance.
  \item \textbf{IR/RGB (preferred).} Companion channel in infrared
    or near-infrared wavelengths outside a declared human visual
    response range, captured by IR-sensitive or NIR-sensitive
    detection, with RT styled emission in the visible range.
    Non-limiting hardware sub-embodiments include: (a) two
    separate emitters, one visible-only and one IR-only, sharing
    scene illumination; (b) a single emitter with visible and IR
    output separated by sensor spectral response; (c) a
    wavelength-tagged single emission with dichroic capture
    recovering the two wavelength ranges.
\end{enumerate}
Alternative realisations of $\mathsf{Sep}$, each subject to a
declared bleed, confound, or geometric-restriction envelope
committed to the protocol digest, include:
\begin{enumerate}
  \setcounter{enumi}{2}
  \item \textbf{Spatial (geometrically-disjoint coverage).} Two or
    more projection sources addressing geometrically disjoint
    regions of the scene, such that each scene point is addressed
    by exactly one projector under a declared coverage rule
    enforced by a committed scene-region mask, occlusion model,
    or surface-partition map, together with a rigid-body or
    known-geometry scene assumption. Non-limiting example:
    rear-projection onto
    a robot or mannequin surface for calibration or probe
    acquisition while front-projection delivers styled content;
    the scene-region mask or surface partition over the rigid-body
    geometry ensures disjoint coverage. The
    geometric-restriction envelope covering scene-geometry
    assumptions, mask/occlusion/partition identity, and any
    boundary-overlap tolerance is committed
    to the protocol digest. This modality does not cover per-eye
    stereo goggles, which are disclosed below as a deployment
    pattern across temporal, polarisation, and chromatic
    modalities rather than as a spatial separation.
  \item \textbf{Code-division on a shared beam path.} Probe and
    style carried by orthogonal modulation codes on a single beam
    path, with demodulated capture recovering the two
    code-separated channels. Requires declared modulation-bandwidth
    and demodulation-sophistication envelopes; shares the beam path
    entirely and eliminates geometric parallax between the
    captures.
  \item \textbf{Polarization (alternative; declared bleed).}
    Orthogonally-polarized probe and style emissions with
    polarization-resolved detection. Scene-induced depolarisation,
    imperfect polariser stack-ups, and viewing-angle-dependent
    bleed require a declared residual-crosstalk envelope. Not
    suitable for co-emission configurations in which the
    companion must be imperceptible to naked-eye viewers without
    declared tolerance for visible bleed.
  \item \textbf{Chromatic / spectral-within-visible (alternative;
    declared confound).} Probe and style emitted in
    non-overlapping bands within the human visual response range,
    with spectrally-resolved detection. Human visual spectral
    sensitivity overlap with any "invisible" band strategy, and
    scene-reflectance-induced spectral shifts, require a declared
    confound envelope committed to the protocol digest.
\end{enumerate}
The separation modality, invariance tolerance, calibration
procedure for the separation channel, and any residual-crosstalk,
bleed, confound, or geometric-restriction envelope are committed
to the protocol digest.

\paragraph{Stereo / per-eye goggles as deployment pattern
(non-limiting).}
In some embodiments, per-eye stereo viewing for viewers equipped
with goggle hardware is achieved as a deployment pattern over the
temporal-interleave, polarization, or chromatic modalities above,
rather than as a standalone separation modality. The goggle filter
specification (polariser orientation, shutter timing, spectral
passband) is committed to the protocol digest alongside the
underlying $\mathsf{Sep}$ modality.

\paragraph{Stable multiplexed-loop acquisition
(non-limiting).}
In some embodiments, paired acquisition is iterated over a
committed cycle upper bound $N$ under a committed per-cycle
evaluation meter, $N$ and the meter committed to the protocol
digest. Each cycle reacquires independent probe-channel ground
truth through the separation channel $\mathsf{Sep}$; the
compounding-drift mechanism of the unpaired iterated method
(Section~\ref{sec:style-corpus-iterated-reset}) does not apply,
because the separation-invariant property supplies an
independent probe anchor at each cycle rather than requiring
each cycle to learn from the prior cycle's own drifted output.
Empirical result quality per cycle is evaluated under the
committed per-cycle meter and is not asserted as monotone or
unconditionally convergent.

\paragraph{Commitment and provenance.}
In some embodiments, $\mathcal{D}_{\mathrm{pair}}$ is committed as
a Merkle-rooted set of paired disclosure atoms, with each paired
atom binding a matched
$(\mathbf{y}^{\mathrm{probe}}_t,\mathbf{y}^{\mathrm{style}}_t)$
observation to a common scene instant and to the run-level
$\chi_t$ at acquisition time. The protocol digest provenance-types
$\mathcal{D}_{\mathrm{pair}}$ as committed physical evidence,
subject to the declared separation-invariance tolerance.

% ---------------------------------------------------------------------
\subsubsection{Slow-teacher / fast-student capability transfer via
physically-captured evidence (non-limiting)}
\label{sec:style-corpus-distilled}
% ---------------------------------------------------------------------

In some embodiments, an expensive styling component $\mathsf{Teacher}$
(whose per-frame inference may require minutes or longer under
declared compute constraints) generates styled emissions offline;
those emissions are projected onto a physical scene held stable
under a declared scene-stability envelope $\mathsf{SceneStab}$
(for example a mannequin, static scene, or slow-moving subject)
and captured through the substrate, producing committed
(emission, captured-physical-response) atoms. A fast styling
component $\mathsf{Student}$, deployable at runtime rates, is then
trained to reproduce the captured physical response from the
input seed, rather than to reproduce the teacher's raw output
tensor. The student therefore inherits physical-channel
transfer-function effects as part of its training signal. For
avoidance of doubt, $\mathsf{Teacher}$ denotes the slow styling
component of this subsubsection and is unrelated to the
teacher-forcing policy referenced in the pose-dynamics-model
training subsection of Section~\ref{sec:pose-rt}.

\paragraph{Scene-stability envelope (non-limiting).}
In some embodiments, $\mathsf{SceneStab}$ is a declared
admissibility bound on scene drift, subject motion, lighting
variation, and inter-frame differences across the teacher's
per-frame inference duration. The envelope may be declared
quantitatively (for example, bounded subject displacement and
bounded inter-frame pixel-difference statistics) or via a
committed subjective stability assessment procedure. The
envelope, the declared teacher inference-cost class, and the
acquisition procedure are committed to the protocol digest.

\paragraph{Student deployment-independence (non-limiting).}
In preferred embodiments, $\mathsf{Student}$ is fully independent
of $\mathsf{Teacher}$ at deployment: once trained, the student
does not invoke the teacher at runtime. In other embodiments,
hybrid deployment may invoke the teacher as a fallback when a
committed student-confidence statistic falls below a committed
threshold, with the fallback policy committed to the protocol
digest.

\paragraph{Committed teacher-student chain.}
In some embodiments, the teacher identity, version hash,
training-seed content, per-frame captured atoms, student identity,
training schedule, and held-out validation results are committed
as a chained record under the run-level $\chi_t$, so that a
verifier can walk from the committed teacher output, through the
committed projection-and-capture atom, to the committed student
training step, and check consistency.

\paragraph{Composition with pose-conditioned acquisition
(non-limiting).}
In some embodiments, when the stable scene is a slow-moving
subject, the acquisition composes with the pose-conditioned RT
embodiments of Section~\ref{sec:pose-rt}: pose-feature traces
$p_t$ extracted during acquisition are committed alongside the
teacher-student chain, and pose-distribution and pose-dynamics
terms may be evaluated on the resulting student training corpus.

\paragraph{Scope note (non-limiting).}
For avoidance of doubt, $\mathsf{Teacher}$ and $\mathsf{Student}$
may be drawn from conventional styling-component architectures
(including but not limited to conditional GANs, diffusion models,
flow-matching generators, or transformer-based generators). The
disclosed technical contribution is the committed teacher-student
chain through physically-captured evidence, the declared
scene-stability envelope, and the resulting physical-provenance
training signal, rather than any particular generator architecture.

\paragraph{Knowledge-distillation and model-compression methods (non-limiting).}
In some embodiments, the teacher-student chain is implemented using
knowledge-distillation methods, including without limitation logit
distillation (Hinton, Vinyals, Dean 2015), feature distillation,
intermediate-representation distillation (FitNets), attention
distillation, relational knowledge distillation, contrastive
representation distillation (CRD), self-distillation, online
distillation, and ensemble distillation. In some embodiments, the
student is additionally trained or post-processed under
model-compression methods including quantization-aware training
(QAT), post-training quantisation (PTQ), magnitude pruning,
structured pruning, lottery-ticket selection, weight clustering,
low-bit-width weight quantisation (including without limitation
8-bit, 4-bit, and 1-bit and 2-bit weight quantisation methods),
mixed-precision training, and sparse-attention pruning. These
methods are non-limiting; the apparatus is compatible with future
distillation and compression methods applied to the committed
teacher-student chain.

% ---------------------------------------------------------------------
\subsubsection{Forward counterfactual generator $\mathsf{G}_{\mathrm{alt}}$
(non-limiting)}
\label{sec:gen-forward-alt}
% ---------------------------------------------------------------------

In some embodiments, a scene-conditioned generator
$\mathsf{G}_{\mathrm{alt}}$ is trained on
$\mathcal{D}_{\mathrm{pair}}$ to predict the styled response that
would have occurred at a given scene instant under a candidate
styled emission $u^{\mathrm{style}}_{\mathrm{cand}}$, given the
committed paired neutral probe
$(u^{\mathrm{probe}},\mathbf{y}^{\mathrm{probe}}_t)$ of that
instant:
\[
  \widehat{\mathbf{y}}^{\mathrm{style}}_t
  =
  \mathsf{G}_{\mathrm{alt}}\!\bigl(
    u^{\mathrm{probe}},\,\mathbf{y}^{\mathrm{probe}}_t,\,
    u^{\mathrm{style}}_{\mathrm{cand}},\,\Pi\bigr).
\]
The generator is scene-conditioned: the committed probe capture
acts as a physically-attested scene embedding, distinguishing
$\mathsf{G}_{\mathrm{alt}}$ from unconditional style generators
that marginalise over scenes.

\paragraph{Training-corpus conditioning modes (non-limiting).}
In some embodiments, $\mathsf{G}_{\mathrm{alt}}$ conditions on the
committed probe capture directly; in other embodiments, it
conditions on a declared feature representation derived via
$\mathsf{RTDec}$, $\mathsf{Pose}$, or another committed feature
extractor applied to the probe capture. The conditioning mode,
feature extractor identity, and feature dimensionality are
committed to the protocol digest.

\paragraph{Use as scene-conditioned twin.}
In some embodiments, $\mathsf{G}_{\mathrm{alt}}$ serves as a
scene-conditioned digital twin distinct from the unconditional
$\mathsf{Twin}$ of Section~\ref{sec:twin-training}: its prediction
is over styled responses rather than arbitrary physical responses,
and its conditioning includes committed scene evidence. The output
$\widehat{\mathbf{y}}^{\mathrm{style}}_t$ is provenance-typed as
synthetic under the same discipline applied to
$C^{\mathsf{Twin}}_{0:T}$ in Section~\ref{sec:twin-training}.

\paragraph{Use as verifier-side forward audit.}
In some embodiments, a verifier holding committed
$(u^{\mathrm{probe}},\mathbf{y}^{\mathrm{probe}}_t,
u^{\mathrm{style}}_t,\mathbf{y}^{\mathrm{style}}_t)$ feeds the
first three to $\mathsf{G}_{\mathrm{alt}}$ and checks whether the
predicted
$\widehat{\mathbf{y}}^{\mathrm{style}}_t$ matches the captured
$\mathbf{y}^{\mathrm{style}}_t$ within a declared tolerance. Mismatch
is evidence of scene change between probe and style captures
(violating the separation-invariance property) or of tampering.

\paragraph{Use for training-data amplification.}
In some embodiments, $\mathsf{G}_{\mathrm{alt}}$ is applied to
unpaired probe captures together with a library of candidate
styled emissions to synthesise plausible styled counterfactuals,
which are provenance-typed as $\mathsf{G}_{\mathrm{alt}}$-generated
and used for training downstream components.

\paragraph{Scope note (non-limiting).}
$\mathsf{G}_{\mathrm{alt}}$ may be drawn from conventional
conditional-generation architectures; the disclosed technical
contribution is the training-data acquisition via
separation-paired capture, the scene-conditioning via committed
physically-attested probe evidence, and the provenance-typing of
outputs, rather than any particular generator architecture.

% ---------------------------------------------------------------------
\subsubsection{Inverse destyling generator $\mathsf{G}_{\mathrm{destyle}}$
(non-limiting)}
\label{sec:gen-inverse-destyle}
% ---------------------------------------------------------------------

In some embodiments, a scene-conditioned generator
$\mathsf{G}_{\mathrm{destyle}}$ is trained on
$\mathcal{D}_{\mathrm{pair}}$ in the inverse direction: given a
committed styled capture
$(u^{\mathrm{style}},\mathbf{y}^{\mathrm{style}}_t)$, it predicts
the response that the probe channel would have captured of the
same scene at the same instant:
\[
  \widehat{\mathbf{y}}^{\mathrm{probe}}_t
  =
  \mathsf{G}_{\mathrm{destyle}}\!\bigl(
    u^{\mathrm{style}},\,\mathbf{y}^{\mathrm{style}}_t,\,
    u^{\mathrm{probe}}_{\mathrm{query}},\,\Pi\bigr),
\]
where $u^{\mathrm{probe}}_{\mathrm{query}}$ is a declared
reference probe emission (typically a canonical neutral probe).
The output is the \emph{synthetic neutral view} of the scene
under active styling.

\paragraph{Use for perception-module input enhancement.}
In some embodiments, a downstream perception module (for example
$\mathsf{Pose}$ extractor, verisimilitude discriminator, or
kinematic-plausibility classifier) that would ordinarily operate
on raw styled captures instead operates on
$\mathsf{G}_{\mathrm{destyle}}$-produced synthetic neutral views,
recovering near-unstyled perception quality under active RT
projection.

\paragraph{Use for continuous baseline estimation.}
In some embodiments, $\mathsf{G}_{\mathrm{destyle}}$ provides a
continuous estimate of the baseline scene appearance
$\mathsf{Baseline}_t$ used in the positive-residual additive-overlay
objective of the pose-RT companion addendum, replacing the
interleaved baseline-block acquisition of that addendum between
scheduled baseline blocks. The baseline-continuous-estimation
policy is committed to the protocol digest.

\paragraph{Use for synthetic-neutral-view secure recording.}
In some embodiments, a recording captured under active RT
projection is accompanied by a synthetic neutral-view version
produced by $\mathsf{G}_{\mathrm{destyle}}$, and the synthetic
version is committed with a declared provenance-type identifying
it as $\mathsf{G}_{\mathrm{destyle}}$-generated rather than
directly captured. The synthetic neutral view may be disclosed
for audit, training, or sharing as a provenance-limited generator
estimate of underlying neutral scene content conditioned on
committed styled-capture evidence, while any committed styled
atoms, if separately retained or selectively opened, remain the
physical-provenance record of the styled run.

\paragraph{Downstream evidentiary scope of synthetic neutral views
(non-limiting).}
For verifier-facing purposes, a $\mathsf{G}_{\mathrm{destyle}}$-
generated synthetic neutral view is evidence only of the committed
styled-capture input, the declared $\mathsf{G}_{\mathrm{destyle}}$
identity and version, and the declared inference procedure; it
shall not be treated as direct physical-provenance evidence of
the scene's neutral appearance, nor as evidence of scene facts
not independently supported by selectively opened styled atoms
or other committed physical evidence. Opening committed styled
atoms under a declared audit authority does not change the
synthetic neutral view's status: the styled atoms remain physical-
provenance evidence of the run, while the synthetic neutral view
remains generator-produced output conditioned on them, merely
corroborated, contradicted, or qualified by comparison against
the opened styled atoms.

\paragraph{Redaction-at-source and selective-disclosure variants
(non-limiting).}
In some embodiments, $\mathsf{G}_{\mathrm{destyle}}$ is applied at
recording time so that the device commits the styled-capture
bundle (for RT operation) and the synthetic neutral view (for
downstream disclosure) but not a separately-redactable styled
original; the configuration is style-private by design. In other
embodiments, the styled capture atoms and the synthetic neutral
view are separately selectively-openable under distinct audit
authorities, so that a verifier with one authority sees only the
synthetic neutral view while a verifier with a higher authority
may additionally open the styled atoms.

\paragraph{Use for governance audit under partial disclosure.}
In some embodiments, a governance-side verifier given only
committed styled-capture evidence recovers an estimated neutral
scene appearance via $\mathsf{G}_{\mathrm{destyle}}$ and runs
downstream integrity checks on the estimated neutral view, under
declared partial-disclosure rules committed to the protocol
digest.

\paragraph{Use for training-data amplification, inverse direction.}
In some embodiments, unpaired styled captures are destyled to
synthesise plausible probe-captures, which are provenance-typed
as $\mathsf{G}_{\mathrm{destyle}}$-generated and used for training
components that expect neutral-scene input.

\paragraph{Scope note (non-limiting).}
$\mathsf{G}_{\mathrm{destyle}}$ may be drawn from conventional
conditional-generation architectures; the disclosed technical
contribution is the separation-paired training-data acquisition,
the scene-conditioning via committed styled evidence, the
provenance-typing of outputs, and the declared operational uses
(perception enhancement, continuous baseline estimation,
synthetic-neutral-view secure recording, governance audit under
partial disclosure), rather than any particular generator
architecture.

% ---------------------------------------------------------------------
\subsubsection{Dual-direction consistency meter (non-limiting)}
\label{sec:gen-dual-consistency}
% ---------------------------------------------------------------------

In some embodiments, both $\mathsf{G}_{\mathrm{alt}}$ and
$\mathsf{G}_{\mathrm{destyle}}$ are trained on the same
$\mathcal{D}_{\mathrm{pair}}$ and deployed together as a
\emph{dual-direction consistency meter}. For a committed frame
$(u^{\mathrm{probe}},\mathbf{y}^{\mathrm{probe}}_t,
u^{\mathrm{style}},\mathbf{y}^{\mathrm{style}}_t)$, the meter
produces three consistency statistics:
\begin{enumerate}
  \item $\delta_{\mathrm{fwd}}
    = \|\mathsf{G}_{\mathrm{alt}}(u^{\mathrm{probe}},
       \mathbf{y}^{\mathrm{probe}}_t, u^{\mathrm{style}},\Pi)
       - \mathbf{y}^{\mathrm{style}}_t\|$
    (forward consistency);
  \item $\delta_{\mathrm{inv}}
    = \|\mathsf{G}_{\mathrm{destyle}}(u^{\mathrm{style}},
       \mathbf{y}^{\mathrm{style}}_t, u^{\mathrm{probe}},\Pi)
       - \mathbf{y}^{\mathrm{probe}}_t\|$
    (inverse consistency); and
  \item $\delta_{\mathrm{sym}}$, a declared pair-symmetry statistic
    evaluating whether the forward and inverse predictions form a
    consistent round-trip when composed.
\end{enumerate}
Tampering, scene change between probe and style captures, or
separation-channel degradation tends to break one, two, or all
three statistics; the meter partition records which
statistic(s) failed, providing a per-frame physical-consistency
signal strictly stronger than forward or inverse consistency alone.
The meter's statistics, thresholds, and partition rule are
committed to the protocol digest.

% ---------------------------------------------------------------------
\subsubsection{Composed styled-training-data configurations
(non-limiting)}
\label{sec:styled-training-data-composed}
% ---------------------------------------------------------------------

In some embodiments, two or more of the mechanisms of
Sections~\ref{sec:style-corpus-curated}--\ref{sec:gen-dual-consistency}
are composed in a single deployment. Non-limiting combined
configurations include:
\begin{enumerate}
  \item paired acquisition (Section~\ref{sec:style-corpus-paired})
    feeding both $\mathsf{G}_{\mathrm{alt}}$ and
    $\mathsf{G}_{\mathrm{destyle}}$ training, with the dual-direction
    consistency meter of
    Section~\ref{sec:gen-dual-consistency} active at runtime;
  \item paired acquisition feeding slow-teacher / fast-student
    distillation (Section~\ref{sec:style-corpus-distilled}), with
    the teacher's styled output serving as the style channel of
    each paired capture;
  \item unpaired iterated acquisition with committed reset
    (Section~\ref{sec:style-corpus-iterated-reset}) producing
    initial training data, which is subsequently augmented by
    paired-acquisition cycles for stability once a seed styling
    component has been established;
  \item curated styled corpus
    (Section~\ref{sec:style-corpus-curated}) serving as declared
    reference for style-distribution loss terms while one or more
    of the above mechanisms supply deployment-matched training
    evidence.
\end{enumerate}
The active combination, mixing policy, and per-component weightings
are committed to the protocol digest.




% =====================================================================
% Virtual views and virtual emissions
% =====================================================================

\subsection{Virtual Views and Virtual Emissions (non-limiting)}
\label{sec:virtual-views-emissions}

In some embodiments, the verisimilitude and twin machinery of
Sections~\ref{sec:verisimilitude-training}, \ref{sec:twin-training},
and \ref{sec:gen-forward-alt} is extended to predicted observations
under detector and emitter configurations not physically instantiated
during the run. Two dual primitives are disclosed: a \emph{virtual
camera}, being a predicted observation under a detector pose,
spectral band, or aperture configuration not physically present; and
a \emph{virtual emitter}, being a predicted observation under an
emission $u^{\mathrm{cand}}$ that was not physically applied at the
declared time index. Each primitive admits three operating modes:
predictive use, loss-function use, and adversarial or competitive
use. The forward counterfactual generator
$\mathsf{G}_{\mathrm{alt}}$ of Section~\ref{sec:gen-forward-alt} and
the external and virtual-view checks paragraph of
Section~\ref{sec:verisimilitude-training} are special cases of this
unified treatment.

\subsubsection{Definitions (non-limiting)}
\label{sec:virtual-views-defs}

\begin{definition}[Virtual camera]
\label{def:virtual-camera}
A \emph{virtual camera} for a configuration $\theta$ is a measurable
map
\[
  \widehat{\mathbf{y}}^{(\mathrm{cam})}_t
  =
  \mathsf{G}_{\mathrm{vcam}}\!\bigl(
    C_{0:t},\,
    \xi^{\mathrm{cam}},\,
    u(t),\,
    \Pi
  \bigr)
\]
from a committed \cba history $C_{0:t}$, a declared
virtual-camera configuration $\xi^{\mathrm{cam}}$ (specifying pose,
intrinsics, spectral band, gating, polarisation, or sub-aperture set),
the emission $u(t)$, and a protocol digest $\Pi$, to a predicted
observation that an additional camera with configuration
$\xi^{\mathrm{cam}}$ would have produced at time index $t$.
$\mathsf{G}_{\mathrm{vcam}}$ may be realised by a digital twin
(Section~\ref{sec:twin-training}), an analogue twin
(Section~\ref{sec:analogue-twin}), a learned 3D representation
including but not limited to NeRF-like radiance fields, point-based
representations, Gaussian primitives, mesh-based representations,
or hybrid implicit--explicit forms, or by a scene-conditioned
generator analogous to $\mathsf{G}_{\mathrm{alt}}$ extended in the
detector axis. The configuration $\xi^{\mathrm{cam}}$, model
identity, version hash, and prediction tolerance are committed to
the protocol digest.
\end{definition}

\begin{definition}[Virtual emitter]
\label{def:virtual-emitter}
A \emph{virtual emitter} for a configuration $\theta$ is a measurable
map
\[
  \widehat{\mathbf{y}}^{(\mathrm{emit})}_t
  =
  \mathsf{G}_{\mathrm{vemit}}\!\bigl(
    C_{0:t},\,
    u^{\mathrm{cand}}(t),\,
    \xi^{\mathrm{emit}},\,
    \Pi
  \bigr)
\]
from a committed \cba history, a candidate emission
$u^{\mathrm{cand}}(t)$ that was not physically applied at time
index $t$, a declared virtual-emitter configuration
$\xi^{\mathrm{emit}}$ (specifying source pose, beam profile,
spectral band, polarisation, modulation schedule, or aperture), and
the protocol digest $\Pi$, to a predicted observation on one or
more physical detectors that the additional emission would have
produced. $\mathsf{G}_{\mathrm{vemit}}$ generalises the forward
counterfactual generator $\mathsf{G}_{\mathrm{alt}}$ of
Section~\ref{sec:gen-forward-alt} from substitution of a single
emission to addition of one or more candidate emissions on top of
$K_{\mathrm{co}}\geq 1$ committed concurrent emission streams under
the co-emission architecture of Section~\ref{sec:co-emission-base}.
The configuration $\xi^{\mathrm{emit}}$, candidate emission
specification, model identity, version hash, and prediction
tolerance are committed to the protocol digest.
\end{definition}

\paragraph{Provenance typing.}
Outputs of $\mathsf{G}_{\mathrm{vcam}}$ and $\mathsf{G}_{\mathrm{vemit}}$
are provenance-typed as synthetic under the same discipline applied
to digital-twin outputs in Section~\ref{sec:twin-training} and
to $\mathsf{G}_{\mathrm{alt}}$-generated outputs in
Section~\ref{sec:gen-forward-alt}, and shall not be treated as
direct physical-provenance evidence of scene facts not independently
supported by physically committed atoms.

\subsubsection{Predictive use (non-limiting)}
\label{sec:virtual-views-prediction}

\paragraph{Virtual-camera prediction.}
In some embodiments, given a learned 3D representation, digital
twin, or analogue twin, $\mathsf{G}_{\mathrm{vcam}}$ predicts the
observation that an additional camera at $\xi^{\mathrm{cam}}$ would
record under the committed emission schedule, generalising the
external and virtual-view checks paragraph of
Section~\ref{sec:verisimilitude-training}.

\paragraph{Virtual-emitter prediction.}
In some embodiments, $\mathsf{G}_{\mathrm{vemit}}$ predicts the
observation that one or more physical detectors would record under
an additional candidate emission $u^{\mathrm{cand}}$ at
$\xi^{\mathrm{emit}}$, generalising
$\mathsf{G}_{\mathrm{alt}}$ from substitutive to additive operation.

\paragraph{Linearity-regime decomposition (non-limiting).}
\label{sec:vemit-linearity}
In some embodiments, the detection pathway operates in an additive,
calibrated, unsaturated regime in which the response to a
multi-emission configuration $\{u^{(k)}\}_{k=1}^{K_{\mathrm{co}}}$
decomposes by superposition into per-stream contributions
$\{\mathbf{y}^{(k)}_t\}$ recoverable under the declared separation
modality $\mathsf{Sep}$ of Section~\ref{sec:co-emission-base}. In
such embodiments, $\mathsf{G}_{\mathrm{vemit}}$ is realised by
predicting the per-stream contribution
$\widehat{\mathbf{y}}^{(\mathrm{cand})}_t$ of a candidate additional
emission and forming the predicted aggregate observation by linear
superposition over the recovered physical contributions and the
predicted candidate contribution. The linearity envelope, including
declared upper bounds on saturation, threshold-crossing, dead-time,
and crosstalk regimes within which superposition is asserted, is
committed to the protocol digest, and a meter monitors
linearity-envelope occupancy throughout the prediction window.

\paragraph{Non-additive regime (non-limiting).}
In some embodiments operating outside the additive regime of
Section~\ref{sec:vemit-linearity}, $\mathsf{G}_{\mathrm{vemit}}$ is
realised as a learned scene-conditioned generator trained on
co-emission paired captures in which the same scene is recorded
under varying multi-emission configurations, and the model learns
the non-linear forward map directly. Training-corpus identity, the
held-out emission configurations used for validation, and the
held-out scene classes are committed to the protocol digest.

\paragraph{Combined virtual-view and virtual-emission prediction
(non-limiting).}
In some embodiments, $\mathsf{G}_{\mathrm{vcam}}$ and
$\mathsf{G}_{\mathrm{vemit}}$ are composed to predict the
observation of a virtual camera at $\xi^{\mathrm{cam}}$ under a
virtual emitter at $\xi^{\mathrm{emit}}$, neither of which was
physically present during the run, with the composition rule and
joint prediction tolerance committed to the protocol digest. As
prediction error compounds under composition of two synthetic
primitives, the joint prediction tolerance committed to the protocol
digest accounts explicitly for the compounded uncertainty.

\paragraph{Light-field hybrid speciality (non-limiting).}
In some embodiments deploying the light-field hybrid embodiment of
Section~\ref{sec:lightfield-hybrid}, virtual cameras correspond to
sub-aperture views $(u^{*},v^{*})$ not physically read out and
virtual emitters correspond to angular-sub-aperture emission
configurations $E(x,y,u^{*},v^{*})$ not physically commanded;
predicted virtual observations are evaluated against measured
sub-apertures or measured emission configurations under the
declared microlens-array pitch, calibration state, and
angular-resolution envelope.

\subsubsection{Loss-function use (non-limiting)}
\label{sec:virtual-views-loss}

\paragraph{Held-out virtual-camera supervision.}
In some embodiments, an agent, controller, decoder, or world model
is trained on input derived from a first physically captured camera
$\mathrm{cam}_1$ and supervised by a loss
\[
  L^{\mathrm{vcam}}_{t}
  \;=\;
  d\!\bigl(
    \widehat{\mathbf{y}}^{(\mathrm{cam}_2)}_t,\;
    \mathbf{y}^{(\mathrm{cam}_2)}_t
  \bigr)
\]
on the predicted virtual-camera output
$\widehat{\mathbf{y}}^{(\mathrm{cam}_2)}_t$ versus a measurement
$\mathbf{y}^{(\mathrm{cam}_2)}_t$ from a second physically captured
camera $\mathrm{cam}_2$ held out from the agent's input, where
$d(\cdot,\cdot)$ is a declared distance functional with declared
scale committed to the protocol digest. In some embodiments, the
second camera is operated in a governance-only role, with its
observations excluded from the agent's memory-formation and
action-selection pathways and routed exclusively to training
supervision and audit logging.

\paragraph{Held-out virtual-emitter supervision.}
In some embodiments, the dual scheme applies on the emitter axis: a
controller is trained to act under a first committed emission
stream and supervised by a loss
\[
  L^{\mathrm{vemit}}_{t}
  \;=\;
  d\!\bigl(
    \widehat{\mathbf{y}}^{(\mathrm{emit}_2)}_t,\;
    \mathbf{y}^{(\mathrm{emit}_2)}_t
  \bigr)
\]
on the predicted virtual-emitter response
$\widehat{\mathbf{y}}^{(\mathrm{emit}_2)}_t$ versus the physically
measured response $\mathbf{y}^{(\mathrm{emit}_2)}_t$ to a second
committed emission stream held out from the controller's reactive
loop. In some embodiments, the emission held out for supervision is
selected from a held-out emission family rotated periodically and
not used during prior training rounds.

\paragraph{Asymmetric controller training.}
In some embodiments, the loss-function uses of this section compose
with the \RT training against the twin of
Section~\ref{sec:rt-against-twin}: the controller's input is a
twin-produced or virtual-camera-produced observation, and the
controller's training loss is evaluated against a held-out physical
camera or held-out physical emitter. Gradients flow through the
twin or through $\mathsf{G}_{\mathrm{vcam}}$ or
$\mathsf{G}_{\mathrm{vemit}}$ to the controller, while the
supervisory signal remains physically grounded.

\paragraph{Cross-virtual-real consistency loss.}
In some embodiments, a consistency loss
\[
  L^{\mathrm{cons}}_t
  \;=\;
  d\!\bigl(
    \mathsf{G}_{\mathrm{vcam}}(C_{0:t}, \xi^{\mathrm{cam}}, u(t), \Pi),\;
    \mathbf{y}^{(\mathrm{cam})}_t
  \bigr)
  +
  d\!\bigl(
    \mathsf{G}_{\mathrm{vemit}}(C_{0:t}, u(t), \xi^{\mathrm{emit}}, \Pi),\;
    \mathbf{y}^{(\mathrm{emit})}_t
  \bigr)
\]
is applied to virtual-view and virtual-emission predictions when
the corresponding physical configurations are subsequently realised
on the device, providing an audit term that physically closes the
loop on the virtual primitives.

\paragraph{Multi-pair extension.}
In some embodiments, the loss-function uses of this section
generalise the multi-pair coordinated training paragraph of
Section~\ref{sec:verisimilitude-training} from cross-pair
consistency between physical pairs to cross-pair consistency
between physical pairs and virtual pairs, with the virtual pair
configurations committed to the protocol digest.

\subsubsection{Adversarial and competitive use (non-limiting)}
\label{sec:virtual-views-adversarial}

\paragraph{Virtual-view verisimilitude discriminator.}
In some embodiments, a verisimilitude functional is extended to
$V^{\mathrm{vcam}}_\theta$, scoring whether windows of predicted
virtual-camera observations are consistent with the empirical
distribution of physical camera observations of the same scene at
the same configuration. Negative training examples include
synthetic outputs drawn from competing 3D representations,
view-synthesis baselines, and twin variants of differing fidelity,
in the manner of the adversarial-autoencoder pattern of
Section~\ref{sec:verisimilitude-training}.

\paragraph{Virtual-emitter verisimilitude discriminator.}
In some embodiments, the dual functional $V^{\mathrm{vemit}}_\theta$
scores whether predicted virtual-emitter responses are consistent
with the empirical distribution of physical responses to the
candidate emission, with negative examples drawn from substitutive
$\mathsf{G}_{\mathrm{alt}}$ variants, additive-superposition
predictions outside their declared linearity envelope, and twin
predictions of differing fidelity.

\paragraph{Per-axis generator-verifier games on virtual primitives.}
In some embodiments, the minimax structure of the GAN-like
adversarial training of Section~\ref{sec:verisimilitude-training} is
instantiated on each axis independently, recognising that the
camera and emitter discriminators are independent verifiers.
A camera-axis game:
\[
  \min_{\mathsf{G}_{\mathrm{vcam}}}
  \max_{V^{\mathrm{vcam}}}
  \;\mathbb{E}\!\bigl[\log V^{\mathrm{vcam}}(\mathrm{real})\bigr]
  + \mathbb{E}\!\bigl[\log(1-V^{\mathrm{vcam}}(\mathsf{G}_{\mathrm{vcam}}))\bigr],
\]
and an independent emitter-axis game:
\[
  \min_{\mathsf{G}_{\mathrm{vemit}}}
  \max_{V^{\mathrm{vemit}}}
  \;\mathbb{E}\!\bigl[\log V^{\mathrm{vemit}}(\mathrm{real})\bigr]
  + \mathbb{E}\!\bigl[\log(1-V^{\mathrm{vemit}}(\mathsf{G}_{\mathrm{vemit}}))\bigr].
\]
In further embodiments, the two games are trained jointly with a
declared weighting term between the two objectives committed to the
protocol digest. The decomposition into per-axis games is
non-limiting.

\paragraph{Re-emission and re-capture audit.}
In some embodiments, virtual primitives are physically validated by
the system subsequently emitting $u^{\mathrm{cand}}$ at
$\xi^{\mathrm{emit}}$ or capturing from $\xi^{\mathrm{cam}}$ on the
physical device, and comparing the previously committed predicted
observation to the freshly captured physical observation. The
declared comparison tolerance and the time interval between
prediction commitment and physical re-capture are recorded in the
protocol digest. In some embodiments, this re-capture is selected
adversarially: the system identifies, by gradient ascent on
$V^{\mathrm{vcam}}$ or $V^{\mathrm{vemit}}$ over admissible
$(\xi^{\mathrm{cam}}, \xi^{\mathrm{emit}}, u^{\mathrm{cand}})$
within meter envelopes, those virtual configurations on which the
generator is least confident, and physically realises those
configurations to maximise audit information per re-capture.

\paragraph{Adversarial robustness training (non-limiting).}
In some embodiments, the agent or controller is trained against
adversarial virtual primitives: $\mathsf{G}_{\mathrm{vcam}}$ and
$\mathsf{G}_{\mathrm{vemit}}$ are perturbed within declared
envelopes to produce worst-case predicted observations, and the
agent's training loss is evaluated under these perturbations,
yielding controllers that are robust to virtual-primitive
hallucination within the declared envelope.

\paragraph{Capacity-accounting scope clause (non-limiting).}
The training, loss, and adversarial uses of this subsection are
directed to physically grounded model improvement, controller
robustness, and verisimilitude auditing under the declared meter
and protocol-digest discipline. They do not constitute, and shall
not be construed as, capacity or scaling claims of the kinds
bounded by the capacity-accounting rule of the present
specification, and do not imply any increase in capacity quantities.

\paragraph{Virtual-interface generalisation (non-limiting).}
Virtual cameras and virtual emitters generalise to any declared
computational sampling, rendering, stimulation, counterfactual, or
projection interface over a physical or simulated apparatus-scene
model. Non-limiting examples include pinhole cameras, rolling-shutter
cameras, event cameras, spectral cameras, light-field cameras,
plenoptic cameras, depth cameras, LiDAR-like samplers, ultrasound
samplers, RF samplers, projector models, structured-light emitters,
laser scanners, holographic emitters, acoustic emitters, microwave
emitters, and hybrid virtual instruments. Scene models may include
NeRF-style radiance fields, occupancy fields, signed-distance fields,
Gaussian splats, mesh models, point clouds, light fields, plenoxels,
hash-grid models, learned latent scene representations, physics
simulators, differentiable renderers, and non-neural parametric or
non-parametric models.

\subsection{Scene-state-contrast training data (non-limiting)}
\label{sec:scene-state-contrast}

In some embodiments, training data for downstream components
(classifiers, segmenters, state-conditional generators, digital
twin conditioning, perception-monitor admissibility rules, and
governance-state gates) is generated not only by style-variable
contrast of the kind disclosed in
Section~\ref{sec:styled-training-data}, but also by
\emph{scene-state contrast}: two or more committed corpora
acquired through the \RK substrate under a matched
emission protocol or declared matching envelope, where the
corpora differ in a declared physical scene-state variable. The
corpora are committed with per-atom state labels and downstream
models are trained, validated, or audited on the state contrast.

Scene-state-contrast is anchored to the substrate: each committed
atom binds an actual emission to an actual physical observation
under a shared time base and protocol digest, with Merkle
commitment and hash-chain binding in the usual manner. The
novelty is not the concept of contrast-based training (which is
generic in supervised and contrastive ML) but the committed
physical-evidence acquisition discipline around scene-state
variables, the provenance-typing of state-conditional model
outputs, the composition with \RK machinery (twin,
analogue twin, perception monitor, co-emission, additive-overlay),
and the privacy discipline around sensitive state corpora.

The following subsubsections describe the primary sub-embodiment
(subject-presence contrast), five adjacent sub-embodiments that
generalise it, and two state-specific discipline subsubsections
(label-integrity and privacy-discipline) that apply across all
sub-embodiments.

% ---------------------------------------------------------------------
\subsubsection{Subject-presence contrast
(non-limiting, lead embodiment)}
\label{sec:subj-presence-contrast}
% ---------------------------------------------------------------------

In some embodiments, two committed corpora are acquired differing
in whether a subject is present in the scene:
\begin{itemize}
  \item An \emph{empty-scene corpus} $\mathcal{D}_{\mathrm{empty}}$
    is acquired by operating the \RK under a declared
    emission protocol on a declared scene or scene-class without
    a subject present; committed atoms record the scene's response
    to emission in the subject's absence.
  \item An \emph{occupied-scene corpus}
    $\mathcal{D}_{\mathrm{occupied}}$ is acquired by operating the
    \RK under the same or declared-matched emission
    protocol on the scene with a subject present; committed atoms
    record the combined scene-plus-subject response.
\end{itemize}
In preferred embodiments reflecting deployed practice, the two
corpora are acquired in \emph{separate sessions} rather than in a
matched entry/exit capture session, to preserve acquisition
cleanliness and prevent cross-contamination between state classes.
In alternative embodiments, matched same-scene acquisition uses
declared entry/exit-event markers and a committed
geometric-stability envelope, producing tightly-paired
per-frame-level state contrast at the cost of more demanding
scene management. The acquisition regime (separate-session or
matched-capture), emission-protocol identity, scene-class
declaration, entry/exit-event markers where applicable, and
per-atom occupancy labels are committed to the protocol digest.

\paragraph{Downstream uses (non-limiting).}
In some embodiments, the contrasted corpora feed one or more of:
\begin{itemize}
  \item a subject-presence binary classifier producing per-frame
    or per-window subject/no-subject decisions;
  \item a subject-segmentation model producing pixel-wise,
    region-wise, or feature-wise subject/background partitions;
  \item a subject-conditional generative model producing empty-
    scene reconstructions from occupied captures (subject removal),
    or the inverse (scene-plus-subject synthesis), in the manner
    of the inverse generator $\mathsf{G}_{\mathrm{destyle}}$ of
    Section~\ref{sec:gen-inverse-destyle} but conditioned on
    occupancy rather than style;
  \item a subject-presence feature for the digital twin of
    Section~\ref{sec:twin-training}, supplying occupancy as a
    declared conditioning variable so that twin predictions
    correctly handle occlusion, reflection, and subject-scene
    interaction patterns;
  \item validation of upstream components (pose extractor,
    verisimilitude discriminator, kinematic-plausibility
    classifier) via physically-grounded occupancy ground truth.
\end{itemize}

% ---------------------------------------------------------------------
\subsubsection{Non-human scene-state contrast (non-limiting)}
\label{sec:nonhuman-scene-state}
% ---------------------------------------------------------------------

In some embodiments, the contrasted scene-state variable does not
involve a human subject. Non-limiting variables include: object
presence/absence; object arrangement class; door/window state;
equipment position or operating mode; calibration-target state;
display, screen, or signage state; occluder state; surface cover
state; wet/dry state; contamination or clean state; lighting
configuration or preset; environmental conditions such as smoke,
dust, or aerosol density (claimed as declared binned or ranged
variants per the continuous-variable discipline below); and
safety-zone occupancy by equipment. Each non-human state variable
is declared by its identity, value domain, acquisition regime, and
matching envelope across corpora, all committed to the protocol
digest.

\paragraph{Continuous and multi-state variables.}
In some embodiments, the contrasted variable is continuous (for
example illuminance, spectrum, sun angle, time of day,
temperature, humidity, subject distance, object pose) rather than
discrete. Such variables are treated as declared binned, ranged,
or regression-conditioned variants with a committed calibration
envelope distinguishing them from the binary state-class case of
subject presence. In other embodiments, the contrasted variable
is multi-state but discrete (subject count; activity class;
equipment operating mode); multi-state variables require more than
two corpora under a matched protocol.

% ---------------------------------------------------------------------
\subsubsection{State-transition contrast (non-limiting)}
\label{sec:state-transition-corpus}
% ---------------------------------------------------------------------

In some embodiments, the contrasted corpora are acquired on the
same scene in temporal before/after order across a declared
transition event, producing a \emph{state-transition corpus}
structurally distinct from separate-session state-class corpora.
Non-limiting transitions include subject entry/exit, object
placement or removal, lighting or environmental change, equipment
movement, calibration update, cleaning, wear accumulation,
tamper event, and weather-exposure transitions. State-transition
corpora support transition prediction, change detection, tamper
detection, and causal-residual training that separate-session
corpora cannot support at the per-instant level.

The transition-event identity, pre-transition and
post-transition atom boundaries, transition-time commitment, and
any declared scene-stability envelope for the pre- and
post-transition intervals are committed to the protocol digest.
State-transition corpora may be composed with separate-session
corpora to provide both transition-grounded and class-level
training evidence.

% ---------------------------------------------------------------------
\subsubsection{Device-state contrast (non-limiting)}
\label{sec:device-state-contrast}
% ---------------------------------------------------------------------

In some embodiments, the contrasted corpora differ in a declared
\emph{device-state variable} rather than a scene-state variable:
same scene, same occupancy, but different projector calibration
state, focus state, detector gain/exposure, reactor operating
point, scan-law parameterisation, emitter aging state, aperture or
filter state, actuator drift state, safety-envelope state, or
analogue-twin matching state of
Section~\ref{sec:analogue-twin}. Device-state-contrast corpora
support training of device-state-conditioned twins, drift
detectors, tamper classifiers, and calibration-update policies.

The device-state variable identity, pre-contrast and post-contrast
device configurations, declared measurement procedure for the
state variable, and any matching envelope across contrasting
corpora are committed to the protocol digest. This sub-embodiment
composes with the trainable parameter-family discipline of the
\RK architecture: device-state variables can be drawn
from the committed parameter families.

% ---------------------------------------------------------------------
\subsubsection{Differential and residual training data (non-limiting)}
\label{sec:residual-diff-training}
% ---------------------------------------------------------------------

In some embodiments, training data is committed not as paired raw
atoms but as committed differential or residual atoms derived from
two contrasting states: for example per-pixel, per-feature, or
per-latent residuals
$\Delta\mathbf{y} = \mathbf{y}^{(A)} - \mathbf{y}^{(B)}$ computed
under a declared residual-extraction rule. Differential atoms
provide a compact training signal and are distinct from paired
raw-atom corpora in that downstream models train directly on the
residual rather than learning to extract it. Per-pixel or
per-frame residuals require matched same-scene or
matched-transition acquisition (for example the matched-capture
regime of Section~\ref{sec:subj-presence-contrast}); separate-
session class-level corpora support distributional residuals,
contrastive losses, or classifier contrasts, not exact
ground-truth residuals.

The residual-extraction rule, declared feature or latent space,
per-atom canonicalisation, and declared usage policy are
committed to the protocol digest. Residual atoms inherit
provenance-typing as derivatives of the two contrasting corpora
and are not independently physical-provenance evidence of either
state.

% ---------------------------------------------------------------------
\subsubsection{Multi-factor factorial contrast
(non-limiting, optional)}
\label{sec:multi-factor-contrast}
% ---------------------------------------------------------------------

In some embodiments, more than one declared state variable is
varied across corpora under a declared factorial or blocked
acquisition design (for example occupancy $\times$ lighting,
subject-identity $\times$ pose, object-arrangement $\times$
weather, device-calibration $\times$ scene-state). The active
design, factor levels, confound envelopes across uncontrolled
factors, and block or replication structure are committed to the
protocol digest. This sub-embodiment is optional and may be
deferred where claim bandwidth or acquisition cost does not
support full factorial coverage.

% ---------------------------------------------------------------------
\subsubsection{State-label integrity (non-limiting)}
\label{sec:state-label-integrity}
% ---------------------------------------------------------------------

Per the commitment discipline of the \RK, per-atom
state labels are themselves committed evidence. In some
embodiments, the label source is one or more of:
\begin{itemize}
  \item \textbf{Operator-declared labels}: the state is asserted
    by the operator at acquisition time, with a committed operator
    identifier and timestamp.
  \item \textbf{Sensor-confirmed labels}: an independent sensor
    (for example a separate occupancy detector, a calibration
    witness, or an environmental sensor) confirms the state,
    with the sensor identity, version, and confidence committed.
  \item \textbf{Externally-witnessed labels}: an independent
    authority witnesses the state declaration, with the witness
    identifier and witnessing protocol committed.
  \item \textbf{Bonded-identity-confirmed labels}: where subject
    identity is relevant, a committed bonded-identity record from
    the \RK's bonded human identity extension may
    confirm presence, consent, or other identity-linked state
    (subject to the identity-firewall discipline of
    Section~\ref{sec:state-privacy-discipline}).
  \item \textbf{Classifier-derived labels}: a prior committed
    classifier produces the state label, with the classifier
    identity, version, training-corpus digest, and confidence
    committed.
\end{itemize}
Label sources of different strength support different audit
claims. In preferred embodiments, the label source, authority,
and audit method are committed per atom or per commitment epoch,
and higher-stakes downstream uses (safety gating, privacy
restrictions, provenance attestations) require stronger label
sources or multiple independent confirmations.

\paragraph{Anti-poisoning and label-audit discipline
(non-limiting).}
In some embodiments, a declared label-audit procedure is
committed, under which a sample of committed atoms is randomly or
adversarially selected and the state label is independently
re-confirmed. Label-audit failure rates are committed as
meter-envelope statistics. A false "empty" label is both a
privacy failure and a safety failure, and the protocol digest
declares the consequences of label-audit failure (corpus
quarantine, retraining, authority downgrade, successor
protocol-digest version).

% ---------------------------------------------------------------------
\subsubsection{Privacy-discipline and three-object sensitivity
envelope (non-limiting)}
\label{sec:state-privacy-discipline}
% ---------------------------------------------------------------------

Subject-presence contrast and related scene-state-contrast
training data raise sensitivities beyond those addressed by the
generic provenance-typing discipline of the filing. Three
sensitive objects require independent handling:

\paragraph{$\mathcal{D}_{\mathrm{empty}}$ sensitivity (non-limiting).}
The empty-scene corpus is not necessarily innocuous. It may
reveal private scene layout, equipment placement, security
posture, absence-of-people patterns, or unobserved-scene
conditions that the scene owner has interest in keeping
private. In some embodiments, $\mathcal{D}_{\mathrm{empty}}$ is
committed under a declared sensitivity envelope distinct from
$\mathcal{D}_{\mathrm{occupied}}$, with separate audit-authority
levels, retention rules, selective-opening rules, and disclosure
granularity. The two corpora are separately Merkle-addressable
and separately redaction-capable, so that an audit of one does
not automatically expose the other.

\paragraph{$\mathcal{D}_{\mathrm{occupied}}$ sensitivity (non-limiting).}
The occupied-scene corpus aggregates personally identifying
information. In some embodiments, per-subject consent records,
bonded-identity commitments where applicable, and
subject-consent-based retention rules are committed alongside
the corpus, with disclosure restricted to authorities authorised
to inspect subject-bearing atoms.

\paragraph{Classifier-output sensitivity (non-limiting).}
Subject-presence classifier outputs reveal when a subject was
present and constitute a surveillance-grade inference in
aggregate. In some embodiments, classifier outputs are
provenance-typed as model inferences over a declared committed
evidence window, with classifier identity, version, training-corpus
digest, threshold/calibration version, confidence score, time
window, and disclosure authority committed. Classifier outputs
shall not be treated as direct physical-provenance evidence of
occupancy or non-occupancy unless supported by selectively
opened atoms or an evidence-bearing companion stream. This is
especially important for "no subject present" determinations,
which are safety- and privacy-sensitive.

\paragraph{Subject-removal reconstruction
provenance (non-limiting).}
In some embodiments, a generator produces a synthetic empty-scene
estimate from an occupied capture (subject removal), trained on
the contrast of $\mathcal{D}_{\mathrm{empty}}$ and
$\mathcal{D}_{\mathrm{occupied}}$ under the same discipline as
$\mathsf{G}_{\mathrm{destyle}}$ of
Section~\ref{sec:gen-inverse-destyle}. Such generator output is
evidence only of the committed occupied-capture input, the
declared generator identity and version, and the declared
inference procedure; it shall not be treated as direct physical-
provenance evidence that the empty scene actually had the
reconstructed appearance. Selective disclosure may permit the
synthetic empty estimate to be shown without opening occupied
atoms, while preserving an escalated audit path to the
underlying occupied evidence for authorised verifiers.

\paragraph{Crypto-companion mode-commitment proof
(non-limiting).}
In some embodiments, a co-emitted committed cryptographic
companion stream (Section~\ref{sec:co-emission-continuous-uses})
with external anchoring proves that the device committed to a
declared state-classifier version, training-corpus digest,
empty-only mode, occupied-data-disabled mode, or disclosure
policy during a capture window. The cryptographic companion
proves \emph{mode commitment}; it does not by itself prove the
factual absence of a subject during the capture window. Factual
absence requires classifier evidence, opened physical atoms, or
an evidence-bearing companion (neutral probe, \TB probe,
or \LI probe). The protocol digest declares the scope of
what the cryptographic commitment proves and the limits of that
proof.

\paragraph{Deployment-scope restrictions (non-limiting).}
In some embodiments, subject-presence classifier outputs are
used only within a declared deployment scope committed to the
protocol digest: safety gating (emission downgrade, hazardous-
capability disablement), data-collection disablement (retention
minimisation when subjects present), selective disclosure
(redact subject-bearing atoms from routine audits), or
policy-triggered fallback to semantic-fallback output. The
deployment-scope restriction is itself a committed claim of the
device: the classifier output is a governance input for the
declared restrictions and is not released as a standalone
surveillance feed outside the declared scope.

\paragraph{Identity firewall (non-limiting).}
Presence, subject-count, subject-segmentation, subject-identity,
and bonded-subject-identity are separate inferences with
separately escalating sensitivity. In some embodiments, the
protocol digest commits independent authority, consent-record,
selective-opening, and provenance-typing discipline for each
layer. A subject-presence inference does not authorise identity
or bonded-identity inference; identity compositions require
separately declared authority, separate bonding records, and
separate audit trails, and in preferred embodiments are placed
under the \RK's bonded human identity extension
rather than the scene-state-contrast subsection.


\section{Application Domains (Non-limiting)}
% ======================================================================

This section describes non-limiting application domains. Any use in
which a controllable emission and sensing system records the \cb
and uses the resulting evidence for verification, perception, and/or
controllable rendering is within scope.

\subsection{Generative media and immersive content}

In some embodiments, \RKs support generation, transformation,
and presentation of media grounded in physical interaction. Subsystems
include controllable rendering and projection mapping (\RT
head selecting emission patterns), hybrid diffusion/flow/consistency
models (physical stochasticity as forward corruption with learned
reverse process), and semantic-conditioned rendering (text prompts
modulating emission policies). In Yoked variants, the scene's response
to rendering informs the next rendering step, creating responsive
immersive content.

\paragraph{3D generative models and view consistency (non-limiting).}
In some embodiments, the system supports 3D reconstruction and
rendering, including NeRF-like radiance field models, point-based
models, and scene graphs. Multi-view constraints may be enforced by
combining channels with distinct invariance profiles and by training
decoders to produce view-consistent outputs under changing control
protocols. In some embodiments, a light-field hybrid embodiment
(Section~\ref{sec:lightfield-hybrid}) provides a concrete hardware
instantiation of sampled light-field evidence for radiance-field
model fitting and parallax-consistent rendering, in which a plenoptic
detector acquires multi-view evidence in a single exposure and a
light-field projector synthesises parallax-consistent output, with
the \cb committing the sampled four-dimensional input--output
correspondence as the evidentiary basis for view consistency.

\paragraph{Audio and multimodal generation (non-limiting).}
In some embodiments, the system generates or transforms audio, speech,
or music using reactor temporal dynamics and audio \cba channels,
optionally coupled to optical channels (for example cymatic media,
meaning media whose surface patterns respond to sound pressure,
observed by a camera). Text or semantic embeddings may condition
emission policies so that physically grounded signals drive audio
synthesis or transformation.

\paragraph{PoliePuter mode: reactor as compute substrate
(non-limiting).}
\label{par:PoliePuter}
In some embodiments consistent with the open-loop sub-case of
Definition~\ref{def:poliputer}, the scene loop is disabled
($S^{\mathrm{scene}} = S_{\mathrm{term}}$) and the reactor subsystem
is used as a physical compute substrate. A learned encoder maps an
input $x$ to a control protocol $U_{0:T}(x)$, the reactor produces
the \cb $C_{0:T} \sim \mathsf{P}_\theta(\cdot \mid S_{\mathrm{term}},
U_{0:T}(x))$, and a learned decoder maps $C_{0:T}$ to a reconstruction
$\hat{x}$. The coupled sub-case of Definition~\ref{def:poliputer}, in
which a continuous abstract or computational scene partner closes a
bidirectional loop with the reactor, is treated separately and is
governed by the Yoked-mode machinery of
Section~\ref{sec:yoked}; the present paragraph concerns only the
open-loop sub-case.

The mathematical structure of the attractor bundle, Fisher--Rao
connection, curvature, holonomy, and characteristic classes is
preserved in the open-loop sub-case to the extent that those
constructions do not depend on scene-coupling dynamics: regime-level
results that are stated for general reactor dynamics under
$S^{\mathrm{scene}} = S_{\mathrm{term}}$ apply to the open-loop
PoliePuter substrate, while results whose enabling derivation invokes
scene coupling (in particular Yoked-mode synchronisation diagnostics,
$\mathrm{CLE}$ and transfer-entropy conditions, the Yoked
authentication warrant of Section~\ref{sec:yoked}, and any
property whose proof requires a coupled scene partner) do not
automatically carry over and are not asserted for the open-loop
sub-case. The following remapping is used where it applies: scene
maps to input data (digital-to-analogue converted or natively
analogue); scene dynamics map to input signal statistics; \cb maps to
computed output; reactor memory maps to working memory or context
buffer; protocol digest maps to execution log; and holonomy maps to
computation audit trail.


\paragraph{Latency and reconfiguration depth (non-limiting).}
This disclosure distinguishes propagation depth from reconfiguration
depth. Propagation depth contributes latency on the scale of cumulative
optical path length. Reconfiguration depth contributes latency on the
scale of the slowest active element per stage. Claims of low-latency
operation herein refer to propagation depth and do not assert that
reconfigured routed ensembles execute in propagation-scale latency.

\paragraph{Multi-reactor architectures (non-limiting).}
\label{par:multi-reactor-topologies}
In some embodiments, $N$ reactors are coupled in a network. The
configuration space of the network is
$\Theta_{\mathrm{net}} = \Theta_1 \times \cdots \times \Theta_N
\times \mathcal{G}$, where $\mathcal{G}$ is the space of coupling
graphs specifying which reactor feeds which, with what strength and
delay. Setting a point
$(\theta_1, \ldots, \theta_N, G) \in \Theta_{\mathrm{net}}$ defines
the computation. Non-limiting topologies include:

\begin{description}[style=nextline, leftmargin=2em]
  \item[Chain / pipeline.] Reactors in sequence,
    $R_1 \to R_2 \to \cdots \to R_N$. Each reactor is one processing
    stage. The chain depth determines a temporal hierarchy: deeper layers
    integrate over longer timescales if each reactor has a distinct time
    constant. This implements deep reservoir computing with physical
    layers.
  \item[Parallel / SIMD.] Same input routed to $N$ reactors with
    different $\theta_i$, outputs aggregated. The collection of outputs
    is a nonlinear feature expansion. Effective reservoir dimension
    equals $\sum_i d_{\mathrm{eff},i}$ plus coupling-induced dimensions
    if aggregation is nonlinear. In some embodiments,
    wavelength-division multiplexing routes different spectral
    components to different reactors simultaneously; in some
    embodiments, time-division multiplexing injects input segments
    sequentially to the same reactor with memory providing context
    across segments.
  \item[Graph / mesh.] Arbitrary coupling topology where a coupling
    matrix $G_{ij}$ specifies strength and delay from $R_i$ to $R_j$.
    Any cycle in $G$ creates recurrence. This topology is isomorphic to
    the networked kernel framework
    (Section~\ref{sec:networks}) with data as input instead of
    scenes; all network dynamics apply.
  \item[Hierarchical / layered.] Graph with strict layer ordering and
    optional skip connections. Each layer aggregates information from the
    previous layer and passes to the next, creating a hierarchy of
    abstraction levels.
  \item[Reservoir-of-reservoirs.] Each reactor is itself a network of
    simpler elements. Two-level hierarchy: micro-dynamics within each
    reactor, macro-dynamics across the network. Coupling between levels
    can introduce additional independent state variables not present in
    either level in isolation; the effective computational capacity of
    the coupled system, measured as the number of accessible linearly
    independent state variables under the relevant fading-memory
    framework and readout regime, may correspondingly exceed a simple
    additive baseline computed from uncoupled isolated readouts for
    suitable readout and coupling regimes. No claim of super-additivity
    in all configurations is asserted; the actual capacity gain is
    regime-specific and is characterised empirically per deployment.
  \item[Recurrent.] Directed cycles in the coupling graph create
    recurrence at the network level. Unlike single-reactor recurrence
    (which arises from the feedback loop's own temporal dynamics),
    network-level recurrence routes information through physically
    distinct media with different nonlinearities, time constants, and
    noise characteristics. In some embodiments, recurrent reactor
    networks implement physical analogues of gated recurrent units:
    one reactor modulates the coupling strength of another, providing
    content-dependent gating of memory and information flow. The
    recurrence depth is limited by signal-to-noise degradation per
    cycle rather than by discrete layer count.
\end{description}

\paragraph{Multi-reactor primitives (non-limiting).}
In some embodiments, the following operations exist only with $N > 1$
coupled reactors and have no single-reactor analogue:

\emph{Route:} dynamically modulate coupling strengths $G_{ij}$ to
direct information flow through specific reactor sub-networks based on
input content, implementing a content-dependent switch fabric
analogous to mixture-of-experts routing. \emph{Compile:} find the
network configuration $(\theta_1, \ldots, \theta_N, G)$ that minimises
a target task loss, using natural gradient descent on the product Fisher
metric of the network's configuration space, where the coupling Fisher
information captures how task loss depends on topology. \emph{Parallelise:} partition the input (spatially, spectrally, or
temporally) and route partitions to separate reactors for simultaneous
processing, with throughput scaling with the number of physical
channels. \emph{Reduce:} combine outputs from multiple reactors via
voting (fault-tolerant), averaging (noise-reducing), or nonlinear
combination (where the aggregator is itself a reactor providing further
computation). \emph{Error-correct:} replicate computation across $K$
reactors with nominally identical but physically distinct parameters
(the reactor-microstructure unclonability property), with agreement indicating correctness and
disagreement indicating fault or attack; in some embodiments, the
attractor manifold itself serves as a codebook for error detection,
where off-manifold outputs indicate computational errors.
\emph{Abstract:} chain reactor stages with different time constants to
create representations at multiple abstraction levels, where early
stages capture fast local features and deep stages capture slow global
features.

\paragraph{Proof-of-physics (non-limiting).}
In some embodiments, when a computation is performed physically by a
specific device or device network, the output inherently carries three
properties: the device's microstructure-specific attractor signature
(the Sign primitive applied to computational output); a
thermodynamic dissipation cost arising from the physical processes of
the reactor (with lower bounds under standard physical assumptions; a
digital simulation may avoid these specific dissipation costs but
is not expected to reproduce the reactor's microstructure-dependent thermal
trace); and a holonomy encoding of the
computational path taken through $\Theta_{\mathrm{net}}$ (what was
computed, not merely the result). This triad provides a basis for
attestation that a specific physical computation was performed by a
specific physical device. Verification checks whether the output lies
on the claimed device's attractor manifold (\TB on the output)
and whether the output's trajectory holonomy is consistent with the
claimed computation's path through $\Theta_{\mathrm{net}}$.

Non-limiting applications include verifiable inference (attesting that a
machine-learning workload was processed by specific physical hardware),
regulatory compliance (attesting that a risk model was evaluated on
certified hardware), and physical provenance for computed artefacts
(the computation's result is empirically attested, under declared
attacker families and meter envelopes, as having been produced by a
specific physical device at a specific time).

\paragraph{Data types as scenes (non-limiting).}
When the input is data rather than a physical scene, different data
types map to different input modalities. Time series map naturally to
temporal reactor dynamics via direct analogue injection. Images are
scan-converted to temporal signals where the scan pattern (raster,
spiral, Hilbert curve, random) determines how spatial information maps
to reactor dynamics; different scan patterns activate different attractor
features. Graphs and networks are encoded as coupling topology: data
graph nodes map to reactors, data graph edges map to couplings, and
graph-level properties (centrality, clustering, community structure)
emerge as attractor features of the reactor network's dynamics.
Tabular data modulates different reactor parameters or input channels,
with the reactor's nonlinear mixing providing automatic feature
interaction. Text and sequences are tokenised and injected sequentially,
with reactor memory providing a physical context window whose length
equals memory persistence time divided by token injection rate.
Audio feeds directly as an analogue signal. Multi-modal inputs route
different modalities to different reactor ports or wavelengths, with
the reactor fusing modalities through nonlinear mixing.

\paragraph{Computational tasks with physical advantage (non-limiting).}
In some embodiments, physical computation offers structural advantages
beyond raw speed. Non-limiting examples include: Boltzmann sampling,
where the reactor's thermal noise is the sampling process itself and
the attractor's basin structure defines the distribution; optimisation
via entrainment, where a cost function is mapped to coupling strengths
and the reactor network settles to its ground-state attractor;
differential equation solving, where the reactor physically integrates
its own dynamics at propagation speed rather than numerical integration
speed; analogue-to-insights processing, where analogue sensor data is
processed in the analogue domain without digital conversion, with
latency equal to propagation time and potentially much lower power
consumption; and privacy-preserving computation, where data need not be
digitally represented and intermediate states may exist only as transient
physical dynamics that are thermodynamically erased after use, with
proof of deletion via the Erase primitive
(Section~\ref{sec:info-operations}).
Cryptographic primitives are a further non-limiting example:
bifurcation structure provides natural one-way functions (the forward
map $x \to C_{0:T}$ given $\theta$ is computable, but inversion without
$\theta$ requires disambiguating exponentially many branches), the
reactor's thermal noise provides a physical random number generator at
GHz rates, and device-specific attractor signatures provide key
material that is grounded in physics rather than computational
hardness assumptions.

\paragraph{Adjoint reactor pairs and physical gradient propagation
(non-limiting).}
In some embodiments, two reactor subsystems are configured as a
\emph{forward--adjoint pair}: reactor~$A$ implements a forward
operator $f_A$ (for example a linear optical transform such as a
matrix multiplication, a convolution, or a Fourier-domain filter),
and reactor~$B$ is configured to physically implement the adjoint
(transpose, conjugate transpose, or generalised inverse) $f_A^\dagger$
of that operator.  Error signals originating from a loss computed on
the output of reactor~$A$ are propagated backward through reactor~$B$,
producing physical gradient estimates without digital Jacobian
computation.  In these embodiments, reactor~$A$ may be linear or
weakly nonlinear---linearity is an advantage, not a limitation,
because it ensures the adjoint relationship $\langle f_A(x),\,
e\rangle = \langle x,\, f_A^\dagger(e)\rangle$ holds exactly for
linear embodiments, and holds to first order under local linearisation
for weakly nonlinear embodiments,
yielding unbiased (or approximately unbiased) gradient propagation.  Non-limiting physical
realisations of the adjoint include: phase-conjugate mirrors that
time-reverse an optical wavefront; transposed free-space optical
interconnects (reversing source and detector planes); reciprocal
fibre-optic couplers traversed in the reverse direction; and
spatial light modulators programmed with the transpose of the
forward modulation pattern.

In some embodiments, a chain of $L$ forward--adjoint pairs implements
layered physical backpropagation: the forward pass propagates data
$x_0 \to x_1 \to \cdots \to x_L$ through reactors $A_1, \ldots,
A_L$, a loss $\mathcal{L}(x_L)$ is computed (digitally or by a
comparator circuit), and the error signal $\nabla_{x_L}\mathcal{L}$
is propagated backward through adjoint reactors $B_L, \ldots, B_1$
to produce per-layer gradient estimates.  Parameter updates to each
forward reactor are then derived from the conjunction of its stored
forward activation and the backward-propagated error at that layer.
In some embodiments, the forward and adjoint reactors share a
physical medium traversed in opposite directions (exploiting optical
reciprocity), reducing the hardware to a single medium per layer.

In some embodiments, the adjoint reactor pair is combined with a
nonlinear reactor that provides attestation: the linear pair performs
the computation, and a separate nonlinear reactor signs the result
via \TB, so that the output carries both a computed result
and a provenance certificate.  The \cb logs both the
forward and adjoint passes, and protocol digests record the
adjoint-pair configuration so that the physical gradient computation
is auditable.  This embodiment is an instance of a broader pattern
in which linear and nonlinear reactors play complementary roles
within a single module or network: linear subsystems for
high-fidelity computation, nonlinear subsystems for authentication
and hardness.  Adjoint-reactor pairs are a non-limiting instance of
the multi-reactor architectures described above, using the Route,
Merge, and Fork coupling primitives to direct forward and adjoint
signal flows through the appropriate reactor subsystems.
In some embodiments, forward--adjoint reactor architectures are integrated
with attestation reactors so that forward computation, backward gradient
propagation, and provenance certification are co-logged within a unified
\cb. This arrangement allows physically executed computation
and physically grounded attestation to be combined within the same module
or network.

\paragraph{Training dimension count (non-limiting).}
In some embodiments, the total number of trainable axes per single
reactor is the number of information-operation primitives
(Section~\ref{sec:info-operations}) plus the number of multi-reactor
primitives defined above. For a network of $N$ reactors, the total
training dimensionality is $N \times P + |\mathcal{G}|$, where $P$ is
the number of per-reactor primitive axes and $|\mathcal{G}|$ is the
dimensionality of the coupling graph parameterisation. In some
embodiments, the coupling graph is parameterised by $O(N^2)$
strength-delay pairs; in sparse topologies, the effective dimensionality
scales as $O(N)$. This count is an addressable-parameter inventory
of the disclosed apparatus and is not asserted as a guaranteed
effective capacity, task-performance bound, or training-scale
result. Effective capacity, generalisation, and task performance
are empirical and embodiment-specific, and any capacity claim
elsewhere in this disclosure is governed by the capacity-accounting
discipline of that subsection rather than by the present axis count.
Under the declared per-reactor conditional independence given
coupling parameters, the network Fisher metric admits the
block-diagonal-plus-coupling decomposition
$\mathcal{F}_{\mathrm{net}}
= \mathrm{blockdiag}(\mathcal{F}_1, \ldots, \mathcal{F}_N)
+ \mathcal{F}_{\mathrm{coupling}}$,
where $\mathcal{F}_{\mathrm{coupling}}$ captures how task loss depends
on topology and inter-reactor coupling parameters and may be
estimated empirically from cross-reactor score covariances; in some
embodiments, the natural gradient defined by this product Fisher
metric is used to navigate the parameter space, with the empirical
efficacy of such navigation being embodiment-specific and not
asserted as a convergence guarantee.

\paragraph{Regime reinterpretation in computational mode (non-limiting).}
In computational mode, the three regimes and Yoked modifier
reinterpret: \TB becomes verified computation (is this output
consistent with computation by a known physical device?); \LI
becomes inference and feature extraction (what does this input data
contain?); \RT becomes generation and transformation
(produce output matching a target distribution); and Yoked operation
becomes unsupervised learning (lock the reactor network's dynamics to
the input's statistical structure, discovering natural modes and
clusters by physical entrainment). Non-limiting interpolations include:
\LI plus \RT as a physical autoencoder where the
attractor dimension $d_{\mathrm{eff}}$ is the bottleneck; \RT plus Yoked as generative entrainment (generating samples
entrained with the data distribution); and \TB plus \RT as signed generation where every output carries proof of
which device created it. Additional non-limiting interpolations
include: \LI plus Yoked as self-supervised learning (the reactor
discovers the input's statistical structure without labels by
physical entrainment, then uses that structure for feature
extraction); \TB plus \LI as verified inference (extract
features from data with simultaneous attestation that the inference
was performed by a specific physical device); and all three regimes
with Yoked coupling engaged, operating simultaneously as a physical
agent that perceives, generates, verifies,
and entrains in a continuous loop.

\paragraph{Scope of computational embodiments (non-limiting).}
In the disclosed computational embodiments, computation is performed by
a physical medium---linear, nonlinear, or a combination---whose
dynamics carry out the computation.  Nonlinear media provide empirical
hardness and device-specific signatures (the reactor-microstructure unclonability property); linear
media provide high-fidelity operators suitable for adjoint-pair
gradient propagation and matrix computation.  Both are constrained by
measured noise, drift, bandwidth, and dissipation inherent to the
physical substrate.
Systems in which computation is performed by software simulation of
dynamics, or in which the ``reactor'' is a purely social, economic, or
abstract process without a physical substrate performing
transformation of physical signals, are not required and are not the
focus of the embodiments described herein.

\paragraph{Haptic and tactile rendering (non-limiting).}
In some embodiments, \RT includes non-visual outputs
such as haptic fields and tactile patterns that are generated and
audited using the same \cba machinery. Non-limiting
embodiments include ultrasonic phased arrays that generate mid-air
haptic sensations, vibrotactile actuators embedded in wearables or
furnishings, electrostatic touch surfaces, and pneumatic or mechanical
actuator arrays. The control protocol specifies spatial and temporal
patterns (including per-element phase and amplitude schedules for
phased arrays), and detectors include microphones, inertial sensors,
contact sensors, cameras, and/or physiological channels (where
available and permitted) to measure delivered fields and responses. In
preferred embodiments, the protocol digest records safety envelopes
and exposure limits, and the \cb binds what was delivered to
what was observed for later audit or dispute resolution.

\paragraph{Immersive chamber and tiled environments (non-limiting).}
In some embodiments, arrays of \RK modules are arranged in a
PolieDeck: a tiled environment (for example walls, ceiling, and floor)
instrumented with multiple projector--detector pairs and optional
reactor loops, so that the room itself functions as a coupled scene for
\LI, \TB, and \RT regimes. In some
embodiments, PolieDeck deployments use trust-weighted witnesses and
transitive proof-of-projection records to corroborate that immersive
content was physically executed in the chamber, with selective opening
used to audit disputed segments with bounded disclosure.

\subsection{World models, planning, and control}

In some embodiments, \RKs provide a physically grounded
interface for planning and control, including learned world models,
model-predictive control, and reinforcement learning with
physically-derived rewards.

\subsection{Materials characterisation and autonomous experimentation
(non-limiting)}

In some embodiments, a \RK operating in \LI mode is
integrated into an automated experimental workflow in which the system
designs illumination protocols, directs structured light at a sample
or material under test, captures the \cb, updates an
internal model of the sample's optical and physical properties, and
selects the next protocol---forming a closed loop that iteratively
refines knowledge of the sample without human intervention. In such
embodiments, the scan law, emission schedule, and decoder heads are
jointly optimised so that each measurement maximises expected
information gain (or reduces expected loss toward a declared target)
conditioned on all previously committed \cba windows.

\paragraph{Optical metamaterial and nanophotonic characterisation
(non-limiting).}
In some embodiments, the sample under test is an optical metamaterial,
metasurface, photonic crystal, or nanophotonic structure whose
far-field emission or transmission profile depends on fabrication
parameters. The \RK operates in \LI mode to map the
relationship between fabrication parameters and optical response,
using structured illumination to probe the structure and committed
\cba data to build a predictive model. In some
embodiments, the feedback loop drives an optimisation cycle in which
fabrication parameters are adjusted between measurement rounds,
creating a closed-loop discovery process in which the system
identifies governing relationships between structure and optical
function from empirical data.  In some embodiments, the sample under
test is a time-varying metamaterial or epsilon-near-zero thin film
(Section~\ref{sec:embodiments}), and the structured illumination
includes temporally modulated probes whose spectral interference
patterns characterise the material's switching dynamics, ENZ
crossing frequency, and nonlinear Kerr response.

\paragraph{Verification of experimental provenance (non-limiting).}
In some embodiments, \TB mode provides tamper-evident records of
the measurement conditions, sample identity, and instrument state for
each experimental run. Protocol digests and meter envelopes committed
during the measurement campaign form an auditable evidence chain, so
that results reported from the closed-loop process can be verified by
third parties via selective opening.

\subsection{Security, privacy, provenance, and compliance}

In some embodiments, \RKs provide physically grounded
evidence for authenticated capture and tamper evidence, liveness checks
and proof-of-projection, privacy-preserving operation (selective
disclosure, secure aggregation, differential privacy), regulatory and
audit use cases, and insurance documentation. Yoked variants add
continuous tamper detection where ongoing scene--reactor coupling
provides a persistent liveness signal that an adversary would need to
maintain in real time.

\paragraph{Journalism and public-interest provenance (non-limiting).}
In some embodiments, \RKs are used in journalism and
investigations to produce tamper-evident capture records for images,
video, and measurements. The \cb, protocol digests, and meter
envelopes are bound into committed records that can be selectively
opened under dispute, making it more difficult to substitute or
fabricate media after capture.

\paragraph{Remote notarisation and contract execution (non-limiting).}
In some embodiments, a \RK supports remote notarisation or
contract execution by attesting that a declared person, object, or
document was physically present in a declared session under a declared
policy. Liveness checks and proof-of-projection protocols reduce replay
and relay attacks, and multi-witness cross-attestation is used for
higher-stakes sessions.

\paragraph{Collateral and inventory monitoring (non-limiting).}
In some embodiments, \RKs verify and monitor collateral or
inventory (for example pledged goods, equipment, or sealed containers)
by executing periodic scan protocols and committing the resulting
evidence. Auditors can request selective openings of committed windows
when anomalies are detected.

\paragraph{Privacy-preserving operation (non-limiting).}
In some embodiments, privacy is supported by one or more of: local
processing, selective disclosure of committed atoms rather than full
traces, secure aggregation of meter outputs, differential privacy
mechanisms in the physical or digital layers, and trusted execution
environments for digital post-processing.

\subsection{Data infrastructure, retrieval, and embedding stores}

In some embodiments, the \cb and derived features are stored as
an evidence archive and may be indexed for retrieval. Retrieval supports
similarity search, anomaly detection, and forensics, with selective
opening of committed atoms for audit. In some embodiments, the system
uses optical correlation, memory media, or reactor dynamics to implement
fast approximate similarity operations on embeddings or signatures
derived from the \cb.

\subsection{Human interaction and personalisation}

\paragraph{On-device adaptation (non-limiting).}
In some embodiments, a personal device includes a compact \RK module and adapts models locally using physically grounded
signals, while exporting only commitments, aggregates, or
privacy-preserving updates.

\paragraph{Assistive interfaces and AR/VR (non-limiting).}
In some embodiments, head-mounted displays provide camera feeds and
pose telemetry to a \RK and receive environment maps,
occlusion masks, calibration parameters, or stylised overlays.
\Cba logs support audit of what was shown and what was
measured.

\subsection{Narrative-charged tokens and economic value discovery
(non-limiting, contextual application)}

\paragraph{Industrial-applicability scope of this subsection
(non-limiting).} The technical contribution disclosed in this
subsection is physical-provenance measurement and verification,
including anti-counterfeiting, authenticity attestation, tamper-
evidence over committed physical evidence windows, and
cryptographically-bound binding of physical tokens to
device-specific optical microstructure signatures. Economic and
collective-decision uses described below are presented as
non-limiting downstream applications of the physical-provenance
machinery; they are not intended as the inventive contribution of
this subsection, and no claim of this disclosure depends on those
downstream uses for its inventive support. The independent claim
hooks of this disclosure rest on the physical-provenance,
authenticity-verification, and tamper-evidence functions disclosed
elsewhere in the present specification; the present subsection
provides contextual deployment description only.

In some embodiments, a narrative-charged microstructure-signed token participates in
\RK verification protocols by binding a physical
token's unclonable optical microstructure signature to declared physically
anchored provenance data through committed evidence windows. The
technical operations performed are: (i)~committed binding of the
token's microstructure signature to time-stamped scans at a designated site;
(ii)~committed binding of attestations from designated guardians
or co-participants to the token's microstructure signature; (iii)~committed
binding of shared material provenance to the token's microstructure signature
through optical microstructure cross-correlation; and (iv)~selective opening
of any of the foregoing committed bindings under dispute or
audit. Each of these operations is a physical-provenance
measurement function; none of them is, by itself, an economic or
collective-decision function.

\paragraph{Downstream economic and collective-decision applications
(non-limiting, contextual).}
In some embodiments, the foregoing physical-provenance machinery
may be used as input to downstream economic or collective-decision
processes external to the \RK itself, for example
revealed-preference aggregation (auction clearing prices, bid
trajectories, holding times, transfer patterns), measurement-side
attention proxies (query frequency, protocol complexity and
duration, multi-party attestation rates, information-theoretic
features), or governance rules using provenance depth and
verification history as inputs to vote-weight assignment or
membership pricing in declared interaction sets. These downstream
applications are described as contextual deployment scenarios; they
are not claimed as inventive contributions of this disclosure,
and the technical anchor for any such downstream application is
the physical-provenance and authenticity-verification machinery
described above and elsewhere in this specification.

\begin{definition}[Narrative-charged microstructure-unclonability artefact (non-limiting)]
A \emph{narrative-charged microstructure-unclonability artefact} is a physical token whose
unclonable optical signature is committed in a \RK
evidence record together with declared physically anchored
provenance data binding the token to a declared event, location,
artefact, relationship, or interaction history. In some
embodiments, such artefacts may be referred to as
``Narravite'' crystals or similar, as a non-limiting naming
convention.
\end{definition}

\paragraph{Interaction-set membership as a verifiable provenance
attestation (non-limiting).}
In some embodiments, the system supports verifiable attestations
that a token (and by extension its holder) participated in a
defined class of events, locations, attestation graphs, or
material lineages. Non-limiting attestation forms include
co-location with other tokens in a scene scan, temporal priority
over a state change established by committed timestamps,
verification by a quorum of designated \RKs, presence
in a named ritual or ecological survey under a declared protocol,
or shared fabrication substrate detectable through optical microstructure
cross-correlation. The technical contribution is the verifiable
attestation; downstream pricing, gating, or governance uses of
such attestations are external to the \RK and are not
claimed as inventive contributions of this disclosure.

\paragraph{Learned anomaly detectors over committed provenance
records (non-limiting, contextual).}
In some embodiments, learned predictors operating over committed
provenance records may flag anomalous patterns where attention or
external usage patterns deviate sharply from the committed
provenance profile, supporting fraud detection and authenticity
audit in systems that combine physical microstructure-unclonability components and commitments. The
technical contribution disclosed in this paragraph is
provenance-anomaly detection over committed physical evidence;
any market-surveillance use of such anomaly signals is a
downstream contextual application external to the \RK
and is not claimed as an inventive contribution of this
disclosure.

\subsection{Scientific, environmental, and industrial sensing}

In some embodiments, \RKs serve in remote sensing and
environmental monitoring (airborne/satellite platforms, buoys,
underwater vehicles), agriculture (crop health, pest detection, yield
estimation), infrastructure inspection, materials characterisation,
geophysical and geological sensing, water-quality monitoring, and
scientific sample tracking. RF, radar, and telescope array embodiments
are non-limiting extensions.

\paragraph{Wildlife camera traps and conservation sensing
(non-limiting).}
In some embodiments, a \RK is integrated with a camera trap
or other conservation sensor to produce verifiable imagery of wildlife
and habitats. Commitments and device identity bindings support audit
that disclosed images were captured by an enrolled device under a
declared protocol and were not substituted or materially altered after
capture.

\paragraph{Pollution source monitoring and environmental compliance
(non-limiting).}
In some embodiments, \RKs provide tamper-evident monitoring
of pollution sources and environmental quality, including industrial
outfalls and distributed air or water sensors. Protocol digests and
meter envelopes document calibration state and sampling conditions, and
multi-device corroboration detects missing data, selective reporting,
or inconsistent histories.

\paragraph{Acoustic biodiversity monitoring (non-limiting).}
In some embodiments, the same evidence and commitment pipeline is
applied to acoustic monitoring (for example microphones or
hydrophones) to produce authenticated recordings of animal
vocalisations or anthropogenic noise, enabling auditable biodiversity
surveys and anti-poaching monitoring.

\paragraph{Environmental DNA (eDNA) chain-of-custody (non-limiting).}
In some embodiments, a \RK is coupled to an eDNA sampling
or analysis workflow (for example a cartridge-based sampler, filtration
cassette, or lab bench instrument), and the sampling act is logged
with protocol digests, device identity, and time/location metadata so
that downstream genetic measurements are bound to a verifiable chain
of custody.

\paragraph{Carbon and ecosystem accounting (non-limiting).}
In some embodiments, verifiable environmental observations support
carbon and ecosystem accounting, including forest-cover and biomass
monitoring, verification of agricultural practices associated with
carbon sequestration (for example cover cropping, no-till, or
rotational grazing), and monitoring of blue-carbon ecosystems (for
example seagrass, mangroves, or kelp) using underwater or aerial
deployments.

\paragraph{Citizen science and community networks (non-limiting).}
In some embodiments, low-cost or smartphone-integrated \RKs
enable citizens and community groups to contribute verifiable local
observations (air and water quality, habitat changes) while preserving
privacy via selective disclosure and secure aggregation; cross-device
corroboration reduces the impact of compromised or misconfigured nodes.

\paragraph{Adaptive sampling, local intervention, and pest control
(non-limiting).}
In some embodiments, a controller adapts sampling schedules and
illumination patterns based on detected changes, and may trigger
bounded interventions (for example alarms, valve actuation, aeration,
or non-lethal deterrence) under safety constraints, while logging
emissions and responses for audit. In some embodiments, \RKs
are integrated with pest or vector-control devices; meters classify
organisms transiting a sensing volume and local interventions are
gated by logged meter output and policy constraints.

\subsection{Medical and biomedical embodiments (non-limiting)}

In some embodiments, \RK architectures are applied to medical or
biomedical sensing and intervention systems. The present disclosure concerns
device architecture, control, logging, and verification mechanisms; clinical
deployment of any embodiment would proceed under the regulatory and
validation requirements applicable to the relevant indication.
Subsystems include imaging and guided scanning, telemedicine and
longitudinal monitoring, closed-loop actuation under safety envelopes,
and ultrasonic actuation with biofeedback and audit.

\subsection{Scientific computing and industrial systems}

In some embodiments, \RKs serve as physical substrates for
fast transforms, filtering, or optimisation, including
reservoir-style computation.

\subsection{Sector-specific deployments}

In some embodiments, \RKs are deployed in agriculture,
energy, logistics, retail, education, gaming, insurance,
telecommunications, and real estate. Yoked operation adds
resonance-based scene characterisation (probing how a building, crop, or
object responds to structured stimulation), liveness detection
(continuous rather than snapshot-based), and distributed coupling proof
for sensor networks.

\subsection{Tooling, ecosystem, and graceful degradation}

In some embodiments, deployments include supporting tools such as
\cba inspection utilities, calibration profilers, SDKs,
simulation toolkits, and certification frameworks. The system is
designed for graceful degradation when reactor hardware drifts or fails.

\paragraph{Cross-substrate equivalence summary (non-limiting).}
The following five sections map core \RK components to five
interpretive substrates. The non-limiting table below summarises the
correspondences; all entries represent non-limiting analogy-level
mappings rather than strict equivalence.

{\small
\begin{longtable}{@{}p{0.13\textwidth}p{0.13\textwidth}p{0.13\textwidth}p{0.15\textwidth}p{0.15\textwidth}p{0.15\textwidth}@{}}
\toprule
RK component & Optical & Neural & Info-theoretic & Control & Crypto \\
\midrule
\endfirsthead
\toprule
RK component & Optical & Neural & Info-theoretic & Control & Crypto \\
\midrule
\endhead
\endfoot
\bottomrule
\endlastfoot
Reactor $Z_\theta$ &
  Scattering medium &
  Latent $\mathbf{z}$ &
  Channel state &
  Plant state &
  microstructure state \\
Emission $\mathbf{e}(t)$ &
  Illumination pattern &
  Input / stimulus &
  Transmitted signal &
  Reference input &
  Challenge \\
Observation $\mathbf{y}_t$ &
  Detected image &
  Network output &
  Received signal &
  Measured output &
  Response \\
Config $\theta$ &
  Actuator settings &
  Weights (analog of $\theta$) &
  Codebook &
  Controller gains &
  Public params / key material \\
Protocol $U_{0:T}$ &
  Scan schedule &
  Training protocol &
  Coding scheme &
  Excitation signal &
  Challenge sequence \\
Meter $\mathbf{m}(t)$ &
  Calibration checks &
  Loss / diagnostics &
  Rate / divergence &
  Error signals &
  Verification score \\
\Cb $C_{0:T}$ &
  Committed disclosure atoms &
  Activations log &
  Codeword sequence &
  I/O trajectory &
  Transcript \\
\end{longtable}
}

% ======================================================================
\section{Optical Primitives Substrate}
\label{sec:optical-primitives}
% ======================================================================

This section catalogues non-limiting optical primitives that may be realised within a \RK module and used as building blocks for verification, perception, and controllable rendering. In the embodiments described in this section, each primitive is accessed as part of a sweepable, parameterised kernel: the primitive is implemented as at least one of (i) a \emph{sweeping kernel} that modulates an emitted field prior to scanning or addressing, (ii) a \emph{stationary kernel} or reactor element that is swept by a scanner or addressed illumination after scanning, or (iii) a multi-stage combination of sweeping and stationary kernels. This discussion is not intended to limit the scope to any single optical effect; it shows that a single closed-loop controllable \RK can traverse a wide range of optical regimes under software control while preserving \cba logging and the Markov-kernel framing.

\subsection{Coherence regimes}
In some embodiments, the system operates under a chosen coherence regime that affects measurement sensitivity, stability requirements, and the structure of the \cb.
Coherence may be temporal, spatial, or both, and the chosen regime is recorded in protocol digests and meters.

\paragraph{Coherent operation (non-limiting).}
In some embodiments, coherent illumination enables interference, phase-sensitive signatures, and high sensitivity to small path-length and alignment changes.
This can strengthen hardness by exposing microstructure and coupling that is difficult to emulate, and it can also increase sensitivity to vibration, thermal drift, and alignment.

\paragraph{Partially coherent operation (non-limiting).}
In some embodiments, partially coherent operation reduces sensitivity to nuisance fluctuations while retaining useful structured signatures.
This supports regimes where stability meters are bounded but tight phase lock is not required.

\paragraph{Incoherent operation (non-limiting).}
In some embodiments, incoherent illumination supports intensity-based measurements with reduced phase sensitivity.
This can improve robustness and simplify modelling, while still supporting multi-channel constraints (for example spectral channels, polarisation, and scan-dependent signatures).

\paragraph{Metering and logging (non-limiting).}
In some embodiments, meters record coherence-related indicators such as source bandwidth estimates, coherence-length proxies, speckle statistics, fringe contrast (when applicable), temperature, vibration, and optical power stability.
Protocol digests record the chosen coherence regime and relevant source configuration.

\paragraph{Speckle-statistics calibration (non-limiting).}
In some embodiments, a scattering-medium reactor is characterised using
speckle-statistics measurements of the kind treated in Goodman's
\emph{Speckle Phenomena in Optics} (2007), including first- and
second-order intensity statistics, speckle contrast, spatial and
temporal correlation length, polarisation-dependent speckle structure,
and memory-effect angular range where applicable. The speckle-statistics
summary is included in the calibration record and is stratified by
wavelength, coherence regime, detector exposure, polarisation state,
scan geometry, reactor temperature, and medium history.

\subsection{Linear canonical transforms and convolution primitives}
In some embodiments, optical propagation and engineered transfer functions implement linear canonical transforms (LCTs) and convolution-like operators that shape how stimuli map into observations.

\paragraph{LCT family (non-limiting).}
In some embodiments, optical propagation is described by an LCT family with a kernel parameterisation.
Fourier, Fresnel, and fractional Fourier transforms are non-limiting special cases.
In some embodiments, a tunable parameter acts as an invariance knob, where stable structure persists over a parameter range while mismatched emulators exhibit characteristic deviations.

\paragraph{Kernel form (non-limiting).}
In some embodiments, a linear canonical transform is written using an integral kernel:
\[
  (\mathcal{T}_\theta S)(x) = \int K_\theta(x,x')\,S(x')\,dx'.
\]
In some embodiments, tunable optical elements implement a programmable LCT engine by sweeping $\theta$ (for example focal length, propagation distance, aperture, or phase-mask settings), and protocol digests record the sweep schedule so that verifiers interpret observations under the correct operator family.

\paragraph{Convolution and point-spread functions (non-limiting).}
In some embodiments, the system is approximately shift-invariant over a region of interest, yielding a convolutional mapping between a stimulus and an effective point-spread function (PSF).
In other embodiments, the mapping is shift-variant, and the protocol digest records the scan law and optical configuration so that verifiers interpret observations under the correct operator family.

\paragraph{Scan-controlled kernels (non-limiting).}
In some embodiments, scanning and addressing cause the effective operator to vary over time, producing a controlled family of kernels indexed by scan coordinate, wavelength, focus, or polarisation.
Meters capture alignment and stability so that the selected transform regime is reproducible and auditable.

\subsection{Correlation and matched filtering}
In some embodiments, correlation operations are used as primitives for alignment, verification, detection, and reconstruction, and are implemented optically, digitally, or in hybrid form.

\paragraph{Cross-correlation (non-limiting).}
In some embodiments, cross-correlation between channels or between predicted and observed windows is used to estimate timing offsets, scan alignment, or scene motion.
Correlation peaks, sidelobe ratios, and consistency across channels are recorded as meter-like summaries.
In some embodiments, a non-limiting cross-correlation form is
\[
  \mathrm{corr}_T(S)(\Delta x) = \int S(x)\,\overline{T(x - \Delta x)}\,dx,
\]
where $T$ is a template (for example an expected response under a logged protocol or a disclosed probe window).

\paragraph{Matched filtering (non-limiting).}
In some embodiments, a known emitted pattern (or a digest-derived template family) is matched-filtered against observations to detect whether expected structure is present.
This supports verisimilitude discrimination, projection-hardness checks, and protocol binding, particularly when selective openings disclose only sparse windows or derived features.
In some embodiments, matched filtering is implemented optically using a correlator architecture (for example a VanderLugt correlator in a 4f optical system) as a non-limiting implementation detail.

\paragraph{Audit compatibility --- correlation and matched filtering (non-limiting).}
In some embodiments, correlation summaries are derived from committed atoms and protocol digests so that a verifier can request openings if correlation-level checks fail or if higher-confidence auditing is required.

\subsection{Interferometry and phase-sensitive primitives}
In some embodiments, phase-sensitive measurement primitives are used to increase sensitivity and to impose constraints that are difficult for emulators to reproduce under meter-bounded operation.

\paragraph{Phase sensitivity (non-limiting).}
In some embodiments, phase $ \phi(t) $ depends on optical path length and refractive index variations.
Small mechanical or thermal changes can produce measurable fringe shifts.
This can strengthen authenticity checks by coupling measurements to fine-grained physical behaviour.

\paragraph{Interferometric configurations (non-limiting).}
Non-limiting embodiments include two-beam interference, common-path interferometry, heterodyne-style phase measurement, or phase retrieval from controlled phase shifts.
Protocol digests record the wavelength, path configuration, modulation schedule, and sampling timing used for phase-sensitive operation.

\paragraph{Homodyne, heterodyne, and balanced detection (non-limiting).}
In some embodiments, phase-sensitive measurement is implemented using a local oscillator (LO) combined with a signal field at a beamsplitter and measured using balanced photodetection. In these embodiments, the LO phase $\phi_{\mathrm{LO}}(t)$ and amplitude may be swept as part of the control protocol $U_{0:T}$ (analogous to a phase scan law), producing quadrature-like observables that are included as channels of $\mathbf{y}_t$. In some embodiments, heterodyne operation is achieved by offsetting the LO frequency so the beat note encodes phase and amplitude over a logged window.

\paragraph{Meters and stability (non-limiting).}
In some embodiments, meters record fringe visibility, vibration proxies, temperature, coherence indicators, and phase-lock status (when applicable).
Verification is conditioned on these meters, and phase-sensitive regimes may be gated or downgraded when stability is insufficient.

\subsection{Polarisation primitives}
In some embodiments, polarisation provides additional channels and constraints for both sensing and verification.

\paragraph{Polarisation channels (non-limiting).}
In some embodiments, the system measures or controls Stokes parameters, polarisation-resolved intensities, or polarisation-dependent responses.
Polarisation channels can increase distinguishability between genuine and emulated channels, and can add invariance profiles across controlled polariser angles.

\paragraph{Polarisation optics and logging (non-limiting).}
In some embodiments, polarisation is controlled by polarisers, waveplates, liquid-crystal modulators, or birefringent elements.
Protocol digests record polarisation settings, modulation schedules, and calibration versions.
Meters record extinction ratios and polarisation stability indicators.

\subsection{Holographic and volumetric storage primitives}
In some embodiments, the system uses memory media that store optical states or mappings in a persistent or semi-persistent manner, supporting both reconstruction and hardness.

\paragraph{Volumetric or holographic storage (non-limiting).}
In some embodiments, volume holograms or interference-based media store patterns or transfer functions.
Retrieval can behave as content-addressable memory, where the output depends sensitively on alignment, wavelength, and incident field structure.

\paragraph{Hardness and auditability (non-limiting).}
In some embodiments, stored media increase analogue hardness because reproducing the same retrieval behaviour, under the stated attacker model and meter tolerances, typically requires matching physical microstructure and alignment.
Write and read protocols are logged via protocol digests and meters, and the resulting records are committed so that later audits can verify that a claimed memory state was actually exercised.

\paragraph{State persistence (non-limiting).}
In some embodiments, persistence times and decay profiles are metered and treated as part of the operating envelope, especially when memory effects influence subsequent emissions or observations.

\subsection{Spectral and wavelength-selective primitives}
\label{sec:spectral-primitives}
In some embodiments, spectral selectivity provides multiplexing and additional constraints.

\paragraph{Wavelength multiplexing (non-limiting).}
In some embodiments, the system emits and observes at multiple wavelengths, either sequentially or with multiplexing.
Spectral signatures can constrain emulation, improve reconstruction, and help verify that an observation is consistent with a logged protocol regime.
In some embodiments, wavelength multiplexing is implemented as a spectral observation primitive using a continuously tunable frequency-comb emitter of the class described in Section~\ref{sec:emitter-catalogue}, supporting, when paired with a second comb or suitable local-oscillator/reference path, dual-comb or multi-heterodyne interrogation of the reactor or scene spectral response as a declared observation channel, with the comb-mode-spacing schedule, heterodyne offset, and detection bandwidth recorded in the protocol digest; any verification or hardness use of such records remains conditioned on declared meter envelopes and empirically measured hardness indices, rather than on comb tuning range or waveform-space size alone.

\paragraph{Dispersive and filtering elements (non-limiting).}
In some embodiments, dispersive optics, tunable filters, diffraction gratings, or spectrally selective materials shape the channel.
In some embodiments, compact photoelastic dispersive structures written into thermoplastic polymer substrates (for example polycarbonate) by ultrafast laser pulses provide broadband, view-angle-independent spectral dispersion in footprints of the order of $10~\mu\mathrm{m} \times 10~\mu\mathrm{m}$ and integrate directly with a CMOS image sensor for per-pixel spectral analysis, as described in Zhang et al., ``Optical dispersion using micro-vortices in thermoplastic polymers for integrated microspectrometers,'' \emph{Nature Electronics} (2026).
Protocol digests record wavelength schedules, bandwidth, and filter states.
Meters record wavelength lock indicators, power stability, and spectral calibration status.

\paragraph{Audit compatibility --- spectral and wavelength-selective (non-limiting).}
In some embodiments, spectral summaries and wavelength schedules are committed and selectively disclosed, enabling verification that disclosed evidence is consistent with the claimed spectral regime.

\subsection{Random scattering and random projection primitives}
In some embodiments, random scattering media and disordered optical paths implement high-dimensional mixing that behaves like a random projection.
This can amplify microstructure sensitivity and increase analogue hardness, because small physical differences can cause measurable changes in the \cb distribution under the same protocol digest.

\paragraph{Random projection viewpoint (non-limiting).}
In some embodiments, an input field or pattern is mapped through a random medium to a mixed observation that behaves approximately like a random feature map.
The resulting readout can be treated as a high-dimensional embedding whose statistics are difficult to emulate without matching the physical medium and alignment.
In some embodiments, this is interpreted as a Johnson--Lindenstrauss-style random projection lens (after suitable normalisation), where distances or inner products are approximately preserved in the embedding, enabling discrimination and learning with small downstream heads.
In some embodiments, a local linearised scattering model is used as an interpretation lens:
\[
  \mathbf{c} = A_{\mathrm{scatter}}(\theta)\,\mathbf{s} + \mathbf{w},
\]
where $\mathbf{s}$ denotes an input pattern or field parameterisation, $\mathbf{c}$ denotes a corresponding observation vector, $A_{\mathrm{scatter}}(\theta)$ denotes an effective mixing operator induced by microstructure and configuration, and $\mathbf{w}$ denotes noise.

\paragraph{Speckle and stability (non-limiting).}
In some embodiments, speckle statistics provide useful signatures, and meters record speckle contrast, stability over time, and sensitivity to scan coordinate.
Protocol digests record illumination conditions and scan settings so that verifiers interpret speckle and correlation summaries under the correct regime.

\paragraph{Audit compatibility --- random scattering (non-limiting).}
In some embodiments, random-scattering outputs are committed as atoms and opened selectively.
In some embodiments, verifiers request openings that test consistency across multiple scan points, wavelengths, or polarisation settings to reduce the probability of emulator overfitting to a single view.

\subsection{Nonlinear optics and saturating media}
In some embodiments, nonlinear or saturating media provide amplitude-dependent transformations that increase distinguishability and make simple linear emulators ineffective.
Non-limiting examples include saturable absorbers, gain media near saturation, Kerr-like effects, photorefractive effects, and other media with intensity-dependent response.

\paragraph{Nonlinear response as a signature (non-limiting).}
In some embodiments, the effective mapping from emitted patterns to observed summaries depends on intensity history and local field distribution.
This can create protocol-dependent signatures that are hard to match with an emulator that lacks the same nonlinearity and internal state.

\paragraph{Metering and safe operation (non-limiting).}
In some embodiments, meters record saturation indicators, thermal indicators, drift, and stability bounds.
Protocol digests record drive intensity schedules and duty cycles.
In some embodiments, nonlinear regimes are gated by safety envelopes and downgraded when meters indicate instability.

\paragraph{Interaction with verification and rendering (non-limiting).}
In some embodiments, nonlinear media strengthen \TB by amplifying microstructure sensitivity under bounded protocols.
In some embodiments, nonlinear regimes are also used in \RT, where the nonlinearity shapes the controllable rendering dynamics while remaining meter-bounded and auditable.

\subsection{Resonance, frequency-domain and cavity primitives}
In some embodiments, resonant structures provide frequency-domain selectivity and memory.
Non-limiting examples include optical cavities, resonant filters, ring resonators, and other frequency-selective elements whose response depends on alignment, temperature, and internal state.

\paragraph{Frequency response and tuning (non-limiting).}
In some embodiments, the system probes a resonance by sweeping frequency, wavelength, or modulation rate and recording a response curve.
The protocol digest records sweep schedules and tuning parameters.
Meters record resonance stability, drift, and lock indicators.

\paragraph{Cavity memory and dynamics (non-limiting).}
In some embodiments, cavity dynamics introduce temporal memory and mixing.
A short burst can influence later response, creating a channel with history dependence.
This can increase analogue hardness because matching the same transient dynamics, under the stated attacker model, materially raises the cost of matching both the structure and the operating regime.

\paragraph{Audit and selective opening (non-limiting).}
In some embodiments, the verifier challenges frequency-domain consistency by requesting openings across multiple probe points, rather than a single point on a response curve.
This reduces the probability that an emulator can fabricate a consistent frequency signature without matching the underlying physical response.

\subsection{Time-of-flight, phase, and timing-anchor primitives}
In some embodiments, timing provides a robust primitive for both sensing and verification.
Non-limiting examples include time-of-flight, phase-delay measurements, modulated illumination, and timing anchors derived from scan schedules and detector timing models.

\paragraph{Timing as a constraint (non-limiting).}
In some embodiments, the \cb includes timing-dependent features that are required to be consistent with the logged protocol and the metered synchronisation state.
This can strengthen verification by forcing consistency across emission timing, scan coordinates, and detector readout timing.

\paragraph{Time-of-flight and phase delay (non-limiting).}
In some embodiments, depth or delay is inferred from phase shifts or time-of-flight signatures.
Protocol digests record modulation frequency, exposure models, and sampling schedules.
Meters record timing jitter, clock stability, and alignment indicators.

In some embodiments, time-of-flight is obtained using pulsed illumination with gated detection or time tagging (TCSPC as a non-limiting example), where $g_{\mathrm{det}}(t)$ encodes gate placement and width and the observation includes a per-step histogram of arrival times. In other embodiments, time-of-flight is obtained using continuous-wave modulation and phase delay, including frequency-chirped or multi-tone modulation (FMCW-style as a non-limiting example) in which a beat frequency or demodulated I/Q pair is recorded as part of $\mathbf{y}_t$ under a logged chirp schedule.

\paragraph{Anchoring and freshness (non-limiting).}
In some embodiments, timing anchors bind evidence windows to externally verifiable time.
This interacts with batching and Merkle commitments: the system commits to windows and anchors commitments so that later openings have a clear freshness interpretation.

\paragraph{GNSS positioning and timing as auxiliary channels (non-limiting).}
In some embodiments, a GNSS receiver (for example GPS, Galileo, GLONASS,
BeiDou, or multi-constellation receivers) is coupled to the \RK and provides position fixes, velocity estimates, and precision
timing references as components of the auxiliary observation vector
$y_{\mathrm{ext}}(t)$.  Position and timing metadata are recorded in the
protocol digest alongside device identity and calibration identifiers,
binding \cba evidence to a declared geospatial location
and an externally traceable time source.

In some embodiments, the GNSS-derived timing reference serves as a
timing anchor that supplements or replaces local oscillator references,
and meters monitor the consistency between GNSS-derived time, local
clock drift estimates, and scan-schedule timestamps.  Detected
discrepancies (for example GNSS signal loss, spoofing indicators, or
timing jumps) are flagged in the meter envelope and may trigger
downgrade to a degraded-confidence operating mode or increased reliance
on alternative timing sources (for example a local atomic clock,
network time protocol, or cross-device timing overlap).

In some embodiments, GNSS-derived position is used in multi-device
corroboration to verify geometric consistency: the claimed positions of
two or more \RKs are checked against cross-attestation
timing, overlap evidence, and signal-propagation constraints.  In some
embodiments, differential GNSS or real-time kinematic (RTK) corrections
are logged as part of the protocol digest so that verifiers can assess
the declared positioning accuracy.

In some embodiments, GNSS pseudorange and carrier-phase observables are
logged as raw auxiliary channels (rather than only processed position
fixes), enabling retrospective reprocessing, integrity analysis, and
detection of multipath or interference signatures.  In such
embodiments, the raw GNSS observables are treated as additional
components of the conditioned observation stream and may be committed
alongside the optical \cb.

\subsection{Structured illumination families and multiplexing}
In some embodiments, structured illumination patterns are selected from families that provide controllable information content, verification signatures, and rendering effects.
Non-limiting families include pseudo-random patterns, coded apertures, Hadamard-like patterns, sinusoidal phase shifts, and multi-scale or multi-frequency patterns.

\paragraph{Pattern families and protocol digests (non-limiting).}
In some embodiments, pattern selection is driven by a seed or by a logged pattern schedule.
Protocol digests record seeds, family identifiers, and schedule parameters so that opened atoms can be checked against the claimed emission regime.

\paragraph{Multiplexing (non-limiting).}
In some embodiments, multiplexing is performed across time, wavelength, polarisation, or spatial tiles.
Multiplexing can increase robustness and reduce acquisition cost, while also increasing the constraint set available for verification.
Meters record cross-channel alignment and stability so multiplexed observations remain comparable.

\paragraph{Orbital angular momentum and higher-order spatial modes
(non-limiting).}
In some embodiments, structured illumination patterns include beams
carrying orbital angular momentum (OAM), Laguerre--Gaussian modes,
Hermite--Gaussian modes, or other higher-order spatial mode families.
OAM beams carry helical phase fronts characterised by an integer
topological charge and provide a theoretically unbounded set of
orthogonal states, each of which can serve as an independent
computational or encoding degree of freedom alongside time, wavelength,
and polarisation. In some embodiments, OAM mode indices are included in
the control protocol $U_{0:T}$ and recorded in protocol digests, so that
verification and reconstruction can condition on the spatial mode
structure of the emission. In some embodiments, mode mixing introduced
by the physical channel (scattering, turbulence, refractive-index
inhomogeneity) is treated as a device-specific and scene-specific
transformation that contributes to the \cb's
distinguishability and hardness properties rather than as noise to be
corrected.

\paragraph{Illumination as a trainable computational resource
(non-limiting).}
In some embodiments, the structured illumination pattern itself---not
only the optical elements in the feedback path---participates in
training or optimisation. The emission schedule $u(t)$ and any
spatial, spectral, or polarisation structure of the emitted field are
treated as learnable parameters that are jointly optimised with
downstream heads (decoders, classifiers, verifiers) via
backpropagation through a differentiable model of the physical channel
or via policy-gradient methods. In such embodiments, the illumination
pattern is not a fixed probe but a co-adapted component of the
computational pipeline, and the protocol digest records the pattern
parameters so that the computational contribution of the illumination
is auditable.

\paragraph{Physical-channel distortion as computation (non-limiting).}
In some embodiments, the mode mixing, scattering, dispersion, and
nonlinear transformation introduced by the physical medium between
emitter and detector are treated not as degradation to be compensated
but as computation performed by the medium itself. When structured
light (including higher-order spatial modes) propagates through a
complex medium, the medium's transfer function acts as a high-dimensional,
physically fixed transformation analogous to the weights of a neural
network layer. In such embodiments, only the readout (decoder or
classifier applied to the \cb) is trained, and the
physical medium provides the rich feature space. This interpretation
unifies the roles of the reactor, the scene, and any intervening
medium: all contribute computational transformations to the \cb, and the distinction between ``device computation'' and
``channel distortion'' is a modelling choice rather than a physical
boundary.

\paragraph{Hybrid verification and sensing (non-limiting).}
In some embodiments, the same pattern family supports both \LI information acquisition and \TB verification.
For example, patterns are chosen to maximise information gain while remaining within regimes where verisimilitude and hardness remain high and auditable.

\subsection{Gain media and optical amplification}
In some embodiments, the module includes optical gain or amplification, for example an active medium that increases signal energy under pump control.
Gain primitives can introduce sensitivity, nonlinearity, and history dependence that strengthen distinguishability and increase analogue hardness, while requiring explicit metering and regime logging.

\paragraph{Gain control and operating regimes (non-limiting).}
In some embodiments, the gain element is controlled by a pump schedule, bias, or drive waveform recorded in the protocol digest.
In some embodiments, the gain is operated in a small-signal regime, a saturated regime, or a near-threshold regime, and the selected regime is treated as part of the protocol and meter envelope.
In some embodiments, gain settings, pump power, wavelength schedules, and polarisation settings are recorded as sweepable parameters included in $\theta$ and/or $U_{0:T}$.

\paragraph{Amplification noise and saturation (non-limiting).}
In some embodiments, amplification introduces additional noise contributions, for example amplified spontaneous emission (ASE) and saturation effects.
A non-limiting intensity-level description is
\[
  I_{\mathrm{out}}(t) = G(t)\,I_{\mathrm{in}}(t) + n_{\mathrm{ASE}}(t),
\]
where $G(t)$ is an effective gain (possibly history dependent) and $n_{\mathrm{ASE}}(t)$ is a non-limiting noise term whose statistics depend on operating regime.
In some embodiments, saturation is treated as a controlled nonlinearity, and meters record saturation indicators and stability bounds so verification and reconstruction remain conditioned on the observed regime.

\paragraph{Metering, stability, and logging --- gain media (non-limiting).}
In some embodiments, meters record gain-related indicators including output power, pump power, gain estimates, saturation proxies, spectral broadening, temperature, and drift.
In some embodiments, these meters gate whether a gain regime is used for verification, sensing, or rendering, and regime downgrades are logged when stability is insufficient.

\paragraph{Audit compatibility --- gain media (non-limiting).}
In some embodiments, gain-regime runs are committed as atoms, and selective opening can reveal probe windows that test consistency of gain behaviour under the logged protocol.
In some embodiments, audits compare expected versus observed amplification statistics (for example correlation and matched filtering under gain-dependent templates) and verify that disclosed windows are consistent with the protocol digest and meter envelope.
In some embodiments, gain primitives increase analogue hardness because emulators would need to reproduce both the amplification behaviour and the associated noise and saturation statistics under the same logged regime.

\subsection{Cavity and eigenmode primitives}
In some embodiments, the module is operated as a cavity or ring-resonator system whose response is governed by resonant modes and eigenstructure.
Cavity and eigenmode primitives provide frequency selectivity, temporal memory, and mode-dependent signatures that can strengthen verification and increase analogue hardness, while requiring explicit stability metering.

\paragraph{Eigenmodes and resonance conditions (non-limiting).}
In some embodiments, the system supports one or more resonant modes whose frequencies depend on geometry, refractive index, coupling, and boundary conditions.
A non-limiting eigenmode form is
\[
  \mathcal{L}_\theta \psi_k = \lambda_k \psi_k,
\]
where $\mathcal{L}_\theta$ denotes a round-trip propagation operator (or other effective cavity operator) under configuration $\theta$.
A non-limiting resonance condition is
\[
  m\lambda \approx 2nL,
\]
where $m$ is an integer mode index, $\lambda$ denotes wavelength (standard optical notation; not the regime-weight symbol), $n$ is an effective index, and $L$ is an effective path length (for example a cavity length).
In some embodiments, the measured response depends on coupling coefficients, loss, and mode competition, and these parameters are treated as sweepable or calibratable components of $\theta$ and/or $U_{0:T}$.

\paragraph{Frequency probing, ring-down, and memory (non-limiting).}
In some embodiments, the cavity is probed by sweeping wavelength, frequency, or modulation rate and recording a response curve or transient.
In some embodiments, ring-down or transient response introduces temporal memory, and a non-limiting time constant is
\[
  \tau \approx \frac{Q}{\omega_0},
\]
where $Q$ is a quality factor and $\omega_0$ is a resonant angular frequency.
In some embodiments, this history dependence increases analogue hardness because an emulator would need to reproduce both steady-state and transient dynamics under the same protocol and meter envelope.

\paragraph{Metering, stability, and logging --- cavity and eigenmode (non-limiting).}
In some embodiments, meters record resonance stability, temperature, vibration proxies, coupling or tuning indicators, lock status (when applicable), and mode-hop indicators.
Protocol digests record sweep schedules, tuning schedules, coupling settings, and calibration versions so that verification and reconstruction are conditioned on the correct cavity regime.

\paragraph{Audit compatibility --- cavity and eigenmode (non-limiting).}
In some embodiments, cavity probes are committed as atoms and selectively opened for audit.
In some embodiments, audits request openings across multiple probe points or time windows so a verifier can check internal consistency of the claimed resonance response rather than a single point.
In networked settings, multiple witnesses can corroborate cavity signatures under shared protocols, and trust-weighted aggregation can reduce the probability of fabricated frequency-domain evidence.

\subsection{Summary (optical substrate)}
This section enumerates non-limiting optical primitives that can be combined with scan protocols, meter-bounded operation, and commitment and selective disclosure mechanisms.

In some embodiments, the primitive set includes coherence regime selection (coherent, partially coherent, incoherent); linear canonical transform and convolutional regimes; correlation and matched filtering; interferometric and phase-sensitive regimes; polarisation-resolved channels; holographic or volumetric storage; wavelength-selective and dispersive channels; random scattering and random projection mixing; nonlinear and saturating media; resonant and cavity frequency-domain regimes; timing and time-of-flight constraints; and structured illumination families with multiplexing.

In some embodiments, optical gain and amplification primitives and cavity or eigenmode primitives provide additional regimes in which microstructure sensitivity, stability, and memory can be tuned and logged.
In some embodiments, photo-reactive and persistent surfaces provide history-dependent state channels whose write/read dynamics are metered and auditable.
In some embodiments, explicit modelling of sensor non-idealities and calibration state versioning supports conditioning verification and reconstruction on the appropriate measurement model.
In some embodiments, cross-domain coupling primitives provide additional constraints by linking optical behaviour to thermal, mechanical, electrical, or acoustic dynamics within a metered operating envelope.

These primitives are non-limiting and may be combined or gated by policy.
In some embodiments, protocol digests record the selected primitive family and configuration, meters record stability and envelope indicators, and commitments enable selective opening to audit that disclosed evidence is consistent with the claimed regime.

\paragraph{Non-limiting summary table.}
{\small
\begin{longtable}{@{}p{0.27\textwidth}p{0.34\textwidth}p{0.34\textwidth}@{}}
  \toprule
  \textbf{Primitive} & \textbf{RK realisation} & \textbf{Use / constraint} \\
  \midrule
\endfirsthead
  \toprule
  \textbf{Primitive} & \textbf{RK realisation} & \textbf{Use / constraint} \\
  \midrule
\endhead
  \bottomrule
\endlastfoot
  Coherence regime selection &
  Choose coherent / partially coherent / incoherent sources; log stability meters &
  Trade sensitivity vs robustness; phase-sensitive signatures vs intensity-only stability \\[0.3em]
  LCT / convolution operators &
  Propagation plus engineered transfer functions; sweepable $\theta$ implements an operator family &
  Programmable optical operators; impose structure and invariances; compare emulator mismatch across $\theta$ \\[0.3em]
  Correlation / matched filtering &
  Cross-correlation and template matching on disclosed windows or predicted responses &
  Timing alignment, scan consistency, and verification checks under selective opening \\[0.3em]
  Interferometry / phase-sensitive probes &
  Interferometers (Mach--Zehnder, Michelson, Sagnac as non-limiting examples) with logged basis and stability &
  Fine-grained path-length sensitivity; hard-to-forge phase constraints under meter envelopes \\[0.3em]
  Polarisation-resolved channels &
  Polarisation optics and multiple detector channels with logged settings &
  Additional constraints and invariance knobs; multi-channel cross-checks \\[0.3em]
  Spectral / dispersive channels &
  Tunable filters, gratings, dispersive media; wavelength schedules in protocol digests &
  Spectral signatures and multiplexing; separate bands for probing, pumping, and presentation \\[0.3em]
  Random scattering / random projections &
  Disordered media; speckle and scattering-matrix viewpoint; meter stability tracking &
  High-dimensional mixing and microstructure sensitivity that amplifies analogue hardness \\[0.3em]
  Nonlinear / saturating media &
  Intensity-dependent elements (saturable absorption, Kerr-like effects, photorefractive media) &
  Nonlinear signatures and history dependence; amplifies differences between genuine and emulated regimes \\[0.3em]
  Gain / amplification primitives &
  Gain media near saturation with metered stability; logged duty cycle and thermal margins &
  Amplify signatures and noise statistics; typically uses explicit stability gating \\[0.3em]
  Cavity / eigenmodes and resonance &
  Ring/cavity configurations with swept tuning and ring-down probing &
  Frequency selectivity and temporal memory; mode-dependent signatures for verification \\[0.3em]
  Timing / time-of-flight anchors &
  Modulation schedules, timing models, and clock stability meters &
  Freshness and causal ordering constraints; detect replay and enforce protocol alignment \\
\end{longtable}
}

% ======================================================================
\section{Neural-Architecture Substrate Mapping}
\label{sec:neural-mapping}
% ======================================================================

This section describes non-limiting correspondences between the \RK formalism and common machine learning architectures. This mapping is not intended to imply that the \RK must be implemented as a particular digital architecture; it shows that a \RK can physically instantiate the computational structure of these architectures as special cases. These correspondences are explanatory embodiments and do not limit the core definition of a \RK.

\subsection{Recurrent architectures (RNN, LSTM, GRU)}
In some embodiments, the \RK loop (emit, interact with scene and/or reactor, detect, control, emit) can be treated as recurrent. Multi-timescale media and controller dynamics may serve roles analogous to gated recurrence. Slow media may behave like long-lived state, while fast dynamics behave like short-lived state. Detector gating and emission envelopes may act as write and read controls.

\paragraph{Gating as multi-timescale media (non-limiting).}
In some embodiments, gated recurrence corresponds to explicit multi-timescale media and explicit control schedules. A non-limiting mapping is:
\begin{center}
\small
\begin{tabular}{ll}
  \toprule
  \textbf{LSTM component} & \textbf{RK physical analogue} \\
  \midrule
  Cell state $c_t$ & Slow-decay media (phosphors, thermal mass, persistent optical states) \\
  Hidden state $h_t$ & Fast-decay media (fluorescence, transient scattering) \\
  Forget gate $f_t$ & Dissipation schedule $\gamma_{\mathrm{slow}}(\theta)$ \\
  Input gate $i_t$ & Illumination envelope $\mathbf{e}(t)$ modulating write strength \\
  Output gate $o_t$ & Detector gating $g_{\mathrm{det}}(t)$ \\
  \bottomrule
\end{tabular}
\end{center}
In these embodiments, ``gates'' are enacted by physical dynamics and by the protocol and metered operating envelope, rather than by a purely digital hidden-state update.

\subsection{Autoencoders and variational inference (AE, VAE, denoising)}
In \LI operation, the physical channel acts as an encoder from scenes to the \cb. Digital heads act as decoders from the \cb (or derived latents) to inferred scene variables. The effective bottleneck is constrained by physical channel capacity, noise, bandwidth, and detector resolution. In stochastic regimes, the channel may function as a variational encoder, and denoising behaviour may arise from physical averaging and controlled probing.

\subsection{Adversarial structure (GAN-like)}
In \TB operation, a verifier meter may play a discriminator-like role distinguishing genuine physical traces from forged or simulated traces. An adversary that trains a simulator or generator to imitate the \cb plays a generator role. The physical \RK provides genuine samples, yielding an adversarial training and evaluation structure in which physics constrains the feasible set of outputs.

\paragraph{Role mapping (non-limiting).}
\begin{center}
\small
\begin{tabular}{ll}
  \toprule
  \textbf{GAN component} & \textbf{\TB analogue} \\
  \midrule
  Generator $G$ & \RK $\mathcal{R}_\theta$ producing physical traces \\
  Discriminator $D$ & Verifier / meter head $V_\psi$ \\
  Adversary & Eve's spoofing policy $\pi_{\mathrm{Eve}}$ \\
  \bottomrule
\end{tabular}
\end{center}

\paragraph{Minimax objective lens (non-limiting).}
In some embodiments, the evaluation game is written in a minimax form:
\[
  \min_{\pi_{\mathrm{Eve}}} \max_{V_\psi}
  \;\E\bigl[\log V_\psi(\mathrm{real})\bigr]
  + \E\bigl[\log(1 - V_\psi(\mathrm{Eve}))\bigr].
\]
This is an interpretation lens: it highlights that ``the generator'' is the physical process and that the adversary's job is to produce synthetic traces that pass a competent meter under a logged protocol and meter envelope.

\subsection{Diffusion, score-based, and flow models}
In some embodiments, repeated application of a configured \RK serves as a forward corruption or mixing process under controlled schedules, for example by increasing scattering, noise, or dissipation. A learned digital head may implement a reverse process or score-like correction operating on the \cb or a latent representation. Flow matching and consistency models may be implemented by learning transport fields between distributions of \cba traces.

\paragraph{Markov chain form (non-limiting).}
In a non-limiting notation, repeated application can be written as a Markov chain:
\[
  x_{k+1} \sim \mathsf{P}_\theta(\cdot\mid x_k, U_{0:T}^{(k)}),
\]
where $x_k$ is an intermediate state or representation (for example a latent, field state, or scene proxy) and $U_{0:T}^{(k)}$ denotes a step-indexed protocol schedule.
Here $\mathsf{P}_\theta$ is used as shorthand for the induced
transition kernel on representations; $x_k$ plays the role of the
scene input at each stage, so the typing is consistent with the
core definition $\mathsf{P}_\theta(\cdot \mid S, U_{0:T})$ when $S$
is instantiated by the intermediate state.  In some embodiments, a
distinct symbol (for example $T_\theta^{(k)}$) is used when the
stage-indexed transition differs materially from the single-pass
kernel.

\subsection{Reservoir computing and random features}
In some embodiments, the reactor acts as a physical reservoir whose internal dynamics induce a high-dimensional, task-adaptive feature map of the stimulus and control protocol.
For example, the \cb and meter summaries may be treated as a feature vector $\phi_\theta(S,U)$ that depends on stimulus $S$, protocol $U$, and physical parameters $\theta$.

\paragraph{Physical random features and induced kernels (non-limiting).}
In some embodiments, the reactor implements physical random features in the sense that inner products in feature space behave as a kernel on inputs:
\[
  k_\theta\bigl((S,U),(S',U')\bigr)
  = \langle \phi_\theta(S,U),\,\phi_\theta(S',U') \rangle,
\]
where $\langle \cdot,\cdot\rangle$ is any non-limiting inner product or similarity measure on extracted features.
Such kernels may be used for classification, regression, nearest-neighbour matching, or as components of a verisimilitude discriminator.

\paragraph{Linear readouts and small heads (non-limiting).}
In some embodiments, a linear readout or small learned head operates on $\phi_\theta(S,U)$ to produce decisions or predictions.
This matches reservoir computing in which a complex dynamical system provides rich features and only the readout is trained or adapted.

\paragraph{Dual verification and computation on a single physical
substrate (non-limiting).}
In some embodiments, the same physical medium that provides
reservoir-like computational features also provides physically
unclonable verification properties. The microstructure that makes a
scattering medium, nonlinear cavity, or multi-mode waveguide useful as
a computational resource (high-dimensional, device-specific transfer
function) is the same microstructure that makes it useful as a microstructure challenge-response medium
(device-specific, hard-to-clone challenge-response behaviour). In such
embodiments, the distinction between \TB (verification) and
\LI or \RT (computation) is not a hardware
distinction but an objective-function distinction: the same \cb produced by the same physical channel is evaluated for
authenticity (\TB), for scene information (\LI), or for
rendering quality (\RT). This dual use is a direct
consequence of the Markov-kernel abstraction, which treats the physical
channel as a single parameterised object
$\mathsf{P}_\theta(\cdot \mid S, U_{0:T})$ whose output supports
multiple downstream objectives without hardware modification.

\paragraph{Audit compatibility --- reservoir computing (non-limiting).}
In some embodiments, the feature map $\phi_\theta$ is computed from committed atoms and meter summaries so that downstream decisions can be audited by selective opening when required.

\subsection{Residual and skip connections (ResNet-like)}
In some embodiments, optical or physical skip paths split a stimulus into a direct branch and a processed branch and then recombine them, yielding residual-like behaviour.
This can be implemented by beam splitters, waveguides, multiplexed paths, or mixed digital and physical routing.

\paragraph{Residual mixing with controllable bypass (non-limiting).}
A non-limiting mixing form is
\[
  E_{\mathrm{out}}^{(k)}
  = b_k\,E_{\mathrm{in}}^{(k)}
    + (1-b_k)\,\mathcal{R}^{(k)}\!\bigl(E_{\mathrm{in}}^{(k)}\bigr),
\]
where $E_{\mathrm{in}}^{(k)}$ and $E_{\mathrm{out}}^{(k)}$ denote the effective field or representation at stage $k$, $\mathcal{R}^{(k)}$ is a non-limiting transformation implemented physically or digitally, and $b_k\in[0,1]$ controls the bypass fraction.

\paragraph{Stability and expressivity (non-limiting).}
In some embodiments, bypass fractions $b_k$ are used to control stability and mixing, aligning with contractivity and commutativity targets in the Markov-metric view.
In some embodiments, bypass paths also carry timing anchors or calibration signals to improve alignment and robustness.

\subsection{Convolutional and equivariant mappings (CNN-like)}
In some embodiments, when an optical or physical transfer function is approximately shift-invariant over a region of interest, the resulting input--output relationship can be modelled as a convolution-like operator.
This connects to convolutional neural networks (CNNs), equivariant mappings, and designed filter banks.

\paragraph{Convolutional form (non-limiting).}
A non-limiting continuous-space convolution form is
\[
  (f * g)(x) = \int f(\xi)\,g(x-\xi)\,d\xi,
\]
with corresponding discrete forms for sampled fields.
In some embodiments, convolution-like behaviour arises from Fourier-optic propagation, pupil filtering, engineered point-spread functions, or structured illumination combined with fixed readout optics.

\paragraph{Equivariance as an engineering target (non-limiting).}
In some embodiments, the mapping $T$ from stimulus to observation is designed to be equivariant to a transformation group $G$ (for example translations, rotations over a limited range, or planar rigid motions), meaning
\[
  T\bigl(\rho(g)\,f\bigr) \approx \rho'(g)\,T(f),
\]
where $\rho$ and $\rho'$ are representations of $G$ acting on inputs and outputs.
Approximate equivariance may be assessed empirically under meter-bounded operation and used to select optical configurations, scan laws, and aggregation rules.

\paragraph{Group convolution viewpoint (non-limiting).}
In some embodiments, convolution is generalised to group convolutions:
\[
  (f *_{G} \psi)(g) = \int_{h\in G} f(h)\,\psi(h^{-1}g)\,dh,
\]
which supports non-limiting embodiments where optical routing and scan protocols implement filtering over transformation parameters (for example angle or pose), rather than only over spatial translations.

\paragraph{Designed operators and controlled scan protocols (non-limiting).}
In some embodiments, operators are designed by choosing illumination patterns and readout optics so that the induced mapping approximates a desired filter family (for example wavelets, steerable filters, or multi-scale banks).
In some embodiments, controlled scan protocols provide the equivalent of receptive-field traversal and multi-scale sampling while preserving \cba logging and audit compatibility.

\subsection{Energy-based and associative memory (Hopfield and Boltzmann-like)}
In embodiments where reactor media evolve toward stable attractors, the physical dynamics can be interpreted as minimising an energy-like functional.
This connects to energy-based models, Hopfield-like associative memories, and Boltzmann-like sampling interpretations.

\paragraph{Energy and relaxational dynamics (non-limiting).}
A non-limiting abstract form is a state $x$ evolving under a potential $E(x)$:
\[
  \dot{x}(t) = -\nabla E\bigl(x(t)\bigr) + w(t),
\]
where $w(t)$ is a non-limiting noise term (including thermal noise, shot noise, or injected structured noise).
In some embodiments, stable attractors correspond to stored patterns, calibrated configurations, or preferred operating regimes.

\paragraph{Associative retrieval (non-limiting).}
In some embodiments, presenting a partial cue via structured illumination causes dynamics to settle into a compatible attractor, yielding associative recall.
The resulting \cba summaries act as readouts of the retrieved state and may be compared to predicted responses for probing or verification.

\paragraph{Training alignment (non-limiting).}
In some embodiments, energy-based losses or contrastive objectives are used so that digital training aligns with the physical relaxational dynamics observed in the reactor under meter-bounded operation.

\subsection{Attention and transformer-like interpretations}
In some embodiments, associative retrieval in a reactor serves a role analogous to attention, where the system forms content-dependent weights over stored or reachable patterns and then produces a weighted retrieval.

\paragraph{Attention-style weighting (non-limiting).}
A non-limiting attention form may be written as
\[
  a_{ij} = \mathrm{softmax}_j\!\left(\frac{\langle q_i, k_j\rangle}{\sqrt{d}}\right),
  \qquad
  z_i = \sum_j a_{ij}\,v_j,
\]
where $q_i$ is a query-like representation, $k_j$ are key-like representations, $v_j$ are value-like representations, and $d$ is a scaling dimension.
In some embodiments, $q_i,k_j,v_j$ are derived from \cba summaries, meter-conditioned features, or reactor-internal states.

\paragraph{Physical correlators and holographic storage (non-limiting).}
In some embodiments, optical correlation, interference, or convolution implements parts of the dot-product or similarity computations that underlie attention-like weighting.
In some embodiments, holographic or interference-based media provide associative storage such that retrieval behaves as content-addressable memory, which can be interpreted as a continuous Hopfield-like mechanism.

\paragraph{Audit and policy constraints (non-limiting).}
In some embodiments, tokenisation, conditioning, and normalisation functions are constrained by meter envelopes and commitment policies, so that transformer-like interpretations remain compatible with selective opening and trust-weight updates.

\subsection{State-space models and continuous-time dynamics}
In some embodiments, the \RK admits a controlled state-space description, with observations arising from the measurement map and the \cba channel.
In the discrete-time formulations below, $u_t := u(t)$ denotes the per-step control of Section~\ref{sec:definitions}, and $y_t$ denotes a scalar or vector observation consistent with $\mathbf{y}_t$.

\paragraph{Discrete-time state-space form (non-limiting).}
A non-limiting form is
\[
  x_{t+1} = f_\theta(x_t, u_t) + w_t,
  \qquad
  y_t = g_\theta(x_t, u_t) + v_t,
\]
where $x_t$ is an extended state, $u_t$ is a control input (including scan law, projection schedule, gains, and alignment commands), $y_t$ is an observation or summary (including the \cb features and meters), and $w_t,v_t$ are non-limiting process and observation noise terms.

\paragraph{Continuous-time form and Neural ODE interpretation (non-limiting).}
In some embodiments, continuous-time dynamics are used:
\[
  \dot{x}(t) = f_\theta\bigl(x(t), u(t), t\bigr),
  \qquad
  y(t) = g_\theta\bigl(x(t), u(t), t\bigr).
\]
When $f_\theta$ is parameterised by a learned function class, this aligns with Neural ODE or continuous-time state-space modelling viewpoints, with meter summaries providing constraints and validation signals.

\paragraph{Control and estimation (non-limiting).}
In some embodiments, standard estimation and control tools (filtering, smoothing, observability and controllability analyses, robust control) are applied, subject to the meter-bounded operating envelope and verification-gated action policies described elsewhere in this document.

\subsection{World models}
In some embodiments, the system is described using a world-model decomposition in which an encoder produces a latent state, a dynamics model predicts latent evolution under actions, and a decoder (or predictor) produces expected observations.
This framing supports planning, control, and verification when the primary observation stream is derived from the \cb and meter summaries.

\paragraph{Latent-state decomposition (non-limiting).}
A non-limiting form is
\[
  z_t = \mathrm{Enc}_\psi(o_t), \qquad
  z_{t+1} = F_\psi(z_t, a_t) + w_t, \qquad
  \widehat{o}_{t+1} = \mathrm{Dec}_\psi(z_{t+1}),
\]
where $o_t$ is an RK-derived observation or summary (for example features extracted from the \cb together with meter vectors and protocol digests), $a_t$ is an action (including RK control components), and $w_t$ is a non-limiting disturbance term.

\paragraph{Meter-bounded and commitment-aware world models (non-limiting).}
In some embodiments, world-model updates and planning are constrained by meter envelopes and by audit policies.
For example, training or adaptation may be conditioned on meter vectors and protocol digests, and may be regularised to preserve verisimilitude and projection-hardness scores.
In some embodiments, when commitments exist without openings, planning operates in a provisional mode and upgrades to a verified mode after selective opening confirms consistency.

\paragraph{Planning and control (non-limiting).}
In some embodiments, planning uses the latent dynamics model for model-predictive control, policy optimisation, or search over action sequences, with costs that include both task costs and evidence-quality terms.
This is compatible with the agent integration section and with verification-gated action policies described elsewhere in this document.

\subsection{Neural radiance fields and differentiable rendering (NeRF-like)}
In some embodiments, \LI and \RT operation are interpreted using scene-representation lenses similar to neural radiance fields (NeRFs), differentiable rendering, and light-field models, with the key distinction that the evidence is produced by a physical swept kernel rather than by purely digital rendering.

\paragraph{NeRF-like representation lens (non-limiting).}
NeRF-style models represent a scene as a continuous field, for example
\[
  F_\psi:(\mathbf{x},\mathbf{d})\mapsto(\mathbf{c},\sigma),
\]
mapping spatial position $\mathbf{x}$ and viewing direction $\mathbf{d}$ to colour $\mathbf{c}$ and density $\sigma$ (or other radiometric quantities). In some embodiments, a \RK operating in \LI mode learns a representation that is analogous in purpose: it predicts or reconstructs scene variables from the \cb under a family of logged control protocols, and it is validated by executing additional physical probes and comparing predicted versus observed \cba windows under the same protocol digests.

\paragraph{Differentiable rendering lens (non-limiting).}
In some embodiments, \RT operation is treated as a physically instantiated inverse-rendering or projection-mapping loop: the controller selects emissions, observes responses, and updates projection parameters to reduce a target appearance loss while remaining within meter envelopes. Gradients may be obtained by differentiable surrogates, finite differences, or policy-gradient methods, and the resulting actuation remains auditable because protocol digests, meter summaries, and committed windows record what was emitted and what was observed.

\paragraph{Explicit scene primitives (non-limiting).}
In some embodiments, scene representations are explicit (for example point clouds, surfels, Gaussian primitives, meshes, or hybrid implicit--explicit forms) and probe protocols are designed to efficiently constrain those primitives (pose, depth, reflectance, and dynamics) under logged scan schedules.

\paragraph{Light-field and ray-space interpretations (non-limiting).}
In some embodiments, a \RK is interpreted as sampling a ray-space or light-field parameterisation, where different control paths correspond to different ray bundles (space, angle, wavelength, polarisation, time-of-flight bins), and the learned model predicts or constrains the induced \cba distribution under those ray samples.

\subsection{Backpropagation through physics and surrogate gradients (non-limiting)}
In some embodiments, training and tuning use gradient-like updates even when physical dynamics are not directly differentiable.

\paragraph{Surrogate backpropagation (non-limiting).}
In some embodiments, a differentiable surrogate $\widehat{\mathsf{P}}_\theta$ is trained on logged input--output pairs and is used to compute gradients for losses defined on physical artifacts (for example hardness meters, verisimilitude meters, calibration penalties, or reconstruction error). Updates are then applied to admissible parameters (digital and/or analogue) subject to actuator projections and safety envelopes.

\paragraph{Perturbation gradients (non-limiting).}
In some embodiments, gradients are approximated by finite differences or stochastic perturbations (for example SPSA), where the controller perturbs control paths or parameters, measures the change in a scalar objective, and applies projected updates under variance-control practices and meter gating.

\paragraph{Adjoint and time-reversal analogues (optional, non-limiting).}
In some embodiments, continuous-time models admit adjoint-style gradient calculations, and reversible or approximately reversible optical subsystems admit time-reversal analogues (for example phase conjugation) that support efficient sensitivity estimation. These are non-limiting implementation patterns and may be combined with digital surrogates and recorded forward trajectories.

\paragraph{Reinforcement learning as a fallback (non-limiting).}
In some embodiments, reinforcement learning treats the physical device as a black-box environment: rewards are derived from meters and task losses, and a policy is updated to select scan protocols, emissions, and tuning actions that improve objectives under logged regimes.

\paragraph{Modern reinforcement-learning algorithm families (non-limiting).}
In some embodiments, the policy-update step uses any of the modern
reinforcement-learning algorithm families, including without limitation
proximal policy optimisation (PPO, Schulman et al. 2017),
trust-region policy optimisation (TRPO), advantage actor-critic
methods (A2C, A3C, and synchronous and asynchronous variants),
soft actor-critic (SAC) and other off-policy
entropy-regularised methods, deep deterministic policy gradients
(DDPG), twin-delayed DDPG (TD3), Q-learning and its deep variants
(DQN, Double DQN, Dueling DQN, Rainbow, prioritised experience
replay), distributional reinforcement learning methods (C51, QR-DQN,
IQN), IMPALA and other distributed actor-learner architectures
(Ape-X, R2D2, SEED-RL), model-based reinforcement learning methods
including MuZero-style learned-model-with-search and Dreamer-style
latent-dynamics-model RL (Hafner et al.), offline reinforcement
learning methods (CQL, IQL, BCQ, AWAC, decision transformer),
group-relative policy optimisation (GRPO) and related verifier-based
post-training methods used for reasoning agents, and
goal-conditioned and hierarchical reinforcement learning methods
(HER, options framework, FeUdal Networks, HIRO). The apparatus is
compatible with any such method applied to scan-policy training,
controller training, or governed-agent training, including future
methods within the same algorithmic envelope.

\subsection{Summary (neural-architecture substrate)}
This section provides a non-limiting mapping between \RK physical mechanisms and common neural-architecture abstractions.
The purpose is explanatory and enablement-oriented: it shows how a single physical substrate (configured by scan protocols, structured illumination, and meter-bounded operation) can instantiate behaviours analogous to widely used computational motifs.

{\small
\begin{longtable}{@{}p{0.27\textwidth}p{0.34\textwidth}p{0.34\textwidth}@{}}
  \toprule
  \textbf{Architecture} & \textbf{RK instantiation} & \textbf{Key physical mechanism} \\
  \midrule
\endfirsthead
  \toprule
  \textbf{Architecture} & \textbf{RK instantiation} & \textbf{Key physical mechanism} \\
  \midrule
\endhead
  \bottomrule
\endlastfoot
  RNN / LSTM / GRU & Core feedback loop with multi-timescale media and gating schedules & Multi-timescale dissipation and metered write/read controls \\
  AE / VAE / denoising & \LI regime with a physical encoder and digital decoder & Channel capacity as bottleneck; noise and controlled probing \\
  GAN-like & \TB meter as discriminator; adversary as generator & Hard-to-forge physical signatures under bounded protocols \\
  Diffusion / score / flow & Physical forward mixing plus learned reverse/score head & Shot noise, scattering, and scheduled dissipation as corruption \\
  Reservoir / random features & Reactor dynamics plus small readout head & High-dimensional mixing; implicit kernel on $(S,U)$ features \\
  ResNet / skip connections & Split/recombine paths with controllable bypass & Beam splitters, multiplexed paths, reference arms \\
  CNN / equivariant & Shift-invariant or group-structured optical operators & PSFs, diffraction, and structured scan-addressed transforms \\
  Energy-based / Hopfield & Relaxation to attractors under structured cues & Physical energy landscape and content-addressable retrieval \\
  Transformers / attention-like & Correlation and associative retrieval over stored patterns & Similarity weighting via correlation/interference and multiplexed channels \\
  State-space / Neural ODE & Controlled continuous-time dynamics with observation map & Physical dynamics with filtering/estimation on the \cb \\
  NeRF / differentiable rendering & \LI/\RT with protocol-conditioned scene models and physical validation & Active probing and meter-bounded physical constraints on predictions \\
  Backprop / surrogate gradients & Surrogate-guided tuning and perturbation-based optimisation & Differentiable twins, finite differences, and metered update envelopes \\
\end{longtable}
}

In some embodiments:
(i) reservoir and random-feature viewpoints arise when \cba and meter-conditioned summaries define a high-dimensional feature map and induced kernel;
(ii) residual and skip-path viewpoints arise when bypass fractions mix direct and processed paths in a controllable manner;
(iii) convolutional and equivariant viewpoints arise when optical transfer functions approximate shift-invariant or group-structured operators under controlled scans;
(iv) energy-based and associative-memory viewpoints arise when reactor media relax toward attractors under structured cues, optionally with noise and injected perturbations;
(v) attention-like viewpoints arise when correlation and associative retrieval implement content-dependent weighting over stored or reachable patterns; and
(vi) state-space and continuous-time viewpoints arise when the system is modelled as controlled dynamics with observations derived from the \cb and meter summaries.

These mappings do not assert that the physical system is identical to any specific neural architecture.
Rather, they provide non-limiting interpretations and design analogies that can guide implementation choices, training objectives, and verification policies while preserving \cba logging, selective disclosure, and the Markov-kernel framing used throughout this document.

% ======================================================================
\section{Information-Theoretic Primitives}
\label{sec:info-theory}
% ======================================================================

This section describes information-theoretic viewpoints and quantities that may be used to characterise \RK systems. These quantities are non-limiting and may be used as design and evaluation tools. They do not replace physical grounding or meter-based evaluation. They provide vocabulary for limits and trade-offs across \TB, \LI, and \RT regimes.  This section establishes formal relationships with information theory as interpretive vocabulary for limits and trade-offs across regimes; it does not constrain the core definitions or embodiments.

\begin{remark}[Notation]
Throughout this section, the symbol $\Xi$ denotes a task-relevant
random variable (stimulus, secret, or label) drawn from a
design-time distribution.  This is distinct from the scene tuple
$S = (S^{\mathrm{scene}}, S^{\mathrm{reactor}}_\theta)$ of
Section~\ref{sec:definitions}: the scene tuple is a physical
configuration; the task variable $\Xi$ is a random variable over
which information-theoretic quantities are optimised. In many
applications the task variable is derived from or indexed by the
scene (for example a label identifying which scene is present), but
the two objects are formally distinct.
\end{remark}

\subsection{Reality Kernel as a noisy channel}
A \RK (RK) module can be modelled as inducing a conditional distribution over the recorded \cb given an input stimulus (or secret) and a control protocol.
Let $\Xi$ denote a stimulus, secret, or task-relevant label drawn from a design-time distribution $\mathsf{P}_{\Xi}$.
Let $U_{0:T}$ denote a control protocol (projection sequence, scan law, actuator schedule, gains, alignment parameters, and any other commanded inputs), and let $C_{0:T}$ denote the resulting \cba record.

The RK induces a channel family parameterised by physical and control parameters $\theta$:
\[
C_{0:T} \sim \mathsf{P}_\theta\!\left(\cdot\mid \Xi, U_{0:T}\right).
\]
When $U_{0:T}$ is treated as given (known to the verifier, logged, or fixed by policy), the induced information transfer from $\Xi$ into $C_{0:T}$ can be quantified by the conditional mutual information
\[
I_\theta\!\left(\Xi;C_{0:T}\mid U_{0:T}\right).
\]
In some embodiments, the protocol itself is randomised (for example, via structured noise or scan randomisation). In such cases $U_{0:T}$ is a random variable, and the same quantity is interpreted as mutual information conditioned on the realised protocol and meter state.

\paragraph{Operational role of $\mathsf{P}_\theta$ (non-limiting).}
The kernel $\mathsf{P}_\theta : \mathcal{S} \times \mathcal{U}^{T+1} \to \mathcal{P}(\mathcal{C})$
serves as a conceptual unification device: it provides a common
probabilistic vocabulary for the three operating regimes and grounds
the information-theoretic quantities that follow.  All operational
computation uses specific statistics derived from the \cb under
declared estimators --- the meter vector, TE estimates, hardness
indices, GP posterior quantities --- rather than the full distribution
$\mathsf{P}_\theta$ itself, which is not directly estimable from
finite data as an element of the infinite-dimensional measure space
$\mathcal{P}(\mathcal{C})$.  In some embodiments, $\mathcal{P}(\mathcal{C})$
is equipped with the weak topology induced by bounded continuous
test functionals of the trajectory (for example, finite-lag moment
functionals of the \cb), and identifiability of $\theta$ from the
induced distribution is understood as a meter-conditioned empirical
target: under protocol family $\mathcal{U}$ and meter envelope
$\mathcal{M}_{\mathrm{accept}}$, two configurations $\theta \neq \theta'$
are distinguishable with probability at least $1-\delta$ given $n(\delta, \varepsilon)$
trajectories, where $n(\delta, \varepsilon)$ is estimated empirically
for the declared estimator family.  The $\varepsilon$-nice kernel
characterisation (Section~\ref{sec:definitions}) is the approximation
regime within which this holds.

\paragraph{Data processing inequality (post-processing cannot add information).}
If a downstream algorithm (including compression, feature extraction, hashing, learned encoders, or any deterministic or stochastic post-processing) produces $Z = g(C_{0:T},U_{0:T})$, then a standard consequence of the data processing inequality is
\[
I_\theta\!\left(\Xi;Z\mid U_{0:T}\right) \le I_\theta\!\left(\Xi;C_{0:T}\mid U_{0:T}\right).
\]
This formalises the design principle that any verification, reconstruction, or attestation pipeline is ultimately limited by the information carried in the recorded \cb under the chosen physical regime and protocol.

\paragraph{Memory, channel dependence, and the DPI (non-limiting).}
The standard data processing inequality above assumes a Markov chain
$\Xi \to C_{0:T} \to Z$.  When the \cb exhibits temporal dependence
(for example due to reactor persistence within the declared memory
depth $\tau_{\mathrm{mem}}$; see the Markov kernel remark in
Section~\ref{sec:definitions}), the inequality continues to hold in
the form stated, since $Z = g(C_{0:T}, U_{0:T})$ is still a
deterministic function of the channel output.  However, the mutual
information $I_\theta(\Xi; C_{0:T} \mid U_{0:T})$ itself may be
inflated by temporal dependencies if $\Xi$ and the channel memory are
not independent under the protocol; in preferred embodiments, the
bundle length $T$ is chosen large enough to average over
within-memory-depth correlations, and the effective information
$I_\theta(\Xi; C_{0:T} \mid U_{0:T})$ is computed with an estimator
that accounts for the declared autocorrelation structure (for example
a block-structured mutual-information estimator with block length
$\tau_{\mathrm{mem}}$).  The resulting bound is referred to as the
\emph{memory-adjusted DPI bound} and is the operative form used in
verification and reconstruction performance statements.

\paragraph{Multi-query DPI composition for adaptive adversaries
(non-limiting).}
The DPI bound stated above applies within a single fixed-horizon run
(non-adaptive setting), where $\hat{\Xi}$ is computed from a single
bundle $C_{0:T}$.  For an adaptive adversary who submits $Q$ queries
sequentially, using each response to choose the next challenge, the
aggregate information gain over $Q$ runs is bounded by a composition
of the per-run DPI bounds.  In some embodiments, this composition is
stated as follows: under the declared query throttle rate and selective
disclosure policy, and under the assumption that the run-to-run mutual
information
$I(C_{0:T}^{(j)}; C_{0:T}^{(j+1)} \mid U_{0:T}^{(j)}, U_{0:T}^{(j+1)})
\leq \rho_{\mathrm{run}}$ is bounded by a declared inter-run
correlation $\rho_{\mathrm{run}}$ (estimable from the AR(1) drift
model and the protocol randomisation structure), the aggregate
information gain satisfies
$I(\Xi; \{C_{0:T}^{(j)}\}_{j=1}^Q \mid \{U_{0:T}^{(j)}\}) \leq
Q \cdot I_\theta(\Xi; C_{0:T} \mid U_{0:T}) + \Delta_{\mathrm{Q,run}}(Q)$,
where $\Delta_{\mathrm{Q,run}}(Q)$ is a nonnegative validation-envelope term
separately upper-bounded under the declared validation regime, with empirical
estimates used to support the declared bound in the operative protocol
regime. In the operative protocol regime, the validation target is that
both pairwise run-to-run dependence and higher-order transcript leakage
across the full collection of runs are empirically driven toward zero;
the operative protocol declares and verifies these conditions empirically
rather than deriving $\Delta_{\mathrm{Q,run}}$ from a closed-form expression
in $\rho_{\mathrm{run}}$. Pairwise run-to-run correlation alone does not
in general control higher-order synergistic information across the full
collection of runs.
The assumption of approximate independence across runs
--- $\rho_{\mathrm{run}} \approx 0$ --- is a declared protocol
requirement, verified by checking that the run-to-run correlation in
\cb statistics does not exceed $\rho_{\mathrm{run}}$
under the declared protocol randomisation; this requirement and its
empirical verification are recorded in the protocol digest.

\paragraph{Transfer entropy: coupling diagnostic vs.\ hardness bound
(non-limiting).}
Transfer entropy $\TE_{S \to R}$ and $\TE_{R \to S}$ are used in two
distinct roles in this specification, and these roles should not be
conflated.  As a \emph{coupling diagnostic} (Yoked operation;
Section~\ref{sec:yoked-convergence}), TE measures the degree to which
the scene and reactor exchange information in both directions; the
relevant property is balance and magnitude, and the primary concern
is statistical power and common-cause confounding.  As a
\emph{component of Yoked-regime security} (Yoked \TB), high
TE indicates that the scene's response is informationally entangled
with the reactor's history, making replay or surrogate attacks harder;
the relevant security parameter is the operative mutual information $I_\theta(S_{0:T}; R_{0:T})$ between scene and reactor under the disclosed coupling and input distribution; the memory-adjusted DPI bound provides an upper bound on the information that any downstream estimator can recover about $S_{0:T}$ from $R_{0:T}$. No converse-style channel-capacity claim (supremum over input distributions) is asserted in this disclosure.  In some embodiments, the hardness contributions from the
Yoked coupling are quantified separately from the hardness
contributions from physical microstructure
(Section~\ref{sec:hardness-index}), and the two estimates are
combined with declared weights in the hardness index; this separation
ensures that a degradation in coupling strength (for example from a
scene change) does not silently collapse the hardness claim without
triggering a meter alarm.

Because TE is used for both the coupling diagnostic and the security
argument, two distinct estimator families with separately declared
calibration parameters are required: one optimised for the
\emph{difference} $\TE_{S \to R} - \TE_{R \to S}$ (bias-cancellation
property of shared embedding is the relevant criterion), and one
optimised for the \emph{absolute} entropy rate and coupling capacity
(bias in the absolute value is the relevant criterion).  Both
estimator families and their protocol parameters (embedding dimension,
lag, window length, $k$-NN count for KSG estimators) are declared
separately in the protocol digest; using a single estimator family for
both purposes is a non-conforming embodiment.

\paragraph{Yoked convergence criterion: necessary vs.\ sufficient
(non-limiting).}
The Yoked convergence criterion (CLE, TE balance, and estimator SNR
gate co-conditions) is a necessary condition for Yoked operation but
is not alone sufficient for the security properties claimed for Yoked
\TB.  The TE criterion can in principle be satisfied by a
signal whose TE value matches the genuine bundle but whose
higher-order statistical structure differs, in a manner analogous to a
hash pre-image attack.  In some embodiments, a higher-order structure
check (for example, the honeypot verification embodiment described in
Section~\ref{sec:honeypot}) is integrated as a mandatory co-condition
of the Yoked convergence criterion: the joint distribution of
consecutive bundle segments must exceed a declared divergence
threshold from the null model of uncoupled processes, computed using
a declared higher-order estimator (e.g., an energy statistic or a
neural mutual-information estimate) committed to the protocol digest.
In embodiments where the honeypot check is not integrated, the
security argument for Yoked \TB explicitly depends on the
higher-order structure checks described in the honeypot embodiment
as a parallel monitoring process.

\subsection{Mutual information and capacity (design-time bounds)}
For a fixed horizon $T$ and RK parameterisation $\theta$, a non-limiting notion of design-time capacity is
\[
C_T(\theta) = \sup_{\mathsf{P}_{\Xi}} I_\theta\!\left(\Xi; C_{0:T}\mid U_{0:T}\right),
\]
where the supremum ranges over admissible stimulus or secret distributions $\mathsf{P}_{\Xi}$ (for example, bounded-energy stimuli, bounded alphabet, or application-specific label priors).
In some embodiments, the admissible set also constrains the control protocol distribution (for example, $\mathsf{P}_{U_{0:T}}$ restricted by scan speed, eye-safety, thermal limits, or meter stability), yielding a constrained capacity
\[
C_T(\theta;\mathcal{U}) = \sup_{\mathsf{P}_{\Xi},\,\mathsf{P}_{U_{0:T}}\in\mathcal{U}} I_\theta\!\left(\Xi; C_{0:T}\mid U_{0:T}\right).
\]

\paragraph{Cross-regime constraint conjecture: hardness and sensing
capacity as complementary quantities (non-limiting, research direction).}
The single-kernel formulation $\mathsf{P}_\theta$ unifies \TB
(authentication), \LI (sensing), and \RT (rendering)
over the same \cb.  A deep implication of this unification, not yet
proved, is that the digital hardness index $k^*_{\mathrm{dig}}(\theta;
\varepsilon)$ in \TB mode and the sensing capacity
$C_T(\theta; \mathcal{U})$ in \LI mode are not independent
quantities: both are bounded above by the same physical resource,
namely the optical information capacity of the channel, which scales
as $N_{\mathrm{modes}} \cdot \log(1 + \mathrm{SNR})$.  The following
is stated as a research conjecture.

\medskip
\noindent\textit{Cross-Regime Constraint Conjecture.}
There exists a decreasing function $h(\cdot)$ such that, for a device
with digital hardness index $k^*_{\mathrm{dig}}(\theta; \varepsilon)$
and sensing capacity $C_T(\theta; \mathcal{U})$ operating under the
same physical channel $\mathsf{P}_\theta$,
\[
  C_T(\theta; \mathcal{U}) \;\leq\; h\!\left(k^*_{\mathrm{dig}}(\theta; \varepsilon)\right),
\]
where $h$ is a decreasing function of the hardness index.
\medskip

Informally: a device that is highly secure in \TB mode (high
$k^*_{\mathrm{dig}}$, hard to predict from any model) is a constrained
sensing device (bounded $C_T$, because the same physical complexity
that prevents modelling also limits extractable scene information).
The argument sketch: (i) $k^*_{\mathrm{dig}}$ grows with the channel's
resistance to statistical characterisation, which is related to
$I_\theta(\theta; C_{0:T})$ via the provisional entropy-rate bound;
(ii) $C_T = \sup I_\theta(\Xi; C_{0:T})$ is also bounded by
$N_{\mathrm{modes}} \cdot \log(1+\mathrm{SNR})$ from optical
information theory; (iii) if both are bounded by the same physical
quantity, an implicit trade-off exists.  The precise form of the
trade-off depends on the relationship between $k^*_{\mathrm{dig}}$
and $I_\theta(\theta; C_{0:T})$, which is established only
provisionally (Section~\ref{sec:info-theory}).  In some
embodiments, this trade-off is treated as an empirical hypothesis:
devices are characterised on both axes $(k^*_{\mathrm{dig}},
C_T)$ across a range of protocol configurations, and the resulting
cloud of operating points is examined for the predicted negative
correlation.  A confirmed negative correlation, stable across device
families and protocol variations, would constitute empirical evidence
for the conjecture and would establish a new design principle:
\emph{hardness and sensing capacity are complementary resources of
the same physical channel, subject to a conservation-like constraint.}

\paragraph{Limager objective (information maximisation, non-limiting).}
In \LI-style operation, the task variable may be a latent scene description $\xi$ (geometry, materials, spectral parameters, pose, or a learned representation), where $\xi$ is distinguished from the extended state $X_t$ of Section~\ref{sec:definitions} (which includes scene, reactor, media, and controller state; $\xi$ is typically a component or projection of $X_t$). A non-limiting objective is to choose protocol parameters to maximise information gained about $\xi$:
\[
\max_{\mathsf{P}_{U_{0:T}}\in\mathcal{U}} I_\theta\!\left(\xi; C_{0:T}\mid U_{0:T}\right),
\]
subject to cost, safety, and stability constraints captured by $\mathcal{U}$ and by meter thresholds. Other objectives (for example, reconstruction accuracy, classification accuracy, or downstream task performance) can be used in addition to or instead of information-theoretic criteria.

\paragraph{Variational mutual-information estimation and information bottlenecks (non-limiting).}
In some embodiments, mutual-information quantities such as $I_\theta(\Xi;C_{0:T}\mid U_{0:T})$ or leakage terms $I(K;C_{0:T}\mid U_{0:T})$ (for a secret $K$) are estimated or bounded using variational objectives. A learned critic may be trained with contrastive losses (InfoNCE-style as a non-limiting example) or mutual-information neural estimation (MINE-style as a non-limiting example) on pairs of stimuli and \cba traces, yielding lower bounds on retained information and upper bounds on unwanted leakage under declared protocol and meter regimes. These bounds may be used at design time to characterise a device or at run time to adapt control policies.
In some embodiments, an explicit information-bottleneck objective is imposed on intermediate latents: a digital encoder is trained to minimise $I(Z;C_{0:T})$ subject to preserving a task or verification objective, thereby trading off utility against invariance and hardness.

\paragraph{Digital hardness index for inference from the \cb (non-limiting).}
To make empirical hardness statements operational, consider a family of attacker models (or hypothesis classes) $\{\mathcal{H}_k\}_{k\ge 1}$ indexed by a resource scale $k$ (for example, parameter count, available FLOPs, wall-clock time, query budget, or training data budget).
Given a predictor $f\in\mathcal{H}_k$ that observes $(C_{0:T},U_{0:T})$ and outputs $\widehat{\Xi}=f(C_{0:T},U_{0:T})$, define a balanced classification error for binary $\Xi\in\{0,1\}$ as
\[
\mathrm{err}_{\mathrm{bal}}(f;\theta)
= \tfrac{1}{2}\Pr_\theta\!\left[f(C_{0:T},U_{0:T})=0 \mid \Xi=1\right]
+ \tfrac{1}{2}\Pr_\theta\!\left[f(C_{0:T},U_{0:T})=1 \mid \Xi=0\right].
\]
For multiclass or structured $\Xi$, this can generalise to an average per-class error or an application-specific loss (for example, label accuracy, key-bit prediction error, or another task-specific misclassification loss).

Define the best achievable error at scale $k$ as
\[
\mathrm{err}^*_k(\theta) = \inf_{f\in\mathcal{H}_k}\ \mathrm{err}_{\mathrm{bal}}(f;\theta),
\]
and define the digital hardness index as the minimal scale required to drive error below a target threshold $\varepsilon$:
\[
k^*_{\mathrm{dig}}(\theta;\varepsilon)=\min\left\{k:\ \mathrm{err}^*_k(\theta)\le \varepsilon\right\}.
\]
In some embodiments, $k^*_{\mathrm{dig}}(\theta;\varepsilon)$ is reported alongside meter summaries (noise, drift, alignment, and stability) so that hardness claims are explicitly conditioned on an observed operating envelope and on the chosen attacker model family.  The meter partition (Section~\ref{sec:security-theory}) restricts hardness computation to evaluation meters $\mathbf{m}^{\mathrm{eval}}$; acceptance meters are excluded from the attacker's discriminative toolkit and from the fidelity criterion against which hardness is assessed.

\paragraph{Analogue hardness and cross-reactor emulation (non-limiting).}
In some embodiments, an attacker attempts to emulate the physical channel using other physical systems, for example by using a different reactor, a coupled reactor network, or a hybrid physical and digital surrogate that maps protocol inputs to output records.
This motivates an analogue hardness notion that complements digital hardness.

Let $R_\theta$ denote a target reactor (or device) operating under a protocol distribution and meter-bounded envelope, and let $\widetilde{R}$ denote an emulator constructed from physical resources (for example another reactor $R_{\theta'}$, transducers, calibration procedures, and optional learned pre- and post-processing).
Let $\mathcal{T}$ denote a transcript for a run, for example
\[
  \mathcal{T} = (\mathrm{rid}, \Pi, \mathbf{m}, \mathrm{com}, \text{opened atoms and proofs as applicable}),
\]
where $\Pi$ is a protocol digest, $\mathbf{m}$ are meter summaries, and $\mathrm{com}$ denotes commitment references.

A non-limiting analogue distinguishability advantage is parameterised
over two separate resource scales: the emulator/attacker construction
scale $m$, governing the physical resources available to the emulator
$\widetilde{R}$, and the verifier/distinguisher scale $r$, governing
the resources available to the distinguishing or verification
procedure $\mathcal{A}$:
\[
  \mathrm{Adv}^{\mathrm{ana}}_{m, r}(\theta)
  = \inf_{\widetilde{R} \in \mathcal{E}_m}
    \sup_{\mathcal{A}\in\mathcal{D}_r}
    \left| \Pr[\mathcal{A}(\mathcal{T})=1 \mid \mathcal{T}\sim R_\theta]
         - \Pr[\mathcal{A}(\mathcal{T})=1 \mid \mathcal{T}\sim \widetilde{R}_m] \right|,
\]
where $\mathcal{E}_m$ is a non-limiting class of emulators
constructed at scale $m$, and $\mathcal{D}_r$ is a non-limiting class
of distinguishers or verification procedures available at scale $r$
(for example bounded witnesses, bounded openings, bounded
cross-prediction checks, and bounded compute). The infimum-supremum
structure makes explicit that the emulator chooses its construction to
minimise distinguishability against the worst-case verifier in the
declared verifier class. The two-scale form is monotone under
declared regularity: increasing $m$ at fixed $r$ does not increase
$\mathrm{Adv}^{\mathrm{ana}}_{m, r}$, and increasing $r$ at fixed $m$
does not decrease it.

The analogue hardness index is defined for a declared verifier
resource budget $r$ as the minimal emulator scale required to reduce
distinguishability below a threshold $\varepsilon$:
\[
  m^*_{\mathrm{ana}}(\theta; \varepsilon, r)
  = \min\{m:\ \mathrm{Adv}^{\mathrm{ana}}_{m, r}(\theta) \le \varepsilon\}.
\]
In some embodiments, the verifier scale $r$ is committed to the
protocol digest alongside the threshold $\varepsilon$, and the
analogue hardness index is reported as the function
$m^*_{\mathrm{ana}}(\theta; \varepsilon, r)$ rather than as a single
scalar; the two-dimensional surface
$(m, r) \mapsto \mathrm{Adv}^{\mathrm{ana}}_{m, r}(\theta)$ may also be
reported for cross-device comparison and for adaptive-attack analysis.

In some embodiments, $m$ is interpreted as a budget over physical degrees of freedom (for example number of coupled subsystems, stability and calibration effort, bandwidth of transducers, achievable timing alignment, or access to reactor--reactor coupling), and $\widetilde{R}_m$ includes emulation via another reactor $R_{\theta'}$ together with protocol-dependent pre-distortion and post-processing.
In some embodiments, analogue hardness is strengthened by protocol randomisation, serial composition of challenges, and two-seed selective opening that limits anticipatory tailoring.
In some embodiments, network witnesses and reactor--reactor coupling edges are used to increase the difficulty of emulation, since the emulator would need to reproduce not only local records but also cross-device compatibility under shared events and coupling.

In some embodiments, hardness scales with serial composition: multiple modules are physically concatenated, each contributing independent microstructure and noise, increasing analogue hardness without shared secrets or joint training, as the joint channel statistics become increasingly harder to emulate, often faster than linearly in practice.

\paragraph{Security oracle game and meter-extraction audit
(non-limiting).}
To make the digital hardness index a falsifiable and
attack-stable security parameter, the attacker access model must be
stated explicitly.  In some embodiments, the following oracle model
is committed to the protocol digest as part of the hardness
evaluation: the adversary is given (i)~white-box access to the
predictive model $\mathcal{Q}$ and the declared device representation
$\mathbf{d}_i$; (ii)~black-box query access to the physical device,
where each query submits a challenge $X$ and receives the response
$Y_i(X)$ (no acceptance or rejection decision is returned by the
device oracle in the default model); and (iii)~a query budget $Q$
declared in the protocol digest.  The adversary wins if they produce
a bundle accepted by the verification procedure that was not
generated by the genuine device.  In some embodiments, a stricter
oracle variant is also evaluated: the adversary additionally receives
acceptance/rejection feedback from the acceptance meter after each
submitted bundle (the ``acceptance-oracle'' variant).  Because
acceptance meters and evaluation meters are statistically correlated
(Section~\ref{sec:security-theory}), the acceptance-oracle variant
provides indirect information about the evaluation meter boundary;
the residual attack advantage from this indirect channel is reported
as part of the hardness evaluation for the declared
acceptance-evaluation meter correlation $\rho_{\mathrm{acc-eval}}$.

In some embodiments, a meter-extraction red-team evaluation is
conducted as part of Eve-mode benchmarking: an attacker family is
given $Q$ adaptive acceptance-oracle queries and attempts to learn
the acceptance meter boundary well enough to construct bundles that
pass verification with probability exceeding a declared threshold.
Success probability as a function of $Q$ is reported in the protocol
digest alongside $k^*_{\mathrm{dig}}$; the query throttle rate is
set so that $Q$ remains below the $Q^*$ at which the success
probability crosses a policy-declared tolerance.

\paragraph{Full hardness curves in protocol digest (non-limiting).}
\label{sec:hardness-curves-and-monotonicity}
The scalar index $k^*_{\mathrm{dig}}(\theta;\varepsilon)$ and the
scalar analogue index $m^*_{\mathrm{ana}}(\theta;\varepsilon, r)$ are
summary statistics of the full hardness curves
$\mathrm{err}^*_k(\theta)$ (error vs.\ model capacity) and
$\mathrm{Adv}^{\mathrm{ana}}_{m, r}(\theta)$ (distinguishing advantage
vs.\ physical scale).  Two devices can have identical scalar indices
while having very different hardness profiles (concentrated
vs.\ diffuse over challenge space, steep vs.\ gradual capacity
curves), so the scalar thresholds alone are not a sufficient
characterisation for all applications.  In preferred embodiments,
the protocol digest includes the full hardness curves (as tabulated
or parametric fits) alongside the scalar thresholds, together with
confidence bands estimated from repeated evaluation windows.  The
full curves are required for calibrating the query throttle
threshold and for comparing devices across model-update cycles.

\paragraph{Worked example hardness instantiation (non-limiting).}
In some embodiments, the digital hardness index $k^*_{\mathrm{dig}}(\theta;\varepsilon)$ is instantiated with a concrete attacker family to make the framework operational without overclaiming security.
A non-limiting attacker family $\{\mathcal{H}_k\}_{k\ge 1}$ may be defined as neural predictors drawn from a template family (for example U-Net-like conditional predictors for inference from the \cb, or GAN-style or diffusion-style predictors for synthesis), with the scale $k$ measuring one or more of: parameter count, FLOPs, wall-clock budget, training data budget, or query budget to the device.

In some embodiments, attacker evaluation and reporting are conditioned on a protocol digest and a meter envelope, rather than asserted unconditionally.
For example, an implementation may report that under protocol class $\Pi$ and within a meter envelope $\mathbf{m}$ observed for the run (for example drift indicators, timing alignment indicators, and stability flags), the best achievable error $\mathrm{err}^*_k(\theta)$ over $f\in\mathcal{H}_k$ remains above a policy threshold unless $k$ exceeds a reported scale.
In some embodiments, the same approach is used for semantic \TB, where the prediction target is a claim $q$ or a label rather than raw sensor traces, and audits open claim-relevant atoms and emission digests to rule out post hoc fabrication.

In some embodiments, hardness reporting is paired with an analogue hardness notion $m^*_{\mathrm{ana}}(\theta;\varepsilon, r)$, reflecting the physical scale required for cross-reactor emulation under the same protocol and meter envelope.
Together, these provide complementary bounds: computational forgery difficulty and physical emulation difficulty.

\subsection{Rate-distortion viewpoint for Limager}
In some embodiments, \LI operation is viewed as reconstruction of scene variables from a physically constrained encoding provided by the \RK channel. A rate-distortion viewpoint expresses trade-offs between reconstruction error and information retained by the channel. In practice, scan protocols and kernel settings may be chosen to move to an operating point that improves task-relevant reconstruction or classification subject to cost and stability constraints. In some embodiments, estimation-focused quantities such as Fisher information and Cram\'er--Rao-style bounds are used as non-limiting tools to characterise how accurately particular scene parameters can be estimated from the \cb under a given protocol and noise model.

\paragraph{Rate-distortion function (non-limiting).}
In some embodiments, the classical rate-distortion function for a source $\Xi$ with distortion measure $d(\Xi,\hat{\Xi})$ is
\begin{equation}
  R(D)
  = \inf_{\substack{p(\hat{\Xi} \mid \Xi) \\ \E[d(\Xi,\hat{\Xi})] \le D}}
  I(\Xi;\hat{\Xi}).
\end{equation}

\paragraph{Physical distortion bound (non-limiting).}
In the RK setting, the physical channel and any decoder together induce an effective conditional law $p(\hat{\Xi}\mid \Xi)$, and achievable distortion is bounded below by a distortion--rate function:
\begin{equation}
  \E[d(\Xi,\hat{\Xi})]
  \ge
  D^*\!\left(I_\theta(\Xi; C_{0:T}\mid U_{0:T})\right),
\end{equation}
where $D^*(R)$ denotes the minimal achievable distortion at information rate $R$ (a non-limiting distortion--rate function).
Since $\Xi \to C_{0:T} \to \hat{\Xi}$ forms a Markov chain, data processing implies $I(\Xi;\hat{\Xi}) \le I_\theta(\Xi; C_{0:T}\mid U_{0:T})$, and the bound above may be read as a physically grounded floor on reconstruction performance under the chosen protocol and meter-bounded regime.

\subsection{Wiretap and advantage gap viewpoint for Truth Beam}
\label{sec:wiretap-advantage}
In some embodiments, \TB verification is viewed through a wiretap-style lens in which a verifier (Bob) has access to a better-conditioned view of the physical channel than an adversary (Eve).
This subsection provides a non-limiting interpretation lens; it does not assert unconditional information-theoretic security.

Let $\Xi$ denote a stimulus or secret, let $Y$ denote a verifier-visible observation (for example the \cb features and meter summaries under a logged protocol), and let $Z$ denote an adversary-visible observation (for example a degraded view due to limited access to the reactor, reduced witness diversity, or restricted openings).
A non-limiting design goal is to increase an advantage gap so that the verifier can reliably discriminate authentic records while an adversary is not expected to efficiently forge records that pass verification under the same protocol digest and meter envelope.

\paragraph{Secrecy-capacity lens (non-limiting).}
In some embodiments, a secrecy-capacity style expression is used as an interpretation lens:
\[
  C_s = \bigl[I(\Xi;Y) - I(\Xi;Z)\bigr]^+,
\]
where $[x]^+=\max(x,0)$.
In this interpretation, increasing $I(\Xi;Y)$ corresponds to improving verifier information (for example via additional witnesses, better meter-bounded stability, or richer multi-channel protocols), while reducing $I(\Xi;Z)$ corresponds to limiting adversary information (for example via two-seed opening, bounded disclosures, protocol randomisation, and rate-limited physical outputs).

In some embodiments, this lens guides protocol design and policy selection by favouring regimes that widen the verifier--adversary gap, while remaining consistent with meter objectives and audit procedures described elsewhere in the specification.

\paragraph{Wiretap-channel coding embodiments (non-limiting).}
In some embodiments, the \TB advantage gap is instantiated using
wiretap-channel coding constructions of the kind associated with Wyner
(1975) and Csisz\'ar--K\"orner (1978), with the verifier-accessible
observation and the bounded-adversary observation modelled respectively
as the main and wiretap channels. In such embodiments, the
secrecy-margin objective used by the meter and training games is a
secrecy-capacity-style design target under declared protocol
randomisation, disclosure policy, witness geometry, and query-throttle
constraints; it is recorded as an operational meter-conditioned margin
rather than as an assertion of unconditional information-theoretic
security.

\paragraph{Reciprocity breaking as a physical origin of Eve degradation (non-limiting).}
In some embodiments, the physical channel deliberately introduces nonreciprocal behaviour so that an adversary is denied an equivalent symmetric view of the verifier-visible interaction in embodiments using the declared nonreciprocal configuration.
Non-limiting mechanisms include optical isolators, Faraday rotators, moving or time-varying media (including epsilon-near-zero thin films under ultrafast temporal modulation as described in Section~\ref{sec:embodiments}), diffusers with effectively irreversible scattering, and decoherence-inducing components.
In these embodiments, reciprocity breaking provides a physical origin of Eve-channel degradation in the wiretap lens, conditioned on meter-bounded stability and logged protocol regimes.

\subsection{Entropy rate of the \cb}
In some embodiments, when a closed-loop process is stationary under a fixed policy, the \cb admits an entropy-rate interpretation. Higher entropy rate, subject to calibration constraints, may correlate with higher empirical hardness because a simulator would need to reproduce richer stochastic structure to avoid detection. This is not a cryptographic guarantee, but it may be used as a calibration statistic.

\paragraph{Entropy rate as calibration statistic (non-limiting).}
In some embodiments, the entropy rate $h$ of the \cb $C_{0:T}$ under a
declared stationary and ergodic process model and policy $\pi$ is
reported as a calibration statistic alongside the empirical hardness
index $k^*_{\mathrm{dig}}$ and the meter summaries. The entropy rate
is used as a regime and drift indicator, not as a security warrant:
high entropy rate is consistent with --- but does not entail ---
empirical hardness, and a forger can in principle match the entropy
rate of the genuine source while differing in higher-order statistical
structure. In some embodiments, sustained departures of the empirical
entropy rate from the declared process model trigger meter-gated
re-characterisation or fallback to a more conservative operating
regime. Estimation of the entropy rate uses declared symbolisation,
per-device and per-envelope stratification, and lower confidence
bounds reported alongside the point estimate. No formal lower bound
on $k^*_{\mathrm{dig}}$ as a function of $h$ is asserted in this
disclosure; entropy is a calibration quantity, not a discriminator
capacity bound.

\paragraph{Conditional min-entropy, extraction, and pseudorandom generation (non-limiting).}
In some embodiments, the security-relevant entropy quantity is the
smooth conditional min-entropy of the \cb on the security domain,
conditioned on all adversary side information $E$:
$H_\infty^{\varepsilon}(C_{0:T} \mid E)$. Where the adversary may hold
quantum side information (for example an entangled probe whose
measurement is deferred), $E$ is treated as a quantum register and
$H_\infty^{\varepsilon}$ is the smooth conditional min-entropy against
quantum side information in the sense of Tomamichel, Schaffner, Smith,
and Renner (2010). Min-entropy supports operationally correct guessing
and extraction statements: an adversary guessing $C_{0:T}$ from $E$ in
a single attempt succeeds with probability at most
$2^{-H_\infty(C_{0:T}\mid E)}$, and a verifier's acceptance set of
size $M$ admits genuine-source mass at most approximately
$M \cdot 2^{-H_\infty(C_{0:T}\mid E)}$ under the relevant
side-information conditioning. Min-entropy does \emph{not} by itself
lower-bound the state size or capacity of a generic discriminator,
and no such discriminator-state lower bound is asserted in this
disclosure.

In some embodiments, the connection between physical-process
min-entropy and derived cryptographic objects is structured in two
information-theoretically separated steps. First, by the leftover hash
lemma (Impagliazzo, Levin, and Luby 1989; conditional and
smooth-min-entropy formulations following Bennett, Brassard,
Cr\'epeau, and Maurer 1995 and the generalisations surveyed in Vadhan,
\emph{Pseudorandomness}, 2012), a source with smooth conditional
min-entropy $H_\infty^{\varepsilon}(C_{0:T} \mid E) \ge m$ can be
\emph{extracted} to a string that is $\varepsilon$-close to uniform
from $E$'s view, using a universal hash function whose seed is
independent of the source and may be public; the extracted string has
length approximately $m - 2\log(1/\varepsilon)$ bits and requires no
computational hardness assumption. Where the adversary may hold
quantum side information, the corresponding quantum-proof formulation
(Tomamichel, Schaffner, Smith, and Renner 2010) is used, with
extraction parameters governed by the same $m - 2\log(1/\varepsilon)$
relation up to the standard quantum-LHL constants. Correlations across
$C_{0:T}$ are absorbed into the block min-entropy term and are not
assumed to be independent; estimating that block min-entropy from data
is a separate empirical problem, conducted under declared
symbolisation, per-device and per-envelope stratification, and lower
confidence bounds. Second, the extracted near-uniform string may serve
as the seed for a computationally secure pseudorandom generator
(PRG); the PRG's stretch and security rest on a computational hardness
assumption separate from and additional to the extraction step, for
example the existence of one-way functions in the sense of H{\aa}stad,
Impagliazzo, Levin, and Luby 1999. The leftover hash lemma provides
the information-theoretic foundation for the extraction; it does not
by itself provide PRG stretch.

\subsection{Error exponents and confidence}
In some embodiments, beyond capacity-style bounds, reliability language characterises how error probabilities decay as observation length increases. In practice, longer windows and stronger protocols may yield improved confidence for verifiers and for decoders, subject to drift and nonstationarity. These concepts may be used to select protocol window lengths, validation budgets, and acceptance thresholds.

\paragraph{Error exponent lens (non-limiting).}
In some embodiments, an error exponent (reliability function) is written as
\begin{equation}
  E(R, \theta)
  = \lim_{n \to \infty} -\frac{1}{n} \log P_e^{(n)}(R),
\end{equation}
where $P_e^{(n)}(R)$ is a minimum achievable error probability at rate $R$ over $n$ channel uses (non-limitingly, $n$ emission--observation bundles or windows of the \cb).
In some embodiments, the RK's temporal structure (multi-step protocols, adaptive control, and reactor memory) is used as a feedback-like resource that improves error decay relative to memoryless use of the channel; classical feedback results (for example Schalkwijk--Kailath for Gaussian channels) provide a non-limiting interpretation lens even when they do not directly increase Shannon capacity.

\subsection{Autopoietic interpretation (optional, non-limiting)}

In some embodiments, the \RK is interpreted through the lens
of autopoietic (self-producing) systems. The following mapping is
non-limiting and does not assert that the \RK is alive or
conscious; it identifies structural correspondences that may guide
design and analysis.

The eight necessary structures of an autopoietic system correspond to
\RK components as follows: a \emph{boundary} (the
device's physical housing and optical apertures define a topological
boundary separating internal reactor dynamics from external scene);
\emph{constitutive processes} (the reactor's nonlinear dynamics
continuously regenerate the attractor structure that defines the
device's identity); \emph{components} (reactor media, emitters,
detectors, and optical elements); a \emph{network of processes}
(the feedback loop coupling emitter, reactor, scene, and detector);
\emph{energy throughput} (optical power, electrical power, and
thermal dissipation); \emph{subordination of component production to
the network} (reactor media properties are shaped by the operating
regime, not independently); \emph{spatial extension} (the device
occupies a specific region of physical space with identifiable
boundaries); and \emph{organizational closure} (the feedback loop is
closed --- the device's output affects its own future input through the
scene, and the device's internal state determines its output,
completing the circular causality).

In some embodiments, organizational closure is the structure most
directly relevant to the specification: it is the formal property that
distinguishes the \RK's closed-loop operation from
open-loop sensing or projection, and it is the property that generates
the attractor structure, Fisher curvature, and holonomy on which all
regime-level claims depend.

\subsection{Connections and non-limiting notes}
In some embodiments, the information-theoretic viewpoint connects to other parts of the specification. Cryptographic overlays provide commitments and selective disclosure for evidence streams. Control-theoretic methods select protocols to improve information gain and robustness. Optical primitives influence channel bandwidth, invariance profiles, and noise structure. The information-theoretic viewpoint is used as a design and analysis lens and is not a claim of reduction-based security.

\paragraph{Connections to other primitive classes.}
In some embodiments, the information-theoretic view provides a shared vocabulary that connects other parts of this disclosure:
\begin{itemize}
  \item \textbf{Cryptographic primitives:} reactor-microstructure uniqueness and unpredictability map to entropy (including min-entropy and R\'enyi-entropy lenses) of channel outputs. Commitment-style protocols can be interpreted by requiring binding to correspond to non-zero information about a committed value for an honest verifier, while hiding corresponds to a degraded-view adversary obtaining negligible information under declared disclosure policies and meter envelopes.
  \item \textbf{Neural architecture primitives:} an information-bottleneck lens, $\min I(X;Z)$ subject to $I(Z;Y)\ge I_0$, provides a non-limiting interpretation of autoencoder-style \LI operation, where $Z$ is a learned representation derived from the \cb and $Y$ is a downstream task variable; reservoir-style computation can be interpreted as exploiting high $I(\Xi;C_{0:T})$ with low-dimensional sufficient statistics.
  \item \textbf{Optical primitives:} channel capacity and rate--distortion trade-offs depend on optical degrees of freedom (spatial, spectral, polarisation, temporal, coherence regime) and on transfer operators (for example LCT/\allowbreak convolutional regimes, scattering matrices, and resonant/\allowbreak eigenmode structure).
  \item \textbf{Control-theoretic primitives:} adaptive control protocols of the form $u(t)=\pi(u_{0:t-1},\allowbreak C_{0:t-1})$ act as feedback and can improve reliability and error decay and can increase usable information rate under memory and constraint conditions, even when Shannon-style capacity is not increased.
  \item \textbf{Thermodynamic primitives:} Landauer-style principles link entropy production to information erasure; in some embodiments, dissipation schedules and irreversibility provide a physical origin for one-wayness and for costs associated with high-capacity operation.
\end{itemize}

\subsection{Summary table (information-theoretic substrate, non-limiting)}
{\small
\begin{longtable}{@{}p{0.27\textwidth}p{0.34\textwidth}p{0.34\textwidth}@{}}
  \toprule
  \textbf{Quantity / lens} & \textbf{RK interpretation} & \textbf{Design use} \\
  \midrule
\endfirsthead
  \toprule
  \textbf{Quantity / lens} & \textbf{RK interpretation} & \textbf{Design use} \\
  \midrule
\endhead
  \bottomrule
\endlastfoot
  Conditional mutual information $I_\theta(\Xi;C_{0:T}\mid U_{0:T})$ &
  Information retained about a task variable or secret $\Xi$ by the \cb under a logged protocol and meter envelope &
  Choose protocols and operating regimes that increase inference power for \LI or increase verifier advantage in \TB \\[0.3em]
  Constrained capacity $C_T(\theta;\mathcal{U})$ &
  Best achievable information transfer given admissible protocol families and safety/stability constraints &
  Compare hardware and protocol design choices; budget horizons and channel allocations across modalities \\[0.3em]
  Rate--distortion $R(D)$ and distortion--rate $D^*(R)$ &
  Reconstruction trade-off between information rate and achievable distortion under the physical channel and decoder &
  Select protocol length, bandwidth, and reconstruction targets; interpret ``no decoder can beat the channel'' limits \\[0.3em]
  Secrecy-capacity lens $[I(\Xi;Y)-I(\Xi;Z)]^+$ &
  Verifier--adversary advantage gap, with Bob observing $Y$ and Eve observing a degraded view $Z$ &
  Guide \TB design toward regimes that widen Bob--Eve gaps using access asymmetries, bounded disclosures, and reciprocity breaking \\[0.3em]
  Entropy rate of $C_{0:T}$ &
  Stochastic richness of logged traces under a stationary or quasi-stationary policy &
  Calibration proxy for empirical hardness and for detecting drift or regime changes \\[0.3em]
  Error exponent $E(R,\theta)$ &
  How quickly error probabilities decay with observation length under a given rate and regime &
  Select window lengths, repetition budgets, and acceptance thresholds for desired confidence under drift constraints \\
\end{longtable}
}

\subsection{Complexity-theoretic interpretation lenses (optional, non-limiting)}
\begin{remark}[Oracle complexity classes (non-limiting)]
In some embodiments, the RK channel $\mathsf{P}_\theta(\cdot \mid S, U_{0:T})$ is treated as a physical oracle (where $S$ here denotes the scene tuple of Section~\ref{sec:definitions}, not the task variable~$\Xi$ of the preceding information-theoretic subsections), defining an oracle-relative complexity class such as $\mathsf{BPP}^{\mathrm{RK}}$: problems solvable in probabilistic polynomial time with oracle access to a physical \RK instance under declared protocol classes and meter envelopes.
This is an interpretation lens only and does not assert a formal separation of standard complexity classes. It provides vocabulary for reasoning about tasks that are easy with physical access (querying the device) but conjecturally hard to simulate with bounded digital resources under the same disclosure policies.
\end{remark}

\begin{remark}[Physical assumptions versus computational conjectures (non-limiting)]
In some embodiments, security and function are grounded in physical assumptions rather than in purely computational conjectures. Non-limiting examples include: one-wayness as a consequence of dissipative irreversibility; true randomness sourced from physical noise (shot noise, thermal noise, or other high-entropy observables) plus extraction; and unclonability sourced from microscopic disorder in reactor media. In some embodiments, analogue relaxation dynamics provide practical heuristic solvers for some optimisation problems without implying that formal complexity-theoretic open problems are resolved. These lenses are intended to clarify that the system's guarantees (where any are asserted) are operational and meter-conditioned rather than reduction-based theorems.
\end{remark}

\begin{remark}[Fundamental problems and physical sidesteps (non-limiting)]
In some embodiments, the \RK framework replaces open \emph{computational} conjectures with \emph{physical} assumptions, while leaving purely mathematical conjectures unchanged. The following catalogue is provided as a non-limiting interpretation lens and does not assert that any formal complexity-theoretic or number-theoretic open problem is solved.

\paragraph{Problems sidestepped (non-limiting).}
\begin{description}
  \item[One-way functions.]
  Computational cryptography assumes one-way functions exist but cannot prove it (a proof that one-way functions exist would imply $\mathsf{P} \neq \mathsf{NP}$).
  In some embodiments, dissipative RK channels provide \emph{physical} one-wayness via thermodynamic irreversibility: information is scattered, absorbed, and thermalised, not merely hidden.

  \item[True randomness.]
  Whether cryptographically strong pseudorandom generators exist is unknown.
  In some embodiments, quantum or other physical noise in RK detectors (for example shot noise, dark counts, thermal noise, or other quantum-limited observables) supplies \emph{actual} entropy, with extraction and health testing providing usable random bits.

  \item[Unclonability.]
  No computational proof establishes unclonable tokens in general.
  In some embodiments, microstructure-based reactor media provide physical unclonability from microscopic disorder, where the ``secret'' is the medium itself and its drift-bounded configuration.

  \item[Halting problem.]
  Turing's theorem shows no algorithm decides whether an arbitrary program halts.
  In some embodiments, physical RK systems terminate under normal operation: bounded energy, finite detector integration time, and dissipation are designed so that each protocol run completes within declared operating constraints, so the halting question does not arise operationally for device queries.

  \item[Kolmogorov complexity.]
  The shortest program producing a string is uncomputable.
  In some embodiments, the RK's dissipative channel provides \emph{physical} compression: the map $\Xi \mapsto C_{0:T}$ destroys information measurably under logged protocols, yielding finite ``physical description length'' even though algorithmic Kolmogorov complexity remains uncomputable.

  \item[G\"odel incompleteness.]
  Sufficiently powerful formal systems contain true but unprovable statements.
  In some embodiments, physical systems are evaluated by observable evidence streams rather than self-referential proofs; incompleteness does not constrain what the device can measure or verify about its channel behaviour under declared protocols.

  \item[Rice's theorem.]
  Non-trivial semantic properties of programs are undecidable.
  In some embodiments, RK observables are physical (detector amplitudes, photon counts, entropy rates, meter summaries), not semantic properties of arbitrary programs.

  \item[NP-hard optimisation.]
  Finding global optima of combinatorial problems is believed hard.
  In some embodiments, analogue physical systems (for example energy-minimising or relaxation dynamics in reactors) provide native heuristic solvers by physical relaxation without resolving $\mathsf{P}$ versus $\mathsf{NP}$.

  \item[Busy Beaver.]
  $\mathrm{BB}(n)$ is uncomputable and grows faster than any computable function.
  In some embodiments, physical systems have bounded state spaces, bounded energy, and finite time horizons; the unbounded growth regimes assumed by Busy Beaver constructions are not expected to manifest under the declared bounded state, energy, and time constraints of device operation.
\end{description}

\paragraph{Problems not sidestepped (non-limiting).}
Number-theoretic conjectures such as the Riemann Hypothesis remain mathematical statements about integers.
In some embodiments, RK modules may exhibit random-matrix-like spectra (for example GOE/GUE-like statistics in strongly scattering regimes), but this does not resolve classical conjectures about $\zeta(s)$ or prime distribution.

\paragraph{The pattern (non-limiting).}
In some embodiments, a \emph{physical assumption} (thermodynamics, quantum noise, bounded resources, microscopic disorder, drift envelopes) replaces a \emph{computational conjecture}.
The RK does not prove $\mathsf{P}\neq \mathsf{NP}$; it grounds operational security and function in physics and in meter-conditioned empirical hardness, providing physical analogues of hard problems and oracle-relative separations rather than solutions to open problems.
\end{remark}

\subsection{Transfer entropy and bidirectional information flow
(non-limiting)}
In some embodiments, transfer entropy is used as a first-class
information-theoretic primitive for quantifying directed information
flow in closed-loop \RK operation, rather than as a
post-hoc diagnostic. Let $S^{\mathrm{scene}}_t$ denote the scene
state at time $t$ (the scene component of the extended state $X_t$
of Section~\ref{sec:definitions}; note that $S$ without superscript
is reserved for the task-relevant random variable per the Notation
remark above). The transfer entropy from scene to reactor $R$ is
\[
  \TE_{S^{\mathrm{scene}}\to R}(t)
  = H(R_{t+1} \mid R_{t-k:t})
  - H(R_{t+1} \mid R_{t-k:t}, S^{\mathrm{scene}}_{t-k:t}),
\]
and symmetrically for $\TE_{R\to S^{\mathrm{scene}}}$.

In Yoked operation (\S\ref{sec:yoked}), the balance
$\TE_{S^{\mathrm{scene}}\to R} / \TE_{R\to S^{\mathrm{scene}}}$
serves as a primary coupling meter.
Balanced transfer entropy (ratio near unity) indicates bidirectional
dynamical coupling. In some embodiments, the ratio is logged
continuously and gating decisions for Yoked claims are conditioned
on it.

\paragraph{Fisher Transfer Entropy (non-limiting).}
In some embodiments, the Fisher information of the transfer entropy
with respect to coupling parameters is tracked.
In some coupled dynamical systems, this quantity (here termed
Fisher Transfer Entropy for convenience)
increases sharply near entrainment transitions (phase-locking onsets
and bifurcations in the coupled system), providing a candidate
detector for regime boundaries in Yoked operation. In some
embodiments, spikes in Fisher Transfer Entropy trigger
recalibration, regime-transition logging, or escalation in
the protocol digest.

% ======================================================================
\section{Control-Theoretic Primitives}
\label{sec:control-theory}
% ======================================================================

This section describes control-theoretic viewpoints and primitives that may be used to design, analyse, and operate \RK systems. These tools are non-limiting and do not change the core definition of a \RK. They provide methods and vocabulary for selecting protocols, ensuring stability, estimating state, and enforcing safety, including under disturbances and adversarial interference.  The viewpoints in this section establish formal relationships with control theory as interpretive vocabulary; they do not constrain the core definitions or embodiments, with the exception of the drifting field and few-shot commissioning subsection, which describes an operational mechanism for device configuration and is relied upon for enablement of configuration-related embodiments.

\subsection{State-space formulation}
In some embodiments, a \RK is modelled as a controlled dynamical system in which an extended state evolves under controls and produces observations recorded as the \cb. A non-limiting extended state includes scene state, reactor state, media state, and fast internal controller state. Controls include scan coordinates, emission settings, and detector settings. Observations include one or more detector channels. This framing supports reasoning about stability, observability, controllability, and protocol selection under chosen families of admissible controls.

\subsection{Controllability and reachability}
In some embodiments, controllability is used to characterise what internal states can be driven by admissible control sequences under physical constraints. Fast components (for example scan mirrors and modulators) may be highly controllable, while slow media (for example persistence layers or thermal states) may be only partially controllable over short horizons. In verification embodiments, limited controllability of slow degrees of freedom may increase empirical hardness because an adversary is not expected to quickly steer the physical state to match a target response distribution within the declared control-rate and access budget.

\subsection{Observability and identifiability}
In some embodiments, observability is used to characterise what aspects of the extended state are inferable from the \cb given known controls. \LI regimes may select protocols to increase observability of task-relevant scene variables. \TB regimes may intentionally reduce invertibility while preserving distinctive cross-channel constraints. In some embodiments, protocol design is treated as improving an observability metric over a chosen state representation, including by selecting scan patterns, spectral settings, or detector gating schedules.

\paragraph{Linearised observability Gramian (non-limiting).}
In some embodiments, a local linearisation is used as an interpretation lens.
Let $x_t$ denote a perturbation about a nominal operating point
$\bar{X}$ of the extended state $X_t$ (Section~\ref{sec:definitions}),
or equivalently a reduced task-relevant projection of $X_t$ within a
meter-bounded regime.  A non-limiting linearised form is:
\[
  x_{t+1} = A x_t + B u_t,
  \qquad
  y_t = C x_t + D u_t,
\]
for matrices $A,B,C,D$ determined by an operating point and meter-bounded regime.
In these embodiments, a finite-horizon observability Gramian is
\[
  W_o = \sum_{t=0}^{T-1} (C A^t)^\top (C A^t).
\]
In some embodiments, scan protocols and sensing schedules are selected to increase an observability metric such as $\det(W_o)$ or $\mathrm{tr}(W_o)$, subject to admissible control sets and safety/stability meters.

\subsection{Stability, dissipation, and boundedness}
In some embodiments, stability refers to bounded operation under bounded inputs. Physical boundedness is supported by dissipation and by safety envelopes on permissible power, duty cycle, scan speed, and parameter ranges. In some embodiments, a Lyapunov-like monitor is used, where the monitored quantity is a physical energy, a surrogate energy, a temperature proxy, or a stability score derived from the \cb. When instability risk is detected, the controller may fall back to conservative protocols, reduced power, neutral patterns, or shutdown.

\paragraph{Control-barrier and reachability safety gates
(non-limiting).}
In some embodiments, actuator commands are filtered through
control-barrier functions, barrier certificates, reachability analysis,
or model-predictive safety filters before being applied to emitters,
scanners, reactor drives, or memory-writing media. The safety filter
defines a forward-invariant safe set over optical power, temperature,
scan speed, mechanical position, stored energy, and meter-state
variables, and projects a proposed command into the admissible set
when necessary. The barrier function, safe-set parameters, solver
tolerances, override logic, and fallback action are recorded in the
protocol digest, and a command that cannot be projected without
violating a hard safety constraint causes abort, reduced-power
continuation, or a declared degraded mode.

\subsection{Filtering and state estimation}
In some embodiments, filtering estimates latent state from the noisy \cb. Non-limiting approaches include extended Kalman filtering, unscented Kalman filtering, particle filtering, and learned filters operating on \cba windows. In some embodiments, the reactor latent $Z_\theta(t)$ is implemented as a learned state estimate updated per emission-observation bundle and used for prediction, control, and meter evaluation.

\subsection{Optimal control, MPC, and information gain}
In some embodiments, protocol selection is treated as optimal control. Stage costs may include reconstruction error, verification error, energy use, drift penalties, and safety penalties. In some embodiments, model-predictive control rolls out candidate protocols using a learned emulator and executes a selected prefix, updating as new \cba data arrives. In \LI regimes, controls may also be chosen to maximise information gain about scene variables, reduce uncertainty, or increase robustness to nuisance variation.

\paragraph{Optimal control objective (non-limiting).}
In some embodiments, protocol selection is written as minimising expected cumulative cost:
\begin{equation}
  \pi^*
  = \arg\min_\pi \;
  \E\Bigl[\sum_{t=0}^{T-1} c(x_t, u_t) + c_T(x_T)\Bigr],
\end{equation}
where $c(\cdot)$ is a stage cost and $c_T(\cdot)$ is a terminal cost.
Different choices of $c$ recover the three operating regimes as special cases, for example:
\begin{itemize}
  \item \textbf{\LI:} $c$ penalises reconstruction error and control effort (and optionally rewards information gain proxies; information-gain terms enter as negative costs under the minimisation convention).
  \item \textbf{\TB:} $c$ penalises low verifier--adversary distinguishability subject to meter envelopes and safety constraints (equivalently, rewards are expressed as negative stage costs).
  \item \textbf{\RT:} $c$ penalises deviation from a target rendering or style objective while penalising meter violations and instability.
\end{itemize}
The Yoked objective $J_{\mathrm{Yoked}}$ is stated in maximisation form (Section~\ref{sec:yoked-control-objective}); this is equivalent to the cost-minimisation formulation above under $J_{\mathrm{task}} = -\sum c(x_t,u_t)$.

\subsection{Adaptive control and drift compensation}
In some embodiments, the system performs periodic calibration probes and adapts lookup tables, gain settings, and other calibration parameters. Where permitted, limited parameter subsets may be updated under rate limits and re-characterisation schedules. In some embodiments, updates to slow physical parameters are gated by approval rules or change-control procedures, and re-characterisation of meter behaviour is performed after material configuration changes.

\subsection{Drifting field and few-shot commissioning}
\label{sec:drifting-field}

In some embodiments, the configuration space $\Theta$ is equipped with
a learned vector field that guides navigation toward operating points
with desired properties. The drifting field $v(\theta)$ assigns to each
point $\theta \in \Theta$ a direction of steepest improvement for a
declared objective drawn from acceptance meters (for example
information gain, coupling strength, or a composite multi-objective;
per the meter partition of Section~\ref{sec:security-theory},
evaluation meters are held out and do not appear in the optimised
objective). The field is learned from
fleet experience: a population of devices characterised under diverse
conditions provides training data from which the field's dependence on
$\theta$ is inferred.

In some embodiments, the drifting field enables few-shot commissioning
of new devices. A new device undergoes a small number of
characterisation runs (for example fewer than ten) that locate its
approximate position in the pre-learned bundle geometry. Once
positioned, the drifting field provides a trajectory from the current
position to any desired operating regime (for example from an initial
low-curvature commissioning point to a high-curvature verification
point near a bifurcation surface). The number of commissioning runs
required scales with the dimensionality of the local fibre, not with
the full complexity of $\Theta$, because the bundle structure constrains
where the device can be.

In some embodiments, the drifting field is implemented as the natural
gradient of a declared objective on the Fisher--Rao metric:
$v(\theta) = \mathcal{F}(\theta)^{-1} \nabla_\theta J$, where $J$ is
the objective and $\mathcal{F}(\theta)$ is the Fisher information
matrix (or its regularised pseudoinverse). This aligns the navigation with the intrinsic
geometry of the configuration space.  Because the physical channel
$\mathsf{P}_\theta$ is stochastic and $\theta_{\mathrm{hw}}$ may include hardware
parameters (lens positions, gain settings, reactor drive levels) that
are not directly differentiable, $\nabla_\theta J$ is in practice
estimated from finite samples---for example via score-function
estimators, surrogate models trained on fleet data, or
finite-difference probes applied to acceptance-meter readings under
small parameter perturbations (see
Section~\ref{sec:hitl-optimisation} for detailed SPSA and
zeroth-order estimation methods). In some embodiments, the drifting
field is pre-computed for a family of objectives and stored as a
lookup table indexed by approximate device position and target regime.
In some embodiments, the drifting field is updated online as the device
acquires more data, refining the trajectory estimate.

\paragraph{Self-calibration stability under drift (non-limiting).}
A device whose sensor model $\mathsf{P}_\theta$ is calibrated using
meter outputs derived from $\mathsf{P}_\theta$ itself faces a
bootstrapping challenge: the instrument is calibrated by its own
readings.  The specification addresses this through four complementary
mechanisms that, taken together, provide bounded tracking error under
physical drift:

(i)~\emph{Inner \LI} (Section~\ref{sec:inner-limager})
uses known calibration targets or reference signals as controlled
scenes, providing external grounding that breaks the
self-referential loop.  The reference response anchors the meter
outputs to a physical standard rather than to the device's own
prior calibration.

(ii)~\emph{Fleet cross-calibration} exchanges committed reference
sweeps under standardised protocols across a population of
devices, so that no single device's drift can propagate
uncorrected.  The fleet provides the external reference that a
single device cannot provide for itself.

(iii)~\emph{Iterated self-modelling convergence (non-limiting theoretical correspondence):} in selected embodiments, a depth curriculum is used to drive self-modelling toward a stable fixed point; convergence properties under this curriculum are discussed in the categorical unification section as a conceptual or diagnostic correspondence and are not required by the operational mechanisms otherwise disclosed.

(iv)~\emph{The meter partition}
(Section~\ref{sec:security-theory}) is arranged so that evaluation
meters used to assess calibration quality are held out from the
optimisation loop, thereby reducing the risk of the drifting field
chasing metrics that drift with the device.

Under active physical drift, the fixed point itself moves, and
convergence is replaced by a tracking problem: the calibration must
track the moving target closely enough that meter outputs remain
meaningful.  In some embodiments, tracking error is bounded by
requiring that the drift rate $\|\dot{\theta}\|$ (measured in
Fisher--Rao geometry) remain below the contraction rate of the
self-calibration loop, so that calibration converges faster than
the device drifts.  When drift rate exceeds this bound, meters
flag an out-of-envelope condition and the device enters a
re-characterisation mode (Inner \LI with reference targets)
until tracking is re-established.  The drift rate, contraction
rate estimate, and re-characterisation transitions are logged in
protocol digests.

\subsection{Adversarial control and robust design}
In some embodiments, \TB operation is treated as a game in which the controller selects challenges and an adversary selects disturbances or forgery attempts under bounded resources. Robust control ideas may be applied by treating the adversary as a bounded disturbance and selecting protocols that maximise distinguishability under worst-case interference. Non-limiting tools include adversarial training over protocol families, conservative acceptance rules when uncertainty rises, and protocol randomisation within safety envelopes.

\paragraph{Stackelberg equilibrium lens (non-limiting).}
In some embodiments, the verification protocol is treated as a Stackelberg game in which a verifier commits to a protocol $u$ (shorthand for $U_{0:T}$) and an adversary best-responds with a forgery pipeline under bounded resources. Let $\mathsf{P}^{\mathrm{Bob}}_\theta(\cdot\mid u)$ denote a \cba distribution induced by a genuine RK under protocol $u$, and let $\mathsf{P}^{\mathrm{Eve}}(\cdot\mid u)$ denote a corresponding distribution induced by an adversary's forgery pipeline for the same declared protocol. In these embodiments, an equilibrium protocol is expressed as
\begin{equation}
  u^*
  = \arg\max_u \min_{\mathrm{Eve}}
  \; D_{\KL}\!\Bigl(\mathsf{P}^{\mathrm{Bob}}_\theta(\cdot\mid u)\,\big\|\,\mathsf{P}^{\mathrm{Eve}}(\cdot\mid u)\Bigr),
\end{equation}
choosing protocols that are robust against an adversary's best response within the assumed attacker family.

\paragraph{$H_\infty$ robust control lens (non-limiting).}
In some embodiments, adversarial interference is modelled as a bounded-energy disturbance acting on sensors, actuators, or the scene channel. In these embodiments, $H_\infty$ robust control provides a non-limiting interpretation lens: the verifier selects protocols that limit worst-case induced gain from disturbances to verification-relevant residuals, and meters gate operation when disturbances exceed envelope assumptions.

\subsection{Safety envelopes and interlocks}
In some embodiments, a safety envelope bounds at least one of emission power, exposure duration, scan speed, flicker, thermal load, and other hazard-relevant constraints. In some embodiments, the envelope is enforced by a separate interlock layer so that even faulty or adversarial policies are blocked from commanding emissions outside permitted bounds in embodiments using the stated interlock configuration. This is compatible with verification gating and with cryptographic overlays, but does not require them.

\subsection{Non-limiting notes (control-theoretic substrate)}
Control-theoretic primitives support all three regimes. \LI uses them for estimation and information gain. \RT uses them for stability and constrained stylisation. \TB uses them to design adversary-robust challenges and safe operation under bounded interference.

\subsection{Summary table (control-theoretic substrate, non-limiting)}
{\small
\begin{longtable}{@{}p{0.27\textwidth}p{0.34\textwidth}p{0.34\textwidth}@{}}
  \toprule
  \textbf{Control primitive} & \textbf{RK realisation} & \textbf{Regime / application} \\
  \midrule
\endfirsthead
  \toprule
  \textbf{Control primitive} & \textbf{RK realisation} & \textbf{Regime / application} \\
  \midrule
\endhead
  \bottomrule
\endlastfoot
  Controllability / reachability &
  Admissible scan laws, emission settings, and slow configuration knobs drive the internal media and observation distribution &
  \LI: explore measurement space; \TB: limited slow controllability can strengthen hardness \\[0.3em]
  Observability / identifiability &
  State components inferable from the \cb under known controls; linearised Gramian $W_o$ as an interpretation lens &
  \LI: choose protocols to increase observability; \TB: preserve distinctive constraints while limiting invertibility \\[0.3em]
  Stability / dissipation &
  Physical boundedness plus meters and interlocks; Lyapunov-like monitors from energy or stability proxies &
  All regimes: safe operation, regime downgrades, and consistent calibration under drift \\[0.3em]
  Filtering / state estimation &
  EKF/UKF/particle filters or learned filters updating a latent $Z_\theta(t)$ from bundles $c_t$ &
  \LI: reconstruct scene variables; \TB: widen Bob--Eve estimation gap via private ports and meters \\[0.3em]
  Optimal control / MPC &
  Protocol optimisation under costs combining task objectives and meter penalties; receding-horizon selection with emulators &
  \LI: information gain and reconstruction; \RT: constrained stylisation; \TB: adversary-robust challenges \\[0.3em]
  Robust / adversarial design &
  Worst-case or game-theoretic selection of protocols; minimax lenses such as Stackelberg KL gaps; $H_\infty$ approximations &
  \TB: protocol families resistant to bounded interference and forgery strategies \\[0.3em]
  Safety envelopes and interlocks &
  Independent enforcement of power, exposure, scan speed, thermal, flicker and other limits &
  All regimes: blocks or screens unsafe commands and constrains the admissible protocol class $\mathcal{U}$ in embodiments using the interlock layer \\
\end{longtable}
}

\subsection{Yoked operation: coupling quality as control
objective (non-limiting)}
\label{sec:yoked-control-objective}
In Yoked operation, the control objective includes a coupling-quality
term that drives the system toward bidirectional dynamical coupling
between device and scene:
\[
\begin{aligned}
  J_{\mathrm{Yoked}}(U_{0:T})
  &= J_{\mathrm{task}}(U_{0:T}) \\
  &\quad + \lambda_{\TE}\,\mathrm{Balance}(\TE_{S\to R}, \TE_{R\to S}) \\
  &\quad + \lambda_{\mathrm{mag}}\,\min(\TE_{S\to R}, \TE_{R\to S}) \\
  &\quad - \lambda_{\CLE}\,|\CLE - \CLE_{\mathrm{target}}|.
\end{aligned}
\]
where $\mathrm{Balance}(\cdot)$ penalises asymmetric information flow,
$\min(\TE_{S\to R},\TE_{R\to S})$ rewards large mutual transfer entropy
in both directions (the ``both large'' Yoked signature), and
the CLE term penalises deviation from a declared target
$\CLE_{\mathrm{target}} < 0$ chosen close to zero so as to maintain
operation at the edge of synchronisation.
The weights $\lambda_{\TE}, \lambda_{\mathrm{mag}}, \lambda_{\CLE} \ge 0$
are non-limiting hyperparameters that trade off symmetry, magnitude, and
convergence depth; in some embodiments they are scheduled or adapted during
a coupling curriculum.  In preferred embodiments, all TE quantities are
reported in a declared information unit (bits or nats) under the declared
estimator family and windowing, CLE quantities are reported in the declared
time base (per-step or per-second), and $J_{\mathrm{task}}$ is normalised
to the same declared scale; the $\lambda$ weights include the necessary
unit conversions so that all four terms are commensurate.  The
normalisation choices are declared in the protocol digest.
In preferred embodiments, the objective $J_{\mathrm{Yoked}}$ is
maximised.

\paragraph{CLE convergence criterion (non-limiting).}
In some embodiments, Yoked operation is declared converged when the
conditional Lyapunov exponent $\CLE$ approaches zero from below, per
Hart's criterion. The sign and magnitude of $\CLE$ are logged as
meter outputs and used to gate Yoked claims.

\paragraph{Non-limiting Balance functional example.}
In some embodiments, a concrete instantiation of
$\mathrm{Balance}(\TE_{S\to R}, \TE_{R\to S})$ is the negative
absolute log-ratio:
$\mathrm{Balance}(a,b) = -|\log((a+\epsilon)/(b+\epsilon))|$,
where $\epsilon > 0$ is a regularisation constant, and the objective
is maximised (penalising asymmetric information flow). Transfer
entropy is estimated using one of the non-limiting estimator families
described in Section~\ref{sec:regimes} (for example
Kraskov--St\"ogbauer--Grassberger).

\paragraph{Arnold tongue analysis (non-limiting).}
In some embodiments, controllability in Yoked operation is
characterised by mapping entrainment regions (Arnold tongues) in the
space of coupling strength versus detuning (frequency ratio between
scene dynamics and reactor dynamics). Protocol design targets
operating points within entrainment tongues where coupling is robust,
and meters flag excursions toward tongue boundaries where coupling
becomes fragile. In some embodiments, tongue boundaries are
estimated empirically from sweep data and used to define the
admissible Yoked operating envelope.

\paragraph{Active-inference interpretation (non-limiting).}
In some embodiments, the \RK loop may be interpreted through
the lens of active inference: the \cb $C_{0:T}$ supplies
sensory evidence, the control protocol $U_{0:T}$ realises epistemic
actions (active sampling to reduce uncertainty) and pragmatic actions
(interventions to achieve goals), and the reactor dynamics implement a
physically grounded generative model whose predictions are compared
against observations via meter scores. Under this interpretation,
\LI operation corresponds to epistemic foraging (actions chosen to
maximise expected information gain), \TB corresponds to
model-evidence comparison (evaluating whether observations are
consistent with the enrolled generative model), and \RT
corresponds to pragmatic action (driving the scene toward a target
state). Yoked operation corresponds to coupled active inference
between two generative models (device and scene) whose joint
free-energy minimum corresponds to the targeted coupled attractor.
This interpretation is offered as a non-limiting explanatory lens that
may aid design intuition and connect the \RK formalism to an
active research literature; it does not alter the definitions, claims,
or formal results elsewhere in this specification, and the \RK formalism is not derived from or dependent on the free-energy
principle. A Markov blanket, in this context, is the minimal set of
variables (active controls and sensory observations) that separates an
agent's internal state from the external world in a probabilistic
graphical model. The Markov blanket interpretation
(described in the related Filing 2 application; optional, non-essential)
provides additional detail on the boundary-probe implications of this
lens.


\section{Cryptographic-Primitive Substrate}
\label{sec:crypto-primitives}
% ======================================================================

This section describes non-limiting cryptographic-primitive constructions that may be implemented using \RK systems. These constructions derive practical value from a combination of (i) standard cryptographic primitives for commitments, signatures, and transport, and (ii) empirical properties of a physical \RK channel, including device-specific microstructure, noise, and history dependence. No reduction-based cryptographic guarantee is required by these embodiments. The constructions are optional overlays and do not alter the core definition of the \cb or the \RK channel.

\subsection{Mapping meter games to analogue cryptographic structures}
In some embodiments, multi-agent meter games are mapped directly to physical configurations of mobile or fixed \RK platforms, including PolieBot configurations, to implement analogue versions of cryptographic structures that exploit real-world asymmetries. In such embodiments, roles analogous to prover, verifier, and adversary are instantiated by physical devices with different access, vantage points, and constraints.

Non-limiting real-world asymmetries that may be exploited include:
\begin{itemize}
  \item \textbf{Access asymmetry:} a verifier device has direct physical access to a reactor, a scene, or a private detector port that an adversary is assumed not to have access to absent physical compromise of the verifier device.
  \item \textbf{Geometry and viewpoint asymmetry:} devices observe the same event from different vantage points, so a forged history would need to satisfy multi-view spatiotemporal constraints.
  \item \textbf{Timing and causality asymmetry:} observable causal effects and bounded propagation delays constrain permissible transcripts without requiring a shared clock.
  \item \textbf{Interference and power asymmetry:} an adversary may be bounded in co-illumination power, duty cycle, or spectral overlap relative to a verifier's sensing configuration.
  \item \textbf{Microstructure asymmetry:} device-specific optical or multi-physics microstructure yields microstructure-unclonable behaviour that is empirically difficult to emulate without comparable hardware.
\end{itemize}

In these embodiments, meter outputs (verisimilitude, hardness, calibration, and spatiotemporal consistency) provide the decision statistics for analogue commitment-open protocols, proof-of-projection protocols, device-bound rate limiting, and web-of-trust updates. The result is a family of physically anchored protocol patterns that behave like cryptographic constructions in operation while being grounded in physical constraints and empirically calibrated hardness.

\subsection{Non-limiting hardness assumptions and attacker model}
In some embodiments, security arguments rely on empirical hardness assumptions concerning physical modelling and forgery rather than on number-theoretic hardness alone. Non-limiting assumptions include:
\begin{itemize}
  \item \textbf{RK microstructure unpredictability:} given access to many observed transcripts from a device, an adversary without the ability to reproduce the corresponding physical measurement event under the declared conditions has difficulty predicting responses to fresh challenges within an acceptance region.
  \item \textbf{RK-distinguishability:} given mixed samples, a competent meter distinguishes physical traces produced by a device family from simulated traces produced by a bounded emulator family, for a chosen operating point and protocol distribution.
  \item \textbf{RK one-wayness by dissipation:} in sufficiently dissipative configurations, reproducing the full joint statistics of the \cb, or inverting a run summary to recover a scene or protocol that matches the same distribution, is empirically difficult without the ability to reproduce the corresponding physical measurement event under the declared conditions.
\end{itemize}

\paragraph{Security parameter (non-limiting).}
In some embodiments, an effective security parameter $n_{\mathrm{sec}}$ scales with total transcript information content. A non-limiting discrete proxy is
\[
  n_{\mathrm{sec}}
  \approx
  \min\!\bigl(\log_2|\mathcal{X}|,\; nT\bigr),
\]
where $|\mathcal{X}|$ is a count of admissible challenges or protocol seeds (for example distinct scan protocols), $n$ is a non-limiting count of reliably extractable bits per response after noise and quantisation, and $T$ is a non-limiting count of independent challenge--response bundles. In some embodiments, $n_{\mathrm{sec}}$ is tied to an estimated entropy of the joint transcript rather than to any single syntactic parameter.

\paragraph{Negligible-bound lens and design-target framing (non-limiting).}
In some embodiments, the standard cryptographic notation $\mathrm{negl}(n_{\mathrm{sec}})$ is used as an interpretation lens to express empirical design targets under stated attacker families and meter envelopes. These bounds are not asserted as theorems; they are targets evaluated and re-evaluated empirically as modelling and attack methods improve.

\paragraph{RK microstructure-hardness lens (non-limiting).}
In some embodiments, for a response space $\mathcal{Y}$ and any efficient adversary $\mathcal{A}$ without the ability to reproduce the corresponding physical measurement event under the declared conditions, a design target is
\[
  \Pr[\mathcal{A} \text{ predicts correctly}]
  \le
  \frac{1}{|\mathcal{Y}|} + \mathrm{negl}(n_{\mathrm{sec}}),
\]
where the probability ranges over fresh challenges and device noise realisations under a declared protocol class and meter envelope.

\paragraph{RK-distinguishability lens (non-limiting).}
In some embodiments, a design target is that no efficient adversary distinguishes genuine device-produced transcripts from efficient simulations with more than negligible advantage:
\[
  \bigl| \Pr[\mathcal{A}(\mathrm{real}) = 1] - \Pr[\mathcal{A}(\mathrm{simulated}) = 1] \bigr|
  \le
  \mathrm{negl}(n_{\mathrm{sec}}).
\]

\paragraph{Unclonability via total-variation gap (non-limiting).}
In some embodiments, a device is treated as implementing a noisy microstructure-unclonable response distribution $f(x)$ on challenges $x\in\mathcal{X}$, and unclonability is expressed by requiring that for any candidate clone $\tilde{f}$, fresh-challenge response distributions are separated in total variation distance on a non-negligible fraction of challenges:
\[
  \Pr_{x \gets \mathcal{X}}\!\left[d_{\mathrm{TV}}\!\bigl(f(x), \tilde{f}(x)\bigr) > \varepsilon\right]
  \ge
  1 - \mathrm{negl}(n_{\mathrm{sec}}).
\]

\paragraph{ML-based modelling as first-class attack strategy (non-limiting).}
In some embodiments, the adversary model explicitly includes machine-learning-based surrogate modelling (emulators trained from transcripts), and empirical hardness indices are reported relative to bounded families of such emulators and their best-achieved performance under the stated protocol distribution and meter envelope.
These are empirical assumptions calibrated to realistic noise models, budgets, and adversary capabilities, and may be re-evaluated over time as modelling methods improve.

\subsection{Commitments to the \cb and selective disclosure}
In some embodiments, the system produces a commitment to the \cb (or to windowed segments) so later parties can verify integrity without receiving the entire trace. The device logs the \cb and produces one or more commitments, including without limitation: a hash commitment to a trace, Merkle roots for windows or atoms, or commitments to summaries together with commitments to underlying atoms.

In some embodiments, selective disclosure is supported by opening requested atoms or windows together with authentication paths, enabling verification against a prior commitment. This supports auditability under bandwidth and privacy constraints.

\paragraph{Hiding and binding lenses (non-limiting).}
In some embodiments, hiding and binding are expressed as design targets for the commitment-and-opening protocol rather than asserted as unconditional guarantees.
Let $M$ denote a committed message variable (or committed statement), and let $\tau$ denote a transcript or commitment view available to an adversary. A non-limiting hiding target is
\[
  I(M;\tau) \le \varepsilon,
\]
for small $\varepsilon$ under the stated disclosure policy and adversary view.
A non-limiting binding target is that it is infeasible to produce two distinct valid openings consistent with the same commitment:
\[
  \Pr\bigl[\mathcal{A} \text{ outputs valid } (m_0, r_0), (m_1, r_1) : m_0 \neq m_1\bigr]
  \le
  \mathrm{negl}(n_{\mathrm{sec}}),
\]
where $r$ denotes non-limiting opening randomness and validity is checked by a verifier using committed atoms, protocol digests, and any disclosed windows.
Note that the hiding target above is information-theoretic (bounded
mutual information), while the binding target uses the computational
convention $\mathrm{negl}(n_{\mathrm{sec}})$.  This hybrid formulation is
intentional: hiding is assessed by information-leakage budgets
consistent with the spec's empirical-hardness posture, while binding
relies on the digital commitment layer's computational properties.
The ``selective disclosure'' described here refers to Merkle-style
audit openings and does not presume any particular formal
selective-opening security notion from the cryptographic literature
unless explicitly claimed and parameterised for a specific
embodiment.

\subsection{microstructure-based identification and authentication}
\label{sec:puf-auth}
In some embodiments, a \RK module functions as a microstructure-unclonable challenge-response oracle. A challenge comprises a control protocol or seed that generates a protocol. A response comprises the measured \cb or a derived signature. Verification is performed by testing membership in an acceptance region, optionally implemented by a learned discriminator or calibrated statistical test.

In some embodiments, device enrolment produces reference data, a learned verifier, or helper data supporting robust matching under noise. In some embodiments, helper data includes secure sketches and/or fuzzy extractor constructions that derive stable keys or identifiers from noisy physical measurements while preserving unclonability and unpredictability in practice.
In some embodiments, microstructure-based authentication is combined with the
device configuration representation
(Section~\ref{sec:device-config-rep}), enabling fleet-level identity
management alongside device-level verification; the irreducible
residual (Section~\ref{sec:error-decomposition}) provides the
security margin that raises the difficulty of model-based forgery of challenge-response
pairs.

\subsection{microstructure-backed signatures and attestations}
In some embodiments, signatures and attestations are implemented using a microstructure-backed process. To attest to a message, the system derives a message-dependent challenge protocol, executes it through the physical kernel, and returns at least one of: a committed transcript, a committed summary, and openings sufficient for verification.

In some embodiments, the secret key is physical possession of a device's microstructure and calibration state, and public verification material comprises a verifier model, reference commitments, or enrolment summaries. In some embodiments, microstructure-backed attestations are combined with conventional digital signatures for compatibility with standard public-key infrastructure.

\subsection{One-way state updates and hash-chain logging}
In some embodiments, the controller maintains a state that is updated step-by-step as a one-way function of prior state and logged measurements. A terminal value serves as a compact fingerprint of a run or session. In some embodiments, the emission policy is driven by the evolving state so that challenges are chained to past observations, improving replay resistance.

The one-way update family may include, without limitation, hash functions, sponge constructions, MACs, PRFs, VRFs, and ZK-friendly algebraic hashes, selected based on implementation constraints and, where used, proof-system cost.

\paragraph{One-wayness and preimage lens (non-limiting).}
In some embodiments, the resulting terminal digest is treated as a hash-like compression of a run transcript. A non-limiting design target is that given a released terminal digest $\chi$, it is infeasible to construct an alternative transcript (or alternative stimulus/protocol pair) that yields the same digest while falling outside a disclosed equivalence relation:
\[
  \Pr\bigl[\mathcal{A}(\chi) = S' : \mathsf{H}_\theta(S') = \chi,\; S' \not\sim S\bigr]
  \le
  \mathrm{negl}(n_{\mathrm{sec}}),
\]
where $\mathsf{H}_\theta$ denotes a non-limiting digest function derived from the one-way update rule, and $\sim$ denotes a non-limiting equivalence relation (for example scenes differing only within an accepted tolerance band).

\subsection{Physically rate-limited outputs}
In some embodiments, cryptographic outputs are rate-limited by requiring completion of a physical run prior to release. For example, to produce a signature-like output on a message, the device executes a run and releases an output after a successful meter acceptance check, in embodiments using that release policy. This ties output rate to physical time and device availability and reduces offline bulk precomputation.

\subsection{Key agreement from correlated measurements}
\label{sec:key-agreement-correlated}
In some embodiments, two parties derive shared secrets from correlated measurements. Two devices (or two parties interacting with a shared channel) execute correlated protocols and obtain correlated \cba-derived bits. They perform reconciliation over a public channel and privacy amplification to derive a shared key. Security depends on an adversary observing a degraded view relative to legitimate parties. This resembles key agreement from correlated randomness rather than conventional public-key exchange.

\paragraph{Key length lens (non-limiting).}
In some embodiments, letting $X$ and $Y$ denote legitimate party views and $Z$ denote an adversary view of the same interaction under authenticated public discussion and a declared source model, a non-limiting key-length design lens for secret-key agreement from correlated randomness in the wiretap / degraded-source special case is
\[
  |K|
  \approx
  I(X;Y) - I(X;Z) - O(\log(1/\varepsilon)),
\]
where $\varepsilon$ is a non-limiting security tolerance parameter and $O(\cdot)$ suppresses constant factors. The general secret-key agreement rate is model-dependent (source model, public-discussion model, authentication assumptions, reconciliation leakage, adversary class); this lens is used as a design guide for the special-case regime and does not assert a generic capacity expression. In some embodiments, this lens is used as a design guide for increasing legitimate correlation (for example via private detector ports or richer protocols) while reducing adversary correlation (for example via bounded disclosures and access asymmetries).

\subsection{Secure aggregation, threshold schemes, and multi-device
attestation}
In some embodiments, multiple devices produce joint outputs using
threshold and aggregation mechanisms. Non-limiting examples include
threshold signatures (a quorum jointly produces a signature),
secure aggregation of meter outputs (individual device scores are
combined without revealing per-device values), and homomorphic
aggregation of encrypted scores. These mechanisms support collective
attestation and fault tolerance in multi-device deployments.

\subsection{Lifecycle: drift, re-enrolment, and revocation}
In some embodiments, device drift and lifecycle are managed explicitly. Registration or enrolment characterises a device and sets verifier material. Operation proceeds under noise and drift constraints. Re-enrolment updates acceptance regions, verifier models, or helper data. Revocation procedures address compromised or retired devices. In some embodiments, physical destruction or irreversible media reset provides strong decommissioning.

\subsection{Non-limiting notes (cryptographic substrate)}
These constructions are intended as optional overlays. They can be mixed and matched and may use conventional cryptographic primitives for transport and standardisation while relying on empirical physical hardness for additional constraint. They do not claim mathematical proof of security. They provide an engineering framework for tamper-evidence, auditability, privacy-aware disclosure, and device-bound scarcity grounded in physical measurement logs.

\paragraph{Witness-indistinguishable proof lens for scene properties (optional, non-limiting).}
In some embodiments, the system supports protocols in which a prover demonstrates a scene property without disclosing a full scene description. This is a non-limiting interactive lens and does not assert full zero-knowledge; formal simulation-based definitions and proofs are not required by these embodiments.
In some embodiments, formal proof systems (for example zk-SNARKs, zk-STARKs, or Bulletproofs) are used to prove predicates over committed summaries or reactor signatures under a declared disclosure policy.
Let $P(S)\in\{0,1\}$ denote a predicate on scenes (for example ``the inspected seal is intact'' or ``the part includes a required fiducial pattern''). A non-limiting protocol sketch is:
\begin{enumerate}
  \item A verifier selects a fresh challenge (for example a protocol seed or scan schedule) and commits to it when required by policy.
  \item A prover with physical access executes the corresponding protocol $U_{0:T}$ on a scene $S$ and returns a committed transcript and/or disclosed window summaries sufficient for the verifier to evaluate $P$.
  \item The verifier checks protocol binding, meter envelope compliance, and that disclosed atoms are consistent with the committed transcript, then accepts if the evidence is consistent with $P(S)=1$ under the declared regime.
\end{enumerate}
In some embodiments, soundness is grounded in RK hardness lenses: without access to a valid witness scene and a compatible physical channel, producing accepted evidence under fresh challenges is empirically difficult for bounded adversaries. In some embodiments, witness indistinguishability is supported by the existence of multiple scenes (or multiple micro-variants) satisfying $P$ that yield statistically similar accepted transcripts under bounded disclosure, so the verifier learns the property outcome but not which witness was used beyond what is inevitably leaked by the disclosed windows.

\paragraph{Family membership proofs via bundle topology (optional,
non-limiting).}
In some embodiments, topological invariants of the attractor bundle
provide a basis for proving family membership without revealing device
identity. Topological invariants (for example characteristic classes or
Euler class of the bundle) are determined by the manufacturing process
and are shared across all devices from the same production family, while
individual holonomy depends on each device's specific microstructure.
In some embodiments, a protocol proves ``this device belongs to family
$F$'' (the topological invariant matches the family's declared
invariant) without revealing ``this is device serial number $N$''
(the specific holonomy is not disclosed). In some embodiments, the
converse is also supported: proving specific device identity without
revealing which family the device belongs to. This separation between
topology (family) and geometry (individual) is a natural property of the
bundle structure rather than a construction requiring additional
cryptographic machinery, though in some embodiments it is combined with
conventional zero-knowledge proof systems for formal guarantees.

\paragraph{True random number generation (TRNG) and min-entropy (non-limiting).}
In some embodiments, a \RK is configured to expose a high-entropy physical noise source (for example shot noise, thermal noise, or quantum-limited observables) and produce raw samples $Z_t$. In some embodiments, an entropy model assumes a conditional min-entropy bound
\[
  H_\infty(Z_t \mid \mathcal{E}) \ge k,
\]
where $\mathcal{E}$ denotes side information available to an adversary. In some embodiments, an extractor $\mathrm{Ext}$ is applied to a window $Z_{1:n}$ to produce output bits that are $\varepsilon$-close to uniform:
\[
  d_{\mathrm{TV}}\bigl(\mathrm{Ext}(Z_{1:n}), U_m\bigr) \le \varepsilon.
\]
In some embodiments, health tests are run continuously on the raw stream to detect degradation or manipulation and to trigger fallback or re-characterisation.

\paragraph{Summary table (non-limiting).}
{\small
\begin{longtable}{@{}p{0.27\textwidth}p{0.34\textwidth}p{0.34\textwidth}@{}}
  \toprule
  \textbf{Primitive} & \textbf{RK instantiation} & \textbf{Key mechanism (lens)} \\
  \midrule
\endfirsthead
  \toprule
  \textbf{Primitive} & \textbf{RK instantiation} & \textbf{Key mechanism (lens)} \\
  \midrule
\endhead
  \bottomrule
\endlastfoot
  Commitments and selective opening &
  Commit atoms/windows of the \cb; open only challenged windows with authentication paths &
  Tamper-evident provenance; hiding/binding targets under bounded disclosure and meter envelopes \\[0.3em]
  microstructure identification &
  Reactor microstructure as challenge--response oracle; acceptance region or learned verifier &
  Unpredictability and unclonability from device-specific disorder and noise \\[0.3em]
  microstructure-backed attestations &
  Message-derived protocol seed; run produces committed transcript or summary plus openings for verification &
  Unforgeability tied to physical access and meter-validated regime \\[0.3em]
  Hash-chain logging &
  One-way chain state $\chi_t$ updated from bundles $c_t$; terminal digest $\chi_T$ anchors session history &
  Replay resistance and tamper evidence; preimage-style targets for alternative transcript search \\[0.3em]
  Rate-limited outputs &
  Release tags/keys only after a physical run passes meters &
  Physical time and device availability bound output rate; reduces offline bulk forgery \\[0.3em]
  Key agreement &
  Correlated RK measurements + reconciliation + privacy amplification &
  Secret-key agreement from correlated randomness with degraded adversary view \\[0.3em]
  WI / ZK-style proof lens &
  Interactive challenge--commit--open protocols proving a property $P(S)$ with bounded disclosure &
  Soundness from empirical hardness under fresh challenges; limited witness leakage via noise and disclosure policy \\[0.3em]
  TRNG &
  High-noise RK regime + extraction + health tests &
  Min-entropy bounds and extractor closeness as design lenses \\[0.3em]
  Multi-witness attestation &
  Threshold signatures / aggregation over commitments and meter summaries &
  Fault tolerance under partial compromise \\
\end{longtable}
}

% ======================================================================

\paragraph{Committed scene representations and forgeability
asymmetry (non-limiting).}
In some embodiments in which the declared target $\phi$ is a
committed primitive-set representation
(Section~\ref{sec:style-parameter}), a deception attempt against
the declared target may require producing (a)~a primitive set
$\{p_k\}_{k=1}^{K}$ consistent with a claimed scene and (b)~a
rendering pathway consistent with the specific enrolled
device's independently calibrated forward-model parameters,
including $K_{\mathrm{spot}}$, $p_{\mathrm{eff}}$,
$w_{\mathrm{det}}$, and associated world-to-device coordinate
transforms drawn from, or referenced by, the enrolled
$\mathbf{d}_i$ and the declared challenge configuration $X$ or
protocol geometry. The first may, in some embodiments, be
performed using rasterisation-based optimisation without
reproducing the enrolled physical forward model, whereas the
second is evaluated under the declared threat model and meter
envelope by the canonical hardness index pair
$(k^*_{\mathrm{dig}}(\theta;\varepsilon),
m^*_{\mathrm{ana}}(\theta;\varepsilon, r))$ of
Section~\ref{sec:security-theory}, and requires that the
primitive-set prediction match the physical measurement not
only photometrically for a single view but also consistently
across held-out scan geometries, meter windows, detector
settings, and temporal persistence behaviour specified by
$U_{0:T}$ and the protocol digest. In such embodiments, the
following separation rules apply: primitive parameters are
optimised under declared scene-space priors and meter
envelopes; device kernels $K_{\mathrm{spot}}$,
$p_{\mathrm{eff}}$, and $w_{\mathrm{det}}$ are held fixed at
their independently calibrated values during verification
rather than fitted as free parameters; and primitive covariance,
density, opacity, appearance-bandwidth, and cardinality bounds
are expressed as a declared representation-complexity envelope,
with spatial support, covariance, and density components
expressed in declared world-coordinate units, including at
least a maximum $K$, covariance eigenvalue or support-radius
bounds, primitive-density or support-overlap limits,
appearance-bandwidth bounds, and per-primitive or aggregate
opacity/radiance-energy bounds, distinct from the device's
device-space PSF. Where these separation rules are satisfied,
the deception-symmetry and unified-trust-root limitations of
Sections~\ref{sec:deception-symmetry} and
\ref{sec:unified-trust-root} specialise to a concrete
evidentiary structure: the committed primitive set is the
claimed-scene commitment; the enrolled physical channel, its
calibration record, and the committed device configuration
representation together supply the verifier's forward model;
and verification requires held-out consistency between
rasterised prediction and physical measurement under the
shared $U_{0:T}$.
\section{Categorical Unification of Substrate Mappings (Optional, Non-limiting Theoretical Framework)}
\label{sec:categorical-unification}
% ======================================================================

\textbf{Note for prosecution.}  This section provides a non-limiting
mathematical framework for understanding relationships between the
substrate mappings described in the preceding sections.  It does not
define additional claim subject matter and does not constitute a
separate embodiment; it is provided as theoretical background for
practitioners familiar with category theory and should be read as an
interpretive tool rather than a technical teaching.  No claim of the
present disclosure depends on any result or conjecture stated in this
section.

The preceding five sections
(Sections~\ref{sec:optical-primitives}--\ref{sec:crypto-primitives})
each map the \RK formalism onto a different established
theoretical framework: optical physics, neural-network architectures,
information theory, control theory, and cryptographic primitives.
Each mapping is constructed independently, and each reveals structure
in the \RK that the others do not. This section presents a
non-limiting theoretical framework, stated as conjectures rather than
proven results, that unifies these five mappings as instances of a
single categorical construction. The framework provides precise
language for transferability, compositional structure, and
self-modelling, and identifies specific mathematical claims that are
testable and falsifiable.

\paragraph{Containment statement.}
These categorical constructions are an optional descriptive lens and
do not introduce additional hardware requirements, operational steps,
or claim limitations beyond those already described in the preceding
sections.  No embodiment requires the categorical framework for
implementation.  The framework may be removed entirely without
affecting the enablement or scope of any disclosed embodiment.
Conjectures~2 and~3 apply specifically to embodiments whose reactor
dynamics satisfy a Turing-completeness condition; this is an
additional design target not achievable with all reactor types, and
the categorical framework introduces no such requirement for existing
embodiments.

\paragraph{Non-load-bearing status of all conjectures.}
\textbf{No core claim in this specification depends on any of
Conjectures~1--5.}  Each conjecture is a research-programme item
whose truth or falsity does not affect the enablement, operation, or
security properties of any described embodiment.  The operational
work that Conjecture~4 (Mod as contraction) was partly intended to
formalise is instead performed by the PoD architecture itself: the
PoD protocol physically anchors complexity growth through the
confirmation requirement; the security/discovery domain partition
prevents self-modelling from eroding verification stability; and the
quorum structure bounds individual divergence.  The practical
theorem---that the PoD architecture keeps the fleet's collective
self-modelling in a bounded regime by design---is independently
grounded in the protocol specification and does not require the
categorical apparatus.  The categorical framework is retained as a
research programme; its conjectures are open problems, not
established results.

\subsection{The weak category of dissipative Hamiltonian systems}
\paragraph{Objects.}
An object of the category $\mathbf{DissHam}$ is a tuple
$(M, \omega, H, \Gamma, \Theta)$ where $M$ is a finite-dimensional
smooth manifold (the phase space), $\omega$ is a symplectic or
pre-symplectic form, $H \colon M \times \Theta \to \R$ is a
family of Hamiltonians parameterised by a configuration space $\Theta$,
and $\Gamma$ is a dissipation structure (a Rayleigh dissipation
function, a Lindblad generator, or more generally any structure that
breaks time-reversal symmetry and drives trajectories toward attractors).
In some embodiments, stochastic forcing (thermal noise, shot noise,
quantum fluctuations) is included as part of $\Gamma$. Every \RK module, every reactor, and every scene whose dynamics are
modelled as a controlled dissipative system is an object of
the weak category $\mathbf{DissHam}$.

\paragraph{Morphisms.}
A morphism $\kappa \colon A \to B$ is a physical coupling: a
specification of how system $A$'s state influences system $B$'s
dynamics, parameterised by coupling type (optical, mechanical,
electrical, informational), coupling strength $g \in [0,\infty)$,
coupling topology (which degrees of freedom of $A$ couple to which
degrees of freedom of $B$), and coupling timescale. The identity
morphism is the zero coupling ($g = 0$). Composition of morphisms
$\kappa_1 \colon A \to B$ and $\kappa_2 \colon B \to C$ is defined by
adiabatic elimination of the mediating system $B$ when its dynamics are
fast relative to $A$ and $C$, yielding an effective coupling
$\kappa_2 \circ \kappa_1 \colon A \to C$.

\paragraph{Self-coupling as endomorphism.}
The \RK feedback loop --- wherein a system observes itself
through a measurement function and feeds the observation back into its
own dynamics --- is a nontrivial endomorphism $\kappa_{\mathrm{self}}
\colon A \to A$, distinct from the identity. The system with feedback is
a different dynamical system from the system without feedback, even
though the underlying phase space is the same. This endomorphism is the
formal counterpart of the projector--camera loop in the anchor
embodiment.

\paragraph{Weak associativity.}
Composition is well-defined physically but may not satisfy strict
associativity in the categorical sense, because the adiabatic
elimination operation depends on timescale ordering: coupling $A$ to $C$
through mediator $B$ differs from coupling $A$ to $C$ through mediator
$D$ even if the effective coupling strengths are the same, because $B$
and $D$ have different dynamics that leave different imprints on the
effective coupling. This means $\mathbf{DissHam}$ may be a weak
category (or bicategory) rather than a strict category. The enriched
framework described below is tolerant of this: the structure may be
more precisely described as an enriched bicategory or
tolerance-enriched category (in the sense of Lawvere's generalised
metric spaces), where hom-objects carry metric structure that quantifies
the degree of associativity failure. Formalising this rigorously is an
open problem identified in this section.

\paragraph{Timescale separation condition for composition (non-limiting).}
The adiabatic elimination that defines morphism composition is valid
when the mediating system $B$ relaxes on a timescale $\tau_B$ that
is much shorter than the timescales $\tau_A$ and $\tau_C$ of the
coupled systems: the operative condition is
$\tau_B / \min(\tau_A, \tau_C) < \eta_{\mathrm{sep}}$ for a declared
timescale separation ratio $\eta_{\mathrm{sep}}$ (a non-limiting
heuristic value is $\eta_{\mathrm{sep}} = 0.1$).  When this condition
holds, the composition is well-defined and the tolerance-enriched
distance between the two compositional orderings
$((\kappa_3 \triangleright \kappa_2) \triangleright \kappa_1)$ and
$(\kappa_3 \triangleright (\kappa_2 \triangleright \kappa_1))$ is
bounded in the enriched metric.  When the timescale condition does
not hold --- that is, when $\tau_B \sim \tau_A$ or $\tau_B \sim \tau_C$
--- adiabatic elimination is not valid, the composition must be
replaced by full coupled simulation, and the categorical structure
does not apply; this is the honest domain of validity of the
$\mathbf{DissHam}$ framework.  In some embodiments, the timescale
separation ratio is estimated from the AR(1) drift model parameters
and committed to the protocol digest.

\subsection{Fisher--Rao enrichment}
In some embodiments, $\mathbf{DissHam}$ is enriched over the category of
Riemannian manifolds (with Cartesian product as monoidal structure): the
hom-set between two objects $A$ and $B$ is not merely a set of couplings
but a Riemannian manifold $\mathrm{Hom}(A,B)$ of all possible
couplings, equipped with a metric derived from how sensitively the
coupled dynamics depend on coupling parameters. This metric is the
Fisher--Rao metric on the coupling parameter space. Composition of
morphisms is a smooth map between the relevant Riemannian manifolds, and
the enrichment coherence conditions (associativity and unitality of
composition up to the manifold structure) are satisfied in the
weak-coupling regime where joint parameter spaces factor as products.

The Fisher--Rao metric therefore appears at every level of the
framework:
\begin{enumerate}[nosep]
  \item On $\Theta$ (the configuration space of a single system),
        as defined in Section~\ref{sec:definitions}.
  \item On the space of attractors within a single system's bundle.
  \item On $\mathrm{Hom}(A,B)$ (the space of couplings between
        systems).
  \item On the space of transformations of couplings
        (meta-couplings, as in a higher-order controller's modulation of system--environment
        interactions).
\end{enumerate}
This ubiquity is the signature of a natural structure: the Fisher--Rao
metric is not imposed externally but arises from the statistical
properties of the physical dynamics at each level. The enriched category
is denoted $\mathbf{DissHam}_{\mathrm{FR}}$.

\subsection{The attractor bundle functor}
\label{sec:att-functor}

\paragraph{Definition (attractor bundle assignment; lax monoidal functor).}
The \emph{attractor bundle assignment}
$\mathrm{Att} \colon \mathbf{DissHam}_{\mathrm{FR}} \to
\mathbf{FibBund}_{\mathrm{FR}}$ maps each dissipative Hamiltonian
system to its attractor bundle and maps each coupling to the induced map
on attractor bundles. Here $\mathbf{FibBund}_{\mathrm{FR}}$ is the
category of fibre bundles over parameter spaces equipped with the
Fisher--Rao connection.

On objects:
\[
  \mathrm{Att}(M, \omega, H, \Gamma, \Theta)
  = \bigl(\Theta,\; \{A_\theta\}_{\theta \in \Theta},\;
    g_{\mathrm{FR}},\; \nabla\bigr),
\]
where $A_\theta$ is the set of attractors at parameter value $\theta$,
$g_{\mathrm{FR}}$ is the Fisher--Rao metric on $\Theta$, and $\nabla$
is the connection on the bundle derived from $g_{\mathrm{FR}}$ (the
$e$-connection or $m$-connection of information geometry).

On morphisms: if $\kappa \colon A \to B$ is a coupling, the induced map
on attractor bundles is
$\mathrm{Att}(\kappa) \colon \mathrm{Att}(A) \times \mathrm{Att}(B) \to
\mathrm{Att}(A \otimes_\kappa B)$, mapping the individual attractor bundles
to the joint attractor bundle of the coupled system, where
$A \otimes_\kappa B$ denotes the coupled system under coupling $\kappa$.
Note that this is not the standard type signature for a functor's action
on morphisms ($F(f):F(A)\to F(B)$); the product domain and
coupled-composite codomain reflect the fact that coupling creates a
\emph{joint} system whose attractor structure depends on both partners.
The map is more precisely a component of the monoidal product structure
(Section~\ref{sec:att-functor}, monoidal product paragraph below)
rather than a strict functorial action on hom-sets.  In some
embodiments, $\mathrm{Att}$ is treated as a lax monoidal functor whose
coherence conditions are satisfied up to the Fisher--Rao distortion
bound of Conjecture~1.

\paragraph{Status and open mathematical questions.}
The functor is well-defined on objects (existence of attractors for
dissipative systems and the Fisher--Rao metric on parameter spaces are
both well-established). Its action on morphisms is well-defined in the
weak-coupling limit, where joint attractors are perturbations of
individual attractors. For strong coupling, the joint attractor structure
can be qualitatively different from either individual's (new attractors
appear, old ones vanish via bifurcation), and the functorial property
--- that composition of couplings maps to composition of attractor-bundle
maps --- is physically plausible (attractors of the coupled system do not
depend on the order in which couplings were activated, in the long-time
limit) but has not been proven in full generality. A known mathematical
obstruction is that attractors change discontinuously at bifurcation
points --- they appear, disappear, change dimension, or undergo
explosive transitions --- so the attractor bundle has singular fibres
precisely where the curvature of the Fisher--Rao connection diverges.
This is consistent with the physical picture (Fisher
information does diverge at phase transitions) but means the functor
$\mathrm{Att}$ maps into a category of singular bundles rather than
smooth bundles. Possible resolutions include restricting
$\mathrm{Att}$ to structurally stable regions and treating bifurcation
loci as boundary strata; extending via the Conley index (which provides
continuation through bifurcations but yields algebraic invariants rather
than bundle fibres); or working with the space of invariant measures
(which varies continuously even through bifurcations, at the cost of
losing geometric structure). The correct resolution is an open
mathematical question whose answer likely depends on whether the
intended application is local (near a fixed configuration) or global
(across the full bifurcation diagram).

\subsection{The five substrate mappings as functors}
Each of the five preceding substrate sections defines, in categorical
language, a functor from $\mathbf{DissHam}_{\mathrm{FR}}$ (or a
subcategory thereof) to a target category specific to that domain:

\begin{center}
\small
\begin{tabular}{lll}
  \toprule
  \textbf{Section} & \textbf{Functor} & \textbf{Target category} \\
  \midrule
  Optical Primitives
    & $F_{\mathrm{opt}}$
    & Category of parameterised optical channels \\[0.3em]
  Neural-Architecture Mapping
    & $F_{\mathrm{nn}}$
    & Category of recurrent network architectures \\[0.3em]
  Information-Theoretic Primitives
    & $F_{\mathrm{info}}$
    & Category of noisy channels with Fisher metric \\[0.3em]
  Control-Theoretic Primitives
    & $F_{\mathrm{ctrl}}$
    & Category of controlled dynamical systems \\[0.3em]
  Cryptographic-Primitive Substrate
    & $F_{\mathrm{crypto}}$
    & Category of interactive proof systems \\
  \bottomrule
\end{tabular}
\end{center}

The repeated appearance of the same formal structures (attractors,
bifurcations, Fisher information, holonomy) across all five substrate
mappings is consistent with this framing: under the proposed categorical
encoding, these structures arise as properties of the source functor
$\mathrm{Att}$, and if the relevant preservation properties of the
specific functors $F_{\mathrm{opt}},\ldots,F_{\mathrm{crypto}}$ are
established, they would be inherited by all five images
rather than coincidentally shared.  Whether these five mappings
are indeed faithful functors and whether they actually preserve the
relevant properties of $\mathrm{Att}$ in the senses needed are open
questions (see the status paragraph above); the framing is offered as
an organising hypothesis, not a proven theorem.

\subsection{Monoidal structure and composition}

In some embodiments, $\mathbf{DissHam}_{\mathrm{FR}}$ has additional
algebraic structure relevant to \RK composition
(Section~\ref{sec:networks}):

\begin{description}[style=nextline, leftmargin=2em]
  \item[Monoidal product ($\otimes$).]
    Parallel composition of systems: given $A$ and $B$, the product
    $A \otimes B$ has phase space $M_A \times M_B$, parameter space
    $\Theta_A \times \Theta_B$, and independent dynamics (no coupling).
    This is the categorical counterpart of the parallel multi-reactor
    topology.
  \item[Traced monoidal structure.]
    Feedback loops are formalised by the trace operation: given a
    morphism $f \colon A \otimes C \to B \otimes C$, the trace
    $\mathrm{Tr}^C(f) \colon A \to B$ feeds $C$'s output back as
    input, yielding the closed-loop system. The \RK feedback
    loop (projector $\to$ scene/reactor $\to$ detector $\to$ controller
    $\to$ projector) is exactly the trace over the optical channel.
  \item[Braiding.]
    In some embodiments, the monoidal product is symmetric (the order
    of parallel composition does not matter up to canonical isomorphism),
    consistent with the physical symmetry that two non-interacting
    systems can be labelled in either order.
\end{description}

The multi-reactor architectures described in the PoliePuter section
--- chain, parallel, graph, hierarchical, recurrent, and
reservoir-of-reservoirs --- correspond to different diagrams in the
traced monoidal category, and the multi-reactor primitives (Route,
Compile, Parallelise, Reduce, Error-correct, Abstract) are natural
transformations between these diagrams.

\subsection{The self-modelling endofunctor and fixed-point conjectures}
\label{sec:self-model-functor}

\paragraph{Definition (self-modelling endofunctor).}
The \emph{self-modelling endofunctor}
$\mathrm{Mod} \colon \mathbf{DissHam}_{\mathrm{FR}} \to
\mathbf{DissHam}_{\mathrm{FR}}$ maps each system $A$ to its best
internal model of itself: the subsystem of $A$ whose dynamics track
$A$'s own attractor structure through the feedback loop (the
self-coupling endomorphism $\kappa_{\mathrm{self}}$). This is the formal
counterpart of the depth curriculum described in
the related Filing 2 application (optional, non-essential), where a
PolieBot iteratively refines its
model of its own coupling boundaries.

On morphisms: given a coupling $\kappa \colon A \to B$, the induced
morphism $\mathrm{Mod}(\kappa) \colon \mathrm{Mod}(A) \to \mathrm{Mod}(B)$
is the coupling between the respective self-model subsystems induced by
$\kappa$---that is, the restriction of $\kappa$ to the self-modelling
degrees of freedom of $A$ and $B$. This is a non-limiting definition;
the precise form depends on how the self-model subsystem is demarcated
within each system's phase space.

\paragraph{Conjectures (open problems, non-limiting).}
The following conjectures are stated as a theoretical framework. They
are specific enough to be tested, proven, or refuted, and they identify
precise mathematical claims that would, if established, unify the
transferability and self-modelling aspects of the \RK
formalism. Proving them would require rigorous construction of
$\mathbf{DissHam}_{\mathrm{FR}}$ as an enriched category (including the
weak-associativity issue), proof that $\mathrm{Att}$ preserves enriched
structure, and analysis of the fixed-point structure --- work that is
beyond the scope of this disclosure but is identified here as a research
programme.

\begin{description}[style=nextline, leftmargin=2em]
  \item[Conjecture~1 (Transferability).]
    $\mathrm{Att}$ is an enriched functor: it preserves the Fisher--Rao
    metric structure on hom-objects, up to bounded distortion. The
    distortion bound depends on the topological distance between the
    attractor bundles of the source and target systems (specifically, on
    the difference in their bifurcation-set topology). If
    $\mathrm{Att}$ preserves enriched structure, then: nearby couplings
    produce nearby attractor-bundle maps (continuity of transfer); the
    metric distortion under $\mathrm{Att}$ is bounded (transfer quality
    is predictable); and the kernel of $\mathrm{Att}$ identifies
    substrate-independent equivalence classes (couplings that produce the
    same attractor-bundle map regardless of substrate).
  \item[Conjecture~2 (Self-consistent self-model, existence).]
    For any object $A$ in $\mathbf{DissHam}_{\mathrm{FR}}$ satisfying
    (a)~$A$'s phase space $M_A$ has sufficient dimensionality to embed a
    model of its own attractor bundle, and (b)~$A$'s dynamics are
    capable of universal computation (the system is dynamically
    Turing-complete), the endofunctor $\mathrm{Mod}$ has a fixed point:
    there exists $A^\ast \cong \mathrm{Mod}(A^\ast)$.
  \item[Conjecture~3 (Self-consistent self-model, structure).]
    The fixed point $A^\ast$ necessarily has an attractor bundle
    $\mathrm{Att}(A^\ast)$ containing structures isomorphic to:
    non-trivial attractors; a bifurcation set of positive codimension;
    positive entropy production (irreversibility); noise-driven
    selection at bifurcations; self-referential feedback modifying the
    attractor structure; and coupling-dependent boundary dissolution.
    The self-model at the fixed point contains representations of all
    six structures (the model correctly describes its own attractors,
    bifurcations, irreversibility, noise, self-reference, and
    boundary-dependence).
  \item[Conjecture~4 (Convergence via depth curriculum).]
    Under a suitable training curriculum (the depth curriculum of
    the related Filing 2 application; optional, non-essential), $\mathrm{Mod}$ is a contraction
    mapping on $\mathbf{DissHam}_{\mathrm{FR}}$, and iterated
    application of $\mathrm{Mod}$ converges to the fixed point $A^\ast$
    at a rate determined by the contractivity constant. In the enriched
    setting, this is a Banach fixed-point theorem: each training
    iteration brings the self-model closer to self-consistency (the
    Fisher--Rao distance between the system and its self-model
    decreases), and the convergence rate is the contraction ratio of
    $\mathrm{Mod}$.

    \textit{Note on viability.}  Conjecture~4 as stated is contested
    and is retained as an open research problem rather than an
    operational claim.  Three structural objections have been
    identified.  (1)~\textbf{Capability expansion.}  A better
    self-model reveals previously unmodelled capability dimensions,
    requiring a strictly more expressive model to represent them.
    $\mathrm{Mod}$ is expansive by construction in parameter space;
    self-modelling discovers new content rather than converging to a
    pre-existing fixed point.  (2)~\textbf{Two-speed dynamics.}
    Self-modelling simultaneously compresses stale structure and
    expands into newly discovered structure.  The net effect on
    Fisher--Rao distance to any putative fixed point is indeterminate;
    a clean contraction is not assured.  (3)~\textbf{Co-evolutionary
    dynamics} (most decisive).  Any realistic operating system models
    interactions with other self-modelling systems that are themselves
    running $\mathrm{Mod}$ and becoming more complex.  There is no
    single fixed point in a fleet of co-evolving agents; the joint
    fixed point of simultaneous fleet-wide self-modelling is a far
    stronger requirement than stated and is unlikely to exist
    (Red Queen dynamics).  The restricted version---locally contractive
    for an isolated system in a fixed environment near a structurally
    stable fixed point---avoids these objections but is too weak to be
    interesting and too far from operating conditions to be useful.
    \textbf{No core claim in this specification depends on
    Conjecture~4.}  The work it was intended to formalise is performed
    instead by the PoD architecture: physical anchoring bounds
    complexity growth; the security/discovery domain partition prevents
    self-modelling from eroding verification stability; quorum
    structure bounds individual divergence.  The practical
    proposition---that the PoD architecture keeps fleet-level
    self-modelling in a bounded regime by design---is independently
    grounded and requires no categorical machinery.
  \item[Conjecture~5 (Universality of natural transformations).]
    The normalised transfer-entropy matrix, synchronisation order
    parameter, and Fisher information profile are enriched natural
    transformations of $\mathrm{Att}$. They are substrate-independent in
    the precise sense that they commute with all morphisms in
    $\mathbf{DissHam}_{\mathrm{FR}}$, up to the distortion bound from
    Conjecture~1.
\end{description}


Even if one or more categorical correspondences hold only approximately
or fail for a given embodiment, the disclosed apparatus, calibration
loops, drift navigation, discrepancy protocol, and network corroboration
remain supported as independently described engineering mechanisms.

\paragraph{Relationship to Lawvere's fixed-point theorem.}
Conjecture~2 is related to, but distinct from, Lawvere's categorical
fixed-point theorem, which establishes that in a cartesian-closed
category with a point-surjective morphism from an object to its
exponential, every endomorphism has a fixed point.
$\mathbf{DissHam}_{\mathrm{FR}}$ is unlikely to be cartesian-closed in
the strict sense (the exponential object $B^A$ --- the ``space of all
ways $A$ can influence $B$'' --- may not exist as an object in the
category). The enriched formulation via Banach contraction would avoid this
difficulty: it requires only that the hom-objects are complete metric
spaces and that $\mathrm{Mod}$ is a contraction, which \emph{if it
holds} gives constructive convergence rather than mere existence.
However, as noted in the viability note to Conjecture~4, the
contraction of $\mathrm{Mod}$ is contested on structural grounds
(capability expansion, two-speed dynamics, and co-evolutionary Red
Queen dynamics).  If those objections apply, the Banach-contraction
route to Conjecture~2 is also unavailable; existence of a fixed point
$A^\ast \cong \mathrm{Mod}(A^\ast)$ would then require a different
proof strategy---for example, topological fixed-point methods
(Schauder, Brouwer, or Kakutani) that require compactness and
continuity but not contraction, and which do not provide constructive
convergence.  Whether such a strategy is available for
$\mathbf{DissHam}_{\mathrm{FR}}$ is an open question.  The
connection to Lawvere remains of theoretical interest because it links
self-reference in physical systems to the same diagonal-argument
structure underlying G\"{o}del's incompleteness theorems, the
halting problem, and Cantor's theorem.

\paragraph{Practical implications for Reality Kernel design.}
If Conjectures~1, 2, 3, and~5 are approximately correct, they have
direct engineering consequences: the drifting field
(Section~\ref{sec:drifting-field}) navigates the Fisher--Rao geometry of
$\Theta$ using the structure that $\mathrm{Att}$ preserves; few-shot
commissioning works because $\mathrm{Att}$'s continuity (Conjecture~1)
implies that nearby configurations produce nearby attractor bundles; and
the five substrate mappings give consistent results because they are
all faithful functors out of the same source category.  These
conjectures provide a theoretical explanation for design choices that
are already empirically motivated elsewhere in this specification.

Conjecture~4 (depth curriculum convergence via contraction of
$\mathrm{Mod}$) is \emph{not} included in the above list.  As noted
in the viability note to Conjecture~4, the contraction claim is
contested on three structural grounds (capability expansion,
two-speed dynamics, and co-evolutionary Red Queen dynamics) and
is recorded as an open and probably hard research problem.  The
empirical observation that depth curricula converge in practice does
not depend on Conjecture~4; it is an operational mechanism (independently monitored by
dedicated metres, Section~\ref{sec:drifting-field}) and requires no
categorical machinery.  Conjecture~4 offers a \emph{possible}
theoretical explanation if it is eventually proved in some restricted
form; it is not a precondition for the observation.

This paragraph is conditional throughout on the remaining conjectures.
Conjecture~2 applies only to embodiments whose reactor dynamics satisfy
a Turing-completeness condition, which is an additional design target
not achievable with all reactor types; verifying this condition for
specific embodiments is a research target and is not required for the
operational mechanisms claimed in this disclosure.

% ======================================================================
\section{Quantum and Non-Classical Coupling (Optional, Non-limiting)}
\label{sec:quantum}
% ======================================================================

In quantum-sensitive embodiments, the disclosed architecture extends the
same committed-bundle, protocol-digest, and multi-regime operator framework
to photon-counting, interferometric, and related non-classical channels,
thereby allowing quantum-sensitive observables to participate in the same
verification, sensing, and controlled-transformation workflows as the
classical embodiments.

In such embodiments, quantum Fisher information provides an additional
metric on the device configuration manifold that can be used to select
protocols enlarging the gap between device-achievable and
attacker-achievable estimation precision under declared budgets.

This section describes non-limiting embodiments that incorporate quantum-sensitive or non-classical effects while preserving the core \RK structure: a swept, controlled operator with configuration parameters $\theta$, a control protocol $U_{0:T}$, and a \cba record $C_{0:T}$ comprising emission-observation bundles. These embodiments are optional enhancements. Core \TB, \LI, and \RT operation may be implemented in purely classical regimes.

\subsection{Overview and quantum classes}
In some embodiments, quantum-sensitive channels are included as additional components of the observation vector $\mathbf{y}_t$. These channels include, without limitation, photon counts, arrival-time distributions, coincidence statistics, homodyne or heterodyne quadrature measurements, interference fringe statistics, and other observables derived from quantum-sensitive detection.

In some embodiments, coincidence statistics include time-tagged correlations and second-order intensity correlations derived from click streams. A non-limiting example is a normalised correlation function $g^{(2)}(\tau)$ computed from time tags or binned counts, which can be treated as a derived channel and logged or committed as part of \cba summaries.

For descriptive convenience, non-limiting quantum-sensitive embodiments are grouped into quantum classes:
\begin{itemize}
  \item \textbf{Q0:} single-photon or few-photon operation using photon counting without requiring entanglement.
  \item \textbf{Q1:} coherent interferometric operation with phase-sensitive observables.
\end{itemize}

\paragraph{Quantum-flavoured hardness metrics (optional, non-limiting).}
In some embodiments, hardness is reported under explicit photon-budget or coherence-budget constraints. Non-limiting auxiliary metrics (distinct from the canonical $(k^*_{\mathrm{dig}}, m^*_{\mathrm{ana}})$ hardness index) include: a minimum photon budget (or minimum integration time) required for a bounded attacker model to meet a fidelity threshold $\tau$ under a declared protocol; and residual mismatch in interferometric statistics (for example visibility or phase-estimation error) under a declared phase sweep and calibration regime. These metrics are empirical, conditioned on meter envelopes and protocol digests, and are optional enhancements rather than requirements for core operation.

\subsection{Quantum class Q0: Single-photon and photon-counting operation}
In Q0 embodiments, a \RK module operates in a photon-starved regime and one or more detectors record photon-count statistics. Sources include, without limitation, attenuated coherent sources and heralded single-photon sources. Detectors include, without limitation, single-photon avalanche diodes and other photon-counting detectors. The system treats resulting click streams and derived statistics as additional channels in the \cb.

\paragraph{Photon-count observation model (non-limiting).}
In some embodiments, a detector channel returns a photon count per step (or per window) that is modelled as a conditional count process. A non-limiting model treats the count $N_t$ (or binned count vector $\mathbf{N}_t$) as
\[
  N_t \;\sim\; \mathrm{Poisson}\!\bigl(\lambda_\theta(S,u(t))\bigr),
\]
where $\lambda_\theta$ depends on the scene, reactor state, and control settings including, without limitation, wavelength, intensity envelope, scan coordinate, and detector gate $g_{\mathrm{det}}(t)$ (integration time and/or time-bin settings). These counts (or derived statistics such as mean/variance over repeated micro-sweeps) are included as components of $\mathbf{y}_t$ and therefore of the \cb.

\paragraph{Time-tagging readout (non-limiting).}
In some embodiments, the photon-counting channel is time-resolved by recording time tags for click events relative to an emitted pulse or modulation cycle (for example time-correlated single-photon counting (TCSPC)). In these embodiments, the per-step observable may be a list of time tags or a time-bin histogram $\mathbf{h}_t$. In some embodiments, the controller sweeps gate widths, bin edges, pulse repetition rates, or modulation frequencies as part of $U_{0:T}$. Time-resolved channels can strengthen both \LI (depth/delay inference) and \TB (fine-grained timing signatures conditioned on metered synchronisation and drift).

\paragraph{Temporal modulation bridge (non-limiting).}
In some embodiments, the classical-regime temporal modulation of
epsilon-near-zero reactor media (Section~\ref{sec:embodiments})
provides a natural bridge to quantum-regime temporal operation.  At
classical intensities, periodic temporal modulation produces Floquet
sidebands; in the photon-starved regime, these correspond to
discrete photon-energy channels whose statistics can be resolved by
TCSPC or spectral detection.  The same physical mechanism therefore
supports both a classical temporal-hardness channel and a
quantum-regime timing-statistics channel, logged in the same
\cba format.

In some embodiments, Q0 operation provides hardness amplification because matching full photon-count statistics across a disordered reactor may be more demanding than matching mean intensities. In some embodiments, shot noise and arrival processes provide entropy sources for protocols. In some embodiments, Q0 supports low-light \LI operation for photosensitive scenes or covert sensing.

\subsection{Quantum class Q1: Coherent interferometric operation}
In Q1 embodiments, the system includes phase-stable interferometric structure such as Mach--Zehnder or Sagnac geometries. Phase modulators and optional delay elements are treated as components of the control protocol $U_{0:T}$. In some embodiments, the scan protocol sweeps phase settings jointly with spatial scan coordinates, producing a phase-resolved \cb.

Outputs such as visibility, phase estimates, and fringe statistics are included as channels of $\mathbf{y}_t$. In some embodiments, calibration sub-protocols maintain phase stability using reference tones or interleaved calibration segments, and the resulting calibration data is logged and optionally committed together with the \cb.

In some embodiments, non-classical sources or stronger quantum constraints (for example squeezed-light enhancements) are used as optional sensitivity amplifiers; such variants are not required for core operation and are not further detailed in this section.

\subsection{Hybrid classical--quantum training and optimisation}
In quantum-sensitive embodiments, the device interface remains classical: detectors output classical data (counts, time tags, quadratures). Meters and decoders are trained on these channels as they would be on any other channel.

Optimisation of quantum-control settings and certain reactor parameters may use gradient-free methods, including without limitation SPSA, evolution strategies, and bandit methods, and may also use score-function estimators. In embodiments with small parametrised interferometers, analytic gradient estimators may be used for those limited parameters while treating the rest of the \RK as an environment.

\subsection{Quantum-regime detector-as-reactor operation (non-limiting)}

In some embodiments, the detector-as-reactor mode described in
Section~\ref{sec:detector-reactor} operates in a regime where quantum
noise sources --- shot noise from photon arrival statistics and, at
very low photon counts, vacuum fluctuations --- are the dominant
contribution to the per-pixel readout variance. In this regime, each
pixel's output is a quantum measurement outcome whose statistics are
set by the photon number distribution of the incident field, the
detector's quantum efficiency (which varies per pixel due to
fabrication), and the electronic noise floor.

In some embodiments, this quantum-noise regime has consequences for both
security and sampling quality. For security: in the photon-starved
regime, and under a declared adversary access model committed to the
protocol digest --- specifying limits on the adversary's verification
photons, calibration photons, observation time, illumination access,
detector access, and observability of calibration ground truth --- an
adversary attempting to characterise the detector's per-pixel response
faces a measurement-budget constraint. Fully characterising the
quantum efficiency, dark current, and timing jitter of each pixel to a
target precision $\delta$ in the shot-noise-limited regime requires a
photon budget per pixel scaling as $\Omega(1/\delta^2)$ under
independent classical or coherent probing; stronger
quantum-metrologically enhanced scaling toward Heisenberg-type
$1/N$-precision is possible in principle for some parameters but
requires nonclassical probe states and compatible measurements and is
not assumed here. Each characterisation measurement consumes part of
the adversary's finite photon and time budget, so probe photons
expended on early pixels are unavailable for later ones; the readout
record is altered in the count-statistics sense (additional samples
recorded) rather than in the wavefunction-collapse sense, and the
security advantage is operational (measurement-budget-limited) rather
than information-theoretic (no-cloning). The advantage is strongest
when the verification protocol uses fewer photons per pixel than would
be needed for full per-pixel characterisation by the adversary under
the declared access model. Adversaries operating outside the declared
access model --- for example with unlimited time, unlimited
illumination access, or full observability of calibration ground
truth --- are not within the scope of this warrant. For sampling: at
very low photon counts, the readout statistics are sub-Poissonian or
super-Poissonian depending on the source characteristics, providing
additional distributional structure that a classical noise model does
not capture.

In some embodiments, the natural metric on the parameter space is the
quantum Fisher information $\mathcal{F}_Q(\theta)$. The quantum Fisher
information $\mathcal{F}_Q$ defines a Riemannian metric on the
manifold of quantum states (the Bures metric, up to a multiplicative
constant; see Provost and Vall\'ee 1980, Braunstein and Caves 1994),
and for a scalar parameterisation the inequality
$\mathcal{F}_Q \geq I(\theta)$ (where $I(\theta)$ is the classical
Fisher information obtained from a particular measurement) holds
under standard regularity conditions on the parameterisation and the
measurement family, with equality attained for a measurement
saturating the symmetric logarithmic derivative (SLD) bound; the
classical limit, in which the relevant family of states commutes and a
common-eigenbasis measurement reduces the problem to classical Fisher
information, is one special case of equality. For multi-parameter
$\theta$, the corresponding matrix inequality
$\mathcal{F}_Q(\theta) \succeq \mathcal{F}(\theta)$ holds for fixed
POVMs, but simultaneous attainment of the multiparameter quantum
Cram\'er--Rao bound is not automatic and is obstructed in general by
non-commuting SLDs and incompatible optimal measurements (Ragy,
Jarzyna, and Demkowicz-Dobrzański 2016, and references therein); the
matrix inequality is used herein as design guidance rather than as an
operative attainable bound. In the quantum-noise regime, the scalar
gap $\mathcal{F}_Q(\theta) - I(\theta)$ (or, in the multi-parameter
case, a scalar functional such as
$\mathrm{tr}[\mathcal{F}_Q - \mathcal{F}]$) is informative when
attainable: protocols supplying the necessary nonclassical resources
and measurements may achieve estimation precision closer to the
quantum Cram\'er--Rao bound than classical protocols, subject to loss,
detector noise, calibration uncertainty, and measurement
compatibility. The attractor bundle formalism extends to this regime
through the connection \emph{induced by} the QFI/Bures metric (in
particular, the Levi-Civita connection of the pullback of the
QFI/Bures metric to the device configuration manifold) in place of the
connection induced by the classical Fisher--Rao metric, and the
curvature of the resulting connection is conjectured to diverge at
bifurcation singularities in the same manner as in the classical case
(this is a quantum-regime extension of the curvature-divergence
property described in Section~\ref{sec:att-functor}; its validity is
an open question related to Conjecture~1). The novelty here is not
the quantum Fisher information itself (which is a standard tool in
quantum estimation theory) and not the QFI/Bures geometry on quantum
state space (which is established; see Provost and Vall\'ee 1980 and
the geometric-phase literature, including Berry 1984 and Simon 1983,
for adjacent fibre-bundle constructions), but the use of the
QFI-induced connection as the connection on an attractor bundle over
a physical device's configuration space, used as a design metric for
selecting protocols that maximise the gap between the device's
achievable precision and an attacker's achievable precision in both
\LI and \TB operation.

\subsection{Optional enhancements and non-limiting notes}
Quantum embodiments are optional and may be mixed with classical channels. In some embodiments, a device supports both classical and photon-counting operation by switching sources and detectors while preserving \cba logging conventions. The Q0/Q1 classification above is used for description only and does not limit combinations or hybrid operation.

\paragraph{Educational-kit and practical-lab quantum embodiments
  (non-limiting).}
In some embodiments, quantum-sensitive operation is implemented at
educational, hobbyist, or practical-lab scale rather than at
full-laboratory scale. Such embodiments comprise without limitation:
attenuated coherent probes derived from continuous-wave or pulsed
laser sources at declared mean-photon-per-pulse ranges (for example
$10^{-2}$ to $1$ photon per pulse, with attenuation calibrated by
committed neutral-density filters or active attenuators); heralded
single-photon sources implemented by spontaneous parametric
down-conversion crystals (for example beta-barium-borate, periodically
poled lithium niobate, or periodically poled potassium titanyl
phosphate) pumped by visible or near-infrared lasers, with herald
events provided by coincidence detection on the idler arm; avalanche
photodiode and single-photon-avalanche-diode detectors operated in
Geiger mode with declared dark-count rates, dead times, and timing
jitter; time-correlated single-photon counting acquisition with
declared picosecond-to-nanosecond timing resolution; spatial-
superposition branches separated by wave plates, polarising beam
splitters, dichroic mirrors, or Mach--Zehnder or Sagnac
interferometric geometries with declared phase-stability budgets;
multi-photon extensions including GHZ-state, cluster-state, and
W-state preparations by successive entangling operations; and
optional stabilisation comprising pilot-tone phase reference,
dispersion compensation in optical-fibre interconnects, and active
thermal or mechanical feedback. Each such practical-lab embodiment
operates under the same committed-evidence, protocol-digest,
verification, and hardness discipline as full-laboratory variants of
the present section, with per-embodiment parameter ranges,
calibration tables, and stabilisation procedures recorded in the
protocol digest.

% ======================================================================


\paragraph{Scope limitation (non-limiting).}
The embodiments of this section support declared sensing, attestation,
and hardness tasks benchmarked under the filing's task procedures. They
do not support, and shall not be construed to support, capacity or
scaling claims of the kinds bounded by
Section~\ref{sec:capacity-accounting-rule}, and do not imply any
increase in capacity quantities.


\section{Methods of Operation (Non-limiting)}

\subsection{Claim-relevant agent-integration embodiments (non-limiting)}

The following embodiments are non-limiting examples relevant to system
and method claims:
\begin{enumerate}
  \item A method for controlling an agent using a \RK as a
    primary observation channel, including updating belief state from
    the \cb and selecting actions that include \RK
    control signals and platform actuation.
  \item A method for verification-gated action execution in which an action
    class is permitted when configured policy thresholds on \RK
    evidence are met, said thresholds comprising one or more of
    verisimilitude, hardness, freshness, or multi-RK corroboration.
  \item A method for offline agent learning using at least one of
    (i)~an RK-calibrated emulator and (ii)~a physical \RK
    operating on synthetic or modified scenes, including periodic
    physical validation probes and confidence-weighted updates from
    simulated experience.
  \item A method for probing an agent internal representation by
    externalising a hypothesis as a physical emission protocol and
    comparing the measured \cb to the predicted \cb.
\end{enumerate}

% ======================================================================


% ======================================================================
\section{Additional Human, Deniability, Physical-Layer, and Agent-Layer Extensions (Non-limiting)}
% ======================================================================


\paragraph{Positioning relative to proof-of-personhood systems
(non-limiting).}
In some embodiments, the human-layer records described below are used
in application spaces adjacent to proof-of-personhood. The
embodiments below share a high-level goal with generic
proof-of-personhood systems: Sybil-resistant human-anchored
attestation as a substrate for downstream voting, governance,
reputation, and resource-allocation applications. They are
distinguished from such systems in at least six architectural respects.
First, the load-bearing evidence object is a committed human--fleet
coupled-trajectory history, not a one-time assertion that a person is
globally unique. Second, the record is derived from HC-Read, HC-Write,
and HC-Verify channels bound into \cba evidence, rather
than from an unqualified biometric, document, account, or credential
alone. Third, recognition and quota consequences are scoped to a
declared fleet or subnet policy catalogue, rather than to a universal
personhood registry. Fourth, token emission and quota accounting are
rate-limited by declared cadence, freshness, meter envelopes, and
interaction windows, rather than by an assumed one-human-one-account
property. Fifth, portability is evidentiary rather than dispositive: a
receiving fleet or subnet applies its own declared thresholds to an
existing portable record and any fresh session it elects to run.
Sixth, adversarial incentives are directed at detector hardening
through declared Eve-style attempts against subnet detectors and
interaction surfaces, not at probing applicants or extracting
border-admission evidence. Accordingly, the embodiments below may be
combined with personhood-adjacent use cases, but their load-bearing
claims remain committed physical evidence, scoped policy, physically
rate-limited outputs, and declared attacker-model accounting.


\subsection{Portable Human Bonding Records (Non-limiting)}
\label{sec:portable-human-bonding}

\begin{definition}[Portable human bonding record (non-limiting)]
A \emph{portable human bonding record} is a peer-portability record
whose authenticated principal is not a device alone, and not a human
identity assertion alone, but a committed human--fleet
coupled-trajectory history. In some embodiments, the record is built
from a cumulative sequence of HC-Read, HC-Write, and HC-Verify
sessions, with auxiliary human-response channels logged as
\cba evidence under the protocol digest and commitment
infrastructure described in Section~\ref{sec:two-seed-protocol}.
\end{definition}

In some embodiments, enrolment comprises $N$ committed human-coupled
sessions executed under declared protocols and meter envelopes. The
resulting human-response signature is derived from \cba
evidence, including committed auxiliary response channels. Non-limiting
response channels include gaze-pattern trajectories, reaction-time
distributions, task-performance fingerprints, physiological-channel
coherence, haptic timing, voice-timing features, and cross-channel
consistency among such signals. In some embodiments, the enrolment
record stores commitments to the session atoms, protocol digests,
meter summaries, time anchors, and any verifier-selected openings
rather than a raw biometric archive.

In some embodiments, verification is performed by any fleet member
that runs a fresh HC session under a declared protocol and compares
the fresh \cba response evidence against the committed
human-response signature within a declared tolerance envelope. In
some embodiments, a passing comparison refreshes the bond, opens the
declared quota window for that fleet or subnet, and appends the fresh
session to the portable record. In some embodiments, a failing
comparison is logged as advisory evidence, causes a downgrade to a
lower bonding level, or requires re-enrolment under the receiving
fleet's declared policy.

In some embodiments, bonding level is graduated by analogy to the
verification levels of Section~\ref{sec:graduated-verification}.
Level~0 is single-session recognition under a low quota and short
freshness window. Level~1 is sustained closed-loop HC history across
multiple sessions under stable meter envelopes. Level~2 is deeper
enrolment including Yoked human-coupled trajectories and
cross-channel coherence checks. Level~3 is a declared full-quota
bonding level requiring the fleet's deepest accepted HC trajectory
class, optional selective openings, and the cadence specified by that
fleet or subnet. Thresholds, session counts, tolerances, and quota
functions are declared per fleet or subnet; no universal
human-identity threshold is asserted.

In some embodiments, a portable human bonding record includes a
staleness window. In some embodiments, a bond not refreshed within
the declared period is treated as stale and falls back to Level~0,
advisory status, or re-enrolment, depending on the declared policy.
In some embodiments, freshness is indexed by the last accepted HC
session, the time anchor of the most recent commitment, and the
meter-envelope status of the session that refreshed the bond.

\paragraph{Scope of human-response signatures (non-limiting).}
Human-response-signature distinctiveness is weaker than optical PUF
uniqueness and is not treated as device-grade unclonability. The
relevant threat family includes twin-class biometric overlap, relay
attacks in which one human rotates across multiple bonded PolieBots,
coerced bonding, coached or rehearsed responses, sensor spoofing, and
synthetic-response attacks generated from observed auxiliary
channels. Longer HC trajectories, Yoked response windows, and richer
auxiliary channels may separate some members of this family under
declared attacker bounds, but there exists a minimum interrogation
complexity below which two humans or two response histories are
functionally equivalent. This is not a failure of the framework; it
bounds the claim. The resolution floor is the cross-fleet tolerance
envelope under the declared attacker family, cadence, and meter
state, and the bond is time-indexed rather than unconditionally
durable.

\paragraph{Plenoptic-parallax liveness primitive (non-limiting).}
In some embodiments, an HC session uses a plenoptic (light-field)
detector to capture the human participant's face or interaction
surface in a single exposure as an auxiliary liveness meter, not as a
stored human-response signature and not as a device-grade identity
primitive. The session admits the interaction only if the captured
four-dimensional light field exhibits angular-parallax relationships
consistent with a physically present three-dimensional subject or
interaction surface under the declared pose, distance, and
illumination conditions. In some embodiments, the angular-parallax
residual is committed as a liveness component of the session meter
envelope, and sessions that fall below a declared
parallax-consistency threshold are downgraded or refused before any
bond-updating step. In some embodiments, the plenoptic liveness check
is composed with the other meters described in this subsection rather
than substituted for them; it geometrically bounds specific planar
spoofing modes (ordinary flat-display replay lacking calibrated
angular multiplexing, printed portrait, planar projection, and planar
retro-reflective replay) and does not replace coached-response,
relay-attack, or synthetic-response mitigations addressed elsewhere. The parallax residual is not asserted
to provide optical reactor-microstructure-grade unclonability, biometric uniqueness, or
human-signature distinctiveness; any human-response signature remains
derived from the committed \cba evidence under the weaker
human-layer scoping stated in this subsection.

\paragraph{Cross-fleet promiscuity (non-limiting).}
In some embodiments, human bonding is promiscuous by default: a human
may hold parallel bonded records with multiple fleets, and each fleet
tracks only its own scoped HC history, quota cadence, and tolerance
envelope. In some embodiments, the portable record may be offered to
a receiving fleet as evidence, but the receiving fleet applies its
own declared thresholds to the existing committed record and any
fresh HC session it elects to run. Subnet-level exclusivity,
disclosure obligations, and cross-subnet accounting are handled by
the subnet policy catalogue described in
Section~\ref{sec:subnet-token-economies}.


\subsection{Subnet-Scoped Identity and Attested-Interaction Token
Economies (Non-limiting)}
\label{sec:subnet-token-economies}

\begin{definition}[Subnet identity scope (non-limiting)]
A \emph{subnet} is a declared identity scope over a fleet, contract
set, or interaction set, associated with a committed policy
catalogue. In some embodiments, the catalogue includes at least:
(i) an identity-exclusivity policy; (ii) a quota function mapping
attested interaction time or accepted evidence windows to emission
rate; (iii) a source/sink contract catalogue with per-contract
pricing; and (iv) a cross-subnet policy specifying whether
proof-of-projection or tPoP records from another subnet count as
evidence, and whether B-PoD handoffs under Section~\ref{sec:bpod}
are permitted.
\end{definition}

In some embodiments, the identity-exclusivity policy is selected from
a declared finite catalogue. Non-limiting examples include: exclusive
within one subnet; mutually exclusive across a declared family of
subnets; disclosure-required across subnets; and isolated, in which
cross-subnet correlation is forbidden except under an expressly
committed audit or dispute policy. In some embodiments, the policy is
committed in the protocol digest or in an append-only policy
catalogue referenced by the subnet's token records.

In some embodiments, token emission is a subnet-scoped, physically
rate-limited PoP or tPoP export denominated in the subnet's own unit.
In some embodiments, the token accrues to a bonded human principal
across the subnet's fleet rather than to a single device. In some
embodiments, ``PoliePunts'' or similar non-limiting labels are used
for such subnet-denominated units, in the same sense that
``Narravite'' is a non-limiting label for narrative-charged microstructure-unclonability
artefacts. The label is not load-bearing; the committed source, sink,
quota, and cross-subnet policies are load-bearing.

In some embodiments, a token source contract specifies which
committed evidence windows produce subnet-denominated credit and at
what rate. Non-limiting source examples include attested
optical-signature scans of new artefacts, books, samples, or
locations; attested presence at declared sites or events; attested
sensor contributions including camera-trap, environmental-monitoring,
scientific-sample-tracking, or industrial-process windows; B-PoD
discrepancy discovery contributions; attested HC-session contribution
under declared protocols; and attested participation in
interaction-set membership events of the type described for
narrative-charged tokens.

In some embodiments, a token sink contract specifies how subnet
tokens are spent, burned, escrowed, or locked. Non-limiting sink
examples include weighted voting on committed fleet-calibrated
latent-space region handles under Section~\ref{sec:fleet-model};
priority-queue access for scarce compute, detector time, or physical
sampling resources; curation or moderation weight; access gates for
restricted subnet actions; prediction-market settlement under
declared-oracle contracts; quadratic voting, conviction voting, or
positive-sum voting mechanisms; and rate-limited messaging boards
over committed latent-space volumes. In some embodiments, a
latent-space region handle is committed as a hash of a centroid and
radius, a cluster identifier, or another declared region descriptor
in the fleet-calibrated embedding, so that the voting target is
cryptographically unambiguous.

In some embodiments, vote weights or access weights are computed from
a declared function of provenance depth, verification history, bond
freshness, interaction-set membership, and token balance. In some
embodiments, the function is committed per subnet and may differ
across subnets. Mechanism names such as positive-sum game, auction,
quadratic vote, conviction vote, prediction market, decentralised
curation rule, or latent-volume message board are descriptive labels
for declared contract mechanisms, not load-bearing theoretical
commitments. Choice of mechanism is a per-subnet product and policy
decision.

\paragraph{Sybil accounting and cross-subnet amplification
(non-limiting).}
In some embodiments, the subnet's declared exclusivity policy
supplies the accounting rule for per-human per-window quota within
that subnet, conditioned on the declared attacker family, cadence
bound, and tolerance envelope of the underlying bonding records.
Under an exclusive policy, accepted records are treated as
non-amplifying only inside the declared subnet family and only while
freshness and tolerance conditions hold. Under a disclosure-required
policy, parallel participation is auditable but may still multiply
influence according to the declared accounting rule. Under a
promiscuous or isolated policy, the relying party bears explicit
responsibility for whether and how third-party subnet records are
counted.

Cross-subnet Sybil strength is bounded by the weakest-exclusivity
subnet in a human's portfolio and by the least restrictive
cross-subnet policy on which the relying subnet depends. Subnets
relying on third-party subnet records inherit those subnets'
attacker-model declarations, refresh cadences, tolerance envelopes,
and disclosure policies. Non-limiting threat families include
cross-subnet collusion, stale-bond replay, relay use of one human
across multiple bonded PolieBots, coerced participation,
weak-disclosure laundering, and mutually inconsistent exclusivity
declarations. The resolution floor is the declared cross-subnet
accounting envelope, not a universal claim of one-human-one-vote
identity.


\subsection{Red-Team Bounty Primitive for Subnet Detector Hardening
(Non-limiting)}
\label{sec:red-team-bounty}

In some embodiments, the Eve hardening principle described for
physical-channel, information-theoretic, meter-training, and semantic
layers is extended to subnet identity and subnet detector operation.
The subnet is the attack target. The red-teamer is the attacker. The
bounty aligns incentives for discovering detector holes in the
subnet's actual declared interaction surfaces rather than in a
separate toy environment.

\begin{definition}[Subnet red-team bounty attempt (non-limiting)]
A \emph{subnet red-team bounty attempt} is a committed adversarial
attempt against a subnet detector, executed through ordinary subnet
interaction surfaces, and later opened against a committed
post-attack subnet state under a declared success/failure oracle.
\end{definition}

In some embodiments, before the attack window opens, the red-teamer
commits to a structured target claim that the subnet is being tricked
into accepting, where the claim is machine-checkable against
committed subnet state. In some embodiments, the pre-attack
commitment also binds the attack class from the subnet's declared
attack taxonomy, the declared subnet-interaction budget for the
attempt, and a nonce. The commitment is time-anchored using the same
run-commitment and transparency-log mechanisms used for committed
evidence records, and in some embodiments follows the two-seed
pattern of Section~\ref{par:two-seed}: the attack is committed before
the post-attack opening condition is known.

In some embodiments, attack execution occurs only through ordinary
subnet interaction surfaces. Non-limiting surfaces include physically
rate-limited token sources, HC sessions, declared token sinks,
proof-of-projection submissions, tPoP exports, B-PoD handoff
interfaces, latent-space voting targets, and selective-opening
dispute flows. In some embodiments, honeypot verification channels of
the type described in Section~\ref{sec:honeypot} are included in the
surface so that attacker effort leaves diagnostically useful traces.

In some embodiments, after the attack window closes, the subnet
publishes a committed post-attack state hash. The red-teamer then
opens the pre-attack commitment and the relevant evidence atoms. The
success/failure oracle is declared in the protocol digest.
Non-limiting oracle options include automated detector output, a
declared reviewer group, an independent verifier quorum, or a hybrid
rule that combines detector output with selectively opened evidence
windows. In some embodiments, the oracle checks whether the target
claim was accepted, whether the acceptance relied on the committed
attack class, and whether the attack stayed within the declared
interaction budget.

In some embodiments, the subnet commits an attack-class taxonomy in
its protocol digest or policy catalogue. Non-limiting attack classes
include identity confusion; synthetic provenance, including
fabricated attested scans of nonexistent artefacts or locations;
relay attack in which one human operates multiple bonded PolieBots
in rotation; quota-laundering through cross-subnet amplification
contrary to a declared exclusivity policy; semantic injection
inducing the subnet to upvote a latent-space region under false
pretences; collusion by a coordinated multi-principal group; and
coercion-simulation, in which a detector is tested for recognition of
a bonded human operating under duress.

In some embodiments, bounty payment is implemented through the same
subnet token rails as an additional declared sink. In some
embodiments, the bounty-signing key family or signing role is
distinct from the key family used for ordinary source, sink, and
detector-operation records. This is a structural separation of
signing roles within the commitment plumbing of this disclosure; it
does not import any companion-filing governance machinery. In some
embodiments, payment is conditioned on successful opening, oracle
outcome, budget compliance, and absence of collusion flags.

In some embodiments, paid attempts and defeated attempts both enter
the subnet's detector-training evidence set with their committed
evidence chain preserved. The resulting subnet detector-hardness
index is conditioned on the observed red-team population,
attack-class diversity, bounty budget, detector version, and
interaction budget rather than treated as unconditional detector
coverage. The resolution floor is the declared red-team coverage
envelope under the attack taxonomy, detector version, interaction
budget, and bounty budget, time-indexed rather than unconditionally
durable. Relying parties may discount subnet trust when the red-team
population is small, the attack taxonomy is narrow, or the bounty
budget is insufficient for the declared attacker family.

\paragraph{Failure modes and mitigations (non-limiting).}
Non-limiting failure modes include red-teamer front-running by
opening the commitment before the subnet state is fixed, mitigated by
requiring opening only after a committed post-attack state hash is
published; subnet-operator patching of the detector during the attack
window, mitigated by freezing detector version to the version
committed at attack start; and collusion between red-teamers and
subnet operators staging fake attacks, mitigated by separated
bounty-signing keys, selective opening of evidence atoms, and
statistical audit of red-team success rates against a declared
reference distribution.

\paragraph{Out-of-scope note (structural, not policy default).}
This primitive tests subnet detectors. As a structural property of
the primitive, a subnet red-team bounty attempt is addressed to a
subnet detector, committed subnet state, declared token surface,
quota rule, source/sink contract, or other subnet interaction
surface, and is never applied to an applicant under any subnet
policy. Applicant-directed evaluation is outside the primitive's
scope by construction. A workflow that presents a bonded human,
prospective member, or cross-subnet applicant as the attack target
is not an instance of this primitive. The primitive is not an
admission mechanism and not a screening mechanism for bonded humans.
Cross-subnet admission, where present, is handled by the receiving
subnet applying its own declared thresholds to the existing portable
record and any admissible committed evidence; the red-team bounty
primitive does not probe applicants.





\subsection{Physical-layer advantage channels: keyed, trained, and
scene-geometric embodiments}
\label{sec:advantage-channels}

The following embodiments, also referred to in some embodiments as
``reality encryption,'' concern physical-layer advantage channels in
which a cooperative observer or decoder (Bob) is made more capable
than a bounded adversarial observer or decoder (Eve) under declared
resource, detector-geometry, and meter-envelope constraints.  These
embodiments extend the Alice/Bob/Eve multi-agent framing of
Section~\ref{sec:meters-training}, the wiretap advantage-gap lens of
Section~\ref{sec:wiretap-advantage}, and the scale-to-threshold Bob/Eve
curricula of the preceding subsection, and they specialise to three
progressively less pre-coordinated embodiment families: (i)~keyed
deterministic-noise obscuration at a declared cadence, (ii)~trained
neural-network obscuration generators optimised against an explicit
Bob/Eve secrecy objective, and (iii)~meter-bounded discovery and
validation of scene-geometric advantage channels without any
pre-shared key or separately declared digital message channel.  In
each case the disclosed inventive matter is an apparatus-and-procedure
configuration, evaluated through committed \cba\ records and declared
meter envelopes; no claim is made that any particular learned policy
will emerge for every scene or every resource bound.

\paragraph{Stage 1: keyed deterministic-noise advantage at declared
cadence (non-limiting).}
In some embodiments, an auxiliary emission subsystem injects a
deterministic obscuration pattern into the scene, the reactor path,
or the detector-facing observation path.  The obscuration pattern is
driven by a committed seed or pre-shared cryptographic key available
to the cooperative party but not to the bounded adversary.  The
cooperative decoder cancels, subtracts, or compensates for the
obscuration pattern in its observation pipeline within a declared
cancellation-latency budget, while the bounded adversary, lacking the
seed or key, is required to invert the obscuration from observations
alone within a declared inversion-time budget.  In some embodiments,
the following quantities are declared as design variables and are
logged in the protocol digest: the obscuration update cadence, the
detector integration time, the programmable-scene refresh rate (for
example an e-ink, LCD, OLED, or micro-LED surface rendering a
time-indexed plaintext or watermark as in the
programmable-scene jamming game of Section~\ref{sec:meters-training}),
the cooperative-party cancellation-latency budget, and the
adversary's inversion-time budget under a declared bounded resource
class.  A non-limiting design posture is to choose the obscuration
cadence and the cooperative-party cancellation latency so that the
implied adversary inversion-time budget, under the declared bounded
resource class, is empirically insufficient to recover the plaintext
within the cadence window.  The obscuration generator in this stage
is deterministic given its seed; it may be a linear-feedback shift
register, a keyed pseudorandom function, a noise profile whose
commitment, encrypted form, identifier, or selective-opening policy
is stored in the protocol digest, an external runtime input drawn
from a sibling device of the same class under a committed tap point,
or a cryptographically specified stream generator, without
limitation.  Where an anchor noise co-illumination embodiment is
present, this stage composes with that embodiment; in embodiments
not including such an anchor extension, the same keyed-obscuration
logic applies to any auxiliary emission path, reactor-facing path,
or detector-facing path whose control schedule or profile commitment
is bound into the protocol digest under the commitment plumbing of
Section~\ref{sec:commitment-plumbing}.  Physical rate-limiting
(Section~\ref{sec:rate-limited-outputs}) provides a cadence floor
that offline precomputation alone does not remove under a declared
attacker model lacking the cooperative key, timely access to the
obscuration realisation, or equivalent physical access to the
emitting path; this makes cadence a meaningful physical-layer
security parameter rather than a purely protocol-level assumption,
and composes with the correlated-measurement key-agreement lens of
Section~\ref{sec:key-agreement-correlated}.

\paragraph{Stage 2: trained neural-network obscuration generator under
an explicit Bob/Eve secrecy objective (non-limiting).}
In some embodiments, the obscuration generator is a trainable
neural-network or parameterised functional whose output is optimised
end-to-end against an explicit Bob/Eve secrecy objective rather than
being fixed by a deterministic seed.  In a non-limiting form, the
obscuration generator~$G_\phi$ produces an obscuration waveform or
emission pattern~$n_t = G_\phi(k, \xi_t)$ from a cooperative-party
key, declared private input, or cooperative-only observation~$k$ and
an internal state~$\xi_t$; a cooperative
decoder~$D^{\mathrm{Bob}}_{\psi_B}$ recovers a plaintext or target
variable from its observation given the key; and a bounded adversary
decoder~$D^{\mathrm{Eve}}_{\psi_E}$ attempts the same recovery from
its declared degraded view without the key.  A non-limiting
generator-side objective, evaluated against a best-response Eve
decoder, is
\[
  \mathcal{L}(\phi, \psi_B, \psi_E)
  = -\,\E\!\left[\log P_{D^{\mathrm{Bob}}_{\psi_B}}
      \bigl(\Xi \mid Y_B, k\bigr)\right]
  + w_{\mathrm{adv}}\,
    \E\!\left[\log P_{D^{\mathrm{Eve}}_{\psi_E}}
      \bigl(\Xi \mid Y_E\bigr)\right],
\]
where $\Xi$ is the plaintext or target variable, $Y_B$ and $Y_E$ are
the cooperative and adversarial observation views, respectively, and
$w_{\mathrm{adv}} \ge 0$ is a non-limiting weighting between Bob
recovery and Eve suppression.  In some embodiments, the optimisation
is minimax, alternating a minimisation over
$(\phi, \psi_B)$ with a best-response maximisation over~$\psi_E$
within a declared bounded resource class for Eve.  In some
embodiments, physical emission parameters of the auxiliary subsystem
(scan law, amplitude gain, spectral profile, polarisation, timing
alignment) are co-trained with~$\phi$ using hardware-in-the-loop
procedures (Section~\ref{sec:hitl-optimisation}) or a differentiable
surrogate, so that training pressure acts on physically realisable
obscuration waveforms rather than on waveforms that cannot be emitted.
In some embodiments, the Stage~2 training game composes directly with
the Bob/Eve scale-to-threshold curricula of the preceding subsection:
adversarial capacity~$c$ for Eve is grown toward a declared threshold
while cooperative capacity for Bob is reduced, with the gap objective
$c^\ast_{\mathrm{Eve}} - w_{\mathrm{adv}}\,c^\ast_{\mathrm{Bob}}$
providing an auditable, capacity-calibrated version of the secrecy
objective above.  Meter envelopes, eye-safety bounds, and power
ceilings constrain the admissible obscuration policies during and
after training.

\paragraph{Stage 3: meter-bounded discovery and validation of
scene-geometric advantage channels without pre-shared keys
(non-limiting).}
In some embodiments, the \RK\ scene is configured as a physical
search space for candidate advantage-channel policies among
role-assigned agents, without requiring a pre-shared seed, a
pre-committed noise profile, or a separately declared digital
message channel between the cooperative parties; the candidate
channel, if any, is the meter-bounded physical scene interaction
itself.  A first policy, assigned the Alice role, controls one or
more emission, projection, scan, wavelength, polarisation, timing,
amplitude, or scene-actuation variables; a second policy or decoder,
assigned the Bob role, receives observations through a declared
detector geometry or observation port; and a third policy or decoder,
assigned the Eve role, receives observations through a different
declared detector geometry, observation port, or resource-bounded
view.  The scene may include, without limitation, specular
reflectors, diffuse reflectors, occluders, shadowing geometry,
saturating or non-linear surfaces, wavelength-selective elements,
polarisation-selective elements, retroreflectors, dynamic or
addressable surfaces, and volumetric scattering media, with their
placements, orientations, and spectral properties committed as part
of the scene configuration metadata.  A non-limiting advantage
objective combines cooperative-party recovery performance,
bounded-adversary recovery performance, meter-envelope compliance,
and physical-safety compliance, and may be optimised through
reinforcement learning, evolutionary search, gradient-based search
over a differentiable scene surrogate, or any combination of these,
subject to declared actuator and emission bounds.  No claim is made
that any particular policy will reliably emerge for every scene or
every resource bound; the inventive matter is the apparatus and
procedure for searching, validating, and committing candidate
policies, not a guaranteed learning outcome.

\paragraph{Discovery, validation, and commitment procedure
(non-limiting).}
In some embodiments, a Stage~3 embodiment executes the following
procedure, producing a committed, auditable evidence record at each
step:
\begin{enumerate}
  \item \textbf{Scene configuration.} Configure the physical scene
    with a declared catalogue of optical affordances (reflectors,
    occluders, saturating surfaces, wavelength- or
    polarisation-selective elements, dynamic surfaces) and commit the
    scene configuration metadata, including affordance placements,
    orientations, and declared spectral and polarisation properties.
  \item \textbf{Role-assigned search.} Execute role-assigned search
    or training over candidate Alice, Bob, and Eve policies under the
    declared advantage objective, bounded actuator envelopes, and
    meter envelopes, with each executed episode producing an
    emission-observation record.
  \item \textbf{Evidence commitment.} Commit for each candidate
    policy the \cb~$C_{0:T}$, the protocol digest, the Bob-side and
    Eve-side detector geometries, the meter outputs, and any
    policy-identifying parameters or parameter digests, using the
    commitment plumbing of
    Section~\ref{sec:commitment-plumbing}.
  \item \textbf{Advantage measurement.} From the committed record,
    measure cooperative-party recovery performance, bounded-adversary
    recovery performance, and the advantage gap between them, using
    a pre-declared evaluation distribution and a pre-declared
    bounded-adversary resource class.
  \item \textbf{Affordance ablation.} Run targeted ablation episodes
    in which selected scene affordances are removed, moved, rotated,
    spectrally altered, polarisation-altered, or replaced with a
    null-affordance control, and measure the advantage gap again on
    each ablated scene.
  \item \textbf{Acceptance criterion.} Accept a candidate advantage
    channel only when (a)~the advantage gap on the unablated scene
    satisfies the declared threshold, (b)~the advantage gap changes
    in the pre-declared direction and by at least a declared
    diagnostic magnitude, or falls below a declared floor, under a
    pre-declared subset of ablations identified as diagnostic of the
    hypothesised physical mechanism, and (c)~meter-envelope and
    safety bounds hold throughout.  The ablation-sensitivity
    requirement functions as a physical-mechanism falsification
    test: an accepted channel must depend on the identified scene
    affordance in the committed evidence, rather than on an unlogged
    digital side-channel or a spurious training artefact.
\end{enumerate}

\paragraph{Viewpoint-specific decoding as the physical origin of the
advantage (non-limiting).}
In some embodiments, the physical origin of the advantage in Stage~3
is that a single physical emission produces different information
channels at different detector geometries, wavelengths, polarisations,
or temporal-integration windows.  Specular reflection from a bounded
region of the scene, for example, can produce a high-intensity return
only along a narrow range of detector viewpoints; occluders can
deny a detector at one declared geometry access to a region that is
fully visible at another declared geometry; saturating surfaces can
exceed the linear range of a detector at one declared geometry while
remaining within range at another; and wavelength- or
polarisation-selective elements can transmit to one declared channel
while attenuating another.  This viewpoint-dependence connects
Stage~3 to the cross-scene regime combinations of
Section~\ref{sec:cross-scene}, and provides a physical account of
why a Bob-side detector geometry can be structurally favoured over an
Eve-side detector geometry under a committed emission sequence,
without any pre-shared key between the cooperative parties.

\paragraph{Relationship to adjacent technical families
(non-limiting).}
Without admitting that any particular reference or family constitutes
prior art against any claim, the embodiments of this subsection may
be understood relative to several adjacent technical families as
follows.  Stigmergy and related indirect-coordination mechanisms in
swarm robotics use persistent low-dimensional environmental markers
as coordination memory; the present embodiments instead use
meter-bounded, high-dimensional, viewpoint-dependent optical or
multi-physics scene affordances recorded in committed \cba.
Emergent-communication methods in multi-agent reinforcement learning,
including reinforced or differentiable inter-agent learning over
declared discrete or differentiable message channels, optimise
messages over an explicit communication action space; the present
embodiments have no declared communication action space or shared
seed and instead search physical scene affordances, validating
candidate channels through committed evidence and scene-affordance
ablation.  Physical adversarial-patch and projector-based
adversarial-attack methods optimise a physical perturbation to
induce classifier misbehaviour in a target vision system; the present
embodiments optimise and validate a cooperative/adversarial recovery
gap under declared detector geometries and committed \cba, rather
than targeting a specific classifier failure mode.  Free-space-optical
and visible-light-communication physical-layer-security methods
design and analyse a predesigned optical link under a channel model
and an eavesdropper model, often with beamforming or precoding at a
known transmitter geometry; the present embodiments treat the scene
affordances themselves---reflectors, occluders, saturation,
polarisation, wavelength selectivity, detector viewpoint---as the
discovered channel, validated through the commitment and ablation
procedure described above rather than through a communications-system
channel model.  These observations are non-limiting and are provided
to clarify inventive scope rather than to narrow it.

\paragraph{Composition with existing embodiments (non-limiting).}
The three stages above are not mutually exclusive.  In some
embodiments, a Stage~1 keyed-noise layer, a Stage~2 trained
obscuration generator, and a Stage~3 scene-geometric search proceed
concurrently, with the three contributions composing into a combined
advantage objective subject to a shared meter envelope and safety
bound.  In some embodiments, the Stage~2 trained generator is
initialised from a Stage~1 deterministic seed and then unfrozen for
end-to-end training; in some embodiments, a Stage~3 scene
configuration that produces a sufficient Bob/Eve advantage gap is
used as a warm start for a Stage~2 training run in which the Stage~2
generator is tuned to widen the gap further.  In all cases the
inventive matter is the apparatus-and-procedure configuration and the
committed evidence record, not a promise of any particular learned
outcome.







% ----------------------------------------------------------------------
% Conditional cross-reference macro for the bonded-human-identity
% extension. Resolves to a live \ref if the BHI extension's label is
% defined in the combined compile, and to a textual reference if not,
% so this extension compiles cleanly against base Filing 1 alone or
% against base Filing 1 + BHI extension.
% ----------------------------------------------------------------------
\makeatletter
\providecommand{\PortableHumanBondingRef}{%
  \@ifundefined{r@sec:portable-human-bonding}%
    {the bonded-human-identity extension}%
    {Section~\ref{sec:portable-human-bonding}}}
\makeatother


\paragraph{Positioning relative to existing commitment infrastructure
(non-limiting).}
In some embodiments, the primitives described below extend the
commitment and selective-disclosure machinery of
Section~\ref{sec:two-seed-protocol} and the hiding and binding lenses
of the same section to support use cases in which the evidentiary
weight of a genuine live session is deliberately scoped at the time
of capture, scoped to a participating party set, or rescaled after
capture under a declared trapdoor authority. The underlying
designated-verifier, commit--reveal randomness, and chameleon-hash
cryptographic primitives are not asserted here as novel in isolation.
The contribution described in this section is their composition with
\cba commitments, the two-seed selective-opening
protocol, disclosure-regime tags in the protocol digest,
proof-of-projection and proof-of-discrepancy audit workflows, and
authority-action discipline. Adjacent anonymity-family primitives,
including anonymous-set attestations, ring signatures,
selective-disclosure credentials, zero-knowledge set-membership
proofs, and plausibly deniable encryption, may be composed with the
records disclosed here as external credential or confidentiality
layers, but are not required for the designated-verifier, mutual
$N$-party liveness, or chameleon-opening embodiments disclosed in
this section. The primitives are composable with the canonical
two-seed selective-opening protocol (Section~\ref{par:two-seed}),
with \cba commitments, with proof-of-projection
records, and with the bonded-human-identity records of
\PortableHumanBondingRef. They do not weaken and are
not a substitute for: positive verification under the \TB
regime (Section~\ref{sec:truth-beam}); provenance negation under
the Anti-\TB lens (Section~\ref{sec:honeypot} and the
surrounding provenance-negation material); proof-of-discrepancy
quorum confirmation under the PoD verifier-selection protocol
(Section~\ref{sec:pod-protocol}); or the human-coupled-trajectory
evidence chain of the bonded-human-identity extension. Each of
those workflows operates on committed evidence under disclosure
policies explicitly declared in the protocol digest; the deniability
primitives below apply only where the protocol digest declares the
relevant disclosure regime. The primitives inherit the
terminological note of the proof-of-projection section: the
deniability, simulatability, and trapdoor properties asserted below
are empirical under declared attacker families and resource
budgets, are time-indexed rather than unconditionally durable, and
do not constitute formal cryptographic proof-system guarantees
unless a specific embodiment explicitly specifies and parameterises
such a guarantee.


\subsection{Designated-verifier liveness via per-block salts
  (non-limiting)}
\label{sec:designated-verifier-liveness}

\begin{definition}[Designated-verifier liveness record
  (non-limiting)]
A \emph{designated-verifier liveness record} is a \cba
commitment whose liveness-attestation property is interpretable only
by a declared verifier~$V$ possessing a pre-session per-block salt
$s_V$ shared with the submitting device~$A$, and whose published
transcript is, under the declared attacker family, adversary view,
and resource budget, not distinguishable by any third party with
advantage above a declared bound from a transcript producible by
$V$ alone using the same salt.
\end{definition}

In some embodiments, before the run commitment is published, device
$A$ and verifier~$V$ establish a fresh per-block salt~$s_V$ over a
declared side channel. Non-limiting side channels include a short
pre-run handshake session, a pre-registered session key, a
key-derivation from a longer-lived shared secret bound to an earlier
committed joint session, and a hash chain rooted at a previously
attested coupling record. The side channel itself is governed by a
declared attacker model, freshness discipline, and compromise
assumptions committed to the protocol digest. The salt is bound
into the run commitment as an additional field:
\[
  \mathrm{com}^{(i)}_{\mathrm{DV}}
  = H\!\left(
    \mathrm{rid}\,\|\,i\,\|\,\Pi_{\mathrm{dig}}\,\|\,\mathbf{m}^{(i)}\,\|\,
    R^{(i),\mathrm{obs}}\,\|\,R^{(i),\mathrm{emit}}\,\|\,s_V
  \right),
\]
and is not disclosed to any party outside $\{A, V\}$ under the
declared disclosure regime.

In some embodiments, the declared verifier identity, verifier-set
digest, or an unlinkable verifier handle $h_V$ is committed in the
protocol digest and, where the disclosure regime permits, bound
into the designated-verifier commitment, so that the same salt
material is not ambiguously reusable across verifier identities,
verifier sets, or parallel sessions. Where public disclosure of
verifier identity would damage the simulatability property above,
$h_V$ is auditor-openable rather than publicly visible, under the
declared disclosure regime and adversary view.

In some embodiments, the liveness property claimed by the record is
stated as two operational components:
\begin{enumerate}[nosep]
  \item \emph{Verifiability by the designated verifier.} $V$
    reconstructs $\mathrm{com}^{(i)}_{\mathrm{DV}}$ from the opened
    bundle atoms, the committed protocol digest, and $s_V$, and
    checks agreement with the published commitment under the
    hiding and binding lenses of
    Section~\ref{sec:two-seed-protocol}.
  \item \emph{Simulatability against third parties.} Under the
    declared attacker family, adversary view, and resource budget,
    a third party~$T \notin \{A, V\}$ observing any subset of the
    published transcript, the run seed, and the opened atoms
    cannot, with advantage above a declared bound, distinguish a
    record produced by a genuine $A$--$V$ session from a record
    $V$ produced unilaterally using an alternative bundle and the
    same $s_V$. The declared advantage bound, the adversary view,
    and the attacker resource budget are committed to the
    protocol digest.
\end{enumerate}

In some embodiments, per-block salts are rotated per run, per
segment, or per atom. Rotation policy, salt-derivation KDF chain,
and salt-anchor identifiers are committed to the protocol digest.
In some embodiments, the salt-anchor identifier references an
earlier committed joint session between $A$ and $V$ so that the
salt chain is itself auditable under selective opening when the
disclosure regime permits. Salt-anchor handles may be unlinkable
tokens rather than long-lived identifiers where cross-session
linkability is itself part of the declared threat model.

\paragraph{Outer-authentication simulation policy (non-limiting).}
In some embodiments, any outer transport signature, device
identifier, registry entry, certificate chain, or publication
channel used to carry a designated-verifier liveness record is
itself marked as designated-verifier scoped, omitted from the
third-party comparison view declared in the protocol digest, or
simulable by $V$ under a declared transcript-simulation policy;
otherwise the simulatability property above is not asserted to
extend to the outer authentication layer. The transcript-simulation
policy, the set of outer layers it covers, and any omitted layers
are committed to the protocol digest.

\paragraph{Composition with the two-seed protocol (non-limiting).}
In some embodiments, the designated-verifier salt is a disclosure
binding separate from the run seed $s_{\mathrm{run}}$ and opening
seed $s_{\mathrm{open}}$ of Section~\ref{par:two-seed}. The two-seed
machinery continues to govern which atoms are opened and when; the
designated-verifier salt governs \emph{to whom} the opened record is
meaningful as a liveness attestation. In some embodiments, a record
may be opened to $V$ under the designated-verifier regime and
independently re-opened under a public-verifier regime only if the
protocol digest declares an explicit re-anchoring protocol at
commit time.

\paragraph{Failure modes and mitigations (non-limiting).}
Non-limiting failure modes include salt leakage to third parties
outside $\{A, V\}$, mitigated by rotation, committed anchoring
side-channel identifiers, and declared salt-validity windows;
compromise of the salt-establishment side channel itself, including
man-in-the-middle capture or replay of handshake material, mitigated
by declared side-channel attacker models, freshness counters, and
externally anchored salt epochs; replay of a prior $\{A, V\}$ salt
or salt-anchor session, mitigated by monotonic counters, freshness
anchors, and committed rotation policies; linkability of multiple
designated-verifier records across salt-anchor identifiers or
rotation chains, mitigated by unlinkable salt-anchor handles and
disclosure-minimised salt-anchor referencing; re-identification of
$A$ from outer signatures, run identifiers, protocol-digest
metadata, meter envelopes, or timing metadata notwithstanding the
hiding of $s_V$, mitigated by disclosure-minimised metadata
profiles and by the outer-authentication simulation policy above;
collusion between $A$ and a third party to fabricate an
$A$-looking transcript without a live session, which is not a
designated-verifier failure and is addressed by the independent
proof-of-projection and quorum-confirmation machinery of the
present filing;
and unilateral fabrication by $V$ of a record that $V$ then
presents to a third party as $A$'s attestation, which is a feature
rather than a bug of the primitive and is addressed at the
application layer by the relying party declining to accept
designated-verifier records from $V$ as $A$-originating
attestations.


\subsection{Mutual $N$-party liveness via commit--reveal
  (non-limiting)}
\label{sec:mutual-n-party-liveness}

\begin{definition}[Mutual $N$-party liveness record (non-limiting)]
A \emph{mutual $N$-party liveness record} is a \cba
commitment whose run seed is derived from a committed $N$-party
nonce protocol such that each participating party independently
witnesses session freshness from its own unpredictable contribution,
while a non-participant observer, under the declared attacker
family, adversary view, and resource budget, cannot with advantage
above a declared bound distinguish a genuine $N$-party session from
a record fabricated by a coalition with equivalent knowledge of the
combined nonce.
\end{definition}

In some embodiments, the protocol proceeds as follows:
\begin{enumerate}[nosep]
  \item \textbf{Commit phase.} Each party
    $k \in \{1, \ldots, N\}$ independently generates a nonce
    $n_k$ from a declared unpredictability source, computes a
    hiding commitment
    \[
      c_k =
      H(\mathsf{tag}_{\mathrm{NPL}}\,\|\,\mathrm{rid}\,\|\,k\,\|\,
        \Pi_{\mathrm{dig}}\,\|\,L_{\mathrm{part}}\,\|\,
        n_k\,\|\,r_k),
    \]
    with fresh randomness $r_k$, where
    $\mathsf{tag}_{\mathrm{NPL}}$ is a declared domain-separation
    tag, $\Pi_{\mathrm{dig}}$ is the protocol digest,
    $L_{\mathrm{part}}$ is a committed canonical participant-list
    digest, and $k$ is interpreted under the declared canonical
    participant ordering. For this commit phase,
    $\Pi_{\mathrm{dig}}$ denotes the pre-run protocol-template
    digest and excludes nonce-derived realised fields except where
    such fields are represented by declared templates or
    placeholders. Each party publishes $c_k$ to the
    joint commit set before a declared commit deadline. The
    unpredictability source, canonical participant ordering,
    participant-list digest, and commit deadline are committed
    to the protocol digest.
  \item \textbf{Reveal phase.} After the commit deadline has
    passed, each party publishes $(n_k, r_k)$. A late reveal, a
    malformed reveal, or a non-reveal is processed according to a
    declared non-reveal policy.
  \item \textbf{Combining.} The combined nonce
    $n_{\mathrm{comb}}$ is computed from the revealed $\{n_k\}$
    under a declared combining rule. Non-limiting combining rules
    include domain-separated hash-concatenation
    \[
      n_{\mathrm{comb}} =
      H(\mathsf{tag}_{\mathrm{NPL\text{-}comb}}\,\|\,\mathrm{rid}\,
        \|\,\Pi_{\mathrm{dig}}\,\|\,L_{\mathrm{part}}\,\|\,
        n_{\pi_{\mathrm{ord}}(1)}\,\|\,\cdots\,\|\,n_{\pi_{\mathrm{ord}}(N)}),
    \]
    where $\pi_{\mathrm{ord}}$ is the committed canonical participant
    ordering (a permutation of the participant indices; distinct from
    trust score $\sigma_{ij}(t)$ and from any other body-text use of
    $\sigma$) and $\mathsf{tag}_{\mathrm{NPL\text{-}comb}}$ is a
    declared combining-rule domain-separation tag; bitwise XOR;
    and threshold rules in which any
    $t$-of-$N$ revealed nonces determine $n_{\mathrm{comb}}$
    under a declared threshold scheme.
  \item \textbf{Run seed derivation.} The run seed
    $s_{\mathrm{run}}$ of the two-seed protocol of
    Section~\ref{par:two-seed} is derived from
    $n_{\mathrm{comb}}$ and the committed protocol digest.
  \item \textbf{Run phase and commitment.} The device executes
    the run under $s_{\mathrm{run}}$ and publishes the run
    commitment. The subsequent opening phase proceeds under the
    two-seed protocol.
\end{enumerate}

\paragraph{Public-beacon variant and run-seed confidentiality
(non-limiting).}
In the variant in which $n_{\mathrm{comb}}$ is publicly revealed
before the run, the combined nonce supplies post-commit freshness
for the run seed; the opening seed $s_{\mathrm{open}}$ remains
independently supplied after the run commitment under the two-seed
protocol. In some embodiments, where run-seed confidentiality until
opening is required for its anti-tailoring role, $n_{\mathrm{comb}}$
is combined with a device-side committed nonce or verifiable-random-
function output that remains undisclosed until selective opening, so
that $s_{\mathrm{run}}$ is pre-bound to the joint freshness but is
not fully determined by the public beacon alone. The chosen variant
(public beacon or device-side combined) is committed to the protocol
digest at commit time.

In some embodiments, the liveness property claimed by the record is:
for each party~$k$ whose nonce $n_k$ was unpredictable to the other
$N-1$ parties at the commit deadline under the declared
unpredictability-source assumptions, the run seed is not fixed
before~$k$'s commitment by any party or coalition lacking $n_k$,
except with advantage bounded under the declared attacker family,
nonce-source quality, reveal policy, and resource budget; therefore
the session is alive from~$k$'s perspective under those declared
bounds. A non-participant observer sees only the combined nonce and
the public transcript, and cannot, with advantage above the declared
bound and under the committed adversary view, distinguish a genuine
$N$-party session from one in which a coalition reconstructed
transcript fields $\{c_k, n_k, r_k\}$ consistent with the combining
rule.

In some embodiments, the primitive is composable with
designated-verifier liveness records by designating one or more of
the $N$ parties as designated verifiers under the machinery of
Section~\ref{sec:designated-verifier-liveness}, with per-block salts
bound into the joint commit set. In some embodiments, the primitive
is composable with proof-of-projection quorum protocols by using
the $N$ parties as an initial freshness-attestation quorum separate
from the verification quorum of the PoD protocol
(Section~\ref{sec:pod-protocol}).

\paragraph{Scope of the liveness claim (non-limiting).}
The primitive binds the run's freshness to $N$~independently
committed unpredictability contributions. It does not on its own
establish that the $N$ parties correspond to $N$ distinct identities;
Sybil resistance, where relevant, is handled by the
bonded-human-identity mechanisms of \PortableHumanBondingRef, by
device-identity
machinery, or by external-identity anchors declared in the protocol
digest.

\paragraph{Failure modes and mitigations (non-limiting).}
Non-limiting failure modes include last-revealer advantage, in which
a party withholds its reveal until after observing other reveals
and then chooses whether to reveal so as to bias
$n_{\mathrm{comb}}$, mitigated by the commit phase preceding
the reveal phase, by non-reveal penalty policies committed in the
protocol digest, and by threshold combining rules that remain
well-defined when a bounded number of parties fail to reveal;
non-uniform or correlated nonce sources that amplify combining-rule
bias, including shared random-number-generator infrastructure or
shared hardware entropy sources, mitigated by extractor-based nonce
conditioning and by committed declaration of independent
unpredictability-source identifiers per participant;
abort-and-retry grinding across repeated sessions, and
threshold-subset grinding in which multiple valid reveal subsets
produce different $n_{\mathrm{comb}}$ values, mitigated by
domain-separated combining hashes, deterministic non-reveal
defaults, abort-rate meters committed in the protocol digest, and
committed tie-breaking rules over the reveal-subset lattice;
rushing during the reveal phase notwithstanding the commit deadline,
mitigated by simultaneous or deadline-enforced reveal publication,
committed reveal-order rules, and externally anchored reveal
timestamps; coordinated pre-agreement in which all $N$ parties
jointly choose a target $n_{\mathrm{comb}}$ before the commit
phase, which defeats the freshness claim against non-participants
and is addressed at the application layer by requiring that at
least one party be independently selected by a non-colluding anchor
(for example a verifier-selected participant drawn from a fleet
under the randomised verifier-selection rule); nonce-source
compromise, mitigated by declared unpredictability sources with
attacker-resource budgets committed in the protocol digest; and
anti-tailoring degradation in the public-beacon variant, addressed
by the device-side committed nonce or VRF composition described
above where anti-tailoring is required.


\subsection{Chameleon-hash retroactive selective evidentiary weight
  (non-limiting)}
\label{sec:chameleon-hash-evidentiary}

\begin{definition}[Chameleon-hash evidentiary record (non-limiting)]
A \emph{chameleon-hash evidentiary record} is a \cba
commitment in which one or more binding fields are instantiated
using a chameleon-hash function
$\mathsf{CH}(\cdot, \cdot; \mathrm{pk})$ under a designated
trapdoor authority's public key $\mathrm{pk}$, such that the record
is treated as binding, under the hiding and binding lenses of
Section~\ref{sec:two-seed-protocol}, against parties without
knowledge of the trapdoor~$\mathrm{tk}$, under the declared
attacker family, adversary view, parameterisation, and resource
budget, and such that the trapdoor authority may, under a declared
authority-action policy, produce an alternative valid opening
consistent with the published commitment.
\end{definition}

In some embodiments, the run commitment is computed as
\[
  \mathrm{com}^{(i)}_{\mathrm{CH}}
  =
  \mathsf{CH}\!\left(
    \mathbf{z}^{(i)},\;
    r^{(i)};\;
    \mathrm{pk}_{\mathcal{T}}
  \right),
\]
where $\mathbf{z}^{(i)}$ denotes the same binding-field tuple used
in $\mathrm{com}^{(i)}_{\mathrm{full}}$ of
Section~\ref{sec:two-seed-protocol}, optionally with layer
identifiers and disclosure-regime tags added under the declared
layer-composition policy; $r^{(i)}$ is opening randomness; and
$\mathrm{pk}_{\mathcal{T}}$ is the public key of a declared
trapdoor authority~$\mathcal{T}$. The corresponding trapdoor
$\mathrm{tk}_{\mathcal{T}}$ is held by the authority under the
declared authority-custody policy.

In some embodiments, the claimed binding-and-plausible-deniability
properties are:
\begin{enumerate}[nosep]
  \item \emph{Third-party binding.} Under the declared attacker
    family and resource budget, a party without
    $\mathrm{tk}_{\mathcal{T}}$ cannot, with advantage above a
    declared bound, produce two distinct valid openings
    $(\mathbf{z}, r) \neq (\mathbf{z}', r')$ consistent with the
    same $\mathrm{com}^{(i)}_{\mathrm{CH}}$. The bound may be
    expressed as $\mathrm{negl}(n_{\mathrm{sec}})$ for a
    conventional digital chameleon-hash instantiation and as an
    empirical acceptance-probability bound for any
    physical-evidence component of the claim; the bound,
    adversary view, parameterisation, and resource budget are
    committed to the protocol digest.
  \item \emph{Authority-side alternative opening.} For a selected
    chameleon-hash instantiation satisfying the declared
    trapdoor-collision algorithm, the trapdoor authority
    $\mathcal{T}$, holding $\mathrm{tk}_{\mathcal{T}}$, can
    compute $r'$ such that
    $\mathsf{CH}(\mathbf{z}', r'; \mathrm{pk}_{\mathcal{T}}) =
    \mathrm{com}^{(i)}_{\mathrm{CH}}$ for an alternative
    $\mathbf{z}'$ drawn from a declared alternative-opening
    policy space. The choice of chameleon-hash instantiation,
    parameter size, hash domain separation, trapdoor-collision
    algorithm, and post-quantum or classical security posture is
    committed to the protocol digest at capture time.
\end{enumerate}

In some embodiments, the trapdoor authority is non-limitingly the
submitting fleet, a designated policy authority separate from the
submitter, the submitting party itself (self-trapdoor), a
committee holding shares of $\mathrm{tk}_{\mathcal{T}}$ under a
threshold secret-sharing or threshold trapdoor-collision-
generation scheme, or a committee using threshold signatures as
an authorisation layer controlling a separately declared trapdoor-
custody mechanism, with a declared quorum of $t$-of-$n$ authority
members required before an alternative opening is treated as
conforming. The selected authority mode, custody mechanism, and
quorum rule are committed to the protocol digest and referenced
in each authority-action record.

\paragraph{Scope of the deniability claim (non-limiting).}
The primitive is evidentiary, not definitional. It supports
retroactive scoping of the weight assigned to a specific opening of
a commitment; it does not erase that a commitment was published at
time~$t$, that a \cba capture occurred in the
committed operating envelope, or that cross-device corroboration
occurred under independent channels. In particular, independent
transitive proof-of-projection records, peer-verifier quorum
confirmations under Section~\ref{sec:pod-protocol}, and externally
anchored transparency-log entries binding the commitment to a time
anchor remain operative against any subsequent alternative-opening
action by the trapdoor authority, and continue to contribute to
provenance determinations under the Anti-\TB lens.

\paragraph{Mixed-layer composition (non-limiting).}
In some embodiments, a chameleon hash is used for the full run
commitment, for selected atoms under Merkle inclusion, or for a
summary digest while leaving lower-level atoms under standard
commitments. The layers at which chameleon and standard commitments
are applied, the declared alternative-opening policy space at each
layer, and the trapdoor-authority identity at each layer are
committed to the protocol digest. In some embodiments, records
using a chameleon hash at any layer carry an explicit
disclosure-regime tag in the protocol digest so that a verifier can
identify which binding regime applies at which level of the record.
Where deniability requires that chameleon-hash records be
indistinguishable from standard-hash records in the public
commitment stream, the regime tag is auditor-openable rather than
publicly visible, and the chameleon instantiation is selected so
that $\mathrm{com}^{(i)}_{\mathrm{CH}}$ values are indistinguishable
from standard-hash values at the commitment-envelope layer under
the declared adversary view. Where the regime tag is
auditor-openable rather than publicly visible, a relying party
whose acceptance policy requires non-trapdoor binding either
requires opening of the regime tag under its acceptance policy or
treats the unopened record as potentially trapdoor-bound and
assigns it the corresponding lower evidentiary weight.

\paragraph{Authority-action audit discipline (non-limiting).}
In some embodiments, any alternative opening produced by the
trapdoor authority is itself committed in a transparency log
anchored externally, under the same external-anchoring convention
used elsewhere in this specification for freshness anchors. In some
embodiments, authority-action frequency is rate-limited under a
declared policy digest and the count of alternative openings
produced against a given commitment is bounded and committed. In
some embodiments, authority-action openings are themselves subject
to selective opening under an auditor-supplied seed analogous to
the two-seed protocol, so that conforming relying-party policies
can reject or downgrade any alternative opening that lacks the
corresponding authority-action audit record, and an off-ledger
alternative opening is treated as a policy violation rather than
as a conforming evidentiary action. Threshold
participation is not treated as evidence of authority honesty. In
embodiments using threshold secret-sharing or threshold signatures,
each share-use, partial-opening contribution, quorum transcript,
signer identity, policy version, and authority-action reason code
is committed as an authority-action record; records lacking such
audit material are treated as lower-weight or inadmissible under
the relying party's declared acceptance policy.

\paragraph{Failure modes and mitigations (non-limiting).}
Non-limiting failure modes include trapdoor compromise, mitigated
by threshold secret-sharing of $\mathrm{tk}_{\mathcal{T}}$,
periodic key refresh under a committed rotation policy, and
authority-custody procedures committed in the protocol digest;
long-term trapdoor exposure and retroactive compromise of
historical records, mitigated by epoch-scoped trapdoor keys,
committed key-rotation and destruction attestations, and
relying-party acceptance policies that discount records whose
governing epoch is no longer under declared custody; inadequate
forward security across key epochs, mitigated by declared
forward-security discipline over epoch transitions and externally
anchored destruction attestations; public-stream distinguishability
between chameleon-hash and standard-hash records where deniability
requires indistinguishable commitment formats, mitigated by
instantiation-choice discipline and by auditor-openable rather
than publicly visible regime tags; quantum-adversary degradation
of classical chameleon-hash assumptions, mitigated by declared
post-quantum instantiations where required and by epoch-scoped
migration policies; post-capture expansion of the alternative-
opening policy space by the trapdoor authority, mitigated by
capture-time commitment of the permitted alternative-opening domain
and by authority-action audit records that expose any deviation;
misuse of the alternative-opening capability by the trapdoor
authority to silently substitute a fabricated opening for a
legitimate one, mitigated by the authority-action audit discipline
above and by relying-party policies that downgrade the evidentiary
weight of chameleon-hash records in contexts where authority
actions are not externally auditable; and confusion between binding
and non-binding regimes by relying parties, mitigated by the
explicit disclosure-regime tagging of the protocol digest and by
relying-party acceptance policies that specify which binding regime
is required for which downstream decision.


\paragraph{Cross-cutting disclosure-regime discipline
(non-limiting).}
In some embodiments, a session does not freely switch between
designated-verifier, $N$-party, and chameleon-hash regimes after
the run commitment is published without a declared re-anchoring
protocol. The disclosure regime in force for each layer of a record
is committed to the protocol digest at commit time. A relying party
accepting a record for a downstream decision consults the
disclosure-regime tags and applies its own declared acceptance
policy for the regime at each layer. In particular:
\begin{enumerate}[nosep]
  \item Designated-verifier records are not acceptable
    proof-of-discrepancy submissions under the PoD quorum protocol
    of Section~\ref{sec:pod-protocol} unless the designated
    verifier is itself part of the declared verifier quorum and
    opens the record to the quorum under the quorum's disclosure
    policy. Admission of such a record does not alter the
    pre-committed PoD parameters
    $(n, q, \tau_{\mathrm{confirm}}, \tau_{\mathrm{device}},
    \operatorname{agg})$, the Byzantine-fault-tolerance
    assumptions (including, where the declared PoD policy is
    Byzantine-fault-tolerant, the constraints $n \ge 3f+1$ and
    $q \ge 2f+1$), or the acceptance and evaluation meter
    partition.
  \item Chameleon-hash records are not acceptable as positive
    attestations under the \TB regime unless the relying
    party's acceptance policy explicitly admits records whose
    binding fields are under a trapdoor authority, in which case
    the policy declares how authority-action records and
    alternative-opening frequencies feed into the acceptance
    decision. Even where such admission is declared, \TB
    positive verification remains based on the underlying
    physical evidence, meter envelope, and independent
    corroboration; the chameleon layer affects disclosure regime
    and evidentiary weighting only, and does not alter the
    underlying physical-verification status of the record.
  \item Designated-verifier and chameleon-hash records do not
    alter the operational semantics of provenance negation under
    the Anti-\TB lens: a record's content remains subject
    to cross-device corroboration under independent channels
    regardless of which deniability regime governs disclosure of
    that record to third parties.
\end{enumerate}

\paragraph{Interaction with bonded-human-identity records
(non-limiting).}
In embodiments in which records from \PortableHumanBondingRef are
present, designated-verifier liveness may be used for
HC-Read, HC-Write, or HC-Verify channels where a human participant
wishes to scope attestation of a specific coupled-trajectory event
to a specific fleet-side counterparty. In some embodiments,
chameleon-hash binding may be used for auxiliary human-response
channels (non-limitingly: physiological coherence summaries,
reaction-time distributions) where an applicable data-protection
policy requires the human participant or a designated authority to
retain retroactive capability to produce alternative openings of
specific response summaries. In both cases, the cumulative
human--fleet coupled-trajectory evidence chain that underpins the
bonded-human-identity record retains its commitment structure, and
the relying subnet's acceptance policy declares which disclosure
regimes are admissible for which HC-channel classes. The human-layer
scope exclusions stated in the bonded-human-identity extension
remain in force; the deniability primitives do not introduce
admissibility paths for any of the channels excluded there. These
uses do not add any new HC-channel class, participant-probing
authority, dream-state pathway, Mandela-probe pathway, or
applicant-directed verification path beyond those expressly
admitted in the bonded-human-identity extension.


\paragraph{Out-of-scope note (structural, not policy default).}
The primitives of this section are commitment-infrastructure
extensions to the present filing. As a structural property:
(i) they do not replace, reduce, or substitute for
proof-of-projection, transitive proof-of-projection,
proof-of-discrepancy quorum confirmation
(Section~\ref{sec:pod-protocol}), or provenance negation under the
Anti-\TB lens;
(ii) they do not grant the trapdoor authority of
Section~\ref{sec:chameleon-hash-evidentiary}, or the designated
verifier of Section~\ref{sec:designated-verifier-liveness}, any
positive-verification role under the \TB regime;
(iii) they govern disclosure regimes for committed evidence and
not training objectives, governance decisions, admission standards,
or behavioural evaluation of agents, each of which is addressed in
the related Filing 2 application on empirical ethics training architecture;
(iv) all security, binding, hiding, and deniability claims stated
in this section remain empirical under declared attacker families,
adversary views, and resource budgets, time-indexed rather than
unconditionally durable, and do not constitute unconditional
cryptographic guarantees unless a specific embodiment explicitly
specifies and parameterises such a guarantee.




\paragraph{HC-Actuate (non-limiting).}
In some embodiments, the \RK issues a prompt, request, or
operator-facing instruction to a human participant through a committed
emission channel and captures the participant's response through the
same detector chain used for scene observation or through a declared
auxiliary tamper-evident channel committed under the same protocol
digest, timebase, and meter policy. In some embodiments, the request
is rendered through a declared emission channel, including any
sensory channel already available for HC-Read or HC-Entrain
operation; this paragraph concerns the commitment of the outbound
request, rendering parameters, and response-capture window rather
than redefining those human-coupled channels. The request-generating
process may be the RK controller or a logged computational partner,
including without limitation a PolieBot or PoliePuter-coupled
process as described in Definition~\ref{def:poliebot},
Definition~\ref{def:poliputer}, and
Section~\ref{sec:agents}; higher-layer multi-agent coordination
remains optional and non-essential to the present paragraph.

In some embodiments, a non-limiting worked flow comprises: (i)~the
request-generating process selects a request payload conditioned on
committed observation $o_{0:t}$, meter vector $\mathbf{m}_{0:t}$,
and, where applicable, a tool-mediated derivation trace; (ii)~the
payload is committed to the protocol digest prior to rendering,
together with the selected rendering-channel parameters and the
expected-response window; (iii)~the payload is rendered through the
committed emission path, and where the rendering uses audio
synthesis, the synthesis voice identity, prosody parameters, and
synthesis model fingerprint, provider/model identifier, or version
hash where available are committed alongside the payload digest;
(iv)~the response may enter through any HC-Write input channel,
including voice, gesture, locomotion, explicit control input,
tool-mediated input, or physiological measurement, and is logged
under the channel identity and detector configuration used for that
response; (v)~request digest, rendering parameters, response atoms,
inter-event timing, and any coupled meter readings are bound into a
single $C_{0:T}$ window under the protocol digest.

In some embodiments, the request--response binding provides a
closed-loop audit object in which the request, rendering parameters,
response atoms, channel tags, and inter-event timing are jointly
checked against the protocol digest and disclosed atoms, so that
post-hoc substitution of a request, response, rendering parameter,
or timing value is detectable as a commitment,
channel-consistency, or meter-consistency failure.

In some embodiments, HC-Actuate composes with channel-separation
primitives of Section~\ref{sec:channel-separation}, with the request
assigned a declared channel tag $u_{\mathrm{ch}}$ and, where
applicable, with the cross-scene regime combinations of
Section~\ref{sec:cross-scene-regime-combinations} in which a
concurrent \TB, \LI, \RT, or cryptographic
companion stream corroborates the request--response event.

In some embodiments, failure to render the request within the
declared emission envelope, failure to capture a response within the
expected response window, detector saturation, channel occlusion,
model-fingerprint mismatch, or response-channel dropout is logged
as an HC-Actuate anomaly and causes the request--response episode
to be treated as incomplete rather than silently accepted.

In some embodiments, where the response is captured through an
auxiliary channel rather than the primary scene-observation detector
chain, the episode is marked as a reduced-assurance HC-Actuate
sub-embodiment unless the auxiliary channel has an independently
declared calibration state, channel tag, timing model, and meter
envelope.

In some embodiments, HC-Actuate episodes are subject to declared
participant-safety, exposure, privacy, and abort envelopes, and an
abort event disables or attenuates the emission path, commits
already-recorded atoms, and marks the episode as aborted in the
protocol digest.

In some embodiments, meters for HC-Actuate are partitioned into
acceptance meters that gate whether the request--response episode
may be relied upon, evaluation meters that are held out for later
audit, and optional shaping meters used only to adapt future
request rendering and not to validate the present response.


\paragraph{Multi-contributor freshness composition (non-limiting).}
\label{par:multi-contributor-freshness}
As a non-limiting extension of the two-seed structure of
Section~\ref{sec:two-seed-protocol} and
Paragraph~\ref{par:two-seed}, in some embodiments a freshness value
used as, or used to derive, $s_{\mathrm{run}}$ or
$s_{\mathrm{open}}$ of Definition~\ref{def:runseed} is composed from two or more
declared contributors via a commit--reveal protocol. In some
embodiments, each contributor publishes a commitment
$K_i = \mathsf{Commit}(r_i;\,\rho_i)$ before the capture window
opens, the commitments are aggregated into the protocol digest
alongside declared contributor identities and role tags (for
example, operator, verifier, participating agent, or provider
class), and the capture window executes.

Where the composed value contributes to $s_{\mathrm{run}}$,
contributor reveal to the seed-computing controller, quorum, or
escrowed derivation process occurs after all commitments are bound
into the protocol digest but before the first consumption of the
run seed by the scan law or emission schedule; disclosure of the
revealed values to other verifying parties may remain withheld
until selective disclosure under the declared opening policy.
Where the composed value contributes to $s_{\mathrm{open}}$,
contributor reveal occurs after publication of the run commitment
or capture-window close anchor, under the declared opening
schedule.

In some embodiments, the composed freshness value is obtained via a
declared composition rule---for example
$s = H(r_1 \Vert r_2 \Vert \cdots \Vert r_k)$ or a bitwise XOR over
fixed-width representations---with the composition rule itself
committed into the protocol digest.

In some embodiments, at least one contributor is a computational
process or agentive partner operating under PolieBot or PoliePuter
coupling; for purposes of this paragraph, the disclosed role of
that partner is limited to contributing or committing a freshness
input, not to broader multi-agent coordination.

In some embodiments, verifier treatment of the composed freshness
value is conditioned on the following declared checks:
(a)~all declared contributor commitments appear in the protocol
digest before the capture-window open anchor;
(b)~freshness assurance under this sub-embodiment is conditioned
on a declared entropy and independence model, and is reduced or
invalidated when the disclosed record shows that the composed value
was determined entirely by one contributor's pre-capture state, by
a single common-control source, or by any source excluded by the
declared freshness policy; and
(c)~the composition rule is recomputed from disclosed atoms during
verification.

In some embodiments, if a declared contributor fails to reveal
within the committed reveal schedule, reveals a value inconsistent
with its commitment, or reveals through a channel failing the
timing-anchor policy, the composed value is marked invalid or
downgraded according to a predeclared missing-contributor rule. In
some embodiments, the missing-contributor rule is selected from
abort, reduced-assurance continuation, threshold composition from a
declared subset of contributors, or recomputation using a
predeclared fallback freshness source, with the selected rule
committed before the capture window opens.

In some embodiments, contributor identities, role tags, commitment
values, reveal schedule, composition rule, canonical serialisation
format, and any truncation or domain-separation tags are committed
in the protocol digest before the corresponding seed is consumed.

In some embodiments, local operator-only contribution,
single-contributor contribution, contribution by a single
computational provider acting as a declared contributor, or
contribution by entities under common control is logged as a
reduced-assurance freshness sub-embodiment rather than as
independent multi-contributor freshness.

In some embodiments, freshness-related meters are partitioned so
that minimum freshness checks gate acceptance, while
contributor-diversity, common-control, and entropy-budget
indicators may be retained as evaluation meters or assurance
annotations rather than used by the operating controller as
shaping signals.


\subsubsection{Logged derivation of session parameters}
\label{sec:agent-derived-session-params}

In some embodiments, session parameters consumed by the controller
of a physical \RK run---including scan-law randomisation
inputs, challenge-selection indices, emission-schedule seeds, and
human-readable presentations of derived seed values, including
wordlist encodings---are computed by one or more participating
agents or controllers through logged tool calls. For the subset of
agent or controller tool calls that derive session parameters for
a physical RK run, the request, response, derivation inputs, and
post-processing record are logged as a declared auxiliary channel
of the \cb rather than treated as an opaque external input.

In some embodiments, the derivation is consumed by, constraining,
or verifying a physical RK run or evidence bundle, with the
parameter-consumption event committed in the protocol digest
$\Pi_{\mathrm{dig}}$ before the capture window opens; the auxiliary
channel binds the tool-call record to a specific physical RK run
through the run-identifier $\mathrm{rid}$ and the timestamp of the
parameter-consumption event.

In a non-limiting worked derivation example, the canonicalisation
rule normalises the tool response to a UTF-8-encoded JSON string
with sorted keys, and the derivation rule is HKDF over the
canonical response with $\Pi_{\mathrm{dig}}$ as salt and a declared
domain-separation tag as info, with the HKDF output parsed as a
32-bit unsigned integer scan-seed, an 8-tuple of exposure-schedule
parameters in declared units, and a 16-bit verifier-cadence value;
other canonicalisation and derivation rules are non-limiting
alternatives.

In some embodiments, each such tool call is specified in advance
by a precommitted tool-call specification
$\mathsf{tc} = (\texttt{name},\,\texttt{args-schema},\,
\texttt{provider-policy},\,\texttt{canonicalisation-rule},\,
\texttt{derivation-rule})$
whose hash is bound into the protocol digest before the capture
window opens, such that substitution of a different tool call after
the fact is detectable on selective opening.
In some embodiments, the provider policy declares a provider
identity, provider class, permitted provider set, local
deterministic tool, or corroboration policy, and the canonicalisation
rule declares how tool responses are normalised before derivation.

In some embodiments, the derivation input to the tool call
comprises one or more of: the protocol digest of a parent run, a
concurrent run whose referenced windows have closed before the
parameter-consumption event, a hash-chain anchor exposed by a
previously committed window, an external blockhash or beacon
value, or a declared composition of the above with
multi-contributor commit--reveal outputs per
Paragraph~\ref{par:multi-contributor-freshness}. In some
embodiments, the derivation trace---inputs, tool identifier, model
version or provider fingerprint, the response or a canonical
response representation sufficient under the declared derivation
rule, together with any response hash, selective opening policy,
and post-processing steps applied before the derived value enters
the protocol---is committed as an auxiliary channel of the \cb
such that a verifier can recompute the derivation from disclosed
atoms.

In some embodiments, the same-session case comprises a
configuration in which an agent participating in a capture session
derives a later session parameter from the chain state of one or
more already-closed prefix windows of that session, and the
derivation trace is committed before the derived parameter is
consumed. In some embodiments, this configuration produces a
bounded same-session consistency loop wherein (a)~session
parameters are verified against the chain state, (b)~chain state is
verified against the committed atoms, and (c)~committed atoms are
verified against the declared session parameters, such that the
three consistency checks are evaluated against the same
protocol-digest, hash-chain, and selective-opening commitment
family rather than against an uncommitted off-chain parameter
source.

In some embodiments, same-session derivation is constrained by a
declared causality envelope including both a time-order constraint
and a consumption constraint: a chain-state value may be read only
from a window whose close anchor precedes the first controller
event that consumes the derived parameter, and this constraint is
enforced by the verifier during selective opening. In some
embodiments, violation of the causality envelope---for example, a
tool-call record whose committed input references a window
timestamped after its consumption event---is surfaced as a
protocol-digest anomaly regardless of the correctness of the
derived value.

Where the derived value is obtained through a semantic or routed
ensemble branch, the semantic-computation attestation record of
Section~\ref{sec:semantic-branch} may be used as the tool-output
attestation object; this subsubsection addresses derivation of
session parameters rather than semantic correctness.

In some embodiments, the tool-call specification further declares
a canonical serialisation format for inputs and outputs, a
domain-separation tag for each derived parameter type, and a
deterministic post-processing rule, so that independent verifiers
recompute the same derived value from disclosed atoms.

In some embodiments, if the tool provider returns an error,
timeout, noncanonical response, response outside the declared
schema, or response whose provider fingerprint does not match the
committed tool-call specification, the derived parameter is
rejected or the run is marked as reduced-assurance under a
predeclared fallback rule.

In some embodiments, nondeterministic tool calls are admitted only
when their nondeterminism is itself declared and committed, for
example by committing provider sampling parameters, provider-side
randomness commitments, or a canonical response-selection rule;
otherwise the nondeterministic output is treated as
non-recomputable and cannot serve as an acceptance-critical
derivation input.

In some embodiments, derivation checks that gate acceptance of the
run are partitioned from evaluation checks that assess provider
reliability, tool-call latency, model drift, or common-mode
failures, so that the participating agent cannot optimise directly
against every derivation-quality surface.


\paragraph{Heterogeneous-provider corroboration of canonical derived
values (non-limiting).}
In addition to multi-RK corroboration of physical evidence, in some
embodiments corroboration of canonical derived values produced
through logged tool-call specifications---including the seed
derivations of Section~\ref{sec:agent-derived-session-params},
tool-call outputs contributing to protocol randomisation, and
commitment values advanced as freshness contributions---is obtained
by executing the same committed tool-call specification on two or
more provider instances having declared diversity or
non-common-control properties, such as distinct organisational
operators, distinct model families, distinct serving clusters, or
separately attested execution environments, whose identities, model
versions, and serving infrastructure fingerprints are committed to
the protocol digest alongside each response.

Provider corroboration verifies agreement on a canonical derivation
under a committed tool-call specification; it does not by itself
add fresh entropy unless a provider is separately declared as a
contributor under Paragraph~\ref{par:multi-contributor-freshness}.

In some embodiments, agreement among provider responses on the
canonical derived value is a non-entropy-bearing
derivation-corroboration meter or freshness-assurance submeter,
while entropy contribution and contributor diversity are accounted
for separately under the multi-contributor freshness composition
rule.

In some embodiments, disagreement is logged as a protocol-digest
anomaly and may be surfaced through a discrepancy-reporting or
audit-escalation channel; the disagreement is not treated as a
proof-of-discrepancy unless the applicable PoD requirements of
Section~\ref{sec:pod-protocol}, including physical re-execution
where required, are separately satisfied.

In some embodiments, provider diversity and contributor diversity
are recorded as distinct fields: provider diversity corroborates
execution of a committed derivation rule, while contributor
diversity is evaluated under the freshness-composition rule of
Paragraph~\ref{par:multi-contributor-freshness}.

In some embodiments, the committed tool-call specification declares
canonicalisation rules for provider responses, including
whitespace, encoding, numerical precision, truncation,
domain-separation tags, and schema validation, so that provider
agreement is assessed on the same canonical value.

In some embodiments, if provider responses disagree, fail to
arrive within the declared response window, or fail schema
validation, the derived value is not silently majority-resolved;
instead the run follows a predeclared rule selected from abort,
reduced-assurance continuation, manual audit, deterministic local
recomputation, or re-execution under a fresh provider set.

In some embodiments, provider-diversity claims are downgraded when
providers share a common model lineage, serving infrastructure,
operator, training corpus, tool backend, or upstream deterministic
implementation likely to create a common-mode failure.

In some embodiments, provider-corroboration meters are partitioned
from acceptance meters so that a provider-comparison mechanism can
be used for held-out audit or reliability scoring without becoming
an optimisation target for the participating agent.

In some embodiments, a single-provider execution of the committed
tool-call specification remains a reduced-assurance sub-embodiment
where the tool trace, canonical response, and post-processing rule
are still committed and recomputable.


\section{Industrial Applicability}
% ======================================================================

The disclosed systems and methods are applicable to at least the
following fields: physical evidence and provenance verification in media
production, journalism, and legal proceedings; active sensing and
computational imaging for metrology, inspection, and autonomous
navigation; projection mapping and immersive display for entertainment,
architecture, and education; hardware security and device authentication
using physically unclonable properties; continuous tamper detection for
critical infrastructure monitoring, evidence custody, and nuclear
safeguards; liveness detection and anti-spoofing for biometric systems;
networked sensing and distributed verification for environmental
monitoring, smart cities, and defence; physical computing and
reservoir computing using optical or hybrid media; and distributed
model improvement and fleet-calibrated device management through
physically verified discrepancy discovery.

% ======================================================================
\section*{Non-limiting Statement}
\label{sec:non-limiting-statement}
% ======================================================================

The embodiments and examples described herein---including all main-body
sections, appendices, and optional modules---are non-limiting. The
scope of protection, once claims are filed, is defined solely by the
claims. No feature, embodiment, example, numerical value, operating
range, material choice, protocol, algorithm, or theoretical
correspondence described in this specification is intended to limit
the scope of the claims unless expressly recited therein. Where the
specification describes a feature as present ``in some embodiments,''
the absence of that feature in other embodiments is equally
contemplated. The terms ``preferred,'' ``advantageous,'' and
``particularly suitable'' indicate non-limiting design preferences
and do not exclude alternative implementations. Dimensions,
quantities, frequencies, and other numerical parameters are
illustrative and may be varied without departing from the principles
described. Theoretical frameworks, substrate mappings, categorical
constructions presented herein are interpretive
tools and do not constrain hardware or software implementation.

% ======================================================================
\section{Declared Open Problems}
\label{sec:declared-open-problems}
\label{sec:appendix-research}

The following are open empirical questions, explicitly not asserted as
established results: survival of simultaneous tone multiplexing; the
effective-rank to task-useful-rank mapping; drift-bias tolerance of
physical gradient methods; and controllability-scaling calibration. These open questions
are not required for the Tier A or Tier B claims of
Section~\ref{sec:claim-tier-structure} and bear only on Tier C capacity
claims within the bounds of Section~\ref{sec:capacity-accounting-rule}.


% ======================================================================

This appendix collects research directions and unproven conjectures
that are intellectually connected to the \RK architecture
but are not claims of this disclosure.  No embodiment outside this
appendix depends on any result stated here.

\subsection{Yoked channel authentication capacity (conjecture)}

The following is an unproven conjecture stated as a research
direction.  Let $\alpha$ denote the coupling strength between scene
and reactor, and let $\mathcal{I}_F(\alpha)$ denote the Fisher
information about $\alpha$ in a bundle $C_{0:T}$:
\[
  \mathcal{I}_F(\alpha)
  = \mathbb{E}\!\left[\left(\frac{\partial}{\partial \alpha}
      \log p(C_{0:T} \mid \alpha)\right)^2\right].
\]
The operating hypothesis is that $\mathcal{I}_F(\alpha) \to \infty$
as $\CLE(\alpha) \to 0^-$, by analogy with susceptibility divergence
at a phase transition (fluctuation--response theorem of statistical
mechanics).  If this divergence holds, then any emulator attempting
to reproduce the coupled dynamics with coupling parameter
$\tilde\alpha \neq \alpha$ incurs an expected log-likelihood gap
that grows without bound as the genuine system approaches the
synchronisation boundary.

This motivates the \emph{Yoked Channel Authentication Capacity
Conjecture}: a Yoked channel operated at
$\CLE \in [-\varepsilon_{\CLE}, 0)$ achieves an authentication
capacity $C_{\mathrm{Yoked}}(\alpha)$ (in the sense of bits of
identity information attributable to the joint attractor rather than
to any product-state emulator of equal physical scale) that is
strictly positive and increases as $\CLE \to 0^-$.  Whether
$C_{\mathrm{Yoked}}$ can be computed in closed form for a
representative class of Yoked systems, and whether it can be
lower-bounded via the mutual information $I(S_{0:T}; R_{0:T})$ in
Yoked mode relative to the maximum mutual information achievable by a
decomposable system of equal physical scale, are open questions.  A
positive result would establish Yoked coupling as a channel primitive
in the information-theoretic sense.

All operational claims in this specification regarding Yoked
operation are conditioned on the empirically measured TE and CLE
proxy suite under declared meter envelopes and are independent of
the truth of this conjecture.

Even if no closed-form Yoked authentication capacity is established, Yoked
operation remains supported by the declared transfer-entropy, conditional
Lyapunov, coupling-quality, and tamper-detection criteria recorded in the
protocol digest.


\subsection{Classifier-mediated Yoked state estimation: cross-scene
generalisation}

The worked embodiment of Section~\ref{sec:yoked-classifier} trains a
coupling state classifier on known analysable reactor pairs and
deploys it against non-analysable scenes.  Several open questions
must be resolved by experiment before this architecture can be relied
upon for deployment against novel scene classes:

Whether device-side bundle data identifies two-sided CLE and TE
across unseen scene classes; how generalisation error scales with
reactor diversity rather than raw reactor count $K$; which reactor
families are needed to cover human-relevant dynamics; whether the
declared supervision window length $T$ is sufficient for CLE
convergence in each declared scene class; what the true run-level
effective sample size is after within-interaction autocorrelation is
accounted for; what end-to-end latency is achievable under on-device
compute constraints; and whether a concrete agentive-scene control
law converges stably or requires empirical bounding.

The minimal adequate experiment for the most important of these ---
cross-scene classifier generalisation --- is as follows.  Train the
classifier on $K \geq 50$ known analysable reactors spanning at
least five distinct scene classes (non-limiting examples: phosphor
screen, scattering plate, resonant mechanical structure, living plant
tissue, artificial vascular phantom).  For each reactor, compute
ground-truth CLE and TE from both device-side and scene-side signals
simultaneously; use device-side bundle data alone as classifier input.
Evaluate on $M \geq 20$ held-out reactors from at least two scene
classes entirely absent from training.  A positive result requires:
area under ROC $> 0.85$ for binary Yoked / not-Yoked classification
on held-out scene classes; Pearson $r > 0.7$ between predicted and
ground-truth CLE on a continuous scale; and latency within the
declared rolling window period.  A negative result requiring
restriction of the embodiment to declared scene families: AUC below
$0.6$ on held-out classes while remaining high within training
classes (indicating the classifier is recognising scene-class
signatures rather than coupling state), or systematic CLE sign-flip
error rate above $0.2$ near the Yoked boundary.

These are experimental questions whose resolution would strengthen
the embodiment but whose absence does not affect any other claim in
this disclosure.

\end{document}
